\documentclass[twoside,12pt,english,openany]{book}

\usepackage{setspace}
% \usepackage{natbib}
\linespread{1.2}

\PassOptionsToPackage{table}{xcolor} % enables \rowcolors for the zebra-striped tables below
%!TEX root = ../main.tex

\newenvironment{abstract}{}{}
\usepackage{abstract}
% \usepackage[T1]{fontenc} % 疑似修改点
% \usepackage{newpxtext,newpxmath} % 疑似修改点
\usepackage{arxiv}
\usepackage[square,numbers]{natbib}






\usepackage[utf8]{inputenc} % allow utf-8 input
\usepackage[T1]{fontenc}    % use 8-bit T1 fonts
% \usepackage{tgpagella}
% \usepackage{mathpazo}
% \usepackage{tgtermes}
% \usepackage{newtxmath}
\usepackage[english]{babel}
\usepackage{dsfont} %
\usepackage{xcolor}

\usepackage{hyperref}
\definecolor{linkgreen}{HTML}{009000}
\definecolor{linkcobalt}{HTML}{3646F0}
\hypersetup{
    colorlinks=true,
    linkcolor=black,
    citecolor=linkgreen,
    filecolor=linkcobalt,
    urlcolor=linkcobalt,
    pdfpagemode=FullScreen,
    }    % hyperlinks
\usepackage{url}            % simple URL typesetting
\usepackage{booktabs}       % professional-quality tables
\usepackage{amsfonts}       % blackboard math symbols
\usepackage{nicefrac}       % compact symbols for 1/2, etc.
\usepackage{microtype}      % microtypography
\usepackage{lipsum}		% Can be removed after putting your text content
\usepackage{graphicx}

\usepackage{doi}
\usepackage{amsmath}
\usepackage{gensymb}
\usepackage{enumitem}
\DeclareMathOperator*{\argmax}{arg\,max}
\DeclareMathOperator*{\argmin}{arg\,min}
\usepackage{pifont}
\usepackage{xparse}
\usepackage{xspace}

\newcommand{\vs}{\textit{vs.}\xspace}
\newcommand{\ie}{\textit{i.e.,}\xspace}
\newcommand{\eg}{\textit{e.g.,}\xspace}
\usepackage{wrapfig}
\usepackage{graphicx}


%\usepackage[landscape,a4paper]{geometry}
\usepackage{colortbl}
\newcommand{\ra}[1]{\renewcommand{\arraystretch}{#1}}
\usepackage{longtable}

\usepackage{titlesec}
% 四级标题
\titleformat{\paragraph}[runin]{\normalsize\bfseries}{\theparagraph}{1em}{}
\titlespacing*{\paragraph}{0pt}{6pt}{4pt}
% 五级标题
\titleformat{\subparagraph}[runin]{\normalsize\itshape}{\thesubparagraph}{1em}{}
\titlespacing*{\subparagraph}{0pt}{4pt}{3pt}
\renewcommand\theparagraph{\thesubsubsection.\arabic{paragraph}}
\renewcommand\thesubparagraph{\theparagraph.\arabic{subparagraph}}
\setcounter{secnumdepth}{4}
\setcounter{secnumdepth}{5} % 允许显示五级标题的序号

%\usepackage[numbers]{natbib}
% \usepackage[margin=1in]{geometry}
\usepackage{tikz}
\usetikzlibrary{trees, tikzmark, shadows, arrows.meta, positioning, calc}
%\usepackage{forest}

% draw the tree
\usepackage[edges]{forest}
\usepackage[normalem]{ulem}
% For timeline
\usepackage[most]{tcolorbox}

% 158 195 231
\definecolor{ai_blue}{RGB}{158, 195, 231}

\definecolor{forestgreen}{RGB}{34,139,34}
\definecolor{steelred}{RGB}{139,30,30}
\newtcolorbox{quotebox}{
    colback=ai_blue!60, 
    colframe=gray!30,
    coltext=black,
    boxrule=0pt,
    arc=4pt,
    left=10pt,
    right=10pt,
    top=10pt,
    bottom=10pt,
    enhanced,
    drop shadow={black!50!white}
}

\usepackage{multicol}
\usepackage{multirow}
\usepackage{cleveref}

\tikzset{%
    parent/.style =          {align=center,text width=2cm,rounded corners=3pt, line width=0.3mm},%, fill=gray!10,draw=gray!80},
    child/.style =           {align=center,text width=2.3cm,rounded corners=3pt, fill=blue!10,draw=blue!80,line width=0.3mm},
    grandchild/.style =      {align=center,text width=2cm,rounded corners=3pt},
    greatgrandchild/.style = {align=center,text width=1.5cm,rounded corners=3pt},
    greatgrandchild2/.style = {align=center,text width=1.5cm,rounded corners=3pt},    
    referenceblock/.style =  {align=center,text width=1.5cm,rounded corners=2pt},
    %%%% Re-define
    % Pretrain Model
    pretrain/.style =           {align=center,text width=2.9cm,rounded corners=3pt, fill=blue!10,draw=blue!80,line width=0.3mm},   
    pretrain_work/.style =           {align=center, text width=2.9cm,rounded corners=3pt, fill=blue!10,draw=blue!0,line width=0.3mm},  
    % Template Mining
    template/.style =           {align=center,text width=2.9cm,rounded corners=3pt, fill=red!10,draw=red!80,line width=0.3mm},   
    template_work/.style =           {align=center,text width=2.9cm,rounded corners=3pt, fill=red!10,draw=red!0,line width=0.3mm},    
    % Answering Mining
    answer/.style =           {align=center,text width=2.9cm,rounded corners=3pt,line width=0.3mm},%, fill= cyan!10,draw= cyan!80,},   
    answer_work/.style =           {align=center,text width=2.9cm,rounded corners=3pt,line width=0.3mm},%, fill= cyan!10,draw= cyan!0,},      
    % Multi-prompt 
    multiple/.style =           {align=center,text width=2.9cm,rounded corners=3pt, fill= orange!10,draw= orange!80,line width=0.3mm},   
    multiple_work/.style =           {align=center,text width=2.9cm,rounded corners=3pt, fill= orange!10,draw= orange!0,line width=0.3mm},        
    % Tuning Strategy
    tuning/.style =           {align=center,text width=2.9cm,rounded corners=3pt, fill= magenta!10,draw= magenta!80,line width=0.3mm},   
    tuning_work/.style =           {align=center,text width=2.9cm,rounded corners=3pt, fill= magenta!10,draw= magenta!0,line width=0.3mm},
    % Assurance
    assurance/.style =           {align=center,text width=2.9cm,rounded corners=3pt, fill= green!10,draw= green!80,line width=0.3mm},   
    assurance/.style =           {align=center,text width=2.9cm,rounded corners=3pt, fill= green!10,draw= green!0,line width=0.3mm},
}


\usepackage{color}			% Color extended names

%DEFINE ENVIRONMENT BLOCK 
% Riddle
\newenvironment{Riddle}{		% Riddle environment
  \begin{center}			% Centering minipage
    \colorbox[HTML]{F8E0E0} {		% Set's the color of minipage
      \begin{minipage}[b]{380px} 	% Starts minipage
 	\textbf{Riddle: }\\ \itshape 	% Set's title and starts italic for text 
      \end{minipage}}			% End minipage
}{\end{center}}				% End Riddle environment


\usepackage{tabularx}
\usepackage{tabularray}
\newcolumntype{C}[1]{>{\centering\arraybackslash}m{#1}}  % centered, auto-wrapping fixed-width column (fraction of \textwidth)
\newcolumntype{L}[1]{>{\raggedright\arraybackslash}m{#1}}  % left-aligned, auto-wrapping fixed-width column
\usepackage{CJKutf8}  % for Chinese text


\usepackage{amssymb}
\usepackage{amsthm}
\usepackage{mathtools}


\tcbuselibrary{theorems,breakable}

% First, adjust the theorem style to match proof spacing
\newtheoremstyle{definition}
  {6pt}        % Space above - matching \topsep6pt from proof
  {6pt}        % Space below - matching \topsep6pt from proof  
  {\normalfont} % Body font
  {}           % Indent amount
  {\bfseries}  % Theorem head font
  {.}          % Punctuation after theorem head
  {.5em}       % Space after theorem head (similar to \labelsep)
  {}           % Theorem head spec

\theoremstyle{definition}
\newtheorem{definition}{Definition}

% Then adjust tcolorbox to preserve these spacings
\definecolor{DefGreen}{HTML}{E8F5E9}
\tcolorboxenvironment{definition}{
  colback=DefGreen, %blue!5!white,
  boxrule=0pt,
  boxsep=0pt,
  % Match internal spacing to typical proof environment
  left=2pt,    % Use same label separation as proof
  right=\labelsep,   % Symmetric padding
  top=12pt,           % Small top padding inside box
  bottom=12pt,        % Small bottom padding inside box
  oversize=2pt,
  sharp corners,
  % Match external spacing to \topsep
  before skip=6pt,   % Same as proof \topsep6pt
  after skip=6pt,    % Same as proof \topsep6pt
  breakable,
}

% Alternative color schemes you might prefer for definitions:
% Option 1: Light gray (more neutral)
% colback=gray!5!white,

% Option 2: Light green (suggesting growth/construction)
% colback=green!5!white,

% Option 3: Light yellow (highlighting important concepts)
% colback=yellow!5!white,




\newcommand{\red}[1]{\textcolor{red}{#1}}
\newcommand{\blue}[1]{\textcolor{red}{#1}}
\newcommand{\bang}[1]{\textcolor{red}{#1}}
\newcommand{\mc}[1]{\textcolor{red}{MC#1}}
\newcommand{\hs}[1]{\textcolor{red}{Huan:#1}}
\newcommand{\kunlun}[1]{{{\textcolor{red}{[Kunlun: #1]}}}}
\newcommand{\xlqi}[1]{{{\textcolor{red}{[Xiaoliang: #1]}}}}
\newcommand{\xf}[1]{{{\textcolor{red}{[Xinfeng: #1]}}}}
\newcommand{\hsr}[1]{{{\textcolor{red}{[Sirui: #1]}}}}
\newcommand{\hbc}[1]{{\textcolor{red}{[haibo: #1]}}}
\newcommand{\hb}[1]{{\textcolor{red}{ #1}}}
\newcommand{\tjhe}[1]{{{\textcolor{red}{[Tanjin: #1]}}}}
\newcommand{\yg}[1]{\textcolor{red}{Yu Gu:#1}}
\newcommand{\jy}[1]{\textcolor{red}{Jiayi:#1}}

\newcommand{\yuyu}[1]{\textcolor{red}{Yuyu:#1}}


\definecolor{customgreen}{RGB}{191,255,191} % BFFFBF in RGB
\definecolor{customblue}{RGB}{228,254,254} % E4FEFE in RGB
\definecolor{customyellow}{RGB}{255,255,191} % FFFFBF in RGB
\definecolor{customred}{RGB}{255,191,191} % FFBFBF in RGB

\definecolor{sealblue}{HTML}{173F67}
\definecolor{sealgold}{HTML}{B08B3E}

% =============
% \DeclareRobustCommand{\FramedGlyph}[1]{%
%   \tikz[baseline=(Letter.base)]{%
%     \node[
%       rounded corners=1.5pt,
%       fill=sealblue,
%       draw=sealgold,
%       double=white,
%       double distance=1.2pt,
%       line width=0.7pt,
%       minimum width=1.02cm,
%       minimum height=1.12cm,
%       inner sep=0pt
%     ] (Frame) {};%    
%     \node[
%       font=\fontsize{29}{29}\selectfont\bfseries,
%       text=white,
%       inner sep=0pt,
%       yshift=-4pt
%     ] (Letter) at (Frame.center) {#1};%
%   }%
% }

% \newcommand{\FramedInitial}[2]{%
%   \lettrine[
%     lines=3,
%     lraise=0.02,
%     lhang=0.00,
%     loversize=0,
%     findent=0.35em,
%     nindent=0.60em,
%     refstring=T
%   ]{\FramedGlyph{#1}}{#2}%
% }

\DeclareRobustCommand{\FramedGlyph}[1]{%
  \raisebox{0.10\baselineskip}{%
    \tikz[baseline=(Letter.base)]{%
      % 固定整个图形的边界，避免 lettrine 误判高度
      \path[use as bounding box] (0,0) rectangle (1.20,1.34);

      % 外框
      \node[
        rounded corners=2pt,
        fill=sealblue,
        draw=sealgold,
        line width=0.9pt,
        minimum width=1.02cm,
        minimum height=1.12cm,
        inner sep=0pt
      ] (Frame) at (0.60,0.67) {};

      % 内部细边框
      \draw[
        white,
        line width=0.45pt
      ]
      ([xshift=2.2pt,yshift=-2.2pt]Frame.north west)
      rectangle
      ([xshift=-2.2pt,yshift=2.2pt]Frame.south east);

      % 字母：向上微调，解决“视觉上偏下”
      \node[
          % font=\initfamily\fontsize{29}{29}\selectfont,
          font=\fontsize{29}{29}\selectfont\bfseries,
        text=white,
        inner sep=0pt,
        yshift=1.6pt
      ] (Letter) at (Frame.center) {#1};
    }%
  }%
}
\newcommand{\FramedInitial}[2]{%
  \lettrine[
    lines=3,
    lraise=0.02,
    lhang=0.00,
    loversize=0,
    findent=0.35em,
    nindent=0.60em,
    refstring=T
  ]{\FramedGlyph{#1}}{#2}%
}


%==== for custom style
% Color Definitions
\usepackage{xcolor}
\definecolor{marron}{RGB}{60,30,10}
\definecolor{darkblue}{RGB}{0,0,80}
\definecolor{lightblue}{RGB}{80,80,80}
\definecolor{darkgreen}{RGB}{0,80,0}
\definecolor{darkgray}{RGB}{0,80,0}
\definecolor{darkred}{RGB}{80,0,0}
\definecolor{shadecolor}{rgb}{0.97,0.97,0.97}
\definecolor{darkblue}{RGB}{37,90,192}
% a very pale mint (≈ #F0FFF4)
\definecolor{LightMint}{HTML}{E8F5E9}
\definecolor{LightCobalt}{HTML}{E9EBFB} % stripe colour of the Notation and Glossary tables
% a darker, richer green for contrast (you can pick any hex)
\definecolor{DarkMint}{HTML}{E8F5E9}   % e.g. “dark green”

% % page background color
% \usepackage{xcolor}
% \usepackage{pagecolor}
%     \definecolor{MyCrimson}{HTML}{880000}
%     \definecolor{MyGrey}{HTML}{080808}
%     \definecolor{MyBackground}{HTML}{FFFFEB}
%     % \definecolor{MyBackground}{HTML}{F0FFF4}    
%     \pagecolor{MyBackground}


% Drop Caps & Special Fonts
\usepackage{lettrine}
\input Acorn.fd
\newcommand*\initfamily{\usefont{U}{Acorn}{xl}{n}}

\usepackage{todonotes}
\newcommand{\jp}[1]{\todo{\textbf{JP}: {#1}}}

% Fonts & Typography
\usepackage{color}            % (needed by many font/color macros)
% \renewcommand{\familydefault}{pplj}  % set Palatino-like text font
\usepackage[T1]{fontenc}
\usepackage{newpxtext}
% \usepackage[
%   final,
%   stretch=10,
%   protrusion=true,
%   tracking=true,
%   spacing=on,
%   kerning=on,
%   expansion=true
% ]{microtype}    % enables subtle typographic refinements (margin kerning, font expansion, etc.).
\setlength{\parskip}{1.3ex plus 0.2ex minus 0.2ex}  % control inter-paragraph spacing and indentation.
\setlength{\parindent}{1em}   % paragraph indentation (“Spanish” style)


% Headers & Footers (fancyhdr)
% \makeatletter
% \def\headrule{{\color{darkgray}\raisebox{-2.1pt}[10pt][10pt]{\leafright} \hrulefill \raisebox{-2.1pt}[10pt][10pt]{~~~\decofourleft \decotwo\decofourright~~~} \hrulefill \raisebox{-2.1pt}[10pt][10pt]{ \leafleft}}} % custom ornamented head rule
% \makeatother
\usepackage{fancyhdr}

\pagestyle{fancy}
% Clear header and footer lines
\fancyhf{} % clears headers and footers
\renewcommand{\headrulewidth}{0pt} % removes header line
\renewcommand{\footrulewidth}{0pt} % removes footer line
\fancyfoot[C]{\thepage}

% redefine "plain" style for chapter pages
\fancypagestyle{plain}{
  \fancyhf{}
  \fancyfoot[C]{\thepage}
}

%––– redefine \chaptermark and \sectionmark to only carry the raw title
\renewcommand{\chaptermark}[1]{\markboth{#1}{}}
\renewcommand{\sectionmark}[1]{\markright{#1}}
\pagestyle{fancy}
\fancyhf{}  
\newcommand{\estcab}[1]{\itshape\textcolor{marron}{\nouppercase #1}}
% even pages: page # and chapter title, both in \estcab
\fancyhead[LE]{\estcab{\thepage}}
\fancyhead[RE]{\estcab{\leftmark}}

% odd pages: section title and page #, both in \estcab
\fancyhead[LO]{\estcab{\rightmark}}
\fancyhead[RO]{\estcab{\thepage}}
% \pagestyle{fancy}   % (uncomment to activate)
% % fancy head
% \fancyhead[LE]{\estcab{Fran Oldstyle}}
% \fancyhead[RE]{\estcab{History of taxonomy}}
% % \fancyhead[CE,CO]{\estcab{\decoone}}
% \fancyhead[LO]{\estcab{\rightmark}} % malo cuando no hay section ~~~ \thesection
% \fancyhead[RO]{\estcab{\leftmark}}

% % \fancyhead[RO]{\bf\nouppercase{ \leftmark}}
% % \fancyfoot[LE]{\bf \thepage ~~ \leafNE}
% % \fancyfoot[RO]{ \leafNE  ~~ \bf \thepage}
% redefine the plain style (used on chapter-start pages) to mimic our fancy style
\fancypagestyle{plain}{%
  \fancyhf{}% clear all fields
  % even pages: page # left, chapter title right
  \fancyhead[LE]{\thepage}%
  \fancyhead[RE]{\leftmark}%
  % odd pages: section title left, page # right
  \fancyhead[LO]{\rightmark}%
  \fancyhead[RO]{\thepage}%
  % optionally re-draw or remove the headrule:
  \renewcommand{\headrulewidth}{0pt}% or 0pt to suppress
}


% % fancy foot
% \fancyfoot[LO]{
% \ornimpar \\ \large \hfill \sffamily\bf \textcolor{darkgray}{\leafNE ~~~ \thepage}
% }
% \fancyfoot[RE]{\ornpar   \\ \large  \sffamily\bf \textcolor{darkgray}{\thepage ~~~ \reflectbox{\leafNE}}  \hfill}

% \newenvironment{Section}[1]
% {\section{\vspace{0ex}#1}}
% {\vspace{12pt}\centering ------- \decofourleft\decofourright ------- \par}


% Decorative Ornaments
\usepackage{fourier-orns}
\newcommand{\ornamento}{\vspace{2em}\noindent \textcolor{darkgray}{\hrulefill~ \raisebox{-2.5pt}[10pt][10pt]{\leafright \decofourleft \decothreeleft  \aldineright \decotwo \floweroneleft \decoone   \floweroneright \decotwo \aldineleft\decothreeright \decofourright \leafleft} ~  \hrulefill \\ \vspace{2em}}}
\newcommand{\ornpar}{\noindent \textcolor{darkgray}{ \raisebox{-1.9pt}[10pt][10pt]{\leafright} \hrulefill \raisebox{-1.9pt}[10pt][10pt]{\leafright \decofourleft \decothreeleft  \aldineright \decotwo \floweroneleft \decoone}}}
\newcommand{\ornimpar}{\textcolor{darkgray}{\raisebox{-1.9pt}[10pt][10pt]{\decoone \floweroneright \decotwo \aldineleft \decothreeright \decofourright \leafleft} \hrulefill \raisebox{-1.9pt}[10pt][10pt]{\leafleft}}}


% Graphics & Floating Figures
\usepackage{graphicx}
\usepackage{wallpaper}
\usepackage{wrapfig}
\usepackage{booktabs}


\usepackage{array}
\usepackage{rotating}
\usepackage{ragged2e}
\usepackage{arydshln}

\RequirePackage{xspace}
\makeatletter
\DeclareRobustCommand\onedot{\futurelet\@let@token\@onedot}
\def\@onedot{\ifx\@let@token.\else.\null\fi\xspace}

\def\eg{\emph{e.g}\onedot} 
\def\ie{\emph{i.e}\onedot} 
\def\cf{\emph{cf}\onedot} 
\def\Cf{\emph{Cf}\onedot}
\def\etc{\emph{etc}\onedot} 
\def\vs{\emph{vs}\onedot}
\def\wrt{w.r.t\onedot} 
\def\dof{d.o.f\onedot}
\def\iid{i.i.d\onedot} 
\def\wolog{w.l.o.g\onedot}
\def\etal{\emph{et al}\onedot}
\makeatother


\definecolor{lightgrey}{HTML}{D3D8F4} % group-heading rows in tables (\rowcolor{lightgrey}): one step darker than the LightCobalt stripe; was RGB(201,203,209)
\definecolor{lightlightgrey}{RGB}{238,238,238}

% Zebra-striped tables: header rows stay plain; every \midrule restarts the
% stripes (white row first), and explicit \rowcolor rows keep their own color.
% Rows above the first \midrule get a negative row count, which clears any
% color each row (no stripes, no \rowcolor bleed into the header). \toprule
% re-arms this so tabularx's trial runs start clean. \cdashline draws its
% dashes as an extra pseudo-row, so the count is stepped back after it.
\colorlet{zebrastripe}{LightCobalt} % same stripe color as the Notation and Glossary tables (#E9EBFB); was RGB(234,241,250)
\makeatletter
\def\zebra@arm{\rowcolors{1}{}{zebrastripe}\global\rownum-10000\relax}
\AddToHook{env/tabular/begin}{\zebra@arm}
\AddToHook{env/tabular*/begin}{\zebra@arm}
\AddToHook{env/tabularx/begin}{\zebra@arm}
\AtBeginDocument{%
  \let\zebra@toprule\toprule
  \def\toprule{\noalign{\global\rownum-10000\relax}\showrowcolors\zebra@toprule}%
  \let\zebra@midrule\midrule
  \def\midrule{\noalign{\global\rownum\@ne}\showrowcolors\zebra@midrule}%
  \let\zebra@cdashline\cdashline
  \def\cdashline#1{\zebra@cdashline{#1}\noalign{\global\advance\rownum\m@ne}\showrowcolors}}
\makeatother


\definecolor{token_blue}{RGB}{84, 120, 140}
\def\blue#1{\textbf{\color{our_blue} #1}} %in Table
\def\red#1{\textbf{\color{our_red} #1}} %in Table
\def\green#1{\textbf{\color{our_green} #1}} %in Table
\def\orange#1{\textbf{\color{our_orange} #1}} %in Table
\def\tokenblue#1{\textbf{\color{token_blue} #1}} %in Table

% \newcommand{\xmark}{\ding{55}}
% \newcommand{\cmark}{\ding{51}}

\def\vx{{\bm{x}}}

\newlength\savewidth\newcommand\shline{\noalign{\global\savewidth\arrayrulewidth \global\arrayrulewidth 1pt}\hline\noalign{\global\arrayrulewidth\savewidth}}
\newcommand{\tablestyle}[2]{\setlength{\tabcolsep}{#1}\renewcommand{\arraystretch}{#2}\centering\footnotesize}

\newcolumntype{x}[1]{>{\centering\arraybackslash}p{#1pt}}
\newcolumntype{y}[1]{>{\raggedright\arraybackslash}p{#1pt}}
\newcolumntype{z}[1]{>{\raggedleft\arraybackslash}p{#1pt}}

\renewcommand{\paragraph}[1]{\vspace{1mm}\noindent\textbf{#1}}

\newcommand\hshline{\noalign{\global\savewidth\arrayrulewidth
  \global\arrayrulewidth 0.5pt}\hline\noalign{\global\arrayrulewidth\savewidth}}

\setlength{\abovecaptionskip}{1pt}

\renewcommand{\paragraph}[1]{\vspace{1.25mm}\noindent\textbf{#1}}

\definecolor{codeblue}{rgb}{0.25, 0.5, 0.5}
\definecolor{codekw}{rgb}{0.35, 0.35, 0.75}
\lstdefinestyle{Pytorch}{
    language = Python,
    backgroundcolor = \color{white},
    basicstyle = \fontsize{9pt}{8pt}\selectfont\ttfamily\bfseries,
    columns = fullflexible,
    aboveskip=1pt,
    belowskip=1pt,
    breaklines = true,
    captionpos = b,
    commentstyle = \color{codeblue},
    keywordstyle = \color{codekw},
}

\definecolor{green}{HTML}{009000}
\definecolor{red}{HTML}{ea4335}
\newcommand{\better}[1]{\textcolor{green}{$\uparrow\,$#1}}
\newcommand{\worse}[1]{\textcolor{red}{$\downarrow\,$#1}}
\newcommand{\betterinv}[1]{\textcolor{green}{$\downarrow\,$#1}}
\newcommand{\worseinv}[1]{\textcolor{red}{$\uparrow\,$#1}}
\newcommand{\eqcontrib}{\clubsuit}
\newcommand{\correspond}{\spadesuit}

\newcommand{\ours}{\textsc{VL-Cogito}\xspace}

\newtcolorbox{promptblock}{
    colback=gray!5,
    colframe=gray!15,
    boxrule=0.5pt,
    arc=3pt,
    left=12pt,
    right=12pt,
    top=8pt,
    bottom=8pt,
    boxsep=8pt,
    breakable
}






\usepackage{CJKutf8}

\newcommand{\zy}[1]{{\color{blue}Zhaoyu: \begin{CJK*}{UTF8}{gbsn}#1\end{CJK*}}}
\newcommand{\wsq}[1]{{\color{orange}Shengqiong: \begin{CJK*}{UTF8}{gbsn}#1\end{CJK*}}}

%==== end custom style




\newcommand{\RespondAuthor}[1]{{\color{red}(#1)}}

\pdfstringdefDisableCommands{%
  \def\RespondAuthor#1{(#1)}%
}

\newcommand{\scott}[1]{{\color{red}[Hao: #1]}}



%========== Add new package here
\usepackage{makecell}
\usepackage{stmaryrd}
% stmaryrd has no bold font; keep its upright symbols in bold math.
\SetSymbolFont{stmry}{bold}{U}{stmry}{m}{n}
% Use supported script sizes for the timeline's 6.8pt labels.
\DeclareMathSizes{6.8}{6.8}{5}{5}
\usepackage{siunitx}
\usepackage{subcaption} % 支持 subfigure
%========== Add new package here

\usepackage{pifont}
\usepackage{xcolor}
\definecolor{my_green}{RGB}{51,102,0}
\definecolor{my_red}{RGB}{204, 0, 0}
\renewcommand{\checkmark}{\textcolor{my_green}{\ding{51}}} % ✔
\newcommand{\crossmark}{\textcolor{my_red}{\ding{55}}} % ✘

\definecolor{lightgray}{RGB}{235,235,235}

\usepackage{stmaryrd}


% ===== Define abstract environment =====
% \newcommand{\prefacename}{Preface}
% \newenvironment{preface}{
%     \setlength{\headheight}{15pt}
%     \vspace*{\stretch{2}}
%     {\noindent \bfseries \Huge \prefacename}
%     \begin{center}
%         % \phantomsection \addcontentsline{toc}{chapter}{\prefacename} % enable this if you want to put the preface in the table of contents
%         \thispagestyle{empty}
%     \end{center}%
% }
% {\vspace*{\stretch{5}}}

% \usepackage[
% backend=biber,
% style=alphabetic,
% ]{biblatex}
% \addbibresource{references.bib} %Imports bibliography file



% \title{A Book on Spatial Intelligence}


\title{Spatial Intelligence\\ \fontsize{12pt}{14pt}\selectfont  Foundations, Architectures, Methodologies, Applications and Challenges}


% Handbook of 


% \title{Spatial Intelligence: Charting the Cognitive Geometry of Intelligent Systems}







% Add this for removing 'A Preprint' And 'Date'
\date{} 
\renewcommand{\headeright}{}
% Below the title rule, above the authors: the book's positioning and project page.
\renewcommand{\undertitle}{%
  {\normalfont\large\itshape A Handbook}\\[1.5mm]%
  {\normalfont\normalsize Explore the project:\enspace{\large\url{https://spatial-agi.github.io}}}}



\author{
% Title-page order (updated 2026-11-09); part_6/author_contribution_list.tex follows this order.
% Each row must stay narrower than the text width: the author block is an unwrapped tabular.
\textbf{Shengqiong Wu}$^{1,\dagger}$, \,
\textbf{Yu Zhao}$^{2}$, \,
\textbf{Tianjie Ju}$^{2}$, \,
\textbf{Jundong Xu}$^{2}$, \,
\textbf{Bobo Li}$^{2}$, \\
\textbf{Jingkang Yang}$^{3}$, \,
\textbf{Shulin Tian}$^{3}$, \,
\textbf{Daoan Zhang}$^{4}$, \,
\textbf{Yanguang Zhao}$^{2}$, \,
\textbf{Xuan Ju}$^{5}$, \\
\textbf{Haoyu Zhao}$^{5}$, \,
\textbf{Yuchen Zhou}$^{2}$, \,
\textbf{Xiangtai Li}$^{3}$, \,
\textbf{Zhedong Zheng}$^{6}$, \,
\textbf{Joy Hsu}$^{7}$, \,
\textbf{Xin Kong}$^{8}$, \\
\textbf{Ruining Li}$^{1}$, \,
\textbf{Cheng Zhang}$^{9}$, \,
\textbf{Chengzu Li}$^{10}$, \,
\textbf{Jiafei Duan}$^{2}$, \,
\textbf{Yu (Bryan) Zhou}$^{11}$, \\
\textbf{Kenneth Li}$^{12}$, \,
\textbf{Chuanxia Zheng}$^{3}$, \,
\textbf{Lei Cui}$^{13}$, \,
\textbf{Furu Wei}$^{13}$, \,
\textbf{Ziwei Liu}$^{3}$, \\
\textbf{Steven K. Feiner}$^{14}$, \,
\textbf{Jiebo Luo}$^{4}$, \,
\textbf{Jianfei Cai}$^{9}$, \,
\textbf{Shuicheng Yan}$^{2}$, \,
\textbf{Ming-Hsuan Yang}$^{15}$, \,
\textbf{Hao Fei}$^{1,*}$
\vspace{.5em}
\\
$^{1}$University of Oxford, 
$^{2}$National University of Singapore,
$^{3}$Nanyang Technological University \\
$^{4}$University of Rochester,
$^{5}$The Chinese University of Hong Kong,
$^{6}$University of Macau\\
$^{7}$Stanford University,
$^{8}$Imperial College London,
$^{9}$Monash University,
$^{10}$University of Cambridge \\
$^{11}$University of California, Los Angeles,
$^{12}$Harvard University,
$^{13}$Microsoft Research \\
$^{14}$Columbia University,
$^{15}$University of California, Merced
}
% Uncomment to override  the `A preprint' in the header
%\renewcommand{\headeright}{Technical Report}
%\renewcommand{\undertitle}{Technical Report}
% \renewcommand{\shorttitle}{\textit{arXiv} Template}






%%% Add PDF metadata to help others organize their library
%%% Once the PDF is generated, you can check the metadata with
%%% $ pdfinfo template.pdf
\hypersetup{
pdftitle={Spatial Intelligence},
pdfsubject={AI},
}

\begin{document}

\maketitle

\vspace{-6mm}




\begin{abstract}
The pursuit of spatial intelligence, i.e., enabling agents to perceive, represent, reason about, and act within complex environments, stands as a foundational pillar in the journey toward Artificial General Intelligence (AGI). From embodied perception to multimodal understanding and interactive world modeling, spatial intelligence spans a broad spectrum of research domains and practical applications. 
This handbook offers a systematic and forward-looking overview of spatial intelligence, bridging insights across computer vision, natural language processing, 3D/4D modeling, robotics, and cognitive reasoning. We organize our exploration into five parts. Part~I motivates why spatial intelligence matters now, traces how the representation of space has evolved, and gives a working definition. Part~II lays the groundwork in \textit{Spatially Aware Representation}, where we study representations across single modalities (2D, 3D, 4D, language, audio, etc.) and their cross-modal alignment. Part~III follows the full pipeline \textit{From Perception to Generation}, including perception (e.g., detection, segmentation), reasoning (e.g., symbolic, geometric, LLM-driven), and generation (e.g., scene, video, world synthesis). Part~IV turns to \textit{Applications}, spanning embodied agents, virtual/augmented realities, geo-spatial systems, biomedical domains, intelligent assistants, and further fields such as construction, sports, and cultural heritage. Finally, Part~V, \textit{Discussions and Futures}, examines what spatial intelligence requires and discusses research gaps, benchmarking, and directions for real-world deployment.
By synthesizing principles from AI, cognitive science, and spatial computing, this handbook aims to serve as both a comprehensive reference for researchers and a practical guide for system developers. We hope it will inspire innovative directions toward building agents and systems that are spatially aware, robustly generalizable, and deeply grounded in real-world interactions.
% The homepage of the project is: \url{https://spatial-agi.github.io}.
\end{abstract}
% Footnotes for the author-list marks; \footnotetext[n] prints the n-th \fnsymbol (2 = dagger, 1 = *) without stepping the footnote counter.
{\renewcommand{\thefootnote}{\fnsymbol{footnote}}%
\footnotetext[2]{Contact: \href{mailto:shengqiongwu@gmail.com}{\texttt{shengqiongwu@gmail.com}}.}%
\footnotetext[1]{Hao Fei leads the project. Contact: \href{mailto:haofei7419@gmail.com}{\texttt{haofei7419@gmail.com}}.}}



\vspace{4mm}


% (epigraph cards moved to preface.tex)



% Page 2: the teaser gets the whole page (the epigraphs moved to the Preface).
\clearpage
\thispagestyle{empty}
\noindent
% \includegraphics[width=0.99\linewidth]{figures/teaser_chart_v16.jpg}
\includegraphics[width=\linewidth]{figures/teaser_v6_color.pdf}\par\vspace{1.5mm}
\noindent{\small\textbf{The handbook as one world.} The map is laid out like the book: representation (Part~II) at the center, the loop from perception through reasoning to generation (Part~III) on the inner ring, and the application domains (Part~IV) along the coast. The stepping stones below trace the five eras of Part~I, the horizon above holds the open questions of Part~V, and each chapter is pinned where it lives in this world.\par}






% \date{\today}
% % \correspondence{Haofei, \email{haofei37@nus.sg}}
% \metadata[Homepage]{\url{https://github.com/GuangSTrip/Awesome-Spatial-Intelligence}}


% keywords can be removed
% \keywords{
%     Large Language Models (LLMs)
%     \and Intelligent Agents
%     \and Modular AI Architectures
%     \and Cognitive Science-Inspired AI
%     \and Perception and Action Systems
%     \and Memory
%     \and World Model
%     \and Emotional AI
%     \and Reinforcement and Reward Systems
%     \and Multi-Agent Collaboration
%     \and Collective Intelligence
%     \and Autonomous Self-Improvement
%     \and AI Safety and Alignment
%     \and Ethical and Trustworthy AI
% }



\clearpage
\chapter*{Preface}

% Epigraphs: three equal-height cards read left to right as a timeline (mathematics, psychology, AI).
\colorlet{quoteaccent}{linkcobalt!70!black}
\tcbset{quotecard/.style={enhanced, colback=LightCobalt, colframe=LightCobalt, boxrule=0pt, arc=1.5pt,
  borderline north={1.4pt}{0pt}{linkcobalt}, left=7pt, right=7pt, top=6pt, bottom=6pt,
  fontupper=\small\itshape, before upper={\raggedright}, before lower={\raggedright},
  lower separated=false, space to upper, middle=5pt}}
\newcommand{\quotecardhead}[2]{{\normalfont\scriptsize\bfseries\color{quoteaccent}#1\enspace\textperiodcentered\enspace\textsc{#2}}\par\vspace{2pt}}
\vspace{-4mm}
\begin{tcbraster}[raster columns=3, raster equal height, raster column skip=3mm]
\begin{tcolorbox}[quotecard]
\quotecardhead{1902}{mathematics}
``To localise an object simply means that we represent to ourselves the movements that must take place to reach it.''
\tcblower
{\normalfont\footnotesize\textsc{Henri Poincar\'{e}}\\\textit{Science and Hypothesis}\strut}
\end{tcolorbox}
\begin{tcolorbox}[quotecard]
\quotecardhead{1983}{psychology}
``The capacities to perceive the visual world accurately, to perform transformations and modifications upon one's initial perceptions, and to be able to re-create aspects of one's visual experience, even in the absence of relevant physical stimuli.''
\tcblower
{\normalfont\footnotesize\textsc{Howard Gardner}\\\textit{Frames of Mind}\strut}
\end{tcolorbox}
\begin{tcolorbox}[quotecard]
\quotecardhead{2025}{artificial intelligence}
``Spatial intelligence is the scaffolding upon which our cognition is built.''
\tcblower
{\normalfont\footnotesize\textsc{Fei-Fei Li}\\\textit{From Words to Worlds}\strut}
\end{tcolorbox}
\end{tcbraster}
\vspace{5mm}


Spatial intelligence lies at the core of human cognition.
We perceive our surroundings, remember places, plan routes, manipulate objects, and coordinate with one another in shared physical and virtual spaces, and we do all of this by building and continuously revising internal models of the world.
Machines are only now acquiring the rudiments of this ability.
Today's AI systems write fluent prose and reason impressively over text, yet they still struggle to judge whether a sofa will fit through a doorway, to keep a scene coherent while moving through it, or to anticipate what happens when an object is pushed.
That gap has become one of the central open problems of contemporary AI.
Advances in vision-language models, 3D and 4D perception, generative modeling, and interactive world models are now converging from different directions on the same goal: machines that perceive, reason, and act in space.
This handbook is our attempt to give that convergence a coherent shape.

The immediate motivation is a long-standing fragmentation.
Research on space in AI has been carried out in largely separate communities, spanning SLAM and state estimation, 3D vision and neural rendering, spatial language grounding, geographic information science, embodied navigation and manipulation, and the psychology and neuroscience of spatial cognition, each with its own representations, benchmarks, and vocabulary.
Progress in one thread has rarely been legible to the others.
We set out to organize the field as a whole.
The organizing commitment of this handbook is \emph{representation-centric}: spatial intelligence is the ability of a system to construct, maintain, and manipulate spatial representations of the world in support of perception, reasoning, prediction, and embodied action.
Reading the field through this single lens, from the symbolic maps of early AI, through probabilistic geometry and learned features, to today's foundation models and executable world states, is what allows historically disjoint literatures to be placed on common ground.

A handful of questions shaped every chapter.
What makes a system spatially intelligent, beyond performing well on any particular task?
How should geometric, semantic, and commonsense structure be integrated within one representation, and aligned across modalities as different as pixels, point clouds, language, and audio?
When does spatial knowledge learned in one environment transfer to another?
How do planning and action emerge from grounded spatial understanding, and how should the field measure genuine progress rather than benchmark saturation?
None of them has a final answer yet.
Each now admits systematic treatment, and the handbook is organized to provide it.

The volume proceeds in five parts.
Part~I lays the foundation: why spatial intelligence has risen to prominence now, how the idea evolved across five representational eras from symbolic spatial computing to world-centered intelligence, and how we define and formalize the concept itself.
Part~II examines spatially aware representation, covering single modalities (2D, 3D, 4D, language, and audio), domain-specific representations, and the cross-modal alignment and fusion that bind them.
Part~III follows the full functional pipeline from perception, through reasoning in its rule-based, learning-based, foundation-model-driven, and hybrid forms, to spatial reconstruction and generation.
Part~IV turns to applications, where these capabilities are assembled into working systems: embodied agents and robots, virtual and extended realities, geo-spatial systems, biomedicine and the natural sciences, intelligent assistants, and a range of further domains.
Part~V closes with open discussions, unresolved challenges, and the directions we believe will matter most in the coming years.
Chapter~\ref{chapter:structure} offers tailored reading paths through this material for students, researchers, engineers, and cross-disciplinary explorers.

We wrote this handbook for several audiences at once: researchers and graduate students in computer vision, robotics, and multimodal learning; engineers building embodied agents, simulation platforms, and spatial computing systems; and scholars in cognitive science, neuroscience, and geography who want to see how their fields' insights are being absorbed into AI.
Wherever possible we favor conceptual clarity over exhaustive enumeration, and we connect each technical thread back to the shared framework so that the whole remains visible through the parts.

A handbook about a fast-moving field is necessarily a snapshot.
Some of the systems discussed here will be surpassed before long, and we regard that as a feature of the field: the framework is designed to outlast its examples.
We call Chapter~\ref{chapter:definition} ``A Definition in Motion'' deliberately.
Spatial intelligence is still being defined by the community that builds it, and our hope is that this volume gives that community a common map, drawn honestly, mid-journey.

\clearpage


\setcounter{tocdepth}{2}
\tableofcontents

%!TEX root = ../main.tex
\clearpage
\chapter*{Notation}
\label{sec:notations-notation}

Symbols are listed in chapter order, with the meanings used in the referenced sections.

\definecolor{verylightgray}{gray}{0.95}

\begin{longtable}{lll}
	\textbf{Symbol} & \textbf{Description} & \textbf{Reference}\\
    % %%% Agent framework notations
    % $\mathcal{W}$ & The world with society systems. & Sec.~\ref{subsec:agent-framework-symbols}\\
    % $\mathcal{S}$ & State space of an environment.  & Sec.~\ref{subsec:agent-framework-symbols}\\
    % $s_t \in \mathcal{S}$ & Environment's state at time $t$. & Sec.~\ref{subsec:agent-framework-symbols}\\
    % $\mathcal{O}$  & Observation space. & Sec.~\ref{subsec:agent-framework-symbols}\\
    % $o_t \in \mathcal{O}$ & Observation at time $t$. & Sec.~\ref{subsec:agent-framework-symbols}\\
    % $\mathcal{A}$  & Agent's action space. & Sec.~\ref{subsec:agent-framework-symbols}\\
    % $a_t \in \mathcal{A}$ & Agent's action output at time $t$. & Sec.~\ref{subsec:agent-framework-symbols}\\
    % $\mathcal{M}$ & Mental states space. & Sec.~\ref{subsec:agent-framework-symbols}\\
    % $M_t \in \mathcal{M}$ & Agent's mental state at time $t$. & Sec.~\ref{subsec:agent-framework-symbols}\\
    % $M_t^{\mathrm{mem}}$  & \textit{Memory} component in $M_t$. & Sec.~\ref{subsec:agent-framework-symbols}\\
    % $M_t^{\mathrm{wm}}$  & \textit{World model} component in $M_t$. & Sec.~\ref{subsec:agent-framework-symbols}\\
    % $M_t^{\mathrm{emo}}$ & \textit{Emotion} component in $M_t$. & Sec.~\ref{subsec:agent-framework-symbols}\\
    % $M_t^{\mathrm{goal}}$ & \textit{Goal} component in $M_t$. & Sec.~\ref{subsec:agent-framework-symbols}\\
    % $M_t^{\mathrm{rew}}$  & \textit{Reward/Learning} signals in $M_t$. & Sec.~\ref{subsec:agent-framework-symbols}\\
    % $\mathrm{L}$ & Agent's learning function. & Sec.~\ref{subsec:agent-framework-symbols}\\
    % $\mathrm{R}$ & Agent's reasoning function. & Sec.~\ref{subsec:agent-framework-symbols}\\
    % $\mathrm{C}$ & Agent's cognition function. & Sec.~\ref{subsec:agent-framework-symbols}\\
    % $\mathrm{E}$ & Action execution (effectors). & Sec.~\ref{subsec:agent-framework-symbols}\\
    % $\mathrm{T}$ & Environment transition. & Sec.~\ref{subsec:agent-framework-symbols}\\
    % %%% other notations (put a brief summary here)
    % % Self Improvement Agent Space
    % $\mathrm{P}$           & Prompt text / prompt set fed to the LLM.                                    & Sec.~\ref{sec:prompt-optim} \\
    % $K_{\text{wf}}$        & Agentic workflow composed of interconnected nodes and control flows.        & Sec.~\ref{sec:workflow-optimization} \\
    % $N$           & Node representing a single LLM invocation used in $K_{\text{wf}}$. & Sec.~\ref{sec:workflow-optimization} \\
    % % Scientific knowledge discovery 
    % $\theta$ & Parameters of the world model $M_t^{\mathrm{wm}}$. & Sec.~\ref{subsec:KL_measure_intelligence} \\
    % $P_\theta$ & Predicted data distribution. & Sec.~\ref{subsec:KL_measure_intelligence} \\
    % $P_\mathcal{W}$ & True data distribution in the real world. & Sec.~\ref{subsec:KL_measure_intelligence} \\
    % $\mathcal{K}$ & Space of known data and information. & Sec.~\ref{subsec:KL_measure_intelligence} \\
    % $\mathcal{U}$ & Space of unknown data and information. & Sec.~\ref{subsec:KL_measure_intelligence} \\
    % ${\bf x}$ & Dataset representing scientific knowledge. & Sec.~\ref{subsec:KL_measure_intelligence} \\
    % ${\bf x}_{\mathrm{K}}$ & Known dataset sampled from $\mathcal{K}$. & Sec.~\ref{subsec:KL_measure_intelligence} \\
    % ${\bf x}_{\mathrm{U}}$ & Unknown dataset sampled from $\mathcal{U}$. & Sec.~\ref{subsec:KL_measure_intelligence} \\
    % $D_0$ & KL divergence from $P_\mathcal{W}$ to $P_\theta$ at time $t=0$. & Sec.~\ref{subsec:KL_measure_intelligence} \\
    % $D_{\mathrm{K}}$ & KL divergence from $P_\mathcal{W}$ to $P_\theta$ after acquiring knowledge. & Sec.~\ref{subsec:KL_measure_intelligence} \\
    % $IQ_t^{\mathrm{agent}}$ & Agent's intelligence at time $t$. & Sec.~\ref{subsec:KL_measure_intelligence} \\
    % $\Delta $ & Subspace of $\mathcal{U}$ for knowledge expansion. & Sec.~\ref{subsec:IQ_increase_w_knowledge} \\
    % ${\bf x}_\Delta$ & Dataset from $\Delta$. & Sec.~\ref{subsec:IQ_increase_w_knowledge} \\
    % $\Theta$ & Space of possible world model parameters $\theta$. & Sec.~\ref{subsec:knowledge_expansion_strategy} \\
    % $\theta^*_{\mathrm{K},t}$ & Optimal world model parameters given the agent's knowledge at time $t$. & Sec.~\ref{subsec:knowledge_expansion_strategy} \\
    % $D_{\mathrm{K},\Theta}^{\rm min}$ & Minimum unknown given the agent's knowledge and $\Theta$. & Sec.~\ref{subsec:knowledge_expansion_strategy} \\
    % % Basic Sequences
    % $\mathbf{x}_{1:n}$ & Input token sequence. & Sec.~\ref{ssec:brain_sec_threat} \\
    % $\mathbf{y}$ & Generated output sequence. & Sec.~\ref{ssec:brain_sec_threat} \\
    % % Jailbreak and Adversarial Attacks
    % $p$ & Probability of generating $\mathbf{y}$ given $\mathbf{x}_{1:n}$. & Sec.~\ref{sssec:jailbreak} \\
    % $\tilde{\mathbf{x}}_{1:n}$ & Perturbed input sequence. & Sec.~\ref{sssec:jailbreak} \\
    % $\mathcal{R}^*$ & Ideal alignment reward (measuring adherence to safety/ethical guidelines). & Sec.~\ref{sssec:jailbreak} \\
    % $\mathbf{y}^\star$ & Jailbreak output induced by perturbations. & Sec.~\ref{sssec:jailbreak} \\
    % $\mathcal{A}$ & a set of safety/ethical guidelines & Sec.~\ref{sssec:jailbreak} \\
    % $\mathcal{T}$ & the distribution or set of possible jailbreak instructions. & Sec.~\ref{sssec:jailbreak} \\
    % $\mathcal{L}^{adv}$ & Jailbreak loss. & Sec.~\ref{sssec:jailbreak} \\
    % % Prompt Injection Attacks
    % $\mathbf{p}$ & Prompt injected into the original input. & Sec.~\ref{sssec:promptinj} \\
    % $\mathbf{x}'$ & Combined (injected) input sequence. & Sec.~\ref{sssec:promptinj} \\
    % $\mathcal{L}^{inject}$ & Prompt injection loss. & Sec.~\ref{sssec:promptinj} \\
    % $\mathbf{p}^\star$ & Optimal injected prompt minimizing $\mathcal{L}^{inject}$. & Sec.~\ref{sssec:promptinj} \\
    % $\mathcal{P}$ & Set of feasible prompt injections. & Sec.~\ref{sssec:promptinj} \\
    % % Hallucination Metrics
    % $e_{x_i}\in\mathbb{R}^{d_e}$ & Embedding of token $x_i$ in a $d_e$-dimensional space. & Sec.~\ref{sssec:hallucination} \\
    % $\mathrm{W}_Q,\mathrm{W}_K,\mathrm{W}_V$ & Projection matrices for query, key, and value. & Sec.~\ref{sssec:hallucination} \\
    % $A_{ij}$ & Attention score between tokens $i$ and $j$. & Sec.~\ref{sssec:hallucination} \\
    % $o_i$ & Contextual representation of token $i$ (weighted sum result). & Sec.~\ref{sssec:hallucination} \\
    % $\delta_{x_i}$ & Perturbation applied to $e_{x_i}$, satisfying $\|\delta_{x_i}\|\leq\epsilon$. & Sec.~\ref{sssec:hallucination} \\
    % $\tilde{e}_{x_i}$ & Perturbed token embedding. & Sec.~\ref{sssec:hallucination} \\
    % $A_{ij}^\Delta$ & Attention score under perturbation. & Sec.~\ref{sssec:hallucination} \\
    % $\tilde{o}_i$ & Updated token representation under perturbation. & Sec.~\ref{sssec:hallucination} \\
    % $\mathcal{H}$ & Hallucination metric. & Sec.~\ref{sssec:hallucination} \\
    % % Misalignment Issues
    % $\mathcal{R}$ & Actual alignment reward of the model's output. & Sec.~\ref{sssec:misalignment} \\
    % $\Delta_{\text{align}}$ & Alignment gap. & Sec.~\ref{sssec:misalignment} \\
    % $\mathcal{L}^{misalign}$ & Misalignment loss. & Sec.~\ref{sssec:misalignment} \\
    % $\lambda$ & Trade-off parameter for the alignment gap in the misalignment loss. & Sec.~\ref{sssec:misalignment} \\
    % % Poisoning Attacks
    % $\mathcal{D}$ & Clean training dataset. & Sec.~\ref{sssec:poison} \\
    % $\tilde{\mathcal{D}}$ & Poisoned training dataset. & Sec.~\ref{sssec:poison} \\
    % $\theta$ & Model parameters. & Sec.~\ref{sssec:poison} \\
    % $\theta^\star$ & Model parameters learned from the poisoned dataset. & Sec.~\ref{sssec:poison} \\
    % $\theta_{\text{clean}}$ & Model parameters obtained using the clean dataset. & Sec.~\ref{sssec:poison} \\
    % $\Delta_\theta$ & Deviation of model parameters due to poisoning. & Sec.~\ref{sssec:poison} \\
    % $t$ & Backdoor trigger. & Sec.~\ref{sssec:poison} \\
    % $\mathcal{B}$ & Backdoor success rate. & Sec.~\ref{sssec:poison} \\
    % $\mathbb{I}$ & Indicator function. & Sec.~\ref{sssec:poison} \\
    % $\mathcal{Y}_{\text{malicious}}$ & Set of undesirable outputs. & Sec.~\ref{sssec:poison} \\
    % % Privacy Concerns
    % $g$ & Function estimating the probability that input $\mathbf{x}$ was in the training set, with range $[0,1]$. & Sec.~\ref{ssec:brain_priv_threat} \\
    % $\eta$ & Threshold for membership inference. & Sec.~\ref{ssec:brain_priv_threat} \\
    % $\mathbf{x}^\star$ & Reconstructed training sample in a data extraction attack. & Sec.~\ref{ssec:brain_priv_threat} \\
    % $\mathbf{p}_{sys}$ & System prompt defining the agent's internal guidelines. & Sec.~\ref{ssec:brain_priv_threat} \\
    % $\mathbf{p}_{user}$ & User prompt. & Sec.~\ref{ssec:brain_priv_threat} \\
    % $\mathbf{p}^\star$ & Reconstructed prompt via inversion. & Sec.~\ref{ssec:brain_priv_threat} \\
$\mathcal{E}$ & Physical or simulated environment. & Sec.~\ref{subsec:spatial-intelligence-formalization} \\
$\mathcal{S}$ & AI system or agent. & Sec.~\ref{subsec:spatial-intelligence-formalization} \\
$\mathcal{T},\;t\in\mathcal{T}$ & Task space and a task in that space. & Sec.~\ref{subsec:spatial-intelligence-formalization} \\
$\mathcal{R}_{\mathcal{S}}$ & Internal spatial representation space of the system. & Sec.~\ref{subsec:spatial-intelligence-formalization} \\
$\Phi_{\mathcal{S}}$ & Mapping from the environment to the internal spatial representation space. & Sec.~\ref{subsec:spatial-intelligence-formalization} \\
$\mathcal{F}_{\mathcal{S}}$ & Set of operators for perception, reasoning, prediction, and action. & Sec.~\ref{subsec:spatial-intelligence-formalization} \\
$\mathcal{O},\;\mathcal{A}$ & Observation space and action space, respectively. & Sec.~\ref{subsec:spatial-intelligence-formalization} \\
$f^{\mathcal{S}}_{\mathrm{perc}}$ & Operator that maps observations to spatial representations. & Sec.~\ref{subsec:spatial-intelligence-formalization} \\
$f^{\mathcal{S}}_{\mathrm{reason}}$ & Reasoning operator inferring relations and constraints in the current representation. & Sec.~\ref{subsec:spatial-intelligence-formalization} \\
$f^{\mathcal{S}}_{\mathrm{pred}}$ & Prediction operator mapping a representation and an action to a future representation. & Sec.~\ref{subsec:spatial-intelligence-formalization} \\
$f^{\mathcal{S}}_{\mathrm{act}}$ & Action operator mapping a spatial representation to an action. & Sec.~\ref{subsec:spatial-intelligence-formalization} \\
$\pi_{\mathcal{S}}$ & Task-conditioned policy for selecting actions. & Sec.~\ref{subsec:spatial-intelligence-formalization} \\
$\mathcal{R}^{\mathcal{S}}_t$ & Internal spatial representation state at time $t$. & Sec.~\ref{subsec:spatial-intelligence-formalization} \\
$a^{\mathcal{S}}_t,\;o_{t+1}$ & Selected action and subsequent observation. & Sec.~\ref{subsec:spatial-intelligence-formalization} \\
$\Psi_{\mathcal{S}}$ & Update function for the internal spatial representation. & Sec.~\ref{subsec:spatial-intelligence-formalization} \\
$\Omega$ & Image grid or domain of pixel coordinates. & Sec.~\ref{subsec:p-2-sec-4-2D-image-and-video-representations} \\
$I,\;V$ & Image and video, respectively. & Sec.~\ref{subsec:p-2-sec-4-2D-image-and-video-representations} \\
$S,\;\mathcal{L}$ & Segmentation map and its set of labels. & Sec.~\ref{subsec:p-2-sec-4-2D-segmentation-maps} \\
$D(u,v)$ & Depth value at pixel $(u,v)$. & Sec.~\ref{subsec:p-2-sec-4-2D-depth-maps} \\
$\mathbf{N}(u,v)$ & Unit surface normal at pixel $(u,v)$. & Sec.~\ref{subsec:p-2-sec-4-2D-surface-normal-maps} \\
$\mathbf{F}(u,v)$ & Optical flow giving horizontal and vertical pixel displacement. & Sec.~\ref{subsec:p-2-sec-4-2D-optical-flow-maps} \\
$\mathcal{K},\;H_k$ & Set of keypoint locations and the heatmap for keypoint $k$. & Sec.~\ref{subsec:p-2-sec-4-2D-2d-keypoints} \\
$P=\{p_i\}$ & Point cloud with 3D coordinates and optional attributes. & Sec.~\ref{sec:rep_3d_point} \\
$M=(V,F)$ & Triangle mesh with vertex set $V$ and face set $F$. & Sec.~\ref{sec:rep_3d_mesh} \\
$f(\mathbf{x})$ & Implicit scalar field whose level set defines a surface; a signed distance function in the SDF case. & Sec.~\ref{sec:rep_3d_implicit} \\
$F_\Theta$ & Neural radiance field mapping position and viewing direction to color and volume density. & Sec.~\ref{sec:rep_3d_nerf} \\
$\mathbf{x},\;\mathbf{d}$ & 3D query position and viewing direction in a radiance field. & Sec.~\ref{sec:rep_3d_nerf} \\
$\sigma,\;\mathbf{c}$ & Volume density and emitted RGB color in a radiance field. & Sec.~\ref{sec:rep_3d_nerf} \\
$\boldsymbol{\mu},\;\boldsymbol{\Sigma}$ & Center and covariance matrix of a 3D Gaussian primitive. & Sec.~\ref{sec:rep_3d_gaussian} \\
$\alpha_i,\;c_i$ & Opacity and color of the $i$-th Gaussian primitive. & Sec.~\ref{sec:rep_3d_gaussian} \\
$\Delta\mathbf{x}(\mathbf{x},t)$ & Position- and time-dependent displacement to a canonical 3D position. & Sec.~\ref{sec:rep_4d_nerf} \\
$\mathcal{G}(t)$ & Time-dependent collection of Gaussian primitives in a dynamic scene. & Sec.~\ref{sec:rep_4d_gaussian} \\
$b_i,\;c_i,\;s_i$ & Bounding box, class label, and confidence score of detection $i$. & Sec.~\ref{sec:6-1-2} \\
$f_\theta,\;\theta$ & Parameterized prediction model and its learned parameters. & Sec.~\ref{sec:keypoint_landmark} \\
$\mathbb{I}(\cdot)$ & Indicator function, equal to one when its condition holds and zero otherwise. & Sec.~\ref{sec:keypoint_landmark} \\
$Z,\;f,\;B,\;d$ & Stereo depth, focal length, camera baseline, and disparity, with $Z=fB/d$. & Sec.~\ref{sec:depth_estimation} \\
$K$ & Camera intrinsic matrix. & Sec.~\ref{sec:height_distance} \\
$R,\;\mathbf{t}$ & Rotation and translation mapping world coordinates into the camera frame. & Sec.~\ref{sec:height_distance} \\
$\mathbf{T}_t$ & Camera-to-world rigid transformation for frame $t$. & Sec.~\ref{sec:visual_odometry} \\
$\Delta\mathbf{T}_{i\to j}$ & Pose of camera frame $j$ expressed in camera frame $i$. & Sec.~\ref{sec:visual_odometry} \\
$\mathcal{G}$ & Scene graph of subject--predicate--object triplets. & Sec.~\ref{sec:scene_graph_generation} \\
$\mathbf{v}(\mathbf{x})$ & Scene flow: 3D displacement of the surface point observed at pixel $\mathbf{x}$. & Sec.~\ref{sec:scene_flow} \\
$\mathcal{T}_k$ & Time-indexed trajectory of tracked object $k$. & Sec.~\ref{sec:object_tracking} \\
$\mathrm{SO}(3),\;\mathrm{SE}(3)$ & Groups of 3D rotations and rigid transformations, respectively. & Sec.~\ref{subsec:p3_s7_1_1} \\
$\mathbf{F}$ & Fundamental matrix relating corresponding image points in two views. & Sec.~\ref{subsec:p3_s7_1_1} \\
$\mathbf{q},\;\mathbf{J}(\mathbf{q})$ & Joint configuration and the Jacobian mapping joint velocities to task-space velocities. & Sec.~\ref{subsec:p3_s7_1_1} \\
$\mathcal{C}_{\mathrm{free}}$ & Feasible configuration space used in reachability and motion planning. & Sec.~\ref{subsec:p3_s7_1_1} \\
$\mathcal{W}$ & Set of contact wrenches used to test grasp force closure. & Sec.~\ref{subsec:p3_s7_1_1} \\
$\mathbf{s}_t,\;\mathbf{u}_t$ & Physical system state and applied action in the dynamics model. & Sec.~\ref{subsec:p3_s7_1_2} \\
$M(\mathbf{q})$ & Configuration-dependent mass matrix. & Sec.~\ref{subsec:p3_s7_1_2} \\
$G=(V,E,w)$ & Planning graph with vertices, edges, and edge weights. & Sec.~\ref{subsec:p3_s7_1_3} \\
$\operatorname{do}(X=x)$ & Intervention that sets variable $X$ to $x$ in a causal model. & Sec.~\ref{subsec:p3_s7_1_6} \\
$z^\star$ & Optimal solution returned by a differentiable optimization layer. & Sec.~\ref{subsec:p3_s7_2_4} \\
$\pi_\theta$ & Learned policy for selecting actions. & Sec.~\ref{subsec:p3_s7_2_5} \\
$r_t,\;\gamma$ & Reward at time $t$ and the discount factor. & Sec.~\ref{subsec:p3_s7_2_5} \\
$Q^*(s,a)$ & Maximum expected discounted return from taking action $a$ in state $s$. & Sec.~\ref{subsec:p3_s7_2_5} \\
$R_t$ & Discounted return from time $t$ onward. & Sec.~\ref{subsec:p3_s7_2_5} \\
$V_\phi,\;A_\phi$ & State-value function and advantage function. & Sec.~\ref{subsec:p3_s7_2_5} \\
$a_{t:t+H}$ & Action sequence predicted over horizon $H$ by a vision-language-action policy. & Sec.~\ref{subsec:p3_s7_3_2} \\
$\tau=(t_1,\ldots,t_L)$ & Sequence of intermediate reasoning steps in a chain of thought. & Sec.~\ref{subsec:p3_s7_3_3} \\
$\mathcal{M}_t$ & External spatial memory after the update at time $t$. & Sec.~\ref{subsec:p3_s7_3_6} \\
$\mathcal{X},\;\mathbf{x}$ & Space of tokenized text prompts and an input prompt. & Sec.~\ref{subsec:p3_s8_1_1} \\
$\mathcal{Y},\;\mathbf{y}$ & Pixel-space image domain and an image in that domain. & Sec.~\ref{subsec:p3_s8_1_1} \\
$\mathcal{Z},\;\mathbf{z}$ & Latent representation space and a latent code. & Sec.~\ref{subsec:p3_s8_1_1} \\
$E,\;D$ & Image encoder and decoder of the latent autoencoder. & Sec.~\ref{subsec:p3_s8_1_1} \\
$\phi,\;\mathcal{E}$ & Text encoder and its conditioning embedding space. & Sec.~\ref{subsec:p3_s8_1_1} \\
$G_\theta$ & Conditional generator or sampler. & Sec.~\ref{subsec:p3_s8_1_1} \\
$p_\theta(\cdot\mid\cdot)$ & Learned conditional distribution over generated outputs or latent codes. & Sec.~\ref{subsec:p3_s8_1_1} \\
$\epsilon_\theta$ & Noise predictor used in diffusion generation. & Sec.~\ref{subsec:p3_s8_1_1} \\
$\mathbf{z}_t$ & Latent code at diffusion or flow time $t$. & Sec.~\ref{subsec:p3_s8_1_1} \\
$v_\theta$ & Learned velocity field used in flow matching. & Sec.~\ref{subsec:p3_s8_1_1} \\
$\mathbf{c},\;\mathcal{C}$ & Structural conditioning signal and its input space. & Sec.~\ref{subsec:p3_s8_1_2} \\
$\kappa,\;\mathcal{C}^{\prime}$ & Condition encoder and the space of encoded conditions. & Sec.~\ref{subsec:p3_s8_1_2} \\
$\mathcal{M}$ & Mask identifying the region to be edited. & Sec.~\ref{subsec:p3_s8_1_3} \\
$\mathcal{R},\;\mathbf{r}$ & Space of 3D representations and a 3D asset in that space. & Sec.~\ref{subsec:p3_s8_2_1} \\
$\pi,\;\pi_v$ & Camera pose used for rendering, with $v$ indexing a view. & Sec.~\ref{subsec:p3_s8_2_1} \\
$\mathrm{Render}$ & Renderer mapping a 3D representation and a camera pose to an image. & Sec.~\ref{subsec:p3_s8_2_1} \\
$\mathcal{L}_{\mathrm{SDS}}$ & Score-distillation objective for optimizing a representation with a frozen diffusion prior. & Sec.~\ref{subsec:p3_s8_2_1} \\
$\mathcal{L}_{\mathrm{rec}}$ & Loss measuring agreement between a reconstruction and its observations. & Sec.~\ref{subsec:p3_s8_2_3} \\
$\lambda,\;\Omega$ & Weight and regularizer used in scene reconstruction. & Sec.~\ref{subsec:p3_s8_2_3} \\
$\Lambda$ & Operator lifting a panorama or a set of panoramic views into a 3D representation. & Sec.~\ref{subsec:p3_s8_2_4} \\
$E_\theta$ & Editing operator mapping a 3D representation and a condition to an edited representation. & Sec.~\ref{subsec:p3_s8_2_5} \\
$\mathbf{m}^{1:T}$ & Sequence of body poses describing a motion. & Sec.~\ref{subsec:p3_s8_4_1} \\
$\mathbf{r}^{1:T}$ & Time-indexed sequence of 3D representations for a dynamic scene. & Sec.~\ref{sec:p3_s8_5} \\
$W_\theta$ & Learned world model predicting the next state from state, action, and observation. & Sec.~\ref{sec:p3_s8_6} \\
$S(t),\;A(t),\;O(t)$ & World state, action, and observation, respectively. & Sec.~\ref{sec:p3_s8_6} \\
$E_\phi,\;P_\theta$ & Observation encoder and latent-state predictor in a representation world model. & Sec.~\ref{subsec:p3_s8_6_rep} \\
$\operatorname{sg}(\cdot)$ & Stop-gradient operator used in representation learning. & Sec.~\ref{subsec:p3_s8_6_rep} \\
$M(t)$ & Latent mental state of another agent, including beliefs, desires, and intentions. & Sec.~\ref{subsec:p3_s8_6_mental} \\
$s_t^{\mathrm{phy}},\;s_t^{\mathrm{ment}}$ & Physical and mental components of a joint world state. & Sec.~\ref{subsec:p3_s8_6_mental} \\
$\rho,\;\mathcal{L}$ & Symbolic scene program and its language. & Sec.~\ref{subsubsec:p3_s8_7_task} \\
$R(\rho)$ & Execution or rendering of symbolic scene program $\rho$. & Sec.~\ref{subsubsec:p3_s8_7_task} \\
$M_t$ & Semantic environment map maintained during warehouse localization and mapping. & Sec.~\ref{subsec:p4-sec9-5-pivotal-task-2-localization-and-mapping} \\
\end{longtable}


    



    
    


%!TEX root = ../main.tex
\clearpage
\chapter*{Glossary}
\label{sec:glossary-glossary}

Terms and abbreviations are listed in the order of their first appearance in the chapters.

\definecolor{verylightgray}{gray}{0.95}

\begin{longtable}{lll}
	\textbf{Term} & \textbf{Description} & \textbf{Reference}\\

    AR & \textbf{Augmented Reality.} Overlaying virtual content registered to a live view of the physical world. & Chap.~\ref{chapter:intro} \\
    
    VR & \textbf{Virtual Reality.} Immersion of the user in a fully computer-generated environment. & Chap.~\ref{chapter:intro} \\
    
    LLMs & \textbf{Large Language Models.} Models trained on large text corpora to process and generate language. & Sec.~\ref{subsec:ai-paradigm-shift} \\
    
    RL & \textbf{Reinforcement Learning.} Learning to select actions through interaction and reward feedback. & Sec.~\ref{subsec:ai-methodological-shift} \\
    
    GIS & \textbf{Geographic Information Systems.} Systems for storing, managing, and analyzing geographically referenced data. & Sec.~\ref{subsec:disciplinary-convergence} \\
    
    VQA & \textbf{Visual Question Answering.} Answering natural-language questions about visual content. & Sec.~\ref{subsec:disciplinary-convergence} \\
    
    VLN & \textbf{Vision-and-Language Navigation.} Navigation using visual observations and natural-language instructions. & Sec.~\ref{subsec:disciplinary-convergence} \\
    
    GIScience & \textbf{Geographic Information Science.} The study of geographic information, its representation, and its analysis. & Sec.~\ref{subsec:era-geometric-probabilistic} \\
    
    SLAM & \textbf{Simultaneous Localization and Mapping.} Joint estimation of an agent's pose and a map of its environment. & Sec.~\ref{subsec:era-geometric-probabilistic} \\
    
    TSDF & \textbf{Truncated Signed Distance Function.} A signed distance function whose values are limited to a fixed range around a surface. & Sec.~\ref{subsec:era-geometric-probabilistic} \\
    
    SDF & \textbf{Signed Distance Function.} A function giving the distance to a surface, with a sign distinguishing its two sides. & Sec.~\ref{subsec:era-geometric-probabilistic} \\
    
    VO & \textbf{Visual Odometry.} Estimation of camera motion from a sequence of images. & Sec.~\ref{subsec:era-learning-based} \\
    
    LiDAR & \textbf{Light Detection and Ranging.} Sensing that uses laser light to measure distances. & Sec.~\ref{subsec:era-learning-based} \\
    
    CT & \textbf{Computed Tomography.} Imaging that reconstructs cross-sectional structure from X-ray measurements. & Sec.~\ref{subsec:era-learning-based} \\
    
    MRI & \textbf{Magnetic Resonance Imaging.} Imaging that uses magnetic fields and radiofrequency signals to depict internal anatomy. & Sec.~\ref{subsec:era-learning-based} \\
    
    CNNs & \textbf{Convolutional Neural Networks.} Neural networks that use convolutional filters to extract local features. & Sec.~\ref{subsec:era-learning-based} \\
    
    BEV & \textbf{Bird's-Eye View.} A top-down representation of a scene. & Sec.~\ref{subsec:era-learning-based} \\
    
    ViT & \textbf{Vision Transformer.} A transformer that processes an image as a sequence of patch tokens. & Sec.~\ref{subsec:era-foundation-models} \\
    
    NeRF & \textbf{Neural Radiance Field.} A continuous representation of scene density and view-dependent radiance for novel-view synthesis. & Sec.~\ref{subsec:era-foundation-models} \\
    
    CoT & \textbf{Chain-of-Thought.} A sequence of intermediate reasoning steps used to solve a problem. & Sec.~\ref{subsec:era-foundation-models} \\
    
    VLA & \textbf{Vision-Language-Action Model.} A model that maps visual observations and language instructions to actions. & Sec.~\ref{subsec:era-foundation-models} \\
    
    VLMs & \textbf{Vision-Language Models.} Models that jointly process visual and linguistic information. & Sec.~\ref{subsec:era-foundation-models} \\
    
    3DGS & \textbf{3D Gaussian Splatting.} A method for representing and rendering scenes using optimized 3D Gaussians. & Sec.~\ref{subsec:era-foundation-models} \\
    
    JEPA & \textbf{Joint Embedding Predictive Architecture.} An architecture that predicts target representations in an embedding space. & Sec.~\ref{subsec:era-world-centered} \\
    
    XR & \textbf{Extended Reality.} An umbrella term for virtual, augmented, and mixed reality. & Sec.~\ref{subsec:p1-sec2-compatibility-with-existing-research-and-paradigms} \\
    
    RGB-D & \textbf{RGB and Depth.} Image data combining color channels with depth measurements. & Sec.~\ref{sec:rep_3d} \\
    
    MLP & \textbf{Multilayer Perceptron.} A feedforward neural network with fully connected layers. & Sec.~\ref{sec:rep_3d_nerf} \\
    
    SfM & \textbf{Structure from Motion.} Recovery of camera poses and 3D structure from overlapping images. & Sec.~\ref{sec:rep_3d_gaussian} \\
    
    DSL & \textbf{Domain-Specific Language.} A programming language designed for a particular class of tasks. & Sec.~\ref{sec:cross_modal_overview} \\
    
    IMU & \textbf{Inertial Measurement Unit.} A sensor unit that measures acceleration and angular velocity. & Sec.~\ref{cross_modality_vision_vision} \\
    
    IoU & \textbf{Intersection over Union.} The size of the overlap between two regions divided by the size of their union. & Sec.~\ref{sec:6-1-2} \\
    
    mAP & \textbf{Mean Average Precision.} Average precision averaged across classes or evaluation categories. & Sec.~\ref{sec:6-1-2} \\
    
    MLLMs & \textbf{Multimodal Large Language Models.} Large language models extended to accept images, video, or audio together with text. & Sec.~\ref{sec:6-1-2} \\
    
    PnP & \textbf{Perspective-n-Point.} Estimation of camera pose from corresponding 3D points and 2D image locations. & Sec.~\ref{sec:keypoint_landmark} \\
    
    MVS & \textbf{Multi-View Stereo.} Dense reconstruction of 3D geometry from images with known camera poses. & Sec.~\ref{sec:depth_estimation} \\
    
    UAV & \textbf{Unmanned Aerial Vehicle.} An aircraft operated without a pilot on board. & Sec.~\ref{sec:height_distance} \\
    
    RANSAC & \textbf{Random Sample Consensus.} Robust model fitting by testing hypotheses from randomly sampled subsets of observations. & Sec.~\ref{sec:visual_odometry} \\
    
    GPS & \textbf{Global Positioning System.} The United States satellite navigation system, one of the constellations of GNSS. & Sec.~\ref{sec:visual_odometry} \\
    
    LPIPS & \textbf{Learned Perceptual Image Patch Similarity.} A perceptual image distance computed from deep visual features. & Sec.~\ref{sec:mapping_slam} \\
    
    PSNR & \textbf{Peak Signal-to-Noise Ratio.} An image quality measure derived from mean squared reconstruction error. & Sec.~\ref{sec:mapping_slam} \\
    
    SSIM & \textbf{Structural Similarity Index.} An image similarity measure based on local luminance, contrast, and structure. & Sec.~\ref{sec:mapping_slam} \\
    
    VAE & \textbf{Variational Autoencoder.} A generative model that learns a probabilistic latent representation of data. & Sec.~\ref{sec:vl_grounding} \\
    
    IK & \textbf{Inverse Kinematics.} Finding joint configurations that achieve a desired pose of an end effector or body part. & Sec.~\ref{sec:p3_s7_1} \\
    
    TAMP & \textbf{Task and Motion Planning.} Joint planning of discrete task actions and feasible continuous motions. & Sec.~\ref{subsec:p3_s7_1_4} \\
    
    MPC & \textbf{Model Predictive Control.} Control that repeatedly optimizes a finite-horizon plan and executes its first action. & Sec.~\ref{subsec:p3_s7_1_4} \\
    
    GNNs & \textbf{Graph Neural Networks.} Neural networks that learn representations by aggregating information over graph connections. & Sec.~\ref{subsec:p3_s7_2_1} \\
    
    POMDP & \textbf{Partially Observable Markov Decision Process.} A sequential decision model in which the agent observes the environment state only indirectly. & Sec.~\ref{subsec:p3_s7_2_5} \\
    
    RAG & \textbf{Retrieval-Augmented Generation.} Generation conditioned on information retrieved from an external collection or memory. & Sec.~\ref{ssec:memory_spatial_reasoning} \\
    
    ToM & \textbf{Theory of Mind.} Reasoning about the beliefs, intentions, and knowledge of other agents. & Sec.~\ref{subsec:p3_s7_3_7} \\
    
    FVD & \textbf{Fr\'echet Video Distance.} A distance between feature distributions of generated and reference videos. & Sec.~\ref{subsec:p3_s7_3_8} \\
    
    FID & \textbf{Fr\'echet Inception Distance.} A distance between Inception feature distributions of generated and reference images. & Sec.~\ref{subsec:p3_s7_3_8} \\
    
    GANs & \textbf{Generative Adversarial Networks.} Generative models trained through competition between a generator and a discriminator. & Sec.~\ref{subsec:p3_s8_1_1} \\
    
    DDIM & \textbf{Denoising Diffusion Implicit Model.} A diffusion sampling formulation that supports deterministic sampling and fewer sampling steps. & Sec.~\ref{subsec:p3_s8_1_3} \\
    
    SDS & \textbf{Score Distillation Sampling.} Optimization of a representation using guidance from a pretrained diffusion model. & Sec.~\ref{subsec:p3_s8_2_1} \\
    
    GNSS & \textbf{Global Navigation Satellite System.} A satellite system providing positioning, navigation, and timing information. & Sec.~\ref{subsec:p4-sec9-2-pivotal-task-2-localization-and-mapping} \\
    
    BIM & \textbf{Building Information Modeling.} Creation and management of digital models describing a building's physical and functional properties. & Sec.~\ref{sec:part-4-section-14-architecture-engineering-and-construction} \\
\end{longtable}

\part{Introduction and Background}
% \section*{\centering \bf Part I: Introduction and Background}
% \addcontentsline{toc}{section}{Part I: Introduction and Background}





\begin{figure*}[h!]
    \centering
    % \includegraphics[width=\linewidth]{figures/part1.png}
    \includegraphics[width=\linewidth]{figures/part1_opener_v2.pdf}
    
    \caption{Part I in one view: (1) why now, as scattered data condense into a grounded world; (2) how it evolved, one room in each of five eras; (3) what it is, a loop of perceiving, reasoning, predicting, and acting.}
    \label{fig:part1-illustration}
\end{figure*}



For decades, AI has advanced by mastering well-bounded tasks: recognizing patterns, answering questions, generating outputs from static inputs.
Intelligence in the natural world unfolds in space, shaped by structure, constraint, motion, and interaction.
As AI systems move from passive prediction to active participation in physical and simulated environments, models trained only on tokens, pixels, or isolated observations struggle to stay coherent across space, time, and action, because they have no explicit notion of where they are, how the world is arranged, or how their interventions change what comes next.
Spatial intelligence supplies that notion, and robust, adaptive, and embodied intelligence depends on it.
It anchors intelligence in world-centered representations that encode geometry, topology, semantics, dynamics, and affordances, and that let a system perceive beyond its current view, reason about regions it has not seen, predict how a configuration will evolve, and act on it with purpose.
Navigation, manipulation, simulation, and long-horizon planning all draw on this one substrate.

This part lays the conceptual ground for treating spatial intelligence as a paradigm.
Chapter~\ref{chapter:intro} explains why it is rising now, through four converging shifts in what AI is expected to do, how it learns, the modalities it handles, and the disciplines that study it, and reads the field's history as five eras, each defined by what it assumed space to be.
Chapter~\ref{chapter:definition} reconciles the views of computer vision, robotics, geographic information science, psychology, cognitive science, and linguistics in a working definition, an AI system and its environment coupled through a world-centered representation, and gives it a light formalization.
Chapter~\ref{chapter:structure} sets out the research pillars that organize the rest of the handbook.

\newpage




\chapter{Spatial Intelligence: A Rising AI Paradigm}
\label{chapter:intro}






% \lettrine[lines=3]{\initfamily\textcolor{darkblue}{A}}{I} 
\FramedInitial{A}{I} 
has entered a stage where scale and capability alone are no longer sufficient.
Despite remarkable advances in perception, language, and generation, modern AI systems continue to struggle with a fundamental challenge: understanding, representing, and acting within the structured physical world.
As intelligent systems increasingly move beyond static inputs and isolated tasks, the question shifts from how well a model predicts tokens or labels to how it forms internal models of space, time, action, and causality.
This shift marks the emergence of \emph{spatial intelligence} as a central organizing paradigm in artificial intelligence.


Spatial intelligence concerns the ability of an intelligent system to build and manipulate internal representations of the world that support perception, reasoning, prediction, and embodied interaction.
It is inherently world-centered, grounded in geometry, motion, and physical constraints, and inseparable from action and experience.
Rather than treating space as an auxiliary signal or downstream output, spatial intelligence positions it as a core substrate through which intelligence is structured.
This perspective cuts across traditional boundaries, unifying problems in computer vision, robotics, autonomous systems, simulation, cognitive science, spatial computing, AR/VR, and geospatial analysis under a shared representational question: how should an intelligent system model the world it inhabits?


This chapter establishes spatial intelligence as a rising AI paradigm by addressing two foundational questions.
Section~\ref{sec:why-spatial-intelligence} examines why spatial intelligence has become critical \emph{now}, tracing its emergence to converging shifts in AI paradigms, methodologies, modalities, and disciplines.
Section~\ref{sec:historical-evolution} then situates this paradigm historically, organizing decades of fragmented progress into a coherent sequence of representational eras.
Together, these perspectives provide the conceptual grounding for the rest of the handbook, framing spatial intelligence not as a collection of tasks or techniques, but as a unifying foundation for world-centered artificial intelligence.









\section{Why Spatial Intelligence? Why Now?}
\label{sec:why-spatial-intelligence}

% \section{From Data-Centric AI to World-Centric AI}


% \lettrine[lines=3]{\initfamily\textcolor{darkblue}{A}}{s} 
\FramedInitial{A}{s}
AI continues to advance, a growing consensus has emerged that data alone is not enough; true intelligence must be grounded in the physical world \cite{park2024model,zhang2025world,long2025survey}. 
This realization marks a fundamental shift from purely \emph{data-centric} paradigms, centered around language or perception, to \emph{world-centric} systems capable of modeling, reasoning, and acting within spatially structured environments. 
In this section, we explore the motivations behind this transition, drawing from the AI paradigm, methodology, and modality to disciplines.




\begin{figure}
    \centering
    % \resizebox{0.92\linewidth}{!}{% Four converging shifts driving spatial intelligence (TikZ, replaces si-shift-1.pdf)
\begin{tikzpicture}[
    shiftpanel/.style={draw, rounded corners=4pt, align=left, text width=6.6cm, inner sep=9pt, line width=0.5pt},
    paneltitle/.style={font=\small\bfseries},
    fromto/.style={font=\small},
    panledesc/.style={font=\footnotesize, color=black!70},
    arr/.style={-{Stealth[length=2.4mm]}, line width=0.9pt, black!65},
]
\node[shiftpanel, fill=ai_blue!18, draw=ai_blue!80!black] (p1) at (0,0) {
    {\small\bfseries AI Paradigm: Task-Specific $\rightarrow$ World-Centric}\\[4pt]
    {\small Narrow, self-contained problems \; $\Longrightarrow$ \; foundation models with global context}\\[3pt]
    {\footnotesize\color{black!70} The unit of competence moves from the single task to the modeled world.}
};
\node[shiftpanel, fill=DefGreen, draw=forestgreen!80, right=0.55cm of p1] (p2) {
    {\small\bfseries Methodology: Static $\rightarrow$ Interactive}\\[4pt]
    {\small Frozen models on offline corpora \; $\Longrightarrow$ \; agents that learn continually from experience}\\[3pt]
    {\footnotesize\color{black!70} Train-freeze-infer gives way to closed-loop learning through interaction.}
};
\node[shiftpanel, fill=customyellow!65, draw=orange!75!black, below=0.5cm of p1] (p3) {
    {\small\bfseries Modality: 2D $\rightarrow$ 3D/4D, Any-to-Any}\\[4pt]
    {\small Single-frame imagery \; $\Longrightarrow$ \; spatiotemporal understanding across sensory channels}\\[3pt]
    {\footnotesize\color{black!70} Perception expands from flat snapshots to dynamic, embodied space-time.}
};
\node[shiftpanel, fill=customred!45, draw=steelred, below=0.5cm of p2] (p4) {
    {\small\bfseries Disciplines: Silos $\rightarrow$ Convergence}\\[4pt]
    {\small CV, NLP, robotics, GIS in isolation \; $\Longrightarrow$ \; a shared world-modeling agenda}\\[3pt]
    {\footnotesize\color{black!70} Fragmented communities align around common spatial substrates.}
};
% converging arrows to a central label
\node[draw, rounded corners=4pt, fill=white, align=center, font=\small\bfseries, inner sep=7pt, line width=0.6pt] (center) at ($(p1.south east)!0.5!(p4.north west)$) {Spatial\\Intelligence};
\draw[arr] (p1.south east) ++(-0.35,0.08) -- (center.north west);
\draw[arr] (p2.south west) ++(0.35,0.08) -- (center.north east);
\draw[arr] (p3.north east) ++(-0.35,-0.08) -- (center.south west);
\draw[arr] (p4.north west) ++(0.35,-0.08) -- (center.south east);
\end{tikzpicture}
}
    % \includegraphics[width=0.99\linewidth]{figures/part1/si-shift-2.png}
    \includegraphics[width=\linewidth]{figures/part1/si_shifts_v2.pdf}
    \caption{Spatial intelligence emerges from converging shifts in AI paradigms, methodologies, modalities, and disciplines, collectively pushing AI toward world modeling and embodied reasoning.}
    \label{fig:p1_s1_1_si_covering_shifts}
\end{figure}


% \subsection{From Task-Specific to World-Centric: A Paradigm Shift}
\subsection{AI Paradigm Shift}
\label{subsec:ai-paradigm-shift}

AI is undergoing a foundational paradigm shift, \ie, from task-specific models toward world-centric general intelligence. 
Traditional AI systems were designed to solve narrowly defined problems, such as image classification, text generation, or speech recognition, each trained and optimized within a well-constrained input-output space. 
However, the emergence of foundation models and large language models (LLMs) has reshaped expectations: AI is now envisioned to generalize across domains, modalities, and tasks in open-world settings. 
This ambition hinges on a deeper capability, the ability to understand and model the world itself, often described as building a ``world model'' \cite{ha2018world}.


\emph{Spatial Intelligence} emerges as a central pillar enabling this capability.
It encompasses the perception and representation of objects, scenes, spatial layouts, motion, and physical interactions, and it serves as a carrier of world-generalizable knowledge that supports reasoning and transfer across tasks and modalities.
Fei-Fei Li has made this stake explicit: ``\emph{Our dreams of truly intelligent machines will not be complete without spatial intelligence}'' \cite{Feifei2025spatial}.
On her account, any AI system that aims to serve in physical environments must possess a rich understanding of spatial structure and dynamics.



Moreover, this transformation also signals a shift in the boundary of context: traditional AI focused on self-contained intelligence, where the model reasons primarily over its own input context.
Today's systems, in contrast, must be understood as \emph{individuals embedded in the world}\footnote{In Context Engineering \cite{hua2025context}, the authors frame the former paradigm as \emph{Context as Instruction}, and cast the latter as \emph{Context as World}.}, agents whose perception and action are grounded in broader spatial, temporal, and social environments.
Beyond the technical evolution, this amounts to a philosophical deepening of AI's cognitive scope.
Li captures the underlying continuity in evolutionary terms: spatial intelligence is, in her words, what ``\emph{links perception with action}'', the capacity that turned seeing into surviving and, eventually, into understanding and doing \cite{Feifei2024ted}.



Without accurate spatial modeling of the world, therefore, there can be no robust reasoning, decision-making, or interaction.
A wide range of emerging tasks, including robotic manipulation, embodied navigation, video understanding, and virtual human interaction, demand spatial grounding as a prerequisite, which makes spatial intelligence a load-bearing requirement for any AI system expected to act in, adapt to, and help within the real world.







\subsection{AI Methodological Shift}
\label{subsec:ai-methodological-shift}

Historically, from the era of classical machine learning \cite{CoughlanY99} to the deep learning revolution \cite{goodfellow2016deep}, and now to the LLM wave \cite{achiam2023gpt,vicuna2023}, progress in AI has largely followed a data-centric trajectory: larger models trained on increasingly vast static corpora.
While this paradigm has delivered remarkable achievements, it also imposes an inherent limitation, in that AI systems are still framed as static input--output function approximators rather than as agents that can perceive, reason, and act within the real world.
The recent surge of LLMs, such as GPT \cite{achiam2023gpt}, PaLM \cite{chowdhery2023palm}, and LLaMA \cite{touvron2023llama}, has only amplified this trajectory.
Although these models exhibit stunning linguistic fluency, symbolic manipulation, and emergent generalization capabilities \cite{yamada2023evaluating,zhang2025mllms}, they remain intrinsically disembodied and spatially ungrounded.
They can generate poetic descriptions of a cup on the table, yet they cannot perceive the cup, interact with it, or build a physical model to explain its behavior.
This exposes the core limitation of the current paradigm: \emph{LLMs might operate on surface-level linguistic statistics without accessing the structured constraints, causal mechanisms, and physical affordances of the real world}.


Fei-Fei Li has repeatedly emphasized this disconnect.
Linguistic competence, she argues, must not be mistaken for understanding: for all their capability, today's language models \emph{``remain wordsmiths in the dark; eloquent but inexperienced, knowledgeable but ungrounded''} \cite{Feifei2025spatial}.
Without spatial intelligence, encompassing physics, motion, and the transformative nature of actions, AI systems cannot develop a grounded, interactive model of reality.
Her insight exposes a deeper methodological gap: true intelligence must ground its internal representations in embodied space, action, and physical experience, and text alone cannot supply that grounding.



Yann LeCun extends this critique at the architectural level, famously dismissing LLMs as ``\emph{an off-ramp, a distraction, a dead end}'' on the path toward human-level intelligence.
In his view, autoregressive language models do not, and cannot, construct models of how the world works; they continue text through pattern completion rather than through prediction grounded in the dynamics of an external world \cite{lecun2022path}.
His proposed alternative, an architecture organized around a learned predictive world model, points squarely toward systems capable of predictive modeling, causal reasoning, and dynamic world interaction.



Rich Sutton, from a reinforcement learning (RL) perspective, presses the same point on the learning process itself.
Sutton's ``bitter lesson'' famously argued that general methods which scale with computation win out over hand-engineered knowledge \cite{sutton2019bitter}; notably, he does not regard today's LLMs as the vindication of that argument.
In his recent view, systems trained once on static human data and then frozen are cut off from the very thing that drives intelligence, namely continual learning from one's own experience \cite{silver2025era}.
Silver and Sutton argue that the field must move from the ``era of human data'' into an ``era of experience'', in which agents learn predominantly from acting in the world and observing consequences.
Learning from static text captures correlations in what humans have already written; understanding the world requires action, consequence prediction, and causal discovery.




Taken together, these critiques, arising independently from visual cognition, neural architectures, and reinforcement learning, signal a profound methodological shift in progress: AI must evolve from a purely LLM-centric paradigm to one of world-integrated intelligence.
It must move beyond static, offline training toward continual, interactive, and adaptive learning, and it must outgrow shallow pattern matching in text to engage with the physical, spatial, and causal structures that govern real environments.
More than an algorithmic upgrade, this demands a redefinition of what intelligence means: agents capable of perceiving the world, predicting its dynamics, interacting within it, and learning continuously from such experiences \cite{bengio2017consciousness,yao2022react,hao2023reasoning}.








\subsection{Modality Shift}
\label{subsec:modality-shift}

Human intelligence is inherently multimodal. 
Any meaningful pursuit of AGI must therefore embrace multimodality as a foundational principle. 
The trajectory of multimodal AI is undergoing a significant shift, from early-stage \emph{modality separation} toward deep \emph{modality integration}. 
In earlier AI systems, perception and modeling were often confined to a single modality \cite{he2016deep}, such as text-only language models, speech recognition engines, or visual classification systems. 
However, with the rise of foundation models, the research paradigm is shifting towards hybrid-modality \cite{lu2019vilbert,tan2019lxmert} and even any-to-any cross-modal understanding and generation \cite{wu2024next}. 
That is, both input and output can now be flexibly composed across arbitrary modalities (\eg, audio-based question answering, video synthesis, image-conditioned dialogue) \cite{yan2021videogpt,yun2021pano,zhang2024dialogstudio}.
Simultaneously, the very nature of modality is expanding in dimensionality and diversity. 
Traditional sensory modalities, \eg, vision, language, and audition, represent only a subset of how humans perceive the world. 
With the rapid advancement of sensor technology, future AI systems are expected to be equipped with a broader range of high-dimensional and multi-channel sensory inputs \cite{li2020review,liu2023multimodal,park2023cross,salih2024odors}, including spectrograms, millimeter-wave radar, infrared thermography, tactile feedback, gustatory signals, and olfactory landscapes. 
These emerging modalities are especially vital in domains such as medicine \cite{topol2019high}, robotics \cite{Levine2016DeepGrasp}, environmental monitoring \cite{alotaibi2024artificial}, and AR \cite{nizam2018review}, offering a more human-aligned perceptual basis for multimodal AI systems. 
This multimodal expansion necessitates a unified modeling capability that operates at the level of holistic world understanding.




Among all modalities, vision remains the most central to human perception. 
One of the most critical trends in modality development is the transition from static imagery to video, and from 2D representations to richer 3D and even 4D (spatiotemporal) structures \cite{saxena2008make3d,peng2023openscene}. 
Traditional image-based recognition tasks operate on single-frame static inputs. 
However, real-world environments demand continuous understanding, prediction, and generation across space and time. 
This drives the growing importance of tasks such as 3D reconstruction, temporal modeling, trajectory prediction, video-based question answering, and dynamic world state tracking \cite{ciaparrone2020deep,azuma2022scanqa}. 
The increasing richness of modalities goes well beyond a technological extension; it is a necessary foundation for building AGI.
To rival or surpass human intelligence, AGI must be equipped with a comparably rich sensorium, enabling it to perceive, reason, and act effectively in dynamic, multimodal environments. 
Thus, the integration and extension of modalities are not ends in themselves, but serve the larger ambition of world modeling, \ie, the central focus of this work.




This point has been underscored by Fei-Fei Li in her discussion of spatial intelligence. Tracing the evolutionary arc from the first light-sensing organisms to primates, she observes that \emph{``nature has created this virtuous cycle of seeing and doing powered by spatial intelligence''} \cite{Feifei2024ted}.
Such capability acts as a bridge between cognition and the physical world, a necessary scaffold for grounding AI in embodied reality; without it, meaningful interaction with the real world is out of reach.
Where traditional intelligence centers on language or symbolic manipulation, spatial intelligence is tightly coupled with physical structure, three-dimensional relationships, and dynamic evolution, making it the cornerstone of embodied, interactive intelligence.
From the perspective of modality development, then, spatial intelligence marks the turning point where multimodal intelligence evolves from surface-level perception to deep, structured understanding, well beyond a mere byproduct of increasing sensory complexity.
The focus of AI research must accordingly shift from modeling and aligning modalities to modeling the structured world that those modalities collectively describe, the path toward a closed-loop paradigm from perception to world modeling, and from world modeling to causal reasoning.







\subsection{Disciplinary Convergence}
\label{subsec:disciplinary-convergence}

Modern AI research is undergoing a profound disciplinary evolution, centered around \textit{modality enrichment} and \textit{spatiophysical coupling}. 
This paradigm shift reflects a convergence of traditionally siloed domains, spanning Computer Vision (CV), Natural Language Processing (NLP), Audio Processing, Robotics, Geographic Information Systems (GIS), and Urban Simulation, toward a unified framework for modeling intelligent agents that are physically embedded in and interact with the real world.


In the vision domain, the research focus has gradually shifted from static image processing to dynamic scene understanding. 
These developments provide a perceptual foundation for large-scale world modeling. 
Language modalities have also evolved beyond static language modeling (\eg, BERT~\cite{devlin2019bert}, GPT~\cite{brown2020language}), toward multimodal semantic fusion tasks such as Visual Question Answering (VQA)~\cite{antol2015vqa}, Vision-and-Language Navigation (VLN)~\cite{gu2022vision}, and embodied dialogue, which have become critical interfaces in the expansion of large language models into the multimodal regime. 
The audio modality has similarly moved beyond traditional speech recognition to include audio event detection, ambient sound understanding, and even 3D auditory modeling, all of which support spatial cognition and situational awareness in real-world interactive environments~\cite{chen2020soundspaces,chen2022soundspaces}.


Meanwhile, the rapid development of robotics and embodied AI is pushing the evolution of AI systems from ``perception-understanding'' pipelines to closed-loop ``perception-action'' systems.
AI is outgrowing its existence as software on a computer, becoming embodied, real-time, and spatially grounded.
In domains such as autonomous driving, robotic manipulation, and wearable AI systems, tasks increasingly require deep spatial reasoning capabilities for layout perception, path planning, and object manipulation~\cite{howard2019evolving,gupta2021embodied}. 
These tasks demand integration not only of vision and language capabilities but also of knowledge from control theory, motion planning, and state estimation, marking a shift toward truly multidisciplinary AI systems.



The expansion of modalities is also reshaping traditional disciplines through interdisciplinary integration. 
In GIS and urban computing, AI models are increasingly combined with spatial coordinates, map representations, and temporal population data, enabling applications in smart city planning, trajectory prediction, and disaster simulation. 
In Earth sciences and remote sensing, AI is employed for multimodal surface analysis, climate modeling, and environmental monitoring, pushing the frontier toward 3D Earth modeling and dynamic ecological understanding~\cite{janga2023review}. 
In architectural design and digital twin systems, AI supports full lifecycle 3D spatial modeling, simulation, and predictive analytics for complex structures and spaces.



Cognitive science and neuro-inspired AI (NeuroAI) are playing an increasingly pivotal role in this trend. 
Foundational research on spatial intelligence in human cognition, such as hippocampus-mediated cognitive maps, spatial memory, and embodied cognition theory, is gradually informing the design of AI spatial perception systems~\cite{o1978hippocampus,barsalou2008grounded}. 
As Fei-Fei Li puts it, spatial intelligence is \emph{``the scaffolding upon which our cognition is built''}, the substrate that anchors intelligence to the physical world~\cite{Feifei2025spatial}.
Thus, the path toward AGI must transcend abstract symbolic computation and advance toward grounded-in-reality agent intelligence.
This trend reshapes the research trajectory of AI and, at the same time, redefines the underlying infrastructure of general intelligence itself, with spatial intelligence standing as a central pillar of the emerging world-centered era.
























\section{Historical Evolution and Milestones}
\label{sec:historical-evolution}

% \scott{TODO: add enough necessary formulations, equations.}

% style-1
% \lettrine[lines=3]{\initfamily\textcolor{darkblue}{T}}{he}
% style-2
% \lettrine[
%   lines=3,
%   findent=0.3em,
%   nindent=0.6em
% ]{\colorbox{darkblue}{\textcolor{white}{\bfseries T}}}{he}
% style-3
% \renewcommand{\LettrineFontHook}
%   {\color{darkblue}\rmfamily\bfseries}
% \lettrine[
%   lines=3,
%   loversize=0.05,
%   lhang=0.05,
%   findent=0.15em,
%   nindent=0.45em
% ]{T}{he}
% style-4
\FramedInitial{T}{he} 
history of spatial intelligence is best understood as a sequence of shifts in \emph{how space is represented} and \emph{what that representation is expected to support}. 
Over decades, the field has repeatedly shifted among explicit structure and learned abstraction, offline reconstruction and online interaction, and local geometry and world-level state.
To make this evolution concrete, we organize the landscape into five eras, each characterized by a dominant representation paradigm and its associated methodological commitments: from early symbolic and cognitive notions of spatial structure, to probabilistic geometric estimation, to learning-based spatial perception, to large foundation representations that are transferable and cross-modal, and finally to world-centered spatial intelligence where representations must be executable, updatable, and tightly coupled with action in dynamic environments \cite{kuipers1978modeling,Thrun2005Probabilistic,eigen2014depth,driess2023palm,Feifei2025spatial}.
Figure~\ref{fig:p1_s1_2_si_evolution} shows how each era would represent the same scene, Table~\ref{tab:p1_s1_2_eras} condenses what each era took space to be and what it left unsolved, and Figure~\ref{fig:p1_s1_2_milestones} places about a hundred representative works from six fields on a common time axis.





\begin{figure}[!t]
\centering
% \resizebox{0.99\linewidth}{!}{% Redesigned representation-centric timeline.
% Editable source: figures/part1/sources/si_evolution.py
% Original TikZ preserved in sources/si_evolution_assets/original_tikz.tex.
\includegraphics{figures/part1/si_evolution_tikz.pdf}
}
% \includegraphics[width=0.99\linewidth]{figures/part1/si_evolution.png}
\includegraphics[width=\linewidth]{figures/part1/si_evolution_v2.pdf}
\vspace{1mm}
\caption{The historical evolution of spatial intelligence across five representation-centric eras, from explicit symbolic structure, through estimated geometry and learned features, to executable world states. The visual examples highlight the dominant representation and its role in each era. Dates indicate approximate shifts in emphasis, not exclusive boundaries; earlier paradigms remain in use.}
\label{fig:p1_s1_2_si_evolution}
\end{figure}

\begin{table}[!t]
\centering
\fontsize{9}{11}\selectfont
\setlength{\tabcolsep}{1.0mm}
\caption{The five eras of spatial intelligence compared along six dimensions: how each treated space, its core representations, what the representation was for, its characteristic methods, the fields that drove it, and what it left unsolved for the next era.}
\label{tab:p1_s1_2_eras}
\begin{tabular}{L{0.145\textwidth} L{0.159\textwidth} L{0.159\textwidth} L{0.159\textwidth} L{0.159\textwidth} L{0.159\textwidth}}
\toprule
 & \textbf{I.\ Symbolic and cognitive} \newline {\scriptsize 1950s--1980s} & \textbf{II.\ Geometric and probabilistic} \newline {\scriptsize 1980s--2012} & \textbf{III.\ Learning-based} \newline {\scriptsize 2012--2019} & \textbf{IV.\ Foundation models} \newline {\scriptsize 2019--2023} & \textbf{V.\ World-centered} \newline {\scriptsize 2023--} \\
\midrule
\textbf{Space treated as} & discrete symbols and relations & metric geometry under uncertainty & learned, task-specific features & shared latent priors & an executable, persistent world state \\
\textbf{Core representations} & cognitive maps, topological graphs, predicates & occupancy grids, landmark maps, pose graphs & feature maps, depth and flow maps, point and voxel features & pretrained embeddings, cross-modal alignment, neural fields & world models, action-conditioned dynamics, spatial memory \\
\textbf{What it is for} & planning and inference by symbol manipulation & a consistent pose and map from noisy sensing & perception objectives: detection, depth, motion & transfer across tasks, modalities, and environments & long-horizon prediction, planning, and interaction \\
\textbf{Characteristic methods} & heuristic search, rule-based planning, logic & Bayesian filtering, nonlinear least squares, multi-view geometry & supervised and self-supervised deep networks & large-scale pretraining, contrastive alignment, instruction tuning & action-conditioned prediction, memory, closed-loop learning \\
\textbf{Main fields} & cognitive science, classical AI, early robotics & robotics, photogrammetry and vision, GIS & computer vision, autonomous driving & multimodal AI, neural rendering, embodied AI & embodied AI, simulation and digital twins, AR/VR \\
\textbf{Left unsolved} & meaning asserted by rules, not inferred from noisy sensing & features and pipelines engineered by hand & one model per task; no persistent world model & no persistent state that action updates & open: long-horizon consistency under motion and partial observability \\
\bottomrule
\end{tabular}
\end{table}














\subsection{Era 1: Symbolic and Cognitive Spatial Computing (1950s--1980s)}
\label{subsec:era-symbolic-cognitive}

\paragraph{Symbolic Space and Cognitive Structure as the Starting Point.}
This era marks the cognitive and symbolic origin of spatial intelligence. 
Under severe limits in computation and data, space was primarily treated as a structured object that could be explicitly described, reasoned about, and manipulated.
Accordingly, researchers favored symbolic rules, relational descriptions, and graph-like structures to encode environments, viewing spatial representation as a verifiable ``spatial description language'' that enables planning and logical inference \cite{kuipers1978modeling,kuipers2000spatial}.
The lasting contribution of this stage is not empirical performance, but two foundational commitments.
First, goal-directed behavior requires an internal world representation.
Second, spatial representations are inherently compositional, consisting of objects, relations, and connectivity, which naturally supports rule-based spatial reasoning \cite{sun1995robust}.


From a cognitive science perspective, this period also established the idea that spatial ability is not mere perception but an internal, operable mental model.
Concepts such as cognitive maps and spatial memory motivated the view that an intelligent system should maintain structured internal states that support navigation, mental imagery, and transformation of spatial relations \cite{eichenbaum1999hippocampus,allen2004human}.
In parallel, early robotics provided an engineering prototype of this view: the classic sense-plan-act pipeline connected symbolic spatial models with planners and execution, demonstrating how an agent can translate limited sensing into purposeful behavior within an environment \cite{nilsson1984shakey,Arkin1998BehaviorBased}.

\paragraph{Representative Milestones in Techniques and Representations.}
A widely cited conceptual root is Tolman's proposal of ``cognitive maps'' (1948), which framed navigation as relying on internal map-like structures rather than stimulus-response associations \cite{tolman1948cognitive}.
In classical AI, the Dartmouth Workshop (1956) crystallized the symbolic paradigm and encouraged formal, rule-based representations of intelligence, providing intellectual soil for early spatial reasoning systems \cite{mccarthy2006proposal,krauss2024artificial}.
Planning and search then became the main computational language for spatial decision making: A* heuristic search (1968)~\cite{Hart1968Astar} offered a principled graph search framework for shortest-path style navigation problems, and STRIPS (1971)~\cite{Fikes1971STRIPS} formalized actions as operators over states, making spatial tasks expressible as sequences of symbolic state transitions \cite{nilsson2014principles,pearl1983heuristics}.



In robotics, Shakey (late 1960s to early 1970s) became a milestone system that integrated perception, symbolic spatial modeling, planning, and acting in a closed pipeline \cite{nilsson1984shakey}.
Although the sensory and actuation capabilities were primitive by modern standards, the system concretely instantiated the era's thesis that spatial intelligence can be built by maintaining symbolic representations and using them to plan actions.
In neuroscience and cognitive modeling, O'Keefe and Nadel's \emph{The Hippocampus as a Cognitive Map} (1978) further systematized the cognitive map account, reinforcing the interpretation that internal spatial representations are functional substrates for memory and navigation \cite{o1978hippocampus}.



\begin{figure}[p]
\centering
\includegraphics[width=\linewidth]{figures/web/si_milestones_tikz.pdf}
\caption{Representative works in spatial intelligence by field and year, set against the five eras of Table~\ref{tab:p1_s1_2_eras} (shaded bands, I--V). Each marker sits at the year of the original work, and the bracketed number cites it; where several works share a year, labels are spread and joined to their markers by short leaders. The time axis is compressed before 1980 and expanded after 2012.}
\label{fig:p1_s1_2_milestones}
\end{figure}

Across these milestones, several representation families repeatedly appeared and stabilized: symbolic spatial models and rule-based spatial descriptions; topological graphs capturing connectivity and place relations; qualitative spatial relations such as containment, adjacency, and support; landmark-based representations for navigation; and grid-world abstractions that later became a convenient substrate for sequential decision making and, eventually, reinforcement learning (RL) formulations \cite{arulkumaran2017deep}.
Together, they define the representation-centric signature of Era 1: spatial intelligence begins with explicit structure, grounded in cognition-inspired internal maps and implemented through symbolic reasoning and sense-plan-act-style control.

\subsection{Era 2: Geometric and Probabilistic Spatial Modeling (1980s--2012)}
\label{subsec:era-geometric-probabilistic}


\paragraph{From Symbolic Descriptions to Metric Geometry under Uncertainty.}
With the maturation of sensors (laser range finders, wheel odometry, early vision), computing infrastructure, and probabilistic estimation theory, spatial representation shifted away from purely symbolic description toward mathematical modeling of continuous geometry with uncertainty.
Space was increasingly treated as a measurable physical entity whose state must be inferred, rather than a rule-based structure whose meaning is asserted.
Consequently, the central question of this era shifted from describing the world in symbols to maintaining spatial consistency under noise, drift, and partial observability.
This transition also re-centered the representation itself: a map is not a static record, but a set of uncertain state variables, typically the joint configuration of pose and environment, which must be updated and refined through Bayesian estimation and optimization \cite{Thrun2005Probabilistic,dellaert2006square}.

In this view, representation is inseparable from the inference machinery that maintains it.
Coordinate frames and metric spaces became the common language that connects sensing, estimation, and control, while probabilistic modeling provided a principled way to fuse heterogeneous measurements and quantify uncertainty.
As a result, a representation-centric principle was consolidated: spatial intelligence requires not only a geometric model of the environment, but also an explicit account of how that model is updated over time as new evidence arrives \cite{smith1990estimating,bar2001estimation,thrun2002robotic}.
This era also laid much of the groundwork for long-term autonomy by formalizing how an agent can localize itself, build maps, and navigate using a consistent geometric state.

In parallel, macro-scale spatial disciplines pushed representation standardization beyond local robotics maps.
Surveying, geodesy, cartography, remote sensing, and GIScience strengthened coordinate reference systems, map projections, topological relations, and spatial databases, making large-scale spatial representations structured, queryable, and interoperable at city-to-Earth scale \cite{goodchild1992geographical}.
At the same time, computer vision developed a rigorous multi-view geometry toolkit that defined how 3D structure and camera motion can be recovered from images, reinforcing the idea that geometry is an estimable latent state rather than an explicit symbolic description \cite{hartley2003multiple}.





\paragraph{Representative Techniques and Milestones.}
A representative milestone in this era is the emergence of probabilistic state estimation as the dominant language for maintaining spatial consistency under noise and partial observability.
Kalman filtering and its nonlinear variants (EKF and related formulations) became widely adopted in robotics and autonomous systems during the 1980s and 1990s, enabling spatial quantities such as pose, velocity, and sensor biases to be represented as random variables with explicit uncertainty propagation \cite{welch1995kalman}.
This shift is representation-centric: it turns ``space'' into a belief state that must be continuously updated, rather than a deterministic symbolic description.
Building on this estimation view, occupancy grid mapping in the mid-1980s provided a canonical probabilistic representation that directly fuses range sensing into a discretized spatial field of free versus occupied space \cite{Moravec1985SonarMaps,Thrun2005Probabilistic}.
Although conceptually simple, occupancy grids became a durable interface between sensing and downstream planning, because they encode geometry together with uncertainty in a form amenable to incremental updates.
Related metric map abstractions, coordinate frames, and sensor models matured in parallel, standardizing the idea that maps are measurable geometric objects anchored to physical reference systems \cite{LaValle2006Planning,hornung2013octomap}.


The SLAM problem then consolidated these ideas into a unified formulation where the agent's pose and map are jointly inferred~\cite{dissanayake2001solution}.
EKF-SLAM, whose stochastic-map foundations were laid in the late 1980s \cite{smith1990estimating} and refined through the 1990s into the early 2000s \cite{bailey2006consistency}, established the coupled pose-landmark state as a central representation, while particle-filter-based SLAM, including FastSLAM (2002)~\cite{montemerlo2002fastslam}, demonstrated that sampling and factorization could scale inference by exploiting conditional independencies.
Beyond their algorithmic contributions, these lines of work stand as milestones in the operationalization of spatial representation: the map becomes a latent variable tied to motion and observation processes, with consistency defined probabilistically.

As systems expanded in scale and duration, graph-based spatial representations rose to prominence in the early 2000s.
Pose graphs and factor graphs~\cite{lu1997globally} recast spatial estimation as a global constraint-satisfaction problem, enabling loop closure and large-scale consistency through nonlinear optimization.
GraphSLAM~\cite{thrun2006graph} and pose-graph optimization frameworks provided a modular representation where constraints from heterogeneous sensors and revisits could be integrated as edges in a unified objective \cite{grisetti2010tutorial}.
This was a conceptual milestone because it decoupled spatial consistency from strict temporal recursion and made map structure explicit as an optimization-ready graph.


In computer vision and photogrammetry, multi-view geometry matured into a systematic toolkit for recovering 3D structure from images.
Bundle adjustment, a core optimization primitive with roots in analog photogrammetry and consolidated for computer vision in the late 1990s, jointly refines camera poses and 3D points to enforce global geometric consistency across views \cite{triggs2000bundle}.
This family of methods made camera-centric coordinate frames, point clouds, and calibrated reconstructions standard spatial representations, and it established geometry and optimization, rather than semantics, as the primary organizing principles of spatial understanding in vision during this era.
Complementary to these estimation and reconstruction pipelines, geometric registration and fusion techniques advanced practical 3D mapping.
Iterative Closest Point (ICP) and its robust refinements enabled consistent aggregation of partial point measurements across viewpoints \cite{besl1992method,chetverikov2002trimmed}.
In parallel, early volumetric representations, including voxel grids and truncated signed distance function style formulations, began to encode surfaces and free space within unified volumetric fields \cite{lorensen1998marching,curless1996volumetric,newcombe2011kinectfusion,niessner2013real}.
While computationally constrained at the time, these representations introduced the field-based view of space that would later become influential for differentiable and implicit modeling.


Finally, at macro scales, GIS and surveying communities standardized spatial representations through coordinate reference systems, projections, geodesy, and spatial databases.
Digital elevation models, map products derived from remote sensing, and queryable geospatial layers enabled spatial reasoning at city, regional, and planetary scales with an emphasis on consistency, interoperability, and long-horizon stability \cite{goodchild1992geographical,li2025giscience}.
Although historically less coupled to embodied agents, these infrastructures are milestones for spatial intelligence because they demonstrate how spatial representation can function as a persistent, structured substrate supporting navigation, planning, and analysis beyond local robotic maps.








\subsection{Era 3: Learning-based Spatial Modeling (2012--2019)}
\label{subsec:era-learning-based}

\paragraph{From Engineered Geometry to Task-driven Learnable Spatial Representations.}
This era marks a fundamental shift in how spatial representations are constructed.
We date its onset to 2012, when AlexNet's ImageNet result made deep learning the dominant paradigm for visual perception \cite{krizhevsky2012imagenet}; the ImageNet dataset itself (2009) and early feature-learning work seeded the transition, but 2012 is when learned representations displaced hand-crafted ones as the default across spatial perception tasks.
With the rise of machine learning, and especially deep learning, spatial modeling moves away from pipelines dominated by hand-crafted geometry, explicit probabilistic inference, and manually designed features.
Instead, spatial structure begins to be encoded implicitly within neural networks through learned feature maps, latent embeddings, and task-optimized intermediate representations.
Space is no longer represented solely as a map, graph, or state vector defined by engineers, but as something that can be learned directly from data through end-to-end optimization~\cite{deng2009imagenet,eigen2014depth,mildenhall2021nerf}.




% \paragraph{Task-driven Learning and Discriminative Spatial Encoding.}
A defining characteristic of this period is that spatial representations are learned in a strongly task-driven manner.
Networks learn spatial cues to solve specific objectives such as object detection, semantic segmentation, depth estimation, visual odometry, or 3D recognition.
As a result, spatial knowledge is often embedded in discriminative feature spaces tailored to particular datasets and benchmarks.
While these representations dramatically improve performance on individual tasks, they rarely constitute a unified or reusable world model.
Spatial intelligence in this era remains fragmented, with representations optimized locally for perception subtasks rather than globally for cross-task reasoning or long-term world understanding~\cite{long2015fcn,ren2015faster,zhou2017unsupervised}.



% \paragraph{Parallel Advances Across Domains without Convergence.}
Another important aspect of this era is the parallel but largely uncoordinated progress across multiple domains.
In computer vision, learning-based approaches revolutionize image, video, and 3D perception.
In robotics, learning begins to augment classical pipelines for perception and control, but geometry-heavy systems remain dominant for reliability.
Remote sensing adopts deep models to learn spatial patterns from multispectral, hyperspectral, SAR, and LiDAR data at large scales.
Medical imaging leverages deep networks to learn volumetric spatial representations from CT and MRI scans.
Autonomous driving stacks integrate learned perception modules with geometric localization and mapping.
Despite shared reliance on learned spatial features, these communities operate with different benchmarks, representations, and evaluation criteria, reinforcing the task-specific nature of spatial modeling in this era~\cite{KITTI-GeigerLU12,ronneberger2015unet,Levine2016DeepGrasp}.




\paragraph{Representative Techniques and Milestones.}
A widely recognized milestone is the introduction of deep convolutional neural networks (CNNs) for visual perception, catalyzed by AlexNet in 2012, which establishes CNNs as the dominant backbone for spatially grounded perception tasks \cite{krizhevsky2012imagenet,rawat2017deep}.
This shift is quickly reinforced by region-based detection frameworks such as R-CNN and its successors from 2014 onward, which embed object-level spatial reasoning directly into learned feature hierarchies \cite{ren2015faster,he2017mask}.
ResNet in 2015 further deepens representational capacity and becomes a de facto foundation for most spatial perception pipelines \cite{he2016deep}.
In three-dimensional perception, PointNet in 2017 marks a milestone by introducing point clouds as a first-class input modality for deep networks \cite{qi2017pointnet}.
This work opens the door to learned spatial representations beyond images, influencing robotics, autonomous driving, and indoor scene understanding.
Concurrently, early voxel-based and sparse convolution methods begin to encode space volumetrically, albeit at limited resolution due to computational constraints \cite{tang2020searching}.




Learning-based depth estimation emerges as another key milestone.
Between 2014 and 2016, monocular and stereo depth prediction networks demonstrated that geometric cues such as disparity and scale can be inferred from data rather than explicit triangulation \cite{eigen2014depth}.
Around the same period, learning-based visual odometry~\cite{Wang2017DeepVO} and ego-motion estimation appeared~\cite{tateno2017cnn}, operating alongside classical SLAM systems.
This coexistence highlights a transitional moment where geometry-driven consistency and learning-driven representation evolve in parallel rather than replacing one another \cite{Wang2017DeepVO,tang2018ba}.
Self-supervised and unsupervised learning of spatial structure became influential from 2016 to 2018.
Methods that exploit photometric consistency, view synthesis, and temporal coherence enable networks to learn depth and motion without dense ground truth supervision \cite{zhou2017unsupervised,godard2019digging}.
These approaches strengthen the idea that spatial representations can be learned from raw sensory streams, reinforcing the data-driven paradigm while remaining largely task-bound.





Autonomous driving perception stacks introduce another representation shift toward bird's-eye-view and multi-view feature fusion.
By the late 2010s, learned BEV representations and cost volumes became central to multi-camera and LiDAR-based perception, enabling spatial reasoning aligned with downstream planning requirements \cite{liu2023bevfusion,yao2018mvsnet,kendall2017end}.
This reflects an increasing emphasis on representations that are geometrically meaningful yet learned.
Toward the end of this era, early attempts at neural scene representations and coordinate-based networks appear.
Works exploring implicit coordinate mappings and continuous scene embeddings between 2018 and 2019 lay conceptual groundwork for later developments in implicit fields and neural rendering \cite{park2019deepsdf,mescheder2019occupancy,sitzmann2019scene}.
Although these models are not yet general or scalable, they signal a growing interest in treating scenes, rather than tasks, as the primary object of representation.



Overall, Era 3 establishes learned spatial representations as the dominant paradigm.
However, these representations remain largely discriminative, task-specific, and fragmented across domains.
They provide powerful spatial priors for perception but stop short of forming unified, reusable world models, setting the stage for the emergence of large foundation spatial modeling in the subsequent era.
















\subsection{Era 4: Large Foundation Spatial Modeling (2019--2023)}
\label{subsec:era-foundation-models}



\paragraph{From Task-Specific Learning to Foundation Spatial Representations.}
The defining shift of this era is the rise of spatial representations trained at scale as foundation-level constructs, a change of kind rather than an incremental improvement on individual spatial tasks.
For the first time, spatial modeling is explicitly shaped through large-scale pretraining over diverse data sources, tasks, and modalities, rather than being optimized only within narrow task boundaries.
Spatial structure is no longer learned as an auxiliary signal for detection, depth, or reconstruction, but is instead distilled into shared latent representations through massive data exposure and objective diversity.
This marks a decisive representation-centric transition, where the primary goal becomes learning general-purpose spatial priors that can be transferred, adapted, and reused across downstream tasks with minimal supervision.



A key enabler of this transition is the combination of large-scale datasets, scalable architectures, and unified training objectives.
Vision Transformers establish architectures capable of modeling global spatial structure over high-resolution inputs, while contrastive and generative pretraining objectives align spatial patterns across vision, language, and generation.
Through training on internet-scale image-text pairs, video corpora, and synthetic data, spatial representations acquire statistical regularities that go beyond any single environment or task.
As a result, space increasingly functions as a shared coordinate system and embedding substrate that supports grounding, retrieval, reasoning, and conditional generation across modalities.
Spatial intelligence in this era is therefore inseparable from scale, as both data diversity and model capacity are essential for forming transferable spatial abstractions.



Another hallmark of this era is the elevation of the scene as a first-class unit of representation under large-scale learning.
Neural scene representations, most notably neural radiance fields (NeRFs)~\cite{mildenhall2021nerf}, treat a scene as a learnable and differentiable function defined over continuous space, enabling optimization from sparse observations.
Their significance lies not in implicitness itself, but in demonstrating that scene-level spatial structure can be recovered, manipulated, and reused through learning.
In parallel, diffusion models introduce powerful generative priors that encode spatial coherence in latent spaces trained on massive datasets, enabling controllable synthesis of images, videos, and three-dimensional content.
Together, these developments elevate space from something to be perceived into a scalable, generative, and composable substrate shaped through foundation-level pretraining.






\paragraph{Representative Techniques and Milestones.}
Before the emergence of vision-specific foundation models, a critical architectural shift had already taken place with the introduction of the Transformer \cite{attention}.
Originally developed for sequence modeling in natural language processing, attention-based architectures demonstrated that large-scale pretraining could yield highly transferable and task-agnostic representations, as exemplified by models such as BERT~\cite{devlin2019bert}.
The key insight was not tied to language itself, but to the ability of self-attention to capture global structure and support scalable representation learning.
This realization set the stage for extending the same architectural paradigm to visual data, where spatial structure could likewise be learned through large-scale pretraining rather than handcrafted inductive biases.
Following this, one central architectural milestone is the introduction of Vision Transformers~\cite{vit} in 2020, which demonstrates that large-scale pretraining on visual data can produce transferable and globally structured spatial representations.
This architectural shift lays the groundwork for foundation visual models that can be adapted across tasks without redesigning spatial pipelines.
Then, cross-modal alignment becomes a defining capability with the release of CLIP in 2021, which aligns visual representations with natural language through large-scale contrastive learning \cite{radford2021learning}.
This alignment provides a universal interface for spatial semantics, enabling zero-shot transfer, retrieval, and language-conditioned spatial reasoning.
Spatial representations learned in this manner escape the confines of a fixed label space and can be flexibly queried and recombined through language.




Concurrently, the emergence of LLMs marks a decisive inflection point in foundation modeling.
Unlike earlier neural models trained at modest scale, LLMs demonstrate that sufficiently large-scale pretraining on diverse data can induce qualitatively new capabilities, including abstraction, composition, and long-horizon reasoning \cite{achiam2023gpt,vicuna2023}.
Crucially, these abilities are not hard-coded but emerge from scale, data diversity, and architectural simplicity.
Language becomes the first modality in which foundation models exhibit general-purpose intelligence, rapidly positioning LLMs as a universal computational substrate that extends beyond NLP into vision, robotics, simulation, and scientific domains.
For spatial intelligence, this shift is pivotal: it introduces a representation space in which spatial structure, relations, and plans can be explicitly described, manipulated, and interfaced with external perceptual or physical systems, thereby elevating language from an auxiliary annotation channel to a central medium for spatial abstraction and control.
Also, Chain-of-Thought (CoT) style reasoning and tool-augmented prompting begin to systematically populate spatial intelligence pipelines \cite{wei2022chain}.
The key milestone is not CoT per se, but the emergence of \emph{structured intermediate representations} induced by LLMs, such as stepwise spatial decomposition, implicit scene graphs, explicit coordinate or layout hypotheses, and program-like plans that interface with mapping, perception, or simulation modules.
In practice, CoT becomes a mechanism to externalize and manipulate spatial states as compositional symbols that can be checked, revised, and grounded via perception or environment feedback, thereby tightening the representation loop from language to spatial structure and back.






In parallel, LLMs begin to serve as coordinators of scene-level reasoning rather than passive text processors.
PaLM-E~\cite{driess2023palm} in 2023 represents a critical milestone by integrating language models with embodied sensory inputs and robotic actions, demonstrating that a single foundation backbone can operate over language, vision, and control signals within a shared representational space.
This integration highlights that scene representations must bridge symbolic reasoning and physical execution, rather than treating them as separate layers.
Also, 2023 marks a clear milestone where Vision-Language-Action (VLA)~\cite{zitkovich2023rt2} models crystallize as a representation-centric paradigm.
The defining contribution of these systems is not embodiment itself, but the treatment of action as a first-class modality within a shared spatial latent interface.
VLA models learn to map perceptual states, such as images, video, and proprioception, together with language instructions, into temporally grounded action tokens or control variables.
This reframes spatial representation from a purely geometric or perceptual construct into an \emph{action-aligned internal state}, where affordances, kinematic constraints, and task-relevant spatial variables are explicitly encoded in executable form.
Enabled by large foundation backbones, this shift establishes a concrete representational bridge from scene-level spatial models to agent-level world states.





Building upon these representational advances, the foundation paradigm begins to visibly reshape embodied and interactive systems.
Early generalist agent models, such as Gato~\cite{reed2022generalist} in 2022, signal an attempt to unify perception, action, and spatial representation within a single pretrained architecture.
From 2022 to 2023, LLM- and VLM-backed embodied systems increasingly adopt foundation spatial representations as priors for online decision making in navigation and manipulation tasks \cite{driess2023palm}.
High-profile humanoid and mobile robot demonstrations in 2023 further highlight this trend, integrating LLM or VLM planners with perception stacks and low-level controllers.
Although many such systems remain engineering-driven rather than standardized academic benchmarks, they mark a qualitative shift: spatial representations turn from offline reconstructions into online task states that must remain coherent under motion, manipulation, and partial observability.
These developments collectively signal the maturation of foundation spatial representations into embodied world models, setting the stage for the transition toward world-centered spatial intelligence.



NeRFs \cite{mildenhall2021nerf}, introduced in 2020, represent another milestone by modeling scenes as continuous functions defined over spatial coordinates and viewing directions.
This formulation enables high-fidelity view synthesis and establishes a functional perspective on scene representation.
Follow-up work extends this idea to signed distance functions and other continuous fields, reinforcing the notion of scene-level spatial modeling as a reusable representation primitive.
Generative modeling advances rapidly during this period as well.
Denoising diffusion probabilistic models~\cite{ho2020denoising}, introduced around 2020, provide a stable and scalable framework for learning generative priors over high-dimensional spatial signals.
By 2022, large-scale text-to-image diffusion systems demonstrate that spatial layout, object relations, and scene structure can be controlled through language-conditioned latent spaces \cite{rombach2022ldm}.
These models turn generative spatial representations into widely accessible foundation priors.



In three-dimensional representation, the introduction of 3D Gaussian Splatting \cite{kerbl20233d} in 2023 significantly improves the practicality and efficiency of differentiable scene representations.
While conceptually aligned with neural rendering, its impact lies in making scene-level spatial representations more scalable and deployable.
This development fits naturally within the broader surge of foundation scene representations rather than defining a separate era.
Closely following this trend, four-dimensional (4D) spatial modeling and generation, which explicitly incorporate temporal dynamics into scene representations, begin to rapidly prevail in 2024 and attract unprecedented attention from the community \cite{kong20253d}.
Building upon continuous scene representations and generative priors established in three-dimensional settings, 4D models aim to represent scenes as time-conditioned functions or dynamic fields that capture motion, deformation, and long-term temporal consistency.
This includes emerging paradigms such as dynamic NeRF variants~\cite{pumarola2021d,li2021neural}, spatiotemporal Gaussian representations~\cite{wu20244d,bian2025gsditadvancingvideogeneration}, video-based neural fields~\cite{liu2025free4d}, and diffusion-based 4D generation models~\cite{zhen2025tesseract}.
The significance of this milestone lies in extending foundation spatial representations from static geometry to evolving worlds, thereby strengthening the role of space as a reusable and temporally grounded latent substrate.
Although still closely tied to the foundation modeling paradigm of this era, the rise of 4D spatial representations foreshadows the subsequent shift toward world-centered spatial intelligence, where temporal evolution and agent interaction become first-class constraints.



Across domains such as remote sensing and medical imaging, large-scale pretraining and multimodal alignment also begin to emerge.
Efforts to pretrain spatial representations on massive satellite imagery corpora or align medical images with textual reports reflect the spread of foundation spatial modeling beyond conventional vision benchmarks \cite{cong2022satmae,sun2022ringmo}.
Simulation and synthetic data generation further support this trend by enabling data scaling for spatial pretraining when real-world annotations are scarce~\cite{Cabon2020VirtualKITTI2,pfaff2025scalable,xu2025stable}.




% Overall, Era 4 establishes spatial representations as foundation assets.
% Space becomes a shared latent substrate that supports alignment, generation, and transfer across modalities and tasks.
% This shift prepares the ground for the subsequent transition toward world-centered spatial intelligence, where these foundation representations are no longer static priors, but executable components embedded within agents acting in the world.







\subsection{Era 5: World-Centered Spatial Intelligence (2023--Present)}
\label{subsec:era-world-centered}





\paragraph{From Spatial Representations to World States.}
This era marks a decisive transition from spatial modeling toward world-centered intelligence.
The defining shift lies in a fundamental redefinition of what spatial representations are required to support, well beyond the introduction of another three-dimensional encoding scheme.
For the first time, spatial representations are systematically constrained to be executable, interactive, continuously updatable, and tightly coupled with agent behavior.
They are no longer designed to serve isolated perception modules or one-shot generation tasks, but to function as persistent internal states that sustain long-running systems embedded in dynamic environments.
As a result, representation quality is no longer judged primarily by reconstruction accuracy or visual fidelity, but by whether it can support ongoing interaction, decision making, and adaptation over extended time horizons~\cite{ha2018world,Feifei2025spatial}.


A central conceptual shift underlying this transition is the expansion of context.
The dominant paradigm moves from context as instruction to context as world.
Rather than reasoning solely over prompts, local observations, or short-horizon inputs, systems are now required to organize cognition around an evolving world state that integrates spatial structure, temporal dynamics, physical constraints, and increasingly social context.
In this view, spatial intelligence operates as an internal world state machine.
It is written by perception, modified by action, maintained by memory, and queried by planning and prediction.
This reframing represents a departure from representations as static input encodings toward representations as operational substrates that mediate perception, reasoning, and action in a unified loop~\cite{genie3,driess2023palm}.


Crucially, this shift also redefines how learning itself is expected to occur.
World-centered spatial intelligence cannot be acquired solely through one-off offline pretraining.
Instead, representations must be shaped and continually refined through embodied interaction with the environment.
Agents are expected to actively explore, test hypotheses, observe consequences, and update their internal world states accordingly.
Spatial intelligence thus becomes inseparable from agency and experience.
The internal representation must encode not only where things are, but how the world evolves under action, which affordances are available, and how constraints shape future possibilities.
Although current systems remain at an early stage of this transition, the trajectory is clear.
By treating spatial intelligence as a persistent, executable, and evolving world substrate, this era points toward a path where artificial agents can maintain coherent world models over long time horizons, learn autonomously through interaction, and ultimately approach human-level and potentially superhuman-level intelligence grounded in the structure and dynamics of the real world~\cite{wang2023voyager,ali2025world,wang2023robogen}.





\paragraph{Representative Techniques and Milestones.}
One prominent signal of this era is the emergence of explicit world modeling frameworks that treat perception, prediction, and memory as components of a unified state transition process.
Yann LeCun has proposed a series of Joint Embedding Predictive Architecture (JEPA) world models from 2022 to 2025, which reframes representation learning around predicting future world states rather than reconstructing sensory inputs, emphasizing abstraction, temporal coherence, and action relevance \cite{assran2023jepa,assran2025vjepa,chen2025vl}.
By shifting the learning objective from pixel-level fidelity to state-level consistency, JEPA positions world dynamics as the primary object of representation.
Following this direction, a growing body of work between 2023 and 2024 explores world models that integrate prediction, memory, and action, progressively moving from short-horizon forecasting toward persistent and interaction-aware world state representations \cite{hafner2023mastering,hao2023reasoning,zhen20243d,yang2026mindjourney,garcin2026beyond}.
These efforts collectively signal a transition from predictive representation learning to explicit world-centered modeling.
Within this trajectory, World Labs' Marble, launched by Fei-Fei Li's team in late 2025 as one of the first generally available generative world models, turns text, image, video, or coarse 3D prompts into persistent, editable, and walkable 3D environments \cite{Feifei2025Marble}.
Its significance for this era is twofold: it demonstrates that world-level spatial representations can be generated rather than only reconstructed, and that the resulting worlds persist as stable, downloadable state rather than transient observations.
Marble thereby moves generative spatial modeling from research demonstrations toward deployable world-building infrastructure, reinforcing the view that spatial representations should be maintained, edited, and acted upon over long horizons.






VLA scaling further accelerates this trend.
RT-2 in 2023 demonstrates that models pretrained on large-scale web data can transfer semantic and spatial knowledge directly into robotic control policies \cite{zitkovich2023rt2}.
Here, spatial representation becomes action-aligned rather than geometry-first.
Internal states must encode affordances and task-relevant spatial variables in forms that are immediately executable, reinforcing the notion that representation quality is measured by behavioral competence.
Agent-centric frameworks also mature during this period.
ReAct-style agents, introduced in 2022, formalize a tight loop between reasoning and acting, where intermediate representations explicitly externalize spatial hypotheses, plans, and state updates \cite{yao2022react}.
Voyager, released in 2023 for the Minecraft environment, further exemplifies this shift by integrating long-term memory, skill libraries, planning, and exploration into a persistent world model that evolves through interaction~\cite{wang2023voyager}.
These systems highlight the importance of maintaining long-horizon spatial consistency rather than optimizing isolated decisions.



Simulation platforms and benchmarks increasingly align with this world-centered view.
The Habitat family of embodied AI benchmarks continues to push evaluation from static perception tasks toward closed-loop interaction, navigation, and manipulation in simulated environments \cite{savva2019habitat,szot2021habitat2,puig2023habitat3,Ramakrishnan2021HM3D}.
This shift in evaluation methodology reinforces the requirement that spatial representations must remain coherent under motion, partial observability, and delayed feedback.
From 2024 onward, interactive world models and action-conditioned video models gain prominence.
These models treat the world state as a predictive object whose evolution depends explicitly on agent actions, unifying spatial prediction and control \cite{cui2025emu35,bruce2024genie}.
Such approaches position world dynamics as first-class representational targets rather than side effects of perception.
More broadly, multimodal agent frameworks emerging between 2024 and 2025 increasingly combine tool use, memory systems, and long-horizon planning.
In these systems, spatial representations act as persistent world states that coordinate perception, reasoning, and action over extended time scales \cite{feng2025survey,yang2025embodiedbench,wang2023voyager}.




From a representational standpoint, this era is characterized by the rise of world-level state representations.
These representations often combine metric maps, semantic abstractions, scene graphs, and memory-augmented structures into unified internal states.
Crucially, they encode not only static layout but also action-conditioned dynamics, affordances, and long-term consistency.
The representation target shifts from isolated scenes to cross-scale world states that span rooms, cities, and even planetary systems.
Such representations are designed to be executable, editable, and safe, enabling planning, intervention, and constraint enforcement within a single coherent framework.









 % \RespondAuthor{Hao \& SQ \& Charles}

\chapter{A Definition in Motion: Revisiting Spatial Intelligence}
\label{chapter:definition}

% {
% \color{red} TODO: add enough necessary formulations, equations.
% }

% \lettrine[lines=3]{\initfamily\textcolor{darkblue}{W}}{hile}
\FramedInitial{W}{hile}
spatial intelligence is rapidly emerging as a foundational paradigm in artificial intelligence, its constituent ideas are by no means new. 
Concepts related to space, geometry, structure, and physical interaction have long existed across multiple specialized research communities, each developing its own assumptions, abstractions, and technical formalisms. 
What is new is the growing realization that these previously fragmented components must now be understood as parts of a unified and evolving whole, rather than the appearance of spatial reasoning per se.
This chapter revisits the notion of spatial intelligence as a moving definition, one that cannot be fixed within a single discipline, representation, or task setting, but must instead be understood through its cross-cutting roles across modern AI systems.










% \section{Spatial Intelligence as a Complex Interdisciplinary Ecosystem}
\section{The Landscape of Spatial Intelligence}
\label{sec:p1-sec2-the-landscape-of-spatial-intelligence}

% \lettrine[lines=3]{\initfamily\textcolor{darkblue}{A}}{t} 
\FramedInitial{A}{t}
its core, spatial intelligence is not a single technique, model, or task category.
Rather, it constitutes a complex interdisciplinary ecosystem, spanning multiple technical paradigms, sensory modalities, scientific disciplines, and application domains.
Different communities approach spatial intelligence with distinct motivations and meanings: for some, it concerns geometric consistency and localization; for others, it is about perception–action coupling, embodied reasoning, or spatial grounding of language.
As a result, any attempt to define spatial intelligence in isolation risks oversimplification.
Understanding spatial intelligence therefore requires acknowledging its breadth, heterogeneity, and layered structure as an evolving research ecosystem rather than a narrowly scoped subfield.







% \subsection{Spatial Intelligence Encompasses Far More Than a Single Capability}
\subsection{Beyond A Single Dimension}
\label{subsec:p1-sec2-beyond-a-single-dimension}

Spatial intelligence is fundamentally \emph{not simple}.
Unlike many AI topics that can be reasonably framed around a dominant abstraction or learning paradigm, spatial intelligence spans a wide spectrum of representational, computational, and conceptual elements.
It touches almost every stage of the perception–cognition–action loop and manifests differently depending on the task, environment, and embodiment of the agent.
This intrinsic complexity is not accidental; it reflects the fact that space itself is a unifying substrate of the physical world, cutting across sensory input, internal representation, and external interaction.


Historically, research on spatial problems in AI has progressed along parallel and often disconnected trajectories.
Simultaneous localization and mapping (SLAM), 3D reconstruction, robotic manipulation, geographic information systems, spatial language understanding, and neural rendering all evolved largely within their own communities, with specialized benchmarks, metrics, and assumptions.
While each of these areas addressed a slice of spatial reasoning, they rarely converged on a shared conceptual framework \cite{feng2025survey}.
As intelligent agents are increasingly expected to operate in open, dynamic, and embodied environments, the limitations of this fragmentation have become increasingly evident.


One major source of complexity arises from the fact that spatial intelligence cuts across \emph{multiple technical paradigms}.
It cannot be reduced to either symbolic rules or neural representations alone.
Instead, it involves the coexistence and interaction of diverse computational approaches:


\paragraph{Cross-Technical Paradigms.}
Spatial intelligence spans heterogeneous technical formalisms, including:
\begin{itemize}[leftmargin=2em,itemsep=0.35em]
    \item \textbf{Rule-based and symbolic systems}, such as topological maps, qualitative spatial calculi (\eg, RCC-8 region connection, interval algebra), and logical planners, which offer interpretability and structured priors \cite{kuipers2000spatial,egenhofer1991reasoning,Randell1992RCC,allen1983maintaining}.
    \item \textbf{Probabilistic and geometric methods}, including SLAM, state estimation, and pose graph optimization, which emphasize metric consistency and uncertainty modeling \cite{Thrun2005Probabilistic,dellaert2017factor}.
    \item \textbf{Neural representations}, such as implicit fields, neural radiance fields (NeRF), and learned spatial embeddings, which enable continuous and data-driven modeling of complex environments \cite{mildenhall2021nerf,kerbl20233d}.
    \item \textbf{Large language/foundation models}, which introduce spatial reasoning through language, symbolic abstraction, and multimodal alignment, but often lack direct physical grounding \cite{driess2023palm,wu2024gpt4vt23dgen}.
\end{itemize}
Equally important is the fact that spatial intelligence is inherently \emph{cross-modal}.
Space is not perceived through a single sensory channel, but emerges from the integration of multiple modalities, each offering partial and complementary information:


\paragraph{Cross-Modal Integration.}
Spatial intelligence operates across a diverse set of sensory and representational modalities:
\begin{itemize}[leftmargin=2em,itemsep=0.35em]
    \item \textbf{Language}, which encodes spatial relations, instructions, and abstract descriptions (\eg, ``behind'', ``inside'', ``near'')~\cite{landau1993whence}.
    \item \textbf{2D visual perception}, including images and video, which provide appearance-level cues and spatial layout hints.
    \item \textbf{3D representations}, such as point clouds, meshes, and volumetric fields, which capture geometric structure and physical extent.
    \item \textbf{4D spatiotemporal signals}, incorporating motion, dynamics, and temporal evolution of scenes.
    \item \textbf{Specialized sensing modalities}, including remote sensing, thermal imagery, radar, and spectral signals, which extend spatial perception beyond human-visible cues~\cite{cong2022satmae,sun2022ringmo}.
\end{itemize}
A further, frequently overlooked axis is \emph{spatial scale}.
Cognitive research has long established that space at different scales engages different psychological processes: Montello's classic analysis distinguishes figural (table-top), vista (single view), environmental (learned through locomotion), and geographical (map-mediated) spaces \cite{montello1993scale}.
The same stratification recurs in AI systems, where an agent manipulating objects on a desk, a robot navigating a building, and a model reasoning over a city or the Earth's surface demand different representations, memory structures, and update mechanisms \cite{feng2025survey}.
Competence at one scale does not transfer automatically to another, which is one reason benchmark results on room-scale tasks say little about urban-scale spatial reasoning.

\paragraph{Cross-Scale Operation.}
Spatial intelligence must therefore operate across nested spatial scales:
\begin{itemize}[leftmargin=2em,itemsep=0.35em]
    \item \textbf{Object and figural scale}, covering manipulation workspaces, shape understanding, and mental transformation of objects.
    \item \textbf{Scene and vista scale}, covering room-level layout, occlusion, and viewpoint reasoning.
    \item \textbf{Environmental scale}, covering buildings, streets, and spaces learned through movement and exploration.
    \item \textbf{Geographic and planetary scale}, covering cities, terrain, and Earth systems mediated by maps and remote sensing \cite{goodchild1992geographical,feng2025survey}.
\end{itemize}
Beyond modality, technique, and scale, spatial intelligence is deeply \emph{cross-disciplinary}.
Its theoretical foundations and practical realizations draw simultaneously from several scientific traditions:



\paragraph{Cross-Disciplinary Foundations.}
Spatial intelligence intersects with:
\begin{itemize}[leftmargin=2em,itemsep=0.35em]
    \item \textbf{Computer Vision}, through scene understanding, reconstruction, and spatial perception.
    \item \textbf{Robotics and Embodied AI}, where spatial reasoning underlies navigation, manipulation, and interaction.
    \item \textbf{Geographic Information Science (GIScience)}, which studies large-scale spatial representation, mapping, and spatial analytics~\cite{goodchild1992geographical,li2025giscience}.
    \item \textbf{Natural Language Processing and Linguistics}, particularly in grounding linguistic expressions in physical space and in the cross-linguistic study of spatial reference \cite{landau1993whence,levinson2003space}.
    \item \textbf{Psychology and Psychometrics}, which pioneered the measurement of human spatial abilities such as mental rotation and spatial visualization \cite{shepard1971mental,uttal2013malleability}.
    \item \textbf{Cognitive Science and Neuroscience}, which investigate human spatial cognition, cognitive maps, and embodied perception~\cite{tolman1948cognitive,stachenfeld2017hippocampus}.
\end{itemize}
Finally, spatial intelligence manifests across a wide spectrum of \emph{application domains}, further complicating any attempt at a narrow definition:



\paragraph{Cross-Application Presence.}
Spatial intelligence is central to applications such as:
\begin{itemize}[leftmargin=2em,itemsep=0.35em]
    \item \textbf{Embodied agents and robotics}, including navigation, manipulation, and human–robot interaction~\cite{Xia2018Gibson,feng2025survey,zitkovich2023rt2}.
    \item \textbf{Autonomous driving and aerial systems}, where spatial prediction and planning are safety-critical~\cite{ettinger2021womd,muhammad2022vision}.
    \item \textbf{AR/VR and virtual humans}, which demand consistent spatial grounding across perception and interaction~\cite{chen2021hightracking,orts2016holoportation,HuLYHCYLY23}.
    \item \textbf{Urban computing, digital twins, and Earth systems}, operating at city and planetary scales~\cite{LuvizonNM17,mu2024robotwin,picard2021autonomousAgents}.
\end{itemize}





Together, these dimensions highlight that spatial intelligence is not a standalone capability but a unifying layer across perception, reasoning, and action. 
This complexity precludes a minimal definition and instead calls for viewing spatial intelligence as a structured, evolving system that emerges across methods, modalities, and domains.






% requires: \usepackage{booktabs}
% optional: \usepackage{array}

\begin{table*}[t]
\centering
\caption{How different disciplines characterize spatial intelligence: primary focus, representations, canonical tasks, and evaluation emphasis.}
\vspace{2mm}
\small
\setlength{\tabcolsep}{6pt}
\renewcommand{\arraystretch}{1.15}

\begin{tabular}{p{2.5cm} p{3cm} p{3cm} p{3cm} p{3.cm}}
\toprule
\textbf{Discipline} &
\textbf{Primary focus} &
\textbf{Representations} &
\textbf{Canonical tasks} &
\textbf{Evaluation emphasis} \\
\midrule

Computer Vision \& Graphics &
Geometry recovery and view-consistent scene appearance &
Depth maps, multi-view geometry, meshes, neural fields (NeRF), 3D Gaussian representations &
Depth/pose estimation, 3D reconstruction, neural rendering, scene synthesis &
Geometric consistency, rendering fidelity, view synthesis quality \\
\addlinespace[2pt]
\hline
\addlinespace[2pt]

Robotics \& Embodied AI &
Perception--action coupling under physical constraints &
Occupancy grids, semantic maps, pose graphs, affordance maps, world models for planning &
SLAM, navigation, manipulation, motion planning, closed-loop control &
Task success, robustness, real-time performance, generalization in dynamic environments \\
\addlinespace[2pt]
\hline
\addlinespace[2pt]

GIS \& Large-scale Spatial Modeling &
Global-scale spatial abstraction, querying, and planning &
Geodetic coordinate systems, map layers, spatial graphs, spatial databases &
Route optimization, spatial analytics, trajectory analysis, urban-scale simulation &
Global consistency, query accuracy, scalability, long-horizon planning utility \\
\addlinespace[2pt]
\hline
\addlinespace[2pt]

Psychology \& Psychometrics &
Measurement of individual spatial abilities and their training &
Factor models of spatial ability (rotation, visualization, orientation), standardized test batteries &
Mental rotation tests, paper folding, perspective taking, spatial visualization assessment &
Test reliability and validity, individual differences, malleability and transfer \\
\addlinespace[2pt]
\hline
\addlinespace[2pt]

Cognitive Science \& Neuroscience &
Mechanisms of human spatial cognition and internal representations &
Cognitive maps, mental imagery, spatial memory representations, neural coding hypotheses &
Mental rotation, navigation studies, spatial memory and attention experiments &
Behavioral validity, explanatory power, neural plausibility \\
\addlinespace[2pt]
\hline
\addlinespace[2pt]

Linguistics \& Spatial Language &
How language encodes, schematizes, and constrains spatial relations &
Spatial prepositions and closed-class relations, frames of reference (intrinsic, relative, absolute) &
Spatial reference resolution, cross-linguistic comparison, grounded instruction understanding &
Descriptive adequacy, cross-linguistic coverage, grounding fidelity \\
\addlinespace[2pt]
\hline
\addlinespace[2pt]

Social \& Human-centered Perspectives &
Spatial behavior in social contexts and shared environments &
Proxemics, social space schemas, spatial cues for interaction, situational context models &
Human--robot interaction, social navigation, spatial language pragmatics, collaborative tasks &
Interaction naturalness, safety, alignment with human norms and intentions \\
\bottomrule
\end{tabular}
\label{tab:p1_s2_1_si-disciplinary-landscape}
\end{table*}





\subsection{Disciplinary Perspectives}
\label{subsec:p1-sec2-disciplinary-perspectives}

Although spatial intelligence is increasingly discussed as one concept, its meaning varies substantially across academic disciplines.
Each field approaches spatial intelligence through its own problem formulations, representational assumptions, and evaluation criteria.
As a result, what counts as spatial intelligence in one domain may only capture a partial aspect of what another domain considers essential.
This section reviews how spatial intelligence is understood from several major disciplinary perspectives, highlighting both their unique emphases and their limitations.




\paragraph{Computer Vision and Graphics.}
From the perspective of computer vision, spatial intelligence is primarily associated with the ability to infer geometric structure from visual input.
Research in this area focuses on problems such as depth estimation, multi-view geometry, camera pose estimation, and 3D scene reconstruction.
The core objective is to recover spatially consistent representations from 2D images or videos, often under strong geometric constraints.
These representations provide powerful spatial priors, enabling reasoning about layout, scale, and relative positioning of objects.


Closely related, computer graphics emphasizes the synthesis and rendering of spatial scenes.
Modern neural rendering approaches, including NeRF and 3D Gaussian Splatting, prioritize appearance fidelity, view consistency, and continuous scene representation.
While these methods excel at producing visually coherent 3D environments, they typically do not model agent-centric interaction, decision-making, or task-driven behavior.
As a result, spatial intelligence in vision and graphics is often appearance-centered and offline, rather than action-oriented.



\paragraph{Robotics and Embodied AI.}
In robotics and embodied AI, spatial intelligence is traditionally defined in terms of an agent’s ability to localize itself, navigate environments, and manipulate objects.
This includes spatial perception, motion planning in three-dimensional space, and interaction with representations such as occupancy grids, semantic maps, and volumetric models.
Here, spatial intelligence is inseparable from embodiment and control, as spatial reasoning directly informs action execution.
The central concern in this discipline is not only representing space, but using spatial representations to support real-time decision-making under physical constraints.
Consequently, spatial intelligence is evaluated through task success, robustness, and adaptability in dynamic environments.
This action-centric perspective highlights aspects of spatial intelligence that are often absent in purely perceptual or representational frameworks.



\paragraph{Geographic Information Systems and Large-Scale Spatial Modeling.}
At a larger spatial scale, geographic information systems view spatial intelligence through the lens of structured spatial data, coordinate systems, and topological relations.
GIS focuses on modeling, querying, and reasoning over large-scale geographic environments using maps, graphs, and spatial databases.
These systems excel at spatial querying, route optimization, and long-term planning across city or regional scales.
However, traditional GIS frameworks are rarely coupled with embodied agents or interactive learning.
Spatial intelligence in this context is largely static, symbolic, and detached from perception or action.
This contrast highlights a scale-dependent interpretation of spatial intelligence, where global consistency and abstraction take precedence over local interaction.



\paragraph{Psychology and Psychometrics.}
Long before AI adopted the term, differential psychology treated spatial ability as a measurable dimension of human intelligence.
Shepard and Metzler's mental rotation experiments showed that the time needed to compare rotated 3D shapes grows linearly with their angular disparity, providing the first chronometric evidence of analog spatial representation \cite{shepard1971mental}, and the factor-analytic tradition from Thurstone's space factor to Carroll's broad visuospatial factor established spatial ability as a distinct, testable construct \cite{carroll1993human}.
The term \emph{spatial intelligence} itself was popularized by Gardner's \emph{Frames of Mind} (1983), which listed it among seven human intelligences and defined it as the capacity to perceive the visual world accurately, to perform transformations upon one's initial perceptions, and to re-create aspects of visual experience even in the absence of the relevant stimuli; the construct, however, predates that label by half a century.
Two findings from this literature matter directly for AI.
First, spatial skills are malleable: a meta-analysis across hundreds of training studies shows they improve with practice and transfer across tasks \cite{uttal2013malleability}, suggesting spatial competence is acquired rather than fixed.
Second, the psychometric task inventory, covering mental rotation, perspective taking, and spatial visualization, has become a direct template for probing spatial understanding in multimodal foundation models, so this measurement tradition now feeds AI benchmarking rather than sitting apart from it \cite{yang2025vsibench}.

\paragraph{Cognitive Science and Neuroscience.}
In cognitive science and neuroscience, spatial intelligence is studied as an emergent property of interacting neural systems rather than as an engineered representation.
Research has identified key brain regions associated with spatial cognition, including the hippocampus for spatial memory and navigation (place cells, and grid cells in entorhinal cortex), the parietal cortex for spatial attention and transformation, and occipito-parietal networks for visuospatial processing \cite{o1978hippocampus,stachenfeld2017hippocampus}.
From this perspective, spatial intelligence arises from neural mechanisms supporting cognitive maps, mental imagery, and spatial reasoning; it is primarily internal and cognitive, centered on perception and imagination rather than external interaction.

\paragraph{Linguistics and Spatial Language.}
A further disciplinary lens comes from linguistics, which asks how language encodes and constrains spatial thought.
Talmy's analysis of how language schematizes space \cite{talmy1983how} and Landau and Jackendoff's account of the asymmetry between rich object naming and the small closed class of spatial prepositions \cite{landau1993whence} established that linguistic descriptions compress spatial structure into coarse, topology-like relations.
Levinson's cross-linguistic work further showed that languages differ systematically in their spatial frames of reference, distinguishing intrinsic, relative, and absolute systems \cite{levinson2003space}.
This vocabulary has migrated directly into AI evaluation, where egocentric versus allocentric reference handling is now a standard axis of spatial benchmarks for multimodal models, making the linguistics of space an operative rather than merely historical influence.




\paragraph{Social and Human-Centered Perspectives.}
Beyond physical geometry, spatial intelligence also carries meaning at the social and mental levels.
Human beliefs, emotions, and behaviors are situated within social environments that are themselves spatially structured.
From this viewpoint, spatial intelligence encompasses the ability to reason about others in shared spaces, interpret spatial cues in social contexts, and align internal mental states with external environments.
This perspective extends spatial intelligence beyond physical navigation toward a unified model of the physical and mental world.
It points to broader scientific and societal opportunities, including AI systems that operate in unstructured real-world settings and interact naturally with humans.
In this sense, spatial intelligence supports a joint modeling of physical space and human cognition, serving as a bridge between external environments and internal mental states.




% Taken together, these disciplinary perspectives illustrate why spatial intelligence cannot be reduced to a single definition.
% Each field captures an essential but partial view.
% Recognizing these differences is a prerequisite for any attempt to unify spatial intelligence under a coherent conceptual framework.





\subsection{Toward A Unified View}
\label{subsec:p1-sec2-toward-a-unified-view}

For decades, research on spatial reasoning and spatial representation has progressed from markedly different perspectives across multiple subfields.
These efforts often evolved in isolation, lacking a shared vocabulary, common research objectives, or directly comparable evaluation benchmarks.
As a result, techniques developed in areas such as SLAM, 3D vision, scene modeling, world modeling, and embodied planning frequently addressed overlapping problems without converging toward a unified conceptual framework.
Spatial intelligence emerges against this backdrop as a response to fragmentation, offering a higher-level abstraction that can organize and align these historically disjoint research trajectories.



Under the emerging spatial intelligence paradigm, a set of shared demands and core concepts has begun to surface.
A central driving force behind this convergence is the growing need for agents that exhibit spatial consistency and physical groundedness, and that can operate continuously within dynamic environments.
Such agents must integrate perception, representation, and action in a coherent manner.
This requirement has catalyzed the fusion of visual perception with geometry-aware modeling, as seen in navigation-oriented neural rendering approaches and interactive NeRF-based systems.
Research at the intersection of neural scene representation and active planning has given rise to a new generation of 3D world models that are not only reconstructive, but also executable and actionable.
Similarly, vision-language models are increasingly extended to operate in three-dimensional environments, enabling multimodal agents to localize, describe, and interact with objects within voxelized or volumetric scenes.



At the conceptual level, a set of shared primitives is gradually taking shape.
Spatial representations, whether instantiated as voxel grids, signed distance fields, neural feature volumes, or topological graph structures, are becoming central components across domains ranging from autonomous driving and AR/VR to assistive robotics.
Fundamental spatial relations, such as containment, adjacency, support, and traversability, are being formalized and embedded into reasoning modules, scene graph systems, and relational neural networks.
These abstractions provide a common representational substrate that enables knowledge transfer and system interoperability across tasks and modalities.



We note that the research community has yet to settle on a single taxonomy for this emerging field.
Recent surveys and position papers organize spatial intelligence along markedly different axes: by task family and benchmark \cite{chen2024spatialvlm}, by the cognitive functions that video-based spatial tasks exercise \cite{yang2025vsibench}, by spatial scale from embodied agents through cities to Earth systems \cite{feng2025survey}, or, in Li's position essay, by the capabilities a world model must possess, namely being generative, multimodal, and interactive \cite{Feifei2025spatial}.
These framings overlap without coinciding, and none has achieved consensus status.
This handbook's representation-centric organization, which follows how spatial representations are built, used, and acted upon, is our deliberate choice among these alternatives: it is the axis along which the historical evolution of Section~\ref{sec:historical-evolution} unfolds most naturally, and the one that lets fragmented communities be compared on common ground.

Taken together, these developments indicate that spatial intelligence is no longer confined to any single domain.
It is evolving into a unifying functional substrate that spans perception, reasoning, and interaction; bridges vision, language, and action; and connects symbolic representations with neural models.
The unification runs deeper than shared technique.
Spatial intelligence aligns and integrates multiple disciplines, including computer vision, robotics, psychology, cognitive science, linguistics, geospatial systems, human-computer interaction, and neuro-inspired AI, and through this alignment it functions as the connective tissue that enables the next generation of generalist AI systems to perceive the world, reason about it, and act within it.























\begin{figure*}
    \centering
    % \resizebox{0.98\linewidth}{!}{% Framework diagram: environment <-> AI system with internal spatial-representation hub
% Redesigned 2026-08-10: (i) fixed overlapping operator boxes; (ii) R_S now sits INSIDE
% the AI system as its central hub, with per-operator arrows; (iii) notation matches the
% Formal Characterization (f^S operators, Phi_S, pi_S, Psi_S, task t in T).
\begin{tikzpicture}[
    bigbox/.style={draw, rounded corners=5pt, align=center, inner sep=8pt, line width=0.6pt},
    opbox/.style={draw, rounded corners=3pt, align=center, text width=2.8cm, inner sep=5pt, font=\footnotesize, fill=white, line width=0.5pt},
    dimchip/.style={draw, rounded corners=3pt, align=center, text width=2.6cm, inner sep=4pt, font=\footnotesize, fill=white, line width=0.5pt},
    lbl/.style={font=\small\bfseries, align=center},
    sub/.style={font=\footnotesize\itshape, color=black!60, align=center},
    arr/.style={-{Stealth[length=2.6mm]}, line width=1pt, black!70},
    arrd/.style={-{Stealth[length=2.2mm]}, line width=0.7pt, black!60, dashed},
]
% ---------- left: environment ----------
\node[bigbox, fill=DefGreen, draw=forestgreen!80, text width=3.4cm, minimum height=6.8cm] (env) at (0,0) {};
\node[lbl, anchor=north] at ($(env.north)+(0,-0.18)$) {Physical World /\\[-1pt] Environment $\mathcal{E}$};
\node[sub, anchor=north, text width=3.1cm] at ($(env.north)+(0,-1.0)$) {geometric, physical,\\ causal constraints};
\node[dimchip, anchor=north] (sc1) at ($(env.north)+(0,-2.0)$) {objects \& scenes};
\node[dimchip, below=0.32 of sc1] (sc2) {rooms \& buildings};
\node[dimchip, below=0.32 of sc2] (sc3) {cities \& terrain};
\node[sub, below=0.16 of sc3] {nested spatial scales};
\node[sub, anchor=south] at ($(env.south)+(0,0.3)$) {tasks $t\in\mathcal{T}$ defined over $\mathcal{E}$};

% ---------- right: AI system ----------
\node[bigbox, fill=ai_blue!16, draw=ai_blue!80!black, minimum width=11.4cm, minimum height=6.8cm] (sys) at (10.0,0) {};
\node[lbl, anchor=north] at ($(sys.north)+(0,-0.18)$) {AI System $\mathcal{S}$};

% central representation hub (maintained INSIDE the system)
\node[bigbox, fill=customyellow!55, draw=orange!75!black, text width=3.3cm, inner sep=5pt] (hub) at (10.0,-0.08) {
    {\small\bfseries Spatial\\[-1pt] Representation $\mathcal{R}_{\mathcal{S}}$}\\[1pt]
    {\scriptsize\itshape\color{black!60} $\Phi_{\mathcal{S}}:\mathcal{E}\rightarrow\mathcal{R}_{\mathcal{S}}$}\\[3pt]
    {\scriptsize \textbf{granularity:} object / scene / relation}\\[2pt]
    {\scriptsize \textbf{modality:} pixels, points, language, \dots}\\[2pt]
    {\scriptsize \textbf{abstraction:} metric / topological / symbolic}\\[2pt]
    {\scriptsize \textbf{concrete vs.\ abstract:} observed / internal model}
};

% four functional operators around the hub
\node[opbox, draw=ai_blue!80!black] (perc) at (6.05,1.35) {Perception\\ $f^{\mathcal{S}}_{\text{perc}}:\mathcal{O}\rightarrow\mathcal{R}_{\mathcal{S}}$};
\node[opbox, draw=ai_blue!80!black] (act)  at (6.05,-1.35) {Embodied Action\\ $f^{\mathcal{S}}_{\text{act}}:\mathcal{R}_{\mathcal{S}}\rightarrow\mathcal{A}$};
\node[opbox, draw=ai_blue!80!black] (reas) at (13.95,1.35) {Reasoning\\ $f^{\mathcal{S}}_{\text{reason}}:\mathcal{R}_{\mathcal{S}}\rightarrow\mathcal{R}_{\mathcal{S}}$};
\node[opbox, draw=ai_blue!80!black] (pred) at (13.95,-1.35) {Prediction\\ $f^{\mathcal{S}}_{\text{pred}}:\mathcal{R}_{\mathcal{S}}\times\mathcal{A}\rightarrow\mathcal{R}_{\mathcal{S}}$};
\node[sub, font=\scriptsize\itshape, anchor=north] at ($(reas.south)+(0,-0.06)$) {elaborates the current state};
\node[sub, font=\scriptsize\itshape, anchor=north] at ($(pred.south)+(0,-0.06)$) {advances the state in time};

% task-conditioned action selection
\node[dimchip, text width=1.9cm] (task) at (6.05,0) {task $t\in\mathcal{T}$};
\draw[arrd] (task.south) -- node[right, font=\scriptsize, color=black!60]{$\pi_{\mathcal{S}}$} (act.north);

% representation update rule
\node[sub, anchor=south] at ($(sys.south)+(0,0.12)$) {state update\; $\mathcal{R}^{\mathcal{S}}_{t+1}=\Psi_{\mathcal{S}}(\mathcal{R}^{\mathcal{S}}_t,\,a^{\mathcal{S}}_t,\,o_{t+1})$};

% ---------- arrows ----------
% perception--action loop, closed through the environment
\draw[arr] ($(env.east)+(0,1.35)$) -- node[above, font=\footnotesize, align=center, pos=0.4]{observations\\ $o_{t}\in\mathcal{O}$} (perc.west);
\draw[arr] (act.west) -- node[below, font=\footnotesize, align=center, pos=0.6]{actions\\ $a^{\mathcal{S}}_{t}\in\mathcal{A}$} ($(env.east)+(0,-1.35)$);
% operators read from / write to the representation hub
\draw[arr] (perc.east) -- (perc.east -| hub.west);
\draw[arr] (act.east -| hub.west) -- (act.east);
\coordinate (reasA) at ($(reas.west)+(0,0.2)$);
\coordinate (reasB) at ($(reas.west)+(0,-0.2)$);
\draw[arr] (reasA -| hub.east) -- (reasA);
\draw[arr] (reasB) -- (reasB -| hub.east);
\coordinate (predA) at ($(pred.west)+(0,0.2)$);
\coordinate (predB) at ($(pred.west)+(0,-0.2)$);
\draw[arr] (predA -| hub.east) -- (predA);
\draw[arr] (predB) -- (predB -| hub.east);
\node[sub, anchor=north] at (6.85,-3.65) {continuous perception--action loop, closed through the environment};
\end{tikzpicture}
}
    % \includegraphics[width=0.99\linewidth]{figures/part1/si_framework.png}
    \includegraphics[width=\linewidth]{figures/part1/si_framework_v2.pdf}
    \vspace{1mm}
    \caption{Spatial intelligence as a closed loop between an AI system $\mathcal{S}$ and its environment $\mathcal{E}$. Perception writes observations into a world-centered spatial representation $\mathcal{R}_{\mathcal{S}}$ maintained inside the system; reasoning elaborates the current state, prediction advances it under a candidate action, and embodied action reads it to act on the environment, whose response is observed in turn. Each module is labeled with its operator from Section~\ref{subsec:spatial-intelligence-formalization}.}
    \label{fig:p1_s2_1_intro}
\end{figure*}





\section{Defining A Spatial Intelligence Framework}
\label{sec:p1-sec2-defining-a-spatial-intelligence-framework}


% \lettrine[lines=3]{\initfamily\textcolor{darkblue}{G}}{iven} 
\FramedInitial{G}{iven}
the inherently hybrid and interdisciplinary nature of spatial intelligence, it becomes necessary to revisit and refine what we mean by ``spatial intelligence'' in a precise and unified manner.
Rather than treating spatial intelligence as a loosely defined collection of capabilities, we aim to articulate a coherent conceptual framework that can accommodate its diverse manifestations across systems, tasks, and environments.
In this section, we therefore propose an explicit definition of spatial intelligence and analyze its core properties in a structured way.

\begin{center}
\begin{tcolorbox}[colback=gray!5,colframe=gray!60]
\textbf{Definition (Spatial Intelligence).}
\emph{Spatial intelligence is the ability of an AI system to construct, maintain, and manipulate spatial representations of the physical world or environment that support perception, reasoning, prediction, and embodied action, to accomplish tasks within a given environmental context.}
\end{tcolorbox}
\end{center}

This definition emphasizes several essential aspects.
First, spatial intelligence is grounded in representations of the world rather than isolated sensory inputs.
Second, it is functional rather than static, supporting perception, reasoning, prediction, and action.
Third, it is inherently situated, meaning that spatial intelligence only manifests meaningfully within an environment where tasks are defined and executed.
In the following, we unpack this definition by examining its key attributes.






\subsection{Core Elements}
\label{subsec:p1-sec2-core-elements}

\subsubsection{Parties}

At a minimum, spatial intelligence involves two fundamental parties: an AI system and the physical world or environment in which it operates.
Neither can be meaningfully defined without the other.

\paragraph{AI System.}
The AI system refers to any computational entity capable of perception, reasoning, and decision-making.
This includes, but is not limited to, large language models, multimodal foundation models, world models, embodied agents, and robotic systems.
Crucially, spatial intelligence is not tied to a specific architecture.
Instead, it characterizes a class of capabilities that may be realized through different model families, training paradigms, and system designs \cite{feng2025embodied,zhang2026theory}.

\paragraph{Physical World or Environment.}
The environment denotes the external context in which the AI system is situated.
This typically refers to the physical world governed by geometric, physical, and causal constraints.
In some cases, the environment may also include social or human-centric settings, where spatial structure interacts with social norms, human intent, and collective behavior \cite{hall1966hidden}.
Spatial intelligence arises precisely from the interaction between the AI system and such structured environments.



\subsubsection{Spatial Representation}

A core component of spatial intelligence lies in the ability to form, maintain, and manipulate world-centered spatial representations.
These representations serve as structured world models that are compositional, generalizable, and reusable across tasks.

Spatial representations may vary along several dimensions:
\begin{itemize}[leftmargin=2em,itemsep=0.35em]
    \item \textbf{Granularity.}
    Representations can operate at different spatial scales, including object-centric representations, scene-centric representations, and relational representations capturing relative geometry or topology.
    
    \item \textbf{Modality.}
    Spatial information may be derived from heterogeneous sensory channels, including pixels, grid maps, point clouds, voxel grids, language descriptions, audio signals, spectral measurements, radar inputs, as well as tactile and force feedback.
    
    \item \textbf{Level of Abstraction.}
    Spatial representations may be geometric, such as metric coordinates and surfaces, topological, such as connectivity graphs, or semantic and symbolic, such as object categories and relational predicates.
    
    \item \textbf{Concrete versus Abstract.}
    Representations may directly encode observable properties of the physical world or serve as abstract internal models that summarize, compress, or reinterpret spatial structure for reasoning and planning.
\end{itemize}

Importantly, spatial intelligence does not require committing to a single representational form.
Instead, it demands the ability to align and translate across multiple representational levels in a world-centered manner.



\subsubsection{Dynamic Interaction}

Spatial intelligence is not a static property acquired through passive observation.
Rather, it emerges through continuous and dynamic interaction with the environment.
This interaction includes navigation, manipulation, exploration, and embodied perception.
An intelligent system must ground its spatial understanding in embodied experience, for example by learning the structure of a room through movement, or improving depth perception through active viewpoint selection.
Interaction also extends to socially situated behaviors.
Maintaining appropriate interpersonal distance, coordinating actions with humans, or navigating shared spaces all require spatial understanding informed by social context.
As a result, spatial representations must be executable and continuously updated through perception-action feedback loops.
We further decompose dynamic interaction into several functional components.



\paragraph{Perception.}
Perception refers to the extraction of spatial structure from multimodal sensory inputs.
This includes recognizing objects and scene elements, estimating distances, identifying spatial layouts, and localizing entities within an environment.

\paragraph{Reasoning and Planning.}
Spatial reasoning involves inference over spatial relations and dynamics.
This includes trajectory prediction, path planning, occlusion reasoning, spatial causality, temporal change modeling, relational inference, and logical or commonsense reasoning about space.
Such reasoning spans both low-level problems, such as shortest-path computation, and high-level abstractions, such as understanding how rearranging objects alters the affordances of a space.

\paragraph{Predicting.}
Prediction in spatial intelligence operates along two axes.
One concerns forecasting future spatial states or events in a manner consistent with physical laws and causal structure.
The other concerns generative capabilities, such as producing spatially aware images, videos, 3D or 4D scenes, and motion trajectories.

\paragraph{Embodied Action.}
Embodied action denotes the capacity to interact with physical or simulated environments while continuously pursuing goals.
This includes pose adjustment, navigation, grasping, and object manipulation.
Typical scenarios include domestic robots and autonomous driving systems.
Such action requires tight integration of perception, reasoning, planning, and prediction, while aligning internal spatial models with task objectives, human intent, and social context.





% \paragraph{Remark.}
% In summary, a well-functioning spatially intelligent AI system must:
% \begin{itemize}[leftmargin=2em,itemsep=0.35em]
%     \item Learn a high-quality and comprehensive spatial representation of the world.
%     \item Support dynamic interaction through accurate perception, reasoning, prediction, and embodied action.
%     \item Maintain consistency between internal world models, external task demands, and environmental feedback.
% \end{itemize}
% These properties collectively distinguish spatial intelligence from purely perceptual or symbolic forms of intelligence.



\vspace{4mm}
\begin{tcolorbox}[
    colback=gray!5,
    colframe=black,
    title=\textbf{What Constitutes a Strong Spatially Intelligent System}
]
A well-designed spatially intelligent AI system should satisfy the following criteria:
\begin{itemize}[leftmargin=2em,itemsep=0.35em]
    \item \textbf{High-quality world representation.}  
    It must acquire a rich, consistent, and generalizable spatial representation of the world that captures geometry, relations, and structure across scales.
    
    \item \textbf{Dynamic interaction capability.}  
    It should support accurate perception, reasoning, prediction, and embodied action through continuous interaction with the environment.
    
    \item \textbf{Alignment between model, task, and environment.}  
    The internal spatial model must remain aligned with external task objectives, environmental constraints, and feedback signals during execution.
\end{itemize}
\end{tcolorbox}







\subsection{Formal Characterization}
\label{subsec:spatial-intelligence-formalization}

To make the definition more precise, we provide a lightweight formalization that captures the core components of spatial intelligence without committing to a specific implementation or learning paradigm.

Let $\mathcal{E}$ denote the environment in which an AI system is situated, representing the physical or simulated world governed by spatial, physical, and causal constraints.
Let $\mathcal{S}$ denote the AI system or agent.
We assume that tasks are defined with respect to the environment and are drawn from a task space $\mathcal{T}$.




\paragraph{World-Centered Spatial Representation.}
We define a world-centered spatial representation as a system-induced mapping from the environment to an internal state space:
\begin{equation}
\Phi_{\mathcal{S}} : \mathcal{E} \rightarrow \mathcal{R}_{\mathcal{S}},
\end{equation}
where $\mathcal{R}_{\mathcal{S}}$ denotes the internal spatial representation space maintained by the AI system $\mathcal{S}$.
The representation $\Phi_{\mathcal{S}}(e) \in \mathcal{R}_{\mathcal{S}}$ may take various forms, including geometric maps, scene graphs, implicit fields, symbolic layouts, or learned latent embeddings.
Crucially, $\mathcal{R}_{\mathcal{S}}$ is not a raw sensory space, but a structured and system-dependent representation space that supports compositionality, generalization, and reuse across tasks.



% \paragraph{World-Centered Spatial Representation.}
% We define a world-centered spatial representation as a mapping from the environment to an internal state space:
% \begin{equation}
% \Phi : \mathcal{E} \rightarrow \mathcal{R},
% \end{equation}
% where $\mathcal{R}$ denotes the space of internal spatial representations.
% The representation $\Phi(e) \in \mathcal{R}$ may take various forms, including geometric maps, scene graphs, implicit fields, symbolic layouts, or learned latent embeddings.
% Crucially, $\mathcal{R}$ is not a raw sensory space, but a structured representation space that supports compositionality, generalization, and reuse across tasks.



% \paragraph{Functional Use of Spatial Representations.}
% Spatial intelligence is not defined by representation alone, but by how representations are functionally manipulated.
% We characterize this through a family of operators acting on $\mathcal{R}$:
% \begin{equation}
% \mathcal{F} = \{ f_{\text{perc}}, f_{\text{reason}}, f_{\text{pred}}, f_{\text{act}} \},
% \end{equation}
% corresponding to perception, reasoning, prediction, and embodied action.
% Each operator maps between representations, actions, or outcomes, for example:
% \begin{align}
% f_{\text{perc}} &: \mathcal{O} \rightarrow \mathcal{R}, \\
% f_{\text{reason}} &: \mathcal{R} \rightarrow \mathcal{R}, \\
% f_{\text{pred}} &: \mathcal{R} \rightarrow \mathcal{R}, \\
% f_{\text{act}} &: \mathcal{R} \rightarrow \mathcal{A},
% \end{align}
% where $\mathcal{O}$ denotes observations and $\mathcal{A}$ denotes action space.

\paragraph{Functional Use of Spatial Representations.}
Spatial intelligence is not defined by representation alone, but by how representations are functionally manipulated by the AI system.
We characterize this through a family of system-conditioned operators acting on $\mathcal{R}_{\mathcal{S}}$:
\begin{equation}
\mathcal{F}_{\mathcal{S}} = \{ f^{\mathcal{S}}_{\text{perc}}, f^{\mathcal{S}}_{\text{reason}}, f^{\mathcal{S}}_{\text{pred}}, f^{\mathcal{S}}_{\text{act}} \},
\end{equation}
corresponding to perception, reasoning, prediction, and embodied action.
Each operator is defined relative to the system $\mathcal{S}$:
\begin{align}
f^{\mathcal{S}}_{\text{perc}} &: \mathcal{O} \rightarrow \mathcal{R}_{\mathcal{S}}, \\
f^{\mathcal{S}}_{\text{reason}} &: \mathcal{R}_{\mathcal{S}} \rightarrow \mathcal{R}_{\mathcal{S}}, \\
f^{\mathcal{S}}_{\text{pred}} &: \mathcal{R}_{\mathcal{S}} \times \mathcal{A} \rightarrow \mathcal{R}_{\mathcal{S}}, \\
f^{\mathcal{S}}_{\text{act}} &: \mathcal{R}_{\mathcal{S}} \rightarrow \mathcal{A},
\end{align}
where $\mathcal{O}$ denotes observations and $\mathcal{A}$ denotes the action space.
The two representation-to-representation operators play distinct roles: $f^{\mathcal{S}}_{\text{reason}}$ elaborates the \emph{current} world state by inferring relations and constraints already implicit in it, whereas $f^{\mathcal{S}}_{\text{pred}}$ advances the state \emph{forward in time}, rolling out how the world would evolve under a hypothesized action (with a null action recovering passive forecasting).




% \paragraph{Task-Oriented Intelligence.}
% Given a task $t \in \mathcal{T}$, spatial intelligence manifests as the ability of the system to select and execute actions based on its internal spatial representation:
% \begin{equation}
% \pi : (\mathcal{R}, t) \rightarrow \mathcal{A},
% \end{equation}
% such that actions are grounded in the current world state and lead to task completion within the environment.
% This highlights that spatial intelligence is inherently situated: the same internal model may support different behaviors depending on task constraints and environmental context.


\paragraph{Task-Oriented Intelligence.}
Given a task $t \in \mathcal{T}$, spatial intelligence manifests as the ability of the AI system $\mathcal{S}$ to select and execute actions based on its internal spatial representation:
\begin{equation}
\pi_{\mathcal{S}} : \mathcal{R}_{\mathcal{S}} \times \mathcal{T} \rightarrow \mathcal{A},
\end{equation}
such that actions are grounded in the system’s current world representation and lead to task completion within the environment.
This formulation highlights that spatial intelligence is inherently situated and agent-dependent: the same environment may induce different behaviors depending on the internal representation and decision mechanisms of $\mathcal{S}$.




% \paragraph{Dynamic Interaction and World State Evolution.}
% Finally, spatial intelligence requires representations to be dynamically updated through interaction.
% We model the evolution of the internal world state as:
% \begin{equation}
% \mathcal{R}_{t+1} = \Psi(\mathcal{R}_t, a_t, o_{t+1}),
% \end{equation}
% where $a_t \in \mathcal{A}$ is the executed action and $o_{t+1} \in \mathcal{O}$ is the subsequent observation.
% This formulation emphasizes that spatial representations are executable, interactive, and continuously revised, rather than static encodings.


\paragraph{Dynamic Interaction and World State Evolution.}
Finally, spatial intelligence requires representations to be dynamically updated through interaction.
We model the evolution of the internal world representation of system $\mathcal{S}$ as:
\begin{equation}
\mathcal{R}^{\mathcal{S}}_{t+1} = \Psi_{\mathcal{S}}(\mathcal{R}^{\mathcal{S}}_t, a^{\mathcal{S}}_t, o_{t+1}),
\end{equation}
where $a^{\mathcal{S}}_t \in \mathcal{A}$ is the action selected by system $\mathcal{S}$ and $o_{t+1} \in \mathcal{O}$ is the subsequent observation.
It emphasizes that spatial representations are executable, interactive, and continuously revised by the system, rather than static encodings of the environment.





This formalization clarifies that spatial intelligence is not tied to a specific modality, algorithm, or architecture.
Instead, it is defined by the existence of a system-induced, world-centered spatial representation space $\mathcal{R}_{\mathcal{S}}$ and the ability of the AI system $\mathcal{S}$ to manipulate, update, and exploit this representation to support perception, reasoning, prediction, and embodied action within an environment.










\subsection{Compatibility with Existing Research and Paradigms}
\label{subsec:p1-sec2-compatibility-with-existing-research-and-paradigms}

The proposed framework is intentionally designed to be inclusive rather than prescriptive.
Rather than introducing a new task taxonomy or replacing existing formulations, it provides a unifying lens through which a wide range of spatial intelligence-related tasks, domains, and research traditions can be consistently described.
From this perspective, most existing work can be understood as instantiating or emphasizing only a subset of the core elements defined in our framework, often driven by domain-specific constraints, data availability, or evaluation protocols.


A large body of research in computer vision, graphics, and extended reality focuses primarily on the \emph{representation layer} of spatial intelligence.
Tasks such as 3D reconstruction, neural rendering, scene synthesis, video generation, and emerging 4D spatiotemporal modeling concentrate on learning expressive spatial representations $\mathcal{R}_{\mathcal{S}}$, often with little or no emphasis on action, task execution, or long-term interaction.
Within our framework, these approaches can be viewed as developing increasingly powerful instantiations of world-centered spatial representations, while largely abstracting away the agent $\mathcal{S}$ and the functional operators that connect representation to behavior.
Their contributions lie in advancing the expressiveness, continuity, and generative capacity of spatial representations, even if they remain largely offline or perception-centric.



Robotics, embodied AI, and embedded systems emphasize a different subset of the framework, placing primary focus on action grounding and closed-loop interaction.
Vision-Language-Action models, navigation agents, and manipulation systems operationalize the mapping from internal representations to actions, often under real-time and physical constraints.
However, in many cases, the underlying spatial representations remain task-specific, shallow, or tightly coupled to particular embodiments.
From the perspective of our framework, these systems instantiate the functional operators and dynamic interaction components, while relying on relatively limited or implicit world representations.
This explains why many embodied systems exhibit strong short-horizon competence but struggle with long-term consistency, world editing, or cross-task transfer.



Classical and modern SLAM systems provide a complementary example.
SLAM explicitly targets the construction and maintenance of a world-centered spatial representation under uncertainty, typically in the form of metric maps, pose graphs, or hybrid geometric-topological structures.
Within our framework, SLAM can be understood as a specialized instantiation of the representation update process, with strong guarantees on geometric consistency and state estimation.
At the same time, SLAM systems often remain agnostic to high-level reasoning, semantic abstraction, or task-conditioned decision making, highlighting again that they occupy a specific slice of the broader spatial intelligence spectrum rather than the full stack.




Viewed through this unified lens, existing spatial intelligence research does not fragment into incompatible paradigms, but instead forms a landscape of partial realizations.
Different communities emphasize representation learning, perception, reasoning, prediction, or action to varying degrees.
Our framework provides a common conceptual structure that makes these emphases explicit, clarifies their limitations, and reveals how advances in one area can be systematically integrated with others.
In this sense, the framework is not a replacement for existing approaches, but a scaffold for understanding how they collectively contribute to, and fall short of, world-centered spatial intelligence.





\chapter{Scope and Structure of This Handbook}
\label{chapter:structure}



% \lettrine[lines=3]{\initfamily\textcolor{darkgreen}{T}}{he}
% study of \textit{Spatial Intelligence} spans a vast and rapidly evolving landscape that intersects geometry, perception, language, and embodied reasoning. As this interdisciplinary field matures, it becomes imperative to map out its foundational pillars, ongoing transformations, and practical applications in a structured and accessible manner. This chapter serves as a cartographic overview of the book: we first introduce the six core research pillars that define the scope of spatial intelligence research—representation, perception, reasoning, generation, manipulation, and application—each of which corresponds to a major Part in the book. We then provide guidance on how different audiences—beginners, researchers, and practitioners—may navigate the content based on their background and objectives. This roadmap is designed not only to clarify the structure of the book but also to reflect the multifaceted nature of spatial intelligence as a scientific and engineering pursuit.

% \lettrine[lines=3]{\initfamily\textcolor{darkblue}{T}}{his}
\FramedInitial{T}{his}
chapter provides a structural overview of the handbook.
We clarify the scope, goals, and organization of the volume, and explain how the different parts jointly form a coherent research landscape for spatial intelligence.
Specifically, we introduce the core research pillars covered in this handbook, outline their relationships, and offer guidance on how readers from different backgrounds may navigate the content.






\section{The Research Pillars of Spatial Intelligence}
\label{sec:p1-sec3-the-research-pillars-of-spatial-intelligence}

% \lettrine[lines=3]{\initfamily\textcolor{darkblue}{W}}{e} 
\FramedInitial{W}{e}
outline the core research pillars that constitute the structural backbone of this handbook. 
They traverse the full pipeline from spatial representation and perception to reasoning and generative capabilities, and ultimately extend into concrete real-world applications. 
Each pillar aligns closely with a corresponding Part or chapter of the handbook, serving as a thematic axis for both foundational methods and emerging challenges in spatial intelligence.



\begin{itemize}[leftmargin=2em,itemsep=0.35em]
    \item \textbf{Representation (Part II: Spatially Aware Representation)} \\
    This pillar investigates how spatial information is encoded and represented across modalities and dimensions. We begin with single-modality representations, including 2D (\eg, visual frames), 3D (\eg, point clouds, meshes), and 4D (\eg, dynamic scenes, spatio-temporal fields), as well as non-visual modalities such as language and audio. Particular attention is given to cross-modal alignment and fusion mechanisms that unify heterogeneous data sources. These representational strategies provide the necessary structural foundation for downstream reasoning, manipulation, and generation.

    \item \textbf{Perception (Part III, Chapter~\ref{sec:perception}: Spatial Perception)} \\
    This pillar focuses on how intelligent systems extract structured knowledge from spatial environments. Topics include object detection, semantic and instance segmentation, depth and surface normal estimation, visual odometry (VO), and ego-motion prediction. These perception modules translate raw sensory inputs into geometric or semantic abstractions, enabling agents to comprehend and navigate the physical world.

    \item \textbf{Reasoning (Part III, Chapter~\ref{sec:reasoning}: Spatial Reasoning)} \\
    Reasoning in spatial intelligence spans multiple paradigms: symbolic reasoning over geometric primitives or logical propositions; geometric reasoning over shapes, topologies, and configurations; neural reasoning using deep architectures; and LLM-driven spatial inference based on large-scale pretrained priors. These methods enable agents to interpret relations, deduce spatial constraints, and plan in dynamic settings.

    \item \textbf{Reconstruction and Generation (Part III, Chapter~\ref{sec:generation}: Spatial Reconstruction and Generation)} \\
    This pillar examines the generative side of spatial intelligence. It encompasses tasks such as 3D reconstruction from observations, 2D/3D scene synthesis, controllable image and video generation, trajectory and motion prediction, and dynamic world creation. These generative processes are often guided by spatial prompts, semantic layouts, or agent intentions, forming a bridge between structured knowledge and creative capability.


    \item \textbf{Applications (Part IV: Applications)} \\
    Finally, the application pillar illustrates how the above methods power real-world systems. Domains include embodied agents and robotics (\eg, navigation, grasping), autonomous driving and aerial agents, warehouse and wearable AI systems, extended reality (XR), including virtual and augmented reality, geospatial AI (\eg, urban planning, satellite analysis), biomedical and scientific computing (\eg, 3D medical reconstruction), intelligent assistants, and further fields such as construction, sports, cultural heritage, marine systems, and retail. These chapters emphasize system-level integration and highlight how spatial AI capabilities translate into societal impact.
\end{itemize}




    % \item \textbf{Manipulation (Part III–IV: Learning to Act in Space)} \\
    % While less often isolated as a standalone topic, spatial manipulation—learning to act within and modify space—is essential to embodied agents. This includes spatial action prediction, affordance learning, scene editing, and interactive control. Whether through robotic embodiment or virtual actuation, this pillar emphasizes the actionable understanding of space, enabling agents to not only perceive and plan, but also intervene \cite{}.









\section{Goal of the Handbook: Charting a Research Ecosystem}
\label{sec:p1-sec3-goal-of-the-handbook-charting-a-research-ecosystem}

\FramedInitial{T}{he} primary goal of this handbook is not to present a catalog of isolated techniques, but to chart an emerging research ecosystem centered on spatial intelligence.
As established in earlier chapters, spatial intelligence is neither a single method nor a narrowly scoped task.
Rather, it is a unifying paradigm that connects representation, perception, reasoning, generation, and action through structured models of the world.

\paragraph{From Isolated Methods to an Integrated Landscape.}
Much of the existing literature on spatial intelligence is fragmented.
Some surveys focus deeply on individual subfields, such as SLAM, neural rendering, or 3D generation, while others adopt a high-level perspective that lacks technical grounding.
This handbook aims to bridge this gap by organizing spatial intelligence as an interconnected landscape, where advances in one component are explicitly related to others.
Our objective is to reveal how seemingly distinct research threads jointly contribute to world-centered intelligence.

\paragraph{Spatial Intelligence as a Connective Substrate.}
A central thesis of this handbook is that spatial intelligence functions as a connective substrate for next-generation AI systems.
It provides the structural glue that links perception with reasoning, language with geometry, and planning with embodied action.
By framing spatial intelligence in this way, the handbook draws connections across computer vision, robotics, simulation, geospatial systems, cognitive science, and neuro-inspired AI.
This perspective enables a shared conceptual vocabulary through which contributions from different communities can be meaningfully aligned.

\paragraph{Identifying Gaps and Open Tensions.}
Beyond synthesizing existing work, this handbook deliberately highlights unresolved gaps in the current research landscape.
These include the disconnect between symbolic and neural spatial representations, the limited grounding of perception in actionable world models, and the difficulty of integrating spatial reasoning into long-horizon embodied agents.
By making these tensions explicit, we aim to clarify where current paradigms fall short and where future breakthroughs are most likely to emerge.

\paragraph{A Blueprint for an Open Research Program.}
Ultimately, this handbook is intended as a blueprint for an open and evolving research program.
Rather than proposing a closed taxonomy or a definitive solution, it offers a structured lens through which spatial intelligence can be studied, extended, and reimagined.
We invite readers not only to learn from the systems and methods surveyed here, but also to rethink spatial intelligence as a first-class organizing principle in the design of generalist intelligent agents.



% This handbook does not merely aim to summarize a set of techniques, \ie, it seeks to map an emerging ecosystem. 
% Spatial intelligence is not a singular method nor a narrow problem formulation; it is a paradigm in the making, spanning multiple axes: from geometry to semantics, from perception to embodiment, and from individual tasks to collaborative agents.

% We argue that existing surveys often fall short in one of two ways: either they are confined within a single domain (\eg, SLAM, NeRF, 3D generation), or they take an overly generic perspective that lacks technical depth and cross-domain synthesis. 
% In contrast, our goal is to present spatial intelligence as an evolving research ecology—composed of interlocking parts, shaped by historical currents, and converging toward shared open problems.


% This handbook positions spatial intelligence as the connective substrate of next-generation AI. It draws from diverse yet converging fields such as computer vision, robotics, simulation, GIS, neuroAI, and cognitive science. It proposes a shared vocabulary—both conceptual and technical—to allow researchers and practitioners to recognize how their contributions fit into a broader picture.

% Crucially, we also aim to identify the missing links: where symbolic and neural representations remain disconnected; where perception has not yet been grounded in spatial structures; where embodied agents still operate with impoverished world models. By highlighting these fault lines, we hope to catalyze new research directions and collaborative bridges.

% Ultimately, this survey aspires to serve as a blueprint—not for a closed taxonomy, but for an open and generative research program. We invite readers to not only learn from existing systems, but to rethink spatial reasoning as a first-class citizen in the architecture of intelligent agents.







\section{How to Read This Handbook}
\label{sec:p1-sec3-how-to-read-this-handbook}


% \lettrine[lines=3]{\initfamily\textcolor{darkgreen}{G}}{iven} 
\FramedInitial{G}{iven} the interdisciplinary nature of spatial intelligence, which spans AI, computer vision, robotics, graphics, GIS, and cognitive science, it is crucial to provide tailored reading paths that help each type of reader gain the most value from this handbook.

% part serves as a practical navigation guide for readers with diverse backgrounds and objectives.



\vspace{1em}

\noindent\textbf{$\blacktriangleright$ For Beginners and Students} \textit{(\eg, undergraduates, junior graduate students, newcomers to spatial AI)}
\begin{itemize}[leftmargin=2em,itemsep=0.35em]
    \item Begin with \textbf{Part I: Introduction and Background} (Chapters~1--3) to grasp the motivation, historical evolution, and research scope.
    \item Then proceed to \textbf{Part II: Spatially Aware Representation} to build a solid understanding of core representations in 2D, 3D, and 4D, including geometry and topology.
    \item Optionally skim \textbf{Part IV: Applications} to get a glimpse of real-world use cases, such as robotics, autonomous driving, and AR/XR.
\end{itemize}

\vspace{0.5em}

\noindent\textbf{$\blacktriangleright$ For Researchers and Academics} \textit{(\eg, PhD students, postdocs, faculty, or interdisciplinary scholars in AI-related fields)}
\begin{itemize}[leftmargin=2em,itemsep=0.35em]
    \item Focus on \textbf{Part III: Spatial Intelligence, from Perception to Generation}, which includes chapters on:
    \begin{itemize}[leftmargin=2em,itemsep=0.35em]
        \item Spatial perception, from detection and segmentation to depth and odometry
        \item Symbolic, geometric, neural, and LLM-driven spatial reasoning
        \item Generative spatial models for scenes, videos, and worlds
    \end{itemize}
    \item Also study \textbf{Part IV: Applications} for cross-domain synergies, particularly in GeoAI, healthcare, urban simulation, and embodied cognition.
    \item Consider \textbf{Part V: Discussions and Futures} for open problems, benchmarking gaps, and research directions.
\end{itemize}

\vspace{0.5em}

\noindent\textbf{$\blacktriangleright$ For System Developers and Engineers} \textit{(\eg, practitioners building simulation environments, digital twin platforms, or perception pipelines)}
\begin{itemize}[leftmargin=2em,itemsep=0.35em]
    \item Prioritize chapters in \textbf{Part III (Generation)} and \textbf{Part IV (Applications)} that introduce:
    \begin{itemize}[leftmargin=2em,itemsep=0.35em]
        \item Modular spatial generation frameworks
        \item Multi-agent simulation engines and control loops
        \item Deployed system domains (\eg, robotic manipulation, warehouse automation, wearable AI, aerial agents)
    \end{itemize}
    \item Use \textbf{Part V: Discussions and Futures} to anticipate where deployment bottlenecks and benchmark gaps lie.
\end{itemize}

\vspace{0.5em}

\noindent\textbf{$\blacktriangleright$ For Cross-disciplinary Explorers} \textit{(\eg, scholars in architecture, biology, geoscience, or cognitive science)}
\begin{itemize}[leftmargin=2em,itemsep=0.35em]
    \item Start with \textbf{Part I (Introduction and Background)} to acquire shared conceptual grounding.
    \item Dive into \textbf{Part IV (Applications)} for domain-aligned chapters and illustrative case studies.
    \item Refer to \textbf{Part II (Representation)} to gain a technical foothold in spatial structures and encodings.
\end{itemize}

\vspace{1em}

\noindent Overall, this handbook is modular in structure and can be read either linearly or selectively depending on the reader's focus. The chapter cross-references are designed to aid flexible navigation and long-term use.




\part{Spatially Aware Representation: Building Blocks of Space}
% ================================
% \section*{\centering \bf Part II: Spatial-aware Representation}
% \addcontentsline{toc}{section}{Part II: Spatial-aware Representation}
% ================================





\begin{figure}
    \centering
    % \includegraphics[width=\linewidth, height=12cm]{example-image-golden}
    % \includegraphics[width=0.99\linewidth]{figures/part2.png}
    \includegraphics[width=\linewidth]{figures/part2_opener_v2.pdf}
    
    \caption{Part II in one view. (1) Encode first: a scene must be encoded before it can be perceived, reasoned about, or acted on. (2) One modality at a time: the same object as pixels, depth, points, meshes, voxels, and Gaussians, and then as language and sound, from explicit geometry to implicit structure (Chapter~\ref{chapter:single_modality}). (3) Into a shared space: the separate encodings are aligned in one space without losing the geometry that makes them spatially meaningful (Chapter~\ref{chapter:cross_modality}).}
    \label{fig:p2_main}
\end{figure}


Everything a spatially intelligent system does, perceiving a scene, reasoning over it, generating or acting within it, passes through some representation of space.
That representation is upstream of the capability: what a system can perceive is bounded by what its encoding can express, and what it can reason about is bounded by what that encoding makes accessible.
Before the later parts turn to perception, reasoning, and generation, this part asks the prior question of how spatial structure is captured at all, how the geometry, topology, and semantics of an environment come to be held in a form a machine can compute with.

The question is approached in two movements.
The first takes each modality on its own terms, asking how images, video, 3D scenes, language, and sound each encode spatial structure and what each renders explicit or leaves implicit.
The second asks how these separate encodings are brought into a shared space, so that signals as different as a pixel grid and a sentence can be aligned without losing the geometry that makes them spatially meaningful.
The order is deliberate, since a modality must be understood by itself before its representation can be fused with others, and it is in that fusion that the partial, sensor-bound view offered by any one signal begins to reinforce the rest.
Together they form the representational substrate on which the remainder of the handbook is built.

% Building upon this, spatial representation
% learning focuses on constructing effective and generalizable spatial abstractions that capture geometric,
% topological, and semantic properties of the environment.

\newpage


\chapter{Single-Modal Spatial Representation}
\label{chapter:single_modality}


% \lettrine[lines=3]{\initfamily\textcolor{darkgreen}{S}}{ingle-modal} 
\FramedInitial{S}{ingle-modal} spatial representations ask how much of a scene's spatial structure can be recovered from one modality alone, without recourse to cross-modal correspondence.
The question matters because every multimodal system is ultimately built from unimodal encoders, and the spatial content a fused representation can carry is bounded by what each modality already captures on its own.

What separates the modalities is not whether they carry spatial information but how overtly they carry it.
Vision presents geometry almost directly, arranging appearance and shape over a spatial grid or the volume a scene occupies.
Language and sound hold the same structure only obliquely, the one as symbolic relations among named entities, the other as the physical trace that propagation leaves on a signal.
Others still are tied to a particular sensor or task, trading generality for fidelity to the data they were built around.
This chapter proceeds along that gradient, from the most explicit geometric encodings of space toward the most implicit, examining each modality on its own terms before the question of how they combine is taken up.





\section{2D}
\label{sec:rep_2d}

\FramedInitial{T}{he} 2D representations are the most direct and information-dense entry point into spatial intelligence: images and videos provide a regular, discretely sampled grid over the visual field, on top of which a family of structured maps can be defined to expose geometric and semantic properties of the underlying scene.
We organize this section around seven such representations (Figure~\ref{fig:p2_s4_1_2d_representations}): the raw \emph{image/video tensor} itself, and six derived maps: \emph{segmentation}, \emph{depth}, \emph{surface normal}, \emph{edge}, \emph{optical flow}, and \emph{2D keypoint} representations.
All seven share the image grid as their domain, which makes them straightforward to fuse with one another and to consume with convolutional or transformer backbones; what differs is what each representation \emph{exposes}: appearance, semantics, geometry, boundaries, motion, or sparse structure.



\begin{figure}[!t]
    \centering
    \includegraphics[width=0.90\linewidth]{figures/part2/2D_representations.pdf}
    \caption{The seven 2D representations of this section, all defined over the image grid $\Omega \subset \mathbb{R}^2$: image/video, segmentation, depth, surface normals, edges, optical flow, and keypoints.}
    \label{fig:p2_s4_1_2d_representations}
\end{figure}

\subsection{Image and Video Representations} 
\label{subsec:p-2-sec-4-2D-image-and-video-representations}

An image can be modeled as a discrete sampling of a continuous visual signal over a 2D spatial domain. Formally, a digital image is represented as a function
\begin{equation}
I: \Omega \rightarrow \mathbb{R}^C, \quad \Omega \subset \mathbb{R}^2,
\end{equation}
where $\Omega$ denotes the image grid, $C$ is the number of channels (\eg, $C=1$ for grayscale, $C=3$ for RGB), and each pixel value encodes intensity or color information at a discrete spatial location \cite{gonzalez2011,szeliski2022}. 
In practice, images are stored as finite tensors $I \in \mathbb{R}^{H \times W \times C}$, where $H$ and $W$ denote image height and width.


A video extends this representation by introducing a temporal dimension.
Given a sequence of frames sampled over time, a video can be written as
\begin{equation}
V: \Omega \times \mathcal{T} \rightarrow \mathbb{R}^C,
\end{equation}
or equivalently as a 4D tensor $V \in \mathbb{R}^{T \times H \times W \times C}$, where $T$ denotes the number of frames \cite{szeliski2022}. 
This formulation preserves spatial structure within each frame while enabling the modeling of temporal dynamics across frames.
From a computational perspective, images and videos are treated as structured arrays of numerical values that can be efficiently processed using linear algebra operations and convolutional kernels. 
This regular grid structure enables weight sharing, local receptive fields, and efficient parallel computation, forming the basis of CNNs and their spatiotemporal extensions~\cite{lecun1998gradient,karpathy2014large}. 
In modern deep learning systems, image and video tensors may undergo normalization, resizing, patch embedding, or tokenization before being fed into neural
architectures, including CNNs and vision transformers~\cite{vit, liu2022convnet}.
For video inputs, spatiotemporal extensions either inflate 2D kernels into 3D convolutions~\cite{tran2015learning, carreira2017quo} or extend self-attention across space and time, as in the family of video transformers~\cite{arnab2021vivit, liu2022videoswin}.






\subsection{Segmentation Maps}
\label{subsec:p-2-sec-4-2D-segmentation-maps}
A segmentation map is a semantic representation that assigns a label or a mask to each pixel, thereby partitioning visual data into meaningful regions such as objects, object parts, or background categories~\cite{szeliski2022,li2023seg_survey}. For a single image defined on $\Omega \subset \mathbb{R}^2$, segmentation can be written as a mapping
\begin{equation}
S: \Omega \rightarrow \mathcal{L},
\end{equation}
where $\mathcal{L}$ is a set of labels. In \emph{semantic segmentation} (SS), $S(u,v)$ denotes the class label of pixel $(u,v)$. 
In \emph{instance segmentation} (IS), the output is a set of masks
$\{M_k\}_{k=1}^{K}$ with associated class labels $\{\ell_k\}$, where each $M_k:\Omega \rightarrow \{0,1\}$ corresponds to an object instance \cite{long2015fcn, he2017mask}.
 \emph{Panoptic segmentation} (PS) unifies SS and IS by predicting a single, non-overlapping labeling of pixels that distinguishes both ``stuff'' and countable object instances \cite{kirillov2019panoptic}.

These tasks extend naturally to videos. Given a sequence of
frames $\{I_t\}_{t=1}^{T}$, \emph{video semantic segmentation} (VSS) predicts per-frame semantic labels $\{S_t\}$, while \emph{video instance segmentation} (VIS) predicts instance masks with temporal consistency across frames, often requiring instance association (tracking) over time \cite{li2023seg_survey}. 
\emph{Video panoptic segmentation} (VPS) further combines VSS and VIS, producing temporally consistent panoptic outputs, and is commonly evaluated with metrics that account for both segmentation quality and tracking consistency~\cite{li2023seg_survey}. 
Closely related
settings include \emph{video object segmentation} (VOS)~\cite{perazzi2016benchmark}, which segments foreground objects across time under semi-supervised or unsupervised protocols, and \emph{referring} segmentation, which predicts masks grounded by natural language expressions in images (RIS)~\cite{wang2022cris} or videos (RVOS)~\cite{seo2020urvos}.
Segmentation maps are typically formed by dense prediction models that produce either hard labels or per-pixel class probabilities $P(u,v,\ell)$ (and, for instance-level tasks, a set of object masks and scores). 
A key advantage of segmentation is its explicit and interpretable structure: it converts low-level pixels into object- and region-level representations that support downstream reasoning and decision-making, and it can be fused naturally with geometry (\eg, depth, normals) and language in multimodal systems~\cite{li2023seg_survey}. 
Consequently, segmentation maps are central to applications such as autonomous driving,
medical image analysis, robotic manipulation, visual editing, and human-computer interaction.


\subsection{Depth Maps}
\label{subsec:p-2-sec-4-2D-depth-maps}
A depth map is a geometric representation that augments a 2D image with per-pixel distance information, describing the spatial layout of a scene from a given viewpoint~\cite{szeliski2022}. 
Formally, given an image domain $\Omega \subset \mathbb{R}^2$, a depth map is a scalar function
\begin{equation}
D: \Omega \rightarrow \mathbb{R},
\end{equation}
where $D(u,v)$ denotes the distance along the camera's optical axis to the first visible surface at pixel $(u,v)$; some sensors instead report radial range along the viewing ray.
Related representations store inverse depth or disparity, which must be converted to the appropriate metric depth or range before geometric back-projection~\cite{hartley2003multiple}.
Depth maps can be obtained through multiple mechanisms, including stereo vision via triangulation~\cite{hartley2003multiple}, active sensors such as structured-light or time-of-flight cameras~\cite{foix2011tof}, and learning-based monocular depth estimation~\cite{eigen2014depth}. 
When combined with known camera intrinsics, a depth map enables the back-projection of image pixels into 3D space, thereby serving as a bridge between 2D image representations and 3D geometry.
A key advantage of depth maps lies in their compact, image-aligned structure, which allows them to be processed using standard convolutional architectures and easily fused with the RGB appearance information. 
Compared to volumetric representations, depth maps are memory-efficient while still encoding meaningful geometric cues. 
As a result, they are widely used in applications such as 3D reconstruction, novel view synthesis~\cite{mildenhall2021nerf}, robotic navigation, and autonomous driving~\cite{KITTI-GeigerLU12}.


\subsection{Surface Normal Maps}
\label{subsec:p-2-sec-4-2D-surface-normal-maps}
A surface normal map is a geometric representation that encodes local surface orientation at each pixel of an image. Given an image domain
$\Omega \subset \mathbb{R}^2$, a surface normal map assigns to each pixel a unit vector
\begin{equation}
\mathbf{N}: \Omega \rightarrow \mathbb{S}^2,
\end{equation}
where $\mathbf{N}(u,v) = (n_x, n_y, n_z)$ represents the direction of the surface normal in camera or world coordinates, and $\mathbb{S}^2$ denotes the unit sphere. Surface normal maps can be derived from depth maps or 3D geometry via local differentiation, or directly estimated from RGB images using learning-based methods~\cite{barron2015shape}.
Classical approaches to normal estimation rely on photometric stereo or geometric reconstruction~\cite{woodham1980photometric}, while modern methods employ CNNs to regress surface normals from monocular images~\cite{eigen2015predicting}. 
Regardless of the estimation procedure, surface normal maps remain image-aligned and
encode rich local geometric information that complements depth and segmentation cues.
Surface normal maps can capture fine-grained shape and orientation details that may be ambiguous or noisy in raw depth estimates.
Because surface normals are invariant to absolute depth scale, they provide robust geometric constraints for scene understanding. 
Consequently, surface normal maps are widely used in applications such as 3D reconstruction, object pose estimation, scene layout analysis, realistic rendering, and as auxiliary supervision in multi-task learning~\cite{eigen2015predicting}.


\subsection{Edge Maps}
\label{subsec:p-2-sec-4-2D-edge-maps}
An edge map is a structural representation that highlights salient intensity or color discontinuities in an image, which often correspond to object boundaries and significant geometric transitions~\cite{canny1986}.
Given an image $I: \Omega \rightarrow \mathbb{R}^C$, an edge map is typically constructed by measuring local image gradients, for example,
\begin{equation}
E(u,v) = \|\nabla I(u,v)\|,
\end{equation}
where $\nabla I$ denotes the spatial gradient operator.
Classical edge detection methods rely on hand-crafted operators such as Sobel, Laplacian of Gaussian, and the Canny detector~\cite{canny1986}, while modern learning-based approaches directly predict the edge probabilities using convolutional neural networks~\cite{xie2015hed}.
Edge maps may be binary, indicating the presence or absence of an edge, or real-valued, encoding edge strength, orientation, or confidence.
The primary advantage of edge maps is their ability to significantly reduce visual complexity while preserving essential geometric structure.
This abstraction makes them relatively robust to changes in illumination and appearance. 
Consequently, edge maps are widely used in image
segmentation, contour detection, shape analysis, line drawing, and as intermediate cues for higher-level tasks such as object recognition and depth estimation~\cite{szeliski2022}.


\subsection{Optical Flow Maps}
\label{subsec:p-2-sec-4-2D-optical-flow-maps}
An optical flow map is a motion representation that encodes the apparent 2D displacement of pixels between consecutive frames in a video sequence~\cite{horn1981optical}. 
Given two image frames $I_t$ and $I_{t+1}$, the optical flow is defined as a vector field
\begin{equation}
\mathbf{F}(u,v) = (f_x(u,v), f_y(u,v)),
\end{equation}
where $(f_x, f_y)$ describes the horizontal and vertical displacement of pixel $(u,v)$ from time $t$ to $t+1$.
Classical optical flow estimation relies on assumptions such as
brightness constancy and spatial smoothness, leading to variational formulations~\cite{horn1981optical,brox2004high}. More recently, deep neural networks have been proposed to directly regress dense flow fields from image pairs or short video clips~\cite{dosovitskiy2015flownet,teed2020raft}.
Regardless of the estimation method, the resulting optical flow map remains image-aligned, preserving spatial correspondence with the input frames.
The main advantage of optical flow maps is their ability to capture rich dynamic information about scene motion, including object motion, camera motion, and motion parallax, while remaining computationally efficient.
As a result, optical flow is a core representation in video analysis and
robotic perception tasks such as motion segmentation, action
recognition, video interpolation, and visual odometry.



\subsection{2D Keypoints}
\label{subsec:p-2-sec-4-2D-2d-keypoints}
A 2D keypoint set is a sparse image representation that encodes the image-plane locations of a finite set of semantically or geometrically meaningful points, such as corners, joints, or landmarks~\cite{szeliski2022}.
Unlike dense maps that assign a value to every pixel, keypoint representations summarize visual content with a small number of salient points, trading pixel-level coverage for compactness and geometric explicitness.
Formally, given an image defined on a domain
$\Omega \subset \mathbb{R}^2$, a keypoint representation is a set of $K$ points
\begin{equation}
\mathcal{K} = \{(x_k, y_k)\}_{k=1}^{K}, \quad (x_k, y_k) \in \Omega,
\end{equation}
where each keypoint corresponds to a specific semantic or geometric entity. 
In learning-based formulations, keypoints are often represented implicitly as confidence or heat maps:
\begin{equation}
H_k : \Omega \rightarrow [0,1],
\end{equation}
with the keypoint location obtained by taking the argmax of $H_k$ or the expected image coordinates under its normalized weights \cite{newell2016stacked}. 
2D keypoints can be obtained using classical feature detectors, such as Harris corners, SIFT, and SURF \cite{harris1988combined,lowe2004distinctive}, or through modern deep neural networks that directly regress keypoint locations or predict dense heat maps \cite{toshev2014deeppose,newell2016stacked}.
Depending on the task, keypoints may represent generic interest points, object-specific landmarks, or articulated structures such as the human body joints.
The detection \emph{task} that produces such a set, including its extension to landmark detection and human pose estimation, is discussed in detail in Section~\ref{sec:keypoint_landmark} of Chapter~\ref{sec:perception}.
A key advantage of 2D keypoint representations is their sparsity and compactness: they significantly reduce the dimensionality of visual data while preserving critical structural information. 
This makes keypoints computationally efficient, robust to background clutter, and well-suited for geometric reasoning and correspondence matching.
Human pose estimation is a specialized and highly structured instance of 2D keypoint detection that focuses on localizing a predefined set of anatomical landmarks of the human body, such as joints and facial keypoints, in images or videos \cite{zheng2022pose}. 
In the 2D setting, the task aims to estimate the image coordinates of body joints, typically represented as a set of keypoints
$\mathcal{P} = \{(x_j, y_j)\}_{j=1}^{J}$, where $J$ denotes the number of joints in a human skeletal model. 
These keypoints are not independent: they are constrained by the underlying human body structure, which is often modeled using kinematic trees, graphical models, or learned
structural priors.






\section{3D}
\label{sec:rep_3d}

\FramedInitial{T}{he} 3D representations lift the 2D pixel grid into the ambient space the scene actually occupies, replacing image-plane appearance with explicit geometry and topology.
They are the substrate on which downstream tasks that require metric reasoning, occlusion handling, viewpoint manipulation, or physical simulation are built, from autonomous driving and robotic manipulation to medical visualization and immersive graphics.
Compared with the depth and surface normal maps of Section~\ref{sec:rep_2d}, which still inhabit the image grid, 3D representations expose geometry directly and can therefore describe parts of the scene that are not visible in any single view.

We organize this section around the eight 3D representations that dominate current practice (Figure~\ref{fig:p2_s4_2_3d_representation_1}, Table~\ref{tab:p2_s4_1_3d_comparison}): RGB-D images, voxels, point clouds, meshes, implicit fields, neural radiance fields (NeRF), 3D Gaussian Splatting (3DGS), and hybrids of the above.
They differ in what they store (surfaces vs.\ volumes, explicit primitives vs.\ neural functions), how they trade off fidelity against memory and rendering speed, and how amenable they are to editing and to gradient-based optimization.

\begin{itemize}[leftmargin=2em,itemsep=0.35em]
    \item \textbf{RGB-D Image}. Pairs a standard RGB image with a per-pixel depth map; view-dependent and cannot capture occluded surfaces from a single frame.
    \item \textbf{Voxel}. A uniform 3D grid of volumetric elements; conceptually simple and well suited to inherently volumetric data, but memory grows as $O(N^3)$ for $N$ cells per axis.
    \item \textbf{Point Cloud}. An unstructured set of 3D points, typically the native output of LiDAR or photogrammetry; faithful to the raw measurement but lacks surface connectivity.
    \item \textbf{Mesh}. Surfaces defined by vertices and faces; the dominant representation in graphics for its GPU-friendliness, at the cost of a fixed connectivity that complicates topology change.
    \item \textbf{Implicit Representation}. Defines surfaces as the level set of a continuous function (\eg, Signed Distance Functions); resolution-agnostic and topology-flexible, but slow to render.
    \item \textbf{NeRF (Neural Radiance Field)}. A neural function over a 5D input domain (3D location + 2D viewing direction) that supports state-of-the-art novel-view synthesis.
    \item \textbf{3DGS (3D Gaussian Splatting)}. Represents a scene as anisotropic 3D Gaussian primitives rendered through a differentiable splatting rasteriser, reaching NeRF-level quality at real-time speeds.
    \item \textbf{Hybrid Methods}. Combinations of the above (\eg, DMTet, Tri-plane) that pair explicit primitives with implicit decoders to balance fidelity, efficiency, and gradient-based optimizability.
\end{itemize}

\begin{figure}[!ht]
  \centering
  \includegraphics[width=0.99\linewidth]{figures/part2/3d_representation.pdf}
  \caption{Examples of 3D representations: RGB-D image, voxel grid, point cloud, triangle mesh, and a signed distance field (SDF) slice.}
  \label{fig:p2_s4_2_3d_representation_1}
\end{figure}


% \bigskip

\begin{table}[!t]
\centering
% \small
% \setlength{\tabcolsep}{4pt} % Adjust column spacing for a better fit
\caption{Qualitative comparison of 3D representations (relative ratings, not measured values).}
\label{tab:p2_s4_1_3d_comparison}
% \fontsize{9}{11}\selectfont
% \setlength{\tabcolsep}{2.5mm}
\resizebox{\linewidth}{!}{
\begin{tabular}{l|ccccc}
\toprule
\textbf{Representation} & \textbf{Fidelity} & \textbf{Memory Cost} & \textbf{Render Speed} & \textbf{Topology Flexibility} & \textbf{Editability} \\
\midrule
\textbf{RGB-D Image}    & Low (2.5D)  & Low                 & High                  & View-dependent           & Low          \\
\textbf{Voxel}          & Low--Medium & High ($O(N^3)$)     & Medium                & High                     & Medium       \\
\textbf{Point Cloud}    & Medium      & Medium (scalable)   & High                  & High                     & Medium       \\
\textbf{Mesh}           & High        & Low (GPU-friendly)  & Very High             & Low (fixed connectivity) & High         \\
\textbf{Implicit (SDF)} & High        & Low                 & Low (volume query)    & High                     & Low          \\
\textbf{NeRF}           & Very High   & Low                 & Low (ray marching)    & High                     & Very Low     \\
\textbf{3DGS}           & Very High   & Medium              & Very High (real-time) & High                     & Medium       \\
\textbf{Hybrid}         & Very High   & Medium              & High                  & Adaptive                 & Medium--High \\
\bottomrule
\end{tabular}
}
\end{table}


\subsection{RGB-D Images}
\label{sec:rep_3d_rgbd}

RGB-D images are a popular 2.5D representation that captures both photometric and geometric information from a single viewpoint.
This data format consists of two co-registered images: a standard RGB color image providing appearance and texture details, and a corresponding depth map where each pixel value encodes the camera-space depth of the corresponding scene point.
Formally, it can be defined as a function mapping a 2D image grid $\Omega$ to a 4D space combining color and depth:
\begin{equation}
I_{rgbd}: \Omega \rightarrow \mathbb{R}^3 \times \mathbb{R}_{\geq 0},
\end{equation}
where the first three channels represent RGB color and the fourth channel encodes the depth of the corresponding scene point. In practice, RGB-D data is stored as a tensor $I \in \mathbb{R}^{H \times W \times 4}$.

The primary advantage of RGB-D data lies in its efficiency and accessibility.
It is captured by low-cost sensors such as Microsoft Kinect or Intel RealSense, and its image-like grid structure enables the direct application of powerful 2D convolutional neural networks (CNNs) to tasks such as semantic segmentation and object detection~\cite{dai2016r,couprie2013indoor,vu2019dada}.
This fusion of color and depth provides a rich, multi-modal input for scene understanding.
However, a significant limitation is its view-dependent nature; an RGB-D image represents only a partial view of the 3D world, suffering from occlusions and providing no information about surfaces hidden from the camera.
Despite this, its real-time acquisition and processing capabilities make it indispensable in robotics for navigation~\cite{desouza2002vision} and grasping~\cite{bicchi2000robotic,saxena2006robotic,zhao2025towards}, in augmented reality (AR) for scene interaction, and as a foundational input for 3D reconstruction pipelines that fuse multiple frames into a complete model~\cite{newcombe2011kinectfusion}.







\subsection{Voxel Grids}
\label{sec:rep_3d_voxel}

The voxel grid is a foundational 3D representation that discretizes space into a uniform grid of volumetric elements, or \textit{voxels} (a portmanteau of ``volume'' and ``pixel'').
Analogous to a pixel in 2D, a voxel grid can be modeled as a function mapping discrete 3D coordinates to a set of attributes:
\begin{equation}
G: \{1, \dots, D\} \times \{1, \dots, H\} \times \{1, \dots, W\} \rightarrow \mathcal{A},
\end{equation}
where $D, H, W$ are the grid dimensions and $\mathcal{A}$ can represent occupancy, color, density, or material properties, thereby enabling a detailed volumetric description of an object.

A key advantage of this representation is its ability to model complex geometries and internal structures, which are often challenging for surface-based methods like polygon meshes to capture.
The conceptual simplicity of voxels, which can be seen as the three-dimensional extension of 2D pixels, facilitates straightforward implementation and manipulation.
Consequently, voxel grids are natively suited for applications involving inherently volumetric data, such as medical imaging (\eg, MRI and CT scans)~\cite{cciccek20163d}, where biological tissues are naturally captured as a 3D array of values.
They are also widely adopted in computer graphics and gaming for tasks like procedural generation and creating destructible environments, as popularized by games like \textit{Minecraft}.

\subsection{Point Clouds}
\label{sec:rep_3d_point}

Point clouds offer a direct and unstructured 3D representation, defined as a set of points in 3D space:
\begin{equation}
P = \{p_i\}_{i=1}^N, \quad p_i = (x_i, y_i, z_i, \mathbf{a}_i) \in \mathbb{R}^3 \times \mathbb{R}^C.
\end{equation}
Each point $p_i$ consists of its geometric coordinates $(x_i, y_i, z_i)$ and an optional $C$-dimensional attribute vector $\mathbf{a}_i$ encoding properties such as color, intensity, or surface normals. They are the native output of active 3D scanning technologies like LiDAR (Light Detection and Ranging) and structured-light scanners~\cite{behroozpour2017lidar}, as well as passive methods like photogrammetry~\cite{schonberger2016structure}.

The primary strength of point clouds lies in their raw and unstructured nature; they represent geometry without imposing any explicit topological connectivity (unlike meshes) or spatial discretization (unlike voxels).
This makes them a faithful and often highly precise representation of the captured data.
However, this very lack of inherent structure poses significant challenges for downstream processing.
The unordered and potentially sparse distribution of points makes tasks like feature extraction, segmentation, and reconstruction non-trivial, necessitating specialized algorithms and deep learning architectures (\eg, PointNet~\cite{qi2017pointnet} and its successors~\cite{qi2017pointnetplus}) that are designed to operate directly on irregular point sets.


\subsection{Meshes}
\label{sec:rep_3d_mesh}

A mesh is a canonical surface representation in computer graphics, defining an object's boundary with a set of geometric and topological primitives. Formally, a triangle mesh is a tuple $M = (V, F)$, where
\begin{equation}
V = \{v_i\}_{i=1}^{|V|} \subset \mathbb{R}^3, \; F\subset V \times V \times V.
\end{equation}
Here, $V$ is a set of vertices representing points in 3D space, and $F$ is a set of faces (typically triangles) that define connectivity between vertices.

The key characteristic of a mesh is its explicit topological structure, which encodes the connectivity between these elements. This structure strikes an effective balance between descriptive power and computational efficiency. Its prevalence is due to several key advantages: meshes can represent complex, smooth surfaces with arbitrary fidelity through adaptive tessellation (\ie, varying the polygon density); moreover, this representation is natively supported by modern Graphics Processing Units (GPUs), enabling highly efficient real-time rendering. The explicit connectivity also facilitates the straightforward computation of essential surface properties like normals and curvature, which are critical for realistic lighting and physical analysis. However, as a purely surface-based representation, standard meshes are ill-suited for modeling volumetric or internal structures, and representing objects with complex or changing topologies can be algorithmically challenging.








\subsection{Implicit Representations}
\label{sec:rep_3d_implicit}

Implicit representations define 3D geometry not as a discrete set of primitives (\eg, points or polygons), but as the level set of a continuous function:
\begin{equation}
f: \mathbb{R}^3 \to \mathbb{R}.
\end{equation}
The surface $\mathcal{S}$ is implicitly defined as the set of points where the function evaluates to a specific value, typically zero:
\begin{equation}
\mathcal{S} = \{\mathbf{x} \in \mathbb{R}^3 | f(\mathbf{x}) = 0\}.
\end{equation}

A particularly powerful and popular form of this is the \textbf{Signed Distance Function (SDF)}.
An SDF maps any spatial coordinate $\mathbf{x}$ to the closest distance to the surface, where the sign indicates whether the point is inside ($f(\mathbf{x}) < 0$) or outside ($f(\mathbf{x}) > 0$) the object.
This formulation inherently guarantees watertight, closed surfaces and provides a smooth, continuous description of the geometry, making it resolution-agnostic.
A key advantage of implicit fields is their ability to elegantly handle complex and dynamic topologies; operations like merging, splitting, or creating holes do not require complex and error-prone re-meshing procedures.
This makes them ideal for physical simulations, such as fluid dynamics and shape manipulation.

This paradigm has experienced a major resurgence with the rise of \textbf{neural fields} or \textbf{neural implicit representations}. Instead of storing the function in a grid, a neural network is trained to approximate it~\cite{mescheder2019occupancy} (\eg, DeepSDF~\cite{park2019deepsdf}).
This approach is memory-efficient and can learn continuous, high-fidelity shape priors from data.
Extensions such as Neural Radiance Fields (NeRF) further generalize implicit fields to jointly model geometry and appearance, and are discussed in detail in Section~\ref{sec:rep_3d_nerf}.
To visualize or use implicit surfaces in traditional rendering pipelines, an explicit mesh representation is often extracted via algorithms like Marching Cubes~\cite{lorensen1998marching}.






\begin{figure}[!ht]
  \centering
  \includegraphics[width=\linewidth]{figures/part2/3D_representations_2.pdf}
  % \vspace{-7pt}
  \caption{Neural and hybrid 3D representations: (a) NeRF, (b) 3DGS, and (c) hybrid methods (DMTet and Tri-plane).}
  \label{fig:p2_s4_2_3d_representation_2}
\end{figure}


\subsection{Neural Radiance Fields}
\label{sec:rep_3d_nerf}

Neural Radiance Fields (NeRF) represent a paradigm shift in novel view synthesis and implicit 3D scene representation \cite{mildenhall2021nerf}.
The core innovation is to represent a static scene as a continuous function with a 5D input domain (3D location $+$ 2D viewing direction), parameterized by a simple Multi-Layer Perceptron (MLP).
This neural network, denoted as $F_\Theta$, maps a 3D spatial coordinate $\mathbf{x}=(x,y,z)$ and a 2D viewing direction $\mathbf{d}=(\theta, \phi)$ to a volume density $\sigma$ and an emitted color $\mathbf{c}=(r,g,b)$.
A crucial design choice is that density $\sigma$ depends only on the location $\mathbf{x}$, while the color $\mathbf{c}$ depends on both location and viewing direction $(\mathbf{x}, \mathbf{d})$.
This allows NeRF to capture complex, non-Lambertian surface effects such as specular reflections and translucency.

To render an image from a novel viewpoint, rays are cast through the scene, and their colors are computed by integrating the predicted color and density values along each ray using principles of classical volume rendering and numerical quadrature.
By optimizing the network's weights to minimize the photometric reconstruction error between rendered and ground-truth images, NeRF can learn a highly detailed and coherent representation of the scene from only a sparse set of posed input images.
The resulting state-of-the-art photorealism has established NeRF as a foundational technique, inspiring a vast body of follow-up research aimed at improving its speed, scalability, and applicability to dynamic scenes~\cite{muller2022instant,barron2021mip,barron2023zip,multinerf2022}.

\subsection{3D Gaussian Splatting}
\label{sec:rep_3d_gaussian}

3D Gaussian Splatting (3DGS) has recently emerged as a groundbreaking alternative to neural radiance fields, achieving state-of-the-art quality for novel view synthesis while enabling real-time rendering \cite{kerbl20233d}.
Unlike the implicit representation of NeRF, 3DGS represents a scene as a set of $N$ anisotropic 3D Gaussian primitives, each parameterized by:
\begin{itemize}[leftmargin=2em,itemsep=0.35em]
    \item \textbf{Position} $\boldsymbol{\mu} \in \mathbb{R}^3$: the 3D center of the Gaussian;
    \item \textbf{Covariance} $\boldsymbol{\Sigma} \in \mathbb{R}^{3\times 3}$: shape and orientation, allowing thin planes through to volumetric wisps;
    \item \textbf{Color}: Spherical Harmonics (SH) coefficients capturing view-dependent appearance;
    \item \textbf{Opacity} $\alpha \in [0,1]$: opacity ($0$ is transparent and $1$ is opaque).
\end{itemize}
The density contributed by a single primitive at a query point $\mathbf{x} \in \mathbb{R}^3$ is
\begin{equation}
G(\mathbf{x}; \boldsymbol{\mu}, \boldsymbol{\Sigma}) = \exp\!\left(-\tfrac{1}{2} (\mathbf{x}-\boldsymbol{\mu})^{\top} \boldsymbol{\Sigma}^{-1} (\mathbf{x}-\boldsymbol{\mu}) \right).
\end{equation}
$\boldsymbol{\Sigma}$ is decomposed into a scaling matrix $\mathcal{S}$ and rotation matrix $\mathcal{R}$, such that
\begin{equation}
\boldsymbol{\Sigma} = \mathcal{R}\mathcal{S}\mathcal{S}^T\mathcal{R}^T,
\end{equation}
with per-primitive scaling and rotation stored as $s \in \mathbb{R}^{N \times 3}$ and quaternions $r \in \mathbb{R}^{N \times 4}$. The key to 3DGS's real-time performance lies in the differentiable splatting rasteriser introduced together with the representation~\cite{kerbl20233d}, which extends earlier differentiable rasterization work on meshes~\cite{liu2019soft}.
Instead of costly per-pixel ray marching, 3DGS projects (or ``splats'') these 3D Gaussians onto the 2D image plane, creating 2D Gaussian splats.
These splats are then sorted and alpha-blended in a highly-parallelized, feed-forward process that is exceptionally well-suited for modern GPU architectures.
Differentiable splatting applies a viewing transform $W$ and Jacobian $J$ to compute the transformed covariance $\boldsymbol{\Sigma}^{\prime} = JW\boldsymbol{\Sigma} W^TJ^T$, enabling novel view rendering. Each pixel color $\mathcal{C}$ is obtained by blending the $N$ overlapping primitives along the ray:
\begin{equation}
\mathcal{C} = \sum_{i=1}^{N}c_i \alpha_i \prod_{j=1}^{i-1} (1-\alpha_j),
\end{equation}
where $c_i$ and $\alpha_i$ denote the color and opacity of the $i$-th primitive.
The entire model, initialized from a sparse point cloud (typically from Structure-from-Motion), is optimized end-to-end against training images.
The optimization process also includes an adaptive density control mechanism to dynamically add, split, or prune Gaussians to best represent the scene's geometry.
This combination of an explicit, flexible primitive with a fast, differentiable rasterizer has set a new standard for performance and quality in 3D scene reconstruction.





\subsection{Hybrid Methods}
\label{sec:rep_3d_hybrid}
To harness the complementary strengths of different 3D representations, hybrid methods have emerged as a powerful paradigm.
These approaches combine two or more representations, integrating the high-fidelity and continuous nature of implicit functions with the computational efficiency and explicit structure of primitives like meshes or voxels.
This section reviews two influential hybrid techniques: Deep Marching Tetrahedra (DMTet) and Tri-plane.

\paragraph{Deep Marching Tetrahedra (DMTet).}
Deep Marching Tetrahedra (DMTet)~\cite{shen2021deep} is an innovative hybrid technique that bridges learned implicit fields with explicit polygonal meshes. 
Its core concept involves using a neural network to predict an implicit representation of a shape, such as a Signed Distance Function (SDF)~\cite{vicini2022differentiable} or an occupancy field. 
A differentiable version of the marching tetrahedra algorithm subsequently extracts a high-quality, explicit triangle mesh surface from this implicit representation.

The primary advantage of DMTet lies in its end-to-end differentiability. 
The entire pipeline, from the implicit field prediction to the final mesh generation, allows gradients to flow from a 2D-image-based loss function back to the 3D shape parameters.
This enables direct optimization of the mesh geometry using 2D supervision, leading to highly detailed and accurate models.
Furthermore, DMTet can learn to produce an adaptive mesh, allocating higher tetrahedral density to regions of fine detail, thus optimizing computational resources while maintaining high fidelity.


\paragraph{Tri-plane.}
The Tri-plane representation~\cite{chan2022efficient} offers an efficient, factorized approach to modeling 3D scenes and objects. Instead of encoding information in a dense 3D volume, it decomposes the 3D space into three orthogonal 2D feature planes (typically aligned with the XY, YZ, and XZ axes). This design significantly reduces the memory and computational complexity associated with 3D representations.

The working principle is to derive the feature for any 3D point by projecting it onto the three 2D planes, querying the features at the corresponding projected coordinates, and aggregating them (\eg, through summation or concatenation). 
This aggregated feature vector is then fed into a small MLP to decode the final property of the point, such as its color, density, or SDF value.
This ``explicit-grid-plus-implicit-decoder'' structure is the essence of its hybrid nature.
It benefits from the efficiency of 2D grid-based feature storage, which can be processed by powerful 2D CNNs, while avoiding the prohibitive memory cost of a full 3D grid.


\begin{table}[!t]
\centering
\caption{Representative methods for each 3D representation family, chosen to show its historical progression and dominant design choices.}
\label{tab:p2_s4_3d_representative_methods}
% \small
\fontsize{9}{11}\selectfont
\setlength{\tabcolsep}{4pt}
\begin{tabular}{L{0.19\textwidth}L{0.4\textwidth} L{0.4\textwidth}}
\toprule
\textbf{Representation} & \textbf{Representative Methods} & \textbf{Key Idea / Contribution} \\
\midrule
\textbf{RGB-D Image} & KinectFusion~\cite{newcombe2011kinectfusion}, DADA~\cite{vu2019dada}, SemanticFusion~\cite{mccormac2017semanticfusion} & Fuses color and per-pixel depth for real-time reconstruction and semantic understanding. \\
\textbf{Voxel Grid} & VoxNet~\cite{maturana2015voxnet}, 3D U-Net~\cite{cciccek20163d}, OGN~\cite{tatarchenko2017octree} & Direct 3D CNNs on regular grids; octree variants reduce memory for high resolution. \\
\textbf{Point Cloud} & PointNet~\cite{qi2017pointnet}, PointNet++~\cite{qi2017pointnetplus}, Point Transformer~\cite{zhao2021point}, PointNeXt~\cite{qian2022pointnext} & Permutation-invariant and attention-based architectures for unordered point sets. \\
\textbf{Mesh} & MeshCNN~\cite{hanocka2019meshcnn}, Pixel2Mesh~\cite{wang2018pixel2mesh}, AtlasNet~\cite{groueix2018papier} & Learning on mesh edges and single-image mesh generation from data. \\
\textbf{Implicit (SDF/occupancy)} & Occupancy Networks~\cite{mescheder2019occupancy}, DeepSDF~\cite{park2019deepsdf}, IM-Net~\cite{chen2019imnet}, NeuS~\cite{wang2021neus} & Neural networks encode continuous surfaces via occupancy or signed-distance fields. \\
\textbf{NeRF} & NeRF~\cite{mildenhall2021nerf}, Plenoxels~\cite{fridovich2022plenoxels}, Instant-NGP~\cite{muller2022instant}, Mip-NeRF 360~\cite{barron2022mip} & 5D radiance fields; later works accelerate training and reduce aliasing. \\
\textbf{3DGS} & 3D Gaussian Splatting~\cite{kerbl20233d}, Mip-Splatting~\cite{yu2024mip}, Scaffold-GS~\cite{lu2024scaffold}, 4DGS~\cite{wu20244d} & Anisotropic Gaussian primitives with differentiable rasterization for real-time rendering. \\
\textbf{Hybrid} & DMTet~\cite{shen2021deep}, Tri-plane~\cite{chan2022efficient}, NeRFactor~\cite{zhang2021nerfactor} & Combines explicit grids or meshes with compact implicit decoders. \\
\bottomrule
\end{tabular}
\end{table}

% \subsection{Comparison, Applications, and Challenges}
% \label{sec:rep_3d_summary}

% Table~\ref{tab:p2_s4_1_3d_comparison} highlights the complementary trade-offs among the eight representations surveyed above, while Table~\ref{tab:p2_s4_3d_representative_methods} collects representative methods.
% RGB-D images, voxels, point clouds, and meshes remain the workhorses of applications that prize acquisition simplicity, volumetric modeling, raw geometric fidelity, or GPU rendering, respectively.
% Neural representations, NeRF and 3DGS in particular, have shifted the frontier toward photorealistic novel-view synthesis from sparse images, but they do so at the cost of editability and, in the case of NeRF, rendering speed.
% Hybrid methods attempt to occupy the middle ground by pairing explicit primitives with compact neural decoders.

% These trade-offs manifest differently across application domains.
% In \textbf{robotics}, RGB-D sensors and point clouds dominate perception and grasping pipelines; meshes and implicit fields are common for collision checking and simulation.
% In \textbf{autonomous driving}, LiDAR point clouds are the canonical input, while voxel grids and bird's-eye-view feature maps serve as intermediate representations for detection and prediction.
% In \textbf{graphics and virtual production}, meshes remain the standard for real-time rendering, while NeRF and 3DGS are increasingly used for view synthesis, digital humans, and scene relighting.
% In \textbf{medical imaging}, voxel grids are natural for CT and MRI data, and implicit surfaces are used for organ and tumor segmentation.

% Several open challenges cut across all of these representations.
% The most pervasive is \textbf{scalability}: as scenes grow from single objects to rooms, buildings, and outdoor environments, memory and computation become bottlenecks for voxels, point clouds, and dense Gaussian primitives alike.
% Even where memory is not the limiting factor, \textbf{editability and composability} lag behind: neural fields and 3DGS are difficult to edit in the intuitive way meshes or voxels allow, and combining several reconstructed assets into a coherent scene remains nontrivial.
% Time compounds both problems, since most representations surveyed here assume static geometry, and extending them to \textbf{dynamic and deformable scenes} (articulated objects, fluids, general deformable media) is still an active research direction.
% Underlying all of these is \textbf{cross-representation conversion} (extracting meshes from neural fields, compressing point clouds into Gaussians, or turning implicit surfaces into editable primitives), which remains a practical barrier to deploying these representations in downstream pipelines.
% Progress will therefore depend not only on better individual representations, but on principled ways to compose and convert between them.



\section{4D}
\label{sec:rep_4d}

\FramedInitial{T}{he} 4D representations extend the 3D representations of Section~\ref{sec:rep_3d} into the temporal dimension, modeling not only the geometry and appearance of a scene at a single instant but also how both evolve over time.
They are the substrate for tasks that involve dynamic content (moving objects, deformable bodies, fluid simulation, novel-view synthesis of casually captured videos), where a sequence of independent 3D reconstructions would either fail to maintain temporal coherence or be prohibitively expensive to recompute frame by frame.

We organize this section around five 4D representations (Figure~\ref{fig:p2_s4_3_4d_representations}), each obtained by adding time to one of the 3D representations of Section~\ref{sec:rep_3d}: RGBD video, point clouds, meshes, neural radiance fields, and deformable 3D Gaussian Splatting.
The voxel and implicit-field families of Section~\ref{sec:rep_3d} admit analogous time extensions but are less commonly used in the 4D literature, and we therefore omit them as standalone subsections here.



\begin{figure}[!th]
    \centering
    \includegraphics[width=\linewidth]{figures/part2/4D_representations.pdf}
     % \vspace{-7pt}
    \caption{Five 4D representations, each a 3D primitive evolving over time: (a) RGBD video, (b) 4D point cloud, (c) deformable mesh, (d) dynamic NeRF, and (e) deformable 3DGS. The RGBD example uses changing camera poses across time to observe the scene from different viewpoints; the top-view camera icons use the same azimuths as the paired RGB and depth renderings. The point-cloud panel distinguishes independent frame-wise samples from persistent point trajectories. Mesh colors indicate vertex displacement; re-extraction is an alternative to preserving vertex and face identities. In (d), rows are views and columns are times. In (e), $R$ and $S$ abbreviate the time-dependent $R_i(t)$ and $S_i(t)$; constant opacity and color describe the illustrated parameterization. All scenes are controlled synthetic illustrations, not experimental reconstructions.}
    \label{fig:p2_s4_3_4d_representations}
\end{figure}



\subsection{RGBD Video}
\label{sec:rep_4d_rgbd}
RGBD videos $\mathcal{V}=\{(I_t, D_t)\}_{t=1}^T$, with $I_t \in \mathbb{R}^{H \times W \times 3}$ and $D_t \in \mathbb{R}^{H \times W}$, are a straightforward extension of RGBD images (see Section~\ref{sec:rep_3d_rgbd}) to the temporal domain.
Equivalently, they can be stored as a 4D tensor $\mathcal{V} \in \mathbb{R}^{T \times H \times W \times 4}$, with the last channel holding depth.
RGBD videos can be captured using depth sensors, such as Microsoft Kinect or Intel RealSense, which provide synchronized RGB and depth data.
These videos are commonly used in applications such as motion capture, gesture recognition, and augmented reality.
There are also many learning-based methods~\cite{wang2023neural,hu2025depthcrafter,chen2025video} that directly predict consistent RGBD videos from monocular video sequences
$\mathcal{I}=\{I_t\}_{t=1}^T$.
However, such a representation may struggle to handle complex dynamic scenes with significant occlusions or topology changes.
If we want to further reconstruct a unified 4D scene in a canonical space, we may also need to provide or estimate camera poses
$
\mathcal{P}=\{P_t\}_{t=1}^T
$
for each frame, as well as point trajectories over time to establish temporal correspondences, \eg via joint pose-trajectory optimization in recent video-based SfM systems~\cite{li2025megasam}.






\subsection{Point Cloud}
\label{sec:rep_4d_point}
Point clouds are another popular representation for 3D scenes (see Section~\ref{sec:rep_3d_point}), and they can be extended to 4D by incorporating temporal information.
A 4D point cloud can be represented as $\mathcal{P}=\{(\mathbf{p}_i(t), \mathbf{c}_i(t))\}_{i=1}^N$, where $\mathbf{p}_i(t) \in \mathbb{R}^3$
is the 3D position of point $i$ at time step $t$, and $\mathbf{c}_i(t) \in \mathbb{R}^3$ is the corresponding color.
To map an uncalibrated, continuous video sequence $\mathcal{I}=\{I_t\}_{t=1}^T, I_t\in \mathbb{R}^{H \times W \times 3}$, to a 4D point cloud representation, there are generally two types of representations:
(1) a sequence of point clouds $\mathcal{P}=\{\mathcal{P}_t\}_{t=1}^T$ in a fixed coordinate system, where each $\mathcal{P}_t=\{(\mathbf{p}_i(t), \mathbf{c}_i(t))\}_{i=1}^{H \times W}$ is an independent point cloud at time step $t$;
(2) a unified point cloud $\mathcal{P}=\{(\mathbf{p}_i(t), \mathbf{c}_i(t))\}_{i=1}^{N}$, where $N$ is the number of persistent point identities, and each point $i$ has a trajectory over time.
The first representation is easier to obtain, as it can be directly derived from the depth maps estimated from each frame using monocular depth estimation methods~\cite{li2025megasam}.
However, it may suffer from temporal inconsistencies and a lack of point correspondences across frames.
The second representation, in contrast, requires explicit temporal correspondences (either ground-truth tracks, optical-flow lifting~\cite{teed2020raft}, or jointly optimized point trajectories~\cite{li2025megasam}), and is therefore harder to acquire, but it offers persistent point identities and consistent trajectories that are useful for downstream 4D reconstruction, motion analysis, and dynamic NeRF/3DGS initialization.

\subsection{Mesh}
\label{sec:rep_4d_mesh}
We can extend the mesh representation (see Section~\ref{sec:rep_3d_mesh}) to 4D by representing a sequence of 3D meshes over time.
A 4D mesh is a sequence $\{\mathcal{M}_t\}_{t=1}^T$ with $\mathcal{M}_t = (\mathbf{V}_t, \mathbf{F}_t)$, where $\mathbf{V}_t \in \mathbb{R}^{N \times 3}$ and $\mathbf{F}_t \in \mathbb{N}^{M \times 3}$ are the vertices and faces of the mesh at time $t$, respectively.
To reconstruct a 4D mesh using deep learning, researchers typically sample a set of points on the mesh surface at the original time step $t=0$, and then learn a deformation field to model the motion of these points over time~\cite{cao2024motion2vecsets}.
Then, the deformed points at each time step can be used to reconstruct watertight manifold meshes using
Marching Cubes~\cite{lorensen1998marching} or Poisson Surface Reconstruction~\cite{kazhdan2006poisson}.








\subsection{NeRF}
\label{sec:rep_4d_nerf}
As discussed in Section~\ref{sec:rep_3d_nerf}, NeRF~\cite{mildenhall2021nerf} is a continuous, volume-based 3D representation that maps a 3D position $\mathbf{x} \in \mathbb{R}^3$ and a 2D viewing direction $\mathbf{d} \in \mathbb{S}^2$ to a color $\mathbf{c} \in \mathbb{R}^3$ and a volume density $\sigma \in \mathbb{R}_{\geq 0}$ via a learned function $F_\Theta:(\mathbf{x}, \mathbf{d}) \rightarrow (\mathbf{c}, \sigma)$.
This function is typically implemented using a neural network, which is trained on a set of input images and their corresponding camera poses.
To extend NeRF to dynamic scenes, we can introduce a temporal dimension $t$ to the input of the function, resulting in a 4D representation
$F_\Theta: (\mathbf{x}, \mathbf{d}, t) \rightarrow (\mathbf{c}_t, \sigma_t)$.
This allows the model to capture changes in the scene over time, enabling the rendering of novel views at different time steps.
In practice, the temporal dimension can be represented in various ways.
For example, a deformation network can take an observed position $\mathbf{x}$ and time $t$ as input and predict a displacement $\Delta\mathbf{x}(\mathbf{x},t)$ that maps it to a canonical position $\mathbf{x}_c=\mathbf{x}+\Delta\mathbf{x}(\mathbf{x},t)$~\cite{pumarola2021d}.
Alternatively, the deformation can be parameterized using a dense
$\text{SE}(3)$ field~\cite{park2021nerfies}, which models the rigid transformations (rotation and translation) of points in the scene over time.
Other approaches~\cite{gao2021dynamic} may also incorporate additional latent codes to represent different time steps, enabling the model to capture more complex temporal variations.







\subsection{Deformable 3D Gaussian Splatting}
\label{sec:rep_4d_gaussian}
Given a static scene representation $\mathcal{G}=\{(\alpha_i, \boldsymbol{\mu}_i, \boldsymbol{\Sigma}_i, c_i)\}_{i=1}^G$, which contains $G$ 3D Gaussians~\cite{kerbl20233d} (see Section~\ref{sec:rep_3d_gaussian}),
where $\alpha_i \in [0, 1]$ is the opacity, $\boldsymbol{\mu}_i\in\mathbb{R}^3$ is the mean (3D position), $\boldsymbol{\Sigma}_i\in\mathbb{R}^{3\times 3}$ is the covariance matrix, and $c_i: \mathbb{S}^2 \rightarrow \mathbb{R}^3$ is the color function that maps a viewing direction to an RGB color, we extend it to a \emph{deformable} 3D Gaussian Splatting representation for dynamic scenes by letting both the mean and the covariance vary in time: $\mathcal{G}(t)=\{(\alpha_i, \boldsymbol{\mu}_i(t), \boldsymbol{\Sigma}_i(t), c_i)\}_{i=1}^G$~\cite{wu20244d}.
The mean is parameterized by a deformation field, $\boldsymbol{\mu}_i(t)=\boldsymbol{\mu}_i(0)+\Delta\boldsymbol{\mu}_i(t)$, where $\boldsymbol{\mu}_i(0)$ is the initial position at $t=0$ and $\Delta\boldsymbol{\mu}_i(t)$ is the time-dependent offset.
The covariance is decomposed in the same way as in the static case, $\boldsymbol{\Sigma}_i(t)=\mathcal{R}_i(t)\mathcal{S}_i(t)\mathcal{S}_i(t)^{T}\mathcal{R}_i(t)^{T}$, with the rotation $\mathcal{R}_i(t)$ and scaling $\mathcal{S}_i(t)$ themselves outputs of the deformation field at time $t$.
The opacity $\alpha_i$ and color function $c_i$ are typically kept constant over time.
In practice, this means that the appearance of the scene does not change significantly over time, while the latest work~\cite{zhang2025gaussian} also explores time-varying colors and opacities.
This representation allows for efficient rendering of dynamic scenes by accounting for the temporal changes in the positions and shapes of the 3D Gaussians.




























\section{Language}
\label{sec:rep_language}

\FramedInitial{U}{nlike} the visual and geometric representations of the previous sections, which encode space explicitly through pixels, points, or surfaces, language represents space \emph{symbolically}, naming entities and stating the relations among them (\eg, ``the chair is next to the table'' or ``the store is across the street from the park'') and thereby capturing relational, topological, and functional structure that metric representations leave implicit. 
The question this section pursues is not how language \emph{describes} space, which is self-evident, but whether models trained only on language come to \emph{encode} it. 
Answering that requires first understanding how language representations are learned at all, so the section moves from the evolution of language representation learning, from masked-language pretraining to large autoregressive models, to the more surprising finding that such sequence-only training induces latent, spatially structured world models rather than mere text statistics. 
The scope here is language in isolation; how language is \emph{aligned} with visual, 3D, and audio modalities is deferred to Chapter~\ref{chapter:cross_modality}.



\begin{figure}[!ht]
    \centering
    \includegraphics[width=\linewidth]{figures/part2/language_representation.pdf}
    \caption{Spatial information implicit in language-only representations. \textit{Left}: natural-language descriptions of spatial relations. \textit{Right}: a probe on Othello-GPT hidden activations recovers the board state from move sequences alone.}
    \label{fig:p2_s4_1_language-representation}
\end{figure}


\subsection{The Evolution of Language Representation Learning}
\label{subsec:p-2-sec-4-the-evolution-of-language-representation-learning}

Language representation learning has evolved rapidly over the past decade, through masked-language pretraining (BERT~\cite{devlin2019bert} and RoBERTa~\cite{liu2019roberta}), permutation-based autoregressive pretraining (XLNet~\cite{yang2019xlnet}), denoising sequence-to-sequence pretraining (T5~\cite{raffel2020exploring}), and autoregressive large language models trained purely on next-token prediction, exemplified by GPT-2 and GPT-3~\cite{radford2019language,brown2020language}.
Across this progression, a recurring observation is that representation quality scales with model size, data size, and the sophistication of the pretraining objective~\cite{srivastava2022beyond}, and that pretrained representations encode richer structure than the surface text statistics they were trained on~\cite{tenney2019bert,hewitt2019structural}.
This last property, that language-only training produces representations carrying structure beyond text, is what makes language representations load-bearing for spatial intelligence, as the next subsection makes explicit.

\subsection{Language Representation Learning Induces Spatial Intelligence}
\label{subsec:p-2-sec-4-language-representation-learning-induces-spatial-intelligence}

Although LLMs are optimized for sequence prediction, recent work suggests that they can develop structured internal representations that resemble world models~\cite{li2021implicit,patel2022mapping,abdou2021can}. 
One particularly pertinent demonstration of this on spatial intelligence is Othello-GPT, a transformer trained solely on sequences of legal moves from the board game Othello~\cite{li2023emergent}.

Despite receiving no prior knowledge of the board layout or game rules, Othello-GPT achieves near-perfect accuracy in predicting legal moves (actions in a physical space). 
Probing experiments reveal that nonlinear~\cite{alain2016understanding,tenney2019bert,belinkov2016probing} classifiers can reliably reconstruct the entire board state from hidden activations~\cite{li2023emergent}. 
Crucially, interventional studies show that manipulating internal representations of specific tiles alters the model’s predictions in a way consistent with counterfactual board states, providing causal evidence of spatial encoding \cite{li2023emergent}.

These findings align with broader research indicating that LMs acquire implicit representations of abstract or grounded concepts. 
For example, models have been shown to encode syntactic structure~\cite{hewitt2019structural}, color spaces \cite{abdou2021can}, and even logical state-tracking~\cite{li2021implicit}. 
Work on chess transcripts similarly shows that LMs can track board positions without explicit supervision~\cite{toshniwal2021learning}. 
Compared to systems like AlphaZero, which are built with strong inductive biases and explicit access to board states~\cite{silver2018general,mcgrath2021acquisition}, Othello-GPT highlights the surprising extent to which sequence-only training can lead to emergent spatial reasoning.

In sum, it has been demonstrated that spatially structured representations can emerge in LMs trained only on symbolic sequences. 
This suggests that autoregressive sequence prediction is sufficient, in certain domains, to induce latent world models that encode spatial relations and can causally drive action predictions.








\section{Audio}
\label{sec:rep_audio}



\FramedInitial{A}{udio} is fundamentally tied to spatial information. 
Unlike vision, which directly captures the geometric layout of a scene, acoustic signals encode spatial structure only \emph{implicitly}, through the physical transformations imposed on sound during propagation.
Recovering space from sound is therefore a matter of reading these transformations backward: from the perceptual cues that human hearing exploits, to the computational representations that formalize them, to the learned models and cross-modal couplings that ground sound in the surrounding geometry. 
This section follows that arc.
Human auditory perception interprets the resulting acoustic cues, shaped by these transformations, to infer the spatial layout of the environment.
The primary cues are binaural: the \emph{Interaural Time Difference (ITD)} and the \emph{Interaural Level Difference (ILD)}, both arising from the geometry of the head and the spatial separation of the ears~\cite{macpherson2002listener,middlebrooks2015sound}.
ITD reflects the difference in arrival time between the two ears and is highly reliable for localizing low-frequency sounds whose long wavelengths exceed head size.
Conversely, ILD captures sound pressure level differences caused by head-shadowing effects and serves as a dominant cue for high-frequency localization where such attenuation is more pronounced.
However, binaural cues alone cannot fully resolve depth, elevation, or front-back ambiguities.

To address these limitations, the auditory system further exploits \emph{monaural spectral cues} and \emph{environmental reverberation}.
Direction-dependent spectral filtering produced by the torso, head, and pinnae introduces characteristic notches and peaks, enabling inference of elevation and disambiguation of positions on the cone of confusion~\cite{gardner19983}.
In addition, the room impulse response (RIR) shaped by early reflections, late reverberation, and surface materials encodes information about source distance, enclosure geometry, and environmental structure~\cite{kuttruff2016room,blauert1997spatial}.
These monaural and environmental cues complement binaural information, jointly supporting robust spatial inference under naturalistic conditions.

Computational models formalize these perceptual cues through structured acoustic representations.
Head-related transfer functions (HRTFs)~\cite{wightman1989headphone} capture listener-specific spectral filtering and are foundational to perceptual research and spatial audio rendering.
In parallel, spherical-harmonics–based ambisonics representations~\cite{zotter2019ambisonics} provide a scene-centric encoding of the sound field around a point in space, with controllable angular resolution.
Together, these representations describe how sound interacts with the body and the environment, enabling principled approaches to spatial reasoning, source localization, and scene reconstruction.
Taken together, the ability of the auditory system to integrate binaural, monaural, and environmental cues demonstrates that sound carries rich implicit structure about the spatial configuration of sources and the surrounding environment.



\begin{figure}[!ht]
    \centering
    \includegraphics[width=\linewidth]{figures/part2/audio_representation.pdf}
    \caption{Three families of acoustic spatial cues: binaural cues (ITD, ILD), monaural spectral cues (direction-dependent filtering by torso, head, and pinnae, summarized by the HRTF), and environmental reverberation (room impulse response).}
    \label{fig:p2_s4_4_audio-representation}
\end{figure}

Recent learning-based approaches aim to capture spatial structure directly from raw audio, going beyond analytic models.
Neural encoders trained on spatially annotated or simulated data learn to infer direction, distance, and environmental geometry from multi-channel waveforms or spectrograms~\cite{chen2022soundspaces,yuksel2025general}.
Learning frameworks further exploit the relation between sound and space: SoundSpaces~\cite{chen2020soundspaces} provides audio-visual simulation for embodied navigation, while VisualEchoes~\cite{gao2020visualechoes} uses image--echo correspondences for self-supervised spatial representation learning.
These learned representations encode not only spatial acoustics but also grounding in surrounding geometry, forming a key component of spatially aware representation.


In real-world environments, spatial audio interacts closely with other sensory modalities.
Audio–visual models align acoustic cues with visual observations to localize sounding objects and understand where events occur~\cite{morgado2018self,afouras2020self,morgado2020learning}.
Embodied agents further leverage auditory signals to perceive layout and navigate toward acoustic targets, as demonstrated in audio–visual navigation~\cite{chen2020soundspaces} and egocentric audio reasoning~\cite{zheng2024bat,biswas2025owl}.
Through such cross-modal grounding, audio provides complementary spatial cues that enhance robustness under visual occlusion and enable reasoning about events beyond the agent's field of view.
This makes audio a powerful complementary modality for spatial representation, especially in scenarios where visual information is partial, noisy, or unavailable.












\section{Domain-Specific Representations}
\label{sec:rep_other}

\FramedInitial{T}{he} representations surveyed so far are organized by modality, each general enough to span many applications. 
A large body of practice, however, develops representations tailored to a specific sensor or downstream task rather than to a modality as such, and these are collected here. 
What unites them is not a shared signal but a shared design logic: each encodes the inductive biases of its data and trades spatial fidelity against semantic compressibility in the way its sensors and target tasks demand. 
This section walks through five such domains: embodied robotics, virtual reality, geospatial and remote sensing, clinical and scientific imaging, and indoor intelligent assistants, examining in each the representations that have become canonical and the constraints that shaped them. 
Beginning with robot perception and planning, systems commonly take as input RGB-D sequences and point clouds (often fused over time into geometric maps that ground observations into a common 3D frame~\cite{newcombe2011kinectfusion, dai2017bundlefusion}), and enable tracking, map building, and manipulation. 
Many such systems also encode visual input into compressed spatial representations, such as keypoints, dense correspondences~\cite{florencemanuelli2018dense}, and point tracks, used to parameterize 6-DoF grasps~\cite{mousavian20196} and object-centric actions.
When available, free space is typically represented with occupancy grids~\cite{Moravec1985SonarMaps, hornung2013octomap} and signed-distance representations~\cite{curless1996volumetric}, and is often additionally augmented with semantic labels such as traversability and cost for planning. 
To further enable action from perception, planners frequently operate on affordance representations~\cite{mo2021where2act} over parts and candidate contact points, encoding feasible interactions such as graspability or predicted contact quality, thereby enabling policies to execute tasks successfully.

Real-time virtual reality systems often represent scenes with triangle meshes and physically based rendering (PBR) materials, where geometry, normals, and texture define appearance, while 6-DoF camera poses keep virtual content registered to the environment. 
To enable flexible interaction, such systems include captured geometry (\eg, spatial meshes, planes, or TSDF and SDF volumes~\cite{curless1996volumetric, newcombe2011kinectfusion}) and depth-based occlusion layers, which are often paired with collision proxies and semantic labels. 
Toward photoreal view synthesis, methods increasingly use neural scene representations (\eg, NeRF-style radiance fields and 3D Gaussian splats~\cite{mildenhall2021nerf, kerbl20233d}), to support novel-view rendering and relighting. 
Additionally, for the purposes of enabling characters in virtual reality systems, such characters are commonly represented with skeletal pose and expression parameters (\eg, SMPL-style body models~\cite{loper2015smpl} or Blendshapes), and their navigation relies on navmeshes combined with acceleration structures (\eg, octrees), to handle moving agents in simulation at high frame rates.

In remote sensing applications, geographic information systems typically represent the world (the Earth) as georeferenced raster fields, for example, optical multispectral and hyperspectral imagery as well as SAR backscatter. 
These fields are aligned to explicit coordinate reference systems and sensor geometry metadata. 
Such systems encode terrain through multiple elevation products (\eg, DEM, DTM, DSM), and compress human infrastructure into vector networks (\eg, roads, power grids, pipelines), often across scales. 
And increasingly, they use LiDAR to not only output raw point clouds, but also to output rasterized height maps that support flexible 3D reasoning. 
These frameworks also commonly operate on multi-temporal stacks and spatiotemporal data with validity masks, which can generate change maps for land-use dynamics and disaster response, as well as feed digital-twin pipelines that generate orthomosaics and photogrammetric point clouds~\cite{schonberger2016structure}, for simulation and virtual modeling.

Clinical and scientific imaging operate on volumetric scans (\eg, CT and MRI), in which 3D voxel grids capture geometry and intensities~\cite{cciccek20163d}, and explicit spatial transforms (\eg, rigid, affine, and nonlinear deformation fields) align subjects and modalities into a shared atlas space. 
In addition, functional measurements (\eg, fMRI BOLD and calcium imaging) are often stored as 4D spatiotemporal tensors, but are frequently compressed into time series and activation maps after parcellation. 
When available, such systems represent structural delineations (\eg, tissue, organ, lesion) as segmentation masks and labeled maps~\cite{ronneberger2015unet}, which are often paired with surface meshes (\eg, cortical surfaces) that enable surface-based topological comparison across subjects. 
For diffusion imaging applications, tensor and orientation fields encode local diffusion directions and magnitudes, while tractography streamlines use these fields to reconstruct fiber pathways and uncover possible structural connectivity. 
These diverse representations across granularities and temporal scales are often mapped onto one another through atlas spaces and parcellations, enabling multimodal fusion and cross-subject analysis.

Intelligent assistants that operate in indoor environments typically ground language in 3D space using multi-scale maps, for example, 3D reconstructions and floorplans anchored by an estimated 6-DoF agent pose~\cite{savva2019habitat, straub2019replica, chang2017matterport3d}, which are often combined with semantic maps that encode navigable routes. 
Context for interactions is captured in structured world models (\eg, scene graphs with typed entities and relations~\cite{armeni20193dsg, wald2020learning, wu2021scenegraphfusion}), which can be augmented with additional semantics like anchors, keypoints, and regions of interest. 
For the purposes of intelligent and persistent assistance, representations that encode some form of memory are necessary, for example, episodic records and long-horizon object trajectories, which are often represented as probabilistic beliefs over object locations in a room. 
In addition, assistant systems can optionally contain geofences and occupancy masks, to help define actionable regions for planning. 

In summary, many domains have developed representations specifically designed for specialized sensors and encoding the inductive biases of the underlying data structure. 
We see that these representations model a tradeoff between (spatial) fidelity and (semantic) compressibility, reflecting how the aforementioned sensors sample space, what priors are built in, and what downstream tasks are to be supported. 



\chapter{Cross-Modal Spatial Representation}
\label{chapter:cross_modality}


% \lettrine[lines=3]{\initfamily\textcolor{darkgreen}{T}}{he} 
\FramedInitial{T}{he}
aim of cross-modal spatial representation is a latent space in which heterogeneous perceptual signals, images, 3D scenes, video, audio, and language are aligned so that the spatial structure each carries is preserved rather than averaged away.
Where the previous chapter asked how much geometry a single modality holds on its own, the concern here is the opposite and harder one: how signals that encode space in incompatible formats (a pixel grid, a world-frame point set, a waveform, a sentence) can be brought into correspondence without the alignment itself distorting what it aligns.

No single recipe achieves this.
Methods differ in what they treat as the unit of alignment, in the objective that binds one modality to another, in how tightly they anchor the space to physical geometry, and in the linguistic interface through which it is queried.
These choices are not independent, and the chapter first makes them explicit as a small set of design axes before turning to the individual cross-modal pairings in which they play out.

\section{Overview}
\label{sec:cross_modal_overview}

\FramedInitial{W}{e} define four axes that structure this cross-modal representation learning problem, which we use to organize the remainder of the chapter. 
The first axis is the \textit{alignment strategy}, which we categorize into contrastive alignment, generative alignment, and programmatic alignment. 
Contrastive alignment learns a shared embedding space by pushing together paired samples and pulling away negative samples. 
Such alignment scales to weakly paired data and enables retrieval and coarse grounding, but often needs additional spatial supervision to understand fine-grained relations grounded to specific visual regions. 
In contrast, generative alignment predicts across modalities (\eg, captioning from vision or language-conditioned decoding into images, 3D, or video), which gives rich supervision and coverage, but also tends to drift from true underlying geometry unless it is further constrained, for example, by multi-view consistency. 
Finally, programmatic alignment aligns visual entities to explicit symbols (\eg, attributes, relation graphs, or concepts in domain-specific languages), which define spatial semantics and naturally interface with program-based, verifiable reasoning, but often requires stronger supervision and trades off expressiveness for structure.

The second axis is \textit{spatial modality}, which ranges from view-aligned 2D to world-aligned 4D. Here, two-dimensional tokens (\eg, pixels, patches) are expressive and broadly applicable, but rely on some form of geometric structure to capture faithful spatial relationships. 
Three-dimensional tokens (\eg, points, meshes, implicit fields, Gaussian primitives) better encode metric scale and rigid-body invariance (\ie SE(3) consistency: rotations and translations of the scene do not change the underlying semantics), so that linguistic queries are grounded in world coordinates rather than in image projections. 
However, they are less ubiquitous and often rely on partial and noisy reconstructions of geometry. 
Extending 3D through time yields 4D traces (\eg, tracks), that bind actions to expressive motion and change, but significantly increase computational complexity and require reliable alignment data, while being sensitive to tracking errors such as identity switches.
These different spatial modalities correspond to and build upon the single-modality representations of Chapter~\ref{chapter:single_modality}.

The third axis is the \textit{language interface}, from free-form text to structured interfaces. 
Free-form text (\eg, natural language) enables open-vocabulary understanding for referring expressions, retrieval, and instruction following, but comes with corresponding ambiguity in terms of frames of reference, scale, and units. 
In contrast, structured interfaces (\eg, programs) reduce such ambiguity by explicitly describing spatial semantics, enabling integration with symbolic reasoning, but are less flexible in capturing the full range of natural language descriptions.

The fourth axis is the \textit{training regime}, which defines both the level of supervision and coverage. Paired supervision, for example, from captioned videos or referring expressions in 3D scenes, gives strong alignment, but is expensive and can be biased toward frequent compositions and co-occurrences. 
Weakly paired internet data, for example, from alt-text or transcripts, enable diversity, but, due to low signal-to-noise ratios, often yield noisier spatial grounding.

Below, we organize the chapter by cross-modal pair, using \emph{spatial modality} (which pair) and \emph{alignment strategy} (which objective) as the primary organisational axes; the \emph{language interface} and \emph{training regime} axes recur as cross-cutting attributes referenced throughout. 
We give an overview of 2D-text alignment (Section~\ref{cross_modality_2d_text}), 3D-text alignment (Section~\ref{cross_modality_3d_text}), audio-text alignment (Section~\ref{cross_modality_audio_text}), and vision-vision alignment (Section~\ref{cross_modality_vision_vision}); each follows the scope and axes defined here.


\section{2D-text}
\label{cross_modality_2d_text}
\FramedInitial{T}{he} most developed cross-modal setting pairs images with language. 
Its defining tension is that a text-aligned 2D representation lives in the image projection, not the world: appearance and layout sit on the pixel grid, yet the relations a caption asserts (left of, above, behind) are only implicit there. 
How faithfully those relations survive alignment turns on two choices: the granularity at which the image is represented, and the objective that binds it to text, and this section is organized around them. 
Models often fall into three broad families of visual representations, across different levels of granularity. First are global embeddings that compress an entire image into a single vector paired with a sentence embedding. 
Here, dual-encoder approaches such as CLIP~\cite{radford2021learning} and ALIGN~\cite{jia2021scaling} show that large-scale contrastive pretraining over noisy web data yields strong zero-shot classification and retrieval performance. 
Second, many works have proposed local tokens tied to image subregions, in which vision-language models, often transformers (\eg, ViLBERT~\cite{lu2019vilbert}, LXMERT~\cite{tan2019lxmert}, UNITER~\cite{chen2020uniter}, OSCAR~\cite{li2020oscar}) input region features extracted by bottom-up detectors (\eg, Faster R-CNN~\cite{ren2015faster}) into cross-modal attention blocks, enabling word-region alignment, and hence fine-grained reasoning over object proposals. 
A third class of methods makes relations explicit; here, models extract intermediate representations that are directly aligned with spatial predicates such as ``left-of'' or ``on.'' NSCL~\cite{mao2019neuro} and LEFT~\cite{hsu2023s} learn hybrid representations that explicitly bind named visual concepts with latent features. 
This representational spectrum, from global to region to relation-based, forms the basis for the objectives below.

Contrastive retrieval serves as a common objective for large-scale pretraining, as it learns a shared metric space, without requiring dense and expensive supervision. 
As representative examples, CLIP~\cite{radford2021learning} and ALIGN~\cite{jia2021scaling} maximize agreement between paired image and text embeddings, enabling robust zero-shot transfer; LiT~\cite{zhai2022lit} shows that contrastively tuning only the text encoder on top of a locked image encoder can further improve such transfer; and CoCa~\cite{yu2022coca} fuses contrastive learning with captioning in a single encoder-decoder architecture to combine discriminative and generative strengths. 
Another line of work adapts contrastive alignment to more fine-grained regions. 
RegionCLIP~\cite{zhong2022regionclip} lifts global contrast to region-text pairs, FILIP~\cite{yao2021filip} proposes late token-token interaction to align words with patches efficiently, X-VLM~\cite{zeng2021multi} learns multi-grained alignment such that entities and attributes can be localized and described consistently across granularities, and ALBEF~\cite{li2021align} conducts image-text contrast with cross-modal fusion and momentum-distilled pseudo-targets to stabilize learning on noisy web data. 

In addition, masked modeling with textual conditioning has been shown to provide a complementary self-supervised signal that encourages cross-modal consistency. LXMERT~\cite{tan2019lxmert} and UNITER~\cite{chen2020uniter} pretrain with masked language modeling, masked region modeling, and image-text matching to build word-region alignment.
VLMo~\cite{bao2022vlmo} generalizes this idea with mixture-of-modality experts, and BEiT-3~\cite{wang2023image} performs masked data modeling over images, text, and image-text pairs in a single multiway transformer.

With more labels, generative objectives can be applied for this representation learning problem, and also complement self-supervised ones. 
Here, methods teach models to produce language conditioned on images, intuitively, to generate tokens that act as spatially grounded explanations. 
As representative examples, SimVLM~\cite{wang2021simvlm} shows that a prefix language-modeling objective trained on weak supervision can improve performance on captioning and visual question answering, BLIP~\cite{li2022blip} adds caption bootstrapping to denoise web supervision, and GIT~\cite{wang2022git} simplifies a scalable image-encoder and text-decoder architecture.

In a more fine-grained manner, grounding objectives can directly optimize spatial correspondence between language and 2D pixels, points, or boxes. Here, text-conditioned detection and phrase grounding unify detection and language. 
MDETR~\cite{kamath2021mdetr} learns a DETR-style~\cite{carion2020end} detector modulated by text, GLIP~\cite{li2022grounded} pretrains on both detection and grounding corpora and self-generates boxes from image-text pairs to scale supervision, and Grounding DINO~\cite{liu2024grounding} augments strong detection backbones with grounded pretraining for open-set detection. 
At the pixel level, DenseCLIP~\cite{rao2022denseclip} adapts image-text matching to pixel-text matching for dense prediction, and GroupViT~\cite{xu2022groupvit} shows that segmentation structure can appear from text supervision alone.

Across these categories, we see that 2D-text representations can emerge from diverse approaches, for example, differing in whether spatial structure is implicit (\eg, in attention over patches) or explicit (\eg, in region tokens), and whether supervision is weakly paired and contrastive, generative, or explicitly grounded. 
Models that learn 2D vision-language representations now increasingly combine all contrastive, captioning, grounding, and masked objectives to solve a wide range of downstream tasks, from (more semantic) referring expression comprehension to (more visually grounded) localization.

\section{3D-text}
\label{cross_modality_3d_text}
\FramedInitial{M}{oving} from the image plane to the scene itself removes the projection gap that limits 2D-text alignment: language is now bound to geometry expressed in world coordinates, so that metric scale and the invariance of meaning under rotation and translation become properties the representation can actually enforce. 
The difficulty shifts accordingly, from recovering relations that pixels hide to obtaining and aligning 3D structure that is often partial, noisy, and expensive to annotate. 
As before, the representation families and the alignment objectives acting on them jointly organize the section. 
Two complementary families of representations have emerged in this domain. First, object-centric and scene-centric encoders learn world-aligned tokens, then align them to language via contrastive or generative objectives. 
ULIP~\cite{xue2023ulip} instantiates this paradigm by aligning point-cloud features to an image-text space, scaling supervision through synthetic triplets.
OpenShape~\cite{liu2023openshape} extends this direction further with curated 3D corpora and hard negative mining to support open-world 3D retrieval, and OpenScene~\cite{peng2023openscene} co-embeds dense features for 3D scenes with text and image pixels in CLIP feature space. 
Additional projection strategies have also been explored, for example, adapting 2D vision-language models to 3D by rendering depth or multi-view images and transferring CLIP's alignment to point clouds; PointCLIP~\cite{zhang2022pointclip} and CLIP2Point~\cite{huang2023clip2point} do so and enable adapters for cross-modal transfer.

A second family binds language to implicit scene fields. Here, LERF~\cite{kerr2023lerf} optimizes a dense 3D language field inside NeRF~\cite{mildenhall2021nerf} by rendering CLIP embeddings across views, enabling open-vocabulary queries that return relevancy volumes. 
CLIP-Fields~\cite{shafiullah2022clip} similarly distills 2D vision-language cues into 3D implicit maps to support semantic search and view localization. 
In addition, replacing NeRF with 3D Gaussian Splatting~\cite{kerbl20233d} yields efficient language fields; for example, LangSplat~\cite{qin2024langsplat} combines CLIP features with Gaussians for fast, precise open-vocabulary querying. 
These methods build on advances in explicit 3D primitives and real-time rendering to maintain high-fidelity representations while training efficiently.

In this 3D-text domain, objectives also follow the representation choices. Contrastive alignment remains common, for example, ULIP trains with InfoNCE~\cite{oord2018representation} over point clouds-images-text, while OpenShape scales both data and hard-negative mining. 
Projection-based approaches transfer contrastive objectives to depth and multi-view renders to align point clouds with text labels. 
Grounded objectives enforce explicit 3D correspondence between language and geometry. 
3D referring expression datasets and models (\eg, ScanRefer~\cite{chen2020scanrefer}, ReferIt3D~\cite{achlioptas2020referit3d}, 3D-VisTA~\cite{zhu20233d}, BUTD-DETR~\cite{jain2022bottom}, NS3D~\cite{hsu2023ns3d}, InstanceRefer~\cite{yuan2021instancerefer}) localize targets from language descriptions and learn to understand complex relations dependent on viewpoints.

Here, geometry-aware representations are critical for aligning language with the physical world. 
Ideal 3D-text representations should respect SE(3) structure (\eg, rotations and translations should not spuriously change semantics), yet cross-modal alignment should still encode view-dependent language (\eg, ``the chair to the left of the table''). 
To this end, priors from equivariant architectures, such as the SE(3)-Transformer~\cite{fuchs2020se} and Vector Neurons~\cite{deng2021vector}, can be integrated into 3D-text models. 
Implicit fields enable different options; while NeRF-based language fields are not strictly equivariant, multi-view consistency and cross-view regularization can improve their spatial coherence. 
Alternatively, Gaussian splatting uses a canonical world frame and enables efficient supervision during language-field optimization. 
Overall, these representation choices and objectives give rise to a wide range of approaches for 3D-text alignment, from tri-modal contrastive pretraining to implicit or splat-based language fields.

\section{Audio-text}
\label{cross_modality_audio_text}
\FramedInitial{S}{ound} aligns to language along a different primary dimension than vision does. What an acoustic signal encodes first is time, the sequence of events as they unfold, and only secondarily, and coarsely, does it carry spatial attributes such as direction of arrival or distance. 
Audio-text representation learning therefore aims for embeddings that are open-vocabulary and temporally grounded, with spatial awareness added where the labels permit. 
The section again moves from the representation granularity to the objectives that align it with text. 
One such direction is to embed audio-text representations at a coarse level, for example, with global encoders mapping an entire audio clip and a full sentence into a joint space, enabling zero-shot retrieval. 
Here, CLAP~\cite{elizalde2023clap} and LAION-CLAP~\cite{wu2023large} scale contrastive language-audio pretraining with curated datasets, and show strong transfer across retrieval benchmarks. 
In addition, AudioCLIP~\cite{guzhov2022audioclip} extends CLIP to a tri-modal image-text-audio setting, to leverage visual-text semantics, and Wav2CLIP~\cite{wu2022wav2clip} distills CLIP's alignment into an audio encoder when audio-text pairs are scarce.

Beyond global representations, a line of work that aims to extract representations at a more fine-grained scale learns temporally local structure, for example, by operating on spectrogram patch tokens with transformer architectures. 
The Audio Spectrogram Transformer (AST)~\cite{gong2021ast} and PaSST~\cite{koutini2021efficient} show that attention over time-frequency patches captures long-range temporal dependencies, and HTS-AT~\cite{chen2022hts} further returns frame-level maps that are useful for event localization. 
These tokenizations make it straightforward to bind subclip semantics to text (\eg, ``a door slams, then footsteps echo''). 
Self-supervised masked objectives, \eg, SSAST~\cite{gong2022ssast}, MAE-AST~\cite{baade2022mae}, and Audio-MAE~\cite{huang2022masked}, additionally pretrain such backbones by reconstructing masked patches, while BEATs~\cite{chen2022beats} adds an iterative tokenizer-encoder scheme that encourages high-level semantic abstraction.

Across advances in representations and objectives, an emergent and increasingly common paradigm is to first pretrain an audio backbone with masked modeling, then align embeddings to text via contrastive losses, and finally, optionally add a captioning head to increase supervision. 
In tasks where spatial semantics is important, incorporating spatial labels, such as direction-of-arrival, helps refine the audio-text representation space. 
By leveraging these approaches, models can learn to align predicates such as ``left-of'' or ``far'' with meaningful acoustic features. 
These methods enable models to learn embeddings that are not only open-vocabulary and temporally grounded, but also (coarsely) spatially aware.

\section{Vision-vision}
\label{cross_modality_vision_vision}
\FramedInitial{A}{} common approach to improving vision-language alignment is to first align vision with vision, intuitively, enforcing geometric consistency such that visual features are anchored to a persistent world model before combining with language.
Concretely, the goal of vision-vision representation learning is to fuse visual inputs (\eg, multiview RGB) into a set of world-aligned scene tokens, such that they are invariant to viewpoint changes and can be later queried by language. 
Here, we discuss three broad families of approaches: lifting 2D features into 3D scene tokens, projecting these 3D tokens back into 2D for supervision, and using audiovisual correspondence to add constraints that further anchor the world model.

Lifting 2D to 3D follows two complementary directions. The most classical and canonical approach first estimates camera poses and geometry with structure-from-motion (SfM) and simultaneous localization and mapping (SLAM), then adds semantics to the generated map. 
Here, commonly used frameworks include COLMAP for SfM~\cite{schonberger2016structure}, KinectFusion and ElasticFusion for dense RGB-D mapping~\cite{newcombe2011kinectfusion,  whelan2015elasticfusion}, and ORB-SLAM2 for SLAM~\cite{mur2017orb}. 
Learning-based reconstructions build upon this direction by back-projecting per-view features into a 3D volume and predicting TSDF or occupancy directly from posed images, for example, as in NeuralRecon~\cite{sun2021neuralrecon}. 
Though components may differ, all of these approaches output view-agnostic scene tokens that can later be guided by linguistic supervision and queries.

Projecting 3D back to 2D provides self-supervision. Differentiable rendering in neural fields anchors world-aligned tokens to multiview observations, and enforces cross-view consistency~\cite{mildenhall2021nerf}. 
Distilling 2D foundation features into 3D fields adds semantic knowledge without manual, expensive 3D labels. Here, works typically transfer 2D features into a co-registered 3D feature volume, which can be rendered to any view~\cite{shen2023distilled, ye2023featurenerf}. 
In addition, distillation is not limited to 2D language targets alone; for example, ImageBind~\cite{girdhar2023imagebind} proposes a multisensor feature target (image-text-audio-depth-thermal-IMU), which enables transfer of cross-sensor semantics even when there are few explicit labels. 
Although these distillations are vision-vision processes, they are explicitly designed for language alignment later. The key benefit is that now vision-language alignment is conducted on a geometry-consistent representation, rather than directly on view-specific features.

Audio-visual alignment can further constrain this vision-vision world model; self-supervised objectives have been designed to enforce consistency between audio and visual motion, which allows models to localize sound sources and disambiguate co-occurring events~\cite{arandjelovic2017look, owens2018audio, zhao2018sound}. 
Such cues provide directional or timing semantics that help anchor linguistic queries to specific world-aligned entities, especially in cases with multiple candidate entities. 
Across this variety of approaches, vision-vision representation learning methods all aim to learn a persistent world model that is geometrically and temporally consistent before adding any linguistic supervision, therefore making downstream cross-modal alignment more data-efficient and spatially faithful to the physical world.



\part{Spatial Intelligence: From Perception to Generation}
\label{part:3}
% ================================
% \section*{\centering \bf Part III: Spatial Intelligence: From Perception to Generation}
% \addcontentsline{toc}{section}{Part III: Spatial Intelligence: From Perception to Generation}
% ================================





\begin{figure}
    \centering
    % \includegraphics[width=\linewidth, height=12cm]{example-image-golden}
    % \includegraphics[width=0.99\linewidth]{figures/part3.png}
    \includegraphics[width=\linewidth]{figures/part3_opener_v2.pdf}
    
    \caption{Part III in one view. (1) Perceive: a given scene is separated into objects, depth, relational structure, and motion. (2) Reason: the understood scene supports inference and planning, here a path that carries the cup around an obstacle to the shelf. (3) Generate: an intent unfolds into images, 3D scenes, video, motion, and whole worlds, so that generation runs the loop in reverse.}
    \label{fig:p3_main}
\end{figure}

If the previous part asked how space is represented, this one asks what a system does with that representation. Under the definition of Chapter~\ref{chapter:definition}, this use resolves into three capabilities that build on one another: perceiving space, reasoning over it, and generating it, the generative side of prediction. These are less separate subfields than three phases of a single loop running from sensory input to spatial artifact.

Perception comes first because everything downstream depends on it: it converts raw observation into structured spatial evidence, answering what is present, where it lies, and how it moves. Reasoning takes that evidence as given and draws consequences from it, inferring relations, planning routes and actions, and predicting how a configuration will change under intervention. Generation closes the loop by running the process in reverse, mapping an intent (a prompt, a constraint, a partial observation) back into concrete spatial content: an image, a shape, a scene, a motion, an entire simulated world. Read in sequence, the three chapters trace the path from a scene that is given, to one that is understood, to one that is made.

Each of the three rests on the representations built in the previous part and supplies capabilities that the applications of the next will draw on; together they are the core of what makes a system spatially intelligent.






\chapter{Spatial Perception}
\label{chapter:3:sec:6}
\label{sec:perception}
% \label{sec:6}


% \lettrine[lines=3]{\initfamily\textcolor{darkgreen}{S}}{patial perception} 
\FramedInitial{S}{patial} perception is the foundational layer of spatial intelligence: it turns raw sensory input into the structured spatial evidence that everything downstream depends on, answering what is present, where it lies, and how it moves. It is what stands between an agent and a scene, and the reliability of every later inference is bounded by the reliability of what perception delivers.

The chapter follows a single progression, from the elemental to the structured. It begins with recognizing the discrete entities that populate a scene, grounds them in metric geometry, then binds them into relational and scene-level structure, extends that structure across time, and finally broadens it beyond vision by fusing complementary modalities. Each step presupposes the ones before it, so that perception builds from isolated percepts toward a dynamic, multimodal model of space that reasoning and generation can act on.



\section{Spatial Entity Recognition}
\label{sec:entity_recognition}

\FramedInitial{B}{efore} space can be measured or organized, its contents must be named. Spatial entity recognition is that entry point: it answers \emph{what is present}, and every later geometric and relational step builds on the entities it names. Recognition operates at several granularities, and the section is ordered accordingly: detecting discrete object instances, delineating dense regions of pixels or points by segmentation, pinning down sparse keypoints and landmarks, and finally assigning a holistic scene label. What unites them is that they recover identity and category, not yet geometry; where each entity sits and how it is oriented is taken up next.


\input{part_3/section_6/p3_sec6_1/p3_sec6_1_2.tex}


\input{part_3/section_6/p3_sec6_1/p3_sec6_1_3.tex}



\input{part_3/section_6/p3_sec6_1/p3_sec6_1_1.tex}



\input{part_3/section_6/p3_sec6_1/p3_sec6_1_4.tex}





% \newpage

\section{Spatial Localization and Position Estimation}
\label{sec:p3-sec6-2-spatial-localization-and-position-estimation}

\FramedInitial{R}{ecognition} says what is present; it does not say where. Spatial localization and position estimation supply the missing geometry, situating recognized entities in a metric frame and recovering how they are oriented within it. The section climbs a ladder of geometric abstraction: from the per-pixel quantities of depth and surface normals, through metric attributes such as height and distance, to full SE(3) pose, first of the moving camera and then of rigid objects in the scene. Throughout, the concern is the static geometry of a single instant; how these quantities evolve is deferred to spatio-temporal perception.


\input{part_3/section_6/p3_sec6_2/p3_sec6_2_1.tex}

\input{part_3/section_6/p3_sec6_2/p3_sec6_2_3.tex}

\input{part_3/section_6/p3_sec6_2/p3_sec6_2_2.tex}

\input{part_3/section_6/p3_sec6_2/p3_sec6_2_4.tex}

\input{part_3/section_6/p3_sec6_2/p3_sec6_2_5.tex}

% \newpage

\section{Spatial Structure Construction}
\label{sec:p3-sec6-3-spatial-structure-construction}

\FramedInitial{W}{ith} entities recognized and their positions fixed, what remains is how they hang together. Spatial structure construction turns independently localized entities into an organized configuration: which objects relate to which, how those relations compose into a scene, and how the scene as a whole is laid out. The section builds outward in scope, from pairwise spatial relations, through the scene graphs that assemble them, to room-level layout; it then lets structure take on time and scale, as the event structure of action detection and the whole-environment structure of mapping and SLAM. Structure in real perception is seldom purely static, and what the section yields is the explicit spatial configuration that downstream reasoning and planning act on.





\input{part_3/section_6/p3_sec6_3/p3_sec6_3_1.tex}

\input{part_3/section_6/p3_sec6_3/p3_sec6_3_2.tex}

\input{part_3/section_6/p3_sec6_3/p3_sec6_3_3.tex}

\input{part_3/section_6/p3_sec6_3/p3_sec6_3_4.tex}

\input{part_3/section_6/p3_sec6_3/p3_sec6_3_5.tex}


% \newpage



\section{Spatio-Temporal Perception}
\label{sec:p3-sec6-4-spatio-temporal-perception}

\FramedInitial{M}{ost} of what precedes has described a single instant. Spatio-temporal perception adds the missing axis, asking not only where entities are but how they move, when events happen, and which identities persist. 
Its three concerns rise from the pixel through the object to the event: scene flow recovers dense, per-point motion between frames; object tracking maintains the identity of an entity as it moves through space and time; and temporal localization fixes when an action or event occurs along a sequence.
Together they give the largely static perception of the preceding sections a temporal thread, the continuity on which motion understanding and long-horizon perception depend.





\input{part_3/section_6/p3_sec6_4/p3_sec6_4_3.tex}


\input{part_3/section_6/p3_sec6_4/p3_sec6_4_1.tex}


\input{part_3/section_6/p3_sec6_4/p3_sec6_4_2.tex}




% \newpage

\section{Multimodal Spatial Perception}
\label{sec:p3-sec6-5-multimodal-spatial-perception}
% \wsq{将视觉与语言、音频、触觉等模态结合，增强空间感知范围与粒度。}

\FramedInitial{T}{he} preceding sections perceive space largely through vision. Yet vision alone is often incomplete: it is blind past occluders, silent about what lies outside the frame, and mute on which entity a person actually means. Multimodal spatial perception closes these gaps by grounding spatial understanding in signals beyond the image. The section takes two complementary couplings: audio-visual spatial localization, which uses sound to place events that vision may miss, and vision-language grounding, which ties linguistic reference to the specific entity or region it denotes. Neither adds a new spatial primitive; each supplies a channel of evidence that makes perception more complete where a single modality falls short.

\input{part_3/section_6/p3_sec6_5/p3_sec6_5_1.tex}

\input{part_3/section_6/p3_sec6_5/p3_sec6_5_2.tex}


\section{Summary and Outlook}
\label{sec:p3_s6_6}

This chapter followed spatial perception from the elemental to the structured: recognizing the entities in a scene, fixing them in metric geometry, binding them into relations, layouts, and maps, carrying that structure through time, and supplementing vision with sound and language. Three observations recur across these tasks.

\textbf{Priors fill gaps that observation leaves open.} Monocular depth, surface normals, and camera motion are underdetermined by a single image: scale, reflectance, and the geometry behind an occluder are not fixed by the pixels (§\ref{sec:depth_estimation}, §\ref{sec:visual_odometry}). Learned models increasingly resolve this ambiguity with strong priors, which makes them confident on familiar scenes and fragile on transparent or reflective materials, moving objects, adverse conditions, and domains their training data under-represent. Distinguishing what a percept observes from what a prior supplies is a requirement the later chapters inherit, and Chapter~\ref{chapter:p5_discussions} returns to it as a question of observability.

\textbf{Foundation models move the bottleneck rather than remove it.} Open-vocabulary detectors and segmenters, zero-shot scene classifiers, and 3D geometry foundation models have widened what perception covers without task-specific training, yet the hardest cases persist at the long tail: small and rare objects, unseen instances in clutter, and metric judgments of distance, size, and angle, on which multimodal models remain weak even when their semantic predictions are strong (§\ref{sec:spatial_relation_extraction}). The accuracy--efficiency tension noted for detection recurs in every task of the chapter.

\textbf{Evaluation rewards the easier part of each problem.} Many benchmarks are synthetic, saturated, or confined to a few curated scenes; tracking metrics do not measure robustness to occlusion length (§\ref{sec:object_tracking}), temporal-localization metrics treat fuzzy boundaries as exact, and grounding benchmarks pose a single query against a single image rather than an interaction (§\ref{sec:vl_grounding}).

What the chapter delivers is structured spatial evidence: entities with positions, relations, trajectories, and maps. Chapter~\ref{sec:reasoning} takes that evidence as input and asks what can be inferred and planned from it, and Chapter~\ref{sec:generation} runs the mapping in the opposite direction.
\newpage
\chapter{Spatial Reasoning}
\label{sec:reasoning}

\FramedInitial{S}{patial} reasoning is the step that turns perceived structure into decisions: given detections, depths, poses, and a task, it works out the relations, the causes, and the actions that follow, the why and how of spatial interaction rather than merely the what and where. It is where perception becomes goal-directed, supporting inference, planning, and problem-solving in physical and virtual environments.

The chapter is organized along the spectrum by which spatial reasoning has actually developed, from explicit symbolic models to learned and foundation-model-driven ones. At one end, reasoning runs over hand-built representations of geometry, physics, and logic, whose structure is inspectable and whose conclusions can be certified. In the middle, it emerges from data, as neural models acquire spatial competence directly rather than being told the rules. More recently, large pre-trained models supply reasoning as a general capability to be elicited and steered. Hybrid approaches recombine these, pairing the semantic reach of foundation models with the precision and verifiability of symbolic and learned components; a closing survey of benchmarks then records how the field measures spatial reasoning as it shifts from passive observation toward active, situated interaction.

\section{Rule-based Symbolic Spatial Reasoning}
\label{sec:p3_s7_1}

\FramedInitial{R}{ule-based} symbolic spatial reasoning covers the methods that operate over \emph{explicit, model-based representations} of space (geometric primitives, physical laws, graph connectivity, formal constraints, action rules, or causal mechanisms) rather than over learned latent representations alone. What sets this family apart is that reasoning happens inside a model that can be inspected: its state and transitions are legible, so claims about feasibility, optimality, or safety can be backed by certificates rather than empirical averages, and when a plan fails it fails legibly, as an infeasible inverse kinematics, a violated constraint, a missing precondition, or an inconsistent causal attribution, not as an opaque correlation.

The section surveys six kinds of explicit structure, ordered from the most concrete to the most abstract: geometric and algebraic models of position, orientation, motion, and feasibility; physical models of contact, friction, and stability; topological and graph abstractions of space that turn planning into search; formal constraint systems whose solvers return certified plans or proofs of infeasibility; rule-based and programmatic logic that sequences and supervises actions; and causal models that explain, predict counterfactually, and prescribe interventions.
% Table~\ref{tab:p3_s7_1_paradigm_comparison} at the end of the section compares the six.

These paradigms are not isolated alternatives. Modern robotic systems routinely compose them: geometric solvers serve as feasibility oracles inside constraint-based search; behavior trees wrap symbolic plans for fault tolerance; causal action languages compile to SAT; topological maps anchor both classical and learning-based navigation policies. The unifying thread is that each paradigm carries explicit structure (a state, an operator, a constraint, a rule, or a mechanism) that can be inspected, composed, and reasoned about.

\input{part_3/section_7/p3_sec7_1/p3_sec7_1_1.tex}


\input{part_3/section_7/p3_sec7_1/p3_sec7_1_2.tex}

\input{part_3/section_7/p3_sec7_1/p3_sec7_1_3.tex}


\input{part_3/section_7/p3_sec7_1/p3_sec7_1_4.tex}

\input{part_3/section_7/p3_sec7_1/p3_sec7_1_5.tex}

\input{part_3/section_7/p3_sec7_1/p3_sec7_1_6.tex}


% \subsection{Summary}
% \label{subsec:p3_s7_1_summary}

% The six paradigms in §\ref{sec:p3_s7_1} build spatial intelligence on explicit, inspectable structure (geometric, mechanical, topological, formal, programmatic, or causal) rather than on learned correlations alone, and that choice buys four things. Plans, schedules, and explanations are artifacts that can be debugged and verified rather than distributions to be sampled. Each paradigm increasingly produces machine-checkable certificates (force-closure conditions, optimality witnesses, $\delta$-complete SMT proofs, unsat cores) on which safety-critical deployment can rely. Explicit interfaces (operators, constraints, rules, causal variables) let the paradigms compose with one another and with learned components without losing their guarantees. And where the model captures the relevant structure, classical methods generalize from a compact specification rather than from a large corpus.

% Table~\ref{tab:p3_s7_1_paradigm_comparison} sets the six side by side; they are complementary rather than redundant, each targeting a structure that the others presuppose (geometry beneath physics), wrap (behavior trees around symbolic plans), or extend (causal mechanisms over constraint solutions).

% \begin{table}[t]
%     \centering
%     \footnotesize
%     \caption{Comparison of the six foundational paradigms of rule-based symbolic spatial reasoning covered in §\ref{sec:p3_s7_1}.}
%     \label{tab:p3_s7_1_paradigm_comparison}
%     \begin{tabularx}{\linewidth}{@{}l X X X X@{}}
%         \toprule
%         \textbf{Paradigm} & \textbf{Core representation} & \textbf{Output / artifact} & \textbf{Notable strength} & \textbf{Main limitation} \\
%         \midrule
%         §\ref{subsec:p3_s7_1_1} Geometric \& algebraic         & $SE(3)$ poses, projective geometry, Jacobians, contact wrenches                  & Frame-aligned coordinates, workspace / reachability sets, feasibility certificates      & Closed-form precision; coordinate-free formulations                          & Brittle near singularities; assumes exact geometry \\
%         §\ref{subsec:p3_s7_1_2} Physical                       & Rigid-body dynamics, contact / friction laws, quasi-static equilibrium           & Force-closure verdict, motion cone, support-polygon test                                & Grounded in mechanics; analytically interpretable                            & Sim-to-real gap; latent parameters (mass, friction) often unobserved \\
%         §\ref{subsec:p3_s7_1_3} Topological / graph-based      & Roadmap / decomposition / sampling graph $G=(V,E,\mathbf{w})$                    & Collision-free path $\pi$, topological map, multi-agent trajectories                    & Combinatorial tractability; mature solvers (A*, CBS, CH)                     & Density vs.\ optimality trade-off; appearance robustness across seasons \\
%         §\ref{subsec:p3_s7_1_4} Constraint-based               & Decision variables + collision / IK / dynamics / logic constraints               & Certified-feasible plan, optimality gap, infeasibility witness                          & Strong correctness guarantees; certificate-producing                         & Mixed-integer scaling; symbol-grounding from perception \\
%         §\ref{subsec:p3_s7_1_5} Rule-based / programmatic      & STRIPS / PDDL operators, FSMs, Behavior Trees, production rules                  & Action plan, reactive policy, monitored execution trace                                 & Auditable; reactive recovery; modular composition                            & Modeling burden; brittle under perception noise without grounding \\
%         §\ref{subsec:p3_s7_1_6} Causal                         & Structural causal models, temporal networks (STN / STNU)                         & Causal model, $do(\cdot)$-effect estimate, counterfactual answer, intervention strategy & Counterfactual / responsibility queries; identifiability theorems            & Grounding causal variables; latent confounding; LiNGAM needs non-Gaussian noise \\
%         \bottomrule
%     \end{tabularx}
% \end{table}

% The same design philosophy exposes the limitations. Explicit models must be authored and maintained, they presume that the relevant structure is observable, and they degrade gracelessly when reality violates their idealizations (perception noise, latent parameters, unmodeled contact, drifting causal mechanisms). These are the conditions under which data-driven methods excel: learned pipelines absorb unobservable variables and perception noise into end-to-end objectives and generalize across regimes that would be intractable to specify by hand.

% The rest of Chapter~\ref{sec:reasoning} takes up spatial reasoning from that complementary angle: §\ref{sec:p3_s7_2} on reasoning learned from data, §\ref{sec:p3_s7_3} on reasoning elicited from foundation models, §\ref{sec:p3_s7_4} on hybrids of the two with the classical paradigms, and §\ref{sec:p3_s7_5} on how the field measures all of them. Throughout, the classical paradigms of §\ref{sec:p3_s7_1} reappear as feasibility oracles, safety shells, structural priors, and explainability layers, and the most effective systems are those that compose the two regimes: explicit structure where it is authoritative, learned structure where it is unavoidable.


% \newpage

\section{Learning-based Spatial Reasoning}
\label{sec:p3_s7_2}

\FramedInitial{L}{\textbf{earning-based spatial reasoning}} refers to the class of methods that acquire and perform spatial reasoning from data, rather than from the hand-authored geometric, mechanical, or symbolic models surveyed in §\ref{sec:p3_s7_1}.
The six subsections that follow are ordered by the \emph{locus of the structural prior}: it begins inside the learned representation (latent feature interaction, explicit programs, hard-coded symmetries, optimization layers in the forward pass), moves to the agent's action-conditioned behavior (reinforcement learning with spatial inductive biases), and ends with joint decision-making across coordinated agents.


\input{part_3/section_7/p3_sec7_2/p3_sec7_2_1.tex}

\input{part_3/section_7/p3_sec7_2/p3_sec7_2_2.tex}

\input{part_3/section_7/p3_sec7_2/p3_sec7_2_3.tex}

\input{part_3/section_7/p3_sec7_2/p3_sec7_2_4.tex}

\input{part_3/section_7/p3_sec7_2/p3_sec7_2_5.tex}

\input{part_3/section_7/p3_sec7_2/p3_sec7_2_6.tex}




% \subsection{Summary}
% \label{subsec:p3_s7_2_summary}

% Read from top to bottom, the six paradigms move the structural prior from the representation to behavior: implicit reasoning leaves it latent, explicit and geometry-aware reasoning lifts it into a program or a symmetry group, optimization layers place it in the forward pass, and the two reinforcement-learning subsections let it shape the policy of one agent and then of a team.
% Table~\ref{tab:p3_s7_2_paradigm_comparison} summarizes the six along the axes used for the symbolic paradigms in §\ref{sec:p3_s7_1}.

% \begin{table}[t]
%     \centering
%     \fontsize{8}{10}\selectfont
%     \setlength{\tabcolsep}{2pt}
%     \caption{Comparison of the six paradigms of learning-based spatial reasoning covered in §\ref{sec:p3_s7_2}, on the same axes as the comparison table of §\ref{sec:p3_s7_1}; the \emph{locus} column gives where each paradigm injects its structural prior.}
%     \label{tab:p3_s7_2_paradigm_comparison}
%     \begin{tabularx}{\textwidth}{@{}p{2.9cm} p{3.3cm} p{3.3cm} X X@{}}
%         \toprule
%         \textbf{Paradigm} & \textbf{Structural prior (\emph{locus})} & \textbf{Canonical inputs / outputs} & \textbf{Strengths} & \textbf{Limitations} \\
%         \midrule
%         §\ref{subsec:p3_s7_2_1} Implicit neural reasoning      & Latent prior: attention, relation networks, FiLM modulate the learned feature space                & Pixel / feature inputs $\to$ answer, caption, or relation label                                  & End-to-end fluency on perception-heavy benchmarks; minimal modeling burden                & Limited inspectability; compositional generalization gaps; shortcut learning under sparse supervision \\
%         §\ref{subsec:p3_s7_2_2} Explicit structured reasoning  & Program- or graph-level prior: parsed layouts, scene graphs, neural-module / LLM-generated programs & Image $+$ compositional question $\to$ executable program $\to$ stepwise inspectable answer       & Compositional generalization; auditability; modular reuse of perception tools              & Parser / program-generator quality is a hard ceiling; brittle on long-tail vocabulary; supervision-hungry \\
%         §\ref{subsec:p3_s7_2_3} Geometry-aware reasoning       & Symmetry-group prior: invariance / equivariance to a chosen group of input transformations         & Image / point cloud / molecule $\to$ classification, pose, or property prediction                & Sample efficiency; out-of-distribution rotation / translation robustness                    & Group choice is a silent error mode; classical equivariant filters scale with $\lvert G \rvert$; brittle when inputs violate the assumed symmetry \\
%         §\ref{subsec:p3_s7_2_4} Differentiable optim.\ layers  & Optimization prior: a constrained convex / SAT problem is solved inside the forward pass            & Image / pose / observation $\to$ optimization-layer call $\to$ optimal solution $+$ KKT gradients & Certified feasibility / structure on the prediction; end-to-end training preserves geometry & Solver wall-clock dominates the forward pass; biased / high-variance gradients near constraint boundaries \\
%         §\ref{subsec:p3_s7_2_5} RL with spatial priors          & Action-conditioned prior: differentiable planning, learned maps, relational policies               & Observation $+$ action history $\to$ action under MDP / POMDP; reward signal                      & Long-horizon credit assignment; generalization across spatial layouts; embodied tasks       & Sample inefficiency vs.\ supervised paradigms; reward-shaping fragility; sim-to-real gap \\
%         §\ref{subsec:p3_s7_2_6} Multi-agent spatial reasoning   & Joint action-conditioned prior: Dec-POMDP, CTDE, graph-MARL, sequence-model coordination           & Per-agent egocentric obs.\ $+$ messages $+$ interaction graph $\to$ joint decentralized actions   & Team-level coordination; topology-aware robustness; resilience under attrition              & Credit assignment under partial observability $+$ non-stationarity; $\mathcal{O}(n^2)$ scaling; partner generalization immature \\
%         \bottomrule
%     \end{tabularx}
% \end{table}

% Every paradigm in the table trades the certificates of §\ref{sec:p3_s7_1} (feasibility proofs, force-closure verdicts, $do(\cdot)$-effect estimates) for end-to-end trainability, and each keeps a slice of the classical structure as its prior; which slice it keeps determines both its strengths and its blind spots.
% As the locus of that prior migrates from representation to behavior, the evaluation regime migrates with it, from offline benchmark accuracy at the top of the table to closed-loop team return at the bottom.
% What no single row supplies is \emph{cross-locus composition}: a geometry-aware backbone (§\ref{subsec:p3_s7_2_3}), an inspectable program-level reasoner (§\ref{subsec:p3_s7_2_2}), a differentiable feasibility layer (§\ref{subsec:p3_s7_2_4}), and an action-conditioned policy (§\ref{subsec:p3_s7_2_5}--§\ref{subsec:p3_s7_2_6}) stacked into one pipeline without losing the guarantee each contributes.
% That composition problem is the bridge to §\ref{sec:p3_s7_3}, where the prior is repositioned inside a large pre-trained model, and to §\ref{sec:p3_s7_4}, where it is distributed across a hybrid of classical and learned components.

% \newpage


\section{Foundation Model-driven Spatial Reasoning}
\label{sec:p3_s7_3}

\FramedInitial{F}{oundation} model-driven spatial reasoning turns to large pre-trained models (LLMs, vision-language models, vision-language-action models, and large vision models) as the engine of spatial inference. The motivation is empirical: such models display a striking facility with spatial concepts and relations even without task-specific fine-tuning, and a large body of work now recruits them as the core of systems for spatial question answering, planning, and explanation.

What separates these approaches is the locus and granularity at which the foundation model intervenes. At one end the model is the reasoning core itself, answering spatial queries directly (LLM/VLM-centric reasoning) or additionally emitting low-level actions (VLA-based reasoning). Its reasoning can be drawn out step by step, as intermediate deductions in language (chain-of-thought) or as an executable program synthesized from the task (code and DSL generation). It can be coupled to external machinery such as planners, renderers, simulators, and perception modules (tool-augmented reasoning), or grounded in external spatial memory (retrieval- and memory-augmented reasoning). And it can be handed control of a closed interaction loop with the environment (agentic reasoning) or paired with a learned neural simulator of dynamics that serves as the spatial substrate (world-model-based reasoning).


\input{part_3/section_7/p3_sec7_3/p3_sec7_3_1.tex}

\input{part_3/section_7/p3_sec7_3/p3_sec7_3_2.tex}

\input{part_3/section_7/p3_sec7_3/p3_sec7_3_3.tex}

\input{part_3/section_7/p3_sec7_3/p3_sec7_3_4.tex}

\input{part_3/section_7/p3_sec7_3/p3_sec7_3_5.tex}

\input{part_3/section_7/p3_sec7_3/p3_sec7_3_6.tex}

\input{part_3/section_7/p3_sec7_3/p3_sec7_3_7.tex}

\input{part_3/section_7/p3_sec7_3/p3_sec7_3_8.tex}








% \subsection{Summary}
% \label{subsec:p3_s7_3_summary}

% The eight subsections vary \emph{where} and \emph{how} the foundation model intervenes in the reasoning loop: as the reasoning core itself (§\ref{subsec:p3_s7_3_1}--§\ref{subsec:p3_s7_3_2}), as a generator of intermediate rationales or programs (§\ref{subsec:p3_s7_3_3}--§\ref{subsec:p3_s7_3_4}), as an orchestrator over tools, memory, and a closed interaction loop (§\ref{subsec:p3_s7_3_5}--§\ref{subsec:p3_s7_3_7}), or beside a learned simulator that supplies the spatial substrate (§\ref{subsec:p3_s7_3_8}).
% Table~\ref{tab:p3_s7_3_paradigm_comparison} summarizes the eight along the axes used for the symbolic paradigms in §\ref{sec:p3_s7_1} and the learning-based paradigms in §\ref{sec:p3_s7_2}.

% \begin{table}[t]
%     \centering
%     \fontsize{8}{10}\selectfont
%     \setlength{\tabcolsep}{2pt}
%     \caption{Comparison of the eight paradigms of foundation-model-driven spatial reasoning covered in §\ref{sec:p3_s7_3}, on the same axes as the comparison tables of §\ref{sec:p3_s7_1} and §\ref{sec:p3_s7_2}.}
%     \label{tab:p3_s7_3_paradigm_comparison}
%     \begin{tabularx}{\textwidth}{@{}p{2.9cm} p{3.3cm} p{3.3cm} X X@{}}
%         \toprule
%         \textbf{Paradigm} & \textbf{Locus of foundation-model intervention} & \textbf{Canonical inputs / outputs} & \textbf{Strengths} & \textbf{Limitations} \\
%         \midrule
%         §\ref{subsec:p3_s7_3_1} LLM/VLM-centric reasoning   & Pre-trained model \emph{is} the reasoning core; intervention via prompt / train / data / model-centric enhancement & Text-only or image$+$query $\to$ qualitative relation, distance, orientation, navigation plan & Zero-shot deployment; rapid prompt iteration; benefits from any backbone improvement                  & Ungrounded under text-only; metric-precision gap under VLM; reference-ambiguity in embodied use \\
%         §\ref{subsec:p3_s7_3_2} VLA-based reasoning         & Pre-trained model also emits low-level actions (action tokens / deltas / waypoints)                                & Image $+$ proprioception $+$ instruction $\to$ executable action sequence                            & End-to-end semantic-to-motor pipeline; cross-embodiment positive transfer via large data            & Semantic--physical misalignment; pseudo-interaction; inference-latency under autoregressive decoding \\
%         §\ref{subsec:p3_s7_3_3} Chain-of-thought (CoT \& variants) & Pre-trained model, with intermediate textual rationale as the prior; five topologies (CoT/ToT/GoT/PoT/CoA) unified as guided search & Query $\to$ rationale trace $\to$ grounded answer (qualitative relation, coordinate, plan, action) & Multi-step decomposition; verifier-friendly intermediate states; search-quality tunable               & Spatial-coordinate hallucination under simple textual CoT; rationale not grounded in truth signal; ToT/GoT cost scales super-linearly in depth \\
%         §\ref{subsec:p3_s7_3_4} Code / DSL generation         & Pre-trained model generates an executable program; environment supplies the truth signal                          & Query$+$API $\to$ program $P$ $\to$ executed output (answer with trace, or environment change)   & Interpretable rationale trace; execution-grounded; composes existing perception or control tools     & Reliance on a fixed human-authored DSL; brittle code execution under upstream-module error; sample-inefficient debug loops \\
%         §\ref{subsec:p3_s7_3_5} Tool-augmented reasoning      & Pre-trained model orchestrates a tool catalogue $\mathcal{T}$ via policy $\pi_\phi$, along execution / representation / validation / memory axes & Query$+$observation$+$tool set $\to$ history-conditioned answer $+$ tool-call trace             & Modular extension of perceptual / geometric / verification capacity; orchestration policy as the learning target & Combinatorial exploration in tool-action space; metric gap in the LLM--tool interface; conceptual overlap with §\ref{subsec:p3_s7_3_4} / §\ref{subsec:p3_s7_3_6} \\
%         §\ref{subsec:p3_s7_3_6} Retrieval- \& memory-augmented & External memory bank itself is the central learning target; Planner--Memory--Retriever triple $(\mathcal{P}, \mathcal{M}, \mathcal{R})$ & Query$+$scene observation$+$memory $\to$ action plan, retrieved subgraph, retrieved past observation & Kilometer-scale environment coverage; personalization; explainable trace of retrieved memories       & Stateless inference under unbounded context; ICL vs.\ cross-attention fusion unresolved; modality mismatch in retrieval (text vs.\ image vs.\ 3D) \\
%         §\ref{subsec:p3_s7_3_7} Agentic reasoning             & Orchestration policy $\pi$ over the perceive--reason--act--reflect cycle becomes the central learning target; tools (§\ref{subsec:p3_s7_3_5}) and memory (§\ref{subsec:p3_s7_3_6}) composed into closed-loop control & Instruction$+$observation $+$ tools $+$ memory $\to$ sensorimotor action sequence, accumulated skill library, updated world model & Closed-loop error correction; long-horizon procedural-knowledge accumulation; explicit Spatial Theory-of-Mind in multi-agent settings & Sim-to-real fidelity gap; long-horizon spatiotemporal coherence; convergence with §\ref{subsec:p3_s7_3_3}--§\ref{subsec:p3_s7_3_6} blurs the agentic boundary \\
%         §\ref{subsec:p3_s7_3_8} World-model--based reasoning  & Learned generative simulator $\mathbb{P}(s_0,\dots,s_T \mid a_0,\dots,a_{T-1}, c)$ replaces / augments the foundation model as the spatial substrate & Context $c$ ($+$ actions if conditioned) $\to$ inferred spatial properties (understanding) or policies $\pi$ (planning) & Predictive-physics reasoning; inference-time mental simulation; 4D / object-centric structural priors & World-model upper bound (Genie-style long-term memory loss); modality coverage gated by pre-training data; likelihood-vs.-physics gap (likelihood necessary but not sufficient) \\
%         \bottomrule
%     \end{tabularx}
% \end{table}

% Down the table the locus of intervention moves from \textbf{model-as-core} (§\ref{subsec:p3_s7_3_1}--§\ref{subsec:p3_s7_3_2}) through \textbf{model-as-rationale-generator} (§\ref{subsec:p3_s7_3_3}--§\ref{subsec:p3_s7_3_4}) and \textbf{model-as-orchestrator} over external components (§\ref{subsec:p3_s7_3_5}--§\ref{subsec:p3_s7_3_7}) to a \textbf{learned simulator as the reasoning substrate} (§\ref{subsec:p3_s7_3_8}), and the principal learning target migrates with it, from model weights to prompt structure to the orchestration policy $\pi$ to the world model $\mathbb{P}$ itself.
% The paradigms are cumulative rather than parallel: from §\ref{subsec:p3_s7_3_3} onward each adds one dimension (programs, tools, memory, reflection, world models) on top of the previous ones, so the ``agentic'' boundary of §\ref{subsec:p3_s7_3_7} lies where the orchestration policy itself becomes a learning target rather than where any new capability appears, which is an open research question rather than a settled definition.
% The open problem that cuts across all eight is the \emph{precision--grounding asymmetry}: the foundation models that supply semantic competence (§\ref{subsec:p3_s7_3_1}--§\ref{subsec:p3_s7_3_4}) lack the metric precision physical manipulation requires, while the external components that supply precision (§\ref{subsec:p3_s7_3_5}--§\ref{subsec:p3_s7_3_8}) inherit the LLM--tool interface as a new bottleneck.
% That asymmetry is the bridge to §\ref{sec:p3_s7_4}, where foundation-model-driven and symbolic or learning-based components are composed to close the gap from both sides at once.

% \newpage
\section{Hybrid Spatial Reasoning}
\label{sec:p3_s7_4}

\FramedInitial{C}{onventional} approaches to spatial reasoning exhibit fundamental limitations.
Systems that rely solely on linguistic or symbolic reasoning typically lack perceptual grounding in real-world visual and physical structures (\eg, geometry, physics, spatial continuity), making them inadequate for handling cross-view transformations, temporal evolution, or physically coherent scenarios.
Conversely, purely perceptual or vision-based systems, such as those combining visual recognition with rule-based geometric reasoning, often fall short in performing high-level abstract reasoning, multi-step planning, long-term memory integration, or interactive decision-making.
These limitations have motivated a growing body of work advocating for a hybrid paradigm of spatial reasoning: integrating heterogeneous modules with distinct capabilities (perception, world modeling, memory, and action) so as to jointly leverage low-level spatial and physical cues together with high-level cognitive reasoning capabilities.
A perceptual or physics-informed module supplies geometry, relations, affordances, and dynamics; a symbolic or neural module plans, infers, and abstracts over them; memory and history extend the pipeline to multi-step, multi-view, and temporally extended tasks; and because no single modality, task, or dataset carries the whole system, the composition transfers across vision-only, language-only, embodied, and multimodal settings.

% 在空间推理 (spatial reasoning) 的研究中，传统单一范式存在显著局限性。以仅依赖语言/符号推理 (symbolic reasoning) 的系统为例，它们往往缺乏对真实世界视觉／物理结构 (geometry / physics / continuity) 的感知与空间记忆，难以处理跨视角 (cross-view)、跨时 (temporal) 和物理连贯 (physical consistency) 的复杂情境；另一方面，仅依赖感知或视觉模块 (例如视觉识别 + 规则几何推理) 的系统，则可能缺乏高阶抽象推理 (abstract reasoning)、规划 (planning)、长程记忆 (long-term memory) 以及交互式决策 (interactive decision) 的能力。
% 因此，很多工作主张一种 “混合 (hybrid)”空间推理 (spatial reasoning) 的研究范式 — 即通过整合异构 (heterogeneous) 模块 (perception, world-model, reasoning, memory, action) —— 以同时兼顾空间结构感知、几何 / 物理一致性、跨视角／跨模态理解、长期规划与推理。
% 这种混合范式的优点包括但不限于：
% 感知 + 抽象推理：底层视觉 / 物理世界模型 (world-model) 提供空间结构与物体关系 (spatial relationships / affordances / dynamics)，而高层符号 /语言 /逻辑推理 (symbolic or neural reasoning) 提供规划、决策、归纳／演绎等能力；
% 长期记忆与交互能力：通过将 memory / history / context 纳入 world-model + reasoning pipeline，使系统能够处理长期任务 (multi-step, multi-view, temporally-extended)；
% 泛化与灵活性 (generalization & flexibility)：混合模块使系统不依赖单一任务或单一数据集 (\eg pure 3D supervision, or pure NL reasoning)，从而更具拓展性 (scalability) 与适应性 (adaptability)。


\subsection{Task Definition}
\label{subsec:p3-sec7-4-task-definition}
A hybrid spatial reasoning system is a pair $(\{M_i\}_{i=1}^{k}, \pi)$: a set of $k$ heterogeneous modules $M_i \in \{M_\text{percept}, M_\text{world}, M_\text{symb}, M_\text{mem}, M_\text{plan}, M_\text{act}\}$, each contributing a distinct inductive bias (perception, world model, symbolic constraints, memory, planning, or action), and an orchestration function $\pi : \mathcal{S} \to \{1, \dots, k\}$ that selects which module acts on the current system state $s_t \in \mathcal{S}$ (perception, internal map, history). Starting from the input $s_0 = x$, the system evolves as
\begin{equation}
  s_{t+1} = M_{\pi(s_t)}(s_t), \qquad y = \mathrm{Read}(s_T),
  \label{eq:p3_s7_4_hybrid}
\end{equation}
and the three composition styles found in practice are special cases of Eq.~\eqref{eq:p3_s7_4_hybrid}: a \emph{pipeline} fixes $\pi$ to a constant module order, an \emph{iterative} system lets $\pi$ revisit modules until a verifier accepts, and an \emph{end-to-end} system makes every $M_i$ differentiable so that the whole trajectory is trained jointly. Hybrid spatial reasoning is the paradigm in which $\pi$ and the module set $\{M_i\}$ are designed together to minimize the precision-grounding asymmetry: the foundation models of §\ref{sec:p3_s7_3} supply semantic competence, and the symbolic and learning-based components of §\ref{sec:p3_s7_1}--§\ref{sec:p3_s7_2} supply metric precision and verifiable structure. Existing systems therefore differ along three axes: \emph{which} modules they compose, \emph{how} they compose them (pipeline, iterative, or end-to-end), and how the orchestration function is realized (hand-coded rules or a learned policy).

\subsection{Methodology}
\label{subsec:p3-sec7-4-methodology}

Five hybridization patterns recur in the literature, distinguished by which module the foundation model is paired with and by what supplies the truth signal.

\begin{figure}
    \centering
    \includegraphics[width=\linewidth]{figures/part3/section7/foundation_rl_hybrid.pdf}
    \caption{The foundation-model + reinforcement-learning hybrid pattern: an LLM/VLM hypothesis generator coupled with an RL controller and with a physics engine or symbolic solver for verification. The RL controller's policy (written $\pi$ in the figure) is $\pi_{\mathrm{RL}}$ in the text; the bottom bar is the orchestration function $\pi$ of Eq.~\eqref{eq:p3_s7_4_hybrid}.}
    \label{fig:p3_s7_4_foundation_rl_hybrid}
\end{figure}

\paragraph{Foundation Models with Reinforcement Learning.}
A prominent hybrid couples a foundation model's semantic competence with reinforcement learning's long-horizon decision-making, often closed by a verifier that supplies executable structure (Figure~\ref{fig:p3_s7_4_foundation_rl_hybrid}).
The foundation model proposes hypotheses, sub-goals, or candidate action sequences; an RL policy $\pi_{\mathrm{RL}}$ commits to and executes them over extended horizons; and the resulting reward, optionally augmented by a physics engine or symbolic checker, updates both the policy and the orchestration.
Embodied-R~\cite{zhao2025embodied} instantiates this division of labor: a large vision--language model handles perception, while a small language model trained by RL, with a reward that scores the logical consistency between its reasoning trace and its answer, performs the slow, deliberative spatial reasoning, evaluated on UrbanVideo-Bench and VSI-Bench.
The group-relative fine-tuning recipes of §\ref{subsec:p3_s7_2_5} (Reason-RFT, OctoNav, SpaceR~\cite{tan2025reasonrft,gao2026octonav,ouyang2025spacer}) are the same hybrid with the verifier folded into a verifiable reward and are not repeated here.
The foundation model contributes broad priors and RL supplies credit assignment and exploration; the cost is the sample inefficiency and reward-specification difficulty inherited from RL (§\ref{subsec:p3_s7_2_5}).

\paragraph{Foundation Models with Explicit Physical Simulation.}
A second pattern pairs the foundation model with a physics engine, so that the truth signal is a simulated trajectory rather than a learned score; the engine side is the physical reasoning of §\ref{subsec:p3_s7_1_2}.
The LLM proposes hypotheses, structural descriptions, or programs, the simulator executes them under conservation laws, kinematic constraints, contact dynamics, and reachability limits, and the discrepancy between simulation and observation is fed back for another round.
LLMPhy~\cite{cherian2024llmphy} is the clearest instance: the LLM writes Python programs that embed estimates of latent physical parameters (friction, stiffness, mass) and of the scene layout, a physics engine executes them, and the reconstruction error against the observed video is returned as feedback, so that building the digital twin becomes a black-box optimization driven by the LLM.
DiffGen~\cite{jin2024diffgen} makes the loop differentiable, generating robot demonstrations by descending gradients through a differentiable simulator, a differentiable renderer, and a vision--language model until the rendered outcome matches the language instruction in embedding space.
PhysVLM~\cite{zhou2025physvlm} moves the physics into the input instead: a robot's reachable workspace is abstracted into a robot-agnostic space-physical reachability map that the VLM consumes alongside the image, trained on Phys100K and evaluated on the EQA-phys benchmark.
At the far end, ChronoLLM~\cite{wang2025chronollm} fine-tunes LLMs to write PyChrono multibody-simulation scripts, from a single pendulum to a full vehicle on deformable terrain; here the model authors the simulation rather than reasoning inside it, which places it at the boundary with the code-generation paradigm of §\ref{subsec:p3_s7_3_4}.
What the simulator contributes is an inductive bias that stays valid for unseen objects and configurations, so the hybrid extrapolates where a purely statistical model would not.

\paragraph{Foundation Models with Symbolic Reasoning.}
A third pattern inserts an explicit symbolic layer (logical constraints, relational predicates, or formal planning specifications) between natural-language reasoning and executable spatial inference, replacing the hand-crafted parsers of §\ref{sec:p3_s7_1} with foundation-model-generated symbolic structure.
Two sub-patterns differ in when the symbolic machinery is used.
In the first it is called at inference time: NS3D~\cite{hsu2023ns3d} has an LLM translate a 3D referring expression into a first-order logic program that a neural executor grounds against the scene's objects and relations on ReferIt3D, and LEFT~\cite{hsu2023s} (§\ref{subsec:p3_s7_3_4}) makes the same programs differentiable so that concept grounding is learned rather than hand-specified.
In the second, symbolic knowledge enters only at training time as a constraint on a neural model: \citet{premsri2024neurosymbolic} train language models for reasoning over spatial language with spatial logic rules as auxiliary losses, and RelateSeg~\cite{lin2025neuro} extracts relations such as $\langle\text{cat}, \text{to-right-of}, \text{person}\rangle$ from open-vocabulary segmentation outputs, writes them as first-order logic, and relaxes them with fuzzy logic into a single auxiliary loss that improves mIoU on four benchmarks without adding parameters.
The Scene Language~\cite{zhang2025scenelanguage} (§\ref{subsec:p3_s8_2_5}) is the representation such hybrids increasingly target: programs for structure, words for semantics, embeddings for appearance.
The symbolic layer buys consistency and compositional guarantees; its price is a closed predicate vocabulary that the neural side constantly exceeds, a point taken up under Challenges and Trends below.

\paragraph{Foundation Models with Spatial Memory and World Models.}
A fourth pattern gives the foundation model a persistent spatial state that it does not have to carry in its context window: an explicit memory (episodic trajectories, topological maps, 3D layouts) or a learned world model that predicts how the scene changes under action.
The hybrid question is the interface between the two.
MSNav~\cite{liu2025msnav} splits zero-shot vision-and-language navigation into a memory module, a spatial module, and a decision module, with a small fine-tuned LLM extracting the landmarks and relations the spatial module needs, and reaches state-of-the-art zero-shot results on R2R and REVERIE.
MindExplore~\cite{li2025towardsvla} stacks a reasoning layer (task-specific chains of thought that emit meta-actions) over an acting layer (a mixture of policy experts with a multimodal diffusion policy) and uses a memory mechanism as the feedback channel between the two, running at 30~FPS in dynamic sand environments.
The 3D spatial memories of §\ref{subsec:p3_s7_3_6} (3DSPMR~\cite{cai2025vision}, Embodied-RAG~\cite{xie2024embodiedrag}, Meta-Memory~\cite{mao2025meta}) and the world-model coupling of §\ref{subsec:p3_s7_3_8} are the memory- and simulator-side components these hybrids draw on; MindJourney's test-time gain on SAT~\cite{yang2026mindjourney} (about eight points of top-1 accuracy from imagined viewpoints alone, with the VLM untouched) is the cleanest quantitative evidence that a world model used as a tool closes part of the gap.
In each case the foundation model keeps high-level inference, instruction following, and semantic grounding, and offloads persistence, geometry, and dynamics to the dedicated module.

\begin{figure}[!ht]
    \centering
    \includegraphics[width=0.8\linewidth]{figures/part3/section7/agentic_hybrid_orchestration.pdf}
    \caption{Illustration of the agentic hybrid orchestration pattern: a foundation-model controller dynamically selects perception, world-modeling, symbolic-reasoning, memory, and action tools within a closed-loop perceive--reason--act--reflect cycle.}
    \label{fig:p3_s7_4_agentic_hybrid}
\end{figure}

\paragraph{Agentic Orchestration of Hybrid Components.}
The fifth pattern makes the orchestration function itself the object of design: a foundation model acts as a deliberative controller that decides, at each step, which of the perception, world-model, symbolic, memory, and action modules to invoke, within a closed perceive--reason--act--reflect loop (Figure~\ref{fig:p3_s7_4_agentic_hybrid}).
It differs from the agentic reasoning of §\ref{subsec:p3_s7_3_7} in what is being scheduled: there, tools and memory; here, additionally symbolic engines and physics simulators, with the scheduler $\pi$ of Eq.~\eqref{eq:p3_s7_4_hybrid} itself a candidate for learning.
VADAR~\cite{marsili2025visual} (§\ref{subsec:p3_s7_3_4}) is the program-synthesis version of such a controller, generating the API it then calls.
``See it. Say it. Sorted''~\cite{zhang2025seeitsayit} is a training-free version for sketch-to-diagram generation: a critic VLM proposes a few qualitative relational edits, several LLMs synthesize candidate SVG programs, and a judge VLM selects among them, which keeps global constraints such as alignment and connectivity intact where numerical estimation would drift.
\citet{manh2025mindmeetspace} arrive at the same module inventory from neuroscience, arguing that such a controller should sit over multimodal sensing, multisensory integration, egocentric--allocentric conversion, a cognitive map, spatial memory, and spatial reasoning.
The gain is compositionality and long-horizon consistency that no static architecture provides; the cost is the latency and attribution problems discussed below.

% 之后
% 具体地，今年来，研究者越来越倾向于将大型基础模型 (foundation models) 与强化学习 (reinforcement learning, RL) 结合，以构建具备 spatial reasoning 能力的混合系统。
% 例如，最新工作 Embodied-R: Collaborative Framework for Activating Embodied Spatial Reasoning in Foundation Models via Reinforcement Learning (简称 Embodied-R) 就是一个典型代表。该工作将大型视觉—语言模型 (VLM) 用于感知 (visual observation) —— 提取视觉特征与空间线索；同时用较小的语言模型 (LM) 作为推理 /决策模块 (reasoner / planner)，并通过强化学习训练 (设计 reward 考量 think-answer 逻辑一致性) 实现 “slow-thinking” (慢思考) 的空间推理。
% 这些工作表明：将基础模型 (VLM / LLM) 与 world-model + RL 结合，是当前实现混合空间推理的最有前景路径之一 — 尤其适合应对复杂、动态、物理 /几何约束强的现实场景 (例如 3D 布局、具身导航、物体交互、场景理解等)。
% 此外，从更理论／宏观层面来看，有分析指出当前基础模型 (foundation models) 在决策 (decision-making) 与空间 /世界智能 (world-model) 方面存在结构性空缺。因此，通过 RL、world-model、模块化 /混合架构 (hybrid architecture) 来弥补这一空缺，有助于未来构建真正具有空间智能 (spatial intelligence) 的通用系统 (general-purpose agents / embodied AI / spatially-aware AI)。

%   通过将混合架构纳入智能体 (agent) 或规划 + 行动 (planning + action) 回路 (loop) 中，我们可以构造更为通用、灵活、可扩展 (scalable) 的空间智能系统。这种 “Agentic Hybrid Orchestration” 的设计核心包括以下几个关键维度：
% LLM / VLM 作为 Orchestrator (协调者／指挥者)：高层语言 /符号 /规划模块 (LLM) 负责决定 “何时感知 (perceive)”、 “如何构建或更新 world-model (scene graph / layout / spatial representation)”、 “何时进行推理 /规划 (plan)” 以及 “何时发起行动 /布局 /生成 / manipulation”。
% 工具选择 (tool selection) 与顺序 (sequencing)：系统可以预定义或动态维护一组 “工具 (tools)” —— 比如视觉感知模块、world-model builder、符号 /语言推理模块、布局 /渲染 模块 —— LLM 负责根据当前任务与上下文选择合适工具并安排调用顺序。
% 反馈驱动 (feedback-driven) 的迭代 (iterative) 推理 (reasoning)：通过 RL (或其他优化方式) 将最终目标 (task success, spatial consistency, physical plausibility, layout realism / utility) 的反馈 (reward / loss) 传播回 orchestrator 和各工具调用决策。
% 这种 Agentic Hybrid Orchestration 的好处是，它将混合架构系统转化为一个动态、交互式、多轮 (multi-step) 的代理 (agent) 回路 (loop)，从而具备长期规划、空间结构一致性、物理可行性、泛化 与适应性 — 正是混合空间推理 (Hybrid Spatial Reasoning) 的目标。

\subsection{Benchmarks}
\label{subsec:p3-sec7-4-benchmarks}
Hybrid systems are evaluated end to end on the benchmark of their target task rather than module by module, and the benchmark is often one the paper introduces.
Among the systems above, Embodied-R reports on UrbanVideo-Bench and VSI-Bench, MindJourney on SAT, NS3D on ReferIt3D, MSNav on R2R and REVERIE, and PhysVLM on RoboVQA and OpenEQA alongside its own EQA-phys; LLMPhy ships TraySim, CLEVRER-LLMPhy, and Real-TraySim, MindExplore ships SandGo-1k and SandThink-21k, and DrawingBench~\cite{kim2025drawingbench} is built specifically for agentic hybrids, scoring an LLM's sequence of low-level mouse actions against eight rule-based criteria for the spatial configuration it was asked to draw.
Physics-side components are still checked against PHYRE and ShapeStacks (§\ref{subsec:p3_s7_1_2}) and symbolic components against the constraint- and rule-based suites of §\ref{subsec:p3_s7_1_4}--§\ref{subsec:p3_s7_1_5}, but no current benchmark attributes an end-to-end score to the modules that produced it; that gap is taken up under Challenges and Trends below and in the cross-paradigm survey of §\ref{sec:p3_s7_5} (Table~\ref{tab:benchmark_comparison}, whose 15-row agentic subset is Table~\ref{tab:spatial_reasoning_benchmark} of §\ref{subsec:p3_s7_3_7}).

\subsection{Challenges and Trends}
\label{subsec:p3-sec7-4-challenges-and-trends}
Hybrid spatial reasoning has produced architectures that combine complementary inductive biases in principled ways, but several recurring limitations originate in the module-composition framework itself rather than in any individual module.

\textbf{Frame, Scale, and Geometric Consistency.}
Unlike generic modular systems, spatial hybrids must reconcile heterogeneous \emph{spatial} representations across modules: perception emits camera or egocentric frames, map and memory modules operate in world or allocentric coordinates, symbolic reasoners abstract to topological relations, and physics engines expect metric units. Misaligned reference frames, scales, or conventions silently corrupt downstream reasoning, and because each module reports confidence in its own terms, small upstream pose, depth, or scale errors propagate and compound across the pipeline without a coherent notion of uncertainty. Maintaining frame, scale, and geometric consistency is thus a first-order obstacle with no analogue in non-spatial hybrid systems.

\textbf{Interfaces, Latency, and Attribution.}
The strength of a hybrid system is bounded by its worst module interface. Symbolic-to-neural translation is particularly brittle when the symbolic side uses closed-vocabulary predicates while the neural side proposes open-vocabulary ones, and the interface design becomes a hidden first-class hyperparameter that current literature seldom isolates. Each module invocation also adds wall-clock latency that compounds across a perceive--reason--act--reflect loop, so high-frequency manipulation and navigation may not be reachable through orchestration alone. And when a hybrid fails, attributing the failure to a specific module is non-trivial; this complicates iteration and is often cited as a practical reason why end-to-end models can outperform hybrid stacks on aggregate benchmarks despite offering weaker per-module guarantees.

The orchestration function $\pi$ is migrating from hand-coded rules to \emph{learnable orchestration policies} trained with reinforcement or imitation learning, which blurs the boundary with §\ref{subsec:p3_s7_3_7}.
\emph{Differentiable hybrid stacks}, built from differentiable physics simulators, differentiable optimization layers (§\ref{subsec:p3_s7_2_4}), and differentiable retrieval, make end-to-end training of a hybrid viable and remove the historical interface penalty, offering a path through the precision--grounding asymmetry that motivates this section.
Evaluation is changing too, toward the \emph{holistic and diagnostic}: in response to the attribution problem, emerging practice pairs aggregate task metrics with per-module diagnostic probes and counterfactual module-swap ablations that isolate where a hybrid system succeeds or fails, a direction consolidated in the cross-paradigm benchmark survey of §\ref{sec:p3_s7_5}.

% \newpage
z\section{Spatial Reasoning Benchmarks}
\label{sec:p3_s7_5}

\FramedInitial{E}{arly} spatial-reasoning benchmarks positioned the model as a detached observer. CLEVR~\cite{johnson2017clevr} and VSR~\cite{liu2023visual} ask for relative geometric relations from a static third-person viewpoint, and their 2D inputs leave depth ambiguous, a weakness that 3DSRBench~\cite{ma20253dsrbench}, MMSI-Bench~\cite{yang2025mmsi}, DSI-Bench~\cite{zhang2025dsi}, and STI-Bench~\cite{li2025sti} now probe directly.
SQA3D~\cite{ma2022sqa3d}, OpenEQA~\cite{majumdar2024openeqa}, and EmbSpatial-Bench~\cite{du2024embspatial} instead place the model inside the scene, so that it must self-localize before it can answer, and ScanQA~\cite{azuma2022scanqa}, ScanRefer~\cite{chen2020scanrefer}, and Multi3DRefer~\cite{zhang2023multi3drefer} ground queries in point clouds and RGB-D data, where the answer is a coordinate rather than a relation.
The reasoning demanded has deepened in parallel: IntPhys~\cite{riochet2018intphys} and Physion~\cite{Bear2021Physion} test whether a model expects gravity, collision, and friction to hold over time, and OpenEQA and RealWorldQA~\cite{xai2024grok15v} ask functional questions (``Can I navigate through this gap?'', ``Where would I find a cloth to clean a spill?'') in which geometry matters only through what it affords.

Data provenance splits the field. Simulators, the sources behind SpatialVLM~\cite{chen2024spatialvlm}, AI2-THOR~\cite{kolve2017ai2thor}, and Habitat~\cite{szot2021habitat2}, yield unlimited, perfectly annotated data; real-world collections such as ScanNet~\cite{dai2017scannet} and OpenEQA supply the clutter, occlusion, and sensor noise that simulation does not reproduce.
The resulting practice is sim-to-real: pre-train on large synthetic corpora for general spatial logic, then fine-tune on the scarcer real data.
Annotated 3D data remain scarce and fine-grained metric estimation remains hard for large models; both limits recur throughout the table below.

% \wsq{总结现有的有关spatial reasoning 相关的benchmarks}


\subsection{Comparative Analysis}
\label{subsec:p3-sec7-5-comparative-analysis}

Table~\ref{tab:benchmark_comparison} compiles 61 spatial-reasoning benchmarks across four categories (physical-focused, rooted in §\ref{subsec:p3_s7_1_2}; vision-based, overlapping with §\ref{subsec:p3_s7_3_1}; 3D-focused, rooted in §\ref{subsec:p3_s7_3_6} and §\ref{subsec:p3_s7_3_8}; and robot-environment, rooted in §\ref{subsec:p3_s7_3_7} and §\ref{sec:p3_s7_4}).
A narrower 15-row subset focused specifically on agentic spatial reasoning is provided in Table~\ref{tab:spatial_reasoning_benchmark} of §\ref{subsec:p3_s7_3_7}; this subset filters Table~\ref{tab:benchmark_comparison} to image-and-video benchmarks that exercise agentic deployments.

\begin{table*}[!t]
  \caption{Spatial-reasoning benchmarks surveyed in this chapter, grouped by what they test. \texttt{Number} is the evaluation-set size where the source paper documents one (train / test when both are given, scene counts for simulators, $\sim$ for approximate figures). \texttt{Perspective} marks explicitly designed ego (first-person) and exo (third-person) settings; \textit{Undeclared} marks attributes the source paper does not state. \texttt{Data Source} and \texttt{Task Type} follow the source papers.}
  \label{tab:benchmark_comparison}
  \centering
  \fontsize{9}{10}\selectfont
  \setlength{\tabcolsep}{1.4mm}
  \resizebox{\linewidth}{!}{
  \begin{tabular}{@{}lcccccccc}
    \toprule
      \multirow{2}{*}{\textbf{Benchmark}}
      & \multirow{2}{*}{\textbf{Number}}
      & \multirow{2}{*}{\textbf{Modality}}
      & \multicolumn{2}{c}{\textbf{View Type}}
      & \multicolumn{2}{c}{\textbf{Perspective}}
      & \multirow{2}{*}{\textbf{Data Source}}
      & \multirow{2}{*}{\textbf{Task Type}} \\
      \cmidrule(lr){4-5}
        \cmidrule(lr){6-7}
        & &
        & \textbf{Single} & \textbf{Multi}
        & \textbf{Ego} & \textbf{Exo}
        & \\
    \midrule
    \rowcolor{lightgrey}\multicolumn{9}{c}{$\bullet$ \textit{Physical Focused Spatial Reasoning}} \\
    {IntPhys}~\cite{riochet2018intphys} & 15k / 4.3k & 3D   & \checkmark & \crossmark    & \crossmark & \checkmark   & Synthetic    & VoE plausibility\\
    {IntPhys 2}~\cite{bordes2025intphys2}& 1,416    & 3D   & \checkmark & \crossmark    & \crossmark & \checkmark       & Synthetic    & VoE plausibility\\
    {Physion}~\cite{Bear2021Physion}  & $\sim$1.2k  & 3D    & \checkmark & \crossmark  & \crossmark & \checkmark   & Synthetic & Physical prediction / choice \\
    {Physion$++$}~\cite{tung2023physion} & $\sim$10k  & 3D  & \checkmark & \crossmark   & \crossmark & \checkmark  & Synthetic  & Physical prediction / choice\\
    {Physical Concepts Dataset}~\cite{piloto2020learn}  & 300k    & 3D   & \checkmark & \crossmark   & \crossmark & \checkmark  & Synthetic   & classification\\
    {Physics-IQ}~\cite{motamed2025generative} & 396  & Video     & \checkmark & \crossmark   & \crossmark & \checkmark  & Real   & Generative physical prediction \\
    {SPhyR}~\cite{siedler2025sphyr} & 1,296 & Image & \checkmark & \crossmark & \multicolumn{2}{c}{\cellcolor{lightgray}\textit{Undeclared}} & Synthetic & Physical prediction / infill \\
    \hline
    \rowcolor{lightgrey}\multicolumn{9}{c}{$\bullet$ \textit{Vision-Based Spatial Reasoning}} \\
    CLEVR \citep{johnson2017clevr} & 853k & Image & \checkmark & \crossmark & \crossmark & \checkmark & Simulation & VQA \\
    VSR \citep{liu2023visual}& 10,972 & Image & \checkmark & \crossmark & \crossmark & \checkmark & Real & VQA \\
    What's Up \citep{KamathHC23a}  & $\sim$8,000 & Image & \checkmark & \crossmark & \crossmark & \checkmark & Real & VQA\\
    RocketScience~\citep{hoehing2025RocketScience} & 482 & Image & \checkmark & \crossmark  & \crossmark & \checkmark & Real & VQA \\
    SpaRE \citep{abs-2504-20648} & 3.4M & Image &  \checkmark & \crossmark & \crossmark & \checkmark & Real/Sim & VQA \\
    RealWorldQA~\cite{xai2024grok15v} & 765 & Image & \checkmark & \crossmark & \checkmark & \crossmark & Real & VQA \\
    3DSRBench~\cite{ma20253dsrbench} & 2,772 & Image & \checkmark  & \crossmark & \crossmark  &  \checkmark & Real/Sim & VQA\\
    InternSpatial-Bench~\cite{deng2025internspatial} &  6,008 & Image  & \checkmark  & \crossmark & \crossmark  &  \checkmark  & Real & VQA \\
    OmniSpatial~\cite{jia2025omnispatial} & $\sim$8,400 & Image & \checkmark & \crossmark & \checkmark & \checkmark   &  Real & VQA\\
    Spatial457~\cite{wang2025spatial457} & 23,752 & Image & \checkmark & \crossmark & \multicolumn{2}{c}{\cellcolor{lightgray}\textit{Undeclared}}  & Real & VQA\\
    All-Angles Bench~\cite{yeh2025seeing} & $\sim$2,100 & Image  & \crossmark & \checkmark & \multicolumn{2}{c}{\cellcolor{lightgray}\textit{Undeclared}} & Real  & VQA \\
    EmbSpatial-Bench~\cite{du2024embspatial} & 3,640 & Image & \checkmark  & \crossmark & \checkmark & \crossmark &  Real & VQA \\
    MMSI-Bench~\cite{yang2025mmsi} & 1,000 & Image & \crossmark & \checkmark  & \multicolumn{2}{c}{\cellcolor{lightgray}\textit{Undeclared}} &  Real & VQA\\
    MindCube~\cite{yin2025spatial} & 21,154 & Image & \crossmark & \checkmark & \multicolumn{2}{c}{\cellcolor{lightgray}\textit{Undeclared}}  & Simulation  & VQA \\
    BLINK~\cite{fu2024blink} & 3,807 & Image & \checkmark & \crossmark & \crossmark & \checkmark & Real & VQA \\
    NuScenes-SpatialQA~\cite{tian2025nuscenesspatialqa} & 3.3M+ & Video & \checkmark  & \checkmark  & \checkmark  & \checkmark  & Real & VQA \\
    VSI-Bench~\cite{yang2025vsibench} & 5,130 & Video & \checkmark & \crossmark & \checkmark & \crossmark & Real  & VQA\\
    VLM4D~\cite{zhou2025vlm4d} & $\sim$1,800 & Video & \checkmark & \crossmark &  \checkmark & \checkmark &  Real & VQA \\
    STI-Bench~\cite{li2025sti} & 2,064 & Video  & \checkmark & \crossmark & \checkmark & \checkmark & Real & VQA\\
    DSI-Bench~\cite{zhang2025dsi} & $\sim$1,700 & Video  & \checkmark & \crossmark & \multicolumn{2}{c}{\cellcolor{lightgray}\textit{Undeclared}} & Real & VQA\\
    SITE~\cite{wang2025site} & 8,068  & Video & \checkmark & \crossmark & \multicolumn{2}{c}{\cellcolor{lightgray}\textit{Undeclared}} & Real/Sim & VQA \\
    RoboSpatial \citep{SongBTTSB25} & 3M & Image/3D & \checkmark & \crossmark & \checkmark & \crossmark & Real/Sim & VQA \\
    OSD \citep{cheng2024spatialrgpt} & \textit{Undeclared} & Image & \checkmark & \crossmark & \crossmark & \checkmark & Real & VQA \\
    SpatialVLM \citep{chen2024spatialvlm} & 2B & Image & \checkmark & \crossmark & \crossmark & \checkmark & Real/Sim & VQA \\
    ProVision \citep{abs-2412-07012} & 65k & Image & \checkmark & \crossmark & \crossmark & \checkmark & Simulation & VQA \\
    WhatIsWhere \citep{LiangZZZ16} & 6,000 & Image & \checkmark & \crossmark & \crossmark & \checkmark & Real & VQA \\
    SURDS~\cite{guo2024surds} & 9,250 & Image & \checkmark & \crossmark & \checkmark & \crossmark & Real & VQA \\
    EarthSpatialBench~\cite{xu2026earthspatial} & 325k+ & Image & \checkmark & \crossmark & \crossmark & \checkmark & Real & VQA \\
    PureSpace~\cite{li2026purespace} & 1,200 & Image & \checkmark & \crossmark & \multicolumn{2}{c}{\cellcolor{lightgray}\textit{Undeclared}} & Simulation & VQA \\
    SAT~\cite{ray2024sat} & 175k / 150 & Image & \crossmark & \checkmark & \checkmark & \checkmark & Simulation & VQA \\
    SpatiaLQA~\cite{xie2026spatialqa} & 9,605 & Image & \checkmark & \crossmark & \crossmark & \checkmark & Real & VQA \\
    Q-Spatial++~\cite{wang2025svqa} & 101 & Image & \checkmark & \crossmark & \crossmark & \checkmark & Real & VQA \\
    \hline
    \rowcolor{lightgrey}\multicolumn{9}{c}{$\bullet$ \textit{3D Focused Spatial Reasoning}} \\
    ScanQA~\cite{azuma2022scanqa} & 41,363 & 3D & \crossmark & \checkmark & \crossmark & \checkmark & Real (ScanNet) & 3D VQA \\
    SQA3D~\cite{ma2022sqa3d} & 33,460 & 3D & \checkmark & \crossmark & \checkmark & \crossmark & Real (ScanNet) & 3D VQA \\
    MMScan QA~\cite{lyu2024mmscan} &  295K & 3D+Image & \checkmark & \checkmark & \crossmark & \checkmark & Real (EmbodiedScan) & 3D VQA \\
    OpenEQA~\cite{majumdar2024openeqa} & 1,632 & Video/3D & \checkmark & \crossmark & \checkmark & \crossmark & Real (HM3D, ScanNet) & Embodied QA (EM / A splits) \\
    Space3D-Bench~\cite{szymanska2024space3dbench} & 1,000 & 3D & \checkmark & \crossmark & \checkmark & \crossmark  & Real(Replica) & 3D QA \\
    Scan2Cap~\cite{chen2021scan2cap} & 39,775 & 3D & \crossmark & \checkmark & \crossmark & \checkmark & Real (ScanNet) & 3D Captioning \\
    MMScan Captioning~\cite{lyu2024mmscan} & 172k & 3D+Image & \checkmark & \checkmark & \crossmark & \checkmark & Real (EmbodiedScan) & 3D Captioning \\
    ScanRefer~\cite{chen2020scanrefer} & 51,583 & 3D & \crossmark & \checkmark & \crossmark & \checkmark & Real (ScanNet) & 3D Grounding \\
    Multi3DRefer~\cite{zhang2023multi3drefer} & 42,735 & 3D & \crossmark & \checkmark & \crossmark & \checkmark & Real (ScanNet) & 3D Grounding \\
    SR3D~\cite{achlioptas2020referit3d} & 83,572 & 3D & \crossmark & \checkmark & \crossmark & \checkmark & Real (ScanNet) & 3D Grounding \\
    NR3D~\cite{achlioptas2020referit3d} & 41,503 & 3D & \crossmark & \checkmark & \crossmark & \checkmark & Real (ScanNet) & 3D Grounding \\
    SpaceEra++~\cite{guan2026spaceera} & 925k & Video/3D & \crossmark & \checkmark & \checkmark & \crossmark & Synthetic & 3D VQA \\
    Ego3D-Bench~\cite{gholami2025ego3d} & 8,600+ & Image & \crossmark & \checkmark & \checkmark & \crossmark & Real & VQA \\
    SpatialScore~\cite{wu2025spatialscore} & 28,093 (Hard: 1,400) & Image/3D & \crossmark & \checkmark & \multicolumn{2}{c}{\cellcolor{lightgray}\textit{Undeclared}} & Real/Sim & VQA \\
    iVISPAR~\cite{mayer2025ivispar} & 300 & 3D Env & \checkmark & \crossmark & \crossmark & \checkmark & Simulation & Interactive reasoning \\
    \hline
    \rowcolor{lightgrey}\multicolumn{9}{c}{$\bullet$ \textit{Robot}} \\
    AI2-THOR~\cite{kolve2017ai2thor} & 120 (Scenes) & 3D Env & \crossmark & \checkmark & \checkmark & \crossmark & Simulation & Nav/Manip \\
    TDW~\cite{gan2020threedworld} & \textit{Infinite} & 3D Env & \crossmark & \checkmark & \checkmark & \checkmark & Simulation & Physical Sim \\
    Habitat 2.0~\cite{szot2021habitat2} & 100 (Scenes) & 3D Env & \crossmark & \checkmark & \checkmark & \crossmark & Simulation & Navigation \\
    iGibson 2.0~\cite{Li2021iGibson2} & 15 (Scenes) & 3D Env & \crossmark & \checkmark & \checkmark & \crossmark & Simulation & Interactive Nav \\
    Habitat 3.0~\cite{puig2023habitat3} & 211 (Scenes) & 3D Env & \crossmark & \checkmark & \checkmark & \crossmark & Simulation & Social Nav \\
    R2R~\cite{anderson2018vision} & 21,567 & Image/Path & \crossmark & \checkmark & \checkmark & \crossmark & Real & Navigation \\
    RxR~\cite{ku2020roomacrossroom} & 126,069 & Image/Path & \crossmark & \checkmark & \checkmark & \crossmark & Real & Navigation \\
    \bottomrule
  \end{tabular}
  }
\end{table*}

\emph{A modality migration} runs from the top of the table to the bottom: the physical-focused benchmarks rely on synthetic 3D renderings, the vision-based block is image- and video-heavy with growing 3D extensions, the 3D-focused block is built on point clouds and RGB-D scans, and the robot block has moved largely to 3D environments.
\emph{The Perspective axis is imbalanced}: a large fraction of vision-based benchmarks carry \textit{Undeclared} ego/exo annotations, the physical-focused benchmarks lean exo (a third-person observer judging physical plausibility), and the robot block is entirely ego; the field is migrating from spectator-style observation to situated cognition faster than benchmark annotations record it.
\emph{The Data-Source spectrum is bifurcated} between predominantly synthetic or simulated (top block and robot block) and predominantly real (middle blocks), with the Real/Sim benchmarks (SpaRE, 3DSRBench, RoboSpatial, SpatialVLM, SITE) as the bridge that the sim-to-real training practice of the opening relies on.

\subsection{Future Directions}
\label{subsec:p3-sec7-5-future-directions}

\paragraph{Unified Spatial-Cognition Benchmark.}
No single benchmark in Table~\ref{tab:benchmark_comparison} exercises a model along all of perception fidelity, symbolic precision, memory persistence, and long-horizon action conditioning simultaneously. 
A unified benchmark covering all four §\ref{sec:reasoning} paradigms (§\ref{sec:p3_s7_1} rule-based, §\ref{sec:p3_s7_2} learning-based, §\ref{sec:p3_s7_3} foundation-model-driven, §\ref{sec:p3_s7_4} hybrid) would close the evaluation-attribution challenge flagged under Challenges and Trends in §\ref{sec:p3_s7_4}.

\paragraph{Cross-Modal Evaluation.}
The current axes (Modality, View, Perspective, Data Source, Task Type) under-specify the cross-modal demand of contemporary spatial reasoning. 
Future benchmarks should explicitly probe modality-to-modality transfer (\eg, 2D-trained models evaluated on 3D queries; text-conditioned policies evaluated on video-grounded tasks), so that the modality-mismatch challenge flagged in §\ref{subsec:p3_s7_3_6} (text-vs.-image-vs.-3D retrieval) becomes a first-class evaluation axis rather than a per-paper case study.

\paragraph{Embodied Long-Horizon Coverage.}
The Robot block in Table~\ref{tab:benchmark_comparison} covers navigation thoroughly but under-represents long-horizon manipulation and multi-room interactive tasks. 
The §\ref{subsec:p3_s7_3_7} 15-row subset (Table~\ref{tab:spatial_reasoning_benchmark}) further restricts attention to image-and-video benchmarks; the embodied long-horizon suites ALFWorld and VirtualHome (§\ref{subsec:p3_s7_3_4}) and cross-embodiment evaluation on Open X-Embodiment (§\ref{subsec:p3_s7_3_2}) are the nearest existing coverage, and standardizing on them would close the long-horizon-spatiotemporal-coherence challenge flagged under Challenges and Trends in §\ref{subsec:p3_s7_3_7}.

\paragraph{From Descriptive to Interactive Spatial Reasoning.}
The Task Type column of Table~\ref{tab:benchmark_comparison} is dominated by visual question answering: the overwhelming majority of benchmarks ask a model to \emph{describe} static spatial relations (``where is the object''), and only the small Robot block evaluates reasoning that is conditioned on the agent's own actions.
Yet the trajectory traced in this chapter, from passive observation to situated interaction, calls for benchmarks that test \emph{interactive} spatial reasoning: predicting how an action changes the spatial state, judging affordances and reachability (``can I fit, reach, or grasp this''), and reasoning counterfactually about manipulation outcomes.
Such evaluation couples the action-conditioned reasoning of §\ref{subsec:p3_s7_2_5} with the embodied hybrid pipelines of §\ref{sec:p3_s7_4}, extending the descriptive benchmarks surveyed here toward the manipulation- and affordance-grounded settings that embodied deployment ultimately demands.



\section{Summary and Outlook}
\label{sec:p3_s7_6}

This chapter followed spatial reasoning along the spectrum on which it developed: inference over explicit geometric, physical, topological, logical, and causal models (§\ref{sec:p3_s7_1}); reasoning acquired from data (§\ref{sec:p3_s7_2}); reasoning elicited from large pre-trained models (§\ref{sec:p3_s7_3}); hybrids that combine these components (§\ref{sec:p3_s7_4}); and the benchmarks that measure them (§\ref{sec:p3_s7_5}).

\textbf{Precision and grounding sit at opposite ends of the spectrum.} Rule-based methods are inspectable and can certify feasibility or safety, but they assume exact, fully specified inputs and a closed vocabulary of predicates, and they break down where perception is noisy or a task falls outside the authored rules. Learning-based methods acquire spatial competence from data, yet often reach high benchmark accuracy through dataset-specific shortcuts that do not survive novel compositions (§\ref{subsec:p3_s7_2_1}), and they need far more interaction than classical planners. Foundation models bring broad semantic competence, but their visual encoders retain only coarse geometry, textual reasoning mistracks coordinates over multi-step chains (§\ref{subsec:p3_s7_3_3}), and vision-language-action models often interpret an instruction correctly and execute it incorrectly (§\ref{subsec:p3_s7_3_2}). Hybrid systems aim to pair the semantic reach of the last family with the metric precision and verifiability of the first two.

\textbf{The failures concentrate at interfaces.} Across tool-augmented, memory-augmented, and hybrid reasoning, the recurring limitations arise less within any single component than where components meet: in forming a query that a memory or tool can answer, in translating between closed and open vocabularies, and in keeping coordinate frames and scales consistent across modules (§\ref{sec:p3_s7_4}).

\textbf{Evaluation lags the shift toward situated reasoning.} The field is moving from questions posed to a detached observer toward agents that must localize themselves, gather evidence, and act, yet no current benchmark exercises perception fidelity, symbolic precision, memory persistence, and long-horizon action together, and few attribute an end-to-end score to the components that produced it (§\ref{sec:p3_s7_5}).

Reasoning in this chapter takes the evidence of Chapter~\ref{chapter:3:sec:6} as given. Chapter~\ref{sec:generation} turns to the inverse problem of producing spatial content, and its world models (§\ref{sec:p3_s8_6}) supply the simulators that world-model-based reasoning (§\ref{subsec:p3_s7_3_8}) consults.

\newpage


\chapter{Spatial Reconstruction and Generation}
\label{sec:generation}


\FramedInitial{S}{patial} reconstruction and generation is the productive counterpart to spatial perception and reasoning: rather than analyzing a scene that is given, it synthesizes spatial content, recovering or creating the geometry, appearance, and dynamics of objects, scenes, and worlds from conditions such as text, images, partial observations, or explicit constraints.
Where perception (Chapter~\ref{chapter:3:sec:6}) maps observations into structured spatial representations and reasoning (Chapter~\ref{sec:reasoning}) draws inferences over them, generation closes the loop by mapping intent back into concrete spatial artifacts, a capability central to embodied agents, content creation, simulation, and world modeling.

The chapter is organized along an axis of increasing output structure and dimensionality, from flat images to volumetric objects and scenes, to video, embodied motion, and full four-dimensional space-time content, and finally to world models and to programmatic generation that emits interpretable, executable structure. What holds this range together is a single recurring backbone: an encoder maps conditions into a latent, a generator samples structure, and a decoder or renderer produces the final artifact, specialized to each output domain. Read this way, the subfields are less seven separate problems than one conditional-generation abstraction instantiated against progressively richer codomains.


\section{2D Spatial Generation}
\label{sec:p3_s8_1}

\FramedInitial{T}{he} 2D spatial generation encompasses a broad spectrum of tasks that require synthesizing or editing images under explicit spatial constraints.
These constraints may be \emph{user-specified} (bounding boxes, masks, sketches, keypoints, depth maps) or \emph{derived} from scene understanding (layout graphs, semantic maps, relational structures), and the overarching goal is to produce images that are not only visually realistic but also \emph{faithful to the specified spatial configuration}: objects appearing in the correct locations, at the correct scales and orientations, while remaining globally coherent in lighting, perspective, and style.
2D generation is where the most mature \emph{semantic-to-pixel prior} in contemporary AI is taught to honor \emph{spatial structure}: it is the proving ground on which large pre-trained visual priors are forced to localize objects, follow geometric scaffolds, respect occlusion ordering, and obey region- or instruction-level edits, all without falling out of the implicit grammar of natural images.

The section follows the 2D pipeline in three steps that share one latent-generator framework.
§\ref{subsec:p3_s8_1_1} reviews foundational text-to-image generation, the semantic-to-pixel prior on which everything else here is built, and the three paradigms (adversarial, autoregressive, and diffusion with its flow-matching generalization) that parameterize it.
§\ref{subsec:p3_s8_1_2} adds a structural condition (layouts, masks, edges, depth, pose, scene graphs, region phrases) and asks how it is injected into a frozen or co-trained generator.
§\ref{subsec:p3_s8_1_3} replaces the empty canvas with an existing image and a mask, so that the generator must change what lies inside the mask while preserving what lies outside it.
% A short Summary (§\ref{subsec:p3_s8_1_summary}) compares the three along shared axes.

\input{part_3/section_8/p3_sec8_1/p3_sec8_1_1.tex}

\input{part_3/section_8/p3_sec8_1/p3_sec8_1_2.tex}

\input{part_3/section_8/p3_sec8_1/p3_sec8_1_3.tex}

% \subsection{Summary}
% \label{subsec:p3_s8_1_summary}

% Table~\ref{tab:p3_s8_1_subfield_comparison} places the three subfields side by side: the conditioning input each consumes, the sampler it instantiates, the mechanisms it surveys, and the benchmarks on which it reports.

% \begin{table}[t]
%     \centering
%     \fontsize{8}{10}\selectfont
%     \setlength{\tabcolsep}{1.0mm}
%     \caption{The three subfields of 2D spatial generation; each sampler is a successive specialization of Eq.~\eqref{eq:t2i_generator}.}
%     \label{tab:p3_s8_1_subfield_comparison}
%     \begin{tabular}{L{0.14\textwidth} L{0.19\textwidth} L{0.19\textwidth} L{0.19\textwidth} L{0.21\textwidth}}
%         \toprule
%         \textbf{Subfield} & \textbf{Input / sampler} & \textbf{Mechanisms} & \textbf{Canonical methods} & \textbf{Benchmarks / metrics} \\
%         \midrule
%         §\ref{subsec:p3_s8_1_1} Text-to-image generation & Text $\mathbf{x}$; $G_\theta:\mathcal{E}\!\to\!\mathcal{Z}$ (Eq.~\eqref{eq:t2i_generator}) & GAN; autoregressive; diffusion / flow matching; backbone, acceleration, and proprietary refinements & StackGAN, AttnGAN; DALL-E, Parti; LDM, Imagen, DiT, SD3 & FID, CLIPScore, HPSv2 / PickScore, VQAScore; GenEval, T2I-CompBench, HRS-Bench, HEIM, WISE \\
%         §\ref{subsec:p3_s8_1_2} Spatial-aware image generation & Text + structure $\mathbf{c}$; $G_\theta:\mathcal{E}\!\times\!\mathcal{C}'\!\to\!\mathcal{Z}$ (Eq.~\eqref{eq:t2i_spatial_objective}) & Side branch; adapter; grounded cross-attention; layout-aware cross-attention; inference-time guidance & ControlNet, T2I-Adapter, GLIGEN, LayoutDiffusion, BoxDiff & Channel adherence (IoU, edge recall, depth RMSE, PCK, relation accuracy); GenEval position; T2I-CompBench spatial \\
%         §\ref{subsec:p3_s8_1_3} Spatial inpainting and editing & Image $I$, mask $\mathcal{M}$, prompt, control; $G_\theta:\mathcal{Z}\!\times\!\{0,1\}^{H\times W}\!\times\!\mathcal{E}\!\times\!\mathcal{C}'\!\to\!\mathcal{Z}$ (Eq.~\eqref{eq:t2i_edit_objective}) & Masked denoising; inversion + rewrite; instruction-tuned; exemplar-guided; drag-based & RePaint, Blended Diffusion; Prompt-to-Prompt, Null-Text; InstructPix2Pix, Emu Edit; IP-Adapter; DragGAN & PSNR / LPIPS, edit locality, harmonization, identity preservation; EditBench, MagicBrush, Emu Edit \\
%         \bottomrule
%     \end{tabular}
% \end{table}

% One move structures the whole section: specialize the conditioning channel of $G_\theta$, from text alone (§\ref{subsec:p3_s8_1_1}) to text plus a structural scaffold (§\ref{subsec:p3_s8_1_2}) to scaffold plus a source image and mask (§\ref{subsec:p3_s8_1_3}).
% Each step is conservative, since the unconditional limit ($\mathbf{c}=\varnothing$ for §\ref{subsec:p3_s8_1_2}; $\mathbf{z}^{\text{in}}=\varnothing,\mathcal{M}=\mathbf{1}$ for §\ref{subsec:p3_s8_1_3}) recovers the preceding subfield, so the three sampler equations form a containment chain rather than a horizontal partition.
% The mechanism axes of the second and third subsections are, moreover, largely orthogonal to the generative core: side branches, adapters, grounded cross-attention, masked denoising, inversion, and drag-based optimization compose with LDM, DiT, and flow-matching backbones interchangeably, and production pipelines routinely stack a spatial-aware injection (Eq.~\eqref{eq:t2i_spatial_objective}) on top of an editing mechanism (Eq.~\eqref{eq:t2i_edit_objective}), so the three subfields are best read as one composable toolbox rather than as three disjoint task families.
% Evaluation has followed the same specialization, from a single FID against COCO for §\ref{subsec:p3_s8_1_1} circa~2020 to a task-specific tuple that adds channel adherence in §\ref{subsec:p3_s8_1_2} and reconstruction quality, edit locality, harmonization, and identity preservation in §\ref{subsec:p3_s8_1_3}.

% The extension of these 2D priors to \emph{volumetric} synthesis (turning a text prompt, an image, or a mask into a textured mesh, a NeRF, or a 3DGS scene) is the subject of §\ref{sec:p3_s8_2}. The conditional latent sampler of Eq.~\eqref{eq:t2i_generator} continues to play a central role there, but now as a \emph{prior to be distilled} via score-distillation sampling, multi-view consistency objectives, or volume-rendering supervision, rather than as the terminal generator.

% \newpage

\section{3D Spatial Reconstruction and Generation}
\label{sec:p3_s8_2}

\FramedInitial{T}{he} 3D spatial reconstruction and generation encompasses the tasks that turn images, video, sensor readings, or high-level prompts into an explicit 3D representation (a textured mesh, a signed-distance or occupancy field, a neural radiance field, or a 3D Gaussian field) of an object, a scene, or a complete navigable environment.
Whereas §\ref{sec:p3_s8_1} taught a 2D semantic-to-pixel before honor spatial structure on the image plane, this section lifts that prior off the plane and into the volumetric substrate on which embodied perception, simulation, and interaction ultimately operate.
This is where appearance and geometry are forced into a single, view-consistent representation: a model must not only produce plausible pixels but place matter correctly in $\mathbb{R}^3$ so that the result can be rendered from arbitrary viewpoints, physically simulated, navigated, and edited. It is the bridge between the flat generative priors of §\ref{sec:p3_s8_1} and the embodied world models, simulators, and digital twins that consume 3D structure downstream.

The five subsections are organized by the direction and source of the map into a shared 3D-representation space $\mathcal{R}$.
§\ref{subsec:p3_s8_2_1} starts with a single asset: 3D object synthesis turns the 2D prior of §\ref{sec:p3_s8_1} into a volumetric generator through 3D-native training, score distillation, multi-view lifting, or feed-forward reconstruction.
§\ref{subsec:p3_s8_2_2} scales that generator to a structured scene of placed objects and a room envelope.
§\ref{subsec:p3_s8_2_3} reverses the direction: reconstruction recovers a representation faithful to sensor observations, and meets synthesis exactly where the observations run out.
§\ref{subsec:p3_s8_2_4} routes through a 360-degree panorama as a structured 2D intermediate that is then lifted to 3D, and §\ref{subsec:p3_s8_2_5} closes the loop by editing an existing representation under user control.
% A Summary (§\ref{subsec:p3_s8_2_summary}) compares the five along shared axes.


\input{part_3/section_8/p3_sec8_2/p3_sec8_2_1.tex}


\input{part_3/section_8/p3_sec8_2/p3_sec8_2_2.tex}

\input{part_3/section_8/p3_sec8_2/p3_sec8_2_3.tex}

\input{part_3/section_8/p3_sec8_2/p3_sec8_2_4.tex}

\input{part_3/section_8/p3_sec8_2/p3_sec8_2_5.tex}


% \subsection{Summary}
% \label{subsec:p3_s8_2_summary}

% Table~\ref{tab:p3_s8_2_subfield_comparison} places the five subfields side by side: the input each consumes, the operator it instantiates over the shared 3D-representation space $\mathcal{R}$, the paradigms it surveys, the methods anchoring each, and the benchmarks on which it reports.

% \begin{table}[t]
%     \centering
%     \fontsize{8}{10}\selectfont
%     \setlength{\tabcolsep}{1.0mm}
%     \caption{The five subfields of 3D spatial reconstruction and generation; each operator acts on the shared 3D-representation space $\mathcal{R}$.}
%     \label{tab:p3_s8_2_subfield_comparison}
%     \begin{tabular}{L{0.13\textwidth} L{0.20\textwidth} L{0.18\textwidth} L{0.19\textwidth} L{0.22\textwidth}}
%         \toprule
%         \textbf{Subfield} & \textbf{Input / operator} & \textbf{Dominant paradigms} & \textbf{Canonical methods} & \textbf{Benchmarks / metrics} \\
%         \midrule
%         §\ref{subsec:p3_s8_2_1} 3D object synthesis & Text / image / multi-view $\to$ asset; $G_\theta^{\text{3D}}:\mathcal{C}'\!\to\!\mathcal{R}$ (Eq.~\eqref{eq:p3_s8_2_1_objective}) & 3D-native generation; SDS distillation; multi-view generation and lifting; feed-forward reconstruction and 3D latent diffusion & GET3D; DreamFusion; Zero-1-to-3, MVDream; LRM, TRELLIS, Hunyuan3D & GT23D-Bench, MATE-3D, Gen3DEval; CLIP alignment, Chamfer / F-score, LPIPS \\
%         §\ref{subsec:p3_s8_2_2} 3D scene synthesis & Text / layout / graph / sparse views / trajectory $\to$ scene tuple; $G_\theta^{\text{scene}}$ (Eq.~\eqref{eq:p3_s8_2_2_objective}) & Layout / scene-graph conditioning; view-to-3D; view-to-view; video-model-based & DiffuScene, GALA3D; DreamScene; LVSM, SEVA; CameraCtrl, ViewCrafter & 3D-FRONT, ScanNet, RealEstate10K; layout adherence, PSNR / LPIPS, FVD, pose error \\
%         §\ref{subsec:p3_s8_2_3} 3D scene reconstruction & Posed or unposed images, depth, point clouds $\to$ 3D; $G_\theta^{\text{rec}}:\mathcal{O}\!\to\!\mathcal{R}^{\text{rec}}$ (Eq.~\eqref{eq:p3_s8_2_3_objective}) & Iterative (SfM--MVS, NeRF, 3DGS); feed-forward; generative complete reconstruction & NeRF, 3DGS; DUSt3R, MASt3R, VGGT; ReconFusion, Bolt3D & DTU, Tanks-and-Temples, Mip-NeRF360; PSNR / SSIM / LPIPS, Chamfer, pose AUC \\
%         §\ref{subsec:p3_s8_2_4} 3D panorama synthesis & Text / NFoV image / layout / viewpoint $\to$ panorama $\to$ 3D; $G_\theta^{\text{pano}}$ $+$ lift $\Lambda$ (Eq.~\eqref{eq:p3_s8_2_4_objective}) & Panoramic image synthesis; panoramic video; panoramic reconstruction; panorama-based scene generation & PanFusion, CubeDiff; 360DVD, Imagine360; PanSplat; LayerPano3D, HunyuanWorld & Structured3D, SUN360, Matterport3D; FID / FAED, WS-PSNR, seam and loop consistency \\
%         §\ref{subsec:p3_s8_2_5} 3D scene editing & Existing asset $+$ edit intent $\to$ edited asset; $E_\theta:\mathcal{R}\!\times\!\mathcal{C}'\!\to\!\mathcal{R}$ (Eq.~\eqref{eq:p3_s8_2_5_objective}) & 2D-lifted iterative; multi-view-consistent guidance; 3D-native primitive; LLM-integrated; feed-forward / training-free & Instruct-NeRF2NeRF; GaussCtrl, VcEdit; 3DSceneEditor; ShapeLLM-Omni; VoxHammer & Edit3D-Bench, 3DEditVerse; Chamfer / F-score, CLIP$_{\text{dir}}$, GPTEval3D \\
%         \bottomrule
%     \end{tabular}
% \end{table}

% The subfields partition not by output domain, since all five produce an element of the same $\mathcal{R}$, but by the direction and source of the map: synthesis (§\ref{subsec:p3_s8_2_1}, §\ref{subsec:p3_s8_2_2}) maps high-level priors into $\mathcal{R}$, reconstruction (§\ref{subsec:p3_s8_2_3}) maps sensor observations into it, panorama synthesis (§\ref{subsec:p3_s8_2_4}) inserts a 2D intermediate before the map, and editing (§\ref{subsec:p3_s8_2_5}) maps $\mathcal{R}$ back into itself.
% Synthesis and reconstruction are not disjoint. As observations grow sparse, the reconstruction operator of Eq.~\eqref{eq:p3_s8_2_3_objective} must hallucinate the unobserved regions and collapses into the generative posterior of Eq.~\eqref{eq:p3_s8_2_1_objective}, so the complete-reconstruction regime of §\ref{subsec:p3_s8_2_3} and the generative regime of §\ref{subsec:p3_s8_2_1} approach one problem from opposite sides, and the formal distinction between the two operators dissolves into a single estimator that consumes whatever observations are available and falls back on a learned prior wherever they are not.
% Across all five subfields the 2D prior of §\ref{sec:p3_s8_1} is reused as a prior to be distilled rather than as a terminal generator (score distillation, multi-view and panorama diffusion, frozen 2D editors lifted into 3D), and all inherit the same failure mode, the 2D$\leftrightarrow$3D consistency gap (Janus artifacts, view drift, cubemap seams, edit flicker) whose suppression motivates the multi-view-consistency machinery surveyed throughout the section.
% Evaluation, finally, has not followed the convergence seen in §\ref{sec:p3_s8_1}: no cross-subfield suite exists, each subfield reports a paradigm-specific tuple over geometry, appearance, and constraint adherence, and a layout-adherence score from §\ref{subsec:p3_s8_2_2} cannot be ranked against a novel-view PSNR from §\ref{subsec:p3_s8_2_3}; a human-aligned protocol that spans the subfields is the clearest open need.

% The temporal extension of these volumetric priors (generating spatially \emph{and} temporally coherent 3D content, where a camera or the scene itself moves through time) is the subject of video-based spatial generation in §\ref{sec:p3_s8_3}.

% \newpage
\section{Video Spatial Generation}
\label{sec:p3_s8_3}

\FramedInitial{V}{ideo} spatial generation is the task of synthesizing a temporally coherent image sequence (a video clip) from a high-level prompt together with optional spatial, temporal, and other conditioning signals.
It has advanced substantially in recent years, evolving from early generative adversarial network (GAN)-based approaches~\cite{goodfellow2014gan} through autoregressive token models to the diffusion- and diffusion-transformer-based systems that now dominate the field~\cite{polyak2024movie,hong2022cogvideo,wan2025wan,sora2024}.
This rapid progress has stimulated a wide range of downstream tasks, extending beyond basic text-to-video generation to encompass fine-grained control mechanisms (spatial, \eg, image, pose, depth, landmark, sketch, and bounding box; temporal, \eg, optical flow, trajectory, camera, and motion; and other conditional controls, \eg, style, audio, bird's-eye view) as well as sophisticated video editing capabilities.

Within the arc of §\ref{sec:generation}, video spatial generation is the \emph{temporal} extension of the priors built in the preceding two sections.
Whereas §\ref{sec:p3_s8_1} learned a 2D semantic-to-pixel prior on a single image plane and §\ref{sec:p3_s8_2} lifted that prior into a static volumetric representation, this section adds the time axis: the model must place matter and appearance correctly \emph{and} make them evolve plausibly as the camera or the scene itself moves through time.
Appearance, geometry, and dynamics must now be \emph{jointly} consistent: a generated clip is judged not only frame by frame but on object permanence, motion plausibility, and multi-view coherence under camera motion, which makes video the natural substrate for the dynamic world models and simulators that consume spatial structure downstream.

Formally, we cast every system in this section as a conditional spatio-temporal sampler that extends the latent-generator framework of §\ref{subsec:p3_s8_1_1}.
Let $\mathcal{V}=[0,1]^{T\times H\times W\times 3}$ be the space of $T$-frame clips and $\mathcal{Z}^{\text{vid}}$ a compressed spatio-temporal latent space supplied by a video autoencoder that extends the image autoencoder of Eq.~\eqref{eq:t2i_autoencoder} along the time axis.
Writing $\phi$ for the text encoder of Eq.~\eqref{eq:t2i_generator} and $\kappa:\mathcal{C}\!\to\!\mathcal{C}'$ for a generic encoder of an auxiliary control signal $\mathbf{c}\in\mathcal{C}$, the core learning problem is to train
\begin{equation}
G_\theta^{\text{vid}}:\mathcal{E}\times\mathcal{C}'\rightarrow\mathcal{Z}^{\text{vid}},\qquad
\hat{\mathbf{z}}^{1:T}\sim p_\theta\!\big(\mathbf{z}^{1:T}\,\big|\,\phi(\mathbf{x}),\,\kappa(\mathbf{c})\big),
\label{eq:p3_s8_3_objective}
\end{equation}
so that the decoded clip $\hat{\mathbf{v}}=D(\hat{\mathbf{z}}^{1:T})$ is photorealistic, temporally coherent, and faithful to both the prompt $\mathbf{x}$ and the control $\mathbf{c}$.
Two limiting cases anchor the framework: with no auxiliary control ($\kappa=\varnothing$) Eq.~\eqref{eq:p3_s8_3_objective} reduces to text-to-video generation, and with a single frame ($T=1$) it recovers the image generator of Eq.~\eqref{eq:t2i_generator} exactly, so video spatial generation is precisely the image prior of §\ref{sec:p3_s8_1} unrolled in time and equipped with a richer conditioning interface.

The subsections are organized by the axis of control each introduces into Eq.~\eqref{eq:p3_s8_3_objective}.
Text-to-video generation (§\ref{subsec:p3_s8_3_1}) is the base case without auxiliary control; spatial control (§\ref{subsec:p3_s8_3_2}) adds a per-frame scaffold such as a pose or depth sequence; temporal control (§\ref{subsec:p3_s8_3_3}) adds a motion or camera signal that governs dynamics rather than appearance; §\ref{subsec:p3_s8_3_4} covers the remaining modalities and the frameworks that fuse several at once; and video editing (§\ref{subsec:p3_s8_3_5}) modifies an existing clip, lifting the 2D editing of §\ref{subsec:p3_s8_1_3} into time.
§\ref{subsec:p3_s8_3_6} compares the five and discusses the challenges they share.



\input{part_3/section_8/p3_sec8_3/p3_sec8_3_1.tex}

\input{part_3/section_8/p3_sec8_3/p3_sec8_3_2.tex}

\input{part_3/section_8/p3_sec8_3/p3_sec8_3_3.tex}

\input{part_3/section_8/p3_sec8_3/p3_sec8_3_4.tex}

\input{part_3/section_8/p3_sec8_3/p3_sec8_3_5.tex}






\subsection{Summary and Critical Challenges and Future Directions}
\label{subsec:p3_s8_3_6}
\label{subsec:p3_s8_3_summary}

Table~\ref{tab:p3_s8_3_subfield_comparison} places the five control axes side by side: the control signal each introduces, the operator it instantiates over the shared sampler of Eq.~\eqref{eq:p3_s8_3_objective}, the paradigms it surveys, the methods anchoring each, and the benchmarks on which it reports.

\begin{table}[t]
    \centering
    \fontsize{8}{10}\selectfont
    \setlength{\tabcolsep}{1.0mm}
    \caption{The five control axes of video spatial generation; each operator specializes the spatio-temporal sampler of Eq.~\eqref{eq:p3_s8_3_objective}.}
    \label{tab:p3_s8_3_subfield_comparison}
    \begin{tabular}{L{0.12\textwidth} L{0.22\textwidth} L{0.18\textwidth} L{0.19\textwidth} L{0.21\textwidth}}
        \toprule
        \textbf{Control axis} & \textbf{Control signal / operator} & \textbf{Dominant paradigms} & \textbf{Canonical methods} & \textbf{Benchmarks / metrics} \\
        \midrule
        §\ref{subsec:p3_s8_3_1} Text-to-video & Text only, $\kappa=\varnothing$; $G_\theta^{\text{vid}}:\mathcal{E}\!\to\!\mathcal{Z}^{\text{vid}}$ (Eq.~\eqref{eq:p3_s8_3_1_objective}) & GAN; autoregressive token; diffusion / diffusion transformer & CogVideo, Sora, Wan, HunyuanVideo, Cosmos & UCF-101, MSR-VTT, VBench, EvalCrafter; FVD, CLIPSIM \\
        §\ref{subsec:p3_s8_3_2} Spatial control & Per-frame scaffold $\kappa_{\text{s}}(\mathbf{c}_{\text{s}}^{1:T})$ (Eq.~\eqref{eq:p3_s8_3_2_objective}); residual injection (Eq.~\eqref{eq:p3_s8_3_2_inject}) & Image, pose, depth, landmark, sketch, bounding box & Animate Anyone, MagicAnimate, ControlVideo, EchoMimic, MagicDrive & nuScenes, pose-annotated video; keypoint / depth / identity / detection adherence \\
        §\ref{subsec:p3_s8_3_3} Temporal control & Motion signal $\kappa_{\text{t}}(\mathbf{m}^{1:T})$ (Eq.~\eqref{eq:p3_s8_3_3_objective}); temporal-branch injection (Eq.~\eqref{eq:p3_s8_3_3_inject}) & Flow, trajectory, camera, motion; latent warping or encoded conditioning & Motion-I2V, Tora, DragNUWA, CameraCtrl, MotionDirector & RealEstate10K, DL3DV-10K, MPI-Sintel; EPE, path / pose error, transfer fidelity \\
        §\ref{subsec:p3_s8_3_4} Other control & Heterogeneous $\kappa_{\text{o}}$ fused by $\mathcal{F}_\psi$ (Eqs.~\eqref{eq:p3_s8_3_4_objective}, \eqref{eq:p3_s8_3_4_inject}) & Style, point, BEV, text, audio; summed adapters, concatenation, unified attention & StyleCrafter, DriveDreamer-2, EMO, VideoComposer, FullDiT & nuScenes, HDTF, VFHQ; style / point / BEV / OCR / lip-sync adherence \\
        §\ref{subsec:p3_s8_3_5} Video editing & Existing clip $+$ edit intent $\to$ edited clip; $E_\theta^{\text{vid}}$ (Eq.~\eqref{eq:p3_s8_3_5_objective}) & Training-free, one-shot tuning, trained instruction editors, attention reuse, structural guidance & Tune-A-Video, TokenFlow, InsV2V, UNIC, EditVerse & DAVIS, LOVEU-TGVE, FiVE, IVEBench; edit-vs-preservation tuple, VLM / MLLM judges \\
        \bottomrule
    \end{tabular}
\end{table}

The axes partition not by output domain, since all five decode a clip from the same $\mathcal{Z}^{\text{vid}}$, but by the control fed to one shared sampler: text-to-video is the $\kappa=\varnothing$ limit, spatial, temporal, and other control progressively enrich $\kappa$, and editing alone reverses the direction of the map, taking an existing clip in place of pure noise.
The injection mechanism converges across axes: the spatial residual branch (Eq.~\eqref{eq:p3_s8_3_2_inject}), the temporal branch (Eq.~\eqref{eq:p3_s8_3_3_inject}), and the multi-condition interface (Eq.~\eqref{eq:p3_s8_3_4_inject}) are one idea, a trainable control pathway added to a pretrained backbone, generalized from one condition to many, and the most recent systems (VideoComposer, FullDiT, UNIC, EditVerse) go further by folding every axis, editing included, into a single token sequence processed by shared attention.
Evaluation has not converged: no suite spans the five axes, each reports generic quality together with a channel-specific adherence score, and the newest editing benchmarks (FiVE, IVEBench) add vision-language and multimodal-LLM judges to improve human alignment.

Video generation inherits every difficulty of the 2D and 3D priors of §\ref{sec:p3_s8_1} and §\ref{sec:p3_s8_2} and adds the time axis, which turns several already-hard problems into open ones; the four below cut across all five axes.

\noindent\textbf{Temporal coherence and long-horizon drift.} The defining difficulty of video over image generation is keeping content stable across frames. The signature artifact is \emph{flicker}, the high-frequency appearance instability between adjacent frames that plagues per-frame and early GAN-based methods (§\ref{subsec:p3_s8_3_1}); over longer spans it gives way to \emph{identity drift}, the gradual mutation of a subject's appearance, color, or fine texture across a clip; and in autoregressive roll-outs to \emph{error accumulation}, the compounding of small per-step mistakes that caps usable clip length. Latent video autoencoders with temporal compression~\cite{blattmann2023stable,blattmann2023align} and dedicated spatio-temporal attention suppress flicker, while autoregressive--diffusion hybrids such as CausVid~\cite{yin2025slow} and MAGI-1~\cite{ai2025magi1autoregressivevideogeneration} combine sequential consistency with diffusion fidelity to push toward minute-scale generation; coherence over arbitrarily long horizons remains unsolved and is the binding constraint on the world-model applications that motivate the section.

\noindent\textbf{Controllability versus fidelity under multi-condition fusion.} As §\ref{subsec:p3_s8_3_2}--§\ref{subsec:p3_s8_3_4} stack heterogeneous conditions (spatial scaffolds, camera paths, audio, style) onto Eq.~\eqref{eq:p3_s8_3_objective}, the conditions begin to conflict. Over-strong conditioning yields rigidity and texture collapse; weak conditioning ignores the control; and naive fusion produces \emph{control-signal dilution} (summed adapters wash one another out), \emph{condition leakage} across frames, and \emph{collapse to the dominant condition} when one modality (\eg, audio in talking-head generation) overwhelms the rest. Compositional and universal frameworks (VideoComposer~\cite{wang2023videocomposercompositionalvideosynthesis}, OmniHuman~\cite{lin2025omnihuman1rethinkingscalinguponestage}, and the unified-attention design of FullDiT~\cite{ju2025fullditmultitaskvideogenerative}) replace per-axis injection with joint token fusion and task-aware positional embeddings, but a principled account of how to weight and arbitrate competing conditions is still missing.

\noindent\textbf{Physical and 3D plausibility.} Even a temporally coherent, well-controlled clip may depict an impossible world, which is the bridge back to §\ref{sec:p3_s8_2}. Generated motion frequently violates object permanence (objects teleport, morph, or vanish), contact and rigid-body dynamics (rigid objects warp non-rigidly; feet slide; occlusions resolve implausibly), and \emph{multi-view consistency} under camera motion, where the same scene is rendered inconsistently from successive viewpoints. Systems presented as steps toward ``world simulators'', Sora~\cite{sora2024} among them, still exhibit these failures on closer inspection. Camera-control methods~\cite{he2024cameractrl,yu2024viewcrafter} and explicitly 3D-aware video diffusion (Diffusion-as-Shader~\cite{gu2025diffusionshader3dawarevideo} and GS-DiT~\cite{bian2025gsditadvancingvideogeneration}) narrow the gap by conditioning generation on geometry, which suggests that the more promising route to physical plausibility runs through the volumetric representations of §\ref{sec:p3_s8_2} rather than through scale alone.

\noindent\textbf{Compute, latency, and evaluation.} Underlying all three is cost: a spatio-temporal latent is an order of magnitude larger than an image latent, so long, high-resolution clips strain memory and make iterative sampling slow, which in turn limits the context length available for coherence and the conditioning budget available for control. Consistency distillation and few-step samplers, together with sparse and causal attention, are the active responses. Evaluation compounds the problem: distributional scores such as FVD and per-frame FID are blind to temporal and control adherence, and automatic metrics correlate weakly with human judgment, so reported progress is hard to compare across methods; benchmarks and human-aligned protocols that jointly score fidelity, temporal coherence, and per-axis control adherence are a central near-term need.

Coherence and physical plausibility both demand longer temporal context, which inflates compute, and every added condition enlarges the evaluation surface, so progress on any one of the four is bounded by the others. Together they mark the last step of the progression traced across §\ref{sec:generation}, from the static 2D prior of §\ref{sec:p3_s8_1} through the static volumetric prior of §\ref{sec:p3_s8_2} to dynamic, navigable models that consume spatial structure over time.

% \newpage


\section{Motion Generation}
\label{sec:p3_s8_4}

\FramedInitial{M}{otion} generation is the temporal, embodied output channel of spatial generation: where §\ref{sec:p3_s8_1}–§\ref{sec:p3_s8_3} synthesize static images, volumetric assets, and controllable video, motion generation produces the \emph{articulated dynamics} that an avatar, character, or robot can actually enact.
A generated motion is a time series of body poses that must be believable in its dynamics, physically plausible (balance, ground contact, non-penetration), and faithful to its driving signal, whether a text description, audio, a scene, or a key frame.
The conditional-sampler template that turns a prompt into pixels (Eq.~\eqref{eq:t2i_generator}) is here instantiated over a sequence of poses in a structured kinematic space, so geometric priors meet physical constraints: the output must respect the biomechanics of an articulated body as well as the semantics of the condition.

The subsections are indexed by the moving agent and then by one cross-cutting operation: human motion (§\ref{subsec:p3_s8_4_1}), animal motion with its cross-species morphology (§\ref{subsec:p3_s8_4_2}), vehicle motion planning, where the body is rigid and the difficulty moves into the multi-agent traffic scene (§\ref{subsec:p3_s8_4_3}), and motion editing (§\ref{subsec:p3_s8_4_4}).
% , with a comparison in §\ref{subsec:p3_s8_4_summary}.




\subsection{Human Motion Generation}
\label{subsec:p3_s8_4_1}
Human motion generation produces natural, varied sequences of human body poses from intentions, environmental stimuli, or explicit conditions, and supports animation, virtual characters in film and games, AR/VR, and human--robot interaction.

\subsubsection{Task Definition}
\label{subsubsec:p3_s8_4_1_task}
Let a motion be a temporal sequence of body poses $\mathbf{m}^{1:T}=(\mathbf{m}^1,\dots,\mathbf{m}^T)$, where each frame $\mathbf{m}^t\in\mathcal{P}$ is a pose in a chosen representation space $\mathcal{P}$ and $\mathcal{M}=\mathcal{P}^T$ is the resulting motion space. Writing $\phi$ for the text encoder of Eq.~\eqref{eq:t2i_generator} and $\kappa:\mathcal{C}\!\to\!\mathcal{C}'$ for an encoder of a driving signal $\mathbf{c}\in\mathcal{C}$ (an action label, a text description, a music or speech track, or a scene), human motion generation learns a conditional sampler
\begin{equation}
\label{eq:p3_s8_4_1_objective}
G_\theta^{\text{mot}}:\mathcal{C}'\!\to\!\mathcal{M},\qquad
\hat{\mathbf{m}}^{1:T}\sim p_\theta\!\left(\mathbf{m}^{1:T}\mid\kappa(\mathbf{c})\right).
\end{equation}
This is the conditional-sampler template of Eq.~\eqref{eq:t2i_generator} instantiated over a temporal pose sequence rather than an image latent: with no driving signal ($\kappa=\varnothing$) it reduces to unconditional motion synthesis, and the choice of $\mathbf{c}$ indexes the five conditioning settings surveyed in §\ref{subsubsec:p3_s8_4_1_method}.

Figure~\ref{fig:p3_s8_4_1_repr} illustrates the pose representations that instantiate $\mathcal{P}$. Broadly, representations fall into keypoint-based and rotation-based categories, and the two are convertible: joint rotations map to spatial keypoints via forward kinematics (FK), and keypoints recover joint rotations via inverse kinematics (IK).

\begin{figure}
    \centering
    \includegraphics[width=\linewidth]{figures/part3/section8/human_motion_repr.png}
    \includegraphics[width=\linewidth]{figures/part3/section8/human_motion_repr_smpl.png}
    \caption{Human motion representations. Top: keypoint-based, as (a) joint positions, (b) joints grouped by kinematic chain, and (c) dense surface markers; and rotation-based, as (d) per-joint rotations in $\mathrm{SO}(3)$ on a kinematic tree. Bottom: the SMPL body model, from (a) the template mesh and skinning weights through (b) shape and (c) pose blend shapes to (d) the skinned, posed mesh.}
    \label{fig:p3_s8_4_1_repr}
\end{figure}

\paragraph{Keypoint-based Representation.} The body is modeled as a collection of keypoints, each a distinct anatomical landmark (joints or salient surface locations) given by its 2D or 3D coordinates in a chosen frame (pixel space or world space), and motion is encoded as a temporal series of such configurations.
Keypoint representations are directly available from many motion-capture systems and are highly interpretable, but many downstream uses (character animation, robotic control) require converting keypoints to joint rotations through the IK problem.
Some works instead augment or replace sparse skeletal keypoints with dense body markers distributed across the body surface~\cite{WuWZZHYT22, TaheriCBT22}, which carry richer information about body shape, limb twisting, and fine geometric detail.

\paragraph{Rotation-based Representation.} A pose is represented by joint angles, the rotations of body segments relative to their parent in a hierarchical kinematic tree. Most studies model these in three dimensions, treating each non-root joint as an element of $\mathrm{SO}(3)$, parameterized as Euler angles, axis--angle vectors, or unit quaternions.

\paragraph{Skinned Multi-Person Linear Model.} Beyond pure joint rotation, statistical mesh models jointly capture pose and body shape, including non-rigid deformations associated with articulation.
The Skinned Multi-Person Linear (SMPL) model~\cite{loper2015smpl} defines a template mesh with fixed topology (6{,}890 vertices, 23 joint rotations plus a root) and uses linear blend skinning augmented by learned blend shapes: identity (shape) blend shapes for inter-subject variation and pose blend shapes for rotation-dependent deformations such as muscle bulging and twisting.



\subsubsection{Methodology}
\label{subsubsec:p3_s8_4_1_method}
Methods for human motion generation are organized along two axes. The first is the \emph{learning paradigm}: regression-based techniques learn a deterministic mapping from conditions to a single target motion, whereas generative approaches model the full conditional distribution of motion using GANs, VAEs, normalizing flows, or diffusion models (the generative paradigms themselves are detailed in §\ref{subsec:p3_s8_1_1}).
Formally, a regression model realizes Eq.~\eqref{eq:p3_s8_4_1_objective} as a point estimate,
\begin{equation}
\label{eq:p3_s8_4_1_regression}
\hat{\mathbf{m}}^{1:T}=f_\theta\!\left(\kappa(\mathbf{c})\right),\qquad
\min_\theta\ \big\|\hat{\mathbf{m}}^{1:T}-\mathbf{m}^{1:T}\big\|,
\end{equation}
which is simple and stable but collapses the inherently one-to-many condition--motion relationship onto an averaged trajectory. Generative models instead sample $\hat{\mathbf{m}}^{1:T}\!\sim\!p_\theta(\mathbf{m}^{1:T}\mid\kappa(\mathbf{c}))$; the now-dominant diffusion instantiation trains a denoiser $G_\theta$ to recover the clean motion at every noise level $n$,
\begin{equation}
\label{eq:p3_s8_4_1_diffusion}
\mathcal{L}=\mathbb{E}_{\mathbf{m}^{1:T},\,n}\,\big\|\,\mathbf{m}^{1:T}-G_\theta(\mathbf{m}^{1:T}_n,\,n,\,\kappa(\mathbf{c}))\,\big\|^2,
\end{equation}
which preserves diversity and supports flexible conditioning, at higher sampling cost.
The second axis is the \emph{conditioning modality} $\mathbf{c}$, which indexes the five settings below; each setting instantiates the same Eq.~\eqref{eq:p3_s8_4_1_objective} and has migrated, broadly, from regression and adversarial models toward diffusion. Figure~\ref{fig:p3_s8_4_1_methods} previews this two-axis organization.



\begin{figure}
    \centering
    \includegraphics[width=\linewidth]{figures/part3/section8/two_axis_human_motion_generation.pdf}
    \caption{Two-axis organization of human motion generation: the conditioning modality (left), encoded by $\kappa$, and the learning paradigm (right), either a regressor that returns one estimate (Eq.~\eqref{eq:p3_s8_4_1_regression}) or a generative model that returns alternative samples, shown with diffusion (Eq.~\eqref{eq:p3_s8_4_1_diffusion}). Branches are alternatives, not consecutive stages.}
    \label{fig:p3_s8_4_1_methods}
\end{figure}

\paragraph{Action-to-Motion.} Given a discrete action label (\eg, ``walk'', ``throw''), these methods generate a corresponding sequence. Early work used GANs with graph convolution and self-attention (SA-GAN~\cite{YuZLYC20}, Kinetic-GAN~\cite{DegardinNLBY022}) for realism and diversity, while VAE/CVAE frameworks (Action2Motion~\cite{GuoZWZSDG020}, ACTOR~\cite{PetrovichBV21}) produced probabilistic, controllable motions. More recent models adopt contrastive latent learning (ODMO~\cite{LuZLR22}), discrete autoregressive sequences (PoseGPT~\cite{LucasBWR22}), and diffusion (MDM~\cite{TevetRGSCB23}, the latent-diffusion hybrid MLD~\cite{ChenJLHFCY23}). As most methods excel only at single actions, MultiAct~\cite{LeeML23} recurrently composes multiple action motions into long sequences via a canonicalization scheme.

\paragraph{Text-to-Motion.} Conditioning on free-form language exploits richer semantic cues than fixed labels allow. Early methods used GANs (Text2Action~\cite{AhnHCYO18}) or joint text--motion embeddings (JL2P~\cite{AhujaM19}, two-stream encoder--decoders~\cite{GhoshCOTS21}); later ones added variational modeling and attention (TEMOS~\cite{PetrovichBV22}, TEACH~\cite{AthanasiouPBV22}).
For zero-shot generalization, MotionCLIP~\cite{TevetGHBC22} and AvatarCLIP~\cite{HongZPCYL22} align the motion latent space with CLIP.
The current frontier combines VQ-VAE tokenization~\cite{GuoZWC22}, autoregressive transformers (T2M-GPT~\cite{zhang2023generating}), and diffusion (FLAME~\cite{KimKC23}), with multimodal conditioning extensions such as MoFusion~\cite{DabralMGT23} fusing text and audio. More recent masked-generative models (MoMask~\cite{GuoMJW024}) push quality and inference speed further by predicting quantized motion tokens with a bidirectional, iteratively-refined masking schedule rather than left-to-right.

\paragraph{Music-to-Dance.} The goal is coherent, expressive dance aligned to a music track. Early approaches framed this as supervised sequence-to-sequence translation via LSTM autoencoders~\cite{TangJM18}, transformers (AI Choreographer~\cite{LiYRK21}), and multi-person frameworks (GroupDancer~\cite{Wang0WXCMCW22}, GDancer~\cite{LePDTT023}).
Because music--dance is inherently one-to-many, generative models model the diverse distribution: MNET~\cite{0001OKT022} with style conditioning, Transflower~\cite{PerezHBHOA21}, Bailando~\cite{SiyaoYGLW0L022} (VQ-VAE + GPT), and EDGE~\cite{TsengCL23} (diffusion). Graph-based methods (ChoreoMaster~\cite{ChenTLZGZ021}, PC-Dance~\cite{GaoPZSZ22}, ChoreoGraph~\cite{AuCJG22}) use motion graphs and shared music--motion embeddings for rhythmic coherence, while curriculum learning, dynamics-prior denoising, and hierarchical pose/motif/choreography modeling target long-term, genre-consistent full-song dances~\cite{HuangHWSZJ21, SunWHLJH22, AristidouYACSC23}.

\paragraph{Speech-to-Gesture.} Co-speech gesture synthesis generates upper-body or full gesture sequences from speech (audio, sometimes text).
Early adversarial methods modeled person-specific gesture patterns~\cite{GinosarBKCOM19}; subsequent work disentangled shared versus motion-specific variation~\cite{KucherenkoHHKK19} and decomposed body, hands, and face~\cite{HabibieXMLSPET21}.
Recent methods employ diffusion~\cite{ZhuLLQLY23, AlexandersonNBH23, AoZL23}, VQ-VAE quantization~\cite{YiLLCWBTB23}, and style or speaker embeddings~\cite{AhujaLNM20} to handle diversity, identity, rhythm, and expression more robustly.

\paragraph{Scene-to-Motion.} Here the condition is a 3D scene, and most methods adopt a multi-stage pipeline that predicts a goal or interaction anchor, plans a trajectory, then generates motion along it. Early works predict a goal from 2D/RGB scene inputs and synthesize motion through trajectory-guided VAEs~\cite{CaoGMCVM20, WangYDL21}, sometimes with scene-aware discriminators. SAMP~\cite{HassanCVSYZB21} models 3D paths and human--object interactions; later methods refine interaction with per-joint contact maps~\cite{D00010HS22} or path-search ($A^\ast$) for trajectory diversity~\cite{0003RLYLD22}, and GOAL~\cite{TaheriCBT22} and SAGA~\cite{WuWZZHYT22} target grasping via two-step goal-to-motion VAEs.
Some assume given goals and stitch sub-goal motions hierarchically for long-term synthesis~\cite{WangXXL021}, while refinement-based frameworks such as CIRCLE~\cite{AraujoLVAWGCL23} initialize by interpolation and refine with scene-geometry-aware modules.


\subsubsection{Benchmarks}
\label{subsubsec:p3_s8_4_1_benchmarks}
Human motion datasets group naturally by conditioning modality; Table~\ref{tab:p3_s8_4_1_datasets} gives their representation, scale, and condition.
Text- and action-conditioned corpora dominate, led by the language-annotated KIT Motion-Language~\cite{PlappertMA16} and HumanML3D~\cite{GuoZZ0JL022}, BABEL~\cite{PunnakkalCAQB21}'s sequence- and frame-level labels over the AMASS~\cite{mahmood2019amass} aggregate, and the action taxonomies of UESTC~\cite{JiXYSSZ18}, NTU-RGB+D~120~\cite{LiuSPWDK20}, and HumanAct12~\cite{GuoZWZSDG020}.
Audio-conditioned resources split into music-to-dance sets (AIST++~\cite{LiYRK21}, AIOZ-GDANCE~\cite{LePDTT023}, and several genre-specific collections) and co-speech gesture sets (Trinity~\cite{FerstlM18}, BEAT~\cite{LiuZIPLZBZ22}, PATS~\cite{AhujaLNM20}), many built by fitting parametric body models to in-the-wild video.
Scene-conditioned datasets (PROX~\cite{HassanCTB19}, GRAB~\cite{TaheriGBT20}, HUMANISE~\cite{WangCLZLH22}, CIRCLE~\cite{AraujoLVAWGCL23}) capture human--scene and human--object interaction with synchronized geometry, from Vicon and IMU rigs to synthetic engines such as GTA-IM~\cite{CaoGMCVM20}. Across all three groups, the cost of high-quality capture forces a trade-off between scale and annotation fidelity.

\begin{table}[htbp]
\centering
\fontsize{8}{10}\selectfont
\setlength{\tabcolsep}{1.2mm}
\caption{Benchmarks for human motion generation, grouped by condition (none, text / action, music, speech, scene/object).}
\label{tab:p3_s8_4_1_datasets}
\begin{tabular}{L{0.18\textwidth} C{0.115\textwidth} C{0.07\textwidth} C{0.08\textwidth} C{0.06\textwidth} C{0.06\textwidth} C{0.13\textwidth} L{0.17\textwidth}}
\toprule
\textbf{Dataset} & \textbf{Representation} & \textbf{Subjects} & \textbf{Sequences} & \textbf{Frames} & \textbf{Length} & \textbf{Condition} & \textbf{Remarks} \\
\midrule
Human3.6M \cite{IonescuPOS14} & Keypoint & 11 & - & 3.6M & 5h & - & 15 actions \\
AMASS \cite{mahmood2019amass} & Rotation & 344 & 11,265 & - & 40h & - & Unifies 15 marker-based MoCap datasets \\
HuMMan \cite{CaiRZLYWFGYPHZL22} & Rotation & 1000 & 400K & 60M & - & - & 500 actions \\
KIT Motion Language \cite{PlappertMA16} & Keypoint & 111 & 3,911 & - & 10.3h & Text & 6.3K text descriptions \\
UESTC \cite{JiXYSSZ18} & Keypoint & 118 & 25.6K & - & 83h & Action & 40 action classes \\
NTU-RGB+D 120 \cite{LiuSPWDK20} & Keypoint & 106 & 114.5K & - & 74h & Action & 120 action classes \\
HumanAct12 \cite{GuoZWZSDG020} & Keypoint & 12 & 1,191 & 90K & 6h & Action & 12 action classes \\
BABEL \cite{PunnakkalCAQB21} & Rotation & 344 & - & - & 43.3h & Text & 260 action classes \\
HumanML3D \cite{GuoZZ0JL022} & Keypoint + Rotation & 344 & 14.6K & - & 28.5h & Text & 44.9K text descriptions \\
DanceWithMelody \cite{TangJM18} & Keypoint & - & 61 & 907K & 1.6h & Music & 4 genres \\
Lee \cite{LeeY0WLYK19} & Keypoint & - & 361K & - & 71h & Music & 3 genres \\
Huang \cite{HuangHWSZJ21} & Keypoint & - & 790 & - & 12h & Music & 3 genres\\
AIST++ \cite{LiYRK21} & Rotation & 30 & 1,408 & 10.1M & 5.2h & Music & 10 genres \\
PMSD \cite{PerezHBHOA21} & Keypoint & 8 & - & - & 3.8h & Music & 4 genres \\
ShaderMotion \cite{PerezHBHOA21} & Keypoint & 8 & - & - & 10.2h & Music & 2 genres \\
Chen \cite{ChenTLZGZ021} & Rotation & - & - & 160K & 9.9h & Music & 9 genres \\
PhantomDance \cite{LiZZS22} & Rotation & - & 260 & 795K & 3.7h & Music & 13 genres \\
MMD-ARC \cite{GaoPZSZ22} & Rotation & - & 213 & - & 11.3h & Music & - \\
MDC \cite{Wang0WXCMCW22} & Rotation & - & 798 & - & 3.5h & Music & 2 genres \\
Aristidou \cite{AristidouYACSC23} & Rotation & 32 & - & - & 2.4h & Music & 3 genres \\
AIOZ-GDANCE \cite{LePDTT023} & Rotation & >4000 & - & - & 16.7h & Music & 7 styles, 16 genres \\
Trinity \cite{FerstlM18} & Keypoint & 1 & 23 & - & 4.1h & Speech & Casual talks\\
TED-Gesture \cite{YoonKJLKL19} & Keypoint & - & 1,295 & - & 52.7h & Text & TED talks \\
Speech2Gesture \cite{GinosarBKCOM19} & Keypoint & 10 & - & - & 144h & Speech & TV shows, lectures \\
TED-Gesture++ \cite{YoonCLJLKL20} & Keypoint & - & 1,766 & - & 97.0h & Speech, Text & Extension of \cite{YoonKJLKL19} \\
PATS \cite{AhujaLNM20} & Keypoint & 25 & - & - & 251h & Speech, Text & Extension of \cite{GinosarBKCOM19} \\
Speech2Gesture-3D \cite{HabibieXMLSPET21} & Keypoint & 6 & - & - & 33h & Speech & Videos from \cite{GinosarBKCOM19}\\
BEAT \cite{LiuZIPLZBZ22} & Rotation & 30 & 2,580 & 30M & 76h & Speech, Text, Emotion & 8 emotions, 4 languages \\
Chinese Gesture \cite{AoGLCL22} & Rotation & 5 & - & - & 4h & Speech, Text & Chinese \\
ZEGGS \cite{GhorbaniFHTC23} & Rotation & 1 & 67 & - & 2.3h & Speech, Style & 19 styles \\
SHOW \cite{YiLLCWBTB23} & Rotation & - & - & - & 27h & Speech & Videos from \cite{GinosarBKCOM19} \\
WBHM \cite{ManderyTDVA15} & Rotation & 43 & 3,704 & 691K & 7.68h & Object & 41 objects \\
PiGraph \cite{Hanrahan16a} & Keypoint & 5 & 63 & 0.1M & 2h & Scene, Object & 30 scenes, 19 objects \\
PROX \cite{HassanCTB19} & Rotation & 20 & 60 & 0.1M & 1h & Scene, Object & 12 indoor scenes\\
i3DB \cite{MonszpartGCYM19} & Keypoint & 1 & - & - & - & Scene, Object & 15 scenes\\
GTA-IM \cite{CaoGMCVM20} & Keypoint & 50 & 119 & 1M & - & Scene & Synthetic, 10 indoor scenes \\
GRAB \cite{TaheriGBT20} & Rotation & 10 & 1,334 & 1.6M & - & Object & 51 objects \\
HPS \cite{GuzovMSP21} & Rotation & 7 & - & 300K & - & Scene & 8 large scenes\\
SAMP \cite{HassanCVSYZB21} & Rotation & 1 & - & 185K & 0.83h & Scene, Object & 7 objects\\
COUCH \cite{ZhangBSGP22} & Rotation & 6 & >500 & - & 3h & Scene, Chairs & 3 chairs\\
HUMANISE \cite{WangCLZLH22} & Rotation & - & 19.6K & 1.2M & - & Scene, Object, Text & 643 scenes \\
CIRCLE \cite{AraujoLVAWGCL23} & Rotation & 5 & >7K & 4.3M & 10h & Scene & 9 scenes \\
\bottomrule
\end{tabular}
\end{table}

\paragraph{Evaluation Metrics.} Evaluating synthesized motion is hard because the condition--motion map is one-to-many, judgments are partly subjective, and quality depends on high-level conditional cues, so the literature reports an ensemble of metrics along three axes.
\textbf{Fidelity} captures naturalness, smoothness, and plausibility: distance-based measures such as NPSS~\cite{GopalakrishnanM19} and NDMS~\cite{TankeZG21}, keypoint accuracy (PCK~\cite{YangR13}), motion-space statistics (AVE, Hellinger distance over speed histograms~\cite{yang2023qpgesture}, Fréchet Distance over pose rotations~\cite{HassanCVSYZB21, 0003RLYLD22}, PSKL~\cite{RuizGM19}), feature-space distributional distances led by FID (with IS~\cite{AhujaLNM20}, MMD~\cite{YuZLYC20}, and related scores), and physical-plausibility proxies for foot sliding and foot--ground contact (FCS~\cite{TsengCL23}, non-collision ratios~\cite{WuWZZHYT22}, contact accuracy~\cite{ZhangT22}), the last still lacking a standardized definition.
\textbf{Diversity} guards against the long-sequence ``freezing'' failure via intra-motion variation and freezing rate~\cite{SunWHLJH22}, inter-motion manifold coverage, and multimodality metrics (n-APD, ADE/FDE and their multimodal variants~\cite{YuanK20}) computed by sampling many motions per condition.
\textbf{Condition consistency} measures alignment to the driving signal: recognition accuracy and R-Precision for text/action, feature-level text--motion distance (Multimodal Distance~\cite{GuoZZ0JL022}, mCLIP and MID~\cite{KimKC23}), audio--motion beat metrics (beat coverage, hit rate~\cite{LeeY0WLYK19}, beat alignment~\cite{LiYRK21}, SRGR~\cite{LiuZIPLZBZ22}, Semantic Score~\cite{AoZL23}), and, for scene conditioning, non-collision and human--scene contact scores~\cite{WuWZZHYT22, ZhangBSGP22}. No single metric is decisive, so methods are compared on a task-appropriate composition of the three.


\subsubsection{Challenges and Trends}
\label{subsubsec:p3_s8_4_1_challenges}
\textbf{Data scarcity} is the most basic constraint: high-quality motion capture is expensive and slow, so dataset scale trades against annotation fidelity, and heterogeneous pose representations and conditioning signals fragment the data that exist; fusing heterogeneous sources through weak supervision and multimodal foundation models is the main response.
\textbf{Semantic alignment} is the second. Motion encodes culture, communication, and context, not only mechanics, and a model that fits joint trajectories while missing their relation to language, music, or environment produces motion that is kinematically clean but perceptually empty; semantic priors from language and foundation models help, but how to measure the result is unresolved.
That is the \textbf{evaluation} problem: every objective metric above has blind spots, none replaces a user study, and the field still needs metrics that are both interpretable and better correlated with human perception.
\textbf{Controllability} matters for practical use as much as in image generation, and joint masking, attribute prompts, and style conditioning are moving toward editing tools with textual, spatial, and temporal handles (§\ref{subsec:p3_s8_4_4}).
\textbf{Interactivity} remains underexplored: most methods animate one person in a static scene, and socially coupled groups (conversation, group dance) and motion grounded in dynamic, actionable scenes need both new corpora and new models.





\subsection{Animal Motion Generation}
\label{subsec:p3_s8_4_2}
Most motion-generation work targets the human body and ignores the diversity of animal behavior, although synthesizing animal motion matters for wildlife conservation, ecology, and biomechanics.
Relative to §\ref{subsec:p3_s8_4_1}, animal motion is harder on two fronts that the human setting never confronts: species vary widely in skeletal topology and morphology, so a single fixed-skeleton representation no longer suffices, and the same textual description (``walk stealthily'') maps to radically different gaits across species.
These two pressures, \ie, data scarcity and morphological heterogeneity, shape every method below, which is why animal motion generation borrows the architectures, datasets, and training recipes of human motion generation while adding species-conditioning machinery on top.



\subsubsection{Task Definition}
\label{subsubsec:p3_s8_4_2_task}
Animal motion generation extends the human sampler of Eq.~\eqref{eq:p3_s8_4_1_objective} with an explicit species (morphology) descriptor. Let $s\in\mathcal{S}$ index a species and let $\mathcal{P}_s$ be its species-dependent pose space (its skeleton topology and joint count depend on $s$), so that $\mathcal{M}_s=\mathcal{P}_s^{T}$ is the corresponding motion space. With $\kappa$ encoding a driving signal $\mathbf{c}\in\mathcal{C}$ (an action class, a text description, an anchor pose, or past motion), the task is to learn
\begin{equation}
\label{eq:p3_s8_4_2_objective}
G_\theta^{\text{amot}}:\mathcal{C}'\!\times\!\mathcal{S}\!\to\!\mathcal{M}_s,\qquad
\hat{\mathbf{m}}^{1:T}\sim p_\theta\!\left(\mathbf{m}^{1:T}\mid\kappa(\mathbf{c}),\,s\right).
\end{equation}
Restricting $\mathcal{S}$ to the single human morphology recovers Eq.~\eqref{eq:p3_s8_4_1_objective} exactly; the open difficulty is that $s$ ranges over a heterogeneous, long-tailed set of skeletons for which paired text--motion data are scarce, so the species condition must do double duty as both a generation control and a regularizer that shares statistical strength across morphologies.


\subsubsection{Methodology}
\label{subsubsec:p3_s8_4_2_method}
Almost all animal motion generation is text-driven, and the methods share a common response to the data-scarcity and morphology problems of Eq.~\eqref{eq:p3_s8_4_2_objective}: transfer representations learned from abundant human motion, then specialize them to a species through \emph{feature modulation}. Concretely, given a backbone feature $\mathbf{h}$ at some layer, species-aware modulation applies a FiLM-style affine transform conditioned on $s$,
\begin{equation}
\label{eq:p3_s8_4_2_modulation}
\tilde{\mathbf{h}}=\gamma(s)\odot\mathbf{h}+\beta(s),
\end{equation}
where $\gamma(s),\beta(s)$ are predicted from the species descriptor, allowing one shared network to adapt its internal representation to diverse anatomical forms. The methods differ in how they obtain the shared representation and how strictly they enforce cross-species plausibility, as Figure~\ref{fig:p3_s8_4_2_methods} summarizes.

\textbf{OmniMotionGPT}~\cite{YangZSWXHBQ024} tackles limited animal data by jointly training motion autoencoders on human and animal corpora and aligning their latent encodings to text via CLIP embeddings; it relies most directly on transfer, and at inference it still depends on human-motion priors, which prevents truly text-only animal generation.
\textbf{Motion Avatar}~\cite{0006W0CZH0F0024} uses SinMDM~\cite{RaabLTABC24} for data augmentation in its ZooGen pipeline (the route by which the Zoo-300K dataset is built) and extends the masked-motion modeling of MoMask~\cite{GuoMJW024} to animals, adding an LLM planner that orchestrates joint motion and 3D-avatar mesh generation; it makes little use of species-specific structure, which is what Eq.~\eqref{eq:p3_s8_4_2_modulation} addresses.
\textbf{AniMo}~\cite{WangRZW25} makes that remedy central, injecting species-aware modulation into a joint-aware spatiotemporal tokenizer and learning, through masked modeling, a mapping from text plus species attributes to motion tokens; its AniMo4D dataset (78{,}149 sequences over 114 species) yields markedly better morphological and behavioral fidelity than transfer-only baselines.
\textbf{X-MoGen}~\cite{abs-2508-05162} pushes toward a single model spanning humans and animals: it embeds motion into a shared space with graph-based VAEs and latent autoencoders, generates via masked modeling, and adds a morphological-consistency module that enforces skeletal plausibility across species, supported by the UniMo4D corpus (${\sim}115$ species, ${\sim}119$K sequences). OmniMotionGPT and Motion Avatar still lean on human priors, whereas AniMo and X-MoGen promote species morphology to a first-class, explicitly regularized conditioning variable, so the four together trace a progression from transfer toward native cross-species modeling.

\begin{figure}[!ht]
    \centering
    \includegraphics[width=\linewidth]{figures/part3/section8/cross_species_animal_motion_generation.pdf}
    \caption{Abstract cross-species motion generation: human and animal corpora support a shared representation, species-dependent modulation (Eq.~\eqref{eq:p3_s8_4_2_modulation}) adapts features, and an action-conditioned sampler produces motion for the selected skeleton. Dog, bird, and horse sequences illustrate alternative target morphologies, not simultaneous outputs. The dashed constraint denotes optional morphology regularization; implementations vary across methods.}
    \label{fig:p3_s8_4_2_methods}
\end{figure}


\subsubsection{Benchmarks}
\label{subsubsec:p3_s8_4_2_benchmarks}
Animal motion resources fall into two groups (Table~\ref{tab:p3_s8_4_2_datasets}).
\textbf{Capture and pose-estimation corpora} supply ground-truth articulated motion but no language: RGBD-Dog~\cite{Kearney0PKC20} and 3DDogs-Wild~\cite{abs-2406-14412} pair RGB-D or RGB video with marker-based mocap for canine pose, and DeformingThings4D~\cite{LiTTZN21} provides synthetic 4D sequences with dense mesh-deformation annotations across humanoid and animal categories.
\textbf{Text-to-motion corpora} drive recent progress: AnimalML3D~\cite{YangZSWXHBQ024} is the first text--animal-motion dataset, and the field then scales sharply through Zoo-300K~\cite{0006W0CZH0F0024} (built by the ZooGen pipeline with SinMDM and LLM annotators), AniMo4D~\cite{WangRZW25} (normalized to a 30-joint cross-species skeleton), and UniMo4D~\cite{abs-2508-05162}, which merges AniMo4D with HumanML3D and KIT-ML. Even the largest text-conditioned sets remain smaller and more morphologically fragmented than their human counterparts of §\ref{subsubsec:p3_s8_4_1_benchmarks}, which is what forces the transfer-and-modulate methodology above.



\begin{table}[!ht]
\centering
\fontsize{8}{10}\selectfont
\setlength{\tabcolsep}{1.0mm}
\caption{Benchmarks for animal motion generation. ``Subjects / species'' counts individual animals for the canine capture sets and species or categories elsewhere; ``Texts / frames'' gives text descriptions for text-conditioned corpora and frames for capture corpora.}
\label{tab:p3_s8_4_2_datasets}
\begin{tabular}{L{0.21\textwidth} L{0.13\textwidth} L{0.13\textwidth} C{0.15\textwidth} C{0.14\textwidth} C{0.14\textwidth}}
\toprule
\textbf{Dataset} & \textbf{Type} & \textbf{Modality} & \textbf{Subjects / species} & \textbf{Sequences} & \textbf{Texts / frames} \\
\midrule
RGBD-Dog \cite{Kearney0PKC20} & Canine & RGB-D, mocap & 5 dogs & 5 motion types & 8,346 frames \\
3DDogs-Wild \cite{abs-2406-14412} & Canine & RGB, mocap & 37 dogs & 143 recordings & - \\
DeformingThings4D \cite{LiTTZN21} & Animal / humanoid & 4D mesh & 31 categories & 1,972 & 122,365 frames \\
AnimalML3D \cite{YangZSWXHBQ024} & Animal / humanoid & 4D & 36 identities & 1,240 & - \\
Zoo-300K \cite{0006W0CZH0F0024} & Animal & 3D mesh & 65 categories & ${\sim}$300K text--motion pairs & - \\
AniMo4D \cite{WangRZW25} & Animal & 4D & 114 species & 78,149 & 185,435 texts \\
UniMo4D \cite{abs-2508-05162} & Animal + human & 4D & ${\sim}$115 species & ${\sim}$119K & - \\
\bottomrule
\end{tabular}
\end{table}

\paragraph{Evaluation Metrics.} Animal motion evaluation reuses the three-axis human protocol of §\ref{subsubsec:p3_s8_4_1_benchmarks}, adapted for species-specific biomechanics.
\emph{Fidelity} is measured by FID against real movement distributions and by MPJPE against ground-truth poses, augmented with physical-plausibility checks for gait-cycle regularity, ground contact, and the absence of foot sliding or limb penetration.
\emph{Condition alignment} carries over R-Precision and Multimodal Distance to score text-/action-to-motion correspondence, and \emph{diversity / multimodality} metrics quantify behavioral variety within and across species. The animal-specific addition is perceptual and biomechanical judgment of naturalness, balance, and anatomical consistency, which matters more for quadrupeds and multi-limbed morphologies than for the single human skeleton.



\subsubsection{Challenges and Trends}
\label{subsubsec:p3_s8_4_2_challenges}
\textbf{Physical realism} is the most basic gap: most methods do not simulate ground-contact forces, quadruped balance, or the fluid dynamics of swimming and flight, so artifacts the human setting has largely learned to suppress reappear, and differently for every body plan.
\textbf{Behavioral diversity} is harder than for humans because one instruction such as ``walk stealthily'' realizes into different gaits, postures, and speeds across species, so the conditional distribution to be captured is far wider while the data to capture it are scarcer.
\textbf{Unseen species and radical morphologies} remain out of reach: species-aware modulation (Eq.~\eqref{eq:p3_s8_4_2_modulation}) and morphological-consistency constraints generalize across the skeletons seen in training, but insects and winged or aquatic creatures violate the shared-skeleton, fixed-topology assumptions on which current methods rest.
AniMo and X-MoGen already move toward treating morphology as a variable rather than a template, and topology-flexible representations are the direction this points.




\subsection{Vehicle Motion Planning}
\label{subsec:p3_s8_4_3}
Where §\ref{subsec:p3_s8_4_1} and §\ref{subsec:p3_s8_4_2} generate the motion of an articulated body, vehicle motion planning generates the trajectory of a (near-)rigid agent moving through a traffic scene.
The body is far simpler, since a vehicle has essentially no internal articulation, but the difficulty moves outward into the environment: a planned trajectory must obey vehicle dynamics, stay on the drivable surface dictated by an HD map, and remain socially compatible with the future motion of many surrounding agents whose intentions are unobserved.
This is the setting where motion generation becomes inseparable from \emph{multi-agent} spatial reasoning and physical feasibility: trajectory forecasting (predicting how others will move) and planning (deciding how the ego vehicle should move) are two readings of the same generative object.

\subsubsection{Task Definition}
\label{subsubsec:p3_s8_4_3_task}
Let a trajectory be a sequence of agent states $\boldsymbol{\tau}^{1:H}=(\mathbf{s}^1,\dots,\mathbf{s}^H)$ over a horizon $H$, where each state $\mathbf{s}^t$ collects position, heading, and velocity.
The driving signal is now a structured scene context $\mathbf{c}=(\mathbf{s}^{-P:0},\,\mathcal{G},\,\mathcal{A})$ comprising the agent's own past states, an HD-map / lane graph $\mathcal{G}$, and the observed histories of neighboring agents $\mathcal{A}$.
With $\kappa$ encoding this context, vehicle motion planning learns
\begin{equation}
\label{eq:p3_s8_4_3_objective}
G_\theta^{\text{veh}}:\mathcal{C}'\!\to\!\mathcal{T},\qquad
\hat{\boldsymbol{\tau}}^{1:H}\sim p_\theta\!\left(\boldsymbol{\tau}^{1:H}\mid\kappa(\mathbf{c})\right)
\quad\text{s.t.}\quad
\mathbf{s}^{t+1}=f_{\text{dyn}}(\mathbf{s}^{t},\mathbf{u}^{t}),
\end{equation}
the conditional-sampler template of Eq.~\eqref{eq:p3_s8_4_1_objective} with the pose sequence replaced by a state trajectory and an added kinematic-feasibility constraint $f_{\text{dyn}}$ (\eg, a bicycle model) that distinguishes this setting from articulated-body generation.
Dropping $\mathcal{A}$ recovers single-agent forecasting; supplying a goal or reward in place of $\mathbf{c}$ turns prediction into planning.


\subsubsection{Methodology}
\label{subsubsec:p3_s8_4_3_method}
Prediction methods vary along two axes: how the scene context is \emph{encoded} and how the multimodal future is \emph{decoded}.
Early encoders rasterize the scene into a bird's-eye-view image, but the field has converged on \emph{vectorized} encoders that represent lanes and agent tracks as polylines: VectorNet~\cite{gao2020vectornet} and LaneGCN~\cite{liang2020lanegcn} build graph representations of the HD map, and transformer encoders such as Wayformer~\cite{nayakanti2023wayformer} and the Motion Transformer (MTR)~\cite{shi2022mtr} attend jointly over map and agents; graph-attention predictors~\cite{ZhangFWH22} model spatial and temporal interaction in the same spirit.
On the decoding side, the defining property is multimodality, since a single context admits many plausible futures, so models predict a set of $K$ modes with probabilities,
\begin{equation}
\label{eq:p3_s8_4_3_multimodal}
p_\theta(\boldsymbol{\tau}\mid\kappa(\mathbf{c}))=\sum_{k=1}^{K}\pi_k\,\mathcal{N}\!\left(\boldsymbol{\tau};\,\mu_k(\kappa(\mathbf{c})),\,\Sigma_k\right),
\end{equation}
rather than the single point estimate of Eq.~\eqref{eq:p3_s8_4_1_regression}, which here would average distinct maneuvers (turning \vs going straight) into an infeasible mean.
The earliest learned predictors modeled social interaction directly in agent-coordinate space, Social-LSTM~\cite{alahi2016social} through pooled recurrent states and Social-GAN~\cite{gupta2018social} through adversarial sampling, and Trajectron++~\cite{salzmann2020trajectron} added dynamic feasibility to a graph-structured latent-variable model.
Goal-based decoders make the modes explicit: DenseTNT~\cite{gu2021densetnt} samples dense candidate endpoints on the map and scores trajectories conditioned on them, exposing the maneuver structure that Eq.~\eqref{eq:p3_s8_4_3_multimodal} captures only implicitly.
Generative decoders replace the fixed mixture with a learned sampler: MotionDiffuser~\cite{jiang2023motiondiffuser} denoises whole multi-agent trajectory sets jointly.

Planning turns the same machinery on the ego vehicle. ChauffeurNet~\cite{bansal2019chauffeurnet} learns the ego policy by imitating expert driving and augments the demonstrations with synthesized perturbations and worst cases, so that the policy stays robust when it drifts off the data manifold.
Figure~\ref{fig:p3_s8_4_3_methods} situates these families along the encode/decode axes.

\begin{figure}[!ht]
    \centering
    \includegraphics[width=\linewidth]{figures/part3/section8/vehicle_motion_planning.pdf}
    \caption{Vehicle motion planning as conditional trajectory generation: a vectorized encoder represents map polylines and agent interactions, and a multimodal decoder generates candidate futures (Eq.~\eqref{eq:p3_s8_4_3_multimodal}). Gray intensity indicates mode weight; blue marks a plan selected using a goal or cost and feasibility constraints, not necessarily the highest-probability mode. Road geometry and trajectories are illustrative.}
    \label{fig:p3_s8_4_3_methods}
\end{figure}

\subsubsection{Benchmarks}
\label{subsubsec:p3_s8_4_3_benchmarks}
Progress is measured on large real-world driving logs that pair agent tracks with HD maps: nuScenes~\cite{caesar2020nuscenes}, Argoverse and Argoverse~2~\cite{Argoverse-ChangLSSBHW0LRH19, wilson2023argoverse}, and the Waymo Open Motion Dataset~\cite{ettinger2021womd}, the last two offering the scale and interaction density that current transformer models require.
Evaluation centers on the multimodal-aware displacement metrics that match Eq.~\eqref{eq:p3_s8_4_3_multimodal}: minimum average and final displacement error over the $K$ predicted modes (minADE$_K$, minFDE$_K$), miss rate (the fraction of scenes where no mode lands within a tolerance of the ground truth), and the soft mean-average-precision of the Waymo protocol, which jointly scores accuracy and the calibration of mode probabilities $\pi_k$.
Off-road and collision rates further test the feasibility constraint $f_{\text{dyn}}$ that purely displacement-based scores ignore.

\subsubsection{Challenges and Trends}
\label{subsubsec:p3_s8_4_3_challenges}
A planner must assign trustworthy probabilities to the futures it enumerates, not merely list them: an over-confident missed mode of Eq.~\eqref{eq:p3_s8_4_3_multimodal} is a safety failure rather than a metric error, so \textbf{calibrated multimodality} is the first open problem.
\textbf{Feasibility and safety} are the second: displacement metrics reward trajectories a vehicle cannot physically execute, so feasibility has to be enforced inside the model ($f_{\text{dyn}}$) rather than checked afterward, and the long tail of near-collision scenarios is exactly where data are scarcest.
The third is \textbf{closed-loop distribution shift}: models trained to imitate logged, open-loop futures drift once their own predictions feed back into the scene, the difficulty ChauffeurNet's worst-case synthesis was built to expose.
Diffusion- and transformer-based scene-level samplers that forecast all agents jointly and run closed-loop are the current direction, and they connect directly to the controllable-video and world-model lines of §\ref{sec:p3_s8_3} and §\ref{sec:p3_s8_6}, in which a driving world model both predicts and is conditioned on the ego plan.







\subsection{Motion Editing}
\label{subsec:p3_s8_4_4}
The subsections above \emph{synthesize} motion from scratch; motion editing instead modifies an \emph{existing} motion so that it satisfies a new constraint while leaving the rest intact. It is the temporal, articulated analogue of the image editing of §\ref{subsec:p3_s8_1_3} and the video editing of §\ref{subsec:p3_s8_3_5}, and its demand is again selective preservation: change a targeted aspect (a trajectory, a joint, a style, a semantic action) while holding the body's identity, timing, and untouched portions fixed, which makes editing the controllability layer that closes the generate-and-refine loop for motion.



\subsubsection{Task Definition}
\label{subsubsec:p3_s8_4_4_task}
Given a source motion $\mathbf{m}^{\text{in}}\!\in\!\mathcal{M}$, an edit region $\mathbf{M}\!\in\!\{0,1\}^{T\times J}$ over time $T$ and joints $J$, and an edit control $\mathbf{c}$ (a text instruction, target keyframes, a desired root trajectory, or a target joint), motion editing learns an operator
\begin{equation}
\label{eq:p3_s8_4_4_objective}
E_\theta:\mathcal{M}\!\times\!\{0,1\}^{T\times J}\!\times\!\mathcal{C}'\!\to\!\mathcal{M},\qquad
\hat{\mathbf{m}}=E_\theta\!\left(\mathbf{m}^{\text{in}},\mathbf{M},\kappa(\mathbf{c})\right),
\end{equation}
which lifts the image-editing operator of Eq.~\eqref{eq:t2i_edit_objective} into the motion domain. As in image and video editing, Eq.~\eqref{eq:p3_s8_4_4_objective} encodes two competing objectives, \emph{editability} (the output must realize the control $\mathbf{c}$ inside $\mathbf{M}$) and \emph{preservation} (it must match $\mathbf{m}^{\text{in}}$ outside $\mathbf{M}$ and remain a plausible, smooth motion), and methods differ chiefly in how they trade these off.



\subsubsection{Methodology}
\label{subsubsec:p3_s8_4_4_method}
Because a pretrained motion diffusion model already samples plausible motion, most editors realize Eq.~\eqref{eq:p3_s8_4_4_objective} without retraining, by \emph{conditioning the sampler} of Eq.~\eqref{eq:p3_s8_4_1_diffusion} rather than learning a dedicated $E_\theta$. The simplest mechanism is \textbf{masked infilling}: at every denoising step the observed entries are overwritten with the (noised) source while only the masked entries are generated,
\begin{equation}
\label{eq:p3_s8_4_4_inpaint}
\hat{\mathbf{m}}=(1-\mathbf{M})\odot\mathbf{m}^{\text{in}}+\mathbf{M}\odot G_\theta\!\left(\mathbf{m}_n,\,n,\,\kappa(\mathbf{c})\right),
\end{equation}
which directly enforces preservation by construction. This is the basis of motion in-betweening, where $1-\mathbf{M}$ marks given keyframes: robust in-betweening~\cite{harvey2020robust} predates the diffusion era, and CondMDI~\cite{cohan2024condmdi} casts flexible keyframe-conditioned in-betweening as masked diffusion, while PriorMDM~\cite{shafir2024priormdm} reuses a frozen motion-diffusion prior for editing and composition tasks such as joining and in-filling. A second family handles \emph{spatial} controls that masking alone cannot express (a target root path, or a specific joint at specific times) through classifier-style guidance: GMD~\cite{karunratanakul2023gmd} and OmniControl~\cite{xie2024omnicontrol} steer the denoiser toward arbitrary spatial constraints at any joint and time. A third, optimization-based family treats the initial diffusion noise itself as the free variable, solving $\min_{\mathbf{z}}\,\mathcal{L}_{\text{edit}}(\mathcal{D}(\mathbf{z}))$ so that the decoded motion satisfies an arbitrary differentiable criterion: DNO~\cite{karunratanakul2024dno} shows this single mechanism serves editing, in-betweening, and refinement as one universal motion prior. Finally, \emph{semantic} editors take a natural-language instruction in place of a geometric mask: MotionFix~\cite{athanasiou2024motionfix} learns text-driven edits from paired source--target motions, and LLM-driven approaches~\cite{goel2024iterative} decompose an instruction into a program of primitive edits applied iteratively. Figure~\ref{fig:p3_s8_4_4_methods} groups these four mechanisms by how they obtain $E_\theta$, paralleling the editing taxonomies of §\ref{subsec:p3_s8_1_3} and §\ref{subsec:p3_s8_3_5}.

\begin{figure}[!ht]
    \centering
    \includegraphics[width=\linewidth]{figures/part3/section8/four_motion_editing_mechanisms.pdf}
    \caption{Four motion-editing mechanisms: (a) reinsert known states while infilling missing frames; (b) guide denoising with root-path or joint targets; (c) optimize initial noise through a fixed decoder using an edit loss; and (d) edit source motion from text, using a learned editor or iterative primitive edits. Rendered people illustrate edits rather than measured model outputs.}
    \label{fig:p3_s8_4_4_methods}
\end{figure}

\subsubsection{Benchmarks}
\label{subsubsec:p3_s8_4_4_benchmarks}
Motion editing inherits the corpora of §\ref{subsubsec:p3_s8_4_1_benchmarks} (HumanML3D~\cite{GuoZZ0JL022} and BABEL~\cite{PunnakkalCAQB21} over AMASS), used in editing protocols where keyframes, joints, or trajectories are held out as the constraint, together with the dedicated MotionFix~\cite{athanasiou2024motionfix} benchmark of paired source--edit--instruction triples that supplies the ground truth most editing tasks lack.
Evaluation must report the editability--preservation pair jointly: editability through control adherence (instruction-conditioned R-Precision, target-joint or trajectory error) and overall realism (FID, foot-skating), and preservation through the deviation from the source outside the edit region.
As in video editing, single scalars correlate weakly with perceived edit quality, so reporting a tuple, and increasingly a learned or VLM-style judge of instruction compliance, is becoming standard.

\subsubsection{Challenges and Trends}
\label{subsubsec:p3_s8_4_4_challenges}
The central tension is the \textbf{editability--preservation trade-off} of Eq.~\eqref{eq:p3_s8_4_4_objective}: pushing harder on a control propagates unwanted change through the kinematic chain (editing one arm disturbs balance and foot contact), so a good editor must respect the body's global coordination, not just the masked entries.
This is compounded by \textbf{sparse-control ambiguity}: a single keyframe or a root path under-determines the full-body motion, and the model must fill the gap plausibly, which the guidance and noise-optimization families address with differing robustness.
Underlying both is an \textbf{evaluation} gap sharper than in synthesis: because the ``correct'' edit is rarely unique, ground-truth-based metrics under-reward valid alternatives, motivating the paired-triplet and judge-based protocols above.
The trend mirrors the rest of §\ref{sec:p3_s8_4} and §\ref{sec:p3_s8_3}: from per-mechanism, training-free editors toward unified models that generate, control, and edit motion within a single conditional sampler, increasingly driven by natural-language instructions rather than geometric masks.








% \subsection{Summary}
% \label{subsec:p3_s8_4_summary}
% Table~\ref{tab:p3_s8_4_subfield_comparison} compares the four subfields as specializations of the human-motion sampler of Eq.~\eqref{eq:p3_s8_4_1_objective}: by the moving object and its condition, the sampler or operator each instantiates, the dominant method families, and the benchmarks on which they report.

% \begin{table}[!ht]
% \centering
% \fontsize{8}{10}\selectfont
% \setlength{\tabcolsep}{1.0mm}
% \caption{The four subfields of motion generation; each sampler or operator specializes Eq.~\eqref{eq:p3_s8_4_1_objective}.}
% \label{tab:p3_s8_4_subfield_comparison}
% \begin{tabular}{L{0.13\textwidth} L{0.19\textwidth} L{0.16\textwidth} L{0.22\textwidth} L{0.22\textwidth}}
% \toprule
% \textbf{Subfield} & \textbf{Moving object \& condition} & \textbf{Sampler / operator} & \textbf{Dominant method families} & \textbf{Benchmarks / metrics} \\
% \midrule
% §\ref{subsec:p3_s8_4_1} Human motion & Articulated body; action / text / music / speech / scene & $G_\theta^{\text{mot}}$ (Eq.~\eqref{eq:p3_s8_4_1_objective}) & Regression or generative (GAN / VAE / diffusion); per-modality conditioning & HumanML3D, KIT-ML, AIST++, BEAT; FID, R-Precision, diversity, beat / contact \\
% §\ref{subsec:p3_s8_4_2} Animal motion & Multi-species body; text $+$ species $s$ & $G_\theta^{\text{amot}}$ (Eq.~\eqref{eq:p3_s8_4_2_objective}) & Human-prior transfer $+$ species-aware modulation; morphological consistency & AnimalML3D, Zoo-300K, AniMo4D; FID, MPJPE, biomechanical plausibility \\
% §\ref{subsec:p3_s8_4_3} Vehicle planning & Rigid agent state; HD map $+$ agent histories & $G_\theta^{\text{veh}}$, $K$-mode (Eqs.~\eqref{eq:p3_s8_4_3_objective},~\eqref{eq:p3_s8_4_3_multimodal}) & Vectorized encoders; goal-based and generative decoders; imitation planners & nuScenes, Argoverse 1/2, WOMD; minADE$_K$, minFDE$_K$, miss rate, soft-mAP \\
% §\ref{subsec:p3_s8_4_4} Motion editing & Existing motion $+$ edit region / control & $E_\theta$ (Eq.~\eqref{eq:p3_s8_4_4_objective}) & Masked infilling, spatial guidance, noise optimization, instruction editing & HumanML3D / BABEL protocols, MotionFix; editability $+$ preservation tuple \\
% \bottomrule
% \end{tabular}
% \end{table}

% The four subfields partition by \emph{what is held fixed}: human, animal, and vehicle generation fix nothing but the condition and sample a fresh trajectory, whereas editing fixes the source motion and changes only a region, standing to the other three exactly as image editing (§\ref{subsec:p3_s8_1_3}) stands to image generation. Underneath that partition the dominant failure mode is shared and physical: whether it surfaces as foot-skating and penetration in human and animal motion, kinematic infeasibility in vehicle trajectories, or balance disruption when one limb is edited, every subfield is limited by how well the generator respects the biomechanics and dynamics of the moving body, which is why physics- and contact-aware modeling recurs in all four Challenges discussions. Evaluation stays fragmented, each subfield reporting its own tuple (fidelity, diversity, condition consistency; calibrated multimodality; editability and preservation), and all are migrating toward learned or perceptual judges as scalar metrics prove weakly correlated with human assessment. Motion generation is thus the point in §\ref{sec:generation} where generative spatial modeling must answer to physical law as well as semantics, which leads directly to the dynamic 4D generation of §\ref{sec:p3_s8_5}.


% \newpage
\section{4D Spatial Generation}
\label{sec:p3_s8_5}




\FramedInitial{G}{enerative} modeling has moved outward from static images~\cite{rombach2022ldm,saharia2022photorealistic} to video~\cite{chen2024videocrafter2,chen2023videocrafter1,ho2022video,harvey2022flexible}, multi-view imagery~\cite{shi2023mvdream,liu2023zero,shi2023zero123++,zhao2024sg,zhao2024chase}, and static 3D assets~\cite{wang2023score,raj2023dreambooth3d}, yet each stops short of a scene that is at once three-dimensional and alive.
A 3D asset holds geometry but not motion; a video holds motion, but no view-consistent geometry one can navigate or simulate.
4D generation is where the two meet: the synthesis of dynamic 3D content whose geometry stays view-consistent as it moves and deforms over time~\cite{singer2023text,zheng2024unified,jiang2023consistent4d,ling2024align,wang2025moving}.
This is the representation an embodied system needs to rehearse how a scene will unfold, and it underwrites photorealistic digital humans~\cite{guo2023vid2avatar,zhao2024sg,zhao2024chase}, interactive AR/VR~\cite{zhao2024hfgs,miao2024pla4d}, and high-fidelity simulation for autonomous driving and embodied AI~\cite{gao2024magicdrive3d,zhao2024efficient}.

Two advances make this tractable.
Temporal 3D representations, deformation-aware fields over NeRFs~\cite{mildenhall2021nerf}, point clouds, and Gaussian splatting~\cite{kerbl20233d,zhao2024efficient}, supply a geometry that can bend and move without losing view consistency, while the extension of generative diffusion into this domain supplies the realism and controllability to drive it.
Together they open dynamic scene synthesis, 4D editing, and interactive agent simulation, at the cost of the difficulty this section returns to throughout: holding long-range temporal coherence while still granting flexible, user-driven control.

Formally, we cast 4D generation as the latent-generator framework of §\ref{subsec:p3_s8_1_1} unrolled along time over a shared 3D representation. Let $\mathcal{R}$ be the 3D-representation space of Eq.~\eqref{eq:p3_s8_2_1_objective} (mesh, NeRF, 3DGS, point cloud, and so on), and let $\kappa:\mathcal{C}\!\to\!\mathcal{C}'$ encode a conditioning signal $\mathbf{c}\in\mathcal{C}$ (text, an image, a video, a static 3D asset, or a combination thereof). 4D generation learns a conditional sampler
\begin{equation}
\label{eq:p3_s8_5_objective}
G_\theta^{\text{4D}}:\mathcal{C}'\!\to\!\mathcal{R}^{[0,T]},\qquad
\hat{\mathbf{r}}^{1:T}\sim p_\theta\!\left(\mathbf{r}^{1:T}\mid\kappa(\mathbf{c})\right),
\end{equation}
that emits a time-indexed sequence of 3D representations $\hat{\mathbf{r}}^{1:T}$ with view-consistent geometry and coherent motion. Two limiting cases anchor the framework: with a single time step ($T=1$) Eq.~\eqref{eq:p3_s8_5_objective} recovers the static 3D generator of Eq.~\eqref{eq:p3_s8_2_1_objective}, and rendering each frame $\mathbf{r}^t$ to a 2D view recovers the video generator of Eq.~\eqref{eq:p3_s8_3_objective}, so 4D generation is precisely the 3D generator of §\ref{sec:p3_s8_2} unrolled in time, or equivalently the video generator of §\ref{sec:p3_s8_3} lifted from per-frame pixels to a shared, renderable geometry.
The subsections are organized by the conditioning input (text, image, video, a static 3D asset, or several at once). Figure~\ref{fig:p3_s8_5_taxonomy} previews this organization together with two methodological strategies that recur across all five, per-scene score distillation and explicit ``generate-then-reconstruct'' supervision; feed-forward generators and rig-based animation of meshes appear alongside them in the subsections where they apply.

\begin{figure}[!ht]
    \centering
    \includegraphics[width=\linewidth]{figures/part3/section8/taxonomy_4d_spatial_generation.pdf}
    \caption{4D spatial generation by conditioning input, (a)--(e), with two representative strategies: (i) score distillation from a frozen diffusion prior (dashed: SDS feedback), and (ii) generating video or multiview targets, then reconstructing a dynamic 3D asset. The strategies are alternatives that may be combined; the chair sequence is illustrative.}
    \label{fig:p3_s8_5_taxonomy}
\end{figure}






\subsection{Text-to-4D Generation}
\label{subsec:p3_s8_5_1}
Text-to-4D generation synthesizes dynamic content directly from natural language, supporting text-driven asset creation and character animation~\cite{HongZPCYL22,chen2024ct4d}. Its progress has been propelled by large language models for semantic understanding and by video foundation models that supply motion priors.


\subsubsection{Task Definition}
\label{subsubsec:p3_s8_5_1_task}
Text-to-4D generation specializes the 4D sampler of Eq.~\eqref{eq:p3_s8_5_objective} to a textual condition $\mathbf{c}=\mathbf{x}$ encoded by the text encoder $\phi$ of Eq.~\eqref{eq:t2i_generator},
\begin{equation}
\label{eq:p3_s8_5_1_objective}
\hat{\mathbf{r}}^{1:T}\sim p_\theta\!\left(\mathbf{r}^{1:T}\mid\phi(\mathbf{x})\right).
\end{equation}
Because language fixes only semantics and leaves geometry and motion grossly under-determined, the sampler must hallucinate both a plausible 3D shape and its temporal evolution from the prompt alone, the source of the semantic-to-motion ambiguity examined in §\ref{subsubsec:p3_s8_5_1_challenges}.



\subsubsection{Methodology}
\label{subsubsec:p3_s8_5_1_method}
Two strategies recur throughout this section, implicit distillation and explicit supervision, implemented across various 4D representations (Figure~\ref{fig:p3_s8_5_taxonomy}).
The more common implicit strategy leverages Score Distillation Sampling (SDS) to indirectly optimize the geometry and motion of a 4D representation, using pre-trained diffusion models as frozen priors~\cite{singer2023text,zheng2024unified}. Formally, it applies a score-distillation update to per-scene parameters $\theta$~\cite{poole2022dreamfusion},
\begin{equation}
\label{eq:p3_s8_5_sds}
\nabla_\theta\mathcal{L}_{\mathrm{SDS}}
=\mathbb{E}_{s,\epsilon,\pi}\!\left[
 w(s)\left(\epsilon_\phi(\mathbf{z}_s;s,\kappa(\mathbf{c}))-\epsilon\right)
 \frac{\partial\mathbf{x}_\theta}{\partial\theta}\right],
\quad \mathbf{z}_s=\alpha_s\mathbf{x}_\theta+\sigma_s\epsilon,
\end{equation}
where $\mathbf{x}_\theta=\mathrm{Render}(\hat{\mathbf{r}}^{1:T}_\theta,\pi)$ is a differentiable rendering at camera $\pi$, $s$ indexes diffusion noise levels, and $\epsilon_\phi$ is a frozen 2D/video diffusion prior. The update differentiates through the renderer while treating the prior noise residual as the supplied gradient signal; this lifts the static score-distillation objective of §\ref{para:p3_s8_2_1_sds} to a temporal sequence.
In contrast, explicitly supervised methods initialize the 4D generation from text-to-3D or text-to-video outputs and use these generated assets as direct supervisory signals~\cite{zhao2024efficient,wu20244d,yang2024deformable,ling2024align,blattmann2023align,xu2024comp4d}. This ``generate-then-reconstruct'' strategy first synthesizes target frames $\hat{\mathbf{y}}$ with a fast generator $\mathcal{G}_{\text{2D}}$ and then fits the 4D asset to them,
\begin{equation}
\label{eq:p3_s8_5_explicit}
\hat{\mathbf{y}}=\mathcal{G}_{\text{2D}}(\mathbf{c}),\qquad
\hat{\mathbf{r}}^{1:T}=\arg\min_{\mathbf{r}^{1:T}}\sum_{\pi,t}\big\|\mathrm{Render}(\mathbf{r}^{t},\pi)-\hat{\mathbf{y}}^{t}_{\pi}\big\|^2,
\end{equation}
replacing score-distillation guidance with explicit pixel targets; reconstruction may itself require per-scene optimization, so a speed advantage depends on the implementation. The two strategies are realized across mesh, NeRF, and 3DGS representations.

\textbf{Mesh-based Methods.} Meshes offer explicit geometric control, which suits precise text-driven shape manipulation and animation.
CT4D~\cite{chen2024ct4d} implements a three-stage ``Generate-Refine-Animate'' pipeline in which a mesh is generated from text, its texture refined, and its vertices animated; TextMesh4D~\cite{dai2025textmesh4d} deforms a generated mesh through a per-face Jacobian deformation field under video-diffusion priors; and AvatarCLIP~\cite{HongZPCYL22} generates a static avatar zero-shot and animates it by registering motion to an SMPL skeleton.

\textbf{NeRF-based Methods.} NeRFs represent continuous geometry and complex temporal variation well.
MAV3D~\cite{singer2023text} introduced multi-stage SDS optimization that refines texture, geometry, and motion with priors from both T2I and T2V diffusion models; 4D-fy~\cite{bahmani20244d} blends 3D-aware, image, and video score signals in a hybrid SDS; and Dream-in-4D~\cite{zheng2024unified} first generates a static asset with a multiview diffusion model and then adds dynamics with a video diffusion model.

\textbf{Gaussian Splatting-based Methods.} 4D Gaussian Splatting (4DGS)~\cite{zhao2024hfgs,wu20244d} renders efficiently and differentiably, which speeds up text-driven dynamic scene synthesis. Early methods followed the two-stage implicit pipeline of Eq.~\eqref{eq:p3_s8_5_sds}: AYG~\cite{ling2024align} generates a static 3D Gaussian model from text via SDS and adds dynamics under a T2V prior, and Comp4D~\cite{xu2024comp4d} uses an LLM to parse a complex prompt into several dynamic objects and their trajectories for a compositional scene.
Later works moved to explicit supervision (Eq.~\eqref{eq:p3_s8_5_explicit}): PLA4D~\cite{miao2024pla4d} aligns the 4D Gaussians pixel by pixel with generated video targets, and 4Real~\cite{yu20244real} uses generated multiview video frames as the reference for photorealistic scene-level generation.
For digital humans, AvatarGO~\cite{cao2024avatargo} generates 4D human--object interactions zero-shot by animating a text-to-3D human with the SMPL-X skinning function.


\subsubsection{Benchmarks}
\label{subsubsec:p3_s8_5_1_benchmarks}
Text-to-4D, like the other subfields of this section, has no standardized benchmark, because large-scale dynamic 3D data with text descriptions is hard to capture and annotate. Training and evaluation instead adapt static 3D repositories (Table~\ref{tab:4d_datasets}): Diffusion4D~\cite{liang2024diffusion4d} filtered Objaverse and Objaverse-XL~\cite{deitke2023objaverse,deitke2023objaversexl} down to 54K high-quality animated assets, and SV4D~\cite{xie2024sv4d} built ObjaverseDy by adaptive sampling and viewpoint adjustment. These are synthetic renderings and lack the multimodal annotations, text in particular, needed to train and evaluate text-conditioned models robustly.

\begin{table}[!ht]
\centering
\fontsize{8}{10}\selectfont
\setlength{\tabcolsep}{1.0mm}
\caption{Datasets used for 4D generation across the subfields of §\ref{sec:p3_s8_5}.}
\label{tab:4d_datasets}
\begin{tabular}{L{0.20\textwidth} L{0.24\textwidth} C{0.20\textwidth} L{0.26\textwidth}}
\toprule
\textbf{Dataset} & \textbf{Content} & \textbf{Scale} & \textbf{Used for} \\
\midrule
Objaverse-XL~\cite{deitke2023objaversexl} & Static and animated 3D objects & 10M+ objects & Source pool for dynamic subsets \\
Diffusion4D~\cite{liang2024diffusion4d} & Animated assets curated from Objaverse / Objaverse-XL & 54K assets & Training image- / video-to-4D \\
ObjaverseDy~\cite{xie2024sv4d} & Animated Objaverse objects, re-rendered & - & Training and evaluating video-to-4D \\
Consistent4D~\cite{jiang2023consistent4d} & Synthetic and real monocular videos & - & Evaluating video-to-4D \\
\bottomrule
\end{tabular}
\end{table}

\paragraph{Evaluation Metrics.} Evaluation across 4D generation combines user studies with metrics borrowed from image and video generation, and the subsections that follow use this same set. Semantic alignment with a prompt is measured by CLIP Score~\cite{radford2021learning} between rendered views and the text; fidelity by PSNR~\cite{hore2010image}, SSIM~\cite{wang2004image}, and LPIPS~\cite{zhang2018unreasonable} where a reference view exists; temporal quality and motion by the Fréchet Video Distance (FVD)~\cite{unterthiner2018fvd}; and overall plausibility and aesthetics by user studies. None of these scores the two properties that make content 4D, cross-view consistency and long-range spatio-temporal coherence.


\subsubsection{Challenges and Trends}
\label{subsubsec:p3_s8_5_1_challenges}
The challenge specific to text-to-4D is semantic-to-motion ambiguity: a prompt such as ``a dancing dog'' admits endlessly many movements, styles, and trajectories, and methods that lean on generic T2V priors struggle to produce diverse, fine-grained motion that matches a nuanced description.
One response uses LLMs for motion planning, as Comp4D~\cite{xu2024comp4d} does when it decomposes a prompt into a structured spatio-temporal plan for several interacting objects; another targets human--object interaction~\cite{cao2024avatargo}, moving toward story-driven dynamic scenes.

% ======================================================================
% ======================================================================


\subsection{Image-to-4D Generation}
\label{subsec:p3_s8_5_2}
Image-to-4D generation synthesizes dynamic content from a single static image, combining the spatial information in the image with temporal priors learned by video models~\cite{yang2024diffusion,sun2024eg4d}. No large image--4D paired dataset exists, which has ruled out direct end-to-end training and pushed the field toward optimization and toward generated supervision.

\subsubsection{Task Definition}
\label{subsubsec:p3_s8_5_2_task}
Image-to-4D generation specializes the 4D sampler of Eq.~\eqref{eq:p3_s8_5_objective} to a single-image condition $\mathbf{c}=I$,
\begin{equation}
\label{eq:p3_s8_5_2_objective}
\hat{\mathbf{r}}^{1:T}\sim p_\theta\!\left(\mathbf{r}^{1:T}\mid\kappa(I)\right),
\end{equation}
where the source view at $t=0$ is fixed by $I$ while every unseen viewpoint and time step must be inferred.
Both the unseen 3D geometry and the entirely absent temporal dynamics have to be inferred from one 2D view, which requires combining the spatial cues of the input with temporal priors from large video models~\cite{ren2023dreamgaussian4d,zhao2023animate124,sun2024eg4d}.



\subsubsection{Methodology}
\label{subsubsec:p3_s8_5_2_method}
Both strategies of §\ref{subsubsec:p3_s8_5_1_method} apply, and most methods use 4D Gaussian Splatting as the representation~\cite{ren2023dreamgaussian4d,yang2024diffusion}.

\textbf{Implicit Optimization with Diffusion Priors.} These methods instantiate Eq.~\eqref{eq:p3_s8_5_sds}, typically combining a multiview diffusion model for geometric consistency with a video diffusion model for temporal coherence. Animate124~\cite{zhao2023animate124} first generates a static NeRF from the image and then refines its dynamics with a video-diffusion SDS loss, and Diffusion\textsuperscript{2}~\cite{yang2024diffusion} reconciles the scores of heterogeneous priors (multiview and video) through Variance-Reducing Sampling, improving both multiview and temporal consistency of the resulting 4D Gaussians.

\textbf{Explicit Supervision with Generated Data.} The slow, iterative nature of SDS motivates the two-stage pipeline of Eq.~\eqref{eq:p3_s8_5_explicit}, in which an image-to-video model first generates an intermediate video that then supervises the 4D reconstruction directly. DreamGaussian4D (DG4D)~\cite{ren2023dreamgaussian4d} generates a static 3D Gaussian model and learns a deformation field to animate it, using the generated video for texture refinement and dynamic supervision, and EG4D~\cite{sun2024eg4d} avoids SDS entirely by chaining image-to-video and image-to-multiview models into coherent image sequences that guide a fast reconstruction.

\textbf{Unified and Tuning-free Generation.} GenXD~\cite{zhao2024genxd} trains one model on both 3D and 4D data, using multiview-temporal modules that separate camera motion from object motion, and Free4D~\cite{liu2025free4d} builds a 4D scene from a single image without fine-tuning any model, distilling pretrained image-to-video and multiview priors into a consistent 4D representation.

\textbf{Specialized Methods for Human Generation.} Dynamic human avatars form a sub-domain in which strong body priors raise fidelity. Human4DiT~\cite{shao2024360} uses a hierarchical 4D transformer that integrates human-specific features, camera parameters, and temporal signals to generate 360-degree human animations from a single image, and Disco4D~\cite{pang2025disco4d} models the body with SMPL-X while representing clothing separately with Gaussian splats, allowing detailed animation of clothed humans.



\subsubsection{Benchmarks}
\label{subsubsec:p3_s8_5_2_benchmarks}
With no image--4D pairs, methods are evaluated on images from Objaverse~\cite{deitke2023objaverse} or on in-the-wild photographs chosen by the authors, which makes comparison across papers unreliable. The metric set is that of §\ref{subsubsec:p3_s8_5_1_benchmarks}, applied to the source view and to pseudo-ground-truth novel views where available.

\subsubsection{Challenges and Trends}
\label{subsubsec:p3_s8_5_2_challenges}
Inferring a complete shape and a plausible motion from one viewpoint is extremely ill-posed, which shows up as geometric inconsistency in unseen views (\eg, the ``Janus problem'') and as unnatural or repetitive motion, since every dynamic detail must be hallucinated.
Domain-specific priors are the main remedy: Human4DiT~\cite{shao2024360} and Disco4D~\cite{pang2025disco4d} regularize geometry and motion with parametric body models and reach much higher fidelity, at the cost of generality.


% ======================================================================
% ======================================================================


\subsection{Video-to-4D Generation}
\label{subsec:p3_s8_5_3}
Video-to-4D generation lifts a single-view video into a dynamic 4D representation that can be viewed from new angles with consistent geometry and motion, turning ordinary video into interactive 4D content. Annotated video-to-4D datasets are hard to curate, so current methods build on pretrained video diffusion priors, either as multi-view supervision or through score distillation.



\subsubsection{Task Definition}
\label{subsubsec:p3_s8_5_3_task}
Video-to-4D generation specializes the 4D sampler of Eq.~\eqref{eq:p3_s8_5_objective} to a monocular video condition $\mathbf{c}=\mathbf{v}^{1:T}$,
\begin{equation}
\label{eq:p3_s8_5_3_objective}
\hat{\mathbf{r}}^{1:T}\sim p_\theta\!\left(\mathbf{r}^{1:T}\mid\kappa(\mathbf{v}^{1:T})\right),
\end{equation}
so that re-rendering $\hat{\mathbf{r}}^{1:T}$ from the input viewpoint reproduces $\mathbf{v}^{1:T}$ while novel views remain geometrically consistent.
Unlike image-to-4D (Eq.~\eqref{eq:p3_s8_5_2_objective}), the temporal dynamics are now \emph{observed} and the residual ambiguity is purely geometric: the unseen off-axis structure.

\subsubsection{Methodology}
\label{subsubsec:p3_s8_5_3_method}
Most methods densify the monocular input into multi-view supervision and reconstruct against it (Eq.~\eqref{eq:p3_s8_5_explicit}), some add score distillation (Eq.~\eqref{eq:p3_s8_5_sds}), and the newest predict the 4D asset feed-forward.

\textbf{Mesh-based Methods.} Meshes are interpretable and efficient to animate. DreamMesh4D~\cite{li2024dreammesh4d} obtains a coarse mesh and drives it with a sparse deformation graph and linear blend skinning (LBS), reaching strong spatio-temporal consistency more cheaply than implicit volumetric methods, and V2M4~\cite{chen2025v2m4} reconstructs an animated, textured mesh asset directly from a monocular video.

\textbf{NeRF-based Methods.} NeRF-based methods model the continuous volumetric content of a video well, and they overcome the single view by generating data. Consistent4D~\cite{jiang2023consistent4d} densifies the monocular input through temporal and multiview interpolation before optimizing a NeRF with SDS, 4Diffusion~\cite{zhang20244diffusion} fine-tunes a video diffusion model with a learnable motion module on synthetic multiview data so that it produces consistent multiview video to guide generation, and SV4D~\cite{xie2024sv4d} completes missing views and frames with a diffusion model and uses the resulting multiview video to supervise the reconstruction.

\textbf{Gaussian Splatting-based Methods.} 4DGen~\cite{yin20234dgen} grounds 4D Gaussian generation in the input video with spatial-temporal pseudo labels; Efficient4D~\cite{pan2024efficient4d} converts the video into a multiview image matrix with a multiview-consistent generator (SyncDreamer) for rapid reconstruction; SC4D~\cite{wu2024sc4d} first fits sparse control Gaussians for coarse shape and motion and then drives dense Gaussians through LBS; and L4GM~\cite{ren2024l4gm} predicts the 4D Gaussians of a whole clip in a single feed-forward pass, trained on rendered animated objects.


\subsubsection{Benchmarks}
\label{subsubsec:p3_s8_5_3_benchmarks}
Video-to-4D has the one dedicated evaluation set in this section, released with Consistent4D~\cite{jiang2023consistent4d}: synthetic and real monocular videos with training and test splits. The metrics are those of §\ref{subsubsec:p3_s8_5_1_benchmarks}, with reconstruction fidelity measured at the input view and novel-view quality measured where synthetic ground-truth multiview video exists.

\subsubsection{Challenges and Trends}
\label{subsubsec:p3_s8_5_3_challenges}
The video constrains motion but not the unseen 3D structure, which leads to geometric distortion, depth ambiguity, and a ``painted-on'' look in novel views.
Explicit geometric structure is the main remedy: DreamMesh4D~\cite{li2024dreammesh4d} and SC4D~\cite{wu2024sc4d} constrain the reconstruction with meshes or sparse control points under LBS, which improves both coherence and efficiency, and combining such structure with the appearance modeling of neural representations is the obvious next step.

% ======================================================================
% ======================================================================

\subsection{3D-to-4D Generation}
\label{subsec:p3_s8_5_4}
3D-to-4D generation animates a static 3D model, so that the geometry is given and only motion and deformation must be synthesized, which makes the large stock of existing 3D assets usable for simulation, virtual reality, and character animation.
Unconditional generators of dynamic neural fields, such as HyperDiffusion~\cite{erkocc2023hyperdiffusion}, which runs diffusion directly on the weights of neural fields, sample 4D content without a given asset; 3D-to-4D instead starts from one, either through explicit mesh deformation and rigging or through neural and physics-based deformation models.

\subsubsection{Task Definition}
\label{subsubsec:p3_s8_5_4_task}
Here the condition is itself a 3D representation $\mathbf{c}=\mathbf{r}_0\in\mathcal{R}$, so only its temporal evolution must be generated. We therefore specialize Eq.~\eqref{eq:p3_s8_5_objective} to a \emph{deformation operator} that takes the static asset and an optional auxiliary driving signal $\mathbf{c}_{\text{aux}}$ (text, audio, or a motion sketch) to a sequence,
\begin{equation}
\label{eq:p3_s8_5_4_objective}
\hat{\mathbf{r}}^{1:T}=\mathcal{D}_\theta\!\left(\mathbf{r}_0,\,\kappa(\mathbf{c}_{\text{aux}})\right),\qquad \hat{\mathbf{r}}^{0}=\mathbf{r}_0,
\end{equation}
with the first frame pinned to the input asset. This places 3D-to-4D in the same family as the editing operators of §\ref{subsec:p3_s8_2_5} and §\ref{subsec:p3_s8_4_4}: it transforms an existing representation rather than synthesizing one from scratch.

\subsubsection{Methodology}
\label{subsubsec:p3_s8_5_4_method}
Methods realize the deformation operator $\mathcal{D}_\theta$ of Eq.~\eqref{eq:p3_s8_5_4_objective} either with explicit mesh structure, increasingly through learned rigging, or with neural and physics-informed deformation fields.

\textbf{Mesh-based Methods.} Meshes give explicit structure and precise control over deformation. Motion3DGAN~\cite{otberdout2022sparse} generates 3D facial animation from a static face mesh by first predicting the displacement of facial landmarks for a target expression and then deforming the full mesh, with a GAN keeping the sequence temporally coherent. A recent line makes arbitrary meshes animatable by predicting the rig itself: MagicArticulate~\cite{song2025magicarticulate} and RigAnything~\cite{liu2025riganything} predict skeletons and skinning weights feed-forward, Puppeteer~\cite{song2025puppeteer} rigs a mesh and then animates it by fitting the rig to videos generated from renderings of the mesh, and AnimateAnyMesh~\cite{wu2025animateanymesh} generates text-driven mesh animation in a single forward pass with a rectified-flow model over a dynamic-mesh latent.

\textbf{NeRF-based Methods.} ElastoGen~\cite{feng2024elastogen} generates physically grounded elastodynamic behavior: it rasterizes a static 3D model and then predicts and refines deformation fields iteratively with recurrent and 3D convolutional modules, so that the motion obeys physical principles.

\subsubsection{Benchmarks}
\label{subsubsec:p3_s8_5_4_benchmarks}
No benchmark with ground-truth dynamic sequences exists; methods are evaluated on a few individual assets, often faces or generic objects, mostly by user study and with the metrics of §\ref{subsubsec:p3_s8_5_1_benchmarks} where a reference animation is available. Physical plausibility, specific to this subfield, can be measured for physics-informed methods such as ElastoGen~\cite{feng2024elastogen} as deviation from a high-fidelity simulator or as violation of conserved quantities such as energy and momentum.

\subsubsection{Challenges and Trends}
\label{subsubsec:p3_s8_5_4_challenges}
Even with the geometry given, the space of plausible animations of a static model is unbounded, so without guidance the task is unconstrained, and current methods specialize to narrow categories such as faces.
Physical priors constrain it (ElastoGen~\cite{feng2024elastogen} grounds motion in elastodynamics), learned rigging makes arbitrary assets animatable, and pairing the asset with text, audio, or motion sketches turns open-ended generation into controllable animation.

% ======================================================================
% ======================================================================


\subsection{Multi-conditional 4D Generation}
\label{subsec:p3_s8_5_5}
Multi-conditional 4D generation steers dynamic scene synthesis with several signals at once: semantic conditions (text), appearance conditions (images, video), structural conditions (skeletal poses, explicit trajectories), and rule-based priors (differentiable physics). The central design problem is to fuse them so that each is respected without one overpowering the others.

\subsubsection{Task Definition}
\label{subsubsec:p3_s8_5_5_task}
We generalize the single-condition cases of Eqs.~\eqref{eq:p3_s8_5_1_objective}--\eqref{eq:p3_s8_5_4_objective} to a set of $K$ conditions, each mapped by its own encoder $\kappa_k$ and fused into the shared embedding:
\begin{equation}
\label{eq:p3_s8_5_5_objective}
\hat{\mathbf{r}}^{1:T}\sim p_\theta\!\left(\mathbf{r}^{1:T}\mid \kappa_1(\mathbf{c}_1),\dots,\kappa_K(\mathbf{c}_K)\right),
\end{equation}
which recovers a single-modality sampler when $K=1$.
Users can then specify not only what is generated (semantics) but how it behaves (motion path, physical behavior, artistic style), which complex applications such as character animation and driving simulation require.

\subsubsection{Methodology}
\label{subsubsec:p3_s8_5_5_method}
Methods fuse the encoders $\kappa_k$ of Eq.~\eqref{eq:p3_s8_5_5_objective}, usually combining score distillation (Eq.~\eqref{eq:p3_s8_5_sds}) with explicit supervision (Eq.~\eqref{eq:p3_s8_5_explicit}) tailored to each condition.

\textbf{Motion and Trajectory Control.} TC4D~\cite{bahmani2024tc4d} lets users define a spline trajectory that guides both the global rigid transformation and the local deformation of an object, with text supplying the semantics. For human animation, skeletons are the natural condition: Beyond Skeletons~\cite{yang2024beyond} integrates text and image conditions with the skeleton to refine motion, and STAR~\cite{chai2024star} uses skeleton-aware SDS and motion retargeting so that avatar movement is both visually coherent and anatomically accurate.

\textbf{Physics-Informed Control.} Sync4D~\cite{fu2024sync4d} and Phy124~\cite{lin2024phy124} integrate differentiable Material Point Method (MPM) simulation, so that the generated content follows realistic dynamics such as elasticity and collisions under video or text guidance.

\textbf{User-Centric and Artistic Control.} Generative Rendering~\cite{cai2024generative} uses a pre-existing dynamic mesh to guide a pretrained text-to-image model, giving precise control over the style and detail of the rendered output, and sketch-based interfaces~\cite{yang2024sketch} let users draw the dynamics they want.

\textbf{Complex, Application-Specific Control.} MagicDrive3D~\cite{gao2024magicdrive3d} generates dynamic driving scenes from 2D bounding boxes, Bird's-Eye-View (BEV) maps, agent trajectories, and text together, producing tightly controlled traffic scenarios.


\subsubsection{Benchmarks}
\label{subsubsec:p3_s8_5_5_benchmarks}
No benchmark exists; building one would require 4D assets annotated with several aligned conditions (\eg, a clip with a description, a skeletal motion sequence, and a trajectory map). Text adherence is scored with the CLIP Score of §\ref{subsubsec:p3_s8_5_1_benchmarks}, while adherence to sketches or trajectories relies on task-specific errors (\eg, distance from the specified trajectory) or on user studies, which remain the only holistic judge of how well a combination of inputs was followed.

\subsubsection{Challenges and Trends}
\label{subsubsec:p3_s8_5_5_challenges}
Semantic guidance from text, geometric constraints from skeletons, and rule-based priors from physics can contradict one another, and keeping each respected without one overpowering the rest while preserving fidelity is the open problem.
Domain-specific stacks such as MagicDrive3D~\cite{gao2024magicdrive3d} sidestep it by fixing the set of inputs for one application, and interactive interfaces such as sketches~\cite{yang2024sketch} point toward iterative creation in which several control channels are adjusted in turn.

% ======================================================================
% ======================================================================


\subsection{Summary, Overarching Challenges and Future Directions}
\label{subsec:p3_s8_5_summary}
Table~\ref{tab:p3_s8_5_subfield_comparison} compares the five subfields as specializations of the 4D sampler of Eq.~\eqref{eq:p3_s8_5_objective}: image- and video-to-4D progressively replace hallucinated structure with observed pixels, 3D-to-4D is given the geometry and generates only motion, and multi-conditional generation fuses several signals at once.

\begin{table}[!ht]
\centering
\fontsize{8}{10}\selectfont
\setlength{\tabcolsep}{1.0mm}
\caption{The five subfields of 4D spatial generation; each sampler or operator specializes Eq.~\eqref{eq:p3_s8_5_objective}.}
\label{tab:p3_s8_5_subfield_comparison}
\begin{tabular}{L{0.13\textwidth} L{0.15\textwidth} L{0.17\textwidth} L{0.25\textwidth} L{0.22\textwidth}}
\toprule
\textbf{Subfield} & \textbf{Conditioning input} & \textbf{Sampler / operator} & \textbf{Dominant method families} & \textbf{Datasets / metrics} \\
\midrule
§\ref{subsec:p3_s8_5_1} Text-to-4D & Text prompt $\mathbf{x}$ & $G_\theta^{\text{4D}}$, $\kappa=\phi$ (Eq.~\eqref{eq:p3_s8_5_1_objective}) & SDS distillation or generate-then-reconstruct; mesh / NeRF / 3DGS; LLM motion planning & Objaverse-derived; CLIP, FVD, PSNR/SSIM/LPIPS, user study \\
§\ref{subsec:p3_s8_5_2} Image-to-4D & Single image $I$ & $G_\theta^{\text{4D}}$, $\kappa(I)$ (Eq.~\eqref{eq:p3_s8_5_2_objective}) & Diffusion-prior SDS; explicit video supervision; unified / tuning-free; human-specific priors & Objaverse / in-the-wild; PSNR/SSIM/LPIPS, FVD, user study \\
§\ref{subsec:p3_s8_5_3} Video-to-4D & Monocular video $\mathbf{v}^{1:T}$ & $G_\theta^{\text{4D}}$, $\kappa(\mathbf{v}^{1:T})$ (Eq.~\eqref{eq:p3_s8_5_3_objective}) & Video-diffusion densification; explicit geometric constraints (mesh / LBS); feed-forward & Consistent4D; PSNR/SSIM/LPIPS, FVD, user study \\
§\ref{subsec:p3_s8_5_4} 3D-to-4D & Static asset $\mathbf{r}_0$ & $\mathcal{D}_\theta$ deformation (Eq.~\eqref{eq:p3_s8_5_4_objective}) & Mesh deformation and learned rigging; physics-informed deformation & Per-asset (faces / objects); FVD, physical plausibility, user study \\
§\ref{subsec:p3_s8_5_5} Multi-conditional & $K$ fused signals & $G_\theta^{\text{4D}}$, $\{\kappa_k\}$ (Eq.~\eqref{eq:p3_s8_5_5_objective}) & Trajectory / skeleton / physics / sketch control; application-specific stacks & No standard suite; CLIP, PSNR/SSIM/LPIPS, FVD, user study \\
\bottomrule
\end{tabular}
\end{table}

\paragraph{Data Scarcity and Generative Priors.}
Large, annotated 4D datasets are the field's main missing ingredient: real dynamic 3D data with multi-view consistency and multimodal annotation (text, motion capture) is expensive and technically hard to collect, which has largely ruled out end-to-end training. The field has therefore converged on distilling pretrained text-to-image, text-to-video, and text-to-3D models, the de facto way to guide 4D synthesis.

\paragraph{From Score Distillation to Explicit Supervision.}
Score distillation~\cite{poole2022dreamfusion} converges slowly, is costly and unstable, and produces multi-faced geometry and flickering textures. The generate-then-reconstruct pipeline of DreamGaussian4D~\cite{ren2023dreamgaussian4d} and EG4D~\cite{sun2024eg4d}, which first generates target videos and then reconstructs against them, is replacing per-prompt SDS, and feed-forward predictors such as L4GM~\cite{ren2024l4gm} and AnimateAnyMesh~\cite{wu2025animateanymesh} take the next step by removing per-scene optimization altogether.

\paragraph{4D Gaussian Splatting as the Default Representation.}
NeRF-based works rendered well but trained and rendered slowly. 3D Gaussian Splatting~\cite{kerbl20233d} and its temporal extension offer real-time, high-fidelity rendering with efficient gradient-based optimization, and their explicit, point-based form is easier to edit and to couple with physical simulation, so 4DGS has become the representation of choice for current methods~\cite{wu20244d, ren2023dreamgaussian4d, ling2024align}.

\paragraph{Consistency, Control, and Physical Plausibility.}
Keeping geometry, appearance, and motion coherent across all viewpoints and over long horizons remains the central difficulty, and control is moving beyond single prompts toward fine-grained manipulation of motion, appearance, and interaction. Physical plausibility is the newest goal: integrating differentiable simulators, as PhysSplat~\cite{zhao2024efficient} and Phy124~\cite{lin2024phy124} do, makes generated dynamics more realistic and predictable. Progress on all three is hard to measure while evaluation relies on small, author-curated test sets and on image and video metrics that ignore cross-view and long-range consistency; 4D-aware benchmarks and metrics are a prerequisite for the rest.

% \newpage


\section{World Modeling}
\label{sec:worldmodeling}
\label{sec:p3_s8_6}

\FramedInitial{D}{espite} its growing popularity across AI research, the notion of a ``world model'' does not yet enjoy a unified or universally agreed-upon definition.
Its conceptual roots trace back to the psychologist Kenneth Craik, who argued in his 1943 monograph \emph{The Nature of Explanation} that the mind ``carries a `small-scale model' of external reality and of its own possible actions within its head''.
In its original form, a world model referred to an individual's internal representation of external reality; contemporary AI has generalized the concept far beyond mental representations to structured models of physical and spatial environments.

Fei-Fei Li has framed the gap between current AI systems and physical understanding in these terms~\cite{Feifei2025spatial}: modern LLMs resemble ``wordsmiths in the dark'', fluent in symbols but divorced from the physical world that gives them meaning, and capable AI must instead understand, act within, and reason about three-dimensional space. On this view world models are new-generation \emph{generative} models that simulate, predict, and imagine the world's dynamics rather than only describe it, which ties world modeling directly to spatial intelligence.



Complementing Li's generative perspective, another influential formulation comes from Yann~LeCun \cite{lecun2022path}. He describes a world model as a system capable of predicting how the world will evolve given its current state and an action.
Formally, it consumes the current state $S(t)$ and an action $A(t)$, and produces a prediction of the next state $S'(t+1)$; in many cases, it also integrates external sensory observations $O(t)$ to update its internal state.
This predictive relation can be written compactly as
\begin{equation}
\label{eq:p3_s8_6_transition}
\hat{S}(t+1) = W_\theta\big(S(t),\, A(t),\, O(t)\big),
\end{equation}
where $W_\theta$ denotes the learned world model parameterized by $\theta$.
Read against the conditional-generation backbone of this chapter (Eq.~\eqref{eq:t2i_generator}), $W_\theta$ occupies the position of the generator, but its decoder is a state-transition predictor rather than a pixel or geometry renderer; this is the sense in which world modeling extends Chapter~\ref{sec:generation} from reconstruction toward prediction.
What exactly constitutes a ``state'' remains open: some work defines it as a compact latent representation in which the world model operates at an abstract semantic level, much as human cognition forecasts events without simulating every physical detail.




Although no single definition of world modeling has been universally accepted, a survey of the literature~\cite{fung2025embodied,li2025comprehensive,ding2025understanding,kong20253d} reveals three recurring paradigms, three realizations of Eq.~\eqref{eq:p3_s8_6_transition} that differ in what they predict.
\textbf{Representation world models} learn abstract latent spaces that capture the state and dynamics of the environment.
\textbf{Generative world models} explicitly simulate or generate future states of the world through data-driven or hybrid mechanisms.
\textbf{Mental world models} aim to infer and represent the mental states of humans or agents, enabling social reasoning and Theory-of-Mind (ToM) capabilities.
Figure~\ref{fig:world-modeling} contrasts them side by side.


 % \cite{}






\begin{figure*}[!t]
    \centering
    \includegraphics[width=\linewidth]{figures/web/world_modeling_tikz.pdf}
    \caption{Three paradigms of world modeling, all instances of the transition $\hat{S}(t{+}1)=W_\theta(S(t),A(t),O(t))$ of Eq.~\eqref{eq:p3_s8_6_transition}, differing in what they predict: (a) a latent state, (b) future observations, and (c) another agent's mental state and next action.}
    \label{fig:world-modeling}
\end{figure*}




\subsection{Representation World Model}
\label{subsec:p3_s8_6_rep}




\subsubsection{Definition}

A representation world model learns an internal latent description of the environment by compressing multimodal observations into structured, predictive embeddings, prioritizing the extraction of physical and semantic regularities over the generation of photorealistic pixels.
Much as LLMs learn abstract linguistic representations through self-supervised pretraining, representation world models are trained on ``world data'' to infer the principles that govern physical and semantic dynamics; the objective is to predict how the world evolves at an abstract level, not to reconstruct pixels.
The latent state is meant to encode objects and their properties (shape, size, appearance, material), their spatial relations (topology, distance, configuration), the dynamics that govern how the scene evolves, and the causal dependence of outcomes on actions. With such a state, an agent can predict over long horizons, reason, and plan without operating in high-dimensional sensory space.





\subsubsection{Multimodal Perception}

A world model draws on several sensory streams, whose encoding is the subject of Chapter~\ref{chapter:3:sec:6}. Vision dominates: images and especially video carry objects, geometry, interaction, and temporal dynamics, and modern systems rely on high-capacity encoders trained with contrastive vision--language objectives for semantic and geometric priors. Audio adds direct evidence about events and the physical state of the surroundings, which language conveys only symbolically. Touch is indispensable where vision is unreliable, occluded, or unavailable (grasping inside a bag, fine manipulation, tool use), yet it remains underexplored because tactile sensors are limited and tactile data are hard to collect at scale.




\subsubsection{World Representation Modeling}

\paragraph{Vision-oriented Modeling.}
A major line of representation world modeling builds latent spaces from visual observations: video frames are encoded into a compact latent space by self-supervised learning, and a transition model predicts future latent states from past latents and actions, which captures the dynamics of the visual world more efficiently than pixel-space prediction.
Concretely, an encoder $E_\phi$ maps observations to latent states and a predictor $P_\theta$ is trained to minimize a prediction loss defined entirely in latent space,
\begin{equation}
\label{eq:p3_s8_6_rep}
\mathcal{L}_{\text{rep}}(\theta,\phi) = \mathbb{E}\!\left[\,\big\lVert P_\theta\big(E_\phi(O_{\le t}),\, A(t)\big) - \operatorname{sg}\!\big(E_{\bar{\phi}}(O_{t+1})\big)\big\rVert^{2}\,\right],
\end{equation}
where $\operatorname{sg}(\cdot)$ denotes the stop-gradient operator and $E_{\bar{\phi}}$ a slowly updated target encoder.
This objective instantiates the transition of Eq.~\eqref{eq:p3_s8_6_transition} with the predicted state $\hat{S}(t+1)$ taken to be an abstract embedding rather than raw pixels (Figure~\ref{fig:world-modeling}(a)), which is precisely what allows representation world models to capture dynamics efficiently while discarding perceptually salient yet predictively irrelevant detail.



The \emph{World Models} architecture~\cite{ha2018world} decomposes the problem into a VAE-based vision model, an RNN-based memory model, and a lightweight controller; the vision and memory components learn latent dynamics from raw pixels without supervision, and the controller then learns its behavior inside the ``dreamed'' environment, which greatly simplifies policy learning.
The Dreamer agents carry this idea into reinforcement learning at scale: Dreamer~\cite{hafner2019dreamer} learns a latent dynamics model from pixels and trains an actor and critic entirely on imagined latent trajectories, and DreamerV3~\cite{hafner2023mastering} applies a single configuration across a wide range of domains. MuZero~\cite{schrittwieser2020mastering} learns a model that predicts only the quantities planning needs (reward, value, and policy) and plans with it by tree search, mastering Atari, Go, chess, and shogi without being given the rules.


Yann~LeCun's autonomous intelligence framework proposes the Joint Embedding Predictive Architecture (JEPA)~\cite{lecun2022path} as the core mechanism for constructing world models.
JEPA learns latent representations and predictions through self-supervised objectives rather than pixel-level reconstruction, and Eq.~\eqref{eq:p3_s8_6_rep} is its generic form.
I-JEPA~\cite{assran2023jepa} first instantiated it on images, predicting the embeddings of masked target regions from a visible context block, and V-JEPA~\cite{bardes2024vjepa} carried the same feature-prediction objective to video, predicting masked spatio-temporal embeddings \emph{in latent space} with no pixel reconstruction, text, or negative examples, and so avoiding the cost and ambiguity of pixel prediction.
V-JEPA~2~\cite{assran2025vjepa} turns this encoder into a world model. Pretrained on more than a million hours of internet video, it is extended with an action-conditioned predictor, V-JEPA~2-AC, post-trained on less than 62 hours of unlabeled robot video; deployed zero-shot on Franka arms, it picks and places objects by planning actions whose predicted embeddings reach a goal image.

The stop-gradient and slowly updated target encoder in Eq.~\eqref{eq:p3_s8_6_rep} exist only to keep the representation from collapsing, and the most recent JEPA work removes them. LeJEPA~\cite{balestriero2025lejepa} shows that an isotropic Gaussian is the optimal embedding distribution for downstream prediction and enforces it with a single sketched regularizer (SIGReg), leaving one trade-off hyperparameter and no teacher--student machinery. LeWorldModel~\cite{maes2026leworldmodel} uses the same regularizer to train an action-conditioned JEPA world model end to end from raw pixels with only two loss terms; at about 15M parameters it trains on a single GPU, plans up to 48 times faster than world models built on foundation-model encoders, and its latents encode physical quantities well enough to flag physically implausible events.
The principle has also reached language: VL-JEPA~\cite{chen2025vl} predicts continuous embeddings of the target text instead of generating tokens, and LLM-JEPA~\cite{huang2026llmjepa} adds an embedding-prediction objective to the next-token training of language models.


DINO-World~\cite{baldassarre2025back} instead freezes DINOv2 as a universal visual encoder, whose features are already semantic, geometric, and object-centric, learns a latent-space predictor from large-scale unlabeled video, and adds action-conditioned prediction to model how interventions change future visual states.



\paragraph{Hybrid World Modeling.}
Another important direction expands representation learning beyond vision by incorporating additional modalities such as LiDAR, camera trajectories, or multi-view geometric signals.
These hybrid models aim to overcome the limitations of purely visual latent spaces by injecting explicit spatial structure, depth cues, and geometric consistency constraints.
For instance, {MUVO}~\cite{bogdoll2025muvo} is a multimodal generative world model that learns jointly from synchronized camera and LiDAR streams.
By integrating LiDAR-derived depth priors into the latent world representation, MUVO achieves more reliable 3D scene prediction, improved object motion modeling, and long-horizon forecasting in autonomous driving environments.
{UniWorld}~\cite{min2023uniworld} introduces a unified pretrained architecture that merges 3D object detection, motion prediction, and scene representation learning.
Through joint training of BEV features, temporal dynamics, and multisensor context, UniWorld produces a structured latent state that simultaneously supports perception and trajectory forecasting.
{Aether}~\cite{zhu2025aether} proposes a geometry-aware unified world model where camera trajectories are treated as explicit action signals.
This allows the model to learn action-conditioned scene transitions and perform predictive visual planning.
By binding egomotion with geometric consistency constraints, Aether generates spatially coherent future predictions and supports counterfactual simulation within a 3D-aware latent space.




\paragraph{Language-oriented World Modeling.}
A final category examines language as a central medium for latent world representation.
This raises an important question: \emph{can LLMs themselves be viewed as world models?}
From a formal perspective, LLMs indeed share structural similarities with classical world models.
They maintain internal states, treat textual context as observations, accept prompts as actions that alter these states, and generate predictions that correspond to inferred future states.
Models such as Dynalang~\cite{lin2023learning} demonstrate that natural language can serve as a universal intermediate representation for encoding environment states, predicting transitions, and expressing causal relations.
LLMs also exhibit emergent knowledge of commonsense physics, intuitive causality, and human behavior derived from massive text corpora.
However, their limitations remain clear: they are trained on symbolic descriptions rather than direct interaction with the physical world; their internal dynamics reflect linguistic statistics more than physical laws; and their world representations tend to be sparse, fragmented, and disconnected from geometric or temporal continuity.
Thus, while LLMs can be interpreted as language-driven world models, their predictive capacity is primarily symbolic and semantic, rather than grounded in spatial or physical understanding.







\subsection{Generative World Model}
\label{subsec:p3_s8_6_gen}

In contrast to representation-based foundation models, generative world models focus on explicitly predicting future states or synthesizing the evolution of the environment over time.
Rather than only learning compact latent encodings, these models aim to simulate how the world might change under different conditions, capturing temporal dynamics and causal structure.
Whereas the representation paradigm of Eq.~\eqref{eq:p3_s8_6_rep} predicts in an abstract latent space, generative world models learn an explicit conditional distribution over future observations and can be unrolled autoregressively over a horizon $H$,
\begin{equation}
\label{eq:p3_s8_6_gen}
p_\theta\big(O_{t+1:t+H} \mid O_{\le t},\, A_{t:t+H-1}\big) = \prod_{k=0}^{H-1} p_\theta\big(O_{t+k+1} \mid O_{\le t+k},\, A(t+k)\big),
\end{equation}
so that sampling from Eq.~\eqref{eq:p3_s8_6_gen} realizes the transition of Eq.~\eqref{eq:p3_s8_6_transition} directly in observation space (Figure~\ref{fig:world-modeling}(b)).
This explicit synthesis is what allows generative world models to be rendered, inspected, and acted within, at the cost of having to model a high-dimensional output distribution.
Broadly, generative world models can be grouped into two main families:
\emph{physics-driven simulation}, which relies on hand-crafted rules and equations,
and \emph{learning-driven world generation}, which learns world dynamics directly from data.
The former can be viewed as the deterministic limit of the latter: it fixes the transition of Eq.~\eqref{eq:p3_s8_6_gen} to a hand-specified rule set instead of a learned distribution, trading sample diversity and coverage for exact reproducibility and verifiability.
A third route, growing quickly since 2024, sits between the two: a learned model writes the rules as code and a simulator executes them (Section~\ref{subsubsec:p3_s8_6_code}).



\subsubsection{Physics-driven Simulation}

Physics-driven simulation builds interactive virtual worlds from human-designed physical, geometric, and logical rules rather than from learned latent spaces. The environment is governed by predefined physical laws, rendering rules, collision logic, and behavioral scripts, specified top-down as explicit mathematical formulations, so the same initial conditions always produce the same outcome; this determinism, together with interpretability and verifiability, is what makes such systems suitable for safety-critical engineering.

Game engines are the canonical example. Physics engines such as Bullet, Havok, PhysX, and Box2D implement Newtonian, rigid-body, soft-body, cloth, particle, and volumetric dynamics, and the engines around them add acceleration structures (\eg, BVH trees), real-time ray tracing, and navigation meshes for path planning, which gives them a natural foothold in spatial-intelligence tasks. Film production relies on the same principle at higher fidelity: pipelines built on Pixar USD, Maya, or Houdini simulate cloth, hair, smoke, fire, and fluids with differential-equation and fluid solvers (\eg, FLIP/MantaFlow) and physically based path tracing. Engineering, architecture, and scientific visualization extend rule-based simulation to structural mechanics, electromagnetics, illumination, and material response through finite element methods, computational fluid dynamics, and photonic simulation.

For embodied AI, physics-driven simulators such as CARLA~\cite{dosovitskiy2017carla}, AirSim, and Isaac Sim~\cite{NVIDIA2021Isaac} remain the standard platforms for testing and training in autonomous driving, navigation, and manipulation, because they model lighting, collisions, dynamics, and sensor noise closely enough to yield data distributions near those of the real world. The same engines underpin digital twins built on platforms such as NVIDIA Omniverse, Unity, and Unreal Engine, which are discussed in Part~IV, Section~\ref{sec:Digital Twins}.





\subsubsection{Learning-driven World Generation}

In contrast to simulation frameworks based on explicit physical rules, learning-driven world generation seeks to acquire world structure, dynamics, and causal regularities directly from large-scale data, in a bottom-up fashion.
These models do not encode Newtonian equations, lighting models, or geometric constraints by hand; instead, world behavior emerges as statistical patterns during training.
For example, a model may never be given the analytic form of the gravity equation, yet by observing countless videos of falling objects, it can still infer the regularity that objects tend to accelerate downward.
The resulting world models are probabilistic and can exhibit complex, emergent behaviors without explicit programming.


Current learning-driven world generation follows two technical directions:
(1) video-centric generative models that learn dynamic world evolution in pixel or latent space,
and (2) 3D/4D asset generation methods that directly model spatially structured worlds in three or four dimensions.




\paragraph{Video-centric Generative World Models.}
Recent advances in video generation have renewed interest in the idea of a ``world simulator''.
Video models are intrinsically suited to encoding spatiotemporal structure, and when they produce sequences that are not only visually coherent but also consistent with basic physical intuition, they begin to resemble implicit world models.
This became especially prominent with the release of OpenAI Sora~\cite{sora2024}, whose generated videos demonstrate strong temporal consistency, long-horizon stability, and apparent physical plausibility, prompting the community to reconsider whether deep video generation might be converging towards a form of physical world modeling.
Subsequently, Google DeepMind's Genie~3~\cite{genie3} extended video generation from passive ``watchable'' content to interactive experiences: the model updates the generated world in response to user actions, enabling agent-in-the-loop, closed-loop world simulation, complementing the action-conditioned scene generators surveyed in §\ref{subsec:p3_s8_2_2}.
World Labs' RTFM~\cite{worldlabs_rtfm} further highlights real-time performance, achieving high frame-rate and low-latency prediction. Systems such as Wan~\cite{wan2025wan}, Emu3.5~\cite{cui2025emu35}, OpenAI Sora~2~\cite{sora2}, and Veo~3~\cite{veo3} continue this trajectory, gradually pushing video generation from mere visual fidelity towards world coherence, physical consistency, and interactivity.
Action-conditioned diffusion has also turned games into testbeds: GameNGen~\cite{valevski2024diffusion} simulates the game DOOM interactively in real time, and DIAMOND~\cite{alonso2024diffusionworldmodelingvisual} trains reinforcement-learning agents entirely inside a diffusion world model of Atari games, while NVIDIA's Cosmos~\cite{nvidia2025cosmos} offers pretrained video world foundation models intended to be post-trained as simulators for robotics and driving.

Video world models are particularly promising in autonomous driving. GAIA-1~\cite{hu2023gaia} proposes a generative world model tailored for driving scenarios, built around a spatiotemporal Transformer that generates future multi-view scenes based on past observations.
By implicitly capturing traffic patterns, pedestrian behavior, scene geometry, and multi-agent motion, GAIA-1 can support closed-loop evaluation, behavior prediction, and synthetic data generation, thereby reducing the cost and risk of traditional real-world testing.



\paragraph{3D/4D Asset Generation as World Modeling.}
A complementary avenue is to build world models that operate directly on structured three- or four-dimensional representations, such as 3D geometry, occupancy fields, or 4D spatiotemporal volumes.
Unlike video-based approaches that primarily model dynamics in pixel space, these methods explicitly parameterize geometry and temporal evolution, making it easier to represent interactive, spatially consistent worlds.


Genesis~\cite{Genesis} exemplifies this direction. It is not only a generative 3D/4D physics simulation platform, but also integrates a vision-language model-driven ``world-building agent'' that uses a tool API to create dynamic 4D environments (geometry plus dynamics).
These generated worlds can be rendered into multiple modalities (RGB, depth, semantics, contact events, etc.), providing rich supervision for training embodied AI systems.
Genesis thus combines a unified physics engine with generative agents for large-scale world construction.


Autonomous driving again serves as a natural application domain for 3D world models.
DriveArena~\cite{yang2025drivearena} builds a closed-loop generative simulation platform, allowing planning modules to iteratively probe and learn within synthesized 3D driving environments.
Its diffusion-model-based multi-agent scene generation can rapidly produce long, interactive driving episodes.
OccWorld~\cite{zheng2024occworld} and OccSora~\cite{wang2024occsora} adopt 4D occupancy grids as the core representation, predicting future 3D occupancy over time from continuous video and sensor streams, thereby providing a unified motion-and-structure world model.
Copilot4D~\cite{zhang2023copilot4d} tokenizes LiDAR scans and forecasts future LiDAR point clouds with discrete diffusion, learning a world model from unlabeled sensor data alone.


Generating \emph{interactive} 3D worlds brings generative world models closest to their conceptual ideal.
Marble~\cite{Feifei2025Marble} from World Labs combines action-conditioned generation, geometric consistency constraints, and multimodal inputs, so that users can act within the generated world and observe the consequences.
Its next-generation model, Atlas~\cite{worldlabs2026atlas}, is an omni world model pretrained from scratch on text, images, video, and 3D: an autoregressive diffusion transformer that conditions every generated image and depth map on an explicit camera pose and grounds them all in a shared 3D context, so that the same model generates camera-controlled video of up to a minute, reconstructs scenes from a handful of images, and exports its output as point clouds or 3D Gaussian splats. World Labs reports that it outperforms specialized open-source reconstruction models, and that it will power future versions of Marble.


\subsubsection{Code-and-Simulator World Models}
\label{subsubsec:p3_s8_6_code}
Between hand-written simulation and learned generation sits a route that has grown quickly since 2024: a language model \emph{writes} the world model as executable code, and a simulator or engine runs it. The learned component is the author of the rules, not the rules themselves,
\begin{equation}
\label{eq:p3_s8_6_code}
\mathcal{P} \sim p_\theta\big(\mathcal{P} \mid c\big), \qquad \hat{S}(t+1) = \mathrm{Sim}\big(S(t),\, A(t);\, \mathcal{P}\big),
\end{equation}
where $c$ is the conditioning context (a description, observations, or a history of interactions), $\mathcal{P}$ the generated program, and $\mathrm{Sim}$ the interpreter or engine that executes it: plain Python, a physics engine such as MuJoCo or PyChrono, or Blender. Compared with Eq.~\eqref{eq:p3_s8_6_transition}, the learned $W_\theta$ is replaced by a program that can be read, tested, and repaired, so the resulting world keeps the determinism and verifiability of physics-driven simulation while gaining the open-endedness of a generative model. The scene programs of Section~\ref{sec:p3_s8_7} rest on the same idea; the difference here is that $\mathcal{P}$ encodes how the world changes under actions, not only how it is laid out.

In one line of work, code is an agent's transition model. Code World Models~\cite{dainese2024codeworld} synthesize Python implementations of reinforcement-learning environments with a generate, improve, and fix search over programs (GIF-MCTS), evaluated on a benchmark of 18 environments, and plan with them faster and more sample-efficiently than with learned dynamics; WorldCoder~\cite{tang2024worldcoder} builds its world model as a program from its own interactions and transfers knowledge to a new environment by editing that program.

In the other, code builds a simulated 3D world. SceneCraft~\cite{hu2024scenecraft} turns a description into Blender Python, planning with a scene graph, converting spatial relations into numerical layout constraints, and checking renders with a vision-language model. GenSim~\cite{wang2023gensim} and RoboGen~\cite{wang2023robogen} have language models write robotic simulation tasks, scenes, and rewards, making the simulator an open-ended source of training data, and the world-building agent of Genesis~\cite{Genesis} follows the same pattern. Recent systems close the loop with the simulator itself. 
Code2Worlds~\cite{zhang2026code2worlds} generates 4D worlds as Blender code refined by a vision-language motion critic; Code2Worlds-2~\cite{zhang2026code2worlds2} compiles such worlds into physically valid MuJoCo scenes and revises them from penetration, contact, and stability signals; and ChronoAgentic~\cite{wang2026chronoagentic} has cooperating agents write and repair PyChrono simulations, meeting the joint semantic and physical correctness criterion on 82.5\% of PhyWorldBench demos against 52.5\% for the best of ten text-to-video models.

The price is coverage. 
A program can express only what its simulator supports, so photorealistic appearance, deformable and fluid materials, and open-ended visual variety still favor learned generators, while code-and-simulator models win where physical correctness has to be checked.


\subsection{Mental World Modeling}
\label{subsec:p3_s8_6_mental}


Physical world models enable agents to understand and predict how the external environment evolves.
However, intelligent behavior in real-world settings also requires an internal model of \emph{other minds}.
Humans naturally construct such models, \ie, often referred to as \emph{mental models}, to reason about others' beliefs, desires, intentions, emotions, and social context.
This ability, commonly associated with \emph{Theory of Mind (ToM)} \cite{astington1995theory}, underlies human social interaction, collaboration, communication, and decision-making.



Craik's ``small-scale models'' of reality are the origin of this notion too~\cite{nersessian1993theoretician,staggers1993mental}, and Johnson-Laird later argued that human reasoning, problem solving, and causal inference rely on abstract internal representations of situations rather than on replaying detailed physical dynamics.
These representations selectively capture the most relevant structure, \ie, objects, relations, causal dependencies, and possible outcomes, while abstracting away low-level perceptual details.
For example, when cycling, a human does not mentally simulate every joint torque or physical trajectory; instead, one imagines ``whether I would fall or not'' when encountering an obstacle.
This high-level, conceptual simulation is a defining property of mental world modeling.



\subsubsection{From Physical to Cognitive Mental World Models}
Human mental simulation operates by combining current observations with stored knowledge to predict plausible futures.
Perception provides the immediate sensory input, which is then transformed into an internal symbolic and probabilistic representation of the world.
Within this representation, the brain runs rapid, approximate simulations that evaluate different hypothetical actions.
These simulations are not physically precise but are cognitively meaningful: they help determine risks, anticipate others' responses, or evaluate alternative courses of action.
In this sense, mental models serve as a compact, abstract, and counterfactual ``world engine'' for decision making.


In artificial systems, the distinction between physical and mental world models mirrors that in human
cognition.
Physical world models predict state transitions in the environment (\eg, dynamics, geometry, object motion), whereas \emph{mental world models} encode the latent mental states of human users or other agents.
These models allow an AI system to reason about what humans know, believe, want, or intend to do, capabilities that embodied agents, collaborative robots, intelligent tutors, and multi-agent planners all need.



Inspired by the Belief-Desire-Intention (BDI)~\cite{rao1995bdi,georgeff1998belief} framework, a mental world model for AI typically includes a small set of structured variables.
\textbf{Beliefs.} These are representations of what the human knows, perceives, or believes about the world.
\textbf{Desires.} These capture user goals, preferences, motivations, and values.
\textbf{Intentions.} These are planned or ongoing actions that the agent expects the human (or another agent) to follow.
Such decomposition enables an AI system to explicitly reason about human rationality, goal pursuit, and action planning.
Formally, a mental world model can be cast as the inference of a latent mental state $M(t)$, comprising the beliefs, desires, and intentions above, from observed behavior, followed by a prediction of the other agent's next action,
\begin{equation}
\label{eq:p3_s8_6_mental}
\hat{M}(t) = \argmax_{M}\, p\big(M \mid O_{\le t}\big), \qquad \hat{A}^{\text{other}}(t+1) = \pi\big(\hat{M}(t)\big).
\end{equation}
This formulation mirrors the transition of Eq.~\eqref{eq:p3_s8_6_transition} but replaces the physical state with the inferred mental state of another agent (Figure~\ref{fig:world-modeling}(c)), so that the predicted quantity is no longer how the environment will evolve but how a human or peer agent is likely to act.

Machine Theory of Mind is now measured directly. ToMi~\cite{le2019tomi} tests false-belief reasoning in short stories with controls against shortcut solutions, BigToM~\cite{gandhi2023bigtom} generates controlled social scenarios with language models to probe whether models infer beliefs from percepts and actions, and MMToM-QA~\cite{jin2024mmtomqa} poses belief and goal questions about a person acting in a household from video and text together. Large language and multimodal models remain unreliable on these tests, whereas MMToM-QA's BIP-ALM, which performs Bayesian inverse planning over a symbolic state with a language model scoring the likelihoods, does better; it is an explicit instance of the inference in Eq.~\eqref{eq:p3_s8_6_mental}.

\paragraph{Dimensions of Mental State Representation.}
Beyond the belief-desire-intention core, a mental world model for embodied agents must span several further layers of internal state and social context. 
It has to read \emph{affect}, recognizing emotional cues and how they shift a person's risk tolerance and willingness to cooperate. 
It has to represent \emph{social dynamics}: the relationships, roles, group structures, and cultural norms that shape behavior. 
And it has to interpret \emph{communication}, both the verbal channel and the non-verbal one carried by tone, gesture, posture, and facial expression. 
These layers are what turn a bare goal estimate into a socially situated account of why a person acts as they do.



Human spatial reasoning is social as well as geometric, which ties mental modeling to spatial intelligence.
To assist or collaborate with humans in homes, factories, classrooms, or vehicles, an agent must infer the intention behind a spatial action (\eg, where a person intends to move next), anticipate trajectories from goals or habits, adapt navigation or manipulation to user preferences, and interpret spatial configurations that reflect what a person believes or misunderstands.


\subsubsection{A Joint Physical-Mental Formulation}
\label{subsubsec:p3_s8_6_mwm}
The benchmarks above still treat the mental state as something inferred next to a world model, not as part of it: Eq.~\eqref{eq:p3_s8_6_mental} estimates another agent's mind, but the environment it acts in is predicted separately. 
Mental World Modeling (MWM)~\cite{fei2026mental} makes the mental state part of the state that the world model of Eq.~\eqref{eq:p3_s8_6_transition} rolls forward. 
We follow its notation, in which subscripts index time, so that $s_t$ and $a_t^{\epsilon}$ play the roles of $S(t)$ and $A(t)$ above.
The case for a joint state is that neither half suffices alone. For a target agent $\epsilon$ whose true action distribution is $p^{\star}$, a physical-only representation fails whenever
\begin{equation}
\label{eq:p3_s8_6_mwm_property}
p^{\star}\big(a_t^{\epsilon} \mid s_t^{\mathrm{phy}}, s_t^{\mathrm{ment}}\big) \neq p^{\star}\big(a_t^{\epsilon} \mid s_t^{\mathrm{phy}}, \tilde{s}_t^{\mathrm{ment}}\big),
\end{equation}
that is, whenever two scenes that are physically identical but differ in what people believe or want ($s_t^{\mathrm{ment}}$ versus $\tilde{s}_t^{\mathrm{ment}}$) lead to different actions; the symmetric condition, with the mental state held fixed and the physical state changed, rules out a mental-only representation.

\paragraph{Joint State.}
The world state is the pair
\begin{equation}
\label{eq:p3_s8_6_mwm_state}
s_t = \big(s_t^{\mathrm{phy}}, s_t^{\mathrm{ment}}\big) \in \mathcal{S}^{\mathrm{phy}} \times \mathcal{S}^{\mathrm{ment}},
\qquad
s_t^{\mathrm{ment}} = \big(\{m_t^{i}\}_{i\in\mathcal{I}},\, \{m_t^{G}\}_{G\in\mathcal{G}},\, R_t^{\mathrm{ment}},\, \alpha_t\big).
\end{equation}
The physical part collects the objects, the embodied characters, their spatial and contact relations, and the environmental conditions of the scene. The mental part holds an individual state $m_t^{i}$ for every agent $i$ in the set $\mathcal{I}$, optional states $m_t^{G}$ for groups $G\in\mathcal{G}$, the social attitudes and role relations $R_t^{\mathrm{ment}}$ between agents, and the atmosphere $\alpha_t$ of the setting. Each individual state extends the belief-desire-intention core introduced above,
\begin{equation}
\label{eq:p3_s8_6_mwm_individual}
m_t^{i} = \big(\mathrm{id}_t^{i},\, b_t^{i},\, q_t^{i},\, g_t^{i},\, \iota_t^{i},\, e_t^{i},\, d_t^{i},\, n_t^{i},\, c_t^{i}\big),
\end{equation}
with identity attributes, beliefs, attention focus, goals, intentions, emotions, dispositions, norms, and behavioral constraints. 
Beliefs, goals, and intentions are the BDI variables; emotions, norms, and the relations $R_t^{\mathrm{ment}}$ turn the affect and social-dynamics layers discussed above into explicit state variables.

\paragraph{Target-Specific Observation.}
The model reasons from the perspective of one target agent $\epsilon\in\mathcal{I}$, which sees only a rendering of the joint state:
\begin{equation}
\label{eq:p3_s8_6_mwm_obs}
o_t^{\epsilon} = \Omega_{\epsilon}(s_t) = \Big(\Omega_{\epsilon}^{\mathrm{phy}}\big(s_t^{\mathrm{phy}};\, \kappa_t^{\epsilon}\big),\; \Omega_{\epsilon}^{\mathrm{ment}}\big(s_t^{\mathrm{phy}}, s_t^{\mathrm{ment}};\, \rho_t^{\epsilon}\big)\Big),
\end{equation}
where $\kappa_t^{\epsilon}$ is the target's perceptual access (position, line of sight, hearing range, available modalities) and $\rho_t^{\epsilon}$ its social and cognitive perspective (role, prior beliefs, relationships to others, awareness of norms, Theory-of-Mind capacity). The mental part of the observation contains the target's view of its own state and its estimates $\widehat{m}^{\,j,(\ell)}_{t,\epsilon}$ of every other agent $j$, nested up to a depth $L$; at the first level they are sampled as
\begin{equation}
\label{eq:p3_s8_6_mwm_tom}
\widehat{m}^{\,j,(1)}_{t,\epsilon} \sim Q^{(1)}_{\epsilon}\big(m_t^{j} \mid o_t^{\epsilon,\mathrm{phy}},\, m_t^{\epsilon},\, R_t^{\mathrm{ment}},\, \alpha_t\big),
\end{equation}
which is the inference of Eq.~\eqref{eq:p3_s8_6_mental}, made stochastic and conditioned on the target's own mind and social context. Spatial access enters the mental model directly through $\kappa_t^{\epsilon}$: what the target can see from where it stands bounds what it can believe, so a poor decision can follow from an object outside its view as much as from a mistaken belief.

\paragraph{Action and Coupled Transition.}
The target proposes $K$ candidate actions $\mathcal{A}_t^{\epsilon} = \Pi_{\epsilon}^{K}(o_t^{\epsilon}) = \{a^{\epsilon}_{t,1},\dots,a^{\epsilon}_{t,K}\}$. Each action has a physical carrier and a mental content, $a_t^{\epsilon} = (a_t^{\epsilon,\mathrm{phy}}, a_t^{\epsilon,\mathrm{ment}})$ (approaching someone, and meaning to comfort them), and the policy factorizes as $\Pi_{\epsilon}(a_t^{\epsilon,\mathrm{phy}}, a_t^{\epsilon,\mathrm{ment}} \mid o_t^{\epsilon}) = \Pi_{\epsilon}^{\mathrm{phy}}(a_t^{\epsilon,\mathrm{phy}} \mid o_t^{\epsilon})\,\Pi_{\epsilon}^{\mathrm{ment}}(a_t^{\epsilon,\mathrm{ment}} \mid a_t^{\epsilon,\mathrm{phy}}, o_t^{\epsilon})$. The world model then predicts the next joint state with a staged transition,
\begin{equation}
\label{eq:p3_s8_6_mwm_transition}
T_{\theta}\big(s_{t+1} \mid s_t, a_t^{\epsilon}\big) = T_{\theta}^{\mathrm{phy}}\big(s_{t+1}^{\mathrm{phy}} \mid s_t^{\mathrm{phy}}, s_t^{\mathrm{ment}}, a_t^{\epsilon,\mathrm{phy}}\big) \cdot T_{\theta}^{\mathrm{ment}}\big(s_{t+1}^{\mathrm{ment}} \mid s_t^{\mathrm{phy}}, s_t^{\mathrm{ment}}, a_t^{\epsilon,\mathrm{phy}}, a_t^{\epsilon,\mathrm{ment}}\big).
\end{equation}
The asymmetry is deliberate. An intention moves nothing until it is carried out physically, so the physical factor ignores $a_t^{\epsilon,\mathrm{ment}}$; the mental factor depends on both components, because the same reassuring words land differently when the speaker is approaching than when walking away. The physical factor still conditions on the mental state, since what objects, space, and visibility allow is joined by what other agents intend to do. Those other agents are not controlled: their actions $a_t^{-\epsilon}$ are marginalized under the simulator's predictive distribution,
\begin{equation}
\label{eq:p3_s8_6_mwm_marginal}
T_{\theta}\big(s_{t+1} \mid s_t, a_t^{\epsilon}\big) = \int T_{\theta}\big(s_{t+1} \mid s_t, a_t^{\epsilon}, a_t^{-\epsilon}\big)\, p\big(a_t^{-\epsilon} \mid s_t\big)\, \mathrm{d}a_t^{-\epsilon},
\end{equation}
with $p(a_t^{-\epsilon} \mid s_t) = \prod_{j\neq\epsilon} \Pi_j\big(a_t^{j} \mid \Omega_j(s_t)\big)$ for simultaneous moves, so that every agent still acts on its own partial observation. The whole pipeline, $s_t \xrightarrow{\Omega_{\epsilon}} o_t^{\epsilon} \xrightarrow{\Pi_{\epsilon}} a_t^{\epsilon} \xrightarrow{T_{\theta}} s_{t+1}$, is Eq.~\eqref{eq:p3_s8_6_transition} with a mental component added to the state and a perspective added to the observation.

\paragraph{Evaluating Candidate Futures.}
Simulating each candidate gives a branch $\hat{s}^{k}_{t+1}$, and the predicted decision is the branch of highest value,
\begin{equation}
\label{eq:p3_s8_6_mwm_select}
\hat{k} = \argmax_{k\in\{1,\dots,K\}}\, r_k .
\end{equation}
In MENTIS, the training-free and fully inspectable baseline of MWM, the branch value $r_k$ is a weighted combination of normalized scores for physical plausibility, mental consistency, and social appropriateness, and a binary safety and legality veto sets it to zero. Each formal component maps to a module: a state parser recovers $\hat{s}_t$, an observation generator implements $\Omega_{\epsilon}$, an action parser structures the candidates, a world-state transitor implements $T_{\theta}$, and an evaluator assigns $r_k$. On situated decision scenarios presented as text, image, and sounding-video stories, experiments across eight LLM-based world models indicate that the explicit mental component is what makes a person's actual decision predictable~\cite{fei2026mental}.


\subsection{Discussions and Challenges}
\label{subsec:p3_s8_6_discussion}

\paragraph{Representation-based World Models as an Alternative Foundation Model Paradigm.}
Representation world models may emerge as a parallel paradigm to current LLMs, functioning as large-scale foundational models trained through massive self-supervised representation learning to capture physical regularities, causal dynamics, and environmental constraints.
Unlike LLMs, which learn the world indirectly through linguistic abstraction, representation world models aim to model the physical world directly from multimodal observations.
Such models could serve video generation, action prediction in VLA systems, 3D and 4D understanding and generation, and decision-making in autonomous driving.



\paragraph{The Necessity of Multi-source Intelligence Integration.}
World models will likely require the integration of heterogeneous informational sources.
On the one hand, LLMs embody language-driven intelligence, while world models emphasize physically grounded intelligence; rather than being mutually exclusive, a truly capable world model must incorporate both forms of reasoning.
On the other hand, accurate world understanding requires far more than visual perception alone: audio, speech, tactile signals, force feedback, and even proprioception provide essential supplementary cues about the dynamics and structure of the physical environment. Yet the current research community focuses overwhelmingly on visual-only modeling.
Building world models that holistically integrate multimodal sensory information remains a major challenge.


\paragraph{Unifying Physical Modeling and Mental Modeling.}
Another promising direction is the integration of physical and mental modeling within a unified world-modeling framework.
Embodied agents operating in real environments must simultaneously reason about the physical evolution of the environment and the mental or behavioral responses of humans (or other agents).
Recent work in multi-agent systems, social navigation, pedagogical reasoning, human-robot collaboration, and interactive LLM agents emphasizes the growing need for AI systems to build internal hypotheses about both external dynamics and human mental states. 
Mental World Modeling~\cite{fei2026mental} is an early formulation of such a joint state (Section~\ref{subsubsec:p3_s8_6_mwm}).
Ultimately, scalable general-purpose world models must simulate not only \emph{what will happen} but also \emph{what others believe, intend, desire, or misunderstand}.


\paragraph{Balance between Generative Fidelity and Physical Correctness.}
A central challenge for learning-driven world models is reconciling generative expressiveness with physical accuracy.
High-capacity video generators often produce visually stunning but physically implausible outcomes, while models optimized for strict physical consistency may underperform in realism or diversity.
Achieving stable long-horizon predictions without temporal drift or compounding errors remains unsolved.
Designing architectures, training objectives, and constraints that jointly enforce physical laws, spatial coherence, and generative flexibility is a key open problem for the next generation of world simulators.


\paragraph{Evaluation, Benchmarks, and Verifiability of World Models.}
Unlike language models, which can be evaluated through perplexity, benchmarks, or task accuracy, world models lack standardized, reliable, and comprehensive evaluation protocols.
Physical correctness, causal consistency, interactive fidelity, multimodal coherence, and long-horizon stability are difficult to quantify in a unified manner.
Furthermore, world models that aim to simulate open-ended environments raise issues of verifiability, reproducibility, and safety: how can we assess whether a world model reliably generalizes to unseen real-world scenarios, or whether its internal dynamics exhibit hidden biases or unsafe behaviors?
Establishing principled evaluation methodologies and shared benchmarks for both representation-based and generative world models is an urgent research direction.

% \newpage
\section{Programmatic and Symbolic Spatial Generation}
\label{sec:p3_s8_7}

\FramedInitial{P}{rogrammatic} and symbolic spatial generation synthesizes 2D and 3D scene structure through explicit rules, executable programs, and symbolic geometric abstractions.
It generates the structural backbone of a scene (spatial layout, hierarchical organization, geometric constraints, executable logic) rather than its raw appearance, which gives it controllability, compositional generalization, global spatial consistency, and interpretability.

Where the largely neural samplers of §\ref{sec:p3_s8_1}–§\ref{sec:p3_s8_6} map a condition directly to pixels, voxels, or a neural latent, programmatic methods emit an intermediate \emph{program} that a deterministic renderer executes into the final scene.
This intermediate layer is what makes the structure auditable, editable, and compositionally reusable, properties that flat neural outputs lack.
Three paradigm families mark the historical and technical progression of the field: \emph{rule- and grammar-based procedural generation}, \emph{symbolic graphics programs and domain-specific languages}, and \emph{neuro-symbolic language-to-program synthesis} driven by large language models.

\subsection{Task Definition}
\label{subsubsec:p3_s8_7_task}

Let $\mathbf{c}\in\mathcal{C}$ denote a conditioning intent, typically a natural-language description, a layout specification, or a partial constraint set, and let $\kappa:\mathcal{C}\!\to\!\mathcal{C}'$ be its encoder.
A programmatic generator $G_\theta$ emits a symbolic program $\rho$ drawn from a graphics language or domain-specific language $\mathcal{L}$, which a deterministic executor $R$ evaluates into a concrete scene $\hat{\mathbf{s}}$ (a raster image, vector graphic, UI specification, or 3D shape):
\begin{equation}
\label{eq:p3_s8_7_objective}
\rho = G_\theta\big(\kappa(\mathbf{c})\big), \qquad \hat{\mathbf{s}} = R(\rho), \qquad \rho \in \mathcal{L}.
\end{equation}
This formulation inherits the conditional-sampler abstraction of Eq.~\eqref{eq:t2i_generator} but replaces its learned decoder $D$ (Eq.~\eqref{eq:t2i_autoencoder}) over a continuous latent space $\mathcal{Z}$ with an interpretable, discrete program $\rho$ and an \emph{exact, non-learned} renderer $R$.
The neural generators of §\ref{sec:generation} are recovered as the limiting case in which $R$ is itself a neural renderer and $\rho$ a continuous latent, whereas the symbolic case constrains $\rho$ to be a human-readable executable artifact.
Two properties follow directly: errors are localizable to specific program statements, and structural constraints can be enforced or verified deterministically on $\rho$ before rendering.



\subsection{Methodology}
\label{subsubsec:p3_s8_7_methodology}

Figure~\ref{fig:p3_s8_7_paradigms} schematizes the three paradigm families that instantiate Eq.~\eqref{eq:p3_s8_7_objective}, contrasting how each obtains the program $\rho$ before the shared deterministic renderer $R$ executes it.

\begin{figure}[t]
    \centering
    \includegraphics[width=\linewidth]{figures/part3/section8/programmatic_spatial_generation.pdf}
    \caption{Three paradigm families for programmatic spatial generation, all producing a program $\rho$ that a deterministic renderer $R$ executes (Eq.~\eqref{eq:p3_s8_7_objective}): (P1) rule- and grammar-based procedural generation, (P2) symbolic graphics programs and domain-specific languages, and (P3) language-to-program synthesis by an LLM planner.}
    \label{fig:p3_s8_7_paradigms}
\end{figure}

\paragraph{Rule- and Grammar-based Procedural Generation.}
The earliest line of work, procedural content generation (PCG)~\cite{parish2001procedural} in computer graphics and game development, synthesizes scenes such as dungeon maps, terrains, architectural plans, or UI layouts from hand-crafted rules or parameterized procedural programs.
Classical formalisms instantiate $G_\theta$ as a fixed production system rather than a learned model.
The Lindenmayer system (L-system)~\cite{lindenmayer1968mathematical}, originally developed to model biological growth, employs parallel rewriting rules to simulate natural growth processes and fractal patterns, making it ideal for self-similar structures such as vegetation and complex terrain.
Formally, the program evolves by repeated parallel application of a production set $\Pi$ to an axiom $\omega$,
\begin{equation}
\label{eq:p3_s8_7_grammar}
\rho_0 = \omega, \qquad \rho_{t+1} = \Pi(\rho_t), \qquad \hat{\mathbf{s}} = R(\rho_T),
\end{equation}
where the generated symbolic string is not the geometry itself but a sequence of Turtle Graphics commands that the interpreter $R$ parses to draw the final shape in 2D or 3D space.
Shape grammars, introduced by \citet{stiny1980introduction}, provide another influential symbolic mechanism by defining geometric primitives and rewrite rules for composing architectural or engineered shapes.
These early paradigms demonstrated a fundamental insight: spatial structures can be synthesized by executing symbolic generative rules, rather than by directly predicting raw geometry.

\paragraph{Symbolic Graphics Programs and Domain-Specific Languages.}
Symbolic graphics programs (SGPs) encode spatial layouts and geometric configurations as formal programs that execute deterministically into images, vector graphics, UI designs, or 3D shapes, bridging abstract scene intent and concrete visual realization. Common languages such as SVG~\cite{bing2025learning}, TikZ/LaTeX~\cite{ellis2018learning}, and HTML/CSS~\cite{chen2025designcoder} already carry rich symbolic layout specifications in real-world use.
In parallel, researchers have developed domain-specific languages (DSLs) tailored for fine-grained spatial and visual synthesis.
Examples include hierarchical GUI-code generators~\cite{zhu2018automatic} that translate visual structures into programmatic UI specifications, and ShapeAssembly~\cite{jones2020shapeassembly}, a 3D assembly language in which objects are represented as executable programs describing part composition, symmetry, and connectivity.
Such DSLs make the advantages of explicit symbolic control concrete: precise geometry, modular design, verifiable constraints, and systematic compositionality.

\paragraph{Neuro-symbolic and Language-to-Program Synthesis.}
A more recent direction couples symbolic programs with learned models, so that $G_\theta$ becomes a trained synthesizer rather than a fixed rule set. The Scene Language~\cite{zhang2025scenelanguage} unifies neural and symbolic components: a program specifying relational and hierarchical layout structure, natural-language labels summarizing entity semantics, and learned embeddings capturing visual identity.
Executable structure, rather than raw geometry, thus becomes a primary representational layer for scenes.
Large language models~\cite{brown2020language,achiam2023gpt} have accelerated the move toward language-to-program spatial generation: instead of predicting images or layouts directly, LLM-based systems act as symbolic planners that synthesize the program by approximately maximizing its conditional likelihood,
\begin{equation}
\label{eq:p3_s8_7_synthesis}
\rho^{\star} = \argmax_{\rho \in \mathcal{L}} \, p_\theta\big(\rho \mid \kappa(\mathbf{c})\big),
\end{equation}
mapping natural-language descriptions to executable layout programs~\cite{aguina2024open, hsu2025programs, yang2024holodeck}.
These programs then guide downstream renderers to produce UI designs, graphic layouts, architectural plans, or 3D scenes with precise structural fidelity, and users can edit or refine the underlying program with fine-grained control.

The three families differ along two cross-cutting axes.
The first is the \emph{representation} of $\rho$: a parallel string-rewriting system (grammar-based generation), an executable program in a general-purpose or domain-specific graphics language (SGPs and DSLs), or a hybrid structure that interleaves code, language, and learned embeddings (neuro-symbolic synthesis).
The second is the \emph{inference procedure} that produces $\rho$: hand-authored production rules, human-written or template-instantiated programs, or learned program synthesis and LLM-based planning. Across all three, the target modality (2D raster, vector graphics, UI, or 3D) is determined solely by the renderer $R$, which keeps the symbolic and rendering concerns cleanly separated.



\subsection{Benchmarks}
\label{subsubsec:p3_s8_7_benchmarks}

Evaluation in this area differs from the distributional metrics of §\ref{sec:generation} because the immediate output is a program, not an image, so quality is assessed along several complementary axes.
\textbf{Program validity} measures whether the emitted $\rho$ parses and executes without error in its target language, a prerequisite that purely neural generators do not face.
\textbf{Rendered-output fidelity} then compares the executed scene $R(\rho)$ against a reference, using pixel- or perceptual-level similarity for raster and vector targets and geometric agreement for 3D shapes.
\textbf{Constraint satisfaction} checks whether explicitly specified spatial relations, counts, alignments, or symmetries hold in the rendered result, a property that symbolic methods are designed to guarantee.
\textbf{Controllability and editability} assess how reliably localized edits to $\rho$ produce the intended localized changes in the scene without unintended global drift. Concretely, these axes are exercised across vector-graphics and SVG synthesis~\cite{bing2025learning}, UI and GUI code generation~\cite{zhu2018automatic, chen2025designcoder}, 3D program induction over part-structured shapes~\cite{jones2020shapeassembly}, evaluated against part-hierarchy benchmarks such as PartNet~\cite{mo2019partnet}, and full scene-program synthesis~\cite{zhang2025scenelanguage, yang2024holodeck}, where the executable program serves simultaneously as the prediction and as the verifiable ground for measurement.

\subsection{Challenges and Trends}
\label{subsubsec:p3_s8_7_challenges}

The field has moved from entirely hand-authored rules toward neuro-symbolic integration with large models, driven by the demand for generators that are more controllable and interpretable than end-to-end neural methods.

\textbf{Generalization.} Automatically synthesized programs still struggle with novel constraints and complex semantic specifications~\cite{mao2025neurosymbolicconcepts}. Program synthesis has to infer the structural logic of a scene rather than memorize surface patterns if the generator is to extrapolate to compositions and constraint combinations absent from training.

\textbf{Efficiency and differentiability.} Neuro-symbolic programs that are both efficient and fully differentiable~\cite{guo2023efficientsymbolicpolicy} would let a rendering-level loss propagate back through the program, so that the whole pipeline, from symbolic planning to neural rendering, could be trained and fine-tuned end to end.

\textbf{Embodied interaction.} Procedurally synthesized interactive environments already supply large-scale embodied worlds~\cite{kolve2017ai2thor, procthor, wang2024pipeline}, and the next step is for agents to generate or validate symbolic scene structure through active exploration rather than from static datasets alone, yielding symbolic representations of physical spaces that can be checked against what the agent observes.



\section{Summary and Outlook}
\label{sec:p3_s8_8}

This chapter has surveyed spatial reconstruction and generation across seven subfields that, taken together, span an axis of increasing output structure and dimensionality: 2D images (§\ref{sec:p3_s8_1}), 3D objects and scenes (§\ref{sec:p3_s8_2}), video (§\ref{sec:p3_s8_3}), embodied motion (§\ref{sec:p3_s8_4}), 4D space-time content (§\ref{sec:p3_s8_5}), world models (§\ref{sec:p3_s8_6}), and interpretable symbolic programs (§\ref{sec:p3_s8_7}). Despite their different output domains, the subfields are unified by a single conditional-generation backbone, established in §\ref{subsec:p3_s8_1_1}: an encoder $\kappa$ maps a condition $\mathbf{c}$ into a latent code, a generator $G_\theta$ samples structure from that code, and a decoder or renderer turns the sampled structure into a concrete artifact (Eq.~\eqref{eq:t2i_autoencoder}, Eq.~\eqref{eq:t2i_generator}). What distinguishes the subfields is not this backbone but the \emph{output space} and the \emph{decoder} that closes it: a pixel decoder for images, a volumetric or splatting renderer for 3D, a spatio-temporal decoder for video, a kinematic sampler for motion, a time-indexed 3D representation for 4D, a state-transition predictor for world models, and an exact, non-learned program executor for symbolic generation. Table~\ref{tab:p3_s8_chapter_comparison} consolidates this view.

\begin{table}[t]
\centering
\caption{The seven spatial-generation subfields of this chapter, all instances of the conditional latent-generator backbone of §\ref{subsec:p3_s8_1_1}.}
\label{tab:p3_s8_chapter_comparison}
\fontsize{8}{10}\selectfont
\setlength{\tabcolsep}{1.0mm}
\begin{tabular}{L{0.17\textwidth} L{0.20\textwidth} L{0.33\textwidth} L{0.22\textwidth}}
\toprule
\textbf{Subfield} & \textbf{Output space} & \textbf{Dominant paradigms} & \textbf{Decoder / renderer} \\
\midrule
§\ref{sec:p3_s8_1} 2D image & Pixel grid (raster) & GAN; autoregressive; latent diffusion; flow matching & Learned pixel decoder \\
§\ref{sec:p3_s8_2} 3D object/scene & Mesh / NeRF / 3DGS / latent & Native 3D diffusion; score distillation; multi-view lift; feed-forward reconstruction & Differentiable / volumetric renderer \\
§\ref{sec:p3_s8_3} Video & Frame sequence (clip latent) & Inflated T2I backbones; spatial / temporal control; editing & Spatio-temporal decoder \\
§\ref{sec:p3_s8_4} Motion & Pose / trajectory sequence & Regression; diffusion samplers; editing operators & Kinematic / skeletal decoder \\
§\ref{sec:p3_s8_5} 4D & Time-indexed 3D representation & Score distillation; generate-then-reconstruct & Deformable / per-frame renderer \\
§\ref{sec:p3_s8_6} World model & Latent state / future observations & Representation; generative; mental world models & Latent transition / observation predictor \\
§\ref{sec:p3_s8_7} Symbolic program & Executable program in a language $\mathcal{L}$ & Procedural grammars; SGP/DSL; LLM program synthesis & Deterministic, non-learned executor \\
\bottomrule
\end{tabular}
\end{table}

\textbf{A shared backbone, specialized by output space.} The progression from §\ref{sec:p3_s8_1} to §\ref{sec:p3_s8_7} is best read not as seven disjoint problems but as one generative abstraction instantiated against progressively richer codomains, and each section casts its own subtasks as conditional specializations of that common sampler. The symbolic case (§\ref{sec:p3_s8_7}) is the limiting instance in which the decoder degenerates to an exact program executor, trading the expressiveness of a learned renderer for interpretability and verifiable control.

\textbf{The 2D prior as a reusable engine.} The image generators of §\ref{sec:p3_s8_1} recur throughout the chapter not as terminal systems but as priors to be distilled or lifted: 3D synthesis distills frozen 2D diffusion through score distillation and multi-view consistency (§\ref{sec:p3_s8_2}), video models inflate text-to-image backbones along the temporal axis (§\ref{sec:p3_s8_3}), and 4D pipelines compose both (§\ref{sec:p3_s8_5}). A single advance in the 2D substrate therefore propagates outward across every higher-dimensional subfield, which concentrates research leverage but also propagates the substrate's failure modes.

\textbf{Physical plausibility and fragmented evaluation.} Across subfields, the recurring failure mode is not low visual fidelity but violated physics and geometry, manifest as multi-view inconsistency in 3D, temporal drift and flicker in video, foot-skating and interpenetration in motion, and dynamic implausibility in 4D. Evaluation remains correspondingly fragmented: each subfield has evolved its own metric tuples, and all are migrating toward learned and vision-language-model judges as scalar similarity scores prove insufficient for spatially structured output. A unified, physically grounded evaluation protocol that spans these subfields remains an open goal.

\textbf{Outlook.} These observations point toward convergence rather than further fragmentation. World modeling (§\ref{sec:p3_s8_6}) already acts as a conceptual capstone that absorbs reconstruction, generation, and dynamics into a single predictive substrate, while programmatic and symbolic generation (§\ref{sec:p3_s8_7}) supplies the controllability and interpretability that purely neural pipelines lack; their integration, neural generators steered and verified by symbolic structure, is a natural trajectory for the field. The capabilities surveyed here also form the generative substrate for the applications of Part~\ref{part:applications}, where embodied agents, digital twins, and interactive environments consume generated spatial content rather than producing it for its own sake.

% \input{part_3/section_4/main}



\part{Applications}
\label{part:applications}
% ================================
% \section*{\centering \bf Part IV: Applications}
% \addcontentsline{toc}{section}{Part IV: Applications}
% ================================




\begin{figure}
    \centering
    % \includegraphics[width=\linewidth, height=12cm]{example-image-golden}
    % \includegraphics[width=0.99\linewidth]{figures/part4.png}
    \includegraphics[width=\linewidth]{figures/part4_opener_v2.pdf}
    
    \caption{Part IV at a glance. (1) Leave the lab: deployment tests spatial intelligence under latency, safety, sensor failure, and the cost of error. (2) Domain by domain: embodied agents, extended reality, geo-spatial and scientific settings, assistants, and further fields, each stressing a different axis. (3) A horizontal foundation: one shared layer beneath science, industry, and daily life, with gaps found in the field feeding back into it.}
    \label{fig:p4_main}
\end{figure}


The preceding parts built spatial intelligence as a set of capabilities; this part asks what those capabilities are worth once they leave the laboratory. Applications are not a victory lap over Parts II and III but a stress test: a method that reconstructs a scene cleanly on a benchmark must, in the field, do so under a latency budget, a safety constraint, a sensor that fails, or a cost of error measured in more than accuracy. What a domain demands of spatial intelligence, and what it reveals about the gaps that remain, is visible only at deployment.

This part is therefore organized by domain rather than by capability, because the interesting content lives in the delta each domain adds over a general treatment. The tour runs from embodied agents that must act in the physical world, through the extended-reality systems that render it back to a user in real time, to the geo-spatial and scientific settings that model it at the scale of planets and of cells, to the assistants that inhabit a person's own environment, and out into a long tail of further fields. Each stresses spatial intelligence along a different axis: real-time action, on-device perception, metric fidelity, human-centered interaction. Read together, they make the case that spatial intelligence is not one more application but a horizontal foundation on which applications across science, industry, and daily life are increasingly built.




\chapter{Spatial Intelligence for Embodied Systems }
\label{sec:embodied}

\FramedInitial{E}{mbodied} systems are the arena in which spatial intelligence stops being a matter of analysis and becomes a matter of consequence. An embodied agent does not merely describe a scene; it occupies one, and every perception it forms and every inference it draws is ultimately answerable to an action it must take in the physical world, under real time budgets, sensor noise, and the cost of being wrong. This chapter turns from the capabilities of the previous part to their deployment, examining how perception, reasoning, and generation are brought to bear inside agents that move, manipulate, and assist.

What makes this more than a restatement is that embodiment changes the problem. The same spatial competence looks different in a manipulator closing a grasp, a car threading traffic, a wearable inferring its user's intent, a drone holding position in wind, or a fleet of robots coordinating on a warehouse floor: each fixes its own sensors, latencies, safety margins, and notion of success. The chapter is organized by these domains rather than by capability, because the interesting content lives in the domain-specific delta, what each setting demands of spatial intelligence that a general treatment does not reveal.


\section{Robotic Manipulation}
\label{sec:p4_s9_1}



 % \RespondAuthor{ }
\subsection{Introduction}
\label{subsec:p4-sec9-1-introduction}

Robotic manipulation has advanced rapidly, driven by strides in perception, learning, and control.
Yet a significant gap remains between competent laboratory systems and robust, general-purpose robots capable of operating in unstructured settings.
A central bottleneck is that manipulation is fundamentally spatial: every decision (where to reach, how to establish contact, how to route a trajectory through clutter, or how to use a tool) hinges on perceiving, representing, and reasoning about three-dimensional structure under uncertainty.
Manipulation is thus a direct test of the \emph{spatial intelligence} developed throughout Parts~I--III: it demands that systems go beyond semantic recognition to build actionable 3D world models, ground tasks in geometric constraints, predict contact outcomes, and plan behaviors that respect the physical laws of the environment.
The perception, reasoning, and generation capabilities that such systems presuppose are surveyed in Chapter~\ref{chapter:3:sec:6}, Chapter~\ref{sec:reasoning}, and Chapter~\ref{sec:generation}; this section focuses not on re-deriving those capabilities but on how they are composed, grounded, and constrained for physical manipulation.

\subsection{System Overview}
\label{subsec:p4-sec9-1-system-overview}
A manipulation system maps raw multimodal observations and a task instruction into executable, constraint-satisfying actions, and this section organizes that mapping into four strategic \textbf{pivots}, each a fundamental shift in how systems bridge perception and physical action.
\textbf{(i) Grounding Action in Spatial Anchors} moves beyond raw pixels to vision priors and sparse spatial tokens (keypoints, dense correspondences) that stay consistent across viewpoints and instances;
\textbf{(ii) Inferring Functional Action Possibilities} shifts from object classification to extracting affordances and tool-use capabilities, bridging perception with contact-rich execution;
\textbf{(iii) Encoding Spatial Reasoning in Generalist Models} transitions from single-task policies to foundation models (VLMs/VLAs) that integrate explicit 3D representations and reasoning traces for cross-embodiment generalization; and
\textbf{(iv) Satisfying Geometric and Logical Constraints} advances from simple trajectory generation to constrained optimization and neuro-symbolic planning (TAMP), ensuring actions respect strict safety limits and long-horizon logical structure.
The four pivots are complementary and mutually reinforcing rather than strictly sequential: anchors and affordances structure the perceptual input, generalist models supply broad priors, and constraint-based planners check feasibility under the modeled geometric, physical, and logical assumptions. Figure~\ref{fig:p4_s9_1_overview} situates them as complementary capabilities within a closed-loop perception-to-action system, and the remainder of this section treats them in this order.

\begin{figure}[t]
    \centering
    \includegraphics[width=\linewidth]{figures/part4/section9/manipulation_overview.pdf}
    \caption{Spatial intelligence for robotic manipulation, illustrated by placing a mug in an open drawer: four complementary spatial capabilities link multimodal inputs to 6-DoF actions, and execution returns new observations for replanning. Evaluation bars and depth/point-cloud graphics are schematic, not measured results.}
    \label{fig:p4_s9_1_overview}
\end{figure}


\subsection{Pivotal Task-1: Grounding Action in Spatial Anchors}
\label{subsec:p4-sec9-1-pivotal-task-1-grounding-action-in-spatial-anchors}

This pivot reframes visuomotor learning from operating over raw pixels to reasoning through explicit \textbf{spatial anchors} (sparse semantic keypoints, dense correspondences, point tracks, or affordance-based contact points) that compress high-dimensional observations into spatial tokens tied to object- or task-centric reference frames and designed to remain comparatively stable across viewpoint, instance, and clutter. As Figure~\ref{fig:p4_s9_1_anchors} indicates, the approach follows a three-step logic: large-scale vision priors first \emph{produce} the anchors, policies are then \emph{conditioned} on them, and the payoff is read off as a \emph{generalization gap} between seen and unseen instances.

\begin{figure}
    \centering
    \includegraphics[width=0.98\linewidth]{figures/part4/section9/robotic_anchors.pdf}
    \caption{Grounding action in spatial anchors: (1) anchors $P_t$ produced or transferred from vision priors across views, (2) an anchor-conditioned policy, and (3) generalization measured as the seen--unseen success gap $\Delta_{\mathrm{gen}}$. Scenes, correspondences, and score bars are schematic, not measured results.}
    \label{fig:p4_s9_1_anchors}
\end{figure}

The dominant methodology leverages large-scale vision priors to generate, propagate, or retrieve these anchors, then trains standard visuomotor learners on the resulting structured inputs. A common strategy prescribes anchor points in a few reference frames and propagates them across large datasets via off-the-shelf correspondence models: P3-PO~\cite{levy2024p3poprescriptivepointpriors} propagates manually specified keypoints using dense features from DINO~\cite{dino}, DINOv2~\cite{dinov2}, or diffusion-based representations~\cite{tang2023emergentcorrespondenceimagediffusion,zhao2025high,zhang2023talefeaturesstablediffusion} (these visual representations are surveyed in Chapter~\ref{chapter:single_modality}). Relatedly, Robo-ABC~\cite{ju2024roboabcaffordancegeneralizationcategories} retrieves affordance contact points from an external memory and transfers them to novel scenes via dense semantic correspondence, while Track2Act~\cite{bharadhwaj2024track2act} learns long-term point tracks from Internet-scale video and uses them to infer object-centric transformations that serve as open-loop manipulation plans. Once anchors are obtained, policy learning conditions action prediction directly on them, through keypoint-conditioned behavior cloning (KALM~\cite{fang2024keypointabstractionusinglarge}), diffusion policies in object-relative frames (AffordDP~\cite{wu2025afforddpgeneralizablediffusionpolicy}), and shared human--robot 3D keypoint spaces (Point Policy~\cite{haldar2025pointpolicyunifyingobservations}). Across these variants, anchoring actions to spatial structure reduces perceptual ambiguity and improves robustness under distribution shift.

Because the central claim is out-of-distribution generalization, evaluation pairs high-fidelity simulation with real-robot deployment rather than relying on either alone. In simulation, affordance-centric methods such as AffordDP~\cite{wu2025afforddpgeneralizablediffusionpolicy} run on IsaacGym~\cite{makoviychuk2021isaac} over articulated objects from GAPartnet~\cite{geng2023gapartnet} whose functional parts (handles, knobs, levers) vary widely in geometry, while imitation suites such as Robomimic~\cite{robomimic2021} and RoboTwin~\cite{mu2024robotwin} isolate trajectory tracking and dual-arm coordination, with P3-PO~\cite{levy2024p3poprescriptivepointpriors} using them for controlled spatial-anchoring ablations. Real-world suites then stress in-the-wild generalization on articulated (drawers, microwaves), deformable (towel folding), and tool-use (sweeping) tasks, as in Point Policy~\cite{haldar2025pointpolicyunifyingobservations} and Track2Act~\cite{bharadhwaj2024track2act}, and a zero-shot cross-category protocol (train on one knife, test on spatulas or novel knives) is formalized by Robo-ABC~\cite{ju2024roboabcaffordancegeneralizationcategories} to measure how well retrieved anchors transfer without extra supervision. Performance is reported as Task Success Rate and, to isolate the contribution of anchoring, a Generalization Gap between seen and unseen instances. Spatially grounded policies outperform pixel baselines under shift in these studies: P3-PO reports large gains on cluttered, novel-object scenes~\cite{levy2024p3poprescriptivepointpriors}, and Robo-ABC improves both affordance-retrieval accuracy and real-robot success without heavy robot-data collection, binding demonstrations to reusable spatial structure more effectively than classic one-shot imitation~\cite{duan2017oneshotimitationlearning}.


\subsection{Pivotal Task-2: Inferring Functional Action Possibilities}
\label{subsec:p4-sec9-1-pivotal-task-2-inferring-functional-action-possibilities}

This pivot shifts from general object recognition to extracting \textbf{object affordances}, \ie, the specific action possibilities implied by an object's functional properties, and uses them as the intermediate representation that bridges raw perception and contact-rich execution. Whereas Pivotal Task-1 treats contact points as viewpoint-stable \emph{anchors} for conditioning a policy, affordances here are studied as \emph{functional possibilities} that specify what action an object supports. As Figure~\ref{fig:p4_s9_1_affordances} shows, the thread runs from how affordances are \emph{represented}, through the \emph{actions} they drive (grasping, placement, tool use), to \emph{outcome-oriented} benchmarks that score achieved function rather than imitated trajectory.

\begin{figure}
    \centering
    \includegraphics[width=0.98\linewidth]{figures/part4/section9/robotic_affordances.pdf}
    \caption{Affordances as functional action possibilities: a task- and context-conditioned affordance representation (left) supports three parallel task examples (6-DoF grasp, stable placement, hook-and-pull tool use), each judged by its functional outcome. Green checks mark intended outcomes in these illustrations, not measured results.}
    \label{fig:p4_s9_1_affordances}
\end{figure}

To capture functional properties beyond whole-object categorization, prior work builds object-centric representations (part-level segmentation~\cite{do2018affordancenet}, dense visual descriptors~\cite{florencemanuelli2018dense}, and keypoint-based encodings~\cite{manuelli2019kpam,qin2020keto,mo2021where2act}) that directly support action synthesis, and more recent vision--language models predict spatial affordance points directly from open-vocabulary prompts (RoboPoint~\cite{yuan2024robopoint}, MOKA~\cite{liu2024moka}). These representations drive learning-based 6-DoF grasping~\cite{sundermeyer2021contact,zhao2025high,jiang2021synergies,zhao2024efficient,yuan2023m2t2,deshpande2025graspmolmo} and stable object placement~\cite{zeng2020transporter,liu2022structformer,zhao2025towards,yuan2023m2t2,zhao2025anyplace} by grounding generation in functional structure rather than raw geometry, and a specialized strand further models tool--behavior relationships, either learning tool--action--effect pairs~\cite{stoytchev2008learning} or predicting specific affordances such as pulling~\cite{tikhanoff2013exploring}, so as to infer end-effector poses that achieve the intended functional outcome.

Benchmarks for affordance methods emphasize understanding object functionality over geometric matching: the property under test is \emph{functional} generalization, not trajectory imitation. For grasping, YCB-based real-world evaluations~\cite{Calli2015YCB}, GraspNet-1Billion~\cite{Fang2020GraspNet}, and simulation suites such as Dex-Net~\cite{Mahler2017DexNet} score 6-DoF stability and contact validity across novel instances and viewpoints; placement benchmarks such as RLBench~\cite{james2020rlbench} and ManiSkill~\cite{mu2021maniskill} test stable stacking and insertion in cluttered, articulated scenes, and M2T2~\cite{yuan2023m2t2} explicitly evaluates valid placement poses conditioned on object geometry and scene context. For articulated objects and tools, Where2Act~\cite{mo2021where2act} and PartNet-Mobility~\cite{xiang2020sapien} score the prediction of actionable parts and interaction parameters for pulling, pushing, and rotating, while tool-use tasks judge whether the inferred end-effector pose produces the intended causal effect (drawing, hooking, pulling) rather than reproducing a specific trajectory. Evaluation therefore centers on outcome-oriented success: stable grasps that lift and hold unseen objects under pose, scale, and appearance variation; placement success, pose accuracy, and post-release stability under clutter and occlusion; and, for articulated and tool interactions, whether the predicted parameters yield the desired functional result.



\subsection{Pivotal Task-3: Encoding Spatial Reasoning in Generalist Models}
\label{subsec:p4-sec9-1-pivotal-task-3-encoding-spatial-reasoning-in-generalist-models}

This pivot moves from specialized, single-task policies toward \textbf{generalist robotic foundation models}. Existing Vision--Language--Action (VLA) models such as RT-2~\cite{zitkovich2023rt2} and OpenVLA~\cite{kim2024openvla} show strong semantic understanding but often lack the spatial common sense needed for precise manipulation or robust generalization to novel viewpoints, so the central challenge is to strengthen spatial grounding during pre-training, enabling explicit reasoning about 3D geometry and generalization across environments and embodiments~\cite{chen2024spatialvlm}. The general foundation-model and spatial-reasoning machinery these policies inherit is treated in Chapter~\ref{sec:reasoning}; here we examine only its manipulation-specific adaptation, which Figure~\ref{fig:p4_s9_1_vla} groups into three families: explicit 3D representations, view-invariant representations, and grounded reasoning traces.

\begin{figure}
    \centering
    \includegraphics[width=0.98\linewidth]{figures/part4/section9/robotic_vla.pdf}
    \caption{Spatial grounding mechanisms for generalist robot policies: (1) explicit 3D features fused with task tokens (3D-LLM, 3D-VLA), (2) view canonicalization (OG-VLA) and spatial action discretization (SpatialVLA), two separate methods rather than stages of one, and (3) grounded reasoning traces that condition action prediction. The bottom row shows evaluation of viewpoint robustness and cross-robot transfer; all maps and overlays are schematic.}
    \label{fig:p4_s9_1_vla}
\end{figure}

The first family builds explicit 3D representations: models such as 3D-LLM~\cite{hong20233d} and 3D-VLA~\cite{zhen20243d} bridge point clouds and large language models, with 3D-VLA adding a 3D generative pre-training objective that lets the model reason about action outcomes in a consistent geometric space, effectively ``imagining'' future states. The second groups view canonicalization and spatial action encoding as distinct mechanisms: OG-VLA~\cite{singh2025og} projects multi-view observations into a point cloud and re-renders canonical orthographic views to reduce sensitivity to camera placement, while SpatialVLA~\cite{qu2025spatialvla} uses Ego3D position encoding for observations and \textbf{Adaptive Action Grids} to discretize spatial movement actions; adapting to new setups can involve re-discretizing the grids. The third injects spatial reasoning through grounded reasoning traces: ``thinking'' models such as Emma-X~\cite{sun2024emma}, MolmoAct series~\cite{lee2025molmoact, fang2026molmoact2}, closely related to TraceVLA~\cite{zheng2024tracevla}, replace direct pixel-to-action regression with intermediate, grounded outputs (2D motion trajectories or relative depth maps) that explicitly condition the final action prediction.

These approaches are evaluated under method-specific zero-shot and adaptation protocols across embodiments, viewpoints, and environments, with key metrics including multimodal goal-generation quality, robustness to extrinsic camera perturbations, and skill transfer across robot morphologies such as from a fixed-base Franka arm to a mobile manipulator~\cite{team2024octo}. The reported gains point the same way: canonical orthographic views keep grounding consistent under viewpoint change (OG-VLA), spatially aligned action grids share 3D structure across setups through discretized movement (SpatialVLA), and intermediate visual or depth-aware reasoning traces act as a regularizer on action prediction (Emma-X, MolmoAct).

\subsection{Pivotal Task-4: Satisfying Geometric and Logical Constraints}
\label{subsec:p4-sec9-1-pivotal-task-4-satisfying-geometric-and-logical-constraints}

This pivot moves planning from heuristic trajectory generation toward formal \textbf{constrained optimization} that enforces geometric feasibility, fundamental physical laws (non-penetration, friction cones, force limits), and high-level logical task structure. By casting manipulation as multi-objective optimization, these methods seek to satisfy modeled hard constraints while optimizing soft task objectives, spanning both high-frequency control and long-horizon decision-making. Figure~\ref{fig:p4_s9_1_constraints} organizes the thread along three fronts (symbolic--geometric planning, contact and trajectory optimization, and constraint-aware generative planning) before synthesizing the roles constraints play in learning-based systems.

\begin{figure}
    \centering
    \includegraphics[width=0.98\linewidth]{figures/part4/section9/robotic_constraints.pdf}
    \caption{Geometric and logical constraints in manipulation: (1) symbolic--geometric planning with revision after infeasibility, (2) contact and trajectory optimization (inset: normal force $f_n$ into the object, tangential force $f_t$ along its surface), and (3) constraint-aware generative planning. The bottom row shows three roles constraints play for learned policies.}
    \label{fig:p4_s9_1_constraints}
\end{figure}

Methodological progress has moved from purely symbolic planners toward neuro-symbolic, learning-augmented optimization frameworks that unify discrete logic with continuous physics. Classical Task and Motion Planning (TAMP) and Logic-Geometric Programming (LGP)~\cite{toussaint2015lgp} provide the foundation by jointly reasoning over symbolic task structure and geometric feasibility, and recent systems augment them with Large Language Models and Vision-Language Models for scalability in open-ended domains; for example, VIZ-COAST~\cite{yan2025vizcoast} uses the spatial reasoning of VLMs to predict refinement failures \emph{before} expensive geometric optimization, substantially reducing search complexity on long-horizon problems. At the level of physical interaction, contact-implicit optimization discovers contact schedules without explicit mode enumeration but is traditionally too slow for real-time control; GPU-parallelized systems such as SPaSM~\cite{chen2025spasm} compile constraint evaluation and gradient computation into CUDA kernels, accelerating joint optimization of object placements and robot trajectories. Beyond deterministic constraints, learning-based and generative models capture uncertainty and complex structure: differentiable physics engines~\cite{freeman2021brax} let gradients flow through physical interactions, stochastic implicit neural signed distance functions~\cite{quintero2024stochastic} model geometric uncertainty directly within the planning manifold, and diffusion-based planners such as JM2D~\cite{jung2025joint} couple model-free diffusion with model-based constraint optimization during sampling, rather than relying only on a sequential plan-then-filter interface.

Such systems are benchmarked on multi-stage re-grasping under uncertainty, forceful tool use, and tight-tolerance sequential interaction, since what they promise is reliability over long, contact-rich horizons; suites such as the Functional Manipulation Benchmark (FMB)~\cite{fmb2023} scoring not only task completion but constraint satisfaction under partial observability, external disturbance, and compounding execution error. In these systems constrained optimization plays three roles: it acts as a \textit{teacher}, generating physically valid trajectories that supervise imitation and offline reinforcement learning; as a \textit{safety shield}, where model-based projections map unconstrained neural or diffusion policy outputs onto constraint-satisfying sets at execution time~\cite{römer2025diffusionpredictivecontrolconstraints}; and as a \textit{structural prior} that prunes physically infeasible regions of the action space and curbs compounding error in long-horizon tasks. Together these roles keep learned and generative planners grounded in geometric and physical reality, reducing physically implausible outputs and supporting interpretable, feasible plans under the modeled assumptions; real-world execution still depends on state estimation, model accuracy, solver convergence, and disturbance handling.

Table~\ref{tab:p4_s9_1_benchmarks} consolidates the principal benchmarks, representative methods, and headline evaluation focus across the four pivots, situating the spatially grounded methods surveyed above within a common evaluation landscape.

\begin{table}[t]
\centering
\fontsize{8}{10}\selectfont
\setlength{\tabcolsep}{1.0mm}
\caption{Representative benchmarks, methods, and headline metrics for the four manipulation pivots, drawn from simulation and real-robot settings; metrics are reported on out-of-distribution object instances and configurations where the underlying work permits.}
\label{tab:p4_s9_1_benchmarks}
\begin{tabular}{L{0.15\textwidth} L{0.32\textwidth} L{0.25\textwidth} L{0.19\textwidth}}
\toprule
\textbf{Pivot} & \textbf{Representative Benchmark(s)} & \textbf{Representative Method(s)} & \textbf{Headline Metric} \\
\midrule
P1: Spatial anchors & IsaacGym~\cite{makoviychuk2021isaac} with GAPartNet~\cite{geng2023gapartnet}, Robomimic~\cite{robomimic2021}, RoboTwin~\cite{mu2024robotwin} + real-world & P3-PO~\cite{levy2024p3poprescriptivepointpriors}, Point Policy~\cite{haldar2025pointpolicyunifyingobservations} & Task SR / Generalization Gap \\
P2: Affordances & GraspNet-1B~\cite{Fang2020GraspNet}, Dex-Net~\cite{Mahler2017DexNet}, Where2Act~\cite{mo2021where2act}, PartNet-Mobility~\cite{xiang2020sapien} & M2T2~\cite{yuan2023m2t2} & Grasp / Placement SR \\
P3: Generalist VLA & ManiSkill~\cite{mu2021maniskill}, Octo eval~\cite{team2024octo} & SpatialVLA~\cite{qu2025spatialvla}, 3D-VLA~\cite{zhen20243d} & Zero-shot cross-embodiment SR \\
P4: Constraints & FMB~\cite{fmb2023} & SPaSM~\cite{chen2025spasm}, JM2D~\cite{jung2025joint} & Constraint satisfaction / long-horizon SR \\
\bottomrule
\end{tabular}
\end{table}

\subsection{Established Tools and Systems}
\label{subsec:p4-sec9-1-established-tools-and-systems}

Progress in spatially grounded manipulation has been inseparable from a maturing ecosystem of simulators, open software stacks, shared datasets, and, increasingly, commercial robot platforms that serve as common development and deployment substrates.

\textbf{Simulation and physics engines.} GPU-parallel engines such as NVIDIA IsaacGym and the newer Isaac Sim / Isaac Lab\footnote{\url{https://developer.nvidia.com/isaac/sim}}~\cite{makoviychuk2021isaac}, articulated-object simulators such as SAPIEN~\cite{xiang2020sapien} (with its PartNet-Mobility assets), and general-purpose physics engines including MuJoCo~\cite{todorov2012mujoco}, PyBullet\footnote{\url{https://pybullet.org}}, and the recent Genesis\footnote{\url{https://github.com/Genesis-Embodied-AI/Genesis}} provide high-throughput, contact-rich environments, while ManiSkill~\cite{mu2021maniskill}, RLBench~\cite{james2020rlbench}, Robomimic~\cite{robomimic2021}, and RoboTwin~\cite{mu2024robotwin} package standardized task suites, demonstrations, and evaluation protocols that have become de facto testbeds for imitation and policy learning.

\textbf{Open software stacks and datasets.} On the software side, the Robot Operating System (ROS\,2) with the MoveIt motion-planning framework\footnote{\url{https://moveit.ai}} remains the de facto middleware for real-robot integration, and Hugging Face's LeRobot\footnote{\url{https://github.com/huggingface/lerobot}} packages datasets, pretrained policies, and training pipelines for low-cost reproduction. Large cross-embodiment efforts such as Open X-Embodiment\footnote{\url{https://robotics-transformer-x.github.io}}, which aggregates demonstrations and the RT-X policies across dozens of robot platforms, together with open generalist checkpoints such as Octo~\cite{team2024octo} and OpenVLA~\cite{kim2024openvla} and NVIDIA's GR00T humanoid foundation models\footnote{\url{https://developer.nvidia.com/isaac/gr00t}}, expose pretrained weights and fine-tuning interfaces that lower the barrier to transferring foundation-model policies onto new morphologies.

\textbf{Commercial and humanoid platforms.} On the deployment side, a rapidly growing set of companies is moving spatially grounded manipulation out of the laboratory. Boston Dynamics\footnote{\url{https://bostondynamics.com}} (the electric Atlas and the Stretch warehouse robot), Tesla Optimus\footnote{\url{https://www.tesla.com/we-robot}}, Figure AI\footnote{\url{https://www.figure.ai}}, 1X Technologies\footnote{\url{https://www.1x.tech}}, Agility Robotics\footnote{\url{https://agilityrobotics.com}} (the Digit humanoid, piloted in logistics), Physical Intelligence\footnote{\url{https://www.physicalintelligence.company}} (the $\pi_0$ generalist policy), and Sanctuary AI\footnote{\url{https://sanctuary.ai}} each pair learned manipulation with humanoid or mobile-manipulator hardware, while firms such as Covariant\footnote{\url{https://covariant.ai}} target warehouse pick-and-place at commercial scale. These platforms are the clearest evidence that the spatial-intelligence capabilities surveyed above are crossing from benchmark tasks into economically deployed systems.

Together these tools convert the representation and reasoning machinery of Chapter~\ref{sec:reasoning} into reproducible, physically grounded manipulation pipelines; their principal limitation remains the fidelity gap between simulated and real contact dynamics, which the following discussion revisits.


\subsection{Challenges}
\label{subsec:p4-sec9-1-challenges}

Despite rapid progress across all four pivots, bringing \emph{spatial intelligence} to robust, general-purpose manipulation remains bounded by a set of intertwined difficulties that cut across perception, policy, and planning. They cluster into four threads: keeping a persistent and calibrated internal model of the scene; closing the loop between contact physics and active sensing; transferring skills across embodiments while sourcing the data to do so; and meeting the reliability budget of a deployed system.

\textbf{Persistent, calibrated world models.} The first requirement is a \emph{persistent 3D belief}: an actionable, temporally consistent world model that survives heavy occlusion, articulation, deformation, and sensor drift, and that is maintained at control rates rather than reconstructed from scratch each frame. Persistence alone is not enough if the model cannot say how much it trusts itself. \emph{Calibrated uncertainty} asks that distributional beliefs over geometry, contact state, friction, and articulation propagate through perception, policy, and planner, so that a decision reflects confidence rather than a bare point prediction and a low-confidence grasp can trigger caution instead of failure.

\textbf{Contact modeling, sensing, and active perception.} Manipulation is ultimately about contact, yet most spatial models still reason about collision-free geometry rather than the physics of touch. Progress on \emph{contact modeling and sensing} hinges on tightly integrating vision with force and tactile feedback to reason about stick--slip, compliance, and incipient failure. Because contact is often unobservable until it happens, this couples naturally to \emph{information-seeking control}: interleaving manipulation with active perception and probing (touch-to-disambiguate, look-while-acting) inside tractable belief-space planners. And because contact is precisely where simulators diverge from reality, \emph{sim-to-real for contact} stays open, demanding better material, friction, and compliance models together with fast system identification for on-robot adaptation.

\textbf{Cross-embodiment transfer and data at scale.} Skills learned on one platform rarely move cleanly to another. \emph{Cross-embodiment and cross-scene transfer} requires aligning action spaces and spatial tokens across grippers, cameras, and environments without hand-tuned frames or retraining. Sustaining that transfer is a data problem as much as an algorithmic one: \emph{data and evaluation at scale} calls for ``long-tail'' data engines that cover rare geometries, materials, and tools with counterfactual and safety labels, and for benchmarks that measure calibration, persistence, recovery, and long-horizon competence rather than success rate alone.

\textbf{Systems reliability.} Finally, a deployed manipulator must hold together as a \emph{system}: meeting latency budgets, absorbing calibration drift, detecting out-of-distribution geometry, and enabling safe fallback and continual learning without catastrophic forgetting. Closing these gaps will require tighter couplings among 3D world models, affordance predictors, language-grounded objectives, and constraint-aware planners, so that spatial representations are not only expressive but \emph{actionable, calibrated, and safe} in unstructured settings.

\newpage
\section{Autonomous Driving}
\label{sec:p4_s9_2}

% NuScenes-SpatialQA: A Spatial Understanding and Reasoning Benchmark for Vision-Language Models in Autonomous Driving
% Gpt-driver: Learning to drive with gpt


\subsection{Introduction}
\label{subsec:p4-sec9-2-introduction}
Autonomous driving sits at the intersection of spatial intelligence, transportation engineering, and human--machine interaction, and it is one of the most demanding real-world testbeds for the spatial capabilities surveyed in Part~\ref{part:3}.
A self-driving vehicle must continuously close a \emph{perceive--localize--decide--act} loop: it has to recognize and track the surrounding traffic scene, estimate its own pose within a metric map, reason about how nearby agents will move, and commit to safe control actions, all under tight latency and safety budgets.
Advances in deep learning, multi-sensor fusion, and vehicle-to-everything (V2X) communication have steadily pushed the boundaries of what vehicles can perceive, predict, and decide, while fundamental questions around robustness, trust, ethics, and regulation remain open.
Perception and reasoning themselves are treated in Chapter~\ref{chapter:3:sec:6} and Chapter~\ref{sec:reasoning}; what is specific to driving is that each capability must hold up at speed, in all weather, and with lives at stake.

Figure~\ref{fig:p4_s9_2_levels} summarizes the SAE J3016 automation levels that frame this domain.
The SAE J3016 standard \footnote{\url{https://www.sae.org/news/2019/01/sae-updates-j3016-automated-driving-graphic}} provides a widely adopted framework for defining automation levels, with most current production systems still operating at Level 2.
Progress toward higher levels, especially full Level 5 autonomy, faces major technical \cite{YurtseverLCT20}, ethical \cite{StahlTM16}, and financial \cite{fisher2020governmental} challenges.



Deep learning has driven most of this progress, from Transformers~\cite{WangLLLXLCLZH22} and graph networks~\cite{ZhangFWH22} that model multi-agent interaction to lifelong learning for open-world adaptation~\cite{ZhuBZZW24}, but its opacity raises safety concerns compounded by sensor noise and the limited reach of vehicle-to-vehicle communication, and reaching Levels 4 and 5 depends as much on redundant hardware, large-scale data pipelines, and verifiable safety assurance as on model architecture.


\begin{figure}[!ht]
    \centering
    \includegraphics[width=0.98\linewidth]{figures/part4/section9/driving_levels.pdf}
    \caption{Functional distinctions among SAE J3016 driving-automation levels. Levels 0--2 retain a human driver; Level 3 requires a human to resume driving when necessary, whereas Levels 4--5 do not rely on human fallback within their operating scope. The stair heights indicate ordered categories, not measured safety, model complexity, or technology readiness.}
    \label{fig:p4_s9_2_levels}
\end{figure}

\subsection{System Overview}
\label{subsec:p4-sec9-2-system-overview}
Modern self-driving stacks realize the \emph{perceive--localize--decide--act} loop as a modular pipeline, sketched in Figure~\ref{fig:p4_s9_2_overview}. 
A front-end \emph{perception} module converts raw, multi-modal sensor streams (camera, LiDAR, radar) into a structured model of the scene: what objects are present, where the drivable surface and lane geometry lie, and which agents are most safety-critical. 
A \emph{localization and mapping} module places the ego-vehicle within a metric, typically HD-map-anchored, representation of that scene, so that perception outputs can be expressed in a stable world frame. 
A \emph{decision and control} module then predicts how the scene will evolve and selects the steering, acceleration, and braking commands that are sent to the actuators and fed back into the environment, closing the loop. Each of the three stages is a distinct spatial-intelligence problem with its own representations, benchmarks, and failure modes, and they are the pivotal tasks below, in pipeline order.

\begin{figure}[!ht]
    \centering
    \includegraphics[width=0.98\linewidth]{figures/part4/section9/autonomous_driving_overview.pdf}
    \caption{Closed-loop spatial intelligence for autonomous driving: complementary sensor streams feed perception and localization and mapping, whose scene and ego/map state drive decision, planning, and trajectory tracking; vehicle--environment interaction returns new observations. The controller also receives ego state; the curved wire crossing is not a junction. Dashed execution logs feed evaluation separately from control. Scenes, maps, and evaluation plots are illustrative, not measured results.}
    \label{fig:p4_s9_2_overview}
\end{figure}

\subsection{Pivotal Task-1: Perception and Detection}
\label{subsec:p4-sec9-2-pivotal-task-1-perception-and-detection}
Perception is the front-end of the pipeline: it turns raw camera, LiDAR, and radar streams into a structured, layered model of the scene that every downstream module consumes, reusing the detection, segmentation, and 3D scene-understanding machinery of Chapter~\ref{chapter:3:sec:6} but specializing it to the driving Operational Design Domain (ODD). It is best read as a set of complementary structured representations, sketched in Figure~\ref{fig:p4_s9_2_perception}: object detection localizes traffic participants, semantic and occupancy prediction labels the surrounding space, lane detection recovers the road geometry that constrains motion, and pedestrian recognition isolates the most safety-critical agents.

\begin{figure}
    \centering
    \includegraphics[width=0.98\linewidth]{figures/part4/section9/driving_perception_tasks.pdf}
    \caption{Four complementary driving perception tasks feeding a shared scene model: (1) object detection, (2) semantics and 3D occupancy, (3) lane geometry, and (4) pedestrian/VRU tracking and motion prediction. Annotations and the dashed pedestrian path are illustrative, not measured predictions.}
    \label{fig:p4_s9_2_perception}
\end{figure}

The base layer is \textbf{object detection}, which identifies and localizes obstacles and is the most direct determinant of collision safety \cite{FengHWD22}. Its central practical constraint is the ODD: environmental factors such as poor weather degrade sensor visibility and model accuracy \cite{KoopmanF19}, so recent work formalizes and iteratively validates the ODD through quantifiable environmental models and scenario catalogues, trusting a detector only within conditions it has demonstrably handled, though the absence of a standardized ODD specification still makes cross-system safety claims hard to compare. Robustness to adverse weather is therefore a first-class concern \cite{zhang2023perception,TahirZACE24,HnewaR21}: rain, snow, fog, and glare corrupt LiDAR, radar, and camera signals alike, and because driving datasets under-represent these conditions, models trained on fair-weather data generalize poorly. Responses range from dedicated robust pipelines, such as multi-stage visibility enhancement coupled with deep CNN tracking \cite{HassaballahKMM21}, to new sensing hardware, \eg, frequency-modulated continuous-wave (FMCW) LiDAR and high-dynamic-range cameras that recover usable signal in glare at higher cost. Moving from the image plane into metric space, \textbf{3D object detection} localizes and tracks objects in the vehicle's surroundings \cite{LiWLXSLQD22}, where the bird's-eye-view (BEV) representation has become dominant because it yields a unified top-down map well suited to planning. BEVFormer \cite{YangCTTZZHLQLZD23} fuses multi-camera and temporal cues with spatiotemporal transformers to approach LiDAR-level accuracy on nuScenes, while Fast-BEV \cite{LiHCCLSLXSOS24} and SparseBEV \cite{LiuTLWW23} trade marginal accuracy for the lightweight, sparse, query-based computation needed for real-time on-vehicle deployment, together tracing the field's balancing of accuracy, latency, and deployability.

Detections tell the planner \emph{where objects are} but not \emph{which space is free}; that dense layer is supplied by \textbf{semantic segmentation} and its 3D extension, occupancy prediction \cite{ZhouBWN20,NestiRNBB22}. The progression runs from 2D pixel labeling with FCN~\cite{long2015fcn} and DeepLab~\cite{chen2017deeplabv2} to volumetric reasoning: the Tri-Perspective View \cite{HuangZZ0L23} lifts RGB into a 3D representation rivalling LiDAR, SurroundOcc \cite{WeiZZZZL23} uses multi-camera 2D--3D attention to capture arbitrary shapes, and occupancy world models such as OccWorld~\cite{zheng2024occworld} extend this to forecasting how the 3D scene evolves, moving beyond the limits of bounding boxes. Because these dense predictors are data-hungry and domain-sensitive, robustness techniques such as the SePiCo pixel-contrast framework \cite{XieLLLHW23} and GAN-based synthetic data \cite{AbdollahiPSMA21} are used to survive scene variability, occlusion, and weather shift. On top of the free-space layer, \textbf{lane detection} recovers the road geometry that defines the lateral reference for lane-keeping and planning; early CNN- and YOLO-style detectors trace markings well in clear conditions but are brittle on curved, worn, or occluded lanes, motivating successors that lift lanes into 3D and use attention over global context to infer occluded segments. 
Finally, \textbf{pedestrian recognition} focuses on vulnerable road users whose behavior can be difficult to predict: CNN-, YOLO-, and SSD-based detectors deliver the required precision and latency, color--thermal fusion handles low-light scenes, and spatial--temporal models couple detection with motion prediction to anticipate crossing behavior.

Across all four layers, perception is benchmarked on large-scale driving datasets, KITTI~\cite{KITTI-GeigerLU12}, the Waymo Open Dataset~\cite{WaymoOpen-SunKDCPTGZCCVHN20}, and nuScenes~\cite{caesar2020nuscenes} for 2D/3D detection, Cityscapes~\cite{cordts2016cityscapes} and BDD100K~\cite{yu2020bdd100k} for segmentation, dedicated lane benchmarks for lane detection, and CityPersons~\cite{zhang2017citypersons} for pedestrians, and scored with mean Average Precision (mAP) and the nuScenes Detection Score (NDS), mean Intersection-over-Union (mIoU), and F1/accuracy, respectively (Table~\ref{tab:p4_s9_2_benchmarks}). Increasingly these metrics are reported on adverse-weather and out-of-distribution splits, precisely to measure behavior within and beyond the declared ODD.

\subsection{Pivotal Task-2: Localization and Mapping}
\label{subsec:p4-sec9-2-pivotal-task-2-localization-and-mapping}
If perception answers \emph{what is around the vehicle}, localization and mapping answer \emph{where the vehicle is} and \emph{what the persistent structure of the world is}. The two are tightly coupled: a metric, HD-map-anchored pose lets perception outputs be expressed in a stable world frame that planning can rely on, and the map itself is built and maintained from the same sensor streams. As Figure~\ref{fig:p4_s9_2_localization} shows, the problem decomposes into pose estimation against a known map (localization), the construction of that map (3D mapping and reconstruction), and the joint solution of both when no map exists (SLAM), with multi-sensor fusion the common thread that holds accuracy together when any single modality degrades.

\begin{figure}
    \centering
    \includegraphics[width=0.98\linewidth]{figures/part4/section9/driving_localization_and_mapping.pdf}
    \caption{Localization and mapping through uncertainty-aware sensor fusion: camera, LiDAR, radar, and GNSS/IMU measurements update the state posterior, yielding a map-relative pose (evaluated by ATE/RPE), an incrementally updated 3D map, and loop closures that feed back only after verification. Pose and map are parallel products of joint estimation; only verified loop constraints return to the estimator. Recall@$N$ evaluates place retrieval, not pose accuracy.}
    \label{fig:p4_s9_2_localization}
\end{figure}

\textbf{Localization} has moved from GPS-dependent positioning to learned matching of live camera or LiDAR observations against a pre-built map, reaching centimeter-level accuracy even in GPS-denied urban canyons; semantic cues such as roads and buildings further let a vehicle re-localize in changed environments without full remapping. Because no single sensor is reliable across occlusion, weather, and GPS dropout, \textbf{sensor-fusion localization} is the dominant paradigm: probabilistic estimators (particle filters, Kalman filters, and NDT scan matching) reconcile GNSS, IMU, and LiDAR, while learned vision--radar fusion compensates for modality-specific weaknesses in rain, fog, and glare. The recurring bottlenecks here are cross-sensor spatial/temporal synchronization and fault tolerance, which motivate Bayesian uncertainty weighting and online-adaptive fusion that recalibrates as conditions change.

\textbf{3D mapping and reconstruction} supply the persistent representation that localization matches against, turning LiDAR and stereo streams into metric models; for the place-recognition and loop-closure that keep maps globally consistent, transformer descriptors such as SeqOT~\cite{ma2022seqot} match sequential LiDAR scans faster than the sensor frame rate, trained on richly annotated corpora such as ApolloScape~\cite{wang2019apolloscape}. The open problems are the cost of processing dense point clouds at city scale and the handling of occluded or incomplete returns, which favor incremental map updating over full re-mapping. When the map is unavailable, \textbf{SLAM} solves localization and mapping jointly, which is critical in GPS-denied corridors; modern driving-oriented vSLAM hybridizes classical geometry with learned features (for example, extending ORB-SLAM~\cite{mur2017orb} with learned features, semantic labels, and learned depth to favor stable landmarks) and recovers from drift through loop closure with bundle adjustment.

These capabilities are evaluated on KITTI Odometry~\cite{KITTI-GeigerLU12}, Oxford RobotCar~\cite{Maddern2017RobotCar}, and ApolloScape~\cite{wang2019apolloscape} for long-horizon SLAM and place recognition, on nuScenes~\cite{caesar2020nuscenes} and Waymo~\cite{WaymoOpen-SunKDCPTGZCCVHN20} for map-relative pose, and on Argoverse~\cite{Argoverse-ChangLSSBHW0LRH19} for the trajectory forecasting that localization feeds, \eg, with the Motion Transformer~\cite{shi2022mtr} (Table~\ref{tab:p4_s9_2_benchmarks}). The standard metrics, \ie, Absolute and Relative Pose Error (ATE/RPE), recall@$N$ for loop closure, and minADE/minFDE for forecasting, are reported under GPS-denied, adverse-weather, and sensor-degradation conditions to stress exactly the robustness the fusion machinery is built to provide.

\subsection{Pivotal Task-3: Decision Making and Control}
\label{subsec:p4-sec9-2-pivotal-task-3-decision-making-and-control}
Decision making and control is the action end of the pipeline: it consumes the scene model from perception and the pose from localization, predicts how the scene will evolve, and commits to the steering, acceleration, and braking that close the loop, drawing on the spatial-reasoning and planning paradigms of Chapter~\ref{sec:reasoning}. As Figure~\ref{fig:p4_s9_2_decision} indicates, two architectural philosophies dominate, differing in how much of the perceive-to-act mapping is learned end-to-end versus decomposed into interpretable modules, and they are increasingly converging.

\begin{figure}
    \centering
    \includegraphics[width=0.98\linewidth]{figures/part4/section9/driving_decision_and_control.pdf}
    \caption{Three alternative architectures for driving decision and control: (1) modular pipelines with an explicit BEV scene, (2) end-to-end sensor-to-control learning, and (3) hybrids with inspectable scene or trajectory intermediates and explicit constraint checks.}
    \label{fig:p4_s9_2_decision}
\end{figure}

The first philosophy is \textbf{learning-based decision making}, centered on (deep) reinforcement learning, which acquires driving policies through trial-and-error interaction and is well suited to the sequential, interactive nature of driving. 
Representative methods illustrate its design axes: the stochastic latent actor-critic (SLAC)~\cite{lee2020stochastic} learns compact latent states from high-dimensional images, PPO~\cite{schulman2017proximalpo} within a nonlinear control framework handles sharp turns and obstacle avoidance, and the DeepSet-Q architecture~\cite{huegle2019dynamic}, building on permutation-invariant set encoders~\cite{zaheer2017deep}, ingests a variable number of surrounding agents for high-level lane-change decisions.
Hybrids with explicit planning, such as Monte Carlo tree search guided by DRL or scene-representation transformers that forecast neighbor intentions, improve sample efficiency and situational awareness in dense traffic. 
The binding difficulty is that safety-critical episodes are rare in logged data, which pushes the field toward high-fidelity simulation for rare-event coverage and toward encoding hard safety constraints directly in the policy rather than trusting the reward signal alone. 
The second philosophy, \textbf{end-to-end decision making}, learns the mapping from sensors to trajectory or control directly, removing hand-designed interfaces between modules. Early instances learned state-to-action policies with DDPG~\cite{lillicrap2015continuous} in simulators such as TORCS, and multi-task backbones such as HybridNets~\cite{vu2022hybridnets} share computation across detection and segmentation without yet closing the loop to control. Current systems are trained by imitation on large driving logs: TransFuser~\cite{chitta2023transfuser} fuses camera and LiDAR features with transformer attention and predicts waypoints for a closed-loop agent, and UniAD~\cite{hu2023uniad} chains tracking, mapping, motion forecasting, and occupancy prediction through query interfaces inside one network optimized for planning. Driving world models such as GAIA-1~\cite{hu2023gaia} generate future multi-view scenes conditioned on ego actions, giving such policies a learned simulator for training and evaluation.
Opacity is the price of end-to-end training: attribution methods such as layer-wise relevance propagation (LRP)~\cite{bach2015pixel} support interpretation but do not by themselves establish safety or regulatory compliance, which is why hybrid designs keep inspectable scene or trajectory intermediates together with explicit constraint checks, whose guarantees extend only as far as the modeled constraints.

Closed-loop driving policies are evaluated chiefly in simulation, \ie, the CARLA urban simulator~\cite{dosovitskiy2017carla} (NoCrash, Longest6), the nuPlan planning benchmark~\cite{caesar2021nuplan}, and highway-interaction simulators built on HighD/INTERACTION (Table~\ref{tab:p4_s9_2_benchmarks}), because rare, safety-critical events cannot be exercised safely on public roads. The headline metrics are the CARLA Driving Score (route completion penalized by infractions), collision and intervention rates, and route completion, complemented by comfort proxies such as jerk and, for the open-loop forecasting sub-problem, minADE/minFDE.

\begin{table}[!th]
\centering
\fontsize{8}{10}\selectfont
\setlength{\tabcolsep}{1.0mm}
\caption{Representative benchmarks for spatial intelligence in autonomous driving, covering perception, localization and mapping, motion prediction, and decision making.}
\label{tab:p4_s9_2_benchmarks}
\begin{tabular}{L{0.22\textwidth} L{0.22\textwidth} L{0.28\textwidth} L{0.19\textwidth}}
\toprule
\textbf{Pivotal task} & \textbf{Benchmark} & \textbf{Representative method} & \textbf{Headline metric} \\
\midrule
Perception (3D detection) & nuScenes & BEVFormer~\cite{YangCTTZZHLQLZD23} & NDS / mAP \\
Perception (3D detection) & Waymo Open & SparseBEV~\cite{LiuTLWW23} & mAP / latency \\
Perception (segmentation) & Cityscapes & SePiCo~\cite{XieLLLHW23} & mIoU \\
Perception (occupancy) & nuScenes Occ3D & SurroundOcc~\cite{WeiZZZZL23} & mIoU \\
Localization \& mapping & KITTI Odometry & Deep / semantic vSLAM & ATE / RPE \\
Prediction (forecasting) & Argoverse & Motion Transformer (MTR)~\cite{shi2022mtr} & minADE / minFDE \\
Decision \& control & CARLA (Longest6) & TransFuser~\cite{chitta2023transfuser} & Driving Score \\
Decision \& control & nuScenes (open-loop planning) & UniAD~\cite{hu2023uniad} & L2 error / collision rate \\
Decision \& control & nuPlan & Learned planners & Closed-loop score \\
\bottomrule
\end{tabular}
\end{table}


\subsection{Established Tools and Systems}
\label{subsec:p4-sec9-2-established-tools-and-systems}
The systems on the road today populate a deployment ladder whose rungs are the SAE levels of Figure~\ref{fig:p4_s9_2_levels}, but what separates one rung from the next is rarely a single breakthrough in perception accuracy.
It is, instead, two coupled engineering choices: how wide an Operational Design Domain (ODD) the system claims, and where legal responsibility for the driving task sits.
Three tiers are in genuine production, \ie, supervised Level~2 assistance, certified Level~3 ``eyes-off'' automation, and driverless Level~4 robotaxis, and, as the following progression shows, moving up the ladder consistently means adding sensing and compute redundancy while \emph{narrowing} the domain in which the autonomy guarantee holds.

\subsubsection{Large-scale Production ADAS (SAE Level 2 / 2+)}
The defining property of Level~2 is that the system performs continuous lateral and longitudinal control, yet the human remains the legal driver and the immediate fallback; consequently, every Level~2 product is fundamentally an exercise in bounding the ODD and enforcing driver attention rather than in achieving autonomy, and the deployed systems differ mainly in how they strike that balance.
One design philosophy gates hands-free operation to pre-validated, HD-mapped limited-access highways and polices it with in-cabin driver monitoring: \textbf{GM Super Cruise}\footnote{\url{https://www.gm.com/innovation/autonomous-driving}}, \textbf{Ford/Lincoln BlueCruise}\footnote{\url{https://www.lincoln.com/technology/bluecruise/}}, and \textbf{BMW Driving Assistant Professional}\footnote{\url{https://www.bmwgroup.com/en/news/general/2023/driver-assistance-systems.html}} all keep the vehicle centered, regulate following distance, and offer assisted lane changes, but escalate alerts and disengage the moment attention drifts or the car leaves a supported segment; their reliability is a direct product of this deliberately narrow envelope.
A second philosophy trades map dependence for scale: \textbf{Tesla Autopilot}\footnote{\url{https://www.tesla.com/en_gb/support/autopilot}} and \textbf{FSD (Supervised)}\footnote{\url{https://www.tesla.com/en_gb/fsd}} rely on a camera-first vision stack refined through fleet-wide over-the-air training, which extends assistance to a far broader and less-bounded set of scenarios but, lacking comparable geofencing, still demands constant supervision and degrades on poor markings, construction zones, and adverse weather. At the conservative end, \textbf{Toyota Safety Sense} declines hands-free operation altogether, functioning as an always-on crash-avoidance layer (automatic emergency braking, lane-keeping, and adaptive cruise) that intervenes only to prevent or mitigate collisions. Despite these differing emphases, the binding limitation is identical across all of them: a bounded ODD coupled with a mandatory human fallback, and it is precisely that fallback which the next tier begins to remove.

\subsubsection{Certified ``Eyes-off'' Conditional Automation (SAE Level 3)}
Level~3 is the first tier at which, inside a defined envelope, the driver may legally divert attention and the manufacturer assumes liability for the driving task. This shift in responsibility, rather than any new perception capability, is what dictates the engineering.
\textbf{Mercedes-Benz Drive Pilot}\footnote{\url{https://www.mercedes-benz.co.uk/passengercars/technology/autonomous-driving.html}} and \textbf{Honda Traffic Jam Pilot}\footnote{\url{https://global.honda/en/tech/Traffic_Jam_Pilot_TJP/}} purchase the ``eyes-off'' guarantee by pairing complementary sensing (cameras, radar, and LiDAR with high-precision, HD-map-anchored localization) with redundant braking, steering, and power paths that retain control through a single-component failure.
Both, tellingly, keep the envelope deliberately small (low-speed or traffic-jam flow on mapped, divided highways), and both center on the same safety mechanism: as conditions approach the ODD boundary, the system issues a timed takeover request and, absent a driver response, executes a minimal-risk maneuver such as a controlled stop with hazard signals.
The narrowness of the Level~3 envelope is therefore not a marketing concession but the direct cost of guaranteeing safe behavior once the human is no longer continuously in the loop: a cost the next tier pays in full by removing the driver entirely.

\subsubsection{Driverless Robotaxis and Autonomous Ride-hailing (SAE Level 4)}
Level~4 removes the human from the loop entirely within a geofence: there is no driver to hand back to, so the vehicle must resolve every in-domain situation itself or bring itself to a safe state, with remote operators supplying only high-level guidance rather than moment-to-moment control.
To meet this far stricter bar, \textbf{Waymo One}\footnote{\url{https://waymo.com/}} and \textbf{Baidu Apollo Go}\footnote{\url{https://www.apollogo.com/}} converge on essentially the same recipe: a dense, overlapping multi-sensor suite (lidar, radar, and camera vision) for redundant coverage, HD-map-anchored localization, and conservative prediction-and-planning that favors lawful, cautious maneuvers through dense urban traffic.
Each makes the problem tractable not by claiming universal competence but by tightly restricting where and when it operates and validating performance city by city before extending the geofence.
Across the three tiers, every step up the ladder buys wider autonomy by narrowing the operating domain and adding redundancy, until at Level~4 the system finally internalizes the fallback that Levels~2 and~3 delegated to the driver.



\subsection{Challenges}
\label{subsec:p4-sec9-2-challenges}
Every tier of the deployment ladder works because it limits where it operates, and general, trustworthy autonomy requires lifting those limits. Four threads, each cutting across perception, localization, and decision, stand in the way: justifying life-critical decisions (\emph{safety and explainability}), dependable behavior across the long tail of road, weather, and traffic conditions (\emph{resilience and robustness}), adaptation to genuinely novel situations from little prior data (\emph{unfamiliar road scenarios}), and continual learning on deployed vehicles without eroding mastered behavior (\emph{enhanced autonomy in complex scenarios}).

\paragraph{Safety and Explainability.}
Because driving decisions are life-critical, the inability to explain why a model selected a given maneuver is a first-order obstacle to certification and public trust.
Post-hoc attribution methods such as layer-wise relevance propagation~\cite{bach2015pixel} and Shapley-value attribution expose which inputs drove a decision, supporting both engineering debugging and regulatory audit, while causal-inference models go further by separating mere correlation from the actual influence of factors such as road state and a maneuver's effect on surrounding agents.
The unresolved difficulty is that these explanations remain largely retrospective and local to individual decisions; converting them into guarantees that provably hold across the full operating domain is still an open problem.

\paragraph{Resilience and Robustness.}
Robustness to the long tail of driving conditions is commonly pursued through curriculum learning, where a policy is trained on a graded sequence from simple navigation to dense mixed-traffic interaction, which improves sample efficiency and training stability relative to confronting hard cases directly.
Its effectiveness, however, hinges on a scenario curriculum that faithfully spans the deployment distribution and on objectives that keep safety and comfort constraints satisfied at every difficulty level.
The practical risk is that, absent a benchmark which deliberately surfaces rare but critical events, strong curriculum performance on common scenes can mask failures precisely on the cases that determine real-world safety.

\paragraph{Unfamiliar Road Scenarios.}
A distinct difficulty is generalization to genuinely novel situations (an unseen sign, an irregular work zone, an atypical vehicle) where data-hungry models degrade but human drivers adapt from a single exposure.
Approaches that target this gap include probabilistic program induction in a Bayesian framework, which represents a new concept as a compact program inferred from few examples, and the injection of compositional, causal, and meta-learning priors so that, for instance, a system can infer that a fallen obstacle will trigger sudden lane changes in adjacent traffic.
The standing obstacle is validating such few-shot behavior to a safety standard, since the very rarity of these events makes their failure modes hard to measure and certify.

\paragraph{Enhanced Autonomy in Complex Scenarios.}
Sustained operation in changing environments motivates lifelong learning~\cite{ZhuBZZW24}, in which a deployed system keeps acquiring new knowledge (a new traffic pattern or road layout) without retraining from scratch and without overwriting prior competence.
The central technical barrier is catastrophic forgetting, which is addressed by memory-augmented architectures and continual-learning regularizers that protect previously learned parameters, together with transfer mechanisms that reuse knowledge across tasks and cities.
The harder, still-open requirement is to perform such updates on-vehicle under real-time and verification constraints, so that an update improving one scenario cannot silently regress a safety-critical behavior elsewhere.

\newpage

\section{Wearable AI System}
\label{sec:p4_s9_3}


\subsection{Introduction}
\label{subsec:p4-sec9-3-introduction}

Wearable devices have become a central sensing modality for embodied and spatial intelligence, delivering continuous, fine-grained measurements of human motion, physiology, and physical interaction. Because they are worn on or integrated into the body, they access proprioceptive and kinesthetic cues, joint rotations, muscle activations, contact forces, and locomotion patterns, which exteroceptive sensors such as RGB-D cameras or LiDAR cannot capture directly. We use the term broadly, covering inertial measurement units (IMUs), surface electromyography (EMG) sensors, smart textiles, haptic interfaces, and wearable cameras: body-mounted technologies that together model human dynamics and embodied interaction in real-world settings. It is this egocentric, body-centric vantage point that makes wearables valuable for spatial intelligence, letting a system reason about how a person moves through, reads, and acts upon their surroundings, and it has driven rapid adoption across health monitoring, sports analytics, rehabilitation~\cite{patel2012review}, industrial safety, and human--robot collaboration.

The signals wearables provide are especially valuable because they are anticipatory. Joint kinematics and muscle activity reveal what a person is doing and, crucially, what they are about to do next, which is what underpins intent recognition, action forecasting, and safe coordination in shared spaces. Wearable IMU arrays, for example, recover whole-body skeletal motion in cluttered or occluded settings where optical capture breaks down~\cite{huang2018dip, vonmarcard2017sparse}, and EMG surfaces pre-activation patterns that precede any visible movement~\cite{atzori2014ninapro}.

These devices are often bracketed with head-mounted XR hardware, yet their role in spatial intelligence differs in kind. XR headsets integrate sensing and display for immersive visualization and mixed-reality interaction, leaning on outward-facing cameras for SLAM and environment mapping. The wearables treated here look inward instead, supplying the biomechanical and physiological state, high-frequency, continuous, and on-body, that an XR headset does not observe.

The spatial-intelligence capabilities that wearables ultimately serve, pose estimation, action recognition, localization, and trajectory prediction, are formulated in Part~III (spatial perception in Chapter~\ref{chapter:3:sec:6}, spatial reasoning in Chapter~\ref{sec:reasoning}). This section does not re-derive them; it concentrates on the body-centric sensing layer, the devices that acquire embodied signals, the systems assembled around them, and the tasks, methods, and failure modes that surface once spatial intelligence must be grounded in continuous on-body measurement, spanning long-horizon activity and gait analysis, intention prediction for human--robot collaboration, and real-time safety in industrial and logistics settings.


\subsection{Wearable Devices}
\label{subsec:p4-sec9-3-wearable-devices}

Body-centric sensing spans a diverse device ecosystem, summarized in Table~\ref{tab:p4_s9_3_devices}.

\paragraph{Motion and Body-State Sensing.}
Inertial measurement units (accelerometers, gyroscopes, and often magnetometers) are the most widely deployed body-worn sensors. Consumer watches and bands such as the Apple Watch\footnote{\url{https://www.apple.com/watch/}}, Samsung Galaxy Watch\footnote{\url{https://www.samsung.com/global/galaxy/galaxy-watch/}}, Fitbit\footnote{\url{https://www.fitbit.com/}}, and Xiaomi Band stream inertial data for activity and gait monitoring, while research and industrial systems place IMUs at specific anatomical sites and fuse them into full-body kinematics~\cite{vonmarcard2017sparse, yi2021transpose}; because they need no line of sight, they keep working under occlusion and poor lighting. Full-body motion-capture suits such as Xsens\footnote{\url{https://www.movella.com/products/xsens}}, Rokoko Smartsuit Pro\footnote{\url{https://www.rokoko.com/}}, and Noitom\footnote{\url{https://www.noitom.com/}} combine many IMUs, sometimes with pressure or stretch sensors, into accurate skeletal pose outside the studio, which has made them standard in biomechanics, sports, robotics, and character animation~\cite{trumble2017totalcapture}. Smart clothing embeds conductive fabrics, stretch sensors, or distributed pressure sensors in garments for unobtrusive, long-duration monitoring of posture, joint angles, and whole-body contact~\cite{luo2021learning}, and physiological sensors (heart rate, galvanic skin response, respiration), foot-mounted inertial or pressure sensors, and ear-worn IMUs add workload, ground-contact timing, and head orientation~\cite{patel2012review}.

\paragraph{Neuromuscular, Egocentric, and Interactive Wearables.}
Electromyography (EMG) armbands and high-density arrays sense the muscle activations that precede visible movement, giving early cues for gesture recognition and intent estimation~\cite{atzori2014ninapro}, and combining EMG with inertial data makes them more robust to arm posture and fatigue~\cite{atzori2016deep}. Wearable cameras add the exteroceptive, first-person view: egocentric RGB, depth, or fisheye streams supply scene geometry and activity context that can be fused with on-body signals~\cite{grauman2022ego4d, adt}. Two device families also act on the body. Haptic gloves such as HaptX\footnote{\url{https://haptx.com/}} and SenseGlove\footnote{\url{https://www.senseglove.com/}} and haptic vests such as bHaptics\footnote{\url{https://www.bhaptics.com/}} return tactile, vibration, or force feedback during manipulation, teleoperation, and training~\cite{pacchierotti2017wearable}, and exoskeletons combine IMUs, torque and force sensors, and joint encoders with actuation, estimating user intent in real time to assist or rehabilitate movement~\cite{young2017state}, which makes them a direct test of embodied spatial intelligence.

\begin{table}[!th]
\centering
\caption{Wearable body-centric sensing devices, by sensed signal and recovered spatial or biomechanical quantity. Haptic and exoskeletal wearables are partly \emph{actuators} that return force and tactile signals to the wearer.}
\label{tab:p4_s9_3_devices}
\fontsize{8}{10}\selectfont
\setlength{\tabcolsep}{1.0mm}
\begin{tabular}{L{0.22\textwidth} L{0.28\textwidth} L{0.42\textwidth}}
\toprule
\textbf{Device family} & \textbf{Primary signal} & \textbf{Spatial / biomechanical quantity recovered} \\
\midrule
IMU wearables            & accel., gyro., (magnet.)          & orientation, acceleration, gait, sparse-to-full body kinematics \\
Motion-capture suits     & multi-IMU (+ pressure/stretch)    & full-body skeletal pose in unconstrained settings \\
EMG / gesture bands      & neuromuscular activation          & muscle intent, hand/arm gestures, pre-movement cues \\
Smart clothing / textile & distributed pressure, strain      & posture, joint angles, whole-body contact patterns \\
Egocentric cameras       & first-person RGB / depth          & scene geometry, activity, environmental context \\
Haptic / force-feedback  & tactile, vibration, force (I/O)   & contact rendering, spatial alerts (feedback channel) \\
Exoskeletons             & IMU, torque, force, joint enc.\   & joint torque, user effort and movement intent \\
Physiological / other    & HR, GSR, respiration; foot press. & workload, exertion, ground contact, head orientation \\
\bottomrule
\end{tabular}
\end{table}


\subsection{System Overview}
\label{subsec:p4-sec9-3-system-overview}

Taken together, the devices above form the front end of a single body-centric perception-to-action loop that defines the wearable AI system as a whole. 
Raw on-body signals are first \emph{sensed} (inertial, neuromuscular, tactile, and egocentric-visual streams), then \emph{modeled} into geometric and biomechanical state (joint kinematics, body pose, contact forces), \emph{interpreted and predicted} into higher-level behavioral and intentional descriptions (activities, gestures, short-horizon motion forecasts), and finally fed back to the wearer or to a collaborating agent through \emph{action} channels (haptic cues, exoskeleton assistance, safety alerts, or shared-control commands). 
The loop then closes: the consequences of feedback alter the wearer's motion, which the sensors observe anew.

Figure~\ref{fig:p4_s9_3_overview} schematizes this loop and situates the three pivotal tasks that organize the remainder of this section as complementary stages of it. 
\textbf{Pivotal Task-1 (Embodied Motion Sensing and Pose Reconstruction)} occupies the sense-to-model stage, turning body-centric signals into stable geometric descriptions of motion. 
\textbf{Pivotal Task-2 (Behavioral Understanding and Intention Modeling)} occupies the interpret-and-predict stage, mapping reconstructed motion and neuromuscular cues to recognized activities and anticipated actions. 
\textbf{Pivotal Task-3 (Safe Collaboration and Embodied Interaction)} closes the loop at the action stage, converting interpreted human state into coordinated, safe, and assistive responses in shared physical environments. 
Because all three stages must run under the power, compute, and reliability constraints of body-worn hardware, lightweight on-device inference is a cross-cutting concern that recurs throughout.


\begin{figure*}[!t]
    \centering
    \includegraphics[width=\linewidth]{figures/part4/section9/wearable_overview.pdf}
    \caption{The wearable AI system as a body-centric loop: on-body signals are modeled into kinematic and biomechanical state (T1), interpreted into activity and intent (T2), and acted on through haptic, exoskeleton, and safety channels (T3), used separately or together. Device and prediction illustrations are schematic.}
    \label{fig:p4_s9_3_overview}
\end{figure*}


\subsection{Pivotal Task-1: Embodied Motion Sensing and Pose Reconstruction}
\label{subsec:p4-sec9-3-pivotal-task-1-embodied-motion-sensing-and-pose-reconstruction}

The first stage of the wearable loop turns raw body-centric signals into a geometric description of how the wearer moves. 
It is a bottom-up reconstruction problem: per-segment inertial and biomechanical measurements must yield joint angles, segment velocities, and ultimately a full-body pose that remains stable over long horizons and in settings where exteroceptive capture fails. 
This task subsumes full-body skeletal pose estimation, fine-grained limb tracking, real-world locomotion and gait analysis, and body-centric localization in GPS-denied spaces; the underlying pose-estimation formulations are developed in Chapter~\ref{chapter:3:sec:6}, so the emphasis here is on reconstruction from on-body sensing. 
Figure~\ref{fig:p4_s9_3_motion} illustrates this reconstruction stack from segment-level signals to full-body pose.

\begin{figure}[!t]
    \centering
    \includegraphics[width=\linewidth]{figures/part4/section9/wearable_motion.pdf}
    \caption{Motion sensing and pose reconstruction from body-worn IMUs: estimators supported by a parametric body model recover root-relative full-body pose, while visual--inertial fusion anchors a world-frame trajectory. Contact or zero-velocity updates apply only when their assumptions hold, and metric scale requires sufficient observability; signals and trajectories are illustrative.}
    \label{fig:p4_s9_3_motion}
\end{figure}

\textbf{From signals to kinematics.} Reconstruction begins with kinematic and dynamic modeling, which provides the physical grounding for interpreting inertial and pressure-based signals. 
Forward and inverse kinematics infer joint angles and segment orientations from inertial measurements, compensating for drift and sensor noise through optimization-based filtering and biomechanical constraints. Classical recursive estimators such as extended and unscented Kalman filters incorporate kinematic priors and temporal smoothness to improve stability in real-world settings~\cite{welch1995kalman, madgwick2011estimation}, while articulated-body dynamics and parametric body models such as SMPL~\cite{loper2015smpl} enforce anatomical plausibility for complex motions such as rapid locomotion or limb coordination. 
Multi-IMU fusion against musculoskeletal models extends these ideas to full-body pose estimation in challenging environments where optical motion-capture systems fail~\cite{huang2018dip, yi2021transpose, yi2022pip}, and recent systems fold in physics-based dynamics and inter-sensor ultra-wideband ranging to tighten accuracy and suppress jitter~\cite{armani2024ultra}.

\textbf{Learning-based reconstruction and sensor fusion.} Where analytic models struggle with sparse or noisy sensing, data-driven approaches learn the mapping from raw inertial streams to pose directly~\cite{huang2018dip, yi2021transpose}, and sensor fusion integrates heterogeneous modalities for drift-compensated estimates. IMU-and-vision fusion resolves ambiguities in inertial tracking by leveraging environmental context from egocentric cameras~\cite{grauman2022ego4d, adt}, while drift-compensation techniques, including zero-velocity updates and constraint-based corrections, stabilize long-horizon wear. 
For body-centric localization, sparse-sensor systems couple inertial pose estimation with simultaneous localization and mapping, so that the wearer's global trajectory is recovered together with body pose rather than left to drift freely~\cite{yi2023egolocate}; visual-inertial odometry that fuses the same inertial signals with first-person imagery offers a complementary route through GPS-denied indoor spaces, assessed on dedicated benchmarks~\cite{Schubert2018TUMVI}.

\textbf{Benchmarks and evaluation.} The kinematic, learning-based, and fusion methods above are measured against the same yardsticks: large-scale motion-capture corpora such as AMASS~\cite{mahmood2019amass}, which provides a canonical surface-and-skeleton reference for full-body pose, alongside dedicated inertial-pose datasets for sparse-IMU reconstruction~\cite{trumble2017totalcapture, huang2018dip}, while body-centric localization is assessed separately on visual-inertial odometry benchmarks~\cite{Schubert2018TUMVI}. Reconstruction quality is reported as mean per-joint position error (MPJPE) and joint-angle error, and localization as absolute trajectory error and accumulated drift (Table~\ref{tab:p4_s9_3_benchmarks}). 
Read against the needs of real deployment, these numbers reveal a clear gap. On curated benchmarks, sparse-IMU systems now recover \emph{root-relative} full-body pose to within roughly five centimeters of optical ground truth~\cite{yi2022pip}, but global translation still drifts over long sequences unless it is anchored by SLAM or ranging~\cite{yi2023egolocate}, and even the root-relative accuracy falls off in the wild, where fast or unseen motions and loose sensor placement push errors past the tolerance that downstream control or clinical gait analysis can absorb. The pose the next stage consumes is thus dependable on average yet brittle at exactly the fast, atypical movements that matter most.


\subsection{Pivotal Task-2: Behavioral Understanding and Intention Modeling}
\label{subsec:p4-sec9-3-pivotal-task-2-behavioral-understanding-and-intention-modeling}

Once motion has been reconstructed, the second stage of the loop interprets it: recognizing what the wearer is doing and anticipating what they will do next. 
Continuous neuromuscular and inertial cues feed activity recognition and segmentation and then short-horizon intention prediction, even in daily-life settings where visual systems face severe occlusion or field-of-view limits. 
Because the recognition and reasoning formulations themselves are treated in Chapter~\ref{sec:reasoning}, the focus here is on how high-frequency wearable streams are modeled to support behavioral inference. 
Figure~\ref{fig:p4_s9_3_behavior} sketches the path from multimodal time-series and body-graph representations to recognized activity and predicted intent.

\begin{figure}[!t]
    \centering
    \includegraphics[width=\linewidth]{figures/part4/section9/wearable_behavior.pdf}
    \caption{Behavioral understanding and intention modeling: IMU and EMG sequences and a kinematic body graph pass through temporal and graph encoders, fused by attention to recognize current activity and forecast short-horizon intent. The EMG pre-activation inset draws anticipatory cues from the observed stream, not from future measurements.}
    \label{fig:p4_s9_3_behavior}
\end{figure}

\textbf{Sequence models for wearable time-series.} The availability of large-scale wearable recordings has driven deep models tailored to high-frequency temporal streams from IMU, EMG, and multimodal sensors. Recurrent networks, temporal convolutional networks, and Transformers are widely applied to gait classification, gesture recognition, and activity segmentation~\cite{ordonez2016deep}: temporal convolutions capture multi-scale dependencies for long-horizon pattern detection~\cite{bai2018tcn}, while attention-based architectures learn complex correlations across sensor channels~\cite{attention}. Because labeled wearable data are scarce and expensive, self-supervised pretraining on unlabeled motion recordings has become the dominant route to transferable representations~\cite{haresamudram2021contrastive}: masked-reconstruction encoders such as LIMU-BERT learn compact IMU features from raw streams~\cite{xu2021limubert}, and multimodal contrastive objectives such as IMU2CLIP align inertial signals with video and language, pointing toward wearable foundation models that generalize across tasks and users~\cite{moon2023imu2clip}.

\textbf{Graph-based body representations.} Because human anatomy is structured as a kinematic tree, graph neural networks offer a natural representation for wearable-derived motion~\cite{kipf2017semisupervised}, encoding relational dependencies between joints and limbs directly in graph space. 
Such formulations fuse IMU placements across body segments, propagate temporal features along skeletal edges, and remain robust to missing or corrupted sensor channels; combining graph convolution with temporal attention jointly models spatial structure and temporal evolution for activity recognition and biomechanical analysis~\cite{yan2018stgcn}.

\textbf{Neuromuscular cues and multimodal fusion for intent.} Beyond recognizing observed motion, EMG-based wearables detect subtle neuromuscular activations that precede visible movement~\cite{atzori2014ninapro}, providing predictive cues for early intention estimation and responsive human-robot collaboration.
Fusing EMG with inertial data sharpens fine-grained gesture and anticipatory-command decoding~\cite{atzori2016deep}, while fusing IMU with egocentric video grounds activity recognition in environmental context~\cite{stein2013combining, grauman2022ego4d}.
A recent milestone shows how far this can scale: a surface-EMG wristband trained across more than a thousand users decodes gestures, clicks, and handwriting from neuromuscular signals alone and, critically, generalizes to new wearers without per-user calibration~\cite{kaifosh2025neuromotor}, addressing the inter-user variability that has long limited EMG interfaces.
Physiological and pressure cues additionally support inference of fatigue, balance, and abnormal gait over extended periods.

\textbf{Benchmarks and evaluation.} Taking the reconstructed motion of the previous stage as input, the sequence, graph, and fusion models above are evaluated on wearable human-activity-recognition corpora~\cite{reiss2012pamap2, chavarriaga2013opportunity}, EMG gesture datasets~\cite{atzori2014ninapro}, and egocentric activity datasets such as Ego4D~\cite{grauman2022ego4d} and EPIC-KITCHENS~\cite{EPIC-KITCHENS}, with recognition reported as accuracy and macro-F1 and intention prediction as anticipation accuracy and prediction lead time (Table~\ref{tab:p4_s9_3_benchmarks}). 
Here the gap to maturity is one of generalization rather than peak score. Balanced locomotion corpora such as PAMAP2 are close to saturated, yet the same architectures score markedly lower on cluttered, null-dominated sets such as Opportunity and lose accuracy again under cross-user and cross-environment evaluation; reliable intention forecasting, meanwhile, still extends only to short, well-rehearsed horizons. This falls well short of the open-world, anticipatory understanding that the safe-collaboration stage below depends on, where the system must read unfamiliar users in unfamiliar settings.


\subsection{Pivotal Task-3: Safe Collaboration and Embodied Interaction}
\label{subsec:p4-sec9-3-pivotal-task-3-safe-collaboration-and-embodied-interaction}

The third stage closes the loop, converting interpreted human state into coordinated, safe, and assistive responses in shared physical environments. 
Contact- and force-level sensing of human--environment interaction supports real-time inference of effort and hazard, which in turn drives feedback for collaboration and shared autonomy. 
The collaboration and shared-control reasoning that this task builds on is developed in Chapter~\ref{sec:reasoning} and, for the manipulation setting specifically, in Section~\ref{sec:p4_s9_1}; here the emphasis is on the wearable feedback and safety layer. 
Figure~\ref{fig:p4_s9_3_safety} depicts this closing stage, from contact sensing to coordinated safety response.

\begin{figure}[!t]
    \centering
    \includegraphics[width=\linewidth]{figures/part4/section9/wearable_safety.pdf}
    \caption{Safe human--robot collaboration with wearables: fused contact, force, and proximity signals, together with intent forecasts (T2) and robot motion context, drive motion prediction and risk evaluation, which select glove, exoskeleton, or cobot responses. Prediction envelopes and risk thresholds are schematic, not measured safety guarantees.}
    \label{fig:p4_s9_3_safety}
\end{figure}

\textbf{Sensing interaction and effort.} Haptic and tactile wearables capture contact forces, surface interactions, and manipulative behaviors, supplying the signals needed to model human--environment interaction and object handling~\cite{HOI4D, luo2021learning}. Exoskeletons equipped with torque and force sensors use these signals to infer user effort and detect hazardous conditions, adjusting assistive strategies accordingly~\cite{young2017state}; probabilistic fusion that integrates biomechanical priors and uncertainty estimates is central to robustness in such safety-critical settings.

\textbf{Coordination, safety, and feedback.} In industrial and shared-autonomy scenarios, on-body inertial and proximity sensors contribute to safety monitoring by detecting collision risks, unsafe postures, and sudden motion anomalies, while haptic channels return spatial alerts and force feedback that close the sensorimotor loop~\cite{pacchierotti2017wearable}. 
Combined with the intention forecasts of Pivotal Task-2, these signals form the basis for predictive safety mechanisms and shared-control frameworks that anticipate human motion rather than merely reacting to it.

\textbf{On-device constraints.} Because feedback must be timely, the action stage is where the wearable system's power, memory, and compute limits bite hardest. Model compression (pruning, quantization, and distillation)~\cite{hinton2015distilling} and edge-optimized architectures~\cite{howard2017mobilenets} enable real-time inference on wearable microcontrollers or smartphone-class processors, and on-device adaptation mitigates inter-user variability without relying on remote servers, keeping safety-critical computation local and low-latency.

\textbf{Benchmarks and evaluation.} Closing the loop, the interaction-sensing, safety, and on-device methods above are assessed on human-object-interaction datasets such as HOI4D~\cite{HOI4D} and, for assistive control, on exoskeleton and teleoperation testbeds, with metrics spanning contact- and force-estimation error, collision-avoidance and safety-violation rates, and intention-anticipation lead time (Table~\ref{tab:p4_s9_3_benchmarks}). 
Evaluation here is the least mature of the three stages: it remains fragmented across hardware-specific testbeds with few shared protocols, so reported gains are hard to compare and seldom stress real safety guarantees. 
Current systems perform convincingly in instrumented demonstrations, but closed-loop, certified safety under the full variability of human partners, the bar that genuine shared autonomy must clear, is still an open target rather than a measured result. The robustness, generalization, and safety gaps that recur at each stage of the loop are what the challenges below take up directly.


\begin{table}[!th]
\centering
\caption{Representative benchmarks for wearable spatial-intelligence tasks, grouped by pivotal task.}
\label{tab:p4_s9_3_benchmarks}
\fontsize{8}{10}\selectfont
\setlength{\tabcolsep}{1.0mm}
\begin{tabular}{L{0.16\textwidth} L{0.20\textwidth} L{0.34\textwidth} L{0.21\textwidth}}
\toprule
\textbf{Pivotal Task} & \textbf{Benchmark / Dataset} & \textbf{Representative method / use} & \textbf{Headline metric} \\
\midrule
T1: Motion \& Pose & AMASS~\cite{mahmood2019amass} & Full-body pose reference (SMPL surface) & MPJPE \\
T1: Motion \& Pose & TotalCapture~\cite{trumble2017totalcapture} & Video+inertial full-body pose & MPJPE \\
T1: Motion \& Pose & DIP-IMU~\cite{huang2018dip} & Sparse-IMU pose reconstruction (DIP, TransPose, PIP) & Joint-angle / position error \\
T1: Localization & TUM-VI~\cite{Schubert2018TUMVI} & Visual-inertial odometry (body-centric) & ATE / drift \\
T2: Activity & PAMAP2~\cite{reiss2012pamap2} & Wearable activity recognition (deep recurrent / TCN) & Accuracy / macro-F1 \\
T2: Activity & Opportunity~\cite{chavarriaga2013opportunity} & On-body activity recognition (DeepConvLSTM) & Weighted / macro-F1 \\
T2: Activity & Ego4D~\cite{grauman2022ego4d} & Egocentric activity recognition & Accuracy / F1 \\
T2: Gesture & Ninapro~\cite{atzori2014ninapro} & EMG gesture / movement classification & Accuracy \\
T2: Gesture & emg2qwerty~\cite{sivakumar2024emg2qwerty} & Large-scale surface-EMG typing & Character error rate \\
T3: Interaction & HOI4D~\cite{HOI4D} & Human-object interaction sensing & Contact / interaction acc. \\
\bottomrule
\end{tabular}
\end{table}


\subsection{Established Tools and Systems}
\label{subsec:p4-sec9-3-established-tools-and-systems}
Commercial wearables have converged on the head and the wrist, and what separates current products is less their sensors than where perception runs and whether anything is shown back to the wearer. Full augmented-reality headsets are covered with the XR platforms in Section~\ref{sec:p4_s10_1}; here the focus is on everyday-wearable glasses, neural input bands, and the research platforms and chips behind them.

\textbf{AI glasses without a display.} By volume, the dominant wearable AI device is a pair of ordinary-looking glasses carrying a camera, microphones, open-ear speakers, and an IMU, with a multimodal assistant running on the paired phone or in the cloud. Meta's line with EssilorLuxottica\footnote{\url{https://www.meta.com/ai-glasses/}} defines the category: the second-generation Ray-Ban Meta (2025) records 3K video from a 12-megapixel camera and runs up to eight hours per charge, and the Oakley Meta Vanguard adapts the design to sport with a 122-degree camera. Display-less models made up 96\% of AI-glasses shipments in the first half of 2026, almost all of them Meta's\footnote{\href{https://counterpointresearch.com/en/insights/global-ai-glasses-shipments-surge-263-percent-yoy-with-display-less-ai-glasses}{Counterpoint Research, global AI glasses shipments, H1 2026}.}, while Xiaomi's AI glasses (2025) and Alibaba's Qwen-based Quark AI glasses\footnote{\url{https://www.quark.cn/}} (November 2025) lead a fast-growing Chinese segment. Google and Samsung previewed Gemini-based Android XR glasses\footnote{\url{https://www.android.com/xr/}} with Gentle Monster and Warby Parker in May 2026, with the first display-less models due in the fall. Spatially, these devices answer questions about what the wearer is looking at one camera frame at a time; most keep no persistent map of the surroundings, so the continuous body and scene tracking of Pivotal Tasks~1 and~2 is still largely absent from shipping products.

\textbf{Display glasses and neural input.} Adding a display makes the glasses an interface, and input then becomes the hard problem. Meta Ray-Ban Display\footnote{\url{https://www.meta.com/ai-glasses/meta-ray-ban-display/}} (September 2025) puts a small monocular display, with a 20-degree field of view, in the right lens and ships with the Meta Neural Band, a surface-EMG wristband that decodes finger gestures and handwriting; it is the commercial form of the generic neuromotor interface described in Pivotal Task-2~\cite{kaifosh2025neuromotor}. In waveguide display glasses, Rokid\footnote{\url{https://global.rokid.com/}} led shipments in the first half of 2026, ahead of Meta and Even Realities\footnote{\url{https://www.evenrealities.com/}}, whose glasses favor a minimal text display over cameras.

\textbf{Research platforms and on-glass compute.} Much of the egocentric data behind the methods of this section was recorded on research glasses rather than consumer products. Meta's Project Aria~\cite{engel2023projectaria}\footnote{\url{https://www.projectaria.com/}} provides glasses with RGB, tracking, and eye-tracking cameras, IMUs, and spatial microphones, together with machine-perception services for trajectories, eye gaze, and hand tracking; its second generation (2025) adds a heart-rate sensor and a contact microphone in the nosepad, alongside on-device perception. The platform underlies Aria Digital Twin~\cite{adt}; Ego-Exo4D~\cite{grauman2024egoexo4d}, which pairs egocentric Aria recordings of skilled activities with synchronized exocentric cameras over 1{,}422 hours; and Nymeria~\cite{ma2024nymeria}, which adds full-body ground-truth motion from a motion-capture suit to 300 hours of everyday activity, the kind of supervision the pose reconstruction of Pivotal Task-1 needs. Ego4D~\cite{grauman2022ego4d} gathered its 3{,}670 hours from several head-mounted cameras, and remains the main benchmark for egocentric activity understanding. On the compute side, Qualcomm's Snapdragon AR1+ Gen~1\footnote{\url{https://www.qualcomm.com/products/mobile/snapdragon/xr-vr-ar}} (2025) runs a one-billion-parameter language model on the glasses without a phone, a first step toward the on-device inference that the latency and power constraints of Pivotal Task-3 call for.

\subsection{Challenges}
\label{subsec:p4-sec9-3-challenges}

Wearable sensing is coupled to the dynamics and variability of the human body in a way that external sensing is not, and that coupling is the source of its hardest problems. The four below track the loop of the System Overview: signal quality at the sense stage, robustness under partial observation, fusion across heterogeneous modalities, and generalization across users.

\paragraph{Signal Instability and Drift.}
On-body signals are inherently unstable. IMU drift, EMG noise, and calibration shifts arise continuously from motion artifacts, sensor displacement, and soft-tissue deformation, and they strike precisely the quantities that pose estimation, trajectory tracking, and dynamics reconstruction integrate over time, so small per-step errors compound into large ones across a long recording.

\paragraph{Incomplete and Intermittent Sensing.}
Sensing is rarely continuous in practice. Sweat, motion, and wireless dropouts cause brief losses of contact and stretches of degraded signal, leaving the system to reconstruct reliable motion from partial and sometimes corrupted segment-level data, a demand that clean-benchmark evaluation rarely exercises.

\paragraph{Heterogeneous Multimodal Fusion.}
The modalities that make wearables powerful are also mismatched. IMU, EMG, pressure, and egocentric-visual streams differ in sampling rate, latency, noise, and coordinate frame, and fusing them under real-world variability, without letting one noisy channel corrupt the whole estimate, remains an open integration problem rather than a solved preprocessing step.

\paragraph{Inter-User Variability.}
Finally, the body itself is the confound. Anatomy, muscle-activation patterns, and movement style vary enough between people that a model tuned on a small or homogeneous population transfers poorly, and the effect is sharpest exactly where wearables matter most, in motion reconstruction, intent prediction, and gait analysis, where individual biomechanics dominate the sensor signature. Personalization, domain adaptation, and larger, more diverse training corpora are the levers for closing this gap, and how far they succeed will decide whether wearable spatial intelligence moves from convincing demonstrations to dependable deployment.

\newpage
\section{Aerial Agents}
\label{sec:p4_s9_4}


% AeroVerse: UAV-Agent Benchmark Suite for Simulating, Pre-training, Finetuning, and Evaluating Aerospace Embodied World Models
% Airvista-ii: An agentic system for embodied uavs toward dynamic scene semantic understanding
% Towards Autonomous UAV Visual Object Search in City Space: Benchmark and Agentic Methodology
% Next-Generation LLM for UAV: From Natural Language to Autonomous Flight

\subsection{Introduction}
\label{subsec:p4-sec9-4-introduction}


Aerial agents (satellites, high-altitude platform stations (HAPS), and unmanned aerial vehicles (UAVs)) acquire rich, dynamic, large-scale observations from the sky and near space and, increasingly, perceive, decide, and act on them autonomously through onboard processing and adaptive control (Table~\ref{tab:air-agents}).
The aerial domain is a demanding testbed for the spatial capabilities surveyed in Part~\ref{part:3}: it stresses perception across extreme scale and viewpoint changes (Chapter~\ref{chapter:3:sec:6}) and decision-making under partial observability and tight latency budgets (Chapter~\ref{sec:reasoning}), across platform tiers that differ in altitude, resolution, endurance, and latency by orders of magnitude.

Satellites supply the global tier. Navigation constellations, namely GPS~\cite{enge1994global,leick2015gps} within the broader Global Navigation Satellite System (GNSS), Russia's GLONASS~\cite{ivanov1992glonass}, the European Union's Galileo, and China's BeiDou~\cite{yang2019introduction,yang2020basic}, provide the positioning, navigation, and timing on which vehicles~\cite{li2022real}, smartphones~\cite{zangenehnejad2021gnss}, network synchronization~\cite{hasan2018gnss}, and every lower aerial tier rely.
Earth-observation satellites operate continuously over years, which suits long-term, wide-area monitoring that neither ground sensors nor short-endurance UAVs can provide: seasonal vegetation dynamics~\cite{malingreau1986global}, wildlife migration~\cite{dobrynin2017integration}, river-basin ecosystem health~\cite{ferreira2013estimating}, deforestation~\cite{tucker2000strategies,finer2018combating}, desertification~\cite{del1998status,collado2002satellite}, biodiversity loss~\cite{turner2015free}, and climate indicators such as land surface temperature, ocean color, and atmospheric composition~\cite{yang2013role,kaufman2002satellite}, together with ecological conservation~\cite{pettorelli2014satellite,boyle2014high}, multispectral and hyperspectral resource discovery~\cite{pettorelli2019satellite}, precision agriculture~\cite{guo2018multi} and crop-yield estimation~\cite{maas1988using,ferencz2004crop}, and the detection of algal blooms~\cite{fang2022global} and pollutant leakage~\cite{pandey2019satellite}.

\begin{table}[!ht]
\centering
\caption{Comparison of satellite, HAPS, and UAV platforms.} 
\label{tab:air-agents}
\small
\setlength{\tabcolsep}{2.5mm}
\begin{tabular}{@{}llll@{}}
\toprule
\textbf{Parameter} & \textbf{Satellite} & \textbf{HAPS} & \textbf{UAV} \\
\midrule
Altitude & 500--36,000 km & 20--50 km & < 10 km \\
Spatial Resolution & m to sub-m & dm to cm & cm-level \\
Temporal Coverage & Hours--days revisit & Persistent (regional) & Real-time (local) \\
Coverage Area & Global/regional & Regional & Localized \\
Endurance & Years & Months & Hours \\
Maneuverability & Minimal & Limited & High \\
Data Latency & High & Medium & Low \\
Payload Capacity & High & Moderate & Low \\
Onboard Processing & Limited & Moderate & High \\
Atmospheric Impact & Minimal & Low & High \\
Deployment Speed & Slow & Medium & Fast \\
Cost Structure & High (low per area) & Medium & Low (high per area) \\
%AI Suitability & Large-scale datasets & Regional monitoring & Real-time inference \\
\bottomrule
\end{tabular}
\end{table}


Differently, UAVs, commonly known as drones, have gained widespread popularity due to their agility, low operational cost, and capability for autonomous operation. Equipped with onboard sensors and processing units, UAVs can perceive their surroundings, plan flight paths, and adapt to dynamic conditions in real time. This unique perspective allows UAVs to capture high-resolution imagery and video from flexible vantage points, enabling frequent revisits and reduced occlusions in complex environments. Compared to satellites, UAVs excel in capturing fine-grained, localized data, and can be rapidly deployed to areas of interest, including hazardous or otherwise inaccessible regions. They now serve aerial photography~\cite{zhao2023ms,shao2022efficient}, cinematography~\cite{pitas2022autonomous,dang2022path}, environmental mapping~\cite{zheng2020university,wang2021each,lin2022joint}, precision agriculture~\cite{goodrich2023placement}, geolocalization~\cite{liu2019lending,shi2019spatial,shi2020optimal,shi2020looking,shi2022beyond,shi2022accurate,shi2022cvlnet,shi2023boosting,song2023learning}, and last-mile delivery~\cite{sorbelli2023wind,dissanayaka2023review}.

Filling the altitude gap between satellite systems and UAVs, emerging platforms such as stratospheric balloons, solar-powered high-altitude platform stations, and small satellite constellations (\eg, CubeSats~\cite{lemmer2017propulsion}) are increasingly being integrated into aerial observation networks.
These high-altitude platform stations (HAPS) offer complementary advantages: HAPS~\cite{gonzalo2018capabilities} provide persistent regional coverage with higher spatial resolution than satellites, while orbital systems~\cite{cunningham2004national} enable global, long-term monitoring at varying temporal and spectral scales.
Together, they form a multi-tiered aerial ecosystem capable of capturing heterogeneous multimedia data, from visible and thermal imagery to hyperspectral and LiDAR scans, across spatiotemporal scales unattainable by any single platform alone.
For instance, solar-powered HAPS such as Airbus Zephyr~\cite{kothe2019flight} or HAWK30 by AeroVironment/HAPSMobile~\cite{fladeland2019demonstrating} can operate autonomously in the stratosphere (about 20 km altitude) for weeks, offering persistent observation and communication relay services with sub-meter resolution, which is superior to most Earth observation satellites and at lower cost than maintaining UAV fleets.

A growing trend is the coordination of heterogeneous aerial agents into cooperative networks, where satellites offer global context, HAPS serve as communication relays or edge computing nodes, and UAVs execute localized missions, enabling real-time, closed-loop intelligence. The utility of these systems is further amplified by advances in artificial intelligence, including machine learning, computer vision, big data analytics, and transfer learning~\cite{wang2024multiple,wang2022learning,zheng2022adaptive}. These technologies empower intelligent processing of massive aerial datasets, enabling automated image and video analysis~\cite{liu2023robust,li2021uav,9854892,zheng2019unsupervised,wang2023context,ju2025video2bev}, real-time object detection and tracking~\cite{zhang2019eye,wen2021detection}, predictive modeling for urban dynamics~\cite{fang2022cognitive,shi2022geometry}, natural language-based control interfaces~\cite{chu2024towards}, and city-scale digital twin simulations~\cite{zheng2022citysim}. Data privacy, airspace integration, energy efficiency, and operational robustness remain open, and are taken up in the Challenges below.


\subsection{System Overview}
\label{subsec:p4_s9_4_overview}
An aerial agent realizes a sensing-to-action loop that spans the multi-tier ecosystem of Table~\ref{tab:air-agents}: satellites supply global context and positioning, HAPS act as persistent regional relays and edge-computing nodes, and UAVs execute localized, low-latency missions at the point of interest. Across this stack, spatial intelligence is exercised through three complementary stages, which constitute the pivotal tasks of this section. Cross-view perception anchors aerial and ground observations to a common geographic frame, so that what a drone sees can be matched to a map or a street-level query (Pivotal Task-1); remote-sensing detection and environmental understanding turn that imagery into structured scene knowledge about objects, layouts, and human activity (Pivotal Task-2); and autonomous navigation and decision making convert perception into flight plans and control actions in unseen environments (Pivotal Task-3). Figure~\ref{fig:p4_s9_4_overview} sketches this pipeline, with a feedback path from action back onto sensing that ties the three stages into one closed loop rather than a feed-forward chain.

\begin{figure}[!ht]
    \centering
    \includegraphics[width=0.98\linewidth]{figures/part4/section9/aerial_overview.pdf}
    \caption{Spatial intelligence for aerial agents as a closed loop from multi-tier aerial sensing (satellite, HAPS, UAV, GNSS, onboard sensors), through three complementary core capabilities, to flight plans with 6-DoF pose targets and downstream flight control. The controller executes the plan; subsequent observations close the sensing loop. The core comprises cross-view geolocalization, detection and environmental understanding, and navigation and decision making. Land-use, water, and infrastructure layers are parallel examples. Platform and scene imagery is illustrative.}
    \label{fig:p4_s9_4_overview}
\end{figure}

\subsection{Pivotal Task-1: Ground-Air Collaborative Perception and Visual Geolocalization}
\label{subsec:p4_s9_4_task1}

A foundational requirement of any aerial system is to place its observations on the map: the integration of ground-level and aerial views for cross-modal understanding enables visual geolocalization, determining the geographic location of a query image by matching it against a reference database. This task anchors the perception stage of the system and is particularly challenging due to the large modality gap between ground street-view images and overhead (aerial or satellite) imagery, an instance of the cross-view correspondence problem treated more generally in Chapter~\ref{chapter:3:sec:6}.
For instance, the \textit{University-1652} dataset~\cite{zheng2020university} has become a benchmark for this task, offering ground, drone and satellite image triplets of campus buildings, thus facilitating research on cross-view matching and localization. Such capabilities are critical for urban search-and-rescue, augmented reality, and infrastructure monitoring, where precise alignment between ground observations and aerial maps is essential.
Progress has come from bridging the modality gap, encoding spatial structure into a viewpoint-invariant representation, and exploiting multiple views per location (Figure~\ref{fig:p4_s9_4_geoloc}).



\textbf{Bridging the modality gap.} Most extant research frames geolocalization as a cross-view image retrieval problem, where the core challenge is to learn a viewpoint-invariant representation that effectively bridges the disparity between images captured from different perspectives. Recent advances in deep learning have led to the widespread application of Convolutional Neural Networks (CNNs) and Vision Transformers (ViTs) for extracting robust visual features. One prominent line of inquiry employs metric learning to construct a shared latent space for images acquired from diverse platforms. Workman~\etal~\cite{workman2015location} show that a CNN pre-trained on the Places dataset~\cite{zhou2014learning} can, by itself, yield highly discriminative features without explicit fine-tuning for the task. The integration of contrastive loss, which minimizes the distance between positive pairs in the feature space, further enhances geolocalization performance~\cite{workman2015wide,lin2015learning}. Liu~\etal~\cite{liu2019stochastic} introduce the Stochastic Attraction and Repulsion Embedding (SARE) loss, which minimizes the Kullback-Leibler divergence between the learned and actual data distributions, and more recently Deuser~\etal~\cite{deuser2023sample4geo} show that hard negative sampling strategies also significantly influence the final discriminative power.

\textbf{Encoding spatial structure.} Another significant body of work addresses the problem of spatial misalignment in ground-to-aerial matching, injecting geometric priors into the representation. Vo~\etal evaluate various network architectures and propose an orientation regression loss to train an orientation-aware network~\cite{vo2016localizing}. Zhai~\etal~\cite{zhai2017predicting} utilize semantic segmentation maps to facilitate semantic alignment, while Hu~\etal~\cite{hu2018cvm} incorporate a NetVLAD layer to extract more discriminative features. Subsequently, Liu~\etal~\cite{liu2019lending} propose a Siamese Network to explicitly integrate spatial cues, such as orientation maps, into the training process, and Shi~\etal~\cite{shi2020optimal} introduce a spatial-aware layer to improve localization performance. Given that each location has multiple training images from different viewpoints, Zheng~\etal~\cite{zheng2020university} train a multi-branch CNN as a base model. Building upon this multi-branch architecture, Wang~\etal~\cite{wang2021each} further partition the feature map into multiple segments to improve matching discriminability, while Lin~\etal~\cite{lin2022joint} incorporate keypoint detection into the learning process. More recently, researchers resort to ViT structures: L2LTR~\cite{yang2021cross} adopts a hybrid CNN-ViT architecture, leveraging CNNs to extract high-level semantic features and ViTs to encode them into global representations; GeoDTR~\cite{zhang2023cross} further enhances geometric reasoning by introducing a ViT-based layout extractor that generates geometry-aware attention maps; and TransGeo~\cite{zhu2022transgeo} advances the paradigm with a two-branch ViT model and a two-stage coarse-to-fine learning strategy.

\begin{figure}[!ht]
    \centering
    \includegraphics[width=0.98\linewidth]{figures/part4/section9/aerial_geoloc.pdf}
    \caption{Cross-view visual geolocalization as viewpoint-invariant metric learning: the illustrated shared encoder embeds ground, drone, and satellite views so that same-place (positive) pairs cluster and different-place negatives separate. Architecture variants may instead use view-specific branches. Arrows within the feature-space sketch indicate attraction and repulsion during representation learning, not geographic motion. Synthetic views illustrate a common site rather than a calibrated image triplet.}
    \label{fig:p4_s9_4_geoloc}
\end{figure}

\begin{table}[t]
\centering
\setlength{\tabcolsep}{3pt}
\caption{Comparison of visual geolocalization datasets. R@K: recall at top $K$; AP: average precision; PR: precision-recall.}
\label{table:geo-dataset}
\small
\resizebox{\linewidth}{!}{
\begin{tabular}{lccccccc}
\toprule
Dataset & Univ.-1652~\cite{zheng2020university} & SUES-200~\cite{zhu2023sues} & CVUSA~\cite{zhai2017predicting} & CVACT~\cite{liu2019lending} & Lin~\etal~\cite{lin2015learning} & Tian~\etal~\cite{tian2017cross} & Vo~\etal~\cite{vo2016localizing} \\
\midrule
\#Train imgs. & 701$\times$71.6 & 120 loc. & 35.5k$\times$2 & 35.5k$\times$2 & 37.5k$\times$2 & 15.7k$\times$2 & 900k$\times$2 \\
Platform & Drone, Ground, Sat. & Drone, Sat. & Ground, Sat. & Ground, Sat. & Ground, 45$^{\circ}$ Aerial & Ground, 45$^{\circ}$ Aerial & Ground, Sat. \\
Imgs./loc. & 54+16.6+1 & drone + 1 & 1+1 & 1+1 & 1+1 & 1+1 & 1+1 \\
Target & Building & Diverse & User & User & Building & Building & User \\
Eval. & R@K, AP & R@K, AP & R@K & R@K & PR, AP & PR, AP & R@K \\
\bottomrule
\end{tabular}}
\end{table}

\textbf{From image pairs to multi-view benchmarks.} The maturation of these methods is closely tied to the datasets that score them (see Table~\ref{table:geo-dataset}). The majority of prior geolocalization datasets are based on image pairs, designed for matching images from two distinct platforms such as terrestrial cameras and satellites. One of the earliest contributions~\cite{lin2015learning} leverages public sources to create 78,000 image pairs spanning a 45-degree bird's-eye perspective and the ground. In a similar vein, Tian~\etal~\cite{tian2017cross} compile image pairs for urban localization, positing that buildings are crucial for this task and thus incorporating building detection into their pipeline. The CVUSA~\cite{zhai2017predicting} and CVACT~\cite{liu2019lending} datasets instead address matching panoramic ground-view images with satellite-view images, which is particularly useful for user localization when GPS data is unavailable. A key distinction is one of objective: the latter two focus on localizing the user who captures the photograph, while the former focus on localizing a specific target within the image, so that multiple views of the target are more beneficial for the target-localization datasets (Lin~\etal, Tian~\etal) because they encourage the model to learn the target's structure. However, most existing datasets typically provide only two views of a given location. The University-1652 dataset~\cite{zheng2020university} addresses this by including a greater number of views per target, thereby fostering more robust viewpoint-invariant features; SUES-200~\cite{zhu2023sues} follows a similar multi-view strategy with drone-captured real-world data. Wang~\etal~\cite{wang2024multiple} further extend University-1652 to multi-weather evaluation (rain, fog, and other complex conditions), while Chu~\etal~\cite{chu2024towards} introduce linguistic annotations and a language-based query setting for a more intuitive user experience.

\textbf{Metrics.} These benchmarks are scored by retrieval metrics that reflect their gallery structure. Several datasets~\cite{liu2019lending,zhai2017predicting,vo2016localizing} adopt Recall@K, defined as 1 if a correct match appears within the top K retrieved images; Recall@K is highly sensitive to the rank of the first correct match and is best suited for test sets containing only a single ground-truth match in the gallery. In contrast, datasets such as University-1652~\cite{zheng2020university} and SUES-200~\cite{zhu2023sues} include multiple ground-truth matches from various viewpoints, where Recall@K fails to comprehensively reflect retrieval performance; they therefore also employ Average Precision (AP), as used by Lin~\etal~\cite{lin2015learning} and Tian~\etal~\cite{tian2017cross}. AP, defined as the area under the Precision-Recall (PR) curve, considers all ground-truth images in the gallery, and recent evaluations usually report both Recall@K and the mean Average Precision (mAP) across all queries.


\subsection{Pivotal Task-2: Remote Sensing Detection and Environmental Understanding}
\label{subsec:p4_s9_4_task2}

Once observations are geolocalized, the next stage turns aerial imagery into structured scene knowledge. Aerial platforms are integral to large-scale environmental monitoring and public safety, facilitating fine-grained object detection in complex scenes, a specialization of the detection and segmentation machinery of Chapter~\ref{chapter:3:sec:6} to the aerial viewpoint. Prominent datasets such as \textit{DOTA}~\cite{xia2018dota}, \textit{DOTA-v2}~\cite{xia2018dota,ding2021object}, and \textit{VisDrone}~\cite{du2019visdrone} provide comprehensive annotations for detecting various objects, including vehicles and pedestrians, under a wide range of weather, illumination, and scale conditions. These benchmarks have spurred the development of robust detection models specifically engineered to address challenges inherent to aerial perspectives, such as extreme aspect ratios, dense object clustering, and diminutive object sizes. A significant application is found in UAV-based wildlife monitoring, where systems are deployed for animal census in conservation areas~\cite{vermeulen2013unmanned}, highlighting the utility of aerial technologies in promoting ecological sustainability. Figure~\ref{fig:p4_s9_4_detection} relates the aerial-specific difficulties of orientation, scale, and density to the detectors built for them and to the human-centric understanding that aerial video enables.

\begin{figure}
    \centering
    \includegraphics[width=0.98\linewidth]{figures/part4/section9/aerial_detection.pdf}
    \caption{Remote-sensing detection from the aerial viewpoint: arbitrary orientation, extreme scale variation, and dense small objects (top) motivate (1) oriented-object detection and (2) multi-scale feature learning, which provide complementary representations for human-centric understanding (bottom). Person re-identification matches identities across views; action examples are alternative classes, not a prescribed sequence; tracking maintains identities across frames. Synthetic scenes illustrate these tasks.}
    \label{fig:p4_s9_4_detection}
\end{figure}

Beyond object detection, as shown in Table~\ref{tab:uav_reid_datasets}, with the proliferation of UAVs in smart-city and other human-centric applications, human behavior understanding~\cite{li2021uav} has emerged as a critical and formidable research challenge, including person re-identification and action recognition. For instance, the UAV person re-identification (re-ID) task~\cite{albaluchi2025uav} is designed to address cross-drone identity matching amid substantial viewpoint and illumination variations, thereby motivating the creation of view-invariant feature representations. In conjunction with resources like \textit{VisDrone}, which contains multi-target tracking sequences from mobile UAVs, these datasets facilitate the design of intelligent video analytics systems capable of performing long-term human behavior analysis and anomaly detection in dynamic environments.


\begin{table}[t]
\centering
\caption{Comparison of UAV-based human and object datasets.}
\label{tab:uav_reid_datasets}
\fontsize{8}{10}\selectfont
\setlength{\tabcolsep}{1.0mm}
\begin{tabular}{L{0.15\textwidth} L{0.26\textwidth} L{0.25\textwidth} L{0.25\textwidth}}
    \toprule
    \textbf{Feature} & \textbf{UAV-Human}~\cite{li2021uav} & \textbf{DOTA-v2}~\cite{xia2018dota,ding2021object} & \textbf{VisDrone}~\cite{du2019visdrone} \\
    \midrule

    \textbf{Primary Focus} & Human Behavior Understanding & Object Detection & Object Detection / Tracking / Crowd Counting \\

    \textbf{Data Type} &  Videos \& Images & Images & Videos \& Images \\

    \textbf{\#Images} & 22,476 frames for pose estimation, 41,290 frames for re-ID, 22,263 frames for attribute recognition & 11,268 & 10,209 \\

    \textbf{\#Videos} & 67,428  & N/A & 288 \\

    \textbf{\#Categories} & 155 action categories & 18 object categories & 10 object categories \\

    \textbf{\#Instances} & 119 subjects for action recognition, 1,144 identities for re-ID & 1,793,658 & > 2.6 million annotated bounding boxes \\

    \textbf{Sensor Modalities} & RGB, Depth, Infrared (IR), Fisheye, Night-Vision & RGB (from Google Earth, aerial) \& Grayscale (from satellites) & RGB \\

    \textbf{Image Resolution} & 1920$\times$1080 (RGB), 640$\times$576 (Depth,IR), 640$\times$480(Fisheye, Night-Vision), & 800$\times$800 to 20,000$\times$20,000 pixels & 2000$\times$1500 (image),  3840$\times$2160 (video) \\

    \textbf{Annotation Format} & Bounding boxes, 17 keypoints for pose & Oriented Bounding Boxes (Quadrilaterals) & Bounding boxes with attributes (\eg, occlusion) \\

    \textbf{Scene Diversity} & Urban and rural, daytime and nighttime & Wide area, diverse landscapes (ports, airports, cities, etc.) & 14 different cities, urban and rural, various weather and lighting \\

    \textbf{Platform} & Drone (DJI Matrice 100) & Satellite, Aerial Platforms & Drone \\

    %\textbf{Unique Features} & Rich multi-modal data for fine-grained human-centric analysis. & Very high-resolution images with oriented bounding boxes for rotated objects. & A comprehensive benchmark with multiple tailored tasks for drone vision. \\
    \bottomrule
  \end{tabular}
\end{table}


Remote sensing object detection is primarily challenged by the prevalence of arbitrarily oriented objects and substantial scale variations~\cite{xia2018dota, ding2021object, sun2022fair1m, liu2016ship, cheng2022anchor, yao2023automated}, and existing research can be broadly organized along two corresponding fronts.
\textbf{Detection of arbitrarily oriented objects.} A significant line of work focuses on developing specialized oriented object detection (OOD) methods. Early approaches often augment horizontal bounding boxes (HBBs) with an angle parameter, employing techniques such as feature refinement in the detector neck~\cite{yang2019scrdet,yang2021r3det}, extraction of rotated regions of interest (RoIs)~\cite{ding2019learning,xie2021oriented}, or designing specific OOD heads~\cite{hou2022shape,han2021redet,ming2021dynamic}. However, these methods frequently suffer from boundary discontinuity and representation ambiguity due to their reliance on angle parameters. To mitigate these issues, subsequent research has explored more flexible \textit{object representations}: some methods describe oriented objects using points~\cite{fu2020point, li2022oriented} or gliding vertexes~\cite{xu2020gliding}, while others model them with Gaussian distributions~\cite{yang2021rethinking, yang2021learning, cheng2022dual} and design corresponding loss functions~\cite{qian2021learning}.
\textbf{Handling substantial scale variations.} Another research direction designs specialized backbones and feature extraction strategies to cope with large scale variations in remote sensing images (RSIs). A common strategy is multi-scale feature learning, achieved through data augmentation~\cite{zhao2019multi, shamsolmoali2021rotation, chen2020stitcher}, multi-scale feature integration~\cite{zhang2019hierarchical,zheng2020hynet,9382268,liu2016ssd}, enhanced feature pyramid networks (FPNs)~\cite{lin2017feature,hou2022refined,zhang2021laplacian,guo2020rotational}, or multi-scale anchor generation~\cite{guo2018geospatial, hou2021self, qiu2019a2rmnet}. Other works focus on ensuring equivalent receptive fields for objects of different orientations~\cite{han2021redet,pu2023adaptive}. More recently, some methods~\cite{Li_2023_ICCV} have employed large kernel convolutions~\cite{liu2022convnet,ding2022scaling,liu2022more} to enlarge the spatial receptive field for large objects, though this can introduce background noise for smaller ones, while others harness multi-scale convolution kernels and attention to address scale variance~\cite{guo2022segnext,dai2023multi,zhang2022multiscale,cai2024poly}.

These method fronts are evaluated on a complementary set of benchmarks and metrics. The VisDrone dataset~\cite{du2019visdrone} provides a comprehensive benchmark for drone-captured imagery, featuring urban and rural scenes with challenges like varying altitudes and occlusions; DOTA-v2~\cite{xia2018dota,ding2021object} focuses on oriented object detection in high-resolution satellite images, emphasizing rotated bounding boxes for objects like planes and ships; and UAV-Human~\cite{li2021uav} targets human behavior understanding from UAVs, with multi-modal sequences for pose and action analysis. Complementary datasets, such as UAVDT~\cite{du2018uavdt} for traffic surveillance and iSAID~\cite{waqas2019isaid} for instance segmentation, broaden the scope of aerial detection tasks and collectively enable cross-dataset fine-tuning to mitigate domain gaps. Evaluation extends standard metrics to address orientation, scale, and density: the primary metric, mean Average Precision (mAP), is computed as the area under the precision-recall curve, typically at an Intersection over Union (IoU) threshold of 0.5 (mAP@0.5) or across 0.5:0.95 for stricter assessment~\cite{lin2014microsoft}, and for oriented detection, as in DOTA-v2, rotated IoU variants ensure alignment with object headings~\cite{xia2018dota}. Per-class and per-size Average Precision exposes imbalances, and small objects ($<32\times32$ pixels) remain the main source of error on drone imagery~\cite{du2019visdrone}. Additional metrics include precision, recall, and F1-score to assess false positives (\eg, background clutter) and misses (\eg, occlusions), while Frames Per Second (FPS) is critical for real-time UAV deployment, targeting 30+ FPS~\cite{yang2023laiyolo}. Video-specific metrics like Multiple Object Tracking Accuracy (MOTA) are used for dynamic scenes in datasets like UAV-Human~\cite{li2021uav}. Modern real-time detectors of the YOLO family~\cite{wang2024yolov9} are increasingly adapted to the aerial setting, though aggregated aerial mAP still trails their COCO-level performance owing to small-object and orientation hardness, keeping open gaps such as multi-modal fusion and zero-shot adaptation active.


\begin{figure}[!ht]
    \centering
    \includegraphics[width=0.98\linewidth]{figures/part4/section9/aerial_navigation.pdf}
    \caption{Autonomous aerial navigation from a language instruction: a UAV observation is grounded into a semantic-topometric graph; the instruction and graph condition an LLM planner that produces pose waypoints with velocity. The displayed route is a top-down projection, with attitude and altitude encoded in the waypoint state. Multi-drone planning is an extension with coordinated altitude separation. SR: success rate; SPL: success weighted by path length; NE: navigation error. Metric glyphs indicate evaluation criteria, not measured scores; scenes and graph geometry are schematic.}
    \label{fig:p4_s9_4_navigation}
\end{figure}


\subsection{Pivotal Task-3: Autonomous Aerial Navigation and Decision Making}
\label{subsec:p4_s9_4_task3}

The final stage closes the loop, converting perception into action. Traditional UAV geolocalization addresses basic aerial navigation by matching a query (\eg, a satellite image or textual description) to a database of historical flight paths or geographic points~\cite{zheng2020university,chu2024towards}. However, a critical distinction exists between such perception-oriented tasks and the overarching goal of autonomous navigation: the latter requires complex decision-making and control to bridge perception to action, drawing on the spatial-reasoning and planning paradigms of Chapter~\ref{sec:reasoning}. In practical scenarios, the objective is often for an aerial agent to autonomously locate a destination, particularly an unexplored target, through environmental exploration and, when available, collaboration with a map. This goal introduces significant challenges in generalization to unseen environments, including novel objects, unmapped regions, and interaction with the environment. For instance, recent benchmarks adapted from general Vision-Language Navigation (VLN) have begun to address this by requiring models to output a detailed flight plan with 6-DoF control signals~\cite{wang2024towards}. Figure~\ref{fig:p4_s9_4_navigation} follows the problem from the perception-to-action gap to foundation-model-driven VLN and multi-agent extensions.



\begin{table}[!ht]
\small
\centering
\caption{Comparison of recent aerial VLN datasets.}
\label{tab:aerial_datasets}
\begin{tabular}{@{}p{2.5cm}lp{2.5cm}p{2cm}p{3cm}p{3.5cm}@{}}
\toprule
\textbf{Dataset} & \textbf{Yr} & \textbf{Env.} & \textbf{Size} & \textbf{Modalities} & \textbf{Key Features} \\
\midrule
UAVA~\cite{albanis2020dronepose} & 2020 & Indoor simulated (MP3D) & 90 traj. & RGB, depth, flow, ego/exo views & Photorealistic UAV assist; vision/robotics tasks; gamified traj. \\
AVDN~\cite{fan2022aerial} & 2022 & Outdoor simulated (xView) & 3,064 traj. with dialog & RGB, dialog & Async human dialogs; attn. for interactive nav. \\
AerialVLN~\cite{liu2023aerialvln} & 2023 & Outdoor simulated (cities) & 8,446 traj. and 25,338 instructions & RGB, lang. instr. & Synthetic VLN; landmark nav in urban scenes. \\
CityNav~\cite{lee2024citynav} & 2024 & Real-world (UK cities) & 32,637 traj./instr. & 3D point clouds, RGB, geo-lang & Real 3D scans (4.65 km²); human demos; scalable VLN. \\
UAV-Need-Help~\cite{wang2024towards} & 2024 & Outdoor simulated (OpenUAV) & 12,149 traj. & RGB, lang. instr. & Continuous 6-DoF traj.; target-oriented. \\
OpenFly~\cite{gao2025openfly} & 2025 & Outdoor real/sim (18 scenes) & 100k traj. & 3D point clouds, RGB, lang. instr. & Large-scale; real-to-sim transfer. \\
AirV2X~\cite{gao2025airv2x} & 2025 & Outdoor simulated (CARLA+AirSim) & 6.73 hours of UAV-ground traj. & 3D point clouds, RGB & V2X collab; drone strategies; reduces occlusions. \\
\bottomrule
\end{tabular}
\end{table}


\textbf{From visual perception to language-guided planning.} Recent years have witnessed significant progress in UAV navigation, with a primary focus on visual perception~\cite{loquercio2018dronet,kang2019generalization,giusti2015machine,smolyanskiy2017toward,fan2020learn,bozcan2020air,majdik2017zurich} and collision avoidance~\cite{xu2025navrl,singla2019memory}. Building on this aerial-perception foundation, Vision-and-Language Navigation (VLN) for UAVs remains an emerging area of research. To address it, recent efforts have leveraged Large Multi-modal Models (LMMs) to harness general knowledge for navigation planning: STMR~\cite{gao2024aerial} builds upon AerialVLN~\cite{liu2023aerialvln} with a zero-shot LLM-based framework utilizing a Semantic-TopoMetric Representation for enhanced spatial reasoning. A parallel line reframes UAV autonomy as a full agentic loop rather than one-shot instruction following. AeroVerse~\cite{yao2024aeroverse} assembles a simulator, real-virtual pre-training corpora, and five downstream tasks, \ie, scene awareness, spatial reasoning, navigational exploration, task planning, and motion decision, to train aerospace embodied world models, and proposes SkyAgentX as a unified perception-reasoning-navigation-planning agent; AirVista-II~\cite{lin2025airvista} pursues zero-shot dynamic-scene semantic understanding via agent-based task scheduling and adaptive keyframe selection; and PRPSearcher~\cite{ji2025uavsearch} casts aerial visual object search over the CityAVOS benchmark as a perception-reasoning-planning problem solved by a multimodal LLM agent. Complementary systems such as NeLV~\cite{yuan2025nelv} operationalize the language-to-flight pipeline end to end, chaining an LLM parser, route and path planners, and a control platform across short-, medium-, and long-range missions. Beyond single-agent setups, some works have begun to explore multi-drone systems; for instance, Wu~\etal~\cite{wu2025aeroduo} deploy two drones flying at different altitudes to facilitate planning.

\textbf{Datasets and simulators.} As shown in Table~\ref{tab:aerial_datasets}, several foundational frameworks and datasets have been established. One of the early works, \ie, UAVA~\cite{albanis2020dronepose}, leverages the panorama images in the Matterport3D dataset~\cite{chang2017matterport3d} to render the indoor environment, while harnessing the camera viewpoints as physically available flight anchors to generate trajectories. AerialVLN~\cite{liu2023aerialvln} introduces a large-scale dataset with an effective baseline, and further infrastructure contributions include datasets such as CityNav~\cite{lee2024citynav}, AVDN~\cite{fan2022aerial}, and OpenFly~\cite{gao2025openfly}. Notably, UAV-Need-Help~\cite{wang2024towards} provides a realistic UAV dynamics simulator, \ie, OpenUAV, alongside its evaluation protocol; OpenUAV is further used in HaL-13k~\cite{wu2025aeroduo} to simulate two drones at different altitudes. Additional observation agents, such as road infrastructures and vehicles with depth sensing, have also been incorporated to enrich environmental understanding~\cite{gao2025airv2x}.

\textbf{Metrics.} Evaluation metrics for UAV aerial navigation, especially VLN, measure how well a model navigates a drone based on visual inputs and language instructions. Common metrics include Success Rate (SR), which checks if the UAV reaches the target location within a specified distance, and Success weighted by Path Length (SPL), which balances success with how efficiently the drone travels (shorter paths score higher). Path Length (PL) measures the total distance traveled, while Navigation Error (NE) calculates the distance between the final position and the goal. Together these metrics assess both accuracy (reaching the target) and efficiency (taking an optimal route), providing a clear way to compare model performance across datasets like UAVA, AerialVLN, and CityNav.


\subsection{Established Tools and Systems}
\label{subsec:p4-sec9-4-established-tools-and-systems}
Beyond the methods and benchmarks above, the aerial domain rests on a real-world layer of simulators, platforms, and infrastructure that turns research prototypes into deployable systems; we group it into three tiers.

\textbf{Simulation environments and benchmarks.} Because flight data is costly and safety-critical to collect, high-fidelity simulation underpins much of the field. The OpenUAV simulator released with UAV-Need-Help~\cite{wang2024towards} supplies realistic UAV dynamics and a continuous 6-DoF evaluation protocol, and is reused to simulate multi-altitude drone teams in HaL-13k~\cite{wu2025aeroduo}; OpenFly~\cite{gao2025openfly} targets large-scale real-to-sim transfer across diverse scenes; and AerialVLN~\cite{liu2023aerialvln} pairs a city-scale synthetic environment with a navigation baseline. Photorealistic indoor rendering for UAV tasks is provided through Matterport3D~\cite{chang2017matterport3d} in UAVA~\cite{albanis2020dronepose}, while V2X-style multi-agent simulation combines CARLA and Microsoft AirSim\footnote{\url{https://microsoft.github.io/AirSim/}} in AirV2X~\cite{gao2025airv2x}.

\textbf{Commercial platforms and vendors.} On the hardware side, consumer and enterprise drones from DJI\footnote{\url{https://www.dji.com/}} (whose Matrice series collects datasets such as UAV-Human) and autonomy-focused vendors such as Skydio\footnote{\url{https://www.skydio.com/}} dominate the UAV tier. In the stratospheric tier, solar-powered HAPS such as Airbus Zephyr~\cite{kothe2019flight} and AeroVironment/HAPSMobile's HAWK30~\cite{fladeland2019demonstrating} provide persistent regional coverage, complemented by small-satellite constellations built on CubeSat platforms~\cite{lemmer2017propulsion} and broader HAPS programs~\cite{gonzalo2018capabilities}.

\textbf{Connectivity and positioning infrastructure.} Aerial autonomy ultimately depends on shared positioning and communication backbones: the GNSS constellations introduced above (GPS~\cite{enge1994global}, GLONASS~\cite{ivanov1992glonass}, Galileo, and BeiDou~\cite{yang2019introduction}) supply global timing and localization, while space-based broadband relays such as Starlink~\cite{michel2022first} and orbital observation systems~\cite{cunningham2004national} increasingly serve as the connectivity layer for coordinated, real-time aerial operations.


\subsection{Challenges}
\label{subsec:p4-sec9-4-challenges}
Aerial autonomy is limited less by any single model than by the gaps between platforms, scales, and operating conditions. Four threads follow the pipeline in order: fusing heterogeneous platforms into a coherent situational model, perceiving small objects and complex activities reliably, navigating unseen and GPS-denied environments, and meeting the ethical and operational constraints under which any of this can be deployed.

\paragraph{Heterogeneous-Platform Fusion.}
A primary obstacle lies in integrating aerial agents with heterogeneous platforms for robust environmental perception. Coordinating a UAV with ground-based sensors or satellite systems to construct a cohesive situational-awareness model remains complex, particularly in dynamic or cluttered settings where the platforms differ in resolution, latency, and viewpoint (Table~\ref{tab:air-agents}). Reconciling these asynchronous, multi-resolution streams into a single consistent representation, without one modality's noise corrupting the fused estimate, is the standing difficulty, and it motivates the V2X-style cooperative-perception efforts that incorporate road infrastructure and ground vehicles as additional observers.

\paragraph{Small-Object and Complex-Activity Understanding.}
Aerial perception is defined by extremes of scale and density. The recognition of small objects, such as identifying debris in disaster zones or distinguishing wildlife in ecological surveys, demands high-resolution imaging and detectors robust to dense clustering, which often strain onboard computational resources. Understanding complex activities, such as discerning human behavior during search-and-rescue missions, compounds the problem by requiring nuanced spatio-temporal and contextual analysis that current systems struggle to deliver consistently. Closing this gap calls for detectors and activity models that remain accurate at extreme scale variation while fitting a UAV's tight compute and power budget.

\paragraph{Navigation in Unseen and GPS-denied Environments.}
Autonomous operation in unseen or GPS-denied environments, such as dense urban canyons or subterranean spaces, further complicates the action stage, as agents must rely on real-time mapping and obstacle avoidance without pre-existing data. The core requirement is generalization: a navigation policy must extend to novel objects, unmapped regions, and genuinely interactive scenarios, rather than replaying matched trajectories. This is precisely the gap that foundation-model-driven VLN and multi-altitude cooperative planning aim to narrow, but validating such behavior to a safety standard in the open world remains open.

\paragraph{Privacy, Airspace, and Operational Robustness.}
Beyond perception and navigation, several ethical and operational challenges persist. Data-privacy concerns arise when aerial agents capture high-fidelity multimedia data, potentially compromising individual rights without stringent safeguards, motivating privacy-preserving approaches such as federated learning. Airspace integration poses another hurdle, as regulatory frameworks struggle to accommodate the proliferation of aerial agents alongside traditional aviation, necessitating harmonized protocols to ensure safety. Energy efficiency remains a critical bottleneck, since prolonged missions demand lightweight yet powerful power sources and motivate renewable solutions such as solar-powered propulsion. Finally, operational robustness is challenged by environmental uncertainties, such as extreme weather or electromagnetic interference, that can disrupt system reliability. Advances in artificial intelligence and sensor fusion that address these constraints, rather than perception accuracy alone, are what will let aerial agents move from bounded demonstrations toward dependable, large-scale deployment.


\newpage

\section{Warehouse and Industrial Logistics}
\label{sec:p4_s9_5}



\subsection{Introduction}
\label{subsec:p4-sec9-5-introduction}

Warehouse and industrial logistics is one of the most mature real-world testbeds for the spatial capabilities surveyed in Part~\ref{part:3}: it demands perception-to-action pipelines that interpret complex, dynamic environments such as dense shelving systems populated with diverse stock keeping units (SKUs), pallets, forklifts, and human workers. 
These systems construct and maintain spatial memory representations (semantic maps, object and pose states) to support downstream planning and manipulation for tasks including inventory counting, order picking, pallet handling, and coordinated fleet operations. 
Perception and reasoning themselves are treated in Chapter~\ref{chapter:3:sec:6} and Chapter~\ref{sec:reasoning}; logistics adds throughput targets, engineered but crowded layouts, and people sharing the floor with fleets of robots. 
Performance is judged in terms of throughput, order accuracy, safety, and operational cost, especially under high SKU diversity, visual clutter, elongated aisles, and human-robot collaboration. 
Dense shelving and manipulation in clutter remain the hardest cases, as the Amazon Picking Challenge made plain~\cite{Correll2016AmazonChallenge}, and they demand that perception, pose estimation, mapping, and multi-agent planning be integrated rather than tuned in isolation.

\begin{figure}[!ht]
    \centering
    \includegraphics[width=0.98\linewidth]{figures/part4/section9/warehouse_overview.pdf}
    \caption{Spatial intelligence for warehouse and industrial logistics: RGB-D, LiDAR, and RFID observations support five complementary capabilities for picking and pallet transport, and execution outcomes update maps and plans. Scene images and diagrams are illustrative.}
    \label{fig:p4_s9_5_overview}
\end{figure}

\subsection{System Overview}
\label{subsec:p4-sec9-5-system-overview}

The complete warehouse automation pipeline, sketched in Figure~\ref{fig:p4_s9_5_overview}, integrates five complementary capabilities into a closed loop. 
\textbf{Perception} converts raw multi-sensor streams into semantic object information (what is present, where, and in what pose); 
\textbf{localization and mapping} grounds that information in a metric spatial map and keeps it consistent over time; 
\textbf{reasoning and planning} consume these representations to allocate tasks and generate optimized fleet schedules; 
\textbf{manipulation and control} physically execute the resulting picking and transport actions; 
and a \textbf{human-robot collaboration and safety} layer adapts behavior dynamically in mixed human-robot environments. 
Integration is achieved via middleware such as ROS\,2 and industrial protocols like Open Platform Communications Unified Architecture (OPC-UA), while digital twin platforms such as NVIDIA Omniverse\footnote{https://www.nvidia.com/en-gb/omniverse/} support simulation and validation before deployment. 
The five capabilities are the pivotal tasks below, taken in pipeline order.



\subsection{Pivotal Task-1: Perception and Recognition}
\label{subsec:p4-sec9-5-pivotal-task-1-perception-and-recognition}
Perception constitutes the sensory foundation of the pipeline: it transforms raw sensor data into actionable knowledge, reusing the detection, segmentation, and pose-estimation machinery of Chapter~\ref{chapter:3:sec:6} but specializing it to the warehouse Operational Design Domain. 
Core capabilities include detecting and classifying objects (packages, pallets, tools), estimating their 3D poses, and tracking their motion and state over time. 
Unlike outdoor environments, warehouses impose distinct perceptual challenges: dense clutter, heavy occlusion, visually similar items, reflective packaging, narrow aisles, and continually changing layouts, which together make robust perception a central and ongoing research focus \cite{Tubis2023IntelligentWarehouse,redmon2016yolo}. 
Formally, perception is a structured mapping from a synchronized multimodal stream $X_t = \{I_t^{\text{RGB}}, I_t^{\text{Depth}}, P_t^{\text{LiDAR}}, R_t^{\text{RFID}}\}$ (together with calibration parameters and a prior map $M$) to a scene graph $G_t = \{o_i, c_i, \mathbf{b}_i, \mathbf{p}_i, \Sigma_i\}$, in which each detected object $o_i$ carries a class label $c_i$, a bounding box or segmentation mask $\mathbf{b}_i$, an estimated pose $\mathbf{p}_i$, and an uncertainty measure $\Sigma_i$ that grounds subsequent manipulation, navigation, and decision making. 
Figure~\ref{fig:p4_s9_5_perception} sketches this representation stack.

\begin{figure}[!ht]
    \centering
    \includegraphics[width=0.98\linewidth]{figures/part4/section9/warehouse_perception.pdf}
    \caption{The warehouse perception stack: 2D detection and segmentation, 6D object pose $T_i$, and a structured scene graph whose nodes carry identity, class, box or mask, pose, and uncertainty $\Sigma_i$. RFID can support identity association; masks and axes are schematic, not calibrated estimates.}
    \label{fig:p4_s9_5_perception}
\end{figure}

\textbf{Detection and recognition} form the base layer. 
Convolutional and transformer-based architectures, \ie, YOLO \cite{redmon2016yolo}, RetinaNet \cite{lin2017focal}, and the transformer-based DETR with deformable attention \cite{Zhu2021DeformableDETR}, have steadily improved accuracy and efficiency under clutter, and are scored on COCO \cite{lin2014microsoft}, whose mean Average Precision (mAP) broken down by object size is particularly informative for the small-SKU and barcode-label cases that dominate warehouse failure modes. 
\textbf{Pose estimation} lifts these 2D detections into the metric space that manipulation requires: PoseCNN \cite{xiang2017posecnn}, DenseFusion \cite{Wang2019DenseFusion}, and CosyPose \cite{Labbe2020CosyPose} fuse visual and geometric cues to infer six-degree-of-freedom (6D) orientation and position even under occlusion, and are evaluated on YCB-Video \cite{xiang2017posecnn} for textured household and industrial objects and on T-LESS \cite{Hodan2017TLESS} for textureless, industry-relevant parts, with Average Distance (ADD) the standard accuracy measure and MOTA/MOTP reporting tracking quality in dynamic scenes. 
The Amazon Picking Challenge \cite{Correll2016AmazonChallenge} couples perception with manipulation in real warehouse scenarios, stressing generalization across novel object sets and thereby anticipating the manipulation task of §\ref{subsec:p4_s9_5_manip} below. 
Representative benchmarks across the section are consolidated in Table~\ref{tab:p4_s9_5_benchmarks}. In industrial practice, perception pipelines are increasingly standardized and modularized; the deployed software stacks that realize them are surveyed in the Established Tools and Systems subsection.

\subsection{Pivotal Task-2: Localization and Mapping}
\label{subsec:p4-sec9-5-pivotal-task-2-localization-and-mapping}
If perception answers \emph{what is around the robot}, localization and mapping answer \emph{where the robot is} and \emph{what the persistent structure of the warehouse is}, forming the spatial backbone on which fleet coordination rests and drawing on the same scene-construction techniques as Chapter~\ref{chapter:3:sec:6}. 
Automated guided vehicles (AGVs) and autonomous mobile robots (AMRs) require centimeter-level accuracy to navigate narrow aisles, dock precisely at shelves, and operate safely alongside workers. 
The problem is conceptually the SLAM problem of autonomous driving, yet it faces warehouse-specific constraints: highly repetitive structures (rows of identical shelves), dynamic obstacles (workers, pallets, forklifts), and low-light or GPS-denied indoor conditions. 
Formally, a sequence of odometry and sensor measurements $Z_{1:t} = \{\text{IMU}, \text{LiDAR}, \text{cameras}, \text{wheel encoders}, \text{UWB beacons}\}$, optionally supplemented by a prior map $M_0$ or fiducial anchors, is fused into a pose graph $\mathcal{X}_t = \{x_1, \ldots, x_t\}$ that estimates the trajectory and an evolving semantic map $M_t$ fusing geometric structure with storage zones and restricted areas. 
Beyond navigation, this map underpins the digital twins (spatial models synchronized with the facility whose simulations feed back into its operation; §\ref{sec:p4_s11_1}) that increasingly govern modern facilities \cite{Cadena2016SLAMSurvey}. Figure~\ref{fig:p4_s9_5_localization} casts the task as a multi-sensor fusion problem.

\begin{figure}
    \centering
    \includegraphics[width=0.98\linewidth]{figures/part4/section9/warehouse_localization.pdf}
    \caption{Localization and mapping in the warehouse as a factor-graph fusion problem: IMU, LiDAR, camera, wheel-encoder, and UWB measurements constrain pose states, optionally with a prior map or beacon anchors, producing poses over warehouse geometry, storage zones, and restricted areas. Repetitive shelving can cause false place matches, so only verified loop candidates add constraints.}
    \label{fig:p4_s9_5_localization}
\end{figure}

A family of \textbf{SLAM algorithms} has been tailored to these conditions. ORB-SLAM3 \cite{Campos2021ORBSLAM3} extends visual SLAM with inertial and multi-camera support for accurate visual-inertial localization; LIO-SAM \cite{Shan2020LIOSAM} provides real-time LiDAR--IMU fusion through factor-graph optimization that stays robust under fast motion and partial occlusion; and RTAB-Map \cite{Labbe2019RTABMap} supports online RGB-D mapping and loop-closure detection with memory management suited to long-term operation, all expressible within the unified factor-graph framework of GTSAM \cite{Dellaert2012GTSAM}. 
\textbf{Multi-modal fusion} is the common thread that holds accuracy together when any single modality degrades, and is also the dominant emerging research direction, alongside learning-based estimation (neural depth prediction, differentiable bundle adjustment, and learned place-recognition embeddings) and the migration of SLAM computation to the edge for fleet-scale latency. 
These capabilities are evaluated on KITTI \cite{Geiger2013KITTI} for outdoor visual and LiDAR odometry and loop closure, on the indoor visual-inertial EuRoC MAV dataset \cite{Burri2016EuRoC} whose challenging motion and illumination mirror warehouse conditions, and on the high-fidelity Replica reconstructions \cite{straub2019replica} for dense mapping and semantic labeling, with Absolute and Relative Pose Error (APE/RPE) and loop-closure precision/recall the standard metrics and, for warehouse use, semantic-zone IoU and map completeness (Table~\ref{tab:p4_s9_5_benchmarks}).


\begin{figure}[!ht]
    \centering
    \includegraphics[width=0.98\linewidth]{figures/part4/section9/warehouse_planning.pdf}
    \caption{Two families for warehouse reasoning and planning over shared task, fleet, and map inputs: combinatorial optimization (assignment and routing) and learning-based coordination (per-agent policies), optionally combined in hierarchical or hybrid designs. Allocations, paths, and schedules are illustrative.}
    \label{fig:p4_s9_5_planning}
\end{figure}



\subsection{Pivotal Task-3: Reasoning and Planning}
\label{subsec:p4-sec9-5-pivotal-task-3-reasoning-and-planning}
Reasoning and planning form the decision-making core of warehouse automation, transforming the perceptual and localization outputs into task assignments and motion plans, and instantiating the spatial-reasoning and planning paradigms of Chapter~\ref{sec:reasoning}. 
This layer determines which robot executes which task, how tasks are sequenced, and how routes and schedules are optimized for safe, efficient, coordinated operation. 
In contrast to autonomous driving, where planning typically focuses on individual vehicles, warehouse reasoning must address multi-agent coordination across fleets of dozens or hundreds of robots in constrained indoor space \cite{Veloso2015CMUCoBot}, maximizing throughput and minimizing latency while maintaining safety under dynamic obstacles, human interactions, and stochastic task arrivals. 
Formally, the perceived states $G_t$, a task list $\mathcal{T}$, the fleet state $F_t$ (positions, capabilities, battery levels), and the warehouse map $M_t$ are mapped to a task-allocation matrix $A$, optimized motion paths $\mathcal{R}$, and temporal schedules $\pi$ that jointly specify when and how tasks execute. Figure~\ref{fig:p4_s9_5_planning} contrasts the two solution families.


The first family is \textbf{combinatorial optimization}. Traditional methods (the Hungarian algorithm, branch-and-bound, and mixed-integer linear programming) remain widely used for task allocation and scheduling in small- to medium-scale deployments, while Vehicle Routing Problem (VRP) formulations \cite{Cordeau2007VRP} and their multi-depot and time-window variants are central to modeling fleet movement. The second family is \textbf{learning-based coordination}: reinforcement learning (RL) and multi-agent RL (MARL) handle dynamic environments and real-time uncertainty \cite{Sartoretti2019PRIMAL,samvelyan2019smac}, with toolkits such as Google OR-Tools \cite{Perron2019ORTools} and PyMARL enabling rapid experimentation across both families. 
Emerging research couples the two: hierarchical and hybrid frameworks combine symbolic reasoning with learned components, graph neural networks and differentiable optimization accelerate multi-agent coordination and transfer policies across layouts, and digital-twin integration enables simulation-informed, predictive re-planning, increasingly under multi-objective criteria that balance throughput, energy, and human safety. 
Planning systems are benchmarked on Solomon's VRP suite \cite{Solomon1987VRP} for static routing and scheduling and on the Open Motion Planning Library (OMPL) \cite{Sucan2012OMPL} for continuous-space motion planning, complemented by synthetic multi-robot testbeds; common metrics are throughput (items/hour), task-completion latency, total travel distance or time, collision rate, and re-planning latency (Table~\ref{tab:p4_s9_5_benchmarks}).

\subsection{Pivotal Task-4: Manipulation and Control}
\label{subsec:p4_s9_5_manip}
Manipulation constitutes the physical execution layer, enabling robots to interact directly with objects, containers, and infrastructure such as bins, shelves, and conveyors. 
In industrial logistics this typically involves robotic arms performing grasping, placing, sorting, or palletizing, as well as AGVs docking to pallets or workstations. Manipulation bridges high-level symbolic or geometric plans with continuous motor control, and therefore demands both precision and robustness to the inherent variability in packaging, geometry, and material properties \cite{Levine2016DeepGrasp}; the grasping and control techniques here connect directly to the robotic-manipulation treatment of §\ref{sec:p4_s9_1} and the planning paradigms of Chapter~\ref{sec:reasoning}. 
Formally, estimated object poses $\mathbf{p}_i$ and task descriptions (grasp targets, placement goals, trajectory constraints) are mapped to low-level motor commands $\mathbf{u}_t$ that ensure stable grasping, precise placement, and compliant, collision-free motion. Figure~\ref{fig:p4_s9_5_manipulation} contrasts modular grasp-and-plan systems with learned visuomotor control.

\begin{figure}
    \centering
    \includegraphics[width=0.98\linewidth]{figures/part4/section9/warehouse_manipulation.pdf}
    \caption{Warehouse manipulation through (a) a modular route, where grasp candidates scored by analytic metrics and learned quality (Dex-Net) feed sampling-based planning (RRT*) and trajectory tracking, and (b) learned visuomotor control from camera observations. Action spaces vary across systems, and the routes can be combined.}
    \label{fig:p4_s9_5_manipulation}
\end{figure}

\textbf{Grasp synthesis} anchors the methodology: Dex-Net combines synthetic grasp data and analytic robustness metrics to train a grasp-quality predictor for candidate selection \cite{Mahler2017DexNet}, while learned visual control can predict grasp success from camera images and candidate task-space motions, then select actions for adaptive manipulation \cite{Levine2016DeepGrasp}. 
\textbf{Motion planning} then realizes the chosen grasp: sampling-based algorithms such as RRT* \cite{Karaman2011RRTStar} remain fundamental for generating collision-free trajectories in high-dimensional configuration spaces, exposed through frameworks like MoveIt \cite{Sucan2010MoveIt}. 
These capabilities are benchmarked on the Amazon Picking Challenge \cite{Correll2016AmazonChallenge}, which established a standardized testbed for picking under clutter and uncertainty, and on the YCB Object and Model Set \cite{Calli2015YCB}, a diverse collection of everyday items with physical models for grasping and manipulation; grasp success rate, cycle time per pick, placement accuracy, and task-completion rate are the headline metrics (Table~\ref{tab:p4_s9_5_benchmarks}). 
Hybrid control frameworks that combine analytical models with data-driven adaptation are emerging as the practical route to generalizing across variable geometries and grasp contexts, complemented by tactile sensing and force control for slip detection and contact regulation.

\subsection{Pivotal Task-5: Human-Robot Collaboration and Safety}
\label{subsec:p4-sec9-5-pivotal-task-5-human-robot-collaboration-and-safety}
Even in highly automated warehouses, human workers remain indispensable for complex, dexterous, or exception-handling tasks that require contextual understanding and flexible decision making, so the final stage of the pipeline is safe and efficient human-robot collaboration (HRC). 
HRC systems must balance productivity with stringent safety guarantees, enabling mixed teams (robots handling repetitive or heavy-load work, humans managing fine-grained, judgment-based activities) to share workspaces without physical or operational interference \cite{Lasota2015HRC}. 
The task maps multimodal human-detection and intent signals (pose, gesture, gaze, and trajectory data from vision, LiDAR, wearables, or motion capture) to adaptive robot behaviors such as dynamic speed modulation, path re-routing, stop-and-wait actions, or task reallocation in response to predicted human movement. Figure~\ref{fig:p4_s9_5_hrc} sketches this predict-then-adapt loop.

\begin{figure}[!ht]
    \centering
    \includegraphics[width=0.98\linewidth]{figures/part4/section9/warehouse_human_robot_collaboration.pdf}
    \caption{The human--robot collaboration loop: multimodal human sensing and intent recognition inform proximity-aware planning, which selects slowing, rerouting, or stopping. Clearance regions and routes are illustrative, not certified safety zones.}
    \label{fig:p4_s9_5_hrc}
\end{figure}

\textbf{Intent recognition} lets robots anticipate human motion: sequence models based on hidden Markov models or recurrent networks predict trajectories and proactively adjust behavior \cite{Mainprice2016Intent}, while \textbf{proxemics-aware planning} maintains socially acceptable distances and re-routes trajectories under human proximity \cite{dragan2013legibility}, with multimodal sensing (vision-based detection fused with wearable or RFID localization) improving reliability. 
These principles are embodied in force- and speed-limited collaborative robots and safety-rated AMRs (see Established Tools and Systems), governed by international standards such as ISO 10218 and ISO/TS 15066. 
Evaluation combines quantitative safety metrics with qualitative interaction assessment: human-motion corpora such as Human3.6M~\cite{IonescuPOS14} supply sequences for intent prediction, and key indicators include intervention frequency (safety-triggered stops per hour), near-miss rate, human-detection accuracy, and compliance with ISO-defined safety zones.

\begin{table}[!th]
\centering
\fontsize{8}{10}\selectfont
\setlength{\tabcolsep}{1.0mm}
\caption{Representative benchmarks for spatial intelligence in warehouse and industrial logistics, spanning the five pivotal tasks.}
\label{tab:p4_s9_5_benchmarks}
\begin{tabular}{L{0.21\textwidth} L{0.21\textwidth} L{0.31\textwidth} L{0.18\textwidth}}
\toprule
\textbf{Pivotal task} & \textbf{Benchmark} & \textbf{Representative method} & \textbf{Headline metric} \\
\midrule
Perception (detection) & COCO~\cite{lin2014microsoft} & YOLO~\cite{redmon2016yolo} / DETR~\cite{Zhu2021DeformableDETR} & mAP \\
Perception (6D pose) & YCB-Video~\cite{xiang2017posecnn} & DenseFusion~\cite{Wang2019DenseFusion} & ADD \\
Perception (textureless pose) & T-LESS~\cite{Hodan2017TLESS} & CosyPose~\cite{Labbe2020CosyPose} & ADD / VSD \\
Localization \& mapping & KITTI~\cite{Geiger2013KITTI} & Visual / LiDAR SLAM & APE / RPE \\
Localization \& mapping & EuRoC MAV~\cite{Burri2016EuRoC} & ORB-SLAM3~\cite{Campos2021ORBSLAM3} (visual-inertial) & APE \\
Mapping (dense) & Replica~\cite{straub2019replica} & RTAB-Map~\cite{Labbe2019RTABMap} & IoU / completeness \\
Reasoning \& planning & Solomon VRP~\cite{Solomon1987VRP} & OR-Tools~\cite{Perron2019ORTools} & total cost / distance \\
Reasoning \& planning & OMPL~\cite{Sucan2012OMPL} & Sampling-based planners & planning time / success \\
Manipulation & Amazon Picking~\cite{Correll2016AmazonChallenge} & Integrated perception--grasping team pipelines & pick success / time \\
Manipulation & YCB Object Set~\cite{Calli2015YCB} & Dex-Net~\cite{Mahler2017DexNet} & grasp success rate \\
\bottomrule
\end{tabular}
\end{table}

\subsection{Established Tools and Systems}
\label{subsec:p4-sec9-5-established-tools-and-systems}
The pivotal tasks above are realized in production through a layered ecosystem of simulators, middleware, optimization services, and robot hardware that together constitute the real-world delta over the methods of Part~\ref{part:3}.

\textbf{Simulation and middleware} form the integration substrate. Digital-twin and physics platforms such as NVIDIA Omniverse Isaac Sim \cite{NVIDIA2021Isaac} and Gazebo provide realistic environments for developing and validating perception, localization, and planning pipelines before deployment, while perception SDKs such as NVIDIA Isaac, the Intel RealSense SDK, and Zebra RFID/barcode platforms  supply end-to-end sensing support in real facilities. Interoperability across heterogeneous hardware is provided by ROS~\cite{Quigley2009ROSIndustrial}, ROS\,2, and ROS-Industrial, with commercial localization stacks such as SLAMcore\footnote{\url{https://www.slamcore.com/}} and BlueBotics ANT\footnote{\url{https://bluebotics.com/}} delivering optimized multi-sensor fusion, map management, and fleet coordination for AGVs and AMRs.

\textbf{Warehouse management and fleet orchestration} sit at the enterprise layer. Large-scale optimization services such as SAP Extended Warehouse Management\footnote{\url{https://www.sap.com/products/scm/extended-warehouse-management.html}}, Blue Yonder WMS\footnote{\url{https://blueyonder.com/}}, and Manhattan Associates\footnote{\url{https://www.manh.com/}} integrate task allocation and scheduling with enterprise resource planning, while open toolkits such as Google OR-Tools \cite{Perron2019ORTools} and PyMARL support the underlying combinatorial and learning-based planners.

\textbf{Robot hardware and manipulation} close the loop physically. High-throughput pick-and-place arms from ABB\footnote{\url{https://new.abb.com/products/robotics}}, FANUC\footnote{\url{https://www.fanuc.com/}}, and KUKA\footnote{\url{https://www.kuka.com/}} are widely deployed and integrated through ROS-based middleware, mobile manipulators such as those from Fetch Robotics\footnote{\url{https://fetchrobotics.com/}} unify navigation and grasping, and motion-planning frameworks such as MoveIt \cite{Sucan2010MoveIt} and NVIDIA cuMotion provide inverse kinematics and trajectory optimization. For shared workspaces, collaborative robots from Universal Robots\footnote{\url{https://www.universal-robots.com/}} and AMRs from Locus Robotics\footnote{\url{https://locusrobotics.com/}} provide the safety-rated platforms on which the HRC layer is built.






\subsection{Challenges}
\label{subsec:p4-sec9-5-challenges}
Deployed warehouse systems work because they fix the SKU sets, layouts, and operating conditions they handle, and each stage of the pipeline has its own obstacle to relaxing those limits: small-object and cluttered perception, perceptual aliasing in repetitive aisles, multi-agent scalability under uncertainty, manipulation of variable objects, and the safety--throughput trade-off in human--robot collaboration.

\paragraph{Cluttered and Small-Object Perception.}
Barcode labels and compact SKU packages often occupy only a few pixels in high-resolution images, leading to frequent misdetections, while occlusion and clutter compound the difficulty as items overlap within bins and shelves. 
A further persistent issue is domain shift: models trained on synthetic or generic data degrade under real-world variation in lighting, reflective surfaces, and packaging. 
Research responses include domain adaptation and self-supervised learning for generalization \cite{Tzeng2017DomainAdversarial}, active perception that deliberately adjusts viewpoints or sensing to resolve occlusion \cite{Bajcsy2018ActivePerception}, and continual learning that lets perception adapt over time without exhaustive retraining, increasingly combined with multi-sensor fusion and neuro-symbolic representations for higher-level causal and semantic reasoning.

\paragraph{Perceptual Aliasing in Repetitive Environments.}
Localization and mapping are degraded by dynamic obstacles (humans, forklifts) that generate transient or misleading landmarks corrupting loop-closure hypotheses \cite{Cadena2016SLAMSurvey}, and above all by perceptual aliasing: repetitive shelving and low-texture surfaces produce ambiguous place recognition and false data associations in feature-based SLAM \cite{Campos2021ORBSLAM3,Labbe2019RTABMap}. 
High-speed AGV and AMR operation further demands low-latency, sub-centimeter pose updates that sensor delays, occlusion, and monocular scale drift can readily violate. Robust multi-modal fusion, learning-based estimation, and tighter integration of SLAM with semantic perception and digital twins are the leading mitigations.

\paragraph{Scalability and Uncertainty in Multi-agent Planning.}
The combinatorial complexity of multi-robot coordination grows exponentially with the number of agents and tasks, rendering exact optimization intractable for large fleets, while real-world uncertainty (new order arrivals, blocked aisles, unpredictable human movement) necessitates adaptive re-planning and decentralized decision making under partial observability \cite{Veloso2015CMUCoBot,Sartoretti2019PRIMAL}. 
Hierarchical and hybrid planners, graph-neural and differentiable coordination, and simulation-informed multi-objective optimization are the active directions toward scalable, resilient fleet behavior.

\paragraph{Manipulation of Variable and Deformable Objects.}
Many SKUs exhibit deformable, reflective, or transparent surfaces that interfere with perception and grasp stability \cite{Calli2015YCB,Mahler2017DexNet}, and deformable handling remains especially difficult because grasp dynamics and contact points resist rigid-body modeling \cite{Levine2016DeepGrasp}. 
Tactile sensing and force control improve slip detection and grasp robustness but face cost and calibration barriers in industrial deployment, motivating hybrid control that combines analytical models with data-driven adaptation for generalization across geometries and grasp contexts.

\paragraph{Safety--Throughput Trade-off in Collaboration.}
Robust human detection under occlusion, variable lighting, and cluttered aisles continues to hinder reliable intent estimation \cite{dragan2013legibility,Mainprice2016Intent}, and the trade-off between throughput and safety remains unresolved: conservative margins reduce efficiency while aggressive planning risks unsafe interaction \cite{Lasota2015HRC}. 
Emerging work emphasizes predictive and explainable safety models that justify decisions in human-interpretable terms, and adaptive safety frameworks that fuse learned intent models with certified control policies for dynamic risk assessment and shared autonomy on the warehouse floor.

\newpage
\chapter{Spatial Intelligence for Virtual Realities}

\FramedInitial{V}{irtual} realities are the media in which spatial intelligence becomes constitutive rather than optional: an extended-reality (XR) device cannot draw a single pixel until it knows where it is, what surrounds it, and where the user is looking and reaching. This makes XR the most demanding real-time deployment target for the perception, reasoning, and generation capabilities of the previous part, specialized to run egocentrically, on-device, and within a hard latency budget.

Beneath everything in XR sits the platform itself, the runtime that localizes the device, maps and registers the surrounding world, senses the user's body, and composites the result in real time; layered on top is the application that has driven the medium hardest, entertainment, spanning live AR and VR experiences and virtual production as well as the offline, high-fidelity reconstruction and performance capture that supply them with authored content.



\section{AR/VR/XR}
\label{sec:p4_s10_1}

\subsection{Introduction}
\label{subsec:p4-sec10-1-introduction}

Augmented, virtual, and mixed reality (collectively, extended reality or XR) form a continuum of media in which a head-worn display couples the user's egocentric viewpoint to a computational world that is registered to, or substituted for, the physical surroundings. 
What distinguishes XR from screen-based media is that spatial intelligence is not an optional feature but a constitutive requirement: before any pixel can be drawn, the system must know where the headset is, what geometry surrounds it, and where the user is looking and reaching. 
XR is therefore the most demanding real-time deployment target for the spatial capabilities surveyed in Part~\ref{part:3}. 
The perception machinery (detection, depth, pose, reconstruction) of Chapter~\ref{chapter:3:sec:6}, the spatial reasoning and planning of Chapter~\ref{sec:reasoning}, and the neural reconstruction and rendering of Chapter~\ref{sec:generation} all reappear here, but specialized to run egocentrically, on-device, and within a hard latency budget.

This section treats the XR \emph{medium itself} as a spatial-computing platform: the runtime engine that localizes the device, maps and registers the world, senses the user's body, and presents a composited result. 
The neighboring section on entertainment (§\ref{sec:entertainment}) builds \emph{on top of} this engine, treating entertainment as an application domain: its real-time AR/VR experiences and virtual production, together with the offline, high-fidelity reconstruction and performance capture that supply them with authored content; here we focus on the platform that makes such experiences run at interactive rates.
Concretely, a headset must continuously answer three coupled questions, \ie, where am I (6-DoF self-localization), what is around me (geometry and semantics), and how do I fuse virtual content into that reality so that it appears physically present, while simultaneously sensing the user and refreshing the display. 
Early systems answered the localization question with outside-in tracking, using external base stations or cameras (SteamVR Lighthouse, the original Oculus constellation), but the dominant modern paradigm is inside-out tracking, in which visual--inertial SLAM runs entirely on the device, as in Meta Quest, Apple Vision Pro, and Microsoft HoloLens. 
The three pivotal tasks below correspond to the stages of this loop.


\begin{figure}[!t]
    \centering
    \includegraphics[width=0.98\linewidth]{figures/part4/section10/xr_overview.pdf}
    \caption{Spatial intelligence for XR as a closed real-time loop: world-facing sensing, tracking and mapping, and registration (T1), user-facing input (T2), and rendering and display (T3). The latency values are illustrative end-to-end targets, not per-stage budgets.}
    \label{fig:p4_s10_1_overview}
\end{figure}


\subsection{System Overview}
\label{subsec:p4-sec10-1-system-overview}

As sketched in Figure~\ref{fig:p4_s10_1_overview}, an XR device closes a tight real-time loop. Egocentric sensors first \textbf{sense} the device and the user (outward cameras and an inertial measurement unit for head pose; inward-facing cameras for eyes and outward- or downward-facing cameras for hands and body). 
The system then \textbf{tracks and maps} the environment, recovering a 6-DoF pose and a persistent geometric model, and \textbf{registers} virtual content into that model with correct anchoring, occlusion, and lighting. 
Finally, it \textbf{renders and displays} the composite back through the headset optics, after which the user moves and the loop repeats. 
The pipeline shares an end-to-end motion-to-photon budget; roughly $20$\,ms at a display refresh of $90$\,Hz or higher are the commonly cited targets, though the right values depend on device and application, and exceeding them degrades presence and comfort. 
The three pivotal tasks map onto this loop directly: Task-1 is the world-facing engine (track, map, register), Task-2 is the user-facing input (egocentric sensing and interaction), and Task-3 is the user-facing output (display and rendering).



\subsection{Pivotal Task-1: Spatial Tracking, Mapping, and Registration}
\label{subsec:p4-sec10-1-pivotal-task-1-spatial-tracking-mapping-and-registration}

The first task is the medium's spatial backbone: it answers where the device is, what surrounds it, and how virtual content is pinned to the world so convincingly that the illusion holds. 
It proceeds in three steps: localize the device, model the environment, and register synthetic content into that model. 
Figure~\ref{fig:p4_s10_1_tracking} sketches this progression from inside-out tracking to a persistent scene map and finally to anchored, occluded, and relit virtual content.

\begin{figure}
    \centering
    \includegraphics[width=0.98\linewidth]{figures/part4/section10/xr_tracking.pdf}
    \caption{The XR world-facing engine: (a) inside-out visual--inertial tracking of 6-DoF pose, (b) mapping of geometry, semantic labels, and persistent anchors, and (c) registration of virtual content with depth-tested occlusion and matched lighting. The room and shading are illustrative, not measured reconstructions.}
    \label{fig:p4_s10_1_tracking}
\end{figure}

\textbf{Inside-out 6-DoF tracking.} The foundational layer estimates the headset's six-degree-of-freedom pose at high rate by fusing visual features with inertial measurements. 
Feature-based visual SLAM with loop closure and relocalization remains the academic standard, exemplified by the ORB-SLAM lineage and its visual-inertial extension ORB-SLAM3~\cite{Campos2021ORBSLAM3}, while tightly-coupled visual--inertial odometry such as VINS-Mono~\cite{qin2018vins} underpins many mobile and headset stacks; direct (photometric) methods such as DSO~\cite{engel2017direct} and the earlier LSD-SLAM~\cite{engel2014lsd} offer an alternative that does not depend on sparse keypoints. 
More recent learned dense systems, including DROID-SLAM~\cite{teed2021droid} and transformer-based geometric matchers~\cite{mast3r}, trade computation for robustness in low-texture and fast-motion regimes. 
These methods are evaluated chiefly on EuRoC~\cite{Burri2016EuRoC} and TUM-VI~\cite{Schubert2018TUMVI}, with absolute trajectory error (ATE) as the headline metric, but the binding constraints in XR are drift accumulation, recovery from tracking loss (the kidnapped-device problem), and the power and thermal envelope of on-device inference.

\textbf{Spatial mapping and scene understanding.} The tracked stream is consolidated into a usable model of the environment: dense depth and mesh, planar surfaces (floor, wall, table), and semantic labels. 
Real-time volumetric fusion in the KinectFusion tradition~\cite{newcombe2011kinectfusion} provides the geometric substrate, on top of which plane and mesh extraction expose the surfaces that gameplay and user interfaces attach to. 
(The offline, high-fidelity reconstruction of such scenes as navigable content is treated in §\ref{sec:entertainment} and Chapter~\ref{sec:generation}; here the map is a means to localization and registration, not an end in itself.) 
Semantic and instance understanding is benchmarked on indoor datasets such as ScanNet~\cite{dai2017scannet} and Matterport3D~\cite{chang2017matterport3d}, with mean intersection-over-union (mIoU) for segmentation, and increasingly on mobile-captured data with laser ground truth such as ARKitScenes~\cite{baruch2021arkitscenes} and high-quality reconstructions such as Replica~\cite{straub2019replica} for depth and geometry evaluation.

\textbf{Persistent and shared anchors.} For content to survive across sessions and to be experienced jointly by co-located users, the map must be persistent and shareable. 
Cloud-hosted and peer-to-peer world maps let a device relocalize against a previously built model, and at city scale a Visual Positioning System anchors content to real landmarks in well-mapped areas. 
The open problems are relocalization under appearance change (lighting, season, rearranged furniture), the memory footprint of on-device maps, and the privacy implications of uploading scans of private spaces, a thread we return to in the challenges.

\textbf{Registration and realistic compositing.} Finally, virtual content must be registered both geometrically and photometrically so that it appears to occupy the scene. 
Geometric registration pins objects to anchors and planes in the SLAM frame; depth-based occlusion lets real geometry correctly hide virtual pixels, as in the per-pixel depth produced for mobile interaction by DepthLab and the ARCore Depth API~\cite{du2020depthlab}; and lighting estimation predicts the dominant light direction and an ambient environment map from a single image so that virtual objects are shaded and shadowed consistently~\cite{gardner2017learning}. 
Small errors here are perceptually unforgiving: a few centimeters of depth error at an object boundary, or a mismatched shadow, immediately reveals the composite as synthetic, which is why robust occlusion and real-time inverse lighting remain active research problems.

\subsection{Pivotal Task-2: Egocentric Human Sensing and Spatial Interaction}
\label{subsec:p4-sec10-1-pivotal-task-2-egocentric-human-sensing-and-spatial-interaction}

The second task closes the loop from the user's body back into the system: it senses the hands, eyes, and posture from inside-out cameras and turns those signals into spatial intent. Low-level egocentric estimates of hands, gaze, and body are fused into selection, manipulation, and locomotion. As Figure~\ref{fig:p4_s10_1_sensing} indicates, the same headset sensors that localize the device are repurposed to track the user, and the resulting estimates drive a gaze-and-gesture interaction model.

\begin{figure}[!ht]
    \centering
    \includegraphics[width=0.98\linewidth]{figures/part4/section10/xr_sensing.pdf}
    \caption{Egocentric human sensing for XR interaction: (a) articulated hand tracking, (b) gaze from inward-facing eye cameras, (c) full-body pose inferred from three tracked points, and (d) gaze-plus-pinch selection, in which gaze supplies the target and a pinch confirms it.}
    \label{fig:p4_s10_1_sensing}
\end{figure}

\textbf{Hand tracking.} Articulated 3D hand pose (commonly $21$ joints, often a full mesh) is regressed from headset cameras under heavy self-occlusion and truncation. 
The standard hand model is MANO~\cite{romero2017mano}, and on-device pipelines such as MediaPipe Hands~\cite{zhang2020mediapipe} combine a palm detector with a landmark regressor to run in real time. 
Methods are benchmarked on FreiHAND~\cite{zimmermann2019freihand} and the two-hand InterHand2.6M~\cite{moon2020interhand2}, with 3D mean per-joint position error (MPJPE) in millimeters as the headline metric; robust two-hand interaction and egocentric truncation remain the hard cases.

\textbf{Eye and gaze tracking.} Estimating 3D gaze from infrared eye images serves three roles: calibration, gaze as a pointing modality, and as the driver for foveated rendering (Task-3). 
Appearance-based estimation in the iTracker and GazeCapture tradition~\cite{krafka2016gazecapture} established that gaze can be regressed from commodity cameras, and headset-specific datasets such as OpenEDS~\cite{garbin2019openeds} provide near-eye infrared imagery with dense eye-region segmentation, with the later OpenEDS2020 track adding gaze prediction (angular error in degrees). 
The decisive interaction consequence is that gaze targets quickly but lacks a natural trigger, a long-standing difficulty known as the Midas-touch problem, which motivates the gaze-plus-pinch pattern discussed below.

\textbf{Full-body pose from sparse signals.} Driving a full-body avatar from a headset and two controllers is an underconstrained inverse problem, since the lower body is unobserved and must be synthesized. 
Recent methods infer plausible full-body motion from these three tracked points, including the transformer-based AvatarPoser~\cite{jiang2022avatarposer} and the diffusion-based AGRoL~\cite{du2023agrol}, while a parallel line recovers pose from sparse body-worn inertial sensors, as in DIP~\cite{huang2018dip}, TransPose~\cite{yi2021transpose}, and the physics-aware PIP~\cite{yi2022pip}. 
These models are trained on large motion corpora such as AMASS~\cite{mahmood2019amass} and evaluated with MPJPE and temporal jitter. 
This sparse, live setting is distinct both from the wearable-IMU body sensing of §\ref{sec:p4_s9_3} (which instruments the body more densely) and from the studio performance capture that drives authored avatars in §\ref{sec:entertainment}; the XR-specific challenge is plausibility from radically incomplete observation.

\textbf{Multimodal interaction and locomotion.} The recovered signals are fused into spatial actions. The foundational idea of combining a deictic gesture with a complementary modality dates to Bolt's ``Put-That-There''~\cite{bolt1980putthatthere}, and the modern instantiation decouples targeting from confirmation: the eyes select while a hand pinch or a voice command commits, sidestepping the Midas-touch problem. 
Locomotion through a virtual space larger than the physical room is handled by redirected walking, which applies imperceptible rotational and translational gains to steer the user within a confined area~\cite{razzaque2001redirected}, with teleportation as the common non-natural fallback. 
Spatial user interfaces, in turn, must choose a reference frame (world-locked, body-locked, or head-locked) and increasingly anchor panels to scene semantics so that content lands on appropriate surfaces.
Learning these interaction behaviors at scale is increasingly grounded in large egocentric corpora such as Ego4D~\cite{grauman2022ego4d}, which supply the first-person video of hands, objects, and activities that egocentric-understanding models are trained and benchmarked on.

\subsection{Pivotal Task-3: Spatial Display and Real-Time Rendering}
\label{subsec:p4-sec10-1-pivotal-task-3-spatial-display-and-real-time-rendering}

The third task presents the composited world back to the user within the perceptual budget that the previous two stages must respect. 
Three questions decide it: where to spend rendering effort, how to hide system latency, and how to display the result comfortably. 
Figure~\ref{fig:p4_s10_1_display} previews this pipeline, from gaze-contingent shading through latency compensation to the optics that place the image before the eye.

\begin{figure}[!ht]
    \centering
    \includegraphics[width=0.98\linewidth]{figures/part4/section10/xr_display.pdf}
    \caption{The XR display and rendering pipeline: (a) gaze-contingent foveated rendering, (b) late reprojection of a rendered frame to the latest pose, and (c) alternative video-passthrough and optical see-through display paths. The motion-to-photon target is illustrative and applies end to end.}
    \label{fig:p4_s10_1_display}
\end{figure}

\textbf{Foveated rendering.} Because acuity falls off sharply outside the fovea, rendering effort can be concentrated at the gaze point. Foveated 3D graphics established multi-resolution shading driven by eye tracking~\cite{guenter2012foveated}, and subsequent work added temporal stability and contrast preservation to make the quality reduction imperceptible~\cite{patney2016foveated}. Eye-tracked foveated rendering now ships on consumer headsets, materially reducing shading cost on thermally constrained hardware.

\textbf{Latency compensation.} Even an efficient renderer cannot meet the motion-to-photon budget on every frame, so the last stage warps the rendered image to the most recent head pose. 
Asynchronous time warp applies a rotational reprojection at display refresh, asynchronous space warp adds depth-aware frame extrapolation, and head-mounted optical-see-through systems use per-pixel late-stage reprojection~\cite{vanwaveren2016timewarp}. 
These techniques keep the world stable under fast head motion and are essential to comfort.

\textbf{Passthrough versus optical see-through.} Two display regimes trade off differently. Video passthrough mixed reality (as on Meta Quest~3 and Apple Vision Pro) renders the real world from cameras, granting full control over occlusion and relighting at the cost of camera-to-display latency, parallax, and limited dynamic range; optical see-through (as on HoloLens~2 and Magic Leap~2) presents the real world with zero latency at its true focal depth (though the virtual overlay still incurs render latency and a fixed focal plane) but cannot occlude real light or match real-world brightness. 
The choice propagates back into the registration problem of Task-1.

\textbf{Depth-correct rendering and visual comfort.} Conventional stereo displays fix the focal plane while vergence varies with content, producing the vergence-accommodation conflict and the eye strain that limits session length. 
Varifocal, light-field, and holographic near-eye displays aim to restore correct focus cues~\cite{maimone2017holographic}, but remain largely at the prototype stage.

\textbf{Neural rendering for telepresence.} Finally, photoreal social presence requires driving a faithful avatar of a remote user from headset-mounted cameras that see only oblique crops of the face. 
The codec-avatar line, beginning with deep appearance models~\cite{lombardi2018deep}, learns a person-specific decoder that is driven on-device, a direct specialization of the neural rendering of Chapter~\ref{sec:generation} to the constrained, real-time XR setting.

\begin{table}[!ht]
\centering
\fontsize{8}{10}\selectfont
\setlength{\tabcolsep}{1.0mm}
\caption{Representative benchmarks for spatial intelligence in XR, spanning the three pivotal tasks.}
\label{tab:p4_s10_1_benchmarks}
\begin{tabular}{L{0.16\textwidth} L{0.25\textwidth} L{0.30\textwidth} L{0.20\textwidth}}
\toprule
\textbf{Pivotal task} & \textbf{Benchmark} & \textbf{Representative method} & \textbf{Headline metric} \\
\midrule
Tracking & EuRoC~\cite{Burri2016EuRoC} & ORB-SLAM3~\cite{Campos2021ORBSLAM3} & ATE \\
Tracking & TUM-VI~\cite{Schubert2018TUMVI} & VINS-Mono~\cite{qin2018vins} & ATE \\
Mapping & ScanNet~\cite{dai2017scannet} & dense fusion / semantics & mIoU \\
Mapping & ARKitScenes~\cite{baruch2021arkitscenes} & mobile RGB-D mapping & depth abs-rel \\
Hand tracking & InterHand2.6M~\cite{moon2020interhand2} & InterNet~\cite{moon2020interhand2} & 3D MPJPE \\
Body (sparse) & AMASS~\cite{mahmood2019amass} & AvatarPoser~\cite{jiang2022avatarposer} & MPJPE / jitter \\
Gaze & OpenEDS~\cite{garbin2019openeds} & appearance-based gaze regressor & angular error \\
Interaction & Ego4D~\cite{grauman2022ego4d} & egocentric understanding & task-specific \\
\bottomrule
\end{tabular}
\end{table}

\subsection{Established Tools and Systems}
\label{subsec:p4-sec10-1-established-tools-and-systems}

A mature ecosystem implements these tasks across mobile AR, standalone headsets, and the web, and it is organized into developer frameworks, cross-platform standards, hardware, and the companies that drive them.

\textbf{SDKs and standards.} At the device level, \textbf{Apple's ARKit}\footnote{\url{https://developer.apple.com/augmented-reality/}} and \textbf{Google's ARCore}\footnote{\url{https://developers.google.com/ar}} expose plane detection, anchors, depth, light estimation, and people occlusion, with ARCore's Geospatial API localizing against world imagery for outdoor anchoring; \textbf{Vuforia}\footnote{\url{https://developer.vuforia.com}} specializes in robust image and object targets for industrial AR. 
Cross-engine abstraction layers such as \textbf{Unity AR Foundation}\footnote{\url{https://docs.unity3d.com/Packages/com.unity.xr.arfoundation@6.0/}} and the \textbf{Meta Presence Platform and Spatial SDK}\footnote{\url{https://developers.meta.com/horizon/develop/spatial-sdk}} present the common subset (planes, anchors, meshes, occlusion, hand and body tracking) over multiple devices, while \textbf{8th Wall and WebXR}\footnote{\url{https://immersiveweb.dev}} bring markerless AR into the browser. These abstractions matter because XR content is expensive to author and must migrate across devices and sessions; this is the role of cross-platform standards, principally \textbf{OpenXR}\footnote{\url{https://www.khronos.org/openxr/}} as a unifying runtime API and glTF together with OpenUSD and \textbf{USDZ}\footnote{\url{https://openusd.org}} as asset and scene-description formats. Without a shared runtime and shared asset formats, the ecosystem fragments into incompatible walled gardens.

\textbf{Hardware.} Current headsets divide along the passthrough versus optical-see-through axis. 
\textbf{Apple Vision Pro}\footnote{\url{https://www.apple.com/apple-vision-pro/}} (2024) and \textbf{Meta Quest~3 and 3S}\footnote{\url{https://www.meta.com/quest/}} use high-resolution video passthrough, with the Quest line the consumer volume leader; \textbf{Microsoft HoloLens~2}\footnote{\url{https://www.microsoft.com/en-us/hololens/}} (optical see-through, with production ended in 2024) and \textbf{Magic Leap~2}\footnote{\url{https://www.magicleap.com}} target enterprise, and \textbf{Varjo}\footnote{\url{https://varjo.com}} and \textbf{Pico}\footnote{\url{https://www.picoxr.com}} serve high-fidelity simulation and the Asian consumer market respectively. 
Lightweight true-AR glasses, such as \textbf{Snap Spectacles}\footnote{\url{https://www.spectacles.com}} and \textbf{Meta's Orion and Project Aria} research prototypes, signal the all-day-glasses horizon toward which the field is converging.

\textbf{Companies and platforms.} Meta leads consumer volume and the most aggressive glasses research bet; 
\textbf{Apple} defines the premium spatial-computing tier and the USD content pipeline; \textbf{Google} supplies world-scale localization through ARCore and Android XR; \textbf{Microsoft} has pivoted from first-party HoloLens hardware toward OpenXR software and enterprise cloud services; \textbf{Niantic}\footnote{\url{https://www.nianticspatial.com}} operates a city-scale Visual Positioning System for persistent, multi-user AR; and Snap, Varjo, and ByteDance's Pico round out the social-AR, high-fidelity, and consumer-standalone niches.

\subsection{Challenges}
\label{subsec:p4-sec10-1-challenges}

XR runs every spatial capability egocentrically and on-device under a perceptual real-time budget, and its open problems follow from that constraint: holding the motion-to-photon budget, fitting perception into the compute and thermal envelope of glasses, keeping tracking and registration robust in the wild, sharing persistent maps without compromising privacy, and making extended wear visually comfortable.

\paragraph{Latency and the Motion-to-Photon Budget.} Meeting the motion-to-photon targets of the System Overview is increasingly difficult as pipelines add wireless or split rendering, cloud offload, and learned perception in the loop, each of which contributes its own delay that the reprojection stage of Task-3 can only partially mask.

\paragraph{On-Device Compute, Thermal, and Power.} All-day glasses face a hard energy and heat envelope: a lightweight form factor leaves little room for the silicon that SLAM, hand and eye tracking, and scene understanding demand. The frontier is partitioning perception across on-device accelerators, companion devices such as wristbands or pucks, and edge offload, without violating the latency or thermal budget.

\paragraph{Robust Tracking and Registration in the Wild.} Visual--inertial SLAM and relocalization degrade in dynamic scenes, textureless rooms, repetitive structures, low light, and outdoors, exactly the conditions that curated benchmarks such as EuRoC and ScanNet underrepresent~\cite{Burri2016EuRoC,dai2017scannet}. 
Persistent Visual Positioning Systems push toward city-scale anchoring, but globally drift-free, cross-device localization, together with convincing occlusion and lighting registration, remains unsolved.

\paragraph{Persistent, Shared Maps and Privacy.} Co-located multi-user AR requires a shared, persistent spatial map, an ``AR cloud'', which collides directly with privacy: always-on egocentric cameras continuously scan private spaces and bystanders. 
Designing maps that support shared anchoring while minimizing the retention of raw imagery, and that admit meaningful consent, is both a research and a policy frontier.

\paragraph{Visual Comfort.} Beyond latency, fixed-focus stereo optics decouple vergence from accommodation, causing eye strain that limits session length, and poorly designed locomotion induces cybersickness. 
Varifocal and holographic displays, gaze-contingent depth-of-field, and accurate interpupillary-distance matching are promising but immature, leaving comfort a gating factor for mainstream adoption.

Because XR subjects the perception, reasoning, and generation capabilities of Part~\ref{part:3} to egocentric, on-device, real-time constraints, it exposes failure modes that offline settings rarely reveal.

\newpage
\section{Entertainment}
\label{sec:entertainment}

\subsection{Introduction}
\label{subsec:p4-sec10-2-introduction}

In entertainment, spatial intelligence is exercised chiefly as the ability to perceive and reconstruct 3D environments, objects, and human bodies, and to embed virtual content into the physical world so that it is geometrically and physically consistent. Spatial entertainment (AR games, VR experiences, location-based attractions, and virtual production for film) differs from screen-based media in that it must compute camera pose, scene geometry, and user state continuously and within strict latency limits, without sacrificing aesthetic quality or narrative coherence; the perception and reconstruction machinery involved is treated in Chapter~\ref{chapter:3:sec:6} and Chapter~\ref{sec:generation}.

The field has followed the same arc as the capabilities it borrows: commodity depth sensing made real-time dense mapping practical for AR games, neural radiance and Gaussian representations made captured sets photoreal and navigable, parametric body models made full-body avatars routine, and world-scale visual positioning let location-based games such as Pok\'emon GO anchor characters to real landmarks across sessions and players. Each of these threads returns in the pivotal tasks below.

\begin{figure}[!ht]
    \centering
    \includegraphics[width=0.98\linewidth]{figures/part4/section10/entertainment_overview.pdf}
    \caption{Spatial intelligence for entertainment as a real-time closed loop: sensor inputs feed world reconstruction, performance capture, and content authoring, whose composited render is presented back to the user.}
    \label{fig:p4_s10_2_overview}
\end{figure}

\subsection{System Overview}
\label{subsec:p4-sec10-2-system-overview}

Entertainment sits on top of the XR platform of §\ref{sec:p4_s10_1}: it does not re-implement that runtime but builds on it, and its distinctive demand is production-quality spatial content, whether rendered live or authored offline. The domain spans two regimes. Real-time interactive experiences (AR games, VR titles, location-based attractions, and live virtual production) close the tight sensing-to-display loop sketched in Figure~\ref{fig:p4_s10_2_overview}, whereas offline high-fidelity creation (scanning sets, capturing performers, building volumetric footage) supplies those experiences with assets.
Across both, three spatial tasks recur as the frequent bottlenecks.
\textbf{World reconstruction} builds the geometric and radiance model of the environment, supplying localization and occlusion for live experiences and navigable captured content for authored ones;
\textbf{embodied performance capture} digitizes the user or performer into a controllable avatar; and
\textbf{spatial content authoring} decides where virtual elements appear and how cameras move so that the result is both consistent and expressive.
The interactive regime runs all three under hard latency budgets, while offline authoring trades that latency for fidelity.



\subsection{Pivotal Task-1: World Reconstruction and Spatial Mapping}
\label{subsec:p4-sec10-2-pivotal-task-1-world-reconstruction-and-spatial-mapping}
The first task builds and maintains an accurate, consistent model of the environment that supports localization, occlusion, and physics for live experiences and navigable content for authored ones, specializing the dense-reconstruction and scene-understanding techniques of Chapter~\ref{chapter:3:sec:6} and Chapter~\ref{sec:generation} to interactive rates. Where §\ref{sec:p4_s10_1} treated mapping as a means to on-device localization, here the reconstruction is itself the deliverable: the navigable set or stage that players and film-makers inhabit.
Figure~\ref{fig:p4_s10_2_reconstruction} contrasts geometric fusion, neural and point-based scene representations, and world-scale localization as complementary capabilities rather than obligatory sequential stages.

\begin{figure}
    \centering
    \includegraphics[width=0.98\linewidth]{figures/part4/section10/entertainment_reconstruction.pdf}
    \caption{Complementary capabilities for entertainment reconstruction, not sequential stages: (a) TSDF fusion of posed RGB-D frames into a mesh, (b) radiance-field or 3D Gaussian appearance models for view synthesis, and (c) visual positioning for persistent place-based content. Samples, surfaces, and place markers are schematic.}
    \label{fig:p4_s10_2_reconstruction}
\end{figure}

\textbf{Geometric reconstruction} forms the base: geometric SLAM and dense methods fuse depth or RGB data into volumetric maps, as in KinectFusion~\cite{newcombe2011kinectfusion} and the reconstructions underlying datasets such as ScanNet~\cite{dai2017scannet}, while DynamicFusion extends this to moving users by tracking non-rigid deformations \cite{newcombe2015dynamicfusion}. 
\textbf{Neural and point-based representations} such as NeRF, instant-NGP, and 3D Gaussian Splatting then represent scene radiance to produce compact, high-fidelity reconstructions that render at interactive rates \cite{mildenhall2021nerf,muller2022instant,kerbl20233d}; dynamic and 4D extensions that deform or time-index the primitives carry this to moving scenes, so a performance or an entire set can be captured as free-viewpoint volumetric video for later replay \cite{yang2024deformable,wu20244d}. At \textbf{global scale}, visual positioning systems (VPS) combine local reconstructions into a worldwide map that lets a mobile device localize from a single image, as in Niantic's Lightship platform. These methods are evaluated on indoor datasets such as ScanNet and Matterport3D, which provide large-scale reconstructions with semantic labels for scene understanding \cite{dai2017scannet,chang2017matterport3d}, and on NeRF-style view-synthesis benchmarks (\eg, synthetic Blender scenes) for photometric fidelity; the latter, however, rarely capture entertainment-specific constraints such as rapid camera motion or cluttered living rooms, so practitioners often fall back on custom internal datasets (Table~\ref{tab:p4_s10_2_benchmarks}). 
The standard metrics are absolute trajectory error (ATE) for localization, surface reconstruction error for geometry, and PSNR/SSIM for neural rendering \cite{mildenhall2021nerf,kerbl20233d}, but in entertainment these are complemented by interaction-centric measures: the accuracy of occlusion masks, the physical plausibility of collisions, and user reports on realism (for instance, whether a virtual character convincingly hides behind a sofa).

\subsection{Pivotal Task-2: Embodied Performance Capture and Avatar Control}
\label{subsec:p4-sec10-2-pivotal-task-2-embodied-performance-capture-and-avatar-control}
The second task estimates full-body pose, shape, and facial expression to drive avatars in real time, specializing the human pose and motion machinery of Chapter~\ref{chapter:3:sec:6} (and the motion-generation methods of Chapter~\ref{sec:generation}) to live interaction. Unlike the sparse three-point body tracking of §\ref{sec:p4_s10_1}, driven from a headset and two controllers, entertainment capture targets studio-grade, fully-observed performance.
As Figure~\ref{fig:p4_s10_2_capture} indicates, approaches range from skeleton-based joint prediction to surface-based model fitting and full volumetric capture.

\begin{figure}
    \centering
    \includegraphics[width=0.98\linewidth]{figures/part4/section10/entertainment_capture.pdf}
    \caption{Three routes to performance capture, each able to drive avatar playback or control: (a) skeleton joints retargeted to a rig, (b) a parametric surface (SMPL) driving a body mesh, and (c) volumetric capture for reconstructed or animatable appearance. Performer illustrations and motion traces are qualitative, not measured scores.}
    \label{fig:p4_s10_2_capture}
\end{figure}

\textbf{Skeleton-based} methods predict joint positions from sensors, while \textbf{surface-based} approaches fit parametric models such as SMPL directly to video \cite{loper2015smpl}, with datasets like Human3.6M and AMASS supplying the supervised motion priors that make this robust \cite{IonescuPOS14,mahmood2019amass}. 
\textbf{Volumetric capture} systems such as Holoportation go further, streaming textured 3D avatars into VR/AR scenes by combining reconstruction with tight time synchronization \cite{orts2016holoportation}. More recent Gaussian-splatting avatars learn drivable, photorealistic humans from multi-view capture, with rigged 3D Gaussians giving expressive control over head avatars \cite{qian2024gaussianavatars} and pose-dependent Gaussian maps modeling full-body avatars with dynamic clothing \cite{li2024animatablegaussians}, narrowing the gap between a captured performance and a real-time renderable asset. These capabilities are benchmarked chiefly on Human3.6M, which offers millions of 3D poses with calibrated cameras \cite{IonescuPOS14}, and on AMASS, which provides over forty hours of motion data for learning priors \cite{mahmood2019amass}; commercial scenarios remain underrepresented, however, since common cases such as players using controllers in small rooms or performers in costumes lack sufficient public benchmark data (Table~\ref{tab:p4_s10_2_benchmarks}).
Quantitative evaluation uses metrics like mean per-joint position error (MPJPE), but for interactive entertainment latency and temporal stability are decisive, as even minor jitter breaks the sense of embodiment, and user studies show that enjoyment depends on the alignment between perceived and actual motion \cite{dongas2023vrrhythm}; production systems therefore pair algorithmic metrics with feedback from expert animators.

\subsection{Pivotal Task-3: Spatially-Aware Content Authoring and Virtual Production}
\label{subsec:p4-sec10-2-pivotal-task-3-spatially-aware-content-authoring-and-virtual-production}
The third task helps creators decide where virtual elements appear and how cameras move, turning the reconstructed world and captured performers into a coherent, expressive scene, and connecting to the spatial-reasoning and planning paradigms of Chapter~\ref{sec:reasoning}. 
Figure~\ref{fig:p4_s10_2_authoring} previews the spectrum from SDK-level authoring primitives to camera-tracked LED-wall virtual production and AI-assisted cinematography.

\begin{figure}
    \centering
    \includegraphics[width=0.98\linewidth]{figures/part4/section10/entertainment_authoring.pdf}
    \caption{Three related spatial authoring capabilities: (a) engine primitives (detected planes and persistent anchors) for asset placement, (b) virtual production, where tracked camera pose updates the LED-wall perspective, and (c) camera planning from narrative and scene constraints to a shot plan. The scene, camera path, and shot settings are illustrative.}
    \label{fig:p4_s10_2_authoring}
\end{figure}

\textbf{Engine-level authoring} provides designers with high-level primitives such as plane detection and anchors through AR SDKs in real-time engines like Unity and Unreal.
\textbf{Virtual production} links physical cameras to 3D scenes on LED walls so that lighting and parallax respond correctly to camera movement, and \textbf{intelligent cinematography} research is exploring AI-assisted camera planning that respects film grammar while adapting to dynamic user actions \cite{azzarelli2025intelligentcine}. Standardized benchmarks for spatial authoring are rare: some studies evaluate camera planning on synthetic storyboards, but most assessment happens internally at studios, and virtual-production workflows are typically tested on complex reference sequences rather than public datasets, which makes objective comparison of authoring tools difficult (Table~\ref{tab:p4_s10_2_benchmarks}). 
Evaluation accordingly mixes system metrics such as rendering latency with expert judgment on cinematic quality, with intelligent-cinematography studies often convening expert panels to analyze framing and continuity \cite{azzarelli2025intelligentcine}; for AR/VR tools usability is paramount, and HCI evaluations measure how quickly a user completes a level-design task or how rich the resulting spatial layouts are under different tool abstractions.
Generative world models are beginning to bypass authoring by synthesizing playable worlds directly: GameNGen~\cite{valevski2024diffusion} simulates the game DOOM interactively in real time with a diffusion model, and Genie~3~\cite{genie3} generates navigable environments that respond to user actions, so that creating content starts to mean conditioning a world model (§\ref{sec:p3_s8_6}).

\begin{table}[!th]
\centering
\fontsize{8}{10}\selectfont
\setlength{\tabcolsep}{1.0mm}
\caption{Representative benchmarks for spatial intelligence in entertainment, spanning the three pivotal tasks.}
\label{tab:p4_s10_2_benchmarks}
\begin{tabular}{L{0.19\textwidth} L{0.22\textwidth} L{0.30\textwidth} L{0.20\textwidth}}
\toprule
\textbf{Pivotal task} & \textbf{Benchmark} & \textbf{Representative method} & \textbf{Headline metric} \\
\midrule
World reconstruction & ScanNet~\cite{dai2017scannet} & Dense TSDF fusion~\cite{newcombe2011kinectfusion} & recon.\ error \\
World reconstruction & Matterport3D~\cite{chang2017matterport3d} & Dense SLAM & recon.\ error / mIoU \\
View synthesis & NeRF-synthetic~\cite{mildenhall2021nerf} & 3D Gaussian Splatting~\cite{kerbl20233d} & PSNR / SSIM \\
Performance capture & Human3.6M~\cite{IonescuPOS14} & SMPL fitting~\cite{loper2015smpl} & MPJPE \\
Performance capture & AMASS~\cite{mahmood2019amass} & Motion-prior learning & MPJPE / jitter \\
Content authoring & (no public benchmark) & Intelligent cinematography~\cite{azzarelli2025intelligentcine} & expert panel / latency \\
\bottomrule
\end{tabular}
\end{table}

\subsection{Established Tools and Systems}
\label{subsec:p4-sec10-2-established-tools-and-systems}

A robust set of tools implements these tasks for games, film, and live events. At the engine level, Unity's AR Foundation\footnote{\url{https://docs.unity3d.com/Packages/com.unity.xr.arfoundation}} offers a unified API over ARKit, ARCore, and other SDKs, exposing essential features such as point clouds and depth meshing, while Google's ARCore Geospatial API\footnote{\url{https://developers.google.com/ar/develop/unity-arf/geospatial/developer-guide}} adds global localization, allowing applications to attach anchors to outdoor landmarks with meter-level accuracy.
These abstractions let designers build AR gameplay around real-world surfaces without managing low-level sensor fusion.

Niantic's Lightship platform packages spatial mapping and multiplayer synchronization into turnkey services: Reality Blending\footnote{\url{https://nianticlabs.com/news/realityblending-announcement}} in Pok\'emon GO uses learned depth to let virtual creatures move behind real-world obstacles, and Pok\'emon Playgrounds\footnote{\url{https://nianticlabs.com/news/pokemon-playgrounds}} extends this to persistent AR, where players leave virtual items at specific locations for others to find, with Niantic's Visual Positioning System (VPS) keeping that content anchored consistently across users and times.

In VR gaming, consumer headsets with inside-out tracking enable room-scale interaction; games like Beat Saber use precise 6-DoF tracking to create rhythm gameplay that doubles as exercise, and studies show a strong link between spatial design choices and player exertion, underscoring the connection between tracking quality and user experience \cite{dongas2023vrrhythm}. 
Virtual production stages have likewise transformed film and sports broadcasting: productions such as \emph{The Mandalorian} (built on ILM's StageCraft) and sports studios such as ESPN's Catalyst Stage\footnote{\url{https://espnpressroom.com/us/press-releases/2023/06/catalyst-stage-disneys-new-immersive-espn-studio-opens-storytelling-and-production-frontiers/}} use LED volumes driven by tracked cameras, where spatial intelligence involves calibrating physical cameras to digital scenes and maintaining low-latency tracking. 
Finally, neural capture tools inspired by NeRF and 3D Gaussian Splatting are entering the pipeline, letting artists scan a location with a phone and import navigable 3D assets into game engines \cite{muller2022instant,kerbl20233d}; while this bridges scouting and production, it introduces new challenges around lighting control and integration.

\subsection{Challenges}
\label{subsec:p4-sec10-2-challenges}

Entertainment asks for production-quality spatial content at interactive rates, and its open problems stem from that demand. The tension between real-time performance and visual fidelity is the most immediate; close behind come robustness once the system leaves controlled spaces, the cognitive load that authoring places on creators, and the human-factors and ethics questions that surface the moment these systems scan real users and real rooms.

\paragraph{Real-time Performance versus Fidelity.} Neural representations have improved visual quality, but many methods remain too computationally heavy for mobile or standalone VR. While techniques such as 3D Gaussian Splatting achieve real-time rendering on high-end hardware \cite{kerbl20233d}, porting them to thermally constrained devices is difficult, and entertainment applications must additionally handle variable network conditions while guaranteeing stable frame rates, since frame drops that are acceptable in research benchmarks can cause physical discomfort in VR.

\paragraph{Robustness in the Wild.} Existing benchmarks such as ScanNet focus on standard indoor spaces, yet entertainment experiences often occur in unpredictable environments such as festivals, dim clubs, or crowded stadiums. World-scale AR platforms address this by training on millions of crowdsourced scans, which raises questions about coverage bias (particularly toward urban areas) and the difficulty of maintaining maps as the world changes; designing models that remain robust under these shifts is an open research area.

\paragraph{Authoring Complexity.} Although modern SDKs are powerful, creating high-quality spatial experiences imposes a high cognitive load, as creators must simultaneously manage geometry, interaction, cinematography, and technical constraints. Intelligent-cinematography tools attempt to assist by suggesting camera paths \cite{azzarelli2025intelligentcine}, but these are not yet standard in creative workflows; there is a pressing need for tools that let designers specify high-level intentions (for example, ``the character should hide behind the sofa'') which the system then realizes via underlying planning modules.

\paragraph{Human Factors and Ethics.} Spatial systems frequently scan users and private spaces. While datasets like AMASS demonstrate the utility of large-scale motion capture, collecting similar data in the wild raises privacy concerns, and from an HCI perspective issues such as cybersickness and fatigue limit the duration of VR and AR sessions, while physically active titles introduce comfort and safety considerations that spatial layout must respect. Addressing these problems requires a collaborative approach spanning algorithm design, user-feedback systems, and clear data governance.

\newpage
\chapter{Spatial Intelligence for Geo-spatial Systems }


\FramedInitial{G}{eo-spatial} systems operate on spatial structure at the scale of assets, cities, and the planet, where the task is to build a faithful model of a physical system and keep it aligned with live observation. This chapter applies the spatial intelligence of the previous part to that task, across three settings that share the same closed loop between model and world but differ in scale and stakes.

The digital twin binds a metric, physically grounded replica to a specific asset through a live bidirectional loop, so that the physical system and its model track each other. GeoAI widens the frame to the planet, where streams of Earth-observation data drive perception and decision-making over terrain, infrastructure, and environmental change. Aerospace systems carry the same loop off the planet, where a spacecraft or rover must perceive, reason about, and act on its surroundings with little or no ground support.


\section{Digital Twins}
\label{sec:p4_s11_1}
\label{sec:Digital Twins}

\subsection{Introduction}
\label{subsec:p4-sec11-1-introduction}

A \textbf{digital twin} is a 3D representation faithful enough to stand in for its physical counterpart in analysis and operation.
Building one requires more than accurate geometric reconstruction and photorealistic appearance under varying illumination; it requires realistic modeling of physical properties and interaction, so that the replica behaves as the original would.
What separates a digital twin from any other 3D asset, however, is not fidelity alone but \emph{coupling}.
Following the now-standard taxonomy of Kritzinger \etal~\cite{kritzinger2018digital}, a \emph{digital model} carries no automated data exchange with the physical object, a \emph{digital shadow} carries an automated one-way flow from the physical to the virtual, and only a \emph{digital twin} closes the loop with an automated, bidirectional flow: the virtual model is continuously synchronized to the physical asset, simulated forward to predict, and used to act back upon it.

A digital twin therefore reuses, but specializes, much of the machinery surveyed earlier in this handbook. The perception stack of Chapter~\ref{chapter:3:sec:6} supplies the captures from which a twin is built, and the neural reconstruction and generation methods of Chapter~\ref{sec:generation} supply its geometry and appearance; what the twin adds is metric accuracy, physical grounding, and a persistent correspondence to a specific real object over its lifecycle.
Entertainment (§\ref{sec:entertainment}) builds the same high-fidelity captures as creative \emph{content}; here they become an \emph{operational engineering model} wired to a live counterpart, and the city- and planet-scale twins that recur below connect to the GeoAI methods of §\ref{sec:p4_s11_2}.
The concept itself traces from Grieves' product-lifecycle lectures, formalized with Vickers~\cite{grieves2017digital}, through its aerospace codification for NASA and U.S. Air Force vehicles~\cite{glaessgen2012digital}, to its industrial maturation~\cite{tao2019digital,wang2023survey}, where the five-dimension model of Tao \etal makes the data connections between physical entity, virtual entity, and services a first-class element of the architecture.

\subsection{System Overview}
\label{subsec:p4-sec11-1-system-overview}

Figure~\ref{fig:p4_s11_1_overview} draws the twin as a closed loop around its physical asset. Sensing (images, LiDAR, RGB-D scans, onboard IoT telemetry) feeds the construction of a replica that carries geometry, appearance, materials, and physical properties; the replica is simulated forward, kept synchronized by live data assimilation, and queried to predict, and its decisions operate the asset, which is then sensed anew.
The three pivotal tasks map onto this loop directly. Task-1 builds the twin (sense and construct), Task-2 makes it behave (model and simulate), and Task-3 keeps it synchronized and turns it to use (assimilate, predict, and act). Table~\ref{tab:p4_s11_1_benchmarks} collects the benchmarks on which the first two are evaluated; synchronization has no shared benchmark.

\begin{figure}
    \centering
    \includegraphics[width=0.98\linewidth]{figures/part4/section11/digital_twins_overview.pdf}
    \caption{The digital twin as a closed loop around its physical asset: construction (T1) and modeling (T2) establish the replica, synchronization (T3) corrects its state from live observations, and forecasts drive decisions through a control path separate from the observation path.}
    \label{fig:p4_s11_1_overview}
\end{figure}

\subsection{Pivotal Task-1: High-Fidelity Twin Construction}
\label{subsec:p4-sec11-1-pivotal-task-1-high-fidelity-twin-construction}

The first task builds the twin from sensor captures, and it must deliver four properties that a generic one-shot reconstruction does not guarantee. The twin must be \emph{metric} (correct at real-world scale), \emph{relightable} (its appearance separable from capture-time illumination), \emph{physically grounded} (carrying mass, material, and friction), and \emph{actionable} (its movable parts and joints encoded so a simulator can drive it).
Figure~\ref{fig:p4_s11_1_construction} sketches the progression from raw captures through geometry and material recovery to a simulatable model.

\begin{figure}
    \centering
    \includegraphics[width=0.98\linewidth]{figures/part4/section11/digital_twins_reconstruction.pdf}
    \caption{Three components of high-fidelity twin construction: (a) metric geometry from multi-view, LiDAR, or RGB-D capture; (b) relightable material channels separated from illumination; and (c) physical articulation, here a prismatic drawer joint with mass and friction. Material swatches and mesh overlays are schematic.}
    \label{fig:p4_s11_1_construction}
\end{figure}

\textbf{Geometry and appearance acquisition.} The earliest replicas were built by photogrammetry: multiple overlapping photographs of an object, terrain, or structure are captured, and a Structure-from-Motion algorithm recovers camera poses and a sparse metric point cloud, as systematized in COLMAP~\cite{schonberger2016structure}, after which dense multi-view stereo and meshing produce a surface.
Real-time RGB-D fusion in the KinectFusion tradition~\cite{newcombe2011kinectfusion} yields metric, watertight geometry over the observed region, and neural implicit-surface methods such as NeuS~\cite{wang2021neus} and the detail-preserving Neuralangelo~\cite{li2023neuralangelo} recover watertight meshes as the zero-level-set of a learned signed-distance field.
A second family optimizes for view synthesis rather than surface quality: NeRF~\cite{mildenhall2021nerf}, its hash-grid acceleration Instant-NGP~\cite{muller2022instant}, and the explicit, real-time 3D Gaussian Splatting~\cite{kerbl20233d} render novel views with high fidelity but historically produced no clean surface and bake appearance into the representation; surface-aligned variants such as SuGaR~\cite{guedon2024sugar} and 2D Gaussian Splatting~\cite{huang20242d} narrow this gap by extracting watertight meshes directly from a Gaussian field.
Because a twin requires metric scale and an extractable surface, the photogrammetric and implicit-surface lines are load-bearing, with meshing operators such as Poisson surface reconstruction~\cite{kazhdan2006poisson} and Marching Cubes~\cite{lorensen1998marching} converting a field or point set into geometry; the radiance-field family enters as an appearance and speed substrate, treated in full in Chapter~\ref{sec:generation}.

\textbf{Inverse rendering and relightable materials.} A captured asset whose appearance is baked at scan time only reproduces the object under the one illumination it was scanned in, because each pixel entangles albedo, geometry, and lighting.
A twin, by contrast, lives in a simulator where viewpoint and illumination change every frame, so its construction must recover separated physically based rendering (PBR) channels (albedo, roughness, normals) together with an estimate of the environment lighting.
Inverse-rendering methods supply exactly this factorization: NeRFactor decomposes a captured radiance field into shape, reflectance, and a single environment light~\cite{zhang2021nerfactor}, while the differentiable, engine-ready NVDiffRec jointly extracts a triangle mesh, an SVBRDF, and a lighting estimate directly optimizable for downstream rendering~\cite{munkberg2022extracting}.
The same factorization has been ported onto explicit Gaussians, as in Relightable 3D Gaussians, which attaches normals and a BRDF to each primitive and ray-traces shadows for real-time relighting~\cite{gao2024relightable}. Relightability is what lets a twin be composited under novel lighting and rendered consistently inside a physics or sensor simulation.

\textbf{Physical-property and articulation estimation.} A static mesh is viewable but not operable: it has no notion of which surfaces move, with what effort, or under what dynamics. Closing this gap requires a movable-part inventory and the kinematic joints (revolute or prismatic, with axes and limits) that a physics engine can integrate, together with physical parameters such as mass and friction. Ditto recovers part geometry and an articulation model jointly from before-and-after observations of an interaction~\cite{jiang2022ditto}, and URDFormer predicts a full URDF, the standard interactive description, from a single real image, turning a photograph into a simulatable articulated twin~\cite{chen2024urdformer}. Dynamical parameters can be identified directly in a differentiable field representation, as in PAC-NeRF, which estimates geometry and physical parameters together through a differentiable material-point simulator~\cite{li2023pacnerf}, and PhysGaussian, which integrates Newtonian dynamics into a Gaussian representation so that what is reconstructed can also be simulated~\cite{xie2024physgaussian}. Such methods lean heavily on articulated-object datasets, principally the PartNet-Mobility models exposed through the SAPIEN simulator~\cite{xiang2020sapien}, the fine-grained part hierarchy of PartNet~\cite{mo2019partnet}, and the cross-category, actionable parts of GAPartNet~\cite{geng2023gapartnet}. The resulting articulated, parameterized model is the input to the behavioral simulation of Task-2.

\textbf{Datasets and metrics.} Construction quality is measured against scanned-object and scene corpora (Table~\ref{tab:p4_s11_1_benchmarks}).
At the object level, the recent Digital Twin Catalog provides a large set of digital-twin-quality scans with millimeter geometry, photorealistic materials, and an explicit relighting evaluation~\cite{dong2025dtc}, complementing the real-scanned OmniObject3D~\cite{wu2023omniobject3d}, the simulation-ready Google Scanned Objects~\cite{downs2022googleGSO}, the artist-authored, PBR-annotated ABO~\cite{collins2022abo}, the large Objaverse collection~\cite{deitke2023objaverse}, and the in-the-wild Common Objects in 3D~\cite{reizenstein2021common}.
At the scene level, twins are evaluated on ScanNet~\cite{dai2017scannet}, Matterport3D~\cite{chang2017matterport3d}, the photorealistic Replica reconstructions~\cite{straub2019replica}, and the commodity-LiDAR ARKitScenes~\cite{baruch2021arkitscenes}, with metric-geometry benchmarks such as Tanks and Temples supplying laser-scan ground truth~\cite{knapitsch2017tanks}.
Geometry is reported with Chamfer distance, F-score at a stated threshold, and normal consistency; appearance with PSNR, SSIM, and LPIPS; and, for twins specifically, an added relighting error that penalizes appearance that fails to transfer across illumination.

\subsection{Pivotal Task-2: Physics-Based Modeling and Behavioral Simulation}
\label{subsec:p4-sec11-1-pivotal-task-2-physics-based-modeling-and-behavioral-simulation}

The second task makes the constructed replica behave.
Geometry is not behavior: a twin must reproduce its asset's dynamics, not merely its shape, so the static model of Task-1 is coupled to a physics, behavioral, or surrogate engine and run forward as an \emph{executable} model whose parameters are fit to the real asset (Figure~\ref{fig:p4_s11_1_simulation}).

\begin{figure}
    \centering
    \includegraphics[width=0.98\linewidth]{figures/part4/section11/digital_twins_simulation.pdf}
    \caption{From a constructed replica to a behaving twin: (a) constructed model, (b) physics engine, and (c) simulated behavior, with system identification fitting the physical parameters $\theta$ to real-asset observations and an optional fast surrogate $f_\phi$. Response traces are illustrative, not measured results.}
    \label{fig:p4_s11_1_simulation}
\end{figure}

\textbf{Physics engines.} The behavioral substrate of a twin is a simulator that integrates its dynamics.
MuJoCo provides accurate rigid-body contact dynamics and is the de facto standard for model-based control~\cite{todorov2012mujoco}; the open-source PyBullet offers a widely used free alternative spanning rigid and soft bodies~\cite{coumans2021}; and SAPIEN specializes in part-level interaction with large articulated objects~\cite{xiang2020sapien}, the natural runtime for the actionable twins of Task-1.
GPU-accelerated platforms couple photorealistic rendering with physics at scale, exemplified by NVIDIA Isaac Sim\footnote{\url{https://developer.nvidia.com/isaac/sim}} on the Omniverse stack and its successor Isaac Lab\footnote{\url{https://developer.nvidia.com/isaac/lab}}, with newer GPU-first simulators such as Genesis\footnote{\url{https://genesis-embodied-ai.github.io/}} and the accelerator-native MuJoCo~MJX pushing throughput further; deformable phenomena such as cloth and fluids are handled by unified particle and material-point (MPM) solvers.

\textbf{Differentiable simulation and system identification.} A twin is only behaviorally faithful if its physical parameters match the real asset, which makes parameter estimation central.
Differentiable simulators expose gradients of simulated outcomes with respect to physical parameters, turning system identification into gradient-based optimization: DiffTaichi compiles differentiable physical simulators from high-level code~\cite{hu2020difftaichi}, Brax provides massively parallel differentiable rigid-body dynamics on accelerators~\cite{freeman2021brax}, and gradSim couples differentiable simulation with differentiable rendering to backpropagate from video pixels all the way to mass, friction, and elasticity~\cite{jatavallabhula2021gradsim}. A twin can thus be calibrated to its counterpart from observation rather than hand-tuned.

\textbf{Surrogate and physics-informed models.} High-fidelity simulation is often too slow to run in the synchronization loop of Task-3, so twins frequently substitute a reduced or learned surrogate.
Physics-informed neural networks embed governing partial differential equations directly in the training loss, solving forward and inverse problems with scarce data~\cite{raissi2019physics}, a strategy situated within the broader programme of physics-informed machine learning~\cite{karniadakis2021physics}.
Reduced-order modeling, identified early as a key enabling technology for digital twins, compresses a high-fidelity simulation to real-time evaluation~\cite{hartmann2018model}, and the hybrid-twin paradigm couples such a physics-based reduced model with a data-driven correction term to absorb unmodeled effects~\cite{chinesta2020virtual}.

\textbf{Real-to-sim for policy learning.} A faithful behavioral twin is also the substrate on which robot policies are trained and validated before deployment, which connects this task to the manipulation methods of §\ref{sec:p4_s9_1} and the reasoning of Chapter~\ref{sec:reasoning}.
The reality gap between simulation and the physical world is bridged by randomizing visual appearance~\cite{tobin2017domain} and dynamics~\cite{peng2018sim} during training, and recent real-to-sim-to-real pipelines build a twin of the workspace from a quick scan and use it to robustify learned policies, as in RialTo~\cite{torne2024reconciling}.
A complementary line argues that an exact twin is not always necessary: ACDC instead generates geometrically and semantically similar ``digital cousins,'' trading exact correspondence for robustness and zero-shot transfer~\cite{dai2024automated}, sharpening by contrast what a true digital twin provides.

\subsection{Pivotal Task-3: Real-Time Synchronization and Predictive Operation}
\label{subsec:p4-sec11-1-pivotal-task-3-real-time-synchronization-and-predictive-operation}

The third task is the one that earns the name \emph{twin}: it maintains the automated bidirectional coupling between model and asset that defines a twin~\cite{kritzinger2018digital}, and turns that coupling to use.
Figure~\ref{fig:p4_s11_1_synchronization} previews the pipeline, from a live sensor stream through data assimilation to prediction and a control signal fed back to the physical asset.

\begin{figure}
    \centering
    \includegraphics[width=0.98\linewidth]{figures/part4/section11/digital_twins_synchronization.pdf}
    \caption{Real-time synchronization and predictive operation: (a) live measurements from the physical asset are (b) assimilated with the model prior into a posterior state of (c) the synchronized twin, which feeds forecasts, what-if, anomaly, and useful-life outputs. The upper feedback path propagates the updated state into the next model prior; the lower path returns decisions or controls to the physical asset. Dashed forecast branches are possible futures, not observed trajectories; all analytical outputs are schematic.}
    \label{fig:p4_s11_1_synchronization}
\end{figure}

\textbf{Data assimilation and state estimation.} Keeping the twin synchronized is a state-estimation problem: the model's latent state must be continually corrected toward the incoming measurements.
The classical recursive estimator is the Kalman filter~\cite{kalman1960new}, and for the nonlinear, high-dimensional systems typical of physical twins the ensemble Kalman filter provides a Monte-Carlo data-assimilation scheme that propagates an ensemble of states to forecast error statistics~\cite{evensen1994sequential}.
At the level of an entire engineered system, a probabilistic graphical model casts data-driven model adaptation as Bayesian state estimation with end-to-end uncertainty quantification~\cite{kapteyn2021probabilistic}, so that predictive twins can be assembled and updated at scale and their predictions carry calibrated confidence rather than point estimates.

\textbf{Predictive operation across domains.} The synchronized twin is queried to anticipate rather than merely to monitor, and the applications follow the asset.
In manufacturing, the twin couples machine geometry to live telemetry for prognostics and health management, estimating remaining useful life and scheduling maintenance before failure~\cite{tao2018digital}. At urban scale, a city digital twin fuses a semantically rich 3D model with sensor and simulation layers to support planning and what-if analysis, as in the Herrenberg study~\cite{dembski2020urban} and national programmes built on CityGML foundations~\cite{soon2017citygml}, a layer that connects directly to the GeoAI methods of §\ref{sec:p4_s11_2}.
In healthcare, patient-specific anatomical twins driven by physiological models support precision diagnosis and treatment planning, with cardiology an early exemplar~\cite{corralacero2020digital} and a growing body of work surveyed across the field~\cite{katsoulakis2024digital}; the medical twin is developed in §\ref{sec:p4_s12_1}.
In aerospace, where the paradigm was first codified, a structural twin mirrors an as-flown airframe to predict fatigue~\cite{glaessgen2012digital}; and in agriculture, field and crop twins fuse geospatial layers with growth models for precision management~\cite{verdouw2021digital}.

\begin{table}[!th]
\centering
\fontsize{8}{10}\selectfont
\setlength{\tabcolsep}{1.0mm}
\caption{Representative benchmarks for twin construction and simulation. Synchronization (Task-3) has no shared benchmark.}
\label{tab:p4_s11_1_benchmarks}
\begin{tabular}{L{0.20\textwidth} L{0.22\textwidth} L{0.30\textwidth} L{0.20\textwidth}}
\toprule
\textbf{Pivotal task} & \textbf{Benchmark / dataset} & \textbf{Representative method} & \textbf{Headline metric} \\
\midrule
Construction (object) & DTC~\cite{dong2025dtc} & neural reconstruction + PBR & geometry + relighting error \\
Construction (object) & OmniObject3D~\cite{wu2023omniobject3d} & multi-view reconstruction & Chamfer / F-score \\
Construction (materials) & ABO~\cite{collins2022abo} & NVDiffRec~\cite{munkberg2022extracting} & material / PSNR \\
Construction (scene) & ScanNet~\cite{dai2017scannet} & KinectFusion-style RGB-D fusion~\cite{newcombe2011kinectfusion} & Chamfer / F-score \\
Construction (metric) & Tanks and Temples~\cite{knapitsch2017tanks} & MVS / neural surfaces & F-score \\
Construction (articulation) & PartNet-Mobility~\cite{xiang2020sapien} & Ditto~\cite{jiang2022ditto} / URDFormer~\cite{chen2024urdformer} & joint / part error \\
Construction (city-scale) & GauU-Scene / MatrixCity & CityGaussianV2~\cite{liu2025citygaussian} & PSNR / geometry F-score \\
Simulation (real-to-sim) & real-world manipulation tasks & RialTo~\cite{torne2024reconciling} / ACDC~\cite{dai2024automated} & policy success rate \\
\bottomrule
\end{tabular}
\end{table}

\subsection{Established Tools and Systems}
\label{subsec:p4-sec11-1-established-tools-and-systems}

\textbf{Capture, authoring, and render runtimes.} The construction stage is served by spatial-capture products such as \textbf{Matterport}\footnote{\url{https://matterport.com/}} for interiors and sites, the mobile scanner \textbf{Polycam}\footnote{\url{https://poly.cam/}}, the photogrammetry pipelines \textbf{RealityScan}\footnote{\url{https://www.realityscan.com/}} (formerly RealityCapture) and Agisoft Metashape\footnote{\url{https://www.agisoft.com/}}, and, as an integrating substrate, NVIDIA Omniverse with \textbf{OpenUSD}\footnote{\url{https://www.nvidia.com/en-us/omniverse/}}, which provides the scene description and real-time rendering on which twins are assembled. Game and render engines such as \textbf{Unreal Engine}\footnote{\url{https://www.unrealengine.com/}} and Unity, the latter packaged for industry as \textbf{Unity Industry}\footnote{\url{https://unity.com/products/unity-industry}}, serve as the interactive runtimes for the resulting models.

\textbf{Industrial twin platforms.} The operational layer is dominated by enterprise platforms that wire CAD geometry to live telemetry. \textbf{Siemens Teamcenter}\footnote{\url{https://plm.sw.siemens.com/en-US/teamcenter/}} drives product twins from a PLM backbone; \textbf{Ansys Twin Builder}\footnote{\url{https://www.ansys.com/products/digital-twin/ansys-twin-builder}} couples reduced-order models to 3D physics; Dassault's 3DEXPERIENCE markets sensor-driven ``virtual twin'' experiences\footnote{\url{https://www.3ds.com/virtual-twin}}; and cloud services including \textbf{PTC ThingWorx}\footnote{\url{https://www.ptc.com/en/products/thingworx}}, \textbf{AWS IoT TwinMaker}\footnote{\url{https://aws.amazon.com/iot-twinmaker/}}, \textbf{Azure Digital Twins}\footnote{\url{https://azure.microsoft.com/en-us/products/digital-twins}}, and \textbf{Bentley iTwin}\footnote{\url{https://www.bentley.com/software/itwin-platform/}} maintain live twin graphs of factories, buildings, and infrastructure, with \textbf{GE Vernova}\footnote{\url{https://www.gevernova.com/software/}} specializing in energy-asset twins.

\textbf{City-scale systems and interoperability standards.} At geo-scale, procedural and streaming platforms such as \textbf{Esri ArcGIS CityEngine}\footnote{\url{https://www.esri.com/en-us/arcgis/products/arcgis-cityengine/overview}}, \textbf{Cesium with 3D Tiles}\footnote{\url{https://cesium.com/}}, \textbf{Google's photorealistic 3D Tiles}\footnote{\url{https://developers.google.com/maps/documentation/tile/3d-tiles}}, and \textbf{NVIDIA Earth-2}\footnote{\url{https://www.nvidia.com/en-us/high-performance-computing/earth-2/}} build and serve city- to planet-scale twins, with national programmes such as Virtual Singapore as flagship deployments. What keeps such twins portable across this fragmented tooling are open standards: OpenUSD for scene description, IFC (ISO 16739) for building information, CityGML for semantic city models, and OGC 3D Tiles for streaming massive geospatial geometry, alongside the framework and terminology standards ISO 23247 and ISO/IEC 30173. The data layer feeding these geo-scale twins is the subject of §\ref{sec:p4_s11_2}.

\subsection{Challenges}
\label{subsec:p4-sec11-1-challenges}

A twin must be high-fidelity, live, and scalable at once, and these three demands pull against one another. The open problems follow from that tension: scaling construction from an object to a city, keeping the twin current, validating that it is accurate rather than merely photorealistic, coupling several physical domains in real time, and making twins interoperable without exposing the spaces they capture.

\paragraph{Scalability to Large and City Scale.} Building a twin of a single object is now routine; building one of a city is not.
City twins must convert heterogeneous source data into semantically structured models (CityGML levels of detail) and stream them efficiently (3D Tiles), and the reconstruction itself must scale, motivating divide-and-conquer pipelines for large-scene Gaussian splatting~\cite{liu2025citygaussian}.
Data integration and cross-format interoperability remain the named bottlenecks at this scale.

\paragraph{Synchronization Latency and Data Assimilation.} A twin is only as useful as it is current, and maintaining real-time correspondence across an edge-fog-cloud hierarchy is hard: assimilation must keep pace with the asset's dynamics, and the ``liveness gap'' between the asset's true state and the twin's estimate degrades every downstream prediction. Scaling data assimilation to high-dimensional physical systems without violating the update budget is an open problem.

\paragraph{Model Fidelity, Validation, and the Photorealism-Accuracy Tradeoff.} Verification, validation, and uncertainty quantification are the maturity frontier for twins used in consequential decisions.
A specific tension recurs: methods such as vanilla Gaussian splatting optimize visual fidelity, not metric accuracy, so a photorealistic reconstruction can still fail a metric tolerance check; surface-aligned variants~\cite{guedon2024sugar,huang20242d} narrow but do not close this gap.

\paragraph{Multi-Physics and Multi-Scale Coupling.} Real assets couple mechanics, thermics, fluids, and electromagnetics across scales, and no universal model spans them; coupled systems risk solver instability, and the reduced-order and physics-informed surrogates~\cite{hartmann2018model,raissi2019physics} that make twins real-time must do so without discarding the cross-domain interactions that matter.

\paragraph{Standardization, Interoperability, Privacy, and Security.} The proliferation of scene and asset formats (USD, IFC, CityGML, 3D Tiles) keeps twins from composing across vendors, and the always-on capture that a live twin requires raises acute concerns: scanning private interiors and urban space implicates bystanders, and a compromised twin server can leak an asset's dynamics, configuration, and credentials~\cite{wang2023survey}. Designing twins that interoperate while minimizing retained raw data is both a research and a policy frontier.

\newpage
\section{GeoAI}
\label{sec:p4_s11_2}


\subsection{Introduction}
\label{subsec:p4-sec11-2-introduction}
Geospatial artificial intelligence (GeoAI) applies machine learning to geographic data (satellite and UAV imagery, LiDAR, GNSS traces, and IoT sensor streams) to model and act on processes that unfold over terrain, cities, and the Earth system.
It makes the geospatial domain a large-scale deployment target for the perception and scene-understanding capabilities surveyed in Part~\ref{part:3}: the detection, segmentation, and remote-sensing techniques of Chapter~\ref{chapter:3:sec:6} (and the aerial and remote-sensing setting of Section~\ref{sec:p4_s9_4}) are here specialized to Earth-observation data and decision workflows at planetary scale.

\subsection{System Overview}
\label{subsec:p4-sec11-2-system-overview}
A GeoAI workflow couples heterogeneous geospatial data to models that turn raw observations into spatially explicit decisions (Figure~\ref{fig:p4_s11_2_overview}).
Following the application taxonomy of the review by Boutayeb \etal~\cite{boutayeb2024geoai}, we treat six domains as pivotal tasks: precision agriculture, the urban built environment, environmental monitoring, water resources, disaster management, and public health.
The six share their inputs, largely the free Sentinel, Landsat, and MODIS archives, and until recently they shared a modeling pattern as well: a task-specific supervised network or tree ensemble trained for one region and one target, which the foundation models at the end of the section are beginning to replace.
Table~\ref{tab:p4_s11_2_benchmarks} samples representative data sources, methods, and reported results across the six.

\begin{figure}[!ht]
    \centering
    \includegraphics[width=0.98\linewidth]{figures/part4/section11/geoai-overview.pdf}
    \caption{GeoAI as a data-to-decision loop: satellite, UAV, LiDAR, GNSS, and IoT observations feed alternative or combined CNN, transformer, and ensemble models, whose outputs serve six application domains and act back, through planning and policy, on the sensed world.}
    \label{fig:p4_s11_2_overview}
\end{figure}

\subsection{Pivotal Task-1: Precision Agriculture}
\label{subsec:p4-sec11-2-pivotal-task-1-precision-agriculture}
Precision agriculture tailors irrigation, fertilization, and pest control to the conditions of each field, or each plant, instead of managing a field uniformly.
GeoAI supplies the maps and predictions this requires through the complementary mapping, yield-estimation, and irrigation tasks of Figure~\ref{fig:p4_s11_2_agriculture}.

\begin{figure}[!ht]
    \centering
    \includegraphics[width=0.98\linewidth]{figures/part4/section11/geoai-agriculture.pdf}
    \caption{Three GeoAI tasks for precision agriculture: (a) crop mapping, (b) yield and leaf-area-index (LAI) estimation, and (c) variable-rate precision irrigation. Maps and fitted points are schematic.}
    \label{fig:p4_s11_2_agriculture}
\end{figure}

\textbf{Crops Mapping.} Crop mapping turns a time series of satellite or UAV imagery into a per-pixel crop-type label, and the accuracy gains in recent work track the richness of the temporal signal.
Han \etal~\cite{han2023spatio} combine a ResNet feature extractor with cross-attention and spatio-temporal self-attention over Sentinel-1 SAR series, reaching 96\% overall accuracy on the Brandenburg benchmark.
3D convolutional networks over Sentinel-2~\cite{gallo2023season} and Landsat~\cite{MOHAMMADI2023272} series treat time as a third convolution axis, and fusing Sentinel-1 SAR with Sentinel-2 optical data maps rice phenological stages within the season~\cite{tiwari2024automated}.
Multi-season and SAR--optical inputs consistently beat single-date classification, in part because SAR sees through the cloud that interrupts optical series.
At the other end of the resolution range, UAV imagery supports row-level segmentation, where U-Net outperforms LinkNet and PSPNet on sugarcane crop lines~\cite{ribeiro2024automated}.

\textbf{Yields Prediction.} Yield prediction regresses a continuous productivity estimate from the same imagery, and the usual design has two stages: extract a vegetation index or phenotype, then regress.
Random forests over Sentinel-2, RTK GPS, and LiDAR variables predict cotton yield~\cite{carneiro2023soil}, and an ensemble over Landsat-8 NDVI and agro-meteorological data reports R$^2$ 0.99 for tea yield on its test set~\cite{ramzan2023multimodal}.
UAV imagery earns its extra cost where the target is a fine-grained phenotype: Mask R-CNN extracts wheat-spike phenotypes that feed a yield regressor~\cite{PengWZYLRLSX23}, and a linear mixed-effect model on fixed-wing UAV multispectral imagery estimates rice leaf-area index and biomass~\cite{rs11111371}.

\textbf{Precision Irrigation.} Irrigation closes the loop, turning maps and estimates into water-delivery decisions~\cite{AbioyeAMBIROOR20}.
One line of work detects irrigation infrastructure at scale, such as center-pivot systems in Sentinel-2 series with Mask R-CNN~\cite{AlbuquerqueCSLB21} or irrigation types in NAIP aerial images with a U-Net and transfer learning~\cite{RaeiANSPA22}.
The other controls water locally: a drone system fuses RGB, thermal-infrared, and GPS data to find the driest zones of a field and sets valve flow with a random forest~\cite{JalajamonyNMF23}.

\subsection{Pivotal Task-2: Urban Planning, Logistics, and Transport}
\label{subsec:p4-sec11-2-pivotal-task-2-urban-planning-logistics-and-transport}
The second task turns GeoAI on the built environment. Figure~\ref{fig:p4_s11_2_urban} presents five complementary tasks on a shared urban context, from built structure and temporal change to services, transport, and green infrastructure.

\begin{figure}[!ht]
    \centering
    \includegraphics[width=0.98\linewidth]{figures/part4/section11/geoai-urban.pdf}
    \caption{GeoAI for the built environment: (a) building and land-use extraction from overhead and street-view imagery; (b) temporal urban change; (c) sensor-supported digital twins and operational services; (d) mobility, routing, and infrastructure analysis; and (e) canopy, parks, green roofs, and biosecurity monitoring. Tasks are complementary rather than a mandatory serial pipeline. Shared map geometry, temporal trends, and the city illustration are schematic; arrows denote time, data flow, or travel direction within the corresponding panels.}
    \label{fig:p4_s11_2_urban}
\end{figure}

\textbf{Urban Data Extraction.} Land use and built form are classified from overhead and street-level imagery, and fusing the two views resolves ambiguities that neither settles alone.
A fine-tuned EfficientNet classifies the 21 land-use classes of the UC Merced set~\cite{AlemK22}, while a dual-branch network that pairs dilated convolutions on Google Earth imagery with an attention-based multi-view network on Street View detects informal urban villages~\cite{ChenFNYGYGL22}.

\textbf{Urban Dynamics.} Temporal models capture how urban form changes over time.
An LSTM with attention simulates urban densification from land-use rasters of 2000, 2010, and 2020~\cite{el2023long}, and a pix2pix GAN predicts land-use and land-cover change by conditioning Landsat imagery on planning indicators~\cite{SunMFWH21}.
A static quantity such as urban solar exposure needs no temporal model; a feed-forward network on building geometry, occupancy, and orientation suffices~\cite{tehrani2024predicting}.

\textbf{Smart Cities.} The service layer turns these representations into live operation: video georeferencing and small-object detection inside an urban digital twin~\cite{HuLYHCYLY23}, LSTM parking-occupancy prediction served to a mobile application~\cite{CanliT21}, and detection of commercial signage in Street View imagery to map urban commerce~\cite{YeWKXC19}.

\textbf{Logistics and Transport.} Movement through the city splits into two problems with different tools.
Flow forecasting is temporal: bus arrival times are predicted from a traffic-density matrix with recurrent and convolutional models~\cite{PanovskiZ20}, a hybrid GRU--LSTM forecasts traffic from OpenStreetMap and GIS features~\cite{ZafarHSCM22}, IODRNN decomposes its outputs incrementally to remove the lag in GNSS-based forecasts~\cite{LuMPXC24}, and itinerary recommendation adds a denoising autoencoder with collaborative learning~\cite{ChenZCWC20}.
Infrastructure mapping is a segmentation problem: open LiDAR point clouds inventory rockfall-prone road slopes~\cite{RuaNAM23}, SegNet on RapidEye imagery updates rural road networks~\cite{kearney2020maintaining}, and U-Net with PointNet++ performs panoptic segmentation of railway environments from multimodal LiDAR~\cite{grandio2023multimodal}.

\textbf{Green Infrastructure and Bio Security.} Vegetation monitoring moves from coverage to function to threat, and each step demands finer spatial and spectral resolution than the last.
A semi-supervised DeepLabv3+ maps urban tree cover nationwide from high-resolution satellite imagery~\cite{guo2023nationwide}; a Faster R-CNN tree inventory identifies seven canopy species and, through geographically weighted regression, their effect on land-surface temperature~\cite{das2022geoai}; and biosecurity surveillance of invasive forest pests uses multispectral aerial imagery, where an instance-segmentation model reached an F1-score of 79\% against 58\% for a tile classifier~\cite{carnegie2023airborne}.

\subsection{Pivotal Task-3: Environmental Monitoring}
\label{subsec:p4-sec11-2-pivotal-task-3-environmental-monitoring}
The third task observes the Earth system itself, across the three fronts of Figure~\ref{fig:p4_s11_2_environment}.

\begin{figure}[!ht]
    \centering
    \includegraphics[width=0.98\linewidth]{figures/part4/section11/geoai-environment.pdf}
    \caption{Three parallel GeoAI tasks for environmental monitoring: (a) climate and vegetation (NDVI/LAI trends), (b) atmospheric chemistry (NO$_2$, SO$_2$, aerosols), and (c) environmental impact of land-cover change. Maps and trends are illustrative, not measured changes.}
    \label{fig:p4_s11_2_environment}
\end{figure}

\textbf{Climate Change.} Long-term land-surface change is read mostly through vegetation records spanning decades of Landsat, Sentinel, and MODIS data, and the methods trade spatial detail against record length.
Lou \etal~\cite{lou2023deep} map land use and land cover on Google Earth Engine with a deep network and find rapid greening of the Qinghai--Tibet Plateau from 1988 to 2018.
Deep-learning-enhanced spatiotemporal fusion reconstructs long-term NDVI in the Danjiang River Basin~\cite{wang2023applying}, random forests retrieve vegetation phenology from MODIS more accurately than rule-based methods~\cite{XinLLLZ20}, and a bidirectional LSTM produces the GLASS 250\,m leaf-area-index product (version 6) from MODIS reflectance~\cite{ma2022development}.

\textbf{Atmospheric Chemistry.} Atmospheric composition is retrieved by learning the map from satellite radiances and meteorology to ground-measured concentrations.
Nitrogen dioxide is the most studied target, estimated hourly from Sentinel-5P, meteorology, and sensor data with XGBoost~\cite{kim2021importance} and daily across China with a random forest and K-means clustering over OMI data~\cite{dou2021estimates}.
A convolutional autoencoder coupled to WRF-Chem models volcanic SO$_2$ dispersion~\cite{opio2023modeling}, AQNet fuses Sentinel-2 and Sentinel-5P to predict NO$_2$, O$_3$, and PM$_{10}$~\cite{rowley2023predicting}, an LSTM forecasts five pollutants from ground-station series~\cite{seng2021spatiotemporal}, and a deep network derives aerosol absorption from the OMI ultraviolet aerosol index, validated against AERONET~\cite{SunVVL21}.

\textbf{Environmental Impacts.} Impact assessment links land-use change to ecological outcomes, and it is judged not only on accuracy but on whether it can name the responsible driver, which is why explainability recurs here.
Neural networks on multi-temporal Landsat land-cover maps assess urban ecological quality~\cite{alqadhi2024applying}; multiscale geographically weighted regression, random forests, and SHAP relate agricultural land to PM$_{2.5}$ concentrations~\cite{gianquintieri2024implementation}; and interpretable ensembles measure the effect of climate change on agricultural land suitability~\cite{ShevchenkoLTBGKNSKM24}.

\subsection{Pivotal Task-4: Water Resources Management and Precipitation Forecasting}
\label{subsec:p4-sec11-2-pivotal-task-4-water-resources-management-and-precipitation-forecasting}
The fourth task addresses four complementary aspects of water resources in Figure~\ref{fig:p4_s11_2_water}.

\begin{figure}[!ht]
    \centering
    \includegraphics[width=0.98\linewidth]{figures/part4/section11/geoai-water.pdf}
    \caption{Four GeoAI tasks for water resources: (a) surface hydrology, (b) groundwater from GRACE storage anomalies and auxiliary data, (c) water quality calibrated against in-situ samples, and (d) radar-based precipitation nowcasting. The panels are parallel tasks, not successive stages of the water cycle; all values are schematic.}
    \label{fig:p4_s11_2_water}
\end{figure}

\textbf{Hydrological Modeling.} Convolutional models dominate local surface geometry.
Space-time image velocimetry on UAV video estimates river velocity and discharge, validated with LiDAR-derived terrain~\cite{islam2023mutual}; U-Net extracts drainage ditches from LiDAR in low-relief terrain~\cite{du2024drainage}; quadtree decomposition with multi-scale CNNs extracts coastlines at lower cost than FCN and U-Net baselines~\cite{sun2024quadtree}; and a reflectance-transformation CNN maps bathymetry from Landsat and Sentinel-2 imagery~\cite{ChanZX0WA24}.
At basin scale, where streamflow depends on several precipitation products (MODIS, CHIRPS, MSWEP), super-ensembles that weight multiple regressors take over~\cite{wegayehu2023super}.

\textbf{Groundwater.} No optical sensor sees the water table, so groundwater models work from proxy variables and gravimetry.
Static delineation of recharge or shallow-groundwater zones uses tree ensembles and deep networks over hydrogeological and remote-sensing variables~\cite{al2024hybrid,sahour2022identification,nhu2024new}.
GRACE/GRACE-FO measures total-water-storage changes; isolating groundwater requires auxiliary estimates of other water-storage components. Storage forecasting uses these gravity-derived anomalies, which a CNN--LSTM forecasts with an RMSE of 44.92\,mm/month~\cite{seo2021predicting} and which machine-learning regressors downscale to catchment resolution~\cite{sabzehee2023enhancing}.

\textbf{Water Quality.} Water quality is regressed from spectral reflectance against in-situ samples, and accuracy follows the optical signature of the constituent.
Pigments with a clear signature are tractable: fusion of data and classifiers through higher-order SVD estimates chlorophyll-a in Lake Nicaragua~\cite{sun2021ensemble}, and on UAV multispectral and hyperspectral data of a karst wetland, transformers and gradient boosting predict phycocyanin, dissolved oxygen, and chlorophyll-a, interpreted with SHAP~\cite{yang2023monitoring}.
Dissolved nutrients remain harder because their optical fingerprint is faint, whether estimated from MODIS with a spatiotemporal deep belief network~\cite{wu2022spatiotemporal} or from Landsat-8 with an ANN~\cite{vakili2020determination}.
Optical water types themselves can be classified from Sentinel-2 with an SVM~\cite{dasilva2021machine}.

\textbf{Precipitation Forecasting.} Precipitation nowcasting predicts the next radar frames from the current sequence.
HST-AFP fuses radar echoes with weather forecasts through attention~\cite{niu2023heterogeneous}, STIN adds humidity, wind, and temperature~\cite{jin2023spatiotemporal}, and PFST-LSTM adds pseudo-flow estimation to follow echo motion~\cite{luo2020pfst}; at longer range, multi-model ensembles post-process numerical weather prediction~\cite{sengoz2023machine}.

\subsection{Pivotal Task-5: Disaster Management}
\label{subsec:p4-sec11-2-pivotal-task-5-disaster-management}
The fifth task covers the disaster lifecycle (hazard or susceptibility mapping, event observation, and damage assessment) for earthquakes, floods, and forest fires, arranged as a grid in Figure~\ref{fig:p4_s11_2_disaster}.

\begin{figure}[!ht]
    \centering
    \includegraphics[width=0.98\linewidth]{figures/part4/section11/geoai-disaster.pdf}
    \caption{GeoAI for disaster management as a lifecycle-by-hazard grid: earthquakes, floods, and forest fires (columns) across hazard or susceptibility mapping, event observation, and post-event assessment (rows). All maps and severity scales are schematic.}
    \label{fig:p4_s11_2_disaster}
\end{figure}

\textbf{Earthquakes.} For earthquakes the lifecycle includes hazard mapping, event-related observation, and damage assessment; these tasks should not be equated with reliable prediction of an exact earthquake.
A CNN over elevation, vector, and attribute data maps seismic risk in north-eastern India~\cite{jena2021earthquake}, and encoder--decoder LSTMs~\cite{xiong2022gnss} and neural networks~\cite{hafeez2021possible} detect ionospheric anomalies in GNSS total-electron-content data around events such as the 2019 Pakistan earthquake.
After the event, LiDAR, orthophotos, and multispectral imagery estimate building vulnerability~\cite{torres2019integration}, and change detection on multi-temporal photogrammetric point clouds classifies structural damage into grades~\cite{zahs2023classification}.

\textbf{Floods.} Floods lean on radar.
Forecasting fuses meteorological, hydrological, and crowdsourced data, either comparing classical learners under cross-validation~\cite{PuttinaovaratH20} or with a ConvLSTM that learns spatial rainfall patterns and their temporal dynamics jointly~\cite{MoishinDPRA21}.
Inundation mapping falls to SAR segmentation, which sees through cloud, as in an attention-augmented U-Net on dual-polarized Sentinel-1 imagery~\cite{LiWCWJG22}.

\textbf{Forest Fires.} Fire has the fullest lifecycle.
Susceptibility is predicted from drone imagery with temperature and humidity sensing~\cite{rathod2023multipurpose}, or from elevation, vegetation indices, and hydrometeorology with a neuro-fuzzy system tuned by an artificial bee colony~\cite{PHAM2024102392}.
Detection is trained on drone RGB and infrared imagery such as the FLAME 2 dataset~\cite{ChenHWOARFCRW22} and can be distributed across sites with federated learning~\cite{SiddiqueADBAA24}.
Burn scars are delineated afterward by a U-Net on PlanetScope RGB and near-infrared imagery, with F1 up to 0.94~\cite{ChoPKKLS23}.

\subsection{Pivotal Task-6: Healthcare}
\label{subsec:p4-sec11-2-pivotal-task-6-healthcare}
The sixth task brings GeoAI to public health, where spatial signals couple disease and environment (Figure~\ref{fig:p4_s11_2_health}).

\begin{figure}[!ht]
    \centering
    \includegraphics[width=0.98\linewidth]{figures/part4/section11/geoai-health.pdf}
    \caption{Three parallel GeoAI tasks for public health: (a) disease surveillance and incidence forecasting, (b) spatial spread modeled by graph message passing over regions, and (c) pollution exposure assessment. The panels are not causally linked; maps and intervals are illustrative, not surveillance results.}
    \label{fig:p4_s11_2_health}
\end{figure}

\textbf{Infectious Diseases and Toxins.} Disease surveillance gains most from leading indicators, since search queries and geotagged posts surface an outbreak before clinical reporting does.
Influenza forecasting combines CDC surveillance data with Google Flu Trends and climate in a multi-stage LSTM~\cite{DBLP:journals/access/VennaTGRMN19}, with geotagged Twitter posts in a neural network tuned by an artificial-tree algorithm~\cite{hu2018prediction}, and with search frequencies and health records in a network-based regression that exploits synchrony across U.S. states~\cite{lu2019improved}; dengue in China is forecast from climate and weekly case counts, with support-vector regression performing best~\cite{guo2017developing}.
FINDER shows the payoff most clearly: it geolocates web-search signals of foodborne illness to restaurants, and the restaurants it flagged were three times more likely to pose food-safety risks than those found by complaint-based inspection~\cite{SadilekCDMSBJRG18}.

\textbf{Pandemics.} A pandemic shifts the question to transmission across geography.
Geospatial similarity learning predicts respiratory-virus spread in new regions~\cite{ZhaoHJCWWYR23}, and spatio-temporal regression tracks COVID-19 in Brazil~\cite{da2021covid}.
Detection networks catch what an epidemic curve misses: gatherings of people and vehicles in satellite imagery~\cite{atek2022geospatial}, proximity between re-identified individuals in video~\cite{ZhangLCZJ23}, and the drop in PM$_{2.5}$ during lockdown, measured from MODIS and Sentinel-5P~\cite{TanLGJ22}.

\textbf{Public Health.} Chronic exposure links where people are to what they breathe.
On the mobility side, mobile health uses wireless devices for clinical and public-health purposes~\cite{kahn2010mobile}, and the MD2K center turns mobile-sensor streams into markers of smoking, craving, and stress~\cite{KumarAAaBCCCEEG15}.
On the pollutant side, a geographically weighted gradient-boosting machine~\cite{zhan2017spatiotemporal} and a neural network over MODIS aerosol optical depth and land-use features~\cite{di2016assessing} map daily PM$_{2.5}$ across China and the United States.
Exposure assessment is only as sharp as the coarser of the two spatial scales.

\begin{table}[!th]
\centering
\fontsize{8}{10}\selectfont
\setlength{\tabcolsep}{1.0mm}
\caption{Representative data sources and methods for GeoAI, one or two per pivotal task. Numbers are each paper's own test set and are not comparable across rows; cells without a number name the metric the paper reports.}
\label{tab:p4_s11_2_benchmarks}
\begin{tabular}{L{0.20\textwidth} L{0.24\textwidth} L{0.28\textwidth} L{0.20\textwidth}}
\toprule
\textbf{Pivotal task} & \textbf{Representative data source} & \textbf{Representative method} & \textbf{Headline metric} \\
\midrule
Precision agriculture & Sentinel-1 SAR & STMA~\cite{han2023spatio} & 96\% OA (Brandenburg) \\
Precision agriculture & Landsat-8 NDVI + meteorology & Ensemble~\cite{ramzan2023multimodal} & R$^2$ 0.99 \\
Urban planning & UC Merced land use & EfficientNet (fine-tuned)~\cite{AlemK22} & accuracy \\
Environmental monitoring & Landsat, 1988--2018 & Deep LUCC mapping on GEE~\cite{lou2023deep} & LUCC accuracy \\
Water resources & GRACE / precipitation & CNN--LSTM~\cite{seo2021predicting} & RMSE 44.92\,mm/month \\
Disaster (flood) & Sentinel-1 SAR & Attention U-Net~\cite{LiWCWJG22} & precision / recall / F1 \\
Disaster (forest fire) & PlanetScope (RGB+NIR) & U-Net~\cite{ChoPKKLS23} & F1 0.94 \\
Healthcare & CDC / Google Flu / climate & Multi-stage LSTM~\cite{DBLP:journals/access/VennaTGRMN19} & RMSE / MAPE \\
\bottomrule
\end{tabular}
\end{table}

\subsection{Established Tools and Systems}
\label{subsec:p4-sec11-2-established-tools-and-systems}

\textbf{Cloud geospatial platforms} provide planetary-scale storage and analysis next to the data. Google Earth Engine\footnote{\url{https://earthengine.google.com/}} offers a petabyte-scale archive of satellite imagery with server-side analysis, Microsoft's Planetary Computer\footnote{\url{https://planetarycomputer.microsoft.com/}} pairs open Earth-observation data with a Python/STAC API, and commercial providers such as Esri ArcGIS\footnote{\url{https://www.arcgis.com/}} and Descartes Labs\footnote{\url{https://www.descarteslabs.com/}} deliver hosted GeoAI pipelines for enterprise and government users.

\textbf{Open-source GIS and ML toolkits} carry the modeling layer. QGIS\footnote{\url{https://qgis.org/}} is the de facto open desktop GIS, while deep-learning libraries tailored to raster geospatial data, such as TorchGeo\footnote{\url{https://github.com/microsoft/torchgeo}} and Raster Vision\footnote{\url{https://rastervision.io/}}, bridge PyTorch to satellite imagery, and Sentinel Hub\footnote{\url{https://www.sentinel-hub.com/}} streamlines on-demand access and processing of Copernicus archives.

\textbf{Open Earth-observation programmes} supply the data foundation: the European Copernicus programme's Sentinel constellation and NASA/USGS Landsat provide free, calibrated, multi-decadal imagery, complemented by gravimetric (GRACE/GRACE-FO), atmospheric (OMI, Sentinel-5P), and high-resolution commercial (WorldView, PlanetScope) sources that recur throughout the methods above.

\subsection{Toward Geospatial Foundation Models and GeoLLMs}
\label{subsec:p4-sec11-2-toward-geospatial-foundation-models-and-geollms}
Almost every method surveyed above is a task-specific supervised model: a CNN, random forest, LSTM, or GAN trained from scratch for one crop, one hazard, or one pollutant. That design is being displaced by the same pretraining-and-transfer paradigm that reshaped the rest of the field (Chapter~\ref{sec:reasoning}), and the geospatial community has begun building its own foundation models. Masked-autoencoding backbones lead the shift: SatMAE~\cite{cong2022satmae} adapts the masked image model to the temporal and multi-spectral structure of satellite imagery, Scale-MAE~\cite{reed2023scalemae} makes the pretext task explicitly scale-aware so a single backbone serves the huge ground-sample-distance range Earth observation spans, and Prithvi~\cite{jakubik2023foundation} packages such a backbone as a generalist model fine-tunable across flood mapping, burn-scar delineation, and crop segmentation. A single pretrained encoder thus stands in for the dozens of bespoke feature extractors that populate the tasks above.

The language side is moving just as fast. RemoteCLIP~\cite{liu2024remoteclip} aligns remote-sensing imagery with text, giving open-vocabulary retrieval and zero-shot classification without per-class labels; GeoChat~\cite{kuckreja2024geochat} grounds a large vision-language model on overhead imagery so an analyst can query a scene, ask for region-level descriptions, and obtain referring localizations in dialogue; and SegEarth-R1~\cite{li2025segearth} pushes further into pixel-level geospatial reasoning, generating segmentation masks from implicit, reasoning-heavy queries rather than fixed class names. These systems inherit the multimodal reasoning and world-model machinery of Part~\ref{part:3} and specialize it to Earth observation, pointing toward GeoAI workflows driven by natural-language instructions rather than hand-built per-task pipelines.

\subsection{Challenges}
\label{subsec:p4-sec11-2-challenges}
Most of the methods above were trained and tested in one region on one sensor, and that pattern shapes the open problems: whether a model transfers across regions and sensors, how heterogeneous and uneven data are reconciled, whether pipelines run at the scale and latency that operations need, and what new failure modes the geospatial foundation models bring.

\paragraph{Generalization across Regions and Sensors.} Few of the studies surveyed test a model outside the region, season, or sensor it was trained on, so their reported accuracies say little about transfer. Crop calendars, building styles, and spectral responses vary from place to place, and a crop map or flood segmenter trained on one country or one satellite can degrade sharply on another. Pretrained backbones such as Prithvi~\cite{jakubik2023foundation} are motivated by exactly this gap, but systematic cross-region benchmarks remain rare.

\paragraph{Data Heterogeneity and Quality.} Geospatial workflows combine imagery, sensor streams, vector layers, and crowdsourced inputs that differ in resolution, revisit time, and reliability, and the volume of these archives makes curation itself demanding~\cite{boutayeb2024geoai}. Open-access data needs quality control before it can serve as training labels. Geospatial data is also strategically sensitive, which raises questions of data sovereignty, and adversarial manipulation of inputs is a threat that none of the methods above addresses.

\paragraph{Scale and Latency.} Operational services such as flood mapping during an event, precipitation nowcasting, and traffic forecasting must ingest new acquisitions continuously and update their models as satellites deliver data, which strains pipelines built around one-off training runs. When several models are chained in one workflow, validation and deployment cost grow with each link, pushing systems toward distributed, modular architectures.

\paragraph{Open Problems for Geospatial Foundation Models.} The foundation models of the previous subsection bring problems of their own: pretraining that fuses radar, optical, and hyperspectral sensors at very different resolutions; evaluation benchmarks that test geospatial reasoning rather than image-level recognition; and grounding language output in physically consistent, georeferenced coordinates, which dialogue systems such as GeoChat~\cite{kuckreja2024geochat} do not yet guarantee.

\newpage
\section{Aerospace Engineering}
\label{sec:p4_s11_3}

\subsection{Introduction}
\label{subsec:p4-sec11-3-introduction}

Aerospace engineering designs and operates systems that work within and beyond the Earth's atmosphere, and its missions are spatial problems at their core.
Rendezvous and docking (RVD), on-orbit servicing (OOS), debris inspection, and planetary landing and roving all depend on estimating three-dimensional structure and relative motion, often of a non-cooperative target that carries no markers, while communication delays rule out continuous ground control.
Both spacecraft autonomy (\eg, guidance, navigation, and control) and human-in-the-loop operations (\eg, situational awareness for ground controllers) therefore rest on spatial perception and reasoning over geometry, relative motion, semantics, and physical constraints.

Aerospace engineering is thus a demanding deployment target for the perception and reasoning capabilities surveyed in Part~\ref{part:3}: the detection, segmentation, pose-estimation, and 3D-reconstruction techniques of Chapter~\ref{chapter:3:sec:6} and the planning and decision-making methods of Chapter~\ref{sec:reasoning} are here specialized to non-cooperative orbital targets, extreme illumination, and long-horizon autonomy under radiation-hardened compute and communication latency.

\subsection{System Overview}
\label{subsec:p4-sec11-3-system-overview}
The value of spatial intelligence in aerospace lies less in any single perception module than in the closed autonomy loop it serves (Figure~\ref{fig:p4_s11_3_overview}): onboard and ground-based sensing (monocular and multi-view cameras, LiDAR, telemetry) feeds a perception-reasoning-decision core, whose commands change the scene that is sensed next.
The four pivotal tasks treated below are complementary stages of this loop.
\textbf{Satellite pose estimation} supplies the geometric anchor, recovering where a target is and how it is oriented;
\textbf{spacecraft detection, segmentation, and parts recognition} enrich that anchor with structural semantics, identifying what the target is and which functional components it exposes;
\textbf{operator agents} lift this perceptual state into high-level reasoning, planning, and human-facing decision support; and
\textbf{planetary robotics} closes the loop on the surface, coupling perception, localization, and planning into physical exploration and sampling.
Table~\ref{tab:p4_s11_3_benchmarks}, after Task-4, consolidates representative benchmarks and metrics across all four.

\begin{figure}[!th]
    \centering
    \includegraphics[width=0.98\linewidth]{figures/part4/section11/aerospace-overview.pdf}
    \caption{Aerospace spatial intelligence as a closed loop from onboard and ground sensing, through a perception--reasoning--decision core, to complementary pose, structural-semantic, planning, and surface-action outputs. Decisions may request additional observations; mission execution produces subsequent measurements. Pose and semantic estimates are information products rather than physical actions. Satellite views, point clouds, and rover imagery are schematic.}
    \label{fig:p4_s11_3_overview}
\end{figure}


\subsection{Pivotal Task-1: Satellite Pose Estimation}
\label{subsec:p4-sec11-3-pivotal-task-1-satellite-pose-estimation}

Vision-based satellite pose estimation predicts the 6-DoF relative position and orientation of a target spacecraft from monocular or multi-view imagery; it is the geometric anchor of the autonomy loop and a prerequisite for rendezvous, docking, and on-orbit servicing.
It specializes the 6-DoF estimation problem of Chapter~\ref{chapter:3:sec:6} to a hostile setting: harsh illumination, specular reflections, little texture, and a severe gap between the synthetic images models are trained on and the real images they meet in orbit.
Figure~\ref{fig:p4_s11_3_pose} traces the pipeline from method to evaluation.

\begin{figure}[!t]
    \centering
    \includegraphics[width=0.95\linewidth]{figures/part4/section11/aerospace-pose.pdf}
    \caption{Satellite pose estimation: a monocular image of a non-cooperative target passes through alternative feature-based, regression, or hybrid pipelines to a target pose relative to the camera. Translation and relative-rotation errors compare estimated and reference poses; 2D reprojection compares detected image keypoints with projected model points. The satellite render and keypoint residuals are schematic; no numerical benchmark scores are implied.}
    \label{fig:p4_s11_3_pose}
\end{figure}

Traditional methods relied on hand-engineered features (\eg, Harris corners, SIFT, SURF, and ORB) or model-based tracking~\cite{d2014poseestimationuncooperativespacecraft}, which are sensitive to noise and domain shift, and deep learning has since become dominant~\cite{proencca2020deepURSO}.
The first deep models regress the pose directly with CNNs~\cite{garcia2021lspnet}, as in the Spacecraft Pose Network (SPN)~\cite{sharma2020neuralposeestimation}, a GoogLeNet-based network~\cite{phisannupawong2020vision}, and a VGG backbone~\cite{sonawani2020assistive}, trained on synthetic data and on laboratory imagery that simulates an on-orbit assembly operation.
SPNv2~\cite{park2024robustspacecraftposeestimation} made multi-task learning the prevailing design: its shared encoder jointly predicts keypoints, a direct pose regression, and foreground segmentation, and it narrows the synthetic-to-real gap by online domain refinement, adapting the network's normalization layers on unlabeled target imagery at deployment.
Hybrid pipelines estimate a coarse pose with a deep model and refine it with geometry-based optimization, keeping the final estimate physically consistent; early work classified discretized viewpoints with a CNN~\cite{sharma2018pose}, and the same coarse-to-fine pattern now extends to 6-DoF tracking with stereo event cameras~\cite{liu2025stereo}.

SPEED~\cite{Sharma2019SpacecraftPE} and SPEED+~\cite{park2022speed+} are the standard testbeds (Table~\ref{tab:ssa_datasets}): both pair a large synthetic training set with laboratory hardware-in-the-loop test images, so that the synthetic-to-real gap is measured rather than assumed, and the Satellite Pose Estimation Challenge built on them fixed common training regimes, input formats, and metrics~\cite{kisantal2020satellitechallenge}.
CubeSat-CDT~\cite{musallam2022cubesatcdt} adds trajectories of a symmetric target, so that temporal methods and pose ambiguity can be evaluated.
Evaluation reports translation and rotation error as defined for SPEED~\cite{Sharma2019SpacecraftPE}, often with 2D reprojection or keypoint error~\cite{garcia2021lspnet}; results are given separately on synthetic and real test sets, and robustness to illumination, occlusion, and background clutter is probed explicitly~\cite{park2024robustspacecraftposeestimation}.

\begin{table}[!ht]
\centering
\caption{Overview of existing pose estimation datasets. \textit{sim} refers to simulated images and \textit{lab} to laboratory data. SPARK's 11 classes are 10 spacecraft and one debris class.}
\label{tab:ssa_datasets}
\fontsize{9}{11}\selectfont
\setlength{\tabcolsep}{1.0mm}
\begin{tabular}{L{0.17\textwidth} C{0.08\textwidth} C{0.11\textwidth} C{0.08\textwidth} C{0.13\textwidth} C{0.10\textwidth} C{0.11\textwidth} C{0.08\textwidth}}
\toprule
\multirow{2}{*}{\textbf{Dataset}} & \multicolumn{2}{c}{\textbf{Image}} & \multirow{2}{*}{\textbf{Classes}} & \multirow{2}{*}{\textbf{Resolution}}  & \multirow{2}{*}{\textbf{Rendering}} & \multirow{2}{*}{\textbf{Trajectories}} & \multirow{2}{*}{\textbf{Public}} \\
\cmidrule(r){2-3}
& \textbf{Synthetic} & \textbf{Non-synthetic} & & & & & \\
\midrule
SPARK~\cite{musallam2021spacecraft} & 150k & - & 11 & $1024 \times 1024$ & sim & \crossmark & \checkmark  \\
SPEED~\cite{Sharma2019SpacecraftPE} & 15k & 305 & 1 & $1920 \times 1200$  & sim + lab & \crossmark & \checkmark \\
SPEED+~\cite{park2022speed+} & 60k & 10k & 1 & $1920 \times 1200$ & sim + lab & \crossmark & \checkmark \\
URSO~\cite{proencca2020deepURSO} & 15k & - & 2 & $1080 \times 960$ & sim & \crossmark & \checkmark \\
SwissCube~\cite{hu2021widedepthrange} & 50k & - & 1 & $1024 \times 1024$ & sim & \checkmark & \checkmark \\
Cygnus~\cite{black2021realCygnus} & 20k & 540 & 1 & $1024 \times 1024$ & sim + real  & \crossmark & \crossmark \\
Prisma12K~\cite{park2019towardsPrisma12K} & 12k & - & 1 & $752 \times 580$ & sim  & \crossmark & \crossmark \\
Prisma25~\cite{park2019towardsPrisma12K} & - & 25 & 1 & -NA- & real  & \crossmark & \crossmark \\
CubeSat-CDT~\cite{musallam2022cubesatcdt} & 14k & 8k & 1 & $1440 \times 1080$ & sim + lab & \checkmark & \checkmark \\
\bottomrule
\end{tabular}
\end{table}


\subsection{Pivotal Task-2: Spacecraft Detection, Segmentation, and Parts Recognition}
\label{subsec:p4-sec11-3-pivotal-task-2-spacecraft-detection-segmentation-and-parts-recognition}

Detection, segmentation, and parts recognition enrich the geometric anchor of Task-1 with structural semantics, telling what the target is and which functional components it exposes; proximity operations, on-orbit servicing, debris inspection, and collision avoidance all depend on this.
The detection and segmentation machinery of Chapter~\ref{chapter:3:sec:6} meets the obstacles of Task-1 here (strong illumination variation, specular metallic surfaces, little texture) and, in addition, a scarcity of annotated real imagery.
The stack runs from coarse presence to fine structure: detection localizes the spacecraft, segmentation delineates its silhouette, and parts recognition resolves its components, as Figure~\ref{fig:p4_s11_3_detection} illustrates.

\begin{figure}
    \centering
    \includegraphics[width=0.95\linewidth]{figures/part4/section11/aerospace-detection.pdf}
    \caption{The ``detect--segment--recognize'' perception stack for non-cooperative spacecraft: (1) object localization, (2) encoder--decoder silhouette segmentation, and (3) recognition of solar panels, the main body, antenna, and thruster. The same synthetic target is used throughout. The silhouette is an illustrative foreground mask, not a measured prediction; arrows indicate increasing semantic detail rather than requiring every implementation to run these stages serially.}
    \label{fig:p4_s11_3_detection}
\end{figure}

Convolutional and transformer architectures now carry all three stages~\cite{shao2023satellite}.
Detection is formulated as bounding-box or keypoint localization, with lightweight saliency-driven detectors for onboard compute, while encoder--decoder networks, pyramid attention modules, and hybrid 2D/3D architectures delineate silhouettes and component boundaries~\cite{zhao20223dnoncooperativespacecraft}.
Parts recognition (solar panels, antennas, thrusters, reflectors) often adds semantic priors or 3D geometric reasoning to stay robust under occlusion.
Two recent lines reduce the dependence on manual labels: prompt-guided Segment Anything (SAM) automates the annotation of space targets~\cite{zhang2024automaticspacetarget}, and diffusion priors support few-shot segmentation and adaptation to unseen spacecraft~\cite{li2025exploitingdiffusion}.

The datasets in Table~\ref{tab:satellite_datasets_datasetwise} differ along a few axes.
Labels range from whole-object masks (SwissCube~\cite{hu2021widedepthrange}, a pose dataset that also provides CubeSat masks, and SPARK~\cite{musallam2021spacecraft}) and object-level boxes (STAR-24K~\cite{zhang2022star24K}) to component labels, from a single component (the solar panels of the Tsinghua dataset~\cite{li2020noncooperativeobjectdetection}) up to nine (Spacecraft-DS~\cite{cao2024spacecraftds}).
Most are rendered (Unity3D for SPARK, Unreal Engine~4 for UESD~\cite{zhao2023intelligentUESD}), whereas Spacecraft-DS uses hardware-in-the-loop capture to narrow the synthetic-to-real gap.
The Adelaide dataset~\cite{dung2021spacecraftdataset} fixed the common three-component scheme (solar panel, main body, antenna), which Shao \etal~\cite{shao2023satellite} extended to video sequences rendered under three illumination conditions and several relative-motion modes; Zhao \etal~\cite{zhao20223dnoncooperativespacecraft} carried the same labels to 3D point clouds, and NCSTP~\cite{liu2025largencstp} raised the scale to 200,000 images of satellites, debris, and rocks.
Detection is scored by mean average precision and segmentation by IoU, as in Chapter~\ref{chapter:3:sec:6}.

\begin{table}[!ht]
\centering
\caption{Summary of satellite recognition datasets. `-' indicates that the exact quantity is unknown; the Zhao \etal dataset consists of point clouds rather than images.}
\label{tab:satellite_datasets_datasetwise}
\fontsize{9}{11}\selectfont
\setlength{\tabcolsep}{1.5mm}
\begin{tabular}{L{0.32\textwidth} C{0.12\textwidth} C{0.16\textwidth} C{0.16\textwidth} C{0.13\textwidth}}
\toprule
\textbf{Dataset}
& \textbf{\#Images}
& \textbf{\#Satellite Models}
& \textbf{Segmentation Level}
& \textbf{\#Components}   \\
\midrule
SwissCube~\cite{hu2021widedepthrange} & 50{,}000 & 1 & Global & 0  \\
UESD~\cite{zhao2023intelligentUESD} & 10{,}000 & 33 & Component & 5 \\
SPARK~\cite{musallam2021spacecraft} & 150{,}000 & 10 & Global & 0  \\
STAR-24K~\cite{zhang2022star24K} & 24{,}451 & - & Global & - \\
NCSTP~\cite{liu2025largencstp} & 200{,}000 & 16 & Component & 4 \\
Tsinghua Satellite Dataset~\cite{li2020noncooperativeobjectdetection} & 18{,}381 & - & Component & 1 \\
Adelaide Satellite Dataset~\cite{dung2021spacecraftdataset} & 3{,}117 & - & Component & 3   \\
Satellite Component Video Dataset~\cite{shao2023satellite} & 32{,}402 & 13 & Component & 3  \\
Spacecraft-DS~\cite{cao2024spacecraftds} & 1{,}602 & 2 & Component & 9 \\
3D component point clouds (Zhao \etal)~\cite{zhao20223dnoncooperativespacecraft} & - & 33 & Component & 3 \\
\bottomrule
\end{tabular}
\end{table}


\subsection{Pivotal Task-3: Operator Agents}
\label{subsec:p4-sec11-3-pivotal-task-3-operator-agents}

Operator agents lift the perceptual state recovered by Tasks 1 and 2 into reasoning, planning, and decision support, monitoring, interpreting, and controlling spacecraft systems to reduce operator workload and allow more autonomous mission execution.
They are the aerospace instance of the planning and decision-making capabilities of Chapter~\ref{sec:reasoning}.
Rooted in rule-based and model-driven architectures, their role has expanded from health monitoring and scheduling to high-level control in dialogue with human experts, a history that runs from symbolic planning through multi-agent coordination to foundation-model reasoning (Figure~\ref{fig:p4_s11_3_operator}).

\begin{figure}
    \centering
    \includegraphics[width=0.95\linewidth]{figures/part4/section11/aerospace-operator.pdf}
    \caption{Three paradigms in the development of spacecraft operator agents: (1) symbolic planning and expert systems, (2) multi-agent coordination across a constellation, and (3) LLM- and RAG-based agents that combine telemetry with retrieved domain knowledge to propose natural-language plans. Inter-panel arrows indicate historical development, not required dataflow. In the rule example, observations and action outcomes support the goal; blue peer links and dashed ground links denote communication. Satellite imagery and rules are illustrative.}
    \label{fig:p4_s11_3_operator}
\end{figure}

Early operator-agent architectures were built on symbolic reasoning, constraint-based planning, and expert-system principles.
The Remote Agent Experiment~\cite{bernard1998designremoteagent} and early autonomous spacecraft-agent prototypes~\cite{pell1997autonomousspacecraftagent} demonstrated onboard planning, diagnosis, and execution in deep-space missions.
As autonomy requirements grew, multi-agent planning frameworks~\cite{colby2016autonomousmultiagent,picard2021autonomousAgents} were introduced to coordinate heterogeneous spacecraft, manage distributed sensing, and schedule jointly across satellite constellations.
Large language models changed the design: recent work~\cite{rodriguez2024languagemodelsspacecraft,carrasco2025visualoperatoragents} uses LLMs as operator agents that interpret telemetry, generate action plans, respond to anomalies, and follow natural-language mission specifications, and fine-tuned LLMs~\cite{zucchelli2025finetunellmSpace,carrasco2024finetuneAutonomousSpacecraft,foutter2024spacellava,carrasco2025llmAutonomousSpacecraft} have been tested in simulators such as Kerbal Space Program (KSP) for guidance, navigation, and control. To meet the interpretability and safety requirements of mission-critical decisions, other studies integrate retrieval-augmented generation (RAG)~\cite{belo2025languageragSpaceSpace}, domain knowledge graphs~\cite{powell2025questionknowledgegraphssatellite}, and structured symbolic reasoning, while hybrid architectures such as MCP-based multi-component agents~\cite{wu2025spacemind} combine small specialized models with foundation models to balance computational cost and reasoning depth, particularly for on-orbit servicing.

Evaluation environments followed the same path.
The Remote Agent Experiment~\cite{bernard1998designremoteagent} supplied real mission scenarios for onboard planning and diagnosis, and multi-agent planning studies~\cite{colby2016autonomousmultiagent} supply constellation-level scheduling, resource-allocation, and communication constraints.
LLM agents are tested in simulation, chiefly in KSP~\cite{carrasco2025llmAutonomousSpacecraft} and in the KSPDG suite of SpaceGym~\cite{allen2023spacegym}, a set of multi-agent pursuit--evasion differential games built in the KSP engine.
Classical agents are scored on planning optimality, constraint satisfaction, fault-diagnosis accuracy, and task completion; LLM agents are additionally judged on adherence to safety constraints, robustness to shifts in telemetry or mission context, and the interpretability of their decisions, which is why uncertainty estimation and human-in-the-loop verification recur throughout this work.


\subsection{Pivotal Task-4: Planetary Robotics}
\label{subsec:p4-sec11-3-pivotal-task-4-planetary-robotics}

Planetary robotics closes the autonomy loop on the surface, coupling the perception, localization, and planning capabilities of Chapters~\ref{chapter:3:sec:6} and~\ref{sec:reasoning} into exploration, navigation, and sampling on the Moon, Mars, and small bodies.
Unlike orbital operations, surface missions face uneven and unstructured terrain, delayed and limited communication, restricted sensing, and long stretches of autonomy without continuous human supervision.
Figure~\ref{fig:p4_s11_3_planetary} layers the autonomy stack from perception and mapping through hierarchical decision-making to vision-guided manipulation.

\begin{figure}
    \centering
    \includegraphics[width=0.99\linewidth]{figures/part4/section11/aerospace-planetary.pdf}
    \caption{The planetary-robotics autonomy stack: visual-odometry feature tracks and 3D terrain mapping support hierarchical decisions and vision-guided manipulation. Mission reasoning handles goals, task decomposition, resources, and replanning; path planning uses global/local maps and constraints; reactive locomotion handles obstacles, stability, and terrain adaptation. Solid downward links convey plans, and dashed upward links convey status. Manipulation yields samples, material observations, environment updates, and task progress. Scenes and map geometry are illustrative.}
    \label{fig:p4_s11_3_planetary}
\end{figure}

Surface robots are organized around a mobility stack that couples perception, localization, mapping, planning, and control, with manipulation and science operations integrated at the task level.
Feature tracking, visual odometry, and stereo reconstruction estimate rover motion and build local 3D models of rugged terrain, a capability flown since the Mars Exploration Rovers~\cite{maimone2007two}.
The Enhanced AutoNav (ENav) system on NASA's Perseverance rover~\cite{toupet2025enav} builds onboard 3D terrain maps, evaluates hazards, and plans safe local traverses under severe computational and communication constraints, and it evaluated terrain over a large majority of the rover's traversed distance.
This autonomy remains embedded in a human-supervised workflow: ground operators specify science goals and route segments, and onboard autonomy resolves local navigation and hazard avoidance.

Research and analogue-testbed systems extend these capabilities to multiple robots, where cooperative visual navigation and shared mapping improve coverage and robustness~\cite{bahraini2021autonomousplanetaryrobots}.
Hierarchical frameworks decompose behavior across temporal and functional layers, balancing reactive locomotion and obstacle avoidance with deliberative path planning, task scheduling, and mission-level reasoning under safety, power, communication, and resource constraints.

Vision-guided manipulation supports sample collection and instrument deployment by coupling image-based pose estimation with position-controlled or compliant manipulation suited to contact with natural terrain~\cite{nickels2010visionplanetaryrobotics}.
Learning-based methods are increasingly explored for terrain classification, slip prediction, and traversability estimation, but model-based perception, localization, and planning remain central: safety requirements are stringent, mission-representative labeled data is scarce, onboard computation is constrained, and purely data-driven methods are hard to validate for mission-critical use.
Evaluation therefore combines physics-based simulators, mission-derived datasets, and laboratory or analogue-field testbeds that model illumination, dust, wheel--soil interaction, and communication delay, and it reports localization accuracy, map quality, hazard-aware traversal, task completion, and energy use.

\begin{table}[!ht]
\centering
\caption{Representative benchmarks and headline metrics for the four pivotal aerospace tasks. The headline column names the standard evaluation axis, not a single best score.}
\label{tab:p4_s11_3_benchmarks}
\fontsize{8}{10}\selectfont
\setlength{\tabcolsep}{1.0mm}
\begin{tabular}{L{0.18\textwidth} L{0.24\textwidth} L{0.30\textwidth} L{0.20\textwidth}}
\toprule
\textbf{Pivotal task} & \textbf{Benchmark} & \textbf{Representative method / system} & \textbf{Headline metric} \\
\midrule
Pose estimation & SPEED / SPEED+~\cite{Sharma2019SpacecraftPE,park2022speed+} & Multi-task / domain-adaptive CNN~\cite{park2024robustspacecraftposeestimation} & Translation + rotation error \\
Pose estimation & CubeSat-CDT~\cite{musallam2022cubesatcdt} & Cross-domain regression~\cite{garcia2021lspnet} & 2D reprojection error \\
Detection \& parts & Adelaide / SPARK~\cite{dung2021spacecraftdataset,musallam2021spacecraft} & Encoder-decoder segmentation~\cite{zhao20223dnoncooperativespacecraft} & mAP \\
Detection \& parts & NCSTP~\cite{liu2025largencstp} & SAM-based prompt segmentation~\cite{zhang2024automaticspacetarget} & mAP / IoU \\
Operator agents & KSPDG / SpaceGym~\cite{allen2023spacegym} & Fine-tuned LLM agent~\cite{carrasco2025llmAutonomousSpacecraft} & Maneuver success rate \\
Operator agents & Deep Space 1 flight experiment~\cite{bernard1998designremoteagent} & Symbolic planner / diagnosis~\cite{pell1997autonomousspacecraftagent} & Task completion \\
Planetary robotics & Mars rover traverses~\cite{toupet2025enav} & Visual odometry~\cite{maimone2007two} & Localization drift \\
\bottomrule
\end{tabular}
\end{table}


\subsection{Established Tools and Systems}
\label{subsec:p4-sec11-3-established-tools-and-systems}

\textbf{Simulators and rendering engines.}
Because real orbital imagery is scarce and flight opportunities are expensive, most perception research is trained and validated in simulation.
General-purpose game and physics engines have been repurposed as space testbeds: Unity3D and Blender drive the synthetic rendering pipelines behind the SPARK~\cite{musallam2021spacecraft} and satellite-component~\cite{shao2023satellite} datasets, while the Kerbal Space Program\footnote{\url{https://www.kerbalspaceprogram.com/}} engine underpins LLM-agent guidance-navigation-control studies~\cite{carrasco2025llmAutonomousSpacecraft}.
For decision-making, SpaceGym\footnote{\url{https://github.com/mit-ll/spacegym-kspdg}}~\cite{allen2023spacegym} packages its OrbitDefender2D gridworld games and the KSPDG suite as standardized non-cooperative game environments, and robotics simulators such as NVIDIA Isaac Sim\footnote{\url{https://developer.nvidia.com/isaac/sim}} and Gazebo\footnote{\url{https://gazebosim.org/}} provide the wheel-soil and sensor models used to prototype planetary-rover autonomy.

\textbf{Public challenges.}
The Satellite Pose Estimation Challenge~\cite{kisantal2020satellitechallenge}, run by Stanford's Space Rendezvous Laboratory and the European Space Agency on its Kelvins platform\footnote{\url{https://kelvins.esa.int/}}, set the protocols that pose estimation now follows, and its SPEED+ edition~\cite{park2022speed+} made the synthetic-to-real gap the explicit test.

\textbf{Flight-proven autonomy systems.}
The NASA Remote Agent~\cite{bernard1998designremoteagent}, flown on Deep Space 1, remains the canonical reference for onboard planning, diagnosis, and execution.
On planetary surfaces, rover visual odometry~\cite{maimone2007two} and the Enhanced AutoNav stack flown on Perseverance~\cite{toupet2025enav} set the operational baseline against which hierarchical-autonomy and learning-augmented methods are measured, and current work pairs such flight-qualified stacks with foundation-model agents~\cite{wu2025spacemind}.


\subsection{Challenges}
\label{subsec:p4-sec11-3-challenges}
Orbital and planetary environments differ from terrestrial ones in ways that terrestrial data does not cover, and the open problems follow the pipeline: the data gap that limits perception, the compute budget that limits deployment, and the trustworthiness requirements that limit certification.

\paragraph{The Unresolved Sim-to-Real Domain Gap.}
Autonomous space systems, such as those performing on-orbit servicing or debris removal, are trained largely on simulated data, and simulators do not reproduce the high-dynamic-range lighting of orbit (\eg, solar glare and deep shadow), intricate material reflections, or sensor degradation such as radiation noise.
Models consequently lose accuracy on real imagery, as the hardware-in-the-loop test sets of the pose-estimation challenges show~\cite{kisantal2020satellitechallenge,park2022speed+}, which keeps domain adaptation and domain randomization central to the field.

\paragraph{Computational Constraints and Real-time Operation.}
Spacecraft hardware is bound by size, weight, and power (SWaP) limits and must use radiation-hardened processors that are far less powerful than ground-based GPUs.
Running 6-DoF pose estimation or SLAM at the update rates that proximity operations require on such computers calls for model compression, lightweight architectures, and hardware-aware optimization that preserve accuracy while cutting latency and power consumption.

\paragraph{Assuring Trustworthiness and Robustness in Extreme Conditions.}
Space missions are safety-critical, so perception must hold up under extreme visual ambiguity, such as estimating the pose of highly symmetric targets like CubeSats~\cite{musallam2022cubesatcdt} or navigating where standard reference signals are unavailable.
Deployed systems also explain little about their own failures, so operators cannot easily trace a perception anomaly (\eg, a pose error caused by sensor saturation) to its root cause in time to intervene.
Formal verification and interpretable models that guarantee predictable behavior across operating extremes remain open problems.

\newpage
\chapter{Spatial Intelligence for Biomedicine and Natural Sciences}


\FramedInitial{B}{iomedicine} and the natural sciences are, at their core, spatial disciplines: their objects are structures in space, from organs and tissues down to cells and molecules, and understanding them means reasoning about shape, arrangement, and interaction in three dimensions. This chapter turns the spatial-intelligence capabilities of the previous part onto that setting, where in both of its cases the unit of value is the same: a metrically faithful 3D model of a specific physical system, kept in correspondence with real measurements.

In medicine, that model is patient-specific, and it drives a closed loop running from imaging through planning to image-guided intervention. In the natural sciences the same perception and reconstruction machinery operates several orders of magnitude smaller, at the molecular and cellular scale, where the structures reconstructed are the very objects of scientific explanation.


\section{Medicine and Clinical Healthcare}
\label{sec:p4_s12_1}

\subsection{Introduction}
\label{subsec:p4-sec12-1-introduction}

Anatomy is a three-dimensional arrangement of organs with shape, volume, and topology; disease is \emph{localized}, so that a lesion is a position and an extent in the body frame rather than a label on an image; and intervention is fundamentally an act of navigation, guiding an instrument to a target through tissue that must not be harmed.

The perception machinery of Chapter~\ref{chapter:3:sec:6} (detection, segmentation, depth, pose) and the reconstruction and generation of Chapter~\ref{sec:generation} reappear here, specialized to volumetric data, to registration across modalities and time, and to a regulatory setting in which a wrong answer harms a patient. 
The unit of analysis is the individual patient and the individual procedure: the molecular and cellular scale belongs to the natural sciences (§\ref{sec:Natural Science}), the generic digital twin, of which the patient-specific medical twin below is a specialization, to §\ref{sec:p4_s11_1}, and the population scale of spatial epidemiology to the GeoAI of §\ref{sec:p4_s11_2}.

\subsection{System Overview}
\label{subsec:p4-sec12-1-system-overview}

Figure~\ref{fig:p4_s12_1_overview} draws clinical spatial intelligence as a closed loop around the patient: multimodal imaging (computed tomography, magnetic resonance, ultrasound, endoscopy, pathology) is acquired and the anatomy perceived, a patient-specific 3D model is built and a plan formed, the plan is executed under image guidance, and the outcome is imaged anew. 
Registration connects the image, patient, and instrument frames throughout this pipeline; the recurring difficulty is keeping the virtual model aligned with the physical patient as tissue moves and deforms. The three pivotal tasks map onto the loop directly (Task-1 perceives, Task-2 models and plans, Task-3 acts), and Table~\ref{tab:p4_s12_1_benchmarks} collects their benchmarks.

\begin{figure}
    \centering
    \includegraphics[width=0.98\linewidth]{figures/part4/section12/medical_overview.pdf}
    \caption{Clinical spatial intelligence with two connected information paths: acquired observations support anatomical perception and patient-specific planning, while live imaging and tracking update the alignment used for guided intervention. Physical intervention changes the patient state, prompting subsequent observation. Anatomical scenes are schematic; pathology may require specimen correspondence.}
    \label{fig:p4_s12_1_overview}
\end{figure}

\subsection{Pivotal Task-1: Volumetric Medical Image Perception}
\label{subsec:p4-sec12-1-pivotal-task-1-volumetric-medical-image-perception}

The first task perceives the anatomy. Medical perception is intrinsically volumetric: it converts a voxel grid into delineated 3D anatomical objects with measurable shape and volume, localizes disease in the body frame, and aligns observations across modalities and time. Figure~\ref{fig:p4_s12_1_perception} illustrates three complementary outputs: segmented organs, a localized candidate finding, and a cross-modal alignment; these tasks need not form a serial pipeline.

\begin{figure}
    \centering
    \includegraphics[width=0.98\linewidth]{figures/part4/section12/medical_perception.pdf}
    \caption{Three input-operation-output views of medical perception: (a) image volumes yield voxel labels, organ masks, and 3D anatomy; (b) detected candidate regions are mapped from voxel indices to physical patient coordinates using image geometry; and (c) moving and reference modalities jointly determine a spatial transform for anatomical alignment. Synthetic CT/MRI-style images and the hypothetical ROI illustrate task structure, not measured clinical findings or a validated registration result.}
    \label{fig:p4_s12_1_perception}
\end{figure}

\textbf{Volumetric segmentation.} Segmentation is the spatial core of medical imaging, and its history is a move from two-dimensional slices to full volumes. 
The encoder-decoder U-Net~\cite{ronneberger2015unet} founded the area, and its volumetric successors, the 3D U-Net~\cite{cciccek20163d} and the V-Net~\cite{milletari2016vnet} (which introduced the now-ubiquitous Dice loss to handle voxel class imbalance), moved it into three dimensions. 
The self-configuring nnU-Net~\cite{isensee2021nnunet} adapts preprocessing, architecture, and training to each dataset and remains the baseline that is hard to beat, while transformer encoders such as UNETR~\cite{hatamizadeh2022unetr} and Swin-UNETR~\cite{hatamizadeh2022swinunetr} add global context. 
At the whole-body scale, TotalSegmentator~\cite{wasserthal2023totalsegmentator} delineates more than a hundred anatomical structures in a single CT, functioning as a spatial atlas in routine use. 
These methods are evaluated by the Dice similarity coefficient for volumetric overlap and by boundary-geometry metrics, the 95th-percentile Hausdorff distance and the average symmetric surface distance, that capture shape errors which overlap alone misses~\cite{taha2015metrics}.

\textbf{Medical foundation models.} The recent shift mirrors the rest of vision: from per-task specialists to promptable foundation models that segment arbitrary anatomy from spatial prompts. 
The Segment Anything Model~\cite{kirillov2023sam} catalyzed the move, and MedSAM~\cite{ma2024medsam} adapted it by fine-tuning on more than a million medical image--mask pairs. 
Because anatomy is volumetric, the decisive step is into 3D: SAM-Med3D~\cite{wang2023sammed3d} replaces the slice-wise encoder with a fully three-dimensional one, and SegVol~\cite{du2024segvol} supports both spatial and textual prompts over hundreds of anatomical categories. 
Much of this work is built on the open MONAI framework~\cite{cardoso2022monai}, whose medical-aware 3D transforms and model zoo have become standard infrastructure.

\textbf{Detection and landmark localization.} Detection localizes a finding in the body frame, and landmark localization fixes canonical anatomical coordinates; both are spatial tasks distinct from whole-image classification. 
Lung-nodule detection is anchored by the LIDC-IDRI reference database~\cite{armato2011lidcidri} and the LUNA16 benchmark~\cite{setio2017luna16}, which standardized evaluation on the free-response operating characteristic (sensitivity against false positives per scan). 
Anatomical landmarks can be found by treating localization as a sequential spatial search, as in the reinforcement-learning agent of Ghesu \etal~\cite{ghesu2016landmark} that navigates the volume to a target.

\textbf{Registration as spatial alignment.} Registration brings two images into a common coordinate frame and is foundational for multimodal fusion, atlas-based labeling, longitudinal tracking, and surgical planning. 
Learning-based methods reframe it as predicting a dense deformation field in a single forward pass: VoxelMorph~\cite{balakrishnan2019voxelmorph} established the unsupervised CNN formulation, and TransMorph~\cite{chen2022transmorph} added a transformer backbone. 
These build on classical diffeomorphic and intensity-based methods, principally the symmetric diffeomorphic SyN~\cite{avants2008syn} and the intensity-based toolbox elastix~\cite{klein2010elastix}, and are measured by the target registration error between corresponding landmarks. 
Segmentation and detection are benchmarked on generalist suites such as the Medical Segmentation Decathlon~\cite{antonelli2022msd}, the brain-tumor benchmark BraTS~\cite{menze2015brats}, the abdominal multi-organ datasets AMOS~\cite{ji2022amos} and BTCV~\cite{landman2015btcv}, the kidney-tumor KiTS19~\cite{heller2021kits19}, and the cardiac ACDC~\cite{bernard2018acdc}.

\subsection{Pivotal Task-2: Anatomical Reconstruction, Modeling, and Surgical Planning}
\label{subsec:p4-sec12-1-pivotal-task-2-anatomical-reconstruction-modeling-and-surgical-planning}

The second task turns the perceived anatomy into a patient-specific 3D model that can be measured, simulated, and planned upon, bridging from pixels to clinical decisions. Figure~\ref{fig:p4_s12_1_modeling} previews this progression from a segmentation to an anatomical mesh, a patient-specific simulation, and a surgical plan.

\begin{figure}
    \centering
    \includegraphics[width=0.99\linewidth]{figures/part4/section12/medical_reconstruction.pdf}
    \caption{Anatomical reconstruction, modeling, and planning: a volumetric segmentation becomes a 3D anatomical mesh, is personalized into a medical digital twin (\eg, a cardiac electrophysiology or blood-flow model), and is used to simulate outcomes and design a surgical or radiotherapy plan. New patient observations update model personalization. Cardiac targeting and radiotherapy dose prediction are alternative planning examples; synthetic anatomy and vector overlays illustrate the workflow rather than clinical results.}
    \label{fig:p4_s12_1_modeling}
\end{figure}

\textbf{3D anatomical reconstruction.} A segmentation becomes a usable model once it is turned into geometry. 
The classical route extracts an isosurface mesh by Marching Cubes~\cite{lorensen1998marching}, optionally constrained by a statistical shape model that encodes anatomical priors as a learned distribution over landmarks, an idea founded by Active Shape Models~\cite{cootes1995asm} and surveyed for medical use by Heimann and Meinzer~\cite{heimann2009ssm}. 
Deep methods deform a template mesh directly from the volume, as in Voxel2Mesh~\cite{wickramasinghe2020voxel2mesh} and the whole-heart reconstruction of Kong \etal~\cite{kong2021meshdeformnet}. 
Implicit neural representations supply a continuous alternative, with DeepSDF~\cite{park2019deepsdf} and Occupancy Networks~\cite{mescheder2019occupancy} as the substrate (Chapter~\ref{sec:generation}) and medical adaptations such as the neural attenuation fields of NAF~\cite{zha2022naf} reconstructing volumes from sparse-view cone-beam CT.

\textbf{Patient-specific modeling and the medical digital twin.} A geometric model becomes a \emph{twin} when it is endowed with patient-specific physiology and kept faithful to a particular patient, a medical specialization of the generic digital twin of §\ref{sec:p4_s11_1}. 
What is clinically specific is that the model is built from the patient's own imaging and physiological signals, governed by mechanistic physiology rather than CAD, and bound by regulatory credibility. 
Cardiac twins are the most mature instance: the vision was set out by Corral-Acero \etal~\cite{corralacero2020digital}, the mechanistic backbone is the multiscale cardiac modeling reviewed by Niederer \etal~\cite{niederer2019cardiacmodels}, and a personalized virtual heart built from scar imaging has been used to guide ablation of ventricular tachycardia~\cite{prakosa2018personalized}. 
Hemodynamic twins compute fractional flow reserve from coronary CT by simulating blood flow~\cite{taylor2013ffrct}, and the paradigm extends to the brain, where a virtual epileptic patient estimates the epileptogenic zone to guide surgery~\cite{wang2023vep}. 
Such models also underwrite in-silico trials, in which virtual patient cohorts partially substitute for physical ones~\cite{viceconti2016avicenna}, a route that regulators have begun to accept, as in the VICTRE imaging trial~\cite{badano2018victre}. 
The broader landscape is surveyed in recent reviews of digital twins for health~\cite{katsoulakis2024digital}.

\textbf{Surgical and treatment planning.} The patient model is finally used to plan. 
Open platforms such as 3D Slicer~\cite{fedorov2012slicer} assemble segmentation, registration, and 3D modeling into a planning workflow, and real-time finite-element engines such as SOFA~\cite{faure2012sofa} support interactive surgical simulation and rehearsal on deformable tissue. 
In radiotherapy, planning is increasingly cast as a learning problem of predicting an achievable dose distribution from anatomy, standardized by the OpenKBP challenge~\cite{babier2021openkbp}. 
Across these settings the plan is a spatial object, a trajectory, a dose map, or a resection boundary, defined on the patient's own anatomy.

\subsection{Pivotal Task-3: Image-Guided Intervention and Surgical Navigation}
\label{subsec:p4-sec12-1-pivotal-task-3-image-guided-intervention-and-surgical-navigation}

The third task brings the plan into the operating room and acts within the patient's live coordinate frame. 
The central problem is patient-to-image registration under real-world deformation and latency: the preoperative model must be aligned to the physical patient, and instruments tracked within that frame, in real time. 
Figure~\ref{fig:p4_s12_1_navigation} previews registered preoperative anatomy and augmented-reality guidance alongside the complementary reconstruction of a deforming endoscopic scene.

\begin{figure}[!ht]
    \centering
    \includegraphics[width=0.98\linewidth]{figures/part4/section12/medical_navigation.pdf}
    \caption{Patient-frame guidance combines registered anatomy and a calibrated tracked-tool pose. New intraoperative observations update anatomical registration; display calibration supports AR presentation. A separate endoscopic branch reconstructs dynamic geometry and registers it to the patient before optional fusion. The robot controller consumes the fused state rather than raw endoscopic imagery. Dashed connections denote optional downstream use. Anatomical scenes, overlays, and silicone-phantom reconstruction are schematic; transforms act on homogeneous coordinates.}
    \label{fig:p4_s12_1_navigation}
\end{figure}

\textbf{Surgical navigation.} Navigation aligns the preoperative model to the patient and tracks instruments within that frame, typically by optical tracking, which is sub-millimeter but requires line of sight, or electromagnetic tracking, which tolerates occlusion but is perturbed by metal. 
The defining difficulty is deformation: in neuronavigation, the brain shifts by up to centimeters once the skull is opened, invalidating the rigid preoperative registration, which is why intraoperative imaging and learning-based compensation are active concerns.

\textbf{Augmented-reality guidance.} Rather than divert the surgeon's gaze to a screen, augmented reality overlays the segmented anatomy directly on the surgical field, a use of the XR runtime treated in §\ref{sec:p4_s10_1}. 
The approach is most mature in spine surgery, where the Augmedics xvision system, the first FDA-cleared augmented-reality navigation platform, projects the patient's CT-derived spine onto a transparent head-mounted display for pedicle-screw placement~\cite{molina2021xvision}.

\textbf{Endoscopic spatial perception.} Minimally invasive surgery presents a hard perception problem: a monocular, moving camera views a deforming, specular, poorly lit scene. 
Surgical workflow and instruments are recognized from video, as standardized by the Cholec80 dataset and the EndoNet architecture for joint tool and phase recognition~\cite{twinanda2017endonet}. 
Geometry is recovered by self-supervised depth and visual odometry, benchmarked by the SCARED challenge~\cite{allan2021scared} and the EndoSLAM dataset~\cite{ozyoruk2021endoslam}, and the deforming surgical scene is reconstructed in 3D by dynamic neural rendering, as in EndoNeRF~\cite{wang2022endonerf}, a clinical specialization of the neural reconstruction of Chapter~\ref{sec:generation}.

\textbf{Surgical robotics.} Robotic systems execute the plan within the patient frame; their manipulation and control are the subject of §\ref{sec:p4_s9_1}, and the concern here is the spatial perception and autonomy they enable. 
The da Vinci Research Kit~\cite{kazanzides2014dvrk} is the standard open research platform, and autonomy is commonly graded on a six-level scale from teleoperation to full autonomy~\cite{yang2017autonomy}. 
Most deployed systems remain teleoperated, but the frontier of autonomous spatial action is exemplified by the Smart Tissue Autonomous Robot, which performs supervised-autonomous soft-tissue suturing using 3D tissue tracking~\cite{shademan2016star} and has been extended to in-vivo intestinal anastomosis~\cite{saeidi2022star}.

\begin{table}[!th]
\centering
\fontsize{8}{10}\selectfont
\setlength{\tabcolsep}{1.0mm}
\caption{Representative benchmarks for spatial intelligence in medicine, spanning the three pivotal tasks.}
\label{tab:p4_s12_1_benchmarks}
\begin{tabular}{L{0.16\textwidth} L{0.28\textwidth} L{0.28\textwidth} L{0.20\textwidth}}
\toprule
\textbf{Pivotal task} & \textbf{Benchmark / dataset} & \textbf{Representative method} & \textbf{Headline metric} \\
\midrule
Segmentation & MSD~\cite{antonelli2022msd} & nnU-Net~\cite{isensee2021nnunet} & Dice / HD95 \\
Segmentation & BraTS~\cite{menze2015brats} & Swin-UNETR~\cite{hatamizadeh2022swinunetr} & Dice \\
Multi-organ & AMOS~\cite{ji2022amos} / BTCV~\cite{landman2015btcv} & UNETR~\cite{hatamizadeh2022unetr} & Dice \\
Detection & LUNA16~\cite{setio2017luna16} & 3D CNN detectors (challenge entries)~\cite{setio2017luna16} & FROC sensitivity \\
Registration & multi-site brain MRI (OASIS, ADNI, ABIDE) & VoxelMorph~\cite{balakrishnan2019voxelmorph} & TRE / Dice \\
Reconstruction & MM-WHS~\cite{zhuang2019mmwhs} & whole-heart mesh~\cite{kong2021meshdeformnet} & surface distance \\
Modeling & patient-specific cardiac & virtual heart~\cite{prakosa2018personalized} & ablation-target agreement \\
Endoscopy & SCARED~\cite{allan2021scared} & EndoNeRF~\cite{wang2022endonerf} & depth error \\
Workflow & Cholec80~\cite{twinanda2017endonet} & EndoNet~\cite{twinanda2017endonet} & phase accuracy \\
\bottomrule
\end{tabular}
\end{table}

\subsection{Established Tools and Systems}
\label{subsec:p4-sec12-1-established-tools-and-systems}

\textbf{Open-source frameworks.} The research stack runs from voxels to a planned model: \textbf{3D Slicer}\footnote{\url{https://www.slicer.org/}} and \textbf{MITK}\footnote{\url{https://www.mitk.org/}} provide segmentation, registration, and planning environments over the lower-level \textbf{ITK} and \textbf{SimpleITK}\footnote{\url{https://itk.org/}} engines; \textbf{MONAI}\footnote{\url{https://monai.io/}} supplies 3D-native deep-learning components and \textbf{nnU-Net}\footnote{\url{https://github.com/MIC-DKFZ/nnUNet}} the standard segmentation pipeline; and the web viewer \textbf{OHIF}\footnote{\url{https://ohif.org/}} renders and annotates volumes in the browser.

\textbf{Navigation and surgical-robotics vendors.} The clinical stack carries the model from preoperative registration to constrained in-OR action. Navigation systems from \textbf{Brainlab}\footnote{\url{https://www.brainlab.com/}}, \textbf{Medtronic} (StealthStation)\footnote{\url{https://www.medtronic.com/}}, and \textbf{Stryker}\footnote{\url{https://www.stryker.com/}} register and track instruments, and \textbf{Augmedics}\footnote{\url{https://augmedics.com/}} adds an augmented-reality head-mounted overlay. Surgical robots span teleoperated soft-tissue systems (\textbf{Intuitive} da Vinci and Ion\footnote{\url{https://www.intuitive.com/}}, \textbf{Medtronic} Hugo, \textbf{CMR Surgical} Versius\footnote{\url{https://cmrsurgical.com/}}) and CT-planned, haptically constrained orthopedic systems (\textbf{Stryker} Mako, \textbf{Zimmer Biomet} ROSA\footnote{\url{https://www.zimmerbiomet.com/}}).

\textbf{Imaging-AI and clinical platforms.} A deployed layer turns spatial findings into care. \textbf{Aidoc}\footnote{\url{https://www.aidoc.com/}} triages critical findings on CT, \textbf{Viz.ai}\footnote{\url{https://www.viz.ai/}} alerts stroke teams to large-vessel occlusions, \textbf{HeartFlow}\footnote{\url{https://www.heartflow.com/}} computes non-invasive fractional flow reserve from coronary CT, \textbf{PathAI}\footnote{\url{https://www.pathai.com/}} quantifies biomarkers on whole-slide pathology, and \textbf{Tempus}\footnote{\url{https://www.tempus.com/}} fuses imaging with genomic and clinical data.

\subsection{Challenges}
\label{subsec:p4-sec12-1-challenges}

Every capability in the clinical loop must be accurate, robust to deformation, and trustworthy enough to act on a patient. The open problems follow from those demands: generalizing across scanners and sites, obtaining labels under privacy constraints, registering deforming tissue during surgery, earning regulatory and clinical trust, and meeting real-time budgets while fusing imaging with non-imaging evidence.

\paragraph{Generalization across Scanners and Sites.} Models trained at one institution degrade at another, because scanner vendor, reconstruction kernel, and acquisition protocol induce a covariate shift that curated benchmarks underrepresent. 
Domain generalization and adaptation are therefore prerequisites for deployment, and leaderboard results routinely overstate performance at a new site.

\paragraph{Data Scarcity, Annotation Cost, and Privacy.} Volumetric expert annotation is slow and expensive, and patient data is locked in privacy-restricted silos. 
Self-supervised pretraining reduces the labeling burden~\cite{huang2023ssl}, and federated learning trains shared models without moving data across institutions~\cite{rieke2020federated}, but label efficiency and privacy-preserving collaboration remain limiting.

\paragraph{Non-Rigid Intraoperative Registration.} The hardest open spatial problem is keeping the model aligned with the patient as tissue moves. Brain shift after craniotomy and soft-tissue deformation during surgery invalidate rigid preoperative registration, and while learning-based deformable registration~\cite{balakrishnan2019voxelmorph} addresses the static case, real-time non-rigid intraoperative alignment under sparse, partial observation is not solved.

\paragraph{Validation, Regulation, and Trust.} A model that changes a clinical decision needs prospective validation and regulatory clearance before use, and at the bedside it needs calibrated uncertainty and explanations that support appropriate rather than blind reliance. Retrospective benchmark accuracy establishes none of these.

\paragraph{Real-Time Constraints and Multimodal Fusion.} Intraoperative guidance must run within a tight latency budget while the scene deforms, a constraint that few surgical-video algorithms yet meet without specialized hardware. 
At the same time, the richest clinical picture comes from fusing imaging with electronic health records and genomics~\cite{acosta2022multimodal}, so spatial perception must ultimately integrate with non-spatial evidence.


\newpage
\section{Natural Science}
\label{sec:Natural Science}
\label{sec:p4_s12_2}


\subsection{Introduction}
\label{subsec:p4-sec12-2-introduction}

A physical system is a field over space; a molecule or a material is a three-dimensional arrangement of atoms together with an electronic field over that arrangement; a protein, a cell, or a tissue is a three-dimensional structure whose function follows from its form.
In every case the scientific quantity of interest, whether an energy, a reaction rate, a biological function, or a material property, is a function of spatial structure under strict symmetry constraints, so that understanding nature is to a large degree a problem of recovering and reasoning about that structure.

The geometric representations of Chapter~\ref{chapter:single_modality} and the reconstruction and generation of Chapter~\ref{sec:generation} reappear here specialized to data that obey physical law: to fields constrained by differential equations, to atomic systems constrained by Euclidean and permutation symmetry, and to biomolecules whose geometry determines their function.
A single capability triad recurs inside every domain and gives the section its spine: spatial \textbf{representation and reconstruction} recovers a continuous field or a 3D structure from sparse, indirect, and noisy measurements; \textbf{dynamics, simulation, and law discovery} learns how that structure evolves and what governs it; and \textbf{design, discovery, and control} either generates new structure or inverts observations for hidden causes. The third leg reads as discovery and control in physics and as generative design in chemistry and life science.

The section is organized by scientific domain (physics, chemistry, and life science); the earth and environmental sciences are treated at geographic scale in the GeoAI of §\ref{sec:p4_s11_2}, the clinical, patient-scale counterpart of the life-science material in §\ref{sec:p4_s12_1}, and the generic digital twin in §\ref{sec:p4_s11_1}.
Figure~\ref{fig:p4_s12_2_overview} previews the section as a cross-scale map in which the three domains share the same vertical triad, and Table~\ref{tab:p4_s12_2_benchmarks} collects a representative benchmark for each capability.

\begin{figure}[!ht]
    \centering
    \includegraphics[width=0.98\linewidth]{figures/part4/section12/natural_science_overview_v2.pdf}
    \caption{Spatial intelligence across physics, chemistry, and life science (columns), each repeating the same capability triad (rows). The lower electron-to-tissue strip indicates illustrative scales rather than exclusive domain boundaries. Fields, molecular and periodic structures, and biomolecular and tissue structures illustrate representation; the equations exemplify wave, quantum, and biological dynamics. Bottom-row design examples are schematic; molecule/material and biomolecular/tissue examples are alternatives rather than literal conversion pipelines. Biological imagery is synthetic.}
    \label{fig:p4_s12_2_overview}
\end{figure}


\subsection{Physics}
\label{subsec:p4-sec12-2-physics}

Spatial intelligence in physics constructs spatial structure and discovers the laws that govern its motion, and it earns its place where first-principles solvers are too costly or the governing equations are unknown: fields are reconstructed from limited observations, their dynamics are modeled and their laws discovered, and the systems are optimized and controlled.
Figure~\ref{fig:p4_s12_2_physics} previews this progression from a sparsely observed field to a learned dynamical rollout and a differentiable-control loop.

\begin{figure}
    \centering
    \includegraphics[width=0.98\linewidth]{figures/part4/section12/natural_science_physics_v2.pdf}
    \caption{The physics triad: (1) reconstruct a field from sparse sensors, (2) fit dynamics and roll out the field, with a separate law-extraction step, and (3) optimize controls through a differentiable simulator. $\mathcal{N}[\mathbf{u};\theta]$ is the learned operator and $\mathcal{N}^*$ a candidate extracted law. The state $\mathbf u$ and control $\mathbf a$ are distinct; the control objective $\mathcal J$ supplies gradients to the optimizer, whereas $\mathcal L_{\rm fit}$ fits observed dynamics. Field contours and temporal shifts are illustrative, not experimental results.}
    \label{fig:p4_s12_2_physics}
\end{figure}

\subsubsection{Spatial Representation and Reconstruction}

Many key quantities in physical systems are available only through sparse and noisy observations, and representation and reconstruction recover continuous spatial fields from such data, the input to law discovery and control.
Embedding physical priors such as governing partial differential equations and symmetry constraints keeps these reconstructions physically consistent in low-data regimes.

A classical task is physical field reconstruction, which aims to recover continuous physical fields such as temperature, pressure, and velocity from sparse sensor readings or boundary measurements; the central challenge lies in ensuring the physical plausibility of the solution under the ill-posed conditions caused by insufficient observations.
\textbf{Physics-informed neural networks} are a classical approach for this task~\cite{raissi2019physics}, embedding partial differential equations such as the Navier--Stokes equations as soft constraints in the loss function and computing physical residuals via automatic differentiation, which enables accurate reconstruction of fluid and thermal fields in data-scarce settings; the same networks serve as real-time surrogates in digital twins (§\ref{sec:p4_s11_1}).
\textbf{Neural fields} implicitly represent continuous physical fields by using coordinate-based networks that map spatial locations to physical quantities.
Sitzmann \etal~\cite{sitzmann2020implicit} proposed SIREN, which employs periodic activation functions to capture high-frequency spatial details with high fidelity and has been applied to solving Poisson equations and reconstructing wave fields.
The general machinery of coordinate-based neural representation is developed in Chapter~\ref{sec:generation}.

Inverse scattering aims to infer the internal structure or material parameters of an object from measured wave-scattering data, and it remains a canonical ill-posed problem in electromagnetics and acoustics~\cite{chen2020review}.
Wei and Chen proposed deep-learning schemes~\cite{wei2019deep} for full-wave nonlinear inverse scattering that directly invert dielectric-permittivity profiles from scattered fields, thereby moving beyond linearized Born-type approximations in practical reconstructions.
SwitchNet~\cite{khoo2019switchnet} introduced a unified neural architecture that learns both the forward map from scatterers to scattered fields and the inverse map from measurements to scatterer profiles, using a switching layer that exploits low-rank structure to reduce model complexity and facilitate training.
Physics-guided losses that inject multiple-scattering near-field priors into the objective further improve robustness.

\subsubsection{Dynamics, Simulation, and Law Discovery}

\textbf{Law discovery} aims to learn the generative structures or explicit laws that govern the evolution of physical systems from observational data, emphasizing the automatic extraction of conservation laws, symmetries, and interpretable mechanisms beyond accurate prediction.
Hamiltonian Neural Networks~\cite{greydanus2019hamiltonian} directly learn the Hamiltonian of a system, thereby enforcing energy conservation and symplectic structure at the model level and providing a principled foundation for interpretable dynamical modeling.
Lagrangian Neural Networks~\cite{cranmer2020lagrangian} parameterize arbitrary Lagrangians with neural networks, learning energy conserving dynamics without canonical coordinates, as validated on the double pendulum and a relativistic particle.
Cranmer \etal~\cite{cranmer2020discovering} combine graph networks with symbolic regression on their internal representations, recovering known force laws and Hamiltonians, and applying the approach to cosmology.
The SINDy algorithm~\cite{brunton2016discovering} employs sparse regression to identify governing differential equations directly from time-series data, enabling data-driven discovery of interpretable physical laws.

\textbf{Efficient physical solvers} focus on predicting the spatio-temporal evolution of physical systems under known governing equations or interaction rules, with the goal of achieving fast and scalable computation while preserving physical consistency beyond the limits of traditional numerical methods.
The Fourier Neural Operator~\cite{li2021fourier} learns mappings between function spaces rather than individual discretized instances, enabling trained models to predict the evolution of partial differential equations at arbitrary spatial resolutions without relying on a fixed computational grid.
DeepONet~\cite{lu2021learning}, grounded in the universal approximation theorem for operators, learns nonlinear mappings from input functions to output functions and thereby approximates the solution operators that govern physical systems.
Interaction Networks~\cite{battaglia2016interaction} introduced the relational reasoning architecture that learns physical interactions between objects, laying the foundation for graph-based physical simulation.
The Graph Network Simulator~\cite{sanchez2020learning} employs graph neural networks to model spatial interactions among particles or objects, achieving high-fidelity simulations of fluids, rigid bodies, and deformable materials, and MeshGraphNets~\cite{pfaff2020learning} extend graph-based models to unstructured meshes, allowing efficient simulation of continuum mechanics on irregular geometries.
For the classical challenge of turbulence modeling, Mohan \etal~\cite{mohan2020spatio} proposed spatio-temporal deep-learning models with physics-informed diagnostics that improve the prediction accuracy of three-dimensional turbulent flows, and Lusch \etal~\cite{lusch2018deep} developed a deep-learning framework for universal linear embeddings of nonlinear dynamics, enabling low-dimensional representations and long-term prediction of nonlinear fluid systems.
Neural representations of the many-body wavefunction are treated with electronic structure in the chemistry subsection.

\subsubsection{Control and Optimization}

Optimization and control in physical systems seek optimal trajectories, shapes, or control policies under governing physical laws and environmental constraints.
Differentiable physics formulates physical simulators as differentiable operators, enabling gradient-based optimization over control inputs, design variables, and physical parameters, and providing a unified framework for optimal control, system identification, and design.
Neural Ordinary Differential Equations~\cite{chen2018neural} provide a differentiable formulation of continuous-time dynamical systems by parameterizing the vector field with a neural network and enabling gradient computation through adjoint sensitivity analysis, serving as a flexible building block for learning and control in continuous time.
ChainQueen~\cite{hu2019chainqueen} demonstrates differentiable simulation of soft and deformable bodies based on the moving least squares material point method, enabling gradient-based optimization in high-degree-of-freedom systems with contact and supporting applications in soft-robot control and co-design.
DiffTaichi~\cite{hu2020difftaichi} generalizes this paradigm through a differentiable programming framework for high-performance physical simulation, allowing efficient construction of differentiable solvers and gradient-driven optimization for fluid control, shape optimization, and material-parameter estimation; the same tools calibrate digital twins (§\ref{sec:p4_s11_1}).
In practice, a mature software substrate of differentiable and physics-informed libraries, such as \textbf{DeepXDE}\footnote{\url{https://github.com/lululxvi/deepxde}} and \textbf{NVIDIA Modulus}\footnote{\url{https://developer.nvidia.com/modulus}} for operator and PINN training, and \textbf{JAX-MD}\footnote{\url{https://github.com/jax-md/jax-md}} and \textbf{DeePMD-kit}\footnote{\url{https://github.com/deepmodeling/deepmd-kit}} for differentiable and learned-potential molecular dynamics, turns this triad into runnable physics.




\subsection{Chemistry}
\label{subsec:p4-sec12-2-chemistry}

Chemistry is intrinsically spatial: a molecule or a material is a three-dimensional arrangement of atoms together with an electronic field over it, and its behavior, whether energy, reactivity, or a bulk property, is a function of that geometry under strict symmetry.
Physical observables must be invariant or equivariant to the Euclidean group of rotations, translations, and reflections and to permutations of identical atoms; for crystalline solids they must also respect periodic lattice symmetry, and for electrons they must obey antisymmetry under exchange.
Spatial intelligence in chemistry is therefore largely the art of building these symmetries into models that represent atomic and electronic structure, simulate its dynamics, and generate new structure under chemical feasibility.
Figure~\ref{fig:p4_s12_2_chemistry} previews this triad from an equivariant representation of a 3D structure to a learned-potential simulation and a 3D generative design.

\begin{figure}
    \centering
    \includegraphics[width=0.98\linewidth]{figures/part4/section12/natural_science_chemistry_v2.pdf}
    \caption{The chemistry triad: (1) symmetry-aware molecular and periodic representations, illustrated by E(3)-equivariant features; (2) a learned invariant interatomic energy $U$ whose coordinate gradients provide forces for molecular dynamics, with related representations supporting property prediction; and (3) complementary 3D generative design by reverse diffusion. Molecular graphs, force directions, and the one-dimensional energy profile are schematic, not a calculated trajectory or energy surface. The denoising example preserves the atom count; generative methods may model coordinates and atom types jointly.}
    \label{fig:p4_s12_2_chemistry}
\end{figure}

\subsubsection{Spatial Representation and Reconstruction}

The central representational problem of chemistry is to map a 3D atomic geometry to energies, forces, and properties while respecting symmetry, and successive representations have captured progressively more of that geometry.
\textbf{Geometry-grounded potentials.} The Behler--Parrinello neural network potential founded the area by decomposing the total energy into atomic contributions computed from rotation-, translation-, and permutation-invariant symmetry functions of each atom's local environment~\cite{behler2007generalized}.
SchNet~\cite{schutt2017schnet} replaced hand-designed descriptors with continuous-filter convolutions over interatomic distances, giving a smooth invariant representation for energy and force prediction, while DimeNet~\cite{klicpera2020directional} added directional message passing that encodes bond angles through spherical harmonics, moving beyond pairwise distances to angular structure, and PaiNN~\cite{schutt2021equivariant} extended message passing to rotationally equivariant vector features for tensorial properties and spectra.
The Deep Potential model~\cite{zhang2018deep} preserves all natural symmetries at linear cost, and the ANI potentials~\cite{smith2017ani} demonstrated transferable organic-chemistry accuracy at force-field cost.
\textbf{Equivariant interatomic potentials.} The decisive recent step was to carry geometric tensors, not merely invariant scalars, through the network.
NequIP~\cite{batzner2022e3} introduced E(3)-equivariant convolutions over geometric tensors and achieved striking data efficiency, often requiring orders of magnitude less training data; Allegro~\cite{musaelian2023learning} made the representation strictly local and thereby scalable to 100 million atoms, and MACE~\cite{batatia2022mace} used higher body-order equivariant messages to reach high accuracy in only two message-passing layers.
A further shift is from per-system fits to \textbf{universal potentials} that span the periodic table: M3GNet~\cite{chen2022universal} couples a graph network with explicit three-body interactions trained on materials relaxations, and CHGNet~\cite{deng2023chgnet} augments the 3D-structure graph with atomic magnetic moments to capture electronic degrees of freedom alongside geometry. These equivariant backbones share machinery with the geometric representations of Chapter~\ref{chapter:single_modality}.
\textbf{Neural electronic structure.} At the quantum scale the spatial object is the many-electron wavefunction itself.
Neural quantum states parameterize that wavefunction with a neural network over spin configurations~\cite{carleo2017solving}, and deep variational ans\"atze such as FermiNet~\cite{pfau2020ferminet} and PauliNet~\cite{hermann2020deep} operate directly on 3D electron coordinates while enforcing exact antisymmetry, solving the electronic Schr\"odinger equation to chemical accuracy from nuclear positions alone.

\subsubsection{Dynamics, Simulation, and Property Prediction}

The potentials of the previous thread are also the simulators: once a model maps a 3D configuration to energies and forces, it drives molecular dynamics at a fraction of the cost of first-principles calculation, and the same representations predict quantum properties such as energies, dipoles, and frontier-orbital gaps directly from coordinates.
The gradient-domain machine-learning force field made this explicit by exploiting energy conservation and the geometry of conformations, and it established the widely used MD17 trajectory benchmark of energies and forces for small molecules~\cite{chmiela2017machine}.
Heterogeneous catalysis pushes the same energy-and-force prediction task to a far larger and more diverse chemical space, spanning many combinations of metal surfaces and adsorbates. 
GemNet-OC~\cite{gasteiger2022gemnet} extends directional message passing to this setting, delivering accurate energies and forces across the surfaces and adsorbates of the Open Catalyst datasets.
In parallel, structure-to-property prediction is standardized on QM9~\cite{ramakrishnan2014quantum}, a dataset of roughly 134{,}000 small organic molecules with density-functional geometries and quantum properties, on which the equivariant networks above are routinely evaluated by the mean absolute error of predicted properties.
This ecosystem rests on open toolkits, including the equivariant-network library \textbf{e3nn}\footnote{\url{https://e3nn.org/}}, the atomistic-simulation environment \textbf{ASE}\footnote{\url{https://wiki.fysik.dtu.dk/ase/}}, and the quantum-chemistry package \textbf{PySCF}\footnote{\url{https://pyscf.org/}}, together with universal-potential foundation models such as MACE-MP, and on reference corpora such as the Materials Project~\cite{jain2013commentary}, the Open Catalyst dataset~\cite{chanussot2021open}, and the MoleculeNet benchmark suite~\cite{wu2018moleculenet}.

\subsubsection{Generation and Inverse Design}

The inverse of property prediction is design: generating new 3D structures that satisfy desired properties under chemical feasibility, by sampling on the appropriate geometric manifold rather than enumerating candidates.
G-SchNet~\cite{gebauer2019symmetry} autoregressively places atoms in 3D Euclidean space with each placement respecting the rotational invariance and local symmetry of those already placed, the first strong 3D-native generator.
Diffusion models then became the dominant paradigm: the equivariant diffusion model~\cite{hoogeboom2022equivariant} jointly denoises atom coordinates and types under E(3) equivariance, GeoDiff~\cite{xu2022geodiff} generates conformers from the molecular graph through roto-translation-equivariant Markov kernels that induce an invariant distribution, and torsional diffusion~\cite{jing2022torsional} restricts the process to the hypertorus of torsion angles, the true intrinsic conformational degrees of freedom, achieving state-of-the-art conformer quality on GEOM-Drugs~\cite{axelrod2022geom} at a fraction of the cost.
For crystalline materials, CDVAE~\cite{xie2022crystal} generates stable structures under periodicity and SE(3) invariance and supports inverse design by optimizing target properties in latent space.
The line has since moved to diffusion directly on the periodic structure: DiffCSP~\cite{jiao2023crystal} jointly denoises lattice and fractional coordinates under periodic equivariance for crystal structure prediction, and MatterGen~\cite{zeni2025generative} scales this into a property-conditioned foundation model for inorganic materials design, with generated candidates validated by experimental synthesis.
The general diffusion machinery underlying these models is developed in Chapter~\ref{sec:generation}; what is chemistry-specific is the constraint manifold (valence and bonding feasibility, energy-minimum geometries, lattice periodicity) on which generation must take place.


\subsection{Life Science}
\label{subsec:p4-sec12-2-life-science}

At the molecular, cellular, and tissue scale of basic biology the object of interest is again a three-dimensional structure: a protein fold, a molecular complex, a spatial map of gene expression, or the wiring of a neural circuit, and in each case function follows from spatial form.
Spatial intelligence here recovers that structure from indirect measurement, models how biomolecules interact and move, and designs new structure.
Figure~\ref{fig:p4_s12_2_lifescience} previews the triad from structure recovery through interaction modeling to generative design.

\begin{figure}
    \centering
    \includegraphics[width=0.98\linewidth]{figures/part4/section12/natural_science_lifescience_v2.pdf}
    \caption{The life-science triad: (1) complementary sequence-to-structure prediction and cryo-EM projection-to-density reconstruction, (2) ligand binding-pose geometry with translation, rotation, and bond torsion, and (3) de novo backbone generation followed by sequence design. These are complementary capabilities, not an obligatory serial pipeline. Protein, pocket, and density renderings, noisy projection examples, and sequence fragments are schematic; they do not constitute experimentally validated structures or designs.}
    \label{fig:p4_s12_2_lifescience}
\end{figure}

\subsubsection{Spatial Representation and Reconstruction}

The defining achievement of spatial intelligence in biology is the recovery of 3D structure from indirect signals.
\textbf{Protein structure prediction} maps a one-dimensional sequence to a three-dimensional fold.
The original AlphaFold predicted inter-residue distance distributions as a learned potential to fold the chain~\cite{senior2020improved}, and AlphaFold2~\cite{jumper2021highly} reached near-experimental accuracy through an end-to-end architecture that reasons explicitly over residue frames and pairwise geometry, the exemplar of spatial intelligence in the natural sciences.
AlphaFold3~\cite{abramson2024accurate} generalized this to the joint structure of proteins, nucleic acids, and small molecules through a diffusion decoder; RoseTTAFold~\cite{baek2021accurate} let one-dimensional, two-dimensional, and three-dimensional information flow between tracks, and ESMFold~\cite{lin2023evolutionary} showed that atomic-level structure emerges from a protein language model without explicit alignments, enabling metagenomic-scale folding.
\textbf{Cryo-electron microscopy} solves the inverse problem of reconstructing a 3D density map from many noisy 2D projections.
RELION cast this as Bayesian refinement~\cite{scheres2012relion}, CryoSPARC enabled rapid ab-initio reconstruction without a starting model~\cite{punjani2017cryosparc}, and cryoDRGN used a coordinate-based neural decoder to reconstruct continuous distributions of conformations and thereby resolve molecular heterogeneity and motion~\cite{zhong2021cryodrgn}.
\textbf{Spatial transcriptomics} places gene expression back into tissue space.
The founding barcoded-array assay preserved the positional information of the transcriptome in tissue sections~\cite{stahl2016visualization}; Tangram learns a spatial mapping that aligns dissociated single cells onto a spatial reference~\cite{biancalani2021deep}, and novoSpaRc reconstructs spatial arrangement from expression alone through optimal transport, assuming that expression-similar cells lie physically close~\cite{nitzan2019gene}.
\textbf{Connectomics} reconstructs neural wiring from electron-microscopy volumes, where flood-filling networks iteratively trace the 3D morphology of individual neurites through the volume~\cite{januszewski2018high}.

\subsubsection{Dynamics, Simulation, and Interaction}

Beyond static structure, spatial intelligence models how biomolecules bind and interact.
Molecular docking predicts the 3D pose in which a ligand binds a protein, and recent methods treat it as a geometric problem on the natural manifold of poses.
DiffDock~\cite{corso2023diffdock} formulates docking as diffusion over the product space of ligand translation, rotation, and torsion, sampling binding poses in their intrinsic degrees of freedom, while EquiBind~\cite{stark2022equibind} regresses the binding location and pose in a single shot using an SE(3)-equivariant network whose equivariance encodes the rotational and translational symmetry of the binding geometry.
Protein--protein and complex geometry is predicted by extending folding models to multiple chains~\cite{evans2022protein}, and AlphaFold3 above provides a generative alternative for protein--ligand and protein--nucleic-acid complexes.
The dynamics of biomolecules in turn rest on the learned interatomic potentials of the chemistry subsection, which accelerate the molecular dynamics that sample conformational change.

\subsubsection{Generation and Design}

The generative leg designs new 3D biological structure.
RFdiffusion~\cite{watson2023novo} adapts structure-denoising diffusion to generate novel protein backbones conditioned on functional and geometric constraints such as motif scaffolding, binders, and symmetric assemblies, and ProteinMPNN~\cite{dauparas2022robust} solves the inverse-folding problem of designing an amino-acid sequence that folds into a given 3D backbone, the standard partner to backbone generation.
Chroma~\cite{ingraham2023illuminating} couples equivariant graph networks with a diffusion process to sample novel all-atom structures and sequences steerable toward desired properties, illuminating regions of protein space never explored by evolution.
As in chemistry, the underlying diffusion machinery is that of Chapter~\ref{sec:generation}; what is biology-specific is that the generated geometry must be foldable and functional, which is why design is paired with structure prediction and, ultimately, wet-lab validation.
The field is supported by open infrastructure, including the AlphaFold Protein Structure Database and AlphaFold3 server\footnote{\url{https://alphafoldserver.com/}}, the ESM Metagenomic Atlas\footnote{\url{https://esmatlas.com/}}, and the open-source folding stack OpenFold\footnote{\url{https://github.com/aqlaboratory/openfold}}, together with reference corpora such as the Protein Data Bank, the CASP blind-assessment benchmark~\cite{kryshtafovych2021critical}, and the PDBBind docking benchmark~\cite{wang2004pdbbind}.

\begin{table}[!th]
\centering
\fontsize{8}{10}\selectfont
\setlength{\tabcolsep}{1.0mm}
\caption{Representative benchmarks for spatial intelligence across the three natural-science domains.}
\label{tab:p4_s12_2_benchmarks}
\begin{tabular}{L{0.11\textwidth} L{0.17\textwidth} L{0.24\textwidth} L{0.19\textwidth} L{0.21\textwidth}}
\toprule
\textbf{Domain} & \textbf{Capability} & \textbf{Benchmark / dataset} & \textbf{Representative method} & \textbf{Headline metric} \\
\midrule
Physics & PDE surrogate & Navier--Stokes / Darcy flow & FNO~\cite{li2021fourier} & relative $L_2$ error \\
Physics & Particle dynamics & Learning-to-Simulate suite (Water / Sand / Goop) & GNS~\cite{sanchez2020learning} & rollout MSE \\
Chemistry & Property prediction & QM9~\cite{ramakrishnan2014quantum} & DimeNet~\cite{klicpera2020directional} & property MAE \\
Chemistry & Force field & MD17~\cite{chmiela2017machine} & NequIP~\cite{batzner2022e3} & energy / force MAE \\
Chemistry & Catalysis & Open Catalyst~\cite{chanussot2021open} & GemNet-OC~\cite{gasteiger2022gemnet} & energy / force MAE \\
Chemistry & Conformer generation & GEOM-Drugs~\cite{axelrod2022geom} & torsional diffusion~\cite{jing2022torsional} & coverage / matching RMSD \\
Life science & Protein folding & CASP~\cite{kryshtafovych2021critical} & AlphaFold2~\cite{jumper2021highly} & GDT-TS \\
Life science & Docking & PDBBind~\cite{wang2004pdbbind} & DiffDock~\cite{corso2023diffdock} & top-1 RMSD $<2$\,\AA{} (\%) \\
Life science & Cryo-EM heterogeneity & EMPIAR single-particle datasets & cryoDRGN~\cite{zhong2021cryodrgn} & resolution (FSC) \\
Life science & Protein design & Motif-scaffolding / binder design & RFdiffusion~\cite{watson2023novo} & in-silico design success rate \\
\bottomrule
\end{tabular}
\end{table}


\subsection{Challenges}
\label{subsec:p4-sec12-2-challenges}

A learned model of nature has to respect the symmetries and conservation laws of its domain and remain correct outside the data it was fitted to. The six problems below follow from that requirement: how to build symmetry in, how to bridge scales, how to solve ill-posed inverse problems, how to learn from scarce and expensive labels, how to generalize to genuinely new systems, and how to make predictions physically consistent and trustworthy.

\paragraph{Symmetry and Equivariance.} Physical observables are constrained by symmetry: Euclidean invariance and equivariance, permutation of identical particles, antisymmetry of electrons, and periodicity of crystals. Whether to build these symmetries into the architecture or to learn them from data is the central design tension, trading the cheaper but less expressive invariant features~\cite{schutt2017schnet} against the richer but costlier equivariant tensor features~\cite{batzner2022e3}, and it recurs from interatomic potentials to protein structure.

\paragraph{Extreme Multi-Scale Structure.} Spatial structure in nature spans many orders of magnitude, from electrons and atoms to molecules and materials, and from biomolecules through cells to tissue. Learned potentials now reach a hundred million atoms~\cite{musaelian2023learning}, still far below cellular scale; no single representation natively bridges electronic-structure accuracy and mesoscale reachability, and the handoffs between scales, through coarse-graining or learned-potential dynamics, are lossy and remain only partially solved.

\paragraph{Ill-Posed Inverse Problems.} Much of natural science is an inverse problem in which a high-dimensional spatial field must be recovered from sparse, indirect, and noisy measurements: field reconstruction from a few sensors, inverse scattering~\cite{wei2019deep}, and the recovery of 3D density from 2D cryo-electron-microscopy projections~\cite{zhong2021cryodrgn} all share this under-determined character and demand strong learned priors to yield a unique and physical solution.

\paragraph{Data Scarcity and the Cost of Ground Truth.} High-fidelity labels are expensive: accurate energies and forces require costly quantum-chemistry calculations, which is why force-field benchmarks such as MD17 cover only small molecules~\cite{chmiela2017machine}, and validated 3D biological structures require slow experiments, so the space of trustworthy training data is small relative to the space of possible systems. This drives the field's emphasis on data efficiency, multi-fidelity learning, and physics-guided objectives.

\paragraph{Generalization Beyond the Training Distribution.} Models must extrapolate to unseen chemistry, novel protein folds, off-equilibrium and reactive geometries, and configurations far from those they were trained on. Discovery is by definition out of distribution, so leaderboard accuracy says little about it; this is why generative materials design is now judged by whether its candidates can actually be synthesized~\cite{zeni2025generative}.

\paragraph{Physical Consistency and Trust.} Predictions that violate conservation laws or come without calibrated uncertainty cannot support a scientific conclusion. Building conservation into the model, as Hamiltonian networks do for energy~\cite{greydanus2019hamiltonian}, and quantifying confidence are what let a learned model extend the scientific record rather than interpolate it.

\newpage
\chapter{Spatial Intelligence for Intelligent Assistance}


\FramedInitial{I}{ntelligent} Assistance is the assistant-facing application of spatial intelligence: an agent must perceive a user's lived environment, remember its changing state, ground natural instructions in that state, and act in ways that are useful, legible, and safe.
Unlike a task-specific robot in a factory, vehicle, or clinic, a personal assistant operates in homes, offices, streets, and wearable or XR interfaces where the scene is private, cluttered, and shaped by individual routines.
This chapter therefore treats assistance as a persistent human-centered loop over spatial memory, language grounding, action, and feedback.
The same spatial model supports tidying and search, close collaboration, and accessibility support for vulnerable users, but each setting changes how much the assistant must model the person as well as the place.





\section{Introduction}
\label{subsec:p4-sec13-1-introduction}\label{sec:p4_s13_1}

An intelligent assistant is a generalist embodied agent that shares the user's own private, unstructured world and acts on the user's behalf. 
Unlike the domain agents of the preceding chapters, which are optimized for a bounded operational envelope (a road network, a warehouse aisle, an operating theater), the assistant must function inside the open-ended geometry of a home, an office, or a sidewalk, in the continuous physical presence of the person it serves. 
Its defining spatial challenge is therefore not a new perception, reasoning, generation, or memory primitive: those capabilities are surveyed in Chapter~\ref{chapter:3:sec:6} (perception), Chapter~\ref{sec:reasoning} (reasoning), Chapter~\ref{sec:generation} (generation), and Section~\ref{subsec:p3_s7_3_6} (retrieval- and memory-augmented spatial reasoning).
The challenge here is rather to build and maintain a persistent, personalized, and queryable spatial-semantic model of the user's environment, to ground human intent onto that model, and to act on it legibly and safely while a person occupies the same space.

Three properties separate the assistant from the embodied agents discussed elsewhere in this part. The environment is the user's own, so the assistant cannot assume an engineered layout and must instead generalize across the long tail of real homes. 
It is also personalized: the same instruction (``tidy up,'' ``take me to the exit'') resolves to different geometric goals for different users, so the spatial model must absorb individual preference. 
And because the human is always co-present, otherwise routine motion becomes a social act that must be predictable, comfortable, and safe. These three demands (an unengineered space, an individual user, and a person sharing that space) recur across every regime in which the assistant operates.

\section{System Overview}
\label{subsec:p4-sec13-1-system-overview}

Across these settings, every assistant closes the same loop. It \textbf{perceives} the environment and the user, \textbf{represents} what it perceives in a persistent spatial-semantic memory, \textbf{grounds} a natural-language or gestural instruction onto that memory, \textbf{reasons} about a goal-directed plan, \textbf{acts} on the physical or informational world, and observes the human's \textbf{feedback}, which re-grounds the next interaction. 
The persistent memory is the part that distinguishes an assistant from a single-shot perception or manipulation system: because the assistant lives in one place over long horizons, it can accumulate a map of where objects belong, where the user tends to be, and how the user prefers tasks to be performed. 
Long-horizon mobile-manipulation benchmarks such as EMMOE instantiate one slice of this loop (language-conditioned planning and action in open environments) but still leave persistent personalization, memory update, and human feedback as broader assistant-level requirements~\cite{li2025emmoe}. Figure~\ref{fig:p4_s13_1_overview} schematizes the loop, with the spatial-semantic memory at its center and the human present at every stage.

\begin{figure}[t]
    \centering
    \includegraphics[width=0.98\linewidth]{figures/part4/section13/assistant_overview.pdf}
    \caption{The intelligent-assistant loop. User and scene observations populate persistent room--object memory; instructions are grounded to specific objects and target states before planning. Alternative outputs illustrate household management, shared-space collaboration, and audio/haptic assistance. Memory links denote spatial associations; the dashed ellipse marks shared interaction space. User responses and observed changes close the feedback loop.}
    \label{fig:p4_s13_1_overview}
\end{figure}

The assistant is deployed in three escalating regimes of human-proximity, which organize the remainder of this chapter. 
In the first, the assistant acts \emph{on} the space while the human is a beneficiary at a distance: it tidies, fetches, and searches. 
In the second, the assistant acts \emph{alongside} the human as a partner that must move and manipulate in shared space without intruding. 
In the third, the assistant acts \emph{for} a vulnerable human, standing in for an impaired capability under hard safety, privacy, and accessibility constraints. 
The three regimes share the loop above but place progressively more weight on modeling the person rather than the place; Table~\ref{tab:p4_s13_1_benchmarks} collects benchmarks for all three.

\section{Pivotal Task-1: Managing the Physical Environment}
\label{subsec:p4-sec13-1-pivotal-task-1-managing-the-physical-environment}

The most autonomous regime asks the assistant to maintain the physical state of a human environment with little real-time supervision. A persistent model of the dwelling supports two complementary operations: grounding instructions for physical action, and retrieving remembered evidence to answer queries. New observations update the model, including after an action changes the environment. Figure~\ref{fig:p4_s13_1_managing} separates these action and recall branches around their shared home memory.

\textbf{Representing the home.} Everything downstream depends on a queryable model of the space, which extends the scene-understanding methods of Chapter~\ref{chapter:3:sec:6} into a persistent, semantically indexed memory. The memory architectures themselves are treated in Section~\ref{subsec:p3_s7_3_6}; here they are the infrastructure that lets an assistant know not only \emph{where} surfaces are, but \emph{what} and \emph{whose}. 
The 3D scene graph \cite{armeni20193dsg} introduced a hierarchical structure that unifies objects, rooms, and cameras with the underlying geometry, and Hydra \cite{hughes2022hydra} brought such graphs into real-time, loop-closure-optimized robot perception so that the memory stays current as the assistant moves. 
To make the memory addressable by free-form language, ConceptFusion \cite{jatavallabhula2023conceptfusion} builds an open-set, multimodal 3D map that answers text, image, or click queries, and ConceptGraphs \cite{gu2024conceptgraphs} fuses foundation-model outputs across views into a captioned, language-queryable 3D scene graph suitable for planning. 
This open-vocabulary memory is what lets an assistant act on objects it was never explicitly trained to recognize.


\begin{figure}[!ht]
    \centering
    \includegraphics[width=0.98\linewidth]{figures/part4/section13/assistant_managing.pdf}
    \caption{Managing the physical environment through a shared home memory. Observations maintain room, object, and relation records; grounded instructions drive manipulation, while queries independently retrieve last-seen evidence. Physical outcomes are re-observed to update memory.}
    \label{fig:p4_s13_1_managing}
\end{figure}

\textbf{Acting on the space: tidy-up and rearrangement.} Restoring order is the canonical management task, framed by Batra \etal \cite{batra2020rearrangement} as bringing an environment from an arbitrary configuration into a target spatial state. 
The benchmark substrate includes the Visual Room Rearrangement (RoomR) challenge \cite{weihs2021visual}, which asks an agent to observe a goal arrangement and then reset one to five displaced objects, and Housekeep \cite{kant2022housekeep}, which removes the explicit goal and instead judges placements against human preferences over nearly two thousand objects. 
Method work has moved from fixed rules to learned commonsense priors. 
TIDEE \cite{sarch2022tidee} tidies a never-before-seen room from pixels using a visuo-semantic out-of-place detector and an associative memory of where receptacles tend to be, and TidyBot \cite{wu2023tidybot} makes the placement \emph{preference-aware}, using few-shot summarization with a large language model to generalize a particular user's tidying rules from a handful of examples. 
The shift these systems mark, from a universal notion of tidiness to one conditioned on the resident, is exactly the personalization property that defines the assistant.

\textbf{Acting on the space: open-vocabulary mobile manipulation.} Tidying a real home requires picking and placing arbitrary objects at arbitrary locations, the task formalized as open-vocabulary mobile manipulation by HomeRobot \cite{yenamandra2023homerobot}, which pairs a simulation and real benchmark with a software stack for the Hello Robot Stretch. 
Notable deployment evidence comes from OK-Robot \cite{liu2024okrobot}, which combines a phone-scanned dense vision-language map with off-the-shelf navigation and grasping to achieve zero-shot pick-and-drop across ten real homes, and from Mobile ALOHA \cite{fu2024mobilealoha}, which shows that low-cost whole-body teleoperation plus imitation learning can acquire dexterous bimanual household skills. 
These systems lean directly on the manipulation primitives of Section~\ref{sec:p4_s9_1}, increasingly instantiated as generalist vision-language-action policies (\eg, OpenVLA \cite{kim2024openvla}, $\pi_0$ \cite{black2024pi0}) that fold perception, grounding, and control into one model; the assistant contribution is to drive them in unstructured homes from open-vocabulary goals rather than in an engineered cell. 
Commercial and pre-commercial examples illustrate the same pressure from different directions: mapping vacuum robots such as iRobot's Roomba, which build semantic room maps for targeted cleaning,\footnote{\url{https://www.irobot.com}} represent the mature narrow end, while emerging home-humanoid platforms such as 1X's Neo\footnote{\url{https://www.1x.tech/neo}} and Figure\footnote{\url{https://www.figure.ai}} show how the field is attempting, so far without a general-purpose result, to generalize mobile manipulation into domestic settings.

\textbf{Searching and recalling the model.} The dual of acting on the space is querying it: an assistant is, among other things, a search engine for the physical world that answers ``where did I leave my keys?'' Object-goal navigation supplies the active half, with goal-oriented semantic exploration (SemExp) \cite{chaplot2020object} building an episodic semantic map and a policy that exploits object co-occurrence to find a target efficiently rather than searching blindly. 
Embodied question answering \cite{das2018embodied} couples navigation with answering, and its open-vocabulary successor OpenEQA \cite{majumdar2024openeqa} spans both an episodic-memory setting, in which the assistant must recall what it has already seen, and an active-exploration setting, over more than a hundred real environments. 
For wearable assistants that record the user's day, the same recall problem becomes natural-language search of long egocentric video, the episodic-memory task defined by Ego4D \cite{grauman2022ego4d}. 
This is the application-facing counterpart of the memory-augmented reasoning systems in Section~\ref{subsec:p3_s7_3_6}.



\section{Pivotal Task-2: Co-habiting and Collaborating}
\label{subsec:p4-sec13-1-pivotal-task-2-co-habiting-and-collaborating}

As the distance between assistant and human shrinks, spatial intelligence must become social intelligence. The assistant must model the human as a co-located agent with intent, viewpoint, and comfort zones, not as a moving obstacle or a bare command source, and act predictably and collaboratively. Running through every part of this task is the reconciliation of reference frames: the assistant's own egocentric view, the human's egocentric view, and a shared allocentric map must be kept aligned.

\begin{figure}[!ht]
    \centering
    \includegraphics[width=0.98\linewidth]{figures/part4/section13/assistant_collaborating.pdf}
    \caption{Collaborating in shared space: (a) person- and context-dependent comfort zones and (b) aligned robot/human views inform an interaction-aware plan; (c) navigation and handover illustrate alternative behaviors. Dashed rings indicate comfort regions, and dashed sight lines associate views with shared objects.}
    \label{fig:p4_s13_1_collaborating}
\end{figure}

\textbf{Socially-aware navigation.} Moving through a populated space is not collision avoidance; it is adherence to the unwritten contracts of proxemics. The classical account is the social force model \cite{helbing1995social}, which casts pedestrian motion as attraction to a goal under repulsion from others, operationalizing Hall's anthropological zones of personal space. 
Data-driven successors learn the interaction directly: Social LSTM \cite{alahi2016social} couples per-person recurrent predictors through a social-pooling layer to forecast jointly plausible trajectories, and deep reinforcement-learning planners such as the attention-based crowd-robot model \cite{chen2019crowdrobot} weigh neighbors to navigate dense crowds while avoiding the ``freezing robot'' failure that purely reactive planners suffer \cite{trautman2010unfreezing}. 
Progress is measured on human-demonstration corpora such as SCAND \cite{karnan2022socially}, which records hours of socially compliant teleoperation for imitation and inverse reinforcement learning. 
The trend is toward anticipatory planning, in which the assistant infers the human's intended path and yields proactively, and toward learnable cost functions that adapt passing distances to individual or cultural preference rather than enforcing one universal margin.

\textbf{Grounded instruction following and spatial dialogue.} A collaborating assistant must resolve spatial language against what it sees, including the deictic references (``put it next to the \emph{other} bowl,'' ``the one over \emph{there}'') that pervade human speech. 
Vision-and-language navigation, introduced with the Room-to-Room benchmark \cite{anderson2018vision}, established the problem of following natural-language route instructions in photorealistic buildings; ALFRED \cite{shridhar2020alfred} extended it to long-horizon household tasks with irreversible state changes, and TEACh \cite{Padmakumar2022TEACh} cast the task as a two-agent dialogue in which a commander and follower must build shared spatial reference over multiple turns. 
Resolving a referring expression to a specific object in a cluttered 3D scene is studied directly by ScanRefer \cite{chen2020scanrefer} and ReferIt3D \cite{achlioptas2020referit3d}, whose same-class distractors force the model to rely on spatial relations rather than appearance. 
These benchmarks build on the reasoning capabilities of Chapter~\ref{sec:reasoning}; the assistant requirement is to maintain a spatial working memory that links current observations to earlier instructions so that deixis can be resolved.

\textbf{Shared autonomy and collaboration.} When full autonomy is brittle, the assistant acts as a third hand under shared control. 
The defining formulation is shared autonomy via hindsight optimization \cite{javdani2018shared}, in which the robot infers a distribution over the user's spatial goal and assists toward it without committing to a single confident prediction, formalizing the earlier policy-blending account of arbitration between human and robot intent \cite{dragan2013policy}. 
Equally important is making the assistant's own motion readable: legible motion \cite{dragan2013legibility} shapes a reaching trajectory so a human can quickly infer its intended target, which is what allows close-quarters collaboration to feel safe. Physical handovers, where giver and receiver must coordinate in space and time \cite{strabala2013handover}, are the elementary collaborative act on which assistive manipulation is built.

\textbf{Perspective-taking and egocentric grounding.} Many assistants observe the world from a wearable, first-person viewpoint, so the frame-alignment problem becomes explicit. 
Ego4D \cite{grauman2022ego4d} and its successors anchor egocentric perception at scale, and dedicated capture platforms such as Project Aria\footnote{\url{https://www.projectaria.com}} supply the calibrated multi-sensor first-person stream such assistants run on, while spatial-VQA backbones such as SpatialVLM \cite{chen2024spatialvlm} give a vision-language model the quantitative 3D reasoning (distances, sizes, directions) needed to answer spatial questions about the user's scene. 
The hardware and pose-estimation side of wearable sensing is treated in Section~\ref{sec:p4_s9_3}, and the augmented-reality runtime that renders guidance into the user's view is treated in Section~\ref{sec:p4_s10_1}; here the contribution is the reasoning that converts a first-person observation into a collaborative action in the human's frame. 
Embodied-AR applications make this alignment problem explicit when directions, warnings, or manipulation cues must be registered into the user's view \cite{yang2025embodiedar}. Figure~\ref{fig:p4_s13_1_collaborating} illustrates the elements of this regime.



\section{Pivotal Task-3: Empowering Vulnerable Users}
\label{subsec:p4-sec13-1-pivotal-task-3-empowering-vulnerable-users}

The highest-impact regime uses spatial intelligence to compensate for a sensory, cognitive, or motor impairment, acting as a prosthesis for the user's perception or agency. The assistant translates the environment into a queryable, navigable, and actionable representation and delivers it through a channel the user can still use (audio, haptics, or physical action). Because the loop closes around a vulnerable person in their own home, safety, privacy, and personalization are design constraints from the start.

\begin{figure}[!ht]
    \centering
    \includegraphics[width=0.98\linewidth]{figures/part4/section13/assistant_empowering.pdf}
    \caption{Personalized assistance combines observed spatial state with user goals and access needs. Consent, confidence, and clearance checks gate audio/haptic cues, stepwise guidance, and shared control; uncertainty or missing permission prompts clarification or a pause. Shared control also receives the user's command. Scenes are illustrative.}
    \label{fig:p4_s13_1_empowering}
\end{figure}

\textbf{Assistance for the blind and visually impaired.} For blind and low-vision users, spatial intelligence turns an unknown and hazardous environment into a queryable one. The oldest strand is sensory substitution, exemplified by the vOICe \cite{meijer1992voice}, which sonifies the spatial structure of an image into sound. 
The modern strand is conversational visual question answering: VizWiz \cite{bigham2010vizwiz} introduced a deployed system in which a blind user photographs their surroundings and asks a spoken question, and the VizWiz Grand Challenge \cite{gurari2018vizwiz} turned the resulting real-world distribution, including blurry and unanswerable images, into the reference benchmark for assistive VQA. 
Indoor wayfinding is addressed by smartphone navigation aids such as NavCog \cite{ahmetovic2016navcog}, which deliver turn-by-turn audio guidance from accurate beacon-based localization, and by three-dimensional spatial audio that places points of interest around the user. 
Commercial assistive applications include Microsoft's Seeing AI for scene and text narration\footnote{\url{https://www.seeingai.com}}, Be My Eyes for volunteer- and model-mediated visual assistance\footnote{\url{https://www.bemyeyes.com}}, and wearable systems from OrCam\footnote{\url{https://www.orcam.com}}, Envision\footnote{\url{https://www.letsenvision.com}}, and biped.ai\footnote{\url{https://biped.ai}} for scene description and obstacle warning. 
These tools move from passive obstacle warning to active interrogation, in which the user asks ``is there a free seat?'' or ``read the sign on the door'' and the assistant retrieves the answer.

\textbf{Cognitive and mobility assistance.} For users with cognitive decline or limited mobility, spatial assistants provide external scaffolding for the activities of daily living. 
Ambient assisted living systems such as the CASAS smart home \cite{cook2013casas} instrument a dwelling to track the resident's location and activity, learning the spatio-temporal patterns whose violation (a fall, a missed meal) signals a need to intervene; wearable accelerometry supplies the canonical fall-detection paradigm \cite{bourke2007evaluation}.

Where monitoring gives way to physical action, the assistant becomes a shared-control effector. 
Smart wheelchairs \cite{simpson2005smart} and wheelchair-mounted assistive arms use collaborative control to enforce spatial safety zones, for example forbidding a chair from driving down a staircase, so that the user retains mobility without constant supervision; the goal-arbitration that underlies such control is the same shared-autonomy machinery introduced above, here cast as a decision-theoretic problem \cite{ghorbel2018decision}. 
Assistive feeding \cite{bhattacharjee2019towards} shows how fine spatial reasoning (where to acquire a bite, how to transfer it to the user's mouth) becomes a caregiving act. 
The platform lineage runs from Robots for Humanity \cite{chen2013robots}, which let a man with quadriplegia perform self-care through a teleoperated PR2, to the affordable Stretch mobile manipulator \cite{kemp2022design}, with physics-based simulation of caregiving tasks consolidated in Assistive Gym \cite{erickson2020assistive}. 
Commercial assistive hardware in this space includes the Kinova JACO arm\footnote{\url{https://www.kinovarobotics.com}} and the Obi feeding device\footnote{\url{https://meetobi.com}}.
This regime overlaps the clinical setting of Section~\ref{sec:p4_s12_1}, but the boundary is clear: there the focus is diagnosis and surgical intervention, whereas here it is daily living in the home. Figure~\ref{fig:p4_s13_1_empowering} summarizes the compensation principle.



\begin{table}[t]
\centering
\fontsize{8}{10}\selectfont
\setlength{\tabcolsep}{1.0mm}
\caption{Representative benchmarks across the three assistant regimes. Where a benchmark has no standard method, the method column names the benchmark paper's own baselines.}
\label{tab:p4_s13_1_benchmarks}
\begin{tabular}{L{0.20\textwidth} L{0.24\textwidth} L{0.26\textwidth} L{0.22\textwidth}}
\toprule
\textbf{Regime / Pivot} & \textbf{Benchmark} & \textbf{Representative method} & \textbf{Headline metric} \\
\midrule
Managing (rearrangement) & Visual Room Rearrangement \cite{weihs2021visual} & TIDEE \cite{sarch2022tidee} & objects restored / episode success \\
Managing (preference tidy) & Housekeep \cite{kant2022housekeep} & LLM-based planner (paper baseline) \cite{kant2022housekeep} & correct-placement rate \\
Managing (mobile manip.) & HomeRobot OVMM \cite{yenamandra2023homerobot} & modular and RL baselines \cite{yenamandra2023homerobot} & overall success rate \\
Managing (search / recall) & EQA \cite{das2018embodied} / OpenEQA \cite{majumdar2024openeqa} & VLM with episodic memory (paper baselines) \cite{majumdar2024openeqa} & LLM-Match score \\
Collaborating (navigation) & SCAND \cite{karnan2022socially} & imitation from demonstrations \cite{karnan2022socially} & demonstration corpus (no leaderboard) \\
Collaborating (instruction) & ALFRED \cite{shridhar2020alfred} / TEACh \cite{Padmakumar2022TEACh} & seq2seq baselines \cite{shridhar2020alfred,Padmakumar2022TEACh} & task / goal-condition success \\
Collaborating (grounding) & ScanRefer \cite{chen2020scanrefer} / ReferIt3D \cite{achlioptas2020referit3d} & ScanRefer baseline \cite{chen2020scanrefer} & Acc@0.25 / 0.5 IoU \\
Empowering (BVI VQA) & VizWiz \cite{gurari2018vizwiz} & VQA baselines \cite{gurari2018vizwiz} & VQA accuracy \\
Empowering (assistive ADL) & Assistive Gym \cite{erickson2020assistive} & RL policies (paper baselines) \cite{erickson2020assistive} & task success / reward \\
\bottomrule
\end{tabular}
\end{table}

\section{Established Tools and Systems}
\label{subsec:p4-sec13-1-established-tools-and-systems}
Deployed assistants have so far approached the persistent, personalized spatial model of this chapter from opposite ends: software agents know a great deal about the user but nothing about the user's space, while home robots occupy the space but act with limited autonomy.

\textbf{Personal agents.} The most widely used assistants now act on a user's behalf in the digital world. OpenAI's ChatGPT agent\footnote{\url{https://openai.com/index/introducing-chatgpt-agent/}} (July 2025) runs on its own virtual computer, browsing and filling in forms across multi-step tasks and asking permission before consequential actions. Meta's Muse\footnote{\href{https://www.cnbc.com/2026/09/08/meta-personal-ai-agents-public-reckoning-privacy-safety.html}{CNBC, 8 September 2026}.} (September 2026), billed as a personal agent for everyday users, likewise runs in an isolated virtual machine and handles errands such as shopping and travel booking, and its request for access to email, calendars, and payments made privacy the main point of public debate around its launch. Voice assistants are being rebuilt on the same models: Amazon's Alexa+\footnote{\url{https://www.aboutamazon.com}}, generally available in the United States since February 2026, adds generative dialogue to the smart-home control that Alexa already provided, and Apple's Siri\footnote{\url{https://www.apple.com/siri/}} is being rebuilt on Google's Gemini models, adding personal context and on-screen awareness. For spatial intelligence, these agents matter mainly as the layer that holds personal context and preference; apart from smart-home control, none of them maintains a model of the physical home.

\textbf{Camera-grounded assistants.} A second line brings the physical world into the conversation through a camera. Gemini Live's camera mode\footnote{\url{https://gemini.google/}}, the product form of Google's Project Astra, lets a user point a phone at the surroundings and ask about what it sees, and Meta AI on camera glasses does the same from the wearer's viewpoint (Section~\ref{sec:p4_s9_3}). These systems answer spatial questions about the current view, and the assistive applications of Pivotal Task-3 build on the same capability, but they keep at most a short history of what they have seen, so the long-horizon spatial memory this chapter calls for remains outside them.

\textbf{Home robots.} Embodied assistants remain the least mature tier. Beyond mapping vacuums, the home humanoids introduced under Pivotal Task-1 are moving toward sale: 1X opened preorders for NEO in October 2025, at \$20{,}000 or a monthly subscription, with first deliveries promised for 2026 and with remote operators teleoperating tasks that its onboard autonomy cannot yet complete, and Figure's third-generation humanoid, Figure~03 (2025), is designed for homes and driven by its Helix vision-language-action model. The reliance on teleoperation shows how far these products are from the autonomous, open-vocabulary mobile manipulation of Pivotal Task-1, and it raises the in-home privacy question discussed below.

\section{Challenges}
\label{subsec:p4-sec13-1-challenges}

The three regimes share a common obstacle: the assistant must be reliable in the one setting that resists engineering, namely a real person's private and unstructured life. 
The challenges below are ordered from the perceptual frontier (generalization, memory) through the human frontier (personalization, social safety) to the deployment frontier (privacy, compute, trust), and each cuts across all three pivotal tasks above. 
Governance benchmarks for embodied agents have begun to measure the deployment end directly, testing whether a system stays controllable, recoverable, and auditable under perturbation rather than only whether it completes tasks \cite{qin2026embodiedgovbench}.

\textbf{Generalization to unstructured homes.} An assistant cannot be retrained per household, yet homes vary without bound in layout, lighting, and contents. 
Zero-shot results such as OK-Robot \cite{liu2024okrobot} are encouraging but remain far from the reliability a daily aide requires, and closing the gap depends on the open-vocabulary representations and sim-to-real transfer surveyed in Part~III rather than on more in-domain data.

\textbf{Long-horizon spatial memory.} Because the assistant persists in one place, it must remember spatial state over days and weeks: where objects were last seen, which routines the user follows, what was already done. 
Real-time scene-graph memories \cite{hughes2022hydra} and episodic-memory benchmarks \cite{majumdar2024openeqa} begin to address this, but the general memory mechanisms are treated in Section~\ref{subsec:p3_s7_3_6}; the assistant-specific unsolved problem is maintaining a consistent, updatable, and forgetting-resistant model over long deployments in a private home.

\textbf{Personalization and preference learning.} The assistant's geometric goals are user-specific, so it must learn an individual's placement rules, comfort distances, and routines from sparse interaction. 
Few-shot preference summarization \cite{wu2023tidybot} and learnable social cost functions are early steps, but acquiring rich preferences without burdening the user remains an open problem.

\textbf{Social compliance and safety.} Operating in continuous contact with people turns motion into a social and safety-critical act. 
Anticipatory, legible navigation and motion \cite{trautman2010unfreezing,dragan2013legibility} and shared-control safety zones \cite{ghorbel2018decision} mitigate risk, yet guaranteeing comfort and physical safety across the diversity of real human behavior is far from settled.

\textbf{Privacy of persistent in-home perception.} An assistant that continuously perceives a private home raises acute privacy concerns, the more so for assistive cameras worn by users who cannot themselves audit what is captured. 
The concern becomes sharper when assistance depends on remote operation or cloud annotation, because the raw spatial stream itself may reveal the user's routines and vulnerabilities. Persistent spatial memory must be reconciled with selective perception, on-device processing, and user control over what is retained.

\textbf{On-device compute and trust.} Wearable and home platforms are power- and thermally constrained, forcing the heavy perception and reasoning stacks of Part~III onto efficient, often on-device, inference. 
Compounding this, assistance for vulnerable users is safety-critical, so calibrated reliability and appropriate user trust, neither over- nor under-reliance, are prerequisites for deployment rather than refinements of it.


\newpage
\chapter{Spatial Intelligence for Others}
\label{sec:p4_s14}

\FramedInitial{O}{ther} application domains of spatial intelligence are too numerous to enumerate, because spatial intelligence is a horizontal capability: any field that must perceive, reconstruct, localize within, or reason about the geometry of the physical world inherits the methods of Part~III. 

This chapter takes five further domains where that inheritance is already visible in deployed systems. Each stresses a capability the earlier chapters do not: registering a building site against its design, tracking many near-identical agents through occlusion, reconstructing fragile objects faithfully, working without clean light or satellite positioning underwater, and binding people to products across a network of cameras.

\section{Architecture, Engineering, and Construction}
\label{sec:part-4-section-14-architecture-engineering-and-construction}

The architecture, engineering, and construction (AEC) sector is built around a single spatial-intelligence problem: continuously registering the messy as-built reality of a site against the structured as-designed model that specifies what it should become. 
This is the same Building-Information-Modeling (BIM) substrate that anchors the city-scale digital twins of Section~\ref{sec:p4_s11_1}, but here the model is the working blueprint of an individual building under construction rather than a synchronized replica of a finished one, and the central task is detecting and acting on the discrepancy between plan and reality.
Figure~\ref{fig:p4_s14_aec} sketches this as-built versus as-designed loop. \textbf{Scan-to-BIM} is its perception front end: dense laser or photogrammetric scans of a structure are semantically parsed into walls, floors, columns, and components, building on the indoor point-cloud understanding of Chapter~\ref{chapter:3:sec:6} and its benchmarks S3DIS \cite{armeni2016s3dis} and ScanNet \cite{dai2017scannet}. 
A specialized sub-problem is vectorizing a scan into an editable floorplan, where synthetic supervision such as Structured3D \cite{zheng2020structured3d} and transformer models that predict all room polygons in parallel \cite{yue2023roomformer} now recover clean geometry from cluttered scans. 
\textbf{Progress monitoring} closes the loop in time: by registering an image-based or scanned as-built point cloud against a four-dimensional IFC-BIM model, deviations between planned and actual construction state can be recognized automatically \cite{golparvarfard2015automated}, the spatiotemporal core of the commercial reality-capture platforms now common on job sites. 
\textbf{Construction robotics} acts on the model directly, with mobile fabricators that localize against the design to print, place, or assemble structural elements at full scale \cite{giftthaler2017mobile}. 
Finally, \textbf{structural inspection} turns the same perception stack toward maintenance, localizing cracks and defects on concrete surfaces from imagery \cite{cha2017deep}, increasingly from a UAV-borne 3D reconstruction of a bridge or facade. Reality-capture and BIM tooling from Matterport,\footnote{\url{https://matterport.com/}} Autodesk,\footnote{\url{https://www.autodesk.com/products/recap/overview}} and Leica,\footnote{\url{https://leica-geosystems.com/products/laser-scanners}} progress-monitoring services such as OpenSpace\footnote{\url{https://www.openspace.ai/}} and Buildots,\footnote{\url{https://buildots.com/}} and on-site robots from Dusty Robotics\footnote{\url{https://www.dustyrobotics.com/}} and Built Robotics\footnote{\url{https://www.builtrobotics.com/}} are representative of the deployed layer.

\begin{figure}[t]
    \centering
    \includegraphics[width=0.98\linewidth]{figures/part4/section14/other_aec.pdf}
    \caption{AEC tasks around an as-designed BIM: (1) semantic parsing of site scans; (2) comparison with the scheduled design state, highlighting a planned element not observed in the scan; (3) model-guided robotic placement; and (4) crack localization linked to a model element. Physical work prompts re-capture of the site. Scenes and point clouds are illustrative.}
    \label{fig:p4_s14_aec}
\end{figure}

\section{Sports Analytics and Performance}
\label{sec:part-4-section-14-sports-analytics-and-performance}

Sports analytics lifts raw broadcast or arena video into a calibrated metric model of the playing surface, then reasons over the resulting positional field. What sets it apart is fine-grained, multi-agent spatiotemporal tracking under severe and persistent occlusion, where many near-identical players and a small fast ball must be followed and assigned consistent identities at sub-second resolution.
Figure~\ref{fig:p4_s14_sports} illustrates three linked capabilities: registered trajectories feed tactical analysis, while calibrated multi-view pose supports biomechanics. \textbf{Tracking and registration} form the geometric foundation: multi-object tracking benchmarks such as SoccerNet-Tracking \cite{cioppa2022soccernet} and the multi-sport SportsMOT \cite{cui2023sportsmot} drive player and ball association, while sports-field registration estimates the homography from the broadcast frame to a metric pitch template so that pixel positions become field coordinates \cite{jiang2020optimizing}. 
\textbf{3D pose and biomechanics} recover athlete motion in three dimensions, with markerless capture systems estimating not only kinematics but joint dynamics from ordinary video, as in OpenCap \cite{uhlrich2023opencap}, typically by triangulating multi-view 2D pose into a 3D skeleton \cite{iskakov2019learnable}. 
\textbf{Tactical and team-space analysis} then treats the positional field itself as the object of study, valuing the space each player controls and forecasting where players will move; deep models of expected possession value \cite{fernandez2019decomposing} and multi-agent trajectory generation \cite{yeh2019diverse} build directly on the spatial-reasoning methods of Chapter~\ref{sec:reasoning}. 
The deployed layer is mature, spanning optical tracking and officiating from Hawk-Eye,\footnote{\url{https://www.hawkeyeinnovations.com/}} player-tracking data from Second Spectrum (now part of Genius Sports)\footnote{\url{https://www.geniussports.com/}} and Stats Perform,\footnote{\url{https://www.statsperform.com/}} and wearable spatial sensing from KINEXON\footnote{\url{https://kinexon-sports.com/}} and Catapult.\footnote{\url{https://www.catapult.com/}}

\begin{figure}[t]
    \centering
    \includegraphics[width=0.98\linewidth]{figures/part4/section14/other_sports.pdf}
    \caption{Sports analytics: (1) player/ball tracking and planar registration; (2) calibrated multi-view pose, kinematics, and model-estimated loads; and (3) tactical analysis from registered trajectories. The tactical view enlarges the dashed pitch region, preserving player IDs. Shading illustrates team control, dotted lines candidate passes, and dashed arrows predicted motion. Scenes and predictions are schematic.}
    \label{fig:p4_s14_sports}
\end{figure}


\begin{figure}[!t]
    \centering
    \includegraphics[width=0.98\linewidth]{figures/part4/section14/other_cultural_heritage.pdf}
    \caption{Three cultural-heritage tasks: (1) multi-view capture and surface reconstruction; (2) geometric and appearance-based matching of fragments for virtual reassembly; and (3) classification of airborne-LiDAR ground returns to reveal candidate structures in bare-earth relief. The tasks are independent; scenes, mesh, and terrain are illustrative.}
    \label{fig:p4_s14_heritage}
\end{figure}


\section{Cultural Heritage and Archaeology}
\label{sec:part-4-section-14-cultural-heritage-and-archaeology}

Cultural heritage and archaeology apply spatial intelligence to the past, running from capture through recovery to discovery. What distinguishes it is high-fidelity, metrically accurate, and usually non-contact reconstruction, coupled to a preservation imperative: the artifact may be fragile, fragmentary, or buried, and the digital record must be faithful enough to stand in for an object that cannot be handled or that no longer exists intact.
Figure~\ref{fig:p4_s14_heritage} organizes the field into three movements. 
\textbf{Digitization} captures artifacts and monuments in three dimensions, a lineage that runs from the landmark in-field laser scanning of the Digital Michelangelo Project \cite{levoy2000digital} through structure-from-motion photogrammetry \cite{schonberger2016structure} to the neural radiance fields \cite{mildenhall2021nerf} and 3D Gaussian splatting \cite{kerbl20233d} now used for photorealistic heritage capture, all drawing on the reconstruction methods of Chapter~\ref{sec:generation}. 
\textbf{Virtual restoration} uses geometric AI to reassemble what time has broken, posing fragment reassembly as a 3D puzzle over fractured parts, with large benchmarks such as Breaking Bad \cite{sellan2022breaking} and the real fresco-reassembly dataset RePAIR \cite{tsesmelis2024repair} driving the methods. 
\textbf{Archaeological prospection} discovers sites at landscape scale, most strikingly through airborne LiDAR whose bare-earth models reveal structures hidden under forest canopy, as in the canopy-penetrating surveys of the Maya lowlands at Caracol \cite{chase2011airborne} and across northern Guatemala \cite{canuto2018ancient}. Organizations such as CyArk\footnote{\url{https://www.cyark.org/}} and the Factum Foundation,\footnote{\url{https://www.factumfoundation.org/}} together with photogrammetry tools including Agisoft Metashape\footnote{\url{https://www.agisoft.com/}} and the Open Heritage 3D archive,\footnote{\url{https://openheritage3d.org/}} anchor the practical layer.



\begin{figure}[t]
    \centering
    \includegraphics[width=0.98\linewidth]{figures/part4/section14/other_marine.pdf}
    \caption{Underwater spatial intelligence under light attenuation/scattering and unavailable submerged GPS fixes. Enhanced visual observations combine with acoustic, inertial, and depth inputs for localization and mapping. Terrain and structure maps support coverage and standoff planning for hull/pipeline inspection; new views close the loop. Blue marks the survey track and selected defect. Scenes and enhancement are illustrative.}
    \label{fig:p4_s14_marine}
\end{figure}

\section{Marine and Underwater Systems}
\label{sec:part-4-section-14-marine-and-underwater-systems}

Underwater spatial intelligence confronts a setting in which the two assumptions that anchor terrestrial methods both fail at once: light no longer propagates cleanly, since wavelength-dependent attenuation and scattering degrade vision with range, and satellite positioning does not penetrate water, so absolute position must be inferred from acoustics, inertia, or the structure being surveyed.
Figure~\ref{fig:p4_s14_marine} depicts the resulting perceive-localize-map-inspect loop. 
\textbf{Navigation and SLAM} solve the GPS-denied localization problem, surveyed comprehensively for autonomous underwater vehicles by Paull \etal \cite{paull2014auv}, with foundational deep-sea visual SLAM demonstrated on the wreck of the Titanic \cite{eustice2005visually} and modern visual-inertial-pressure localization benchmarked by datasets such as AQUALOC \cite{ferrera2019aqualoc}. 
\textbf{Perception and enhancement} contend with the optical medium itself: physically grounded models remove the water from an image to recover true scene color and range \cite{akkaynak2019sea}, learned enhancement restores contrast in real time for downstream perception \cite{islam2020fast}, and benchmarks such as UIEB standardize the evaluation \cite{li2020underwater}, while large annotated image databases like FathomNet \cite{katija2022fathomnet} enable ocean-scale detection. 
\textbf{Mapping and inspection} reconstruct the seafloor, habitats, and submerged structures, with structure-relative perception and coverage planning enabling autonomous inspection of complex ship hulls \cite{hover2012advanced}. The operational layer spans low-cost platforms such as the BlueROV2,\footnote{\url{https://bluerobotics.com/store/rov/bluerov2/}} survey vehicles from Kongsberg\footnote{\url{https://www.kongsberg.com/discovery/autonomous-and-uncrewed-solutions/auv/hugin/}} and research institutions including WHOI\footnote{\url{https://www.whoi.edu/}} and MBARI.\footnote{\url{https://www.mbari.org/}}




\begin{figure}[t]
    \centering
    \includegraphics[width=0.98\linewidth]{figures/part4/section14/other_retail.pdf}
    \caption{Retail and smart-space intelligence across calibrated ceiling cameras: (1) shopper--product association updates a virtual basket; (2) detected shelf contents are compared with a planogram; (3) ground-plane tracks yield shopper-flow and visit-density maps. Dotted lines link scene details; the geometry and density are illustrative.}
    \label{fig:p4_s14_retail}
\end{figure}

\section{Retail and Smart Spaces}
\label{sec:part-4-section-14-retail-and-smart-spaces}

Retail and smart-space intelligence reduces to one persistent spatial problem: jointly associating, across a network of calibrated overlapping cameras and over time, which person is interacting with which product at which shelf location. 
Unlike the warehouse logistics of Section~\ref{sec:p4_s9_5}, where the environment is engineered for the robot, here the space is a crowded store full of untracked people, and its distinguishing challenge is dense multi-object tracking with reliable human-object association through frequent mutual occlusion.
Figure~\ref{fig:p4_s14_retail} shows how the major applications share this core. 
\textbf{Multi-camera tracking} provides the spatial backbone, with benchmarks such as WILDTRACK \cite{chavdarova2018wildtrack} and MMPTRACK \cite{han2023mmptrack} and methods that project per-camera features onto a shared ground plane to detect and track people across views \cite{hou2020multiview}; this is the machinery behind cashierless checkout, where a shopper's path and pickups must be reconstructed without a checkout line. 
\textbf{Shelf monitoring} turns perception toward the products, detecting densely packed items for inventory and planogram compliance, supported by the dense-detection benchmark SKU-110K \cite{goldman2019precise} and by methods that formulate compliance as matching the observed shelf layout against the ideal one as a sub-graph isomorphism \cite{tonioni2017product}. 
\textbf{In-store analytics} aggregates the same tracks into shopper-flow heatmaps and layout insights, building on cross-camera re-identification \cite{zheng2015scalable} to stitch one customer's trajectory across non-overlapping views.
The deployed layer includes cashierless systems such as Amazon's Just Walk Out\footnote{\url{https://www.justwalkout.com/}} and platforms from Trigo,\footnote{\url{https://www.trigoretail.com/}} AiFi,\footnote{\url{https://www.aifi.com/}} and Zippin,\footnote{\url{https://www.getzippin.com/}} alongside shelf-analytics vendors such as Trax.\footnote{\url{https://traxretail.com/}}


\begin{table}[t]
\centering
\fontsize{8}{10}\selectfont
\setlength{\tabcolsep}{1.0mm}
\caption{Representative benchmarks for the five domains of this chapter.}
\label{tab:p4_s14_benchmarks}
\begin{tabular}{L{0.16\textwidth} L{0.30\textwidth} L{0.28\textwidth} L{0.20\textwidth}}
\toprule
\textbf{Domain} & \textbf{Benchmark / dataset} & \textbf{Representative method} & \textbf{Headline metric} \\
\midrule
AEC & S3DIS \cite{armeni2016s3dis} & point-cloud semantic segmentation & mIoU \\
AEC & Structured3D \cite{zheng2020structured3d} & RoomFormer \cite{yue2023roomformer} & room / corner / angle F1 \\
Sports & SoccerNet-Tracking \cite{cioppa2022soccernet} / SportsMOT \cite{cui2023sportsmot} & multi-object trackers & HOTA \\
Sports & broadcast field registration & learned-error optimization \cite{jiang2020optimizing} & IoU (whole / part) \\
Heritage & Breaking Bad \cite{sellan2022breaking} / RePAIR \cite{tsesmelis2024repair} & learned fragment reassembly & rotation / translation error \\
Marine & UIEB \cite{li2020underwater} & Water-Net \cite{li2020underwater} & PSNR / SSIM \\
Marine & AQUALOC \cite{ferrera2019aqualoc} & visual-inertial-pressure SLAM & trajectory error \\
Retail & WILDTRACK \cite{chavdarova2018wildtrack} & MVDet \cite{hou2020multiview} & MODA \\
Retail & SKU-110K \cite{goldman2019precise} & dense detector with EM-merger \cite{goldman2019precise} & AP \\
\bottomrule
\end{tabular}
\end{table}



\section{Outlook}
\label{sec:part-4-section-14-outlook}

The five domains above, whose benchmarks are collected in Table~\ref{tab:p4_s14_benchmarks}, are a sample, not a census. 
Spatial intelligence is equally at work in fields this survey does not treat at length: virtual try-on and 3D e-commerce, immersive education and high-fidelity training simulators, security and surveillance, energy and infrastructure inspection, precision agriculture beyond the remote-sensing view of Section~\ref{sec:p4_s11_2}, and many others.
What unites them, and what unites them with the headline applications of the preceding chapters, is that each takes the perception, reasoning, and generation capabilities of Part~III and specializes them to the geometry, sensing modality, and constraints of a particular slice of the physical world.

\part{Discussions and Futures}
% ================================
% \section*{\centering \bf Part V: Future and Challenges}
% \addcontentsline{toc}{section}{Part V: Future and Challenges}
% ================================

\begin{figure}
    \centering
    \includegraphics[width=\linewidth]{figures/part5_opener_v2.pdf}
    
    \caption{Part V in one view. (1) What it requires: a household robot must keep one consistent account of its kitchen while it turns away, people move things, and a stack of cans may topple. (2) Where the gaps are: the loop from perception to action still breaks, benchmark success diverges from deployment, and consistency erodes over long horizons. (3) What comes next: open problems that set the directions for future work.}
    \label{fig:p5_main}
\end{figure}

Parts~II to~IV described how spatial representations are constructed, how they support perception, reasoning, and generation, and how they are deployed in applications ranging from household robots and surgical systems to city-scale digital twins.
However, current systems often perceive or act well at a single moment, while most of the open problems arise when a system must maintain a consistent picture of the world as it moves, revisits places, and acts over extended periods.

A household robot that puts away groceries must remember the location of a cupboard after turning away from it, register a chair that has been pushed into its path, track which bags it has already emptied, and judge whether a stack of cans will remain stable if another can is added.
Methods described earlier in this handbook handle each of these steps reasonably well in isolation.
The difficulty lies in maintaining a single account of the kitchen that all steps read from and write to, and in keeping that account correct while the robot and the people around it change the scene.
In the formalism of Chapter~\ref{chapter:definition}, this account is the world state $\mathcal{R}^{\mathcal{S}}_t$, and keeping it correct is the role of the update $\Psi_{\mathcal{S}}$.
It is also the problem that Section~\ref{sec:historical-evolution} identifies as open for the current era, long-horizon consistency under motion and partial observability.

Chapter~\ref{chapter:p5_discussions} asks what makes such a world state usable: what can be inferred from limited observations, what must be represented, how beliefs are kept current, when predictions can be trusted, and how spatial knowledge is acquired through interaction in shared space.
We present the evidence on each question and state which answers remain open.
Chapter~\ref{chapter:p5_challenges} turns to gaps in current methods and benchmarks, the difficulties of deployment, and directions for future work. 
Chapter~\ref{chapter:p5_conclusion} concludes by revisiting the definition of spatial intelligence introduced in Chapter~\ref{chapter:definition}.

\newpage

\chapter{Foundations: What Spatial Intelligence Requires}
\label{chapter:p5_discussions}



% \lettrine[lines=3]{\initfamily\textcolor{darkgreen}{O}}{pen}
% Vision-based ranger involves leveraging geometric principles to calculate distance.
% A vision-based ranger could use one or multiple cameras.

\FramedInitial{A}{dvances} in spatial perception, reasoning, and generation have enabled AI systems to recover geometry from images, answer questions about layout, and simulate complex environments, as Parts~II to~IV documented.
Success on these tasks does not establish that a system maintains a reliable model of the physical world.
A model can reconstruct a room to millimeter accuracy, answer most questions about a photograph, or continue a video convincingly, and still fail when observations are incomplete, when the environment changes while it is not looking, or when its actions depart from the conditions under which it was trained.
The failures reported throughout this handbook cluster at these points, which suggests a shared deficiency rather than a collection of unrelated task-level shortcomings.

This chapter asks what makes an internal model of physical space usable for continued reasoning and action.
By usable we mean that the model can be queried for what a task requires, extended when new evidence arrives, and relied upon when its predictions are acted on; we do not presuppose that such a model must be an explicit 3D reconstruction, nor that it must be a video generator.
The question decomposes into five, each posed at a different level.
What can be inferred from limited observations, and what remains underdetermined (Section~\ref{sec:p5_observability})?
Which information must a representation retain, and for which operations (Section~\ref{sec:p5_sufficiency})?
How can a world state be maintained and revised over time (Section~\ref{sec:p5_persistence})?
When can its predictions be trusted (Section~\ref{sec:p5_validity})?
How is spatial knowledge acquired and used through interaction, in environments shared with other agents (Section~\ref{sec:p5_agency})?
The five are interdependent: action produces new evidence, which re-enters at the first level and constrains all of the others.


\section{Spatial Observability and Identifiability}
\label{sec:p5_observability}

\FramedInitial{S}{patial} understanding begins with an incomplete view of the world.
A system never has direct access to the world state $\mathcal{R}^{\mathcal{S}}_t$ of Chapter~\ref{chapter:definition}; it has images, video, depth, or sound, from which that state must be inferred.
Images and videos carry rich evidence about appearance and structure, but they do not uniquely determine every geometric, physical, or functional property of a scene.
A single image fixes the scale of a scene only through the prior knowledge brought to it (Section~\ref{sec:height_distance}), an occluding surface hides whatever lies behind it, and the mass of a box or the state of a lock leaves no trace in a photograph.
A usable model must therefore keep track of which of its contents are observed and which are inferred.

Whether a spatial inference fails because the model cannot extract the evidence available, or because the evidence itself is insufficient, is a distinction current evaluation rarely draws.
The former is a limitation of learning or computation, and better models or more data can be expected to reduce it.
The latter is a problem of identifiability, which more capable inference alone does not remove.
Asked whether a sofa seen in one photograph will fit through a door seen in another, a multimodal model typically produces a confident answer with a coherent explanation, yet without a common scale between the two images the question has no determinate answer, and a model with perfect geometric reasoning should say so.
Evaluation that counts every wrong answer as a model failure conceals this class of cases, while evaluation that accepts every ambiguity as irreducible would excuse failures that are not.

Additional modalities constrain different aspects of the underlying state, and their contributions are neither interchangeable nor universally sufficient.
A second viewpoint supplies parallax, which fixes relative depth and reveals part of what the first view occluded.
A depth sensor supplies distances directly, with an accuracy and a scale behavior that depend on the sensing principle (Section~\ref{sec:depth_estimation}).
Motion adds constraints on dynamic geometry; sound carries direction, distance, and room structure (Section~\ref{sec:av_localization}); contact and force reveal mass, friction, and compliance, which no camera observes (Section~\ref{subsec:p3_s7_1_2}).
Assessing a modality therefore requires specifying which ambiguities it resolves, under which assumptions, and which properties remain unidentifiable once it has been added.
Part~II treated the modalities as alternative representations of space (Chapter~\ref{chapter:single_modality}); here they appear as complementary constraints on one underdetermined state.

Observability is not fixed by the initial input.
Moving to another viewpoint, changing the illumination, or pushing an object can reveal previously ambiguous properties, and some properties, such as a drawer's operability, may require physical interaction.
Theory of Space and VABench bear directly on this point: models perform better with supplied observations than when they must gather evidence themselves, and allowing a model to move its camera and look again raises its manipulation success substantially \cite{zhang2026theory,zhang2026vabench}.
Spatial intelligence therefore depends partly on recognizing when an additional observation would change what is knowable.
How an agent should choose such observations is a question of decision making, taken up in Section~\ref{sec:p5_agency}, which properties are identifiable depends on the agent's available actions.

Distinguishing uncertainty from fundamental underdetermination, and stating what evidence would resolve it, are abilities that current systems are not asked to demonstrate.
Static question accuracy cannot distinguish correct reasoning under insufficient evidence from incorrect reasoning despite sufficient evidence, and controls as simple as removing the image or shuffling the input remain rare (Section~\ref{sec:vl_grounding}).
A benchmark designed around identifiability would include questions that are unanswerable from the given input and reward a model for saying so, pair each ambiguous case with the modality or action that resolves it, and score whether the model asks for it.
Given these limits, which obtainable properties should a spatial representation preserve?


\section{Representation and State Sufficiency}
\label{sec:p5_sufficiency}
\label{sec:p5_cognitive}

\FramedInitial{T}{he} same physical environment can be represented through images, point clouds, occupancy grids, scene graphs, neural fields, or the latent state of a learned model, and Part~II surveyed each of these forms (Chapter~\ref{chapter:single_modality}).
They differ less in how much information they hold than in which aspects of the environment they preserve, which they abstract away, and how they organize what remains.
Scene graphs typically emphasize entities and their relations, Gaussian representations primarily encode appearance and geometry, and a latent state retains what its training objective rewarded; each can be extended with further structure, but none is informative in isolation from the questions and operations it is meant to support.

That spatial knowledge and spatial reasoning are separable is shown by large models that acquire the former without the latter.
A transformer trained only on Othello move sequences forms an internal representation of the board (Section~\ref{sec:rep_language}), and language models encode the geographic positions of cities in their activations without ever being shown a map \cite{gurnee2023language}.
The same models remain well below human performance on perspective taking, mental rotation, and building a model of a place from a sequence of views \cite{cai2025holistic}, and chain-of-thought prompting, which helps on arithmetic and code, tends to degrade their accuracy on visual spatial benchmarks \cite{kancheti2026cotdegrades}.
Performance recovers when the model is allowed to draw a coarse cognitive map, render an imagined view, or call a geometric tool before answering (Sections~\ref{subsec:p3_s7_3_5} and~\ref{subsec:p3_s7_3_8}).
Encoding spatial information and computing with it are therefore distinct requirements, and reasoning performance depends partly on whether the representation in use supports the operations the task demands.
This does not show that language is unsuited to spatial reasoning; it suggests that a representation in which distances and viewpoints are described rather than held tends to accumulate error each time they are manipulated.

Spatial tasks differ not only in the information they require but in the computations they demand of it.
Judging a relative position needs a reference frame and a relation; metric reasoning needs distances, angles, and a scale; changing viewpoint needs a rigid or projective transformation applied to 3D structure; navigation needs topological connectivity and traversability; manipulation needs local geometry together with contact constraints and material properties.
A representation suited to semantic relations can be insufficient for metric transformation, and a detailed geometric reconstruction can still omit the properties on which physical interaction depends, a gap Section~\ref{sec:p4_s11_1} summarized as geometry is not behavior.
The effectiveness of a representation is therefore a matter of its compatibility with the downstream operations.
The recurring finding of Part~III that hybrid systems are limited by their weakest interface (Section~\ref{sec:p3_s7_4}) is an instance of this: the interface is where one representation's operations are asked to run on another representation's content.

These observations suggest a task-relative criterion of sufficiency.
Spatial-state sufficiency depends on the queries, predictions, and actions it must support; a representation is sufficient for a family of tasks if every decision the family requires can be made from it as well as from the world itself.
Information irrelevant to one task may be essential to another, so sufficiency cannot be judged from the representation alone.
This creates a tension between compact states optimized for an immediate objective and general-purpose states intended to remain useful across tasks that are not yet known.
A latent world model can succeed at the tasks it was trained for while modeling the environment's dynamics unfaithfully \cite{gonzalezlaiz2026limits}, which satisfies sufficiency for those tasks and gives no assurance for others; a dense geometric reconstruction, by contrast, serves many queries but remains silent about physical properties that are not visible.

Whether a finite shared state can preserve the distinctions required by an open-ended range of tasks, or whether a spatial system should instead maintain adaptable representations whose detail and structure depend on downstream demands, remains open.
Shared spatial latents read by many task heads, task-specific representations maintained in parallel, and interfaces through which modules request information at the precision they need (Section~\ref{sec:p3_s7_4}) represent three answers, and the available evidence does not yet favor one, particularly when future tasks are unknown.

Evaluating a spatial representation therefore requires examining the information it preserves and the operations it supports, rather than its expressiveness in the abstract.
Reconstruction quality and recognition accuracy are the usual proxies, and both are insufficient: a photometrically accurate replica can fail a metric tolerance check and encode nothing about which drawers open (Section~\ref{sec:p4_s11_1}), and a correct answer can rest on language priors rather than on a correct belief about the scene.
Sufficiency has to be tested across tasks, by asking whether a representation fixed for one purpose supports queries, transformations, and action predictions it was not built for.
Whether a state that is sufficient at one moment remains so as the world changes and objects leave the field of view is a separate question, taken up next.


\section{Persistent World States and Belief Updating}
\label{sec:p5_persistence}
\label{sec:p5_worldstate}

\FramedInitial{U}{nderstanding} a scene at one moment differs from maintaining an understanding of a changing world.
An assistant asked where the keys were left must have observed them, retained their location after the user walked away, distinguished that remembered location from the current view, and allowed for the possibility that someone moved them in the meantime.
Egocentric benchmarks built around such queries show that an explicit memory, updated as video arrives and tracking objects relative to the moving user, substantially improves leading multimodal models \cite{wang2026keepitinmind}, whereas most commercial glasses answer from a single frame and keep no map (Section~\ref{sec:p4_s9_3}).
The household robot in this part's introduction faces the same requirement at every step.
In the formalism of Chapter~\ref{chapter:definition} this is the role of the update $\Psi_{\mathcal{S}}$, the component that Parts~II to~IV examined least.

Access to a long observation history does not ensure a consistent world state.
Video world models steered through long explorations fail to render revisited places consistently even as image quality improves \cite{wang2026worldprobe}; multimodal models given several views of a room rarely detect that one view was captured after an object moved \cite{cao2026samespace}; and over routes spanning several streets, models answer by recognizing landmarks and matching text, with little evidence of path integration \cite{mahendran2026kilometervision}.
A model may retain every earlier frame and still fail to recognize a revisited place, preserve an object's identity, or decide which remembered properties have become outdated.
Research on spatial memory has concentrated on recall, through longer contexts, better retrieval, and richer memory structures (Section~\ref{subsec:p3_s7_3_6}), while these failures concern revision.
Persistence therefore involves maintaining and revising beliefs about the world, not merely retaining observations.

A persistent world state must distinguish absence of observation from evidence of absence, maintain uncertainty during occlusion, and reconcile new observations with prior beliefs.
When a new observation conflicts with a stored one, the system must decide whether the environment changed, the earlier estimate was wrong, or the two observations refer to different entities, each requiring different state updates.
Classical systems partly address this problem: SLAM detects loop closures and corrects the map upon revisiting a known place (Section~\ref{sec:mapping_slam}), object trackers maintain identities through occlusion, and dynamic scene graphs timestamp their edges.
Learned systems often struggle to make such revisions reliably over extended interactions.
Models retain outdated beliefs after the world changes \cite{zhang2026theory}, append-only memories accumulate contradictory entries, and agents drift over long episodes as their context fills (Section~\ref{subsec:p3_s7_3_7}).
Increasing capacity alone may not suffice, since a memory that only grows keeps each outdated belief beside its replacement.
A write policy that decides what enters the state, with what confidence, and when an entry is revised or removed, is a prerequisite that current systems leave largely implicit.

Where such a state should be kept is a question of architecture rather than of representation format.
Explicit structures such as maps, scene graphs, and splats (Section~\ref{sec:rep_3d}) persist by construction and make selected state variables accessible and correctable, but they must anticipate every type of object, property, and change they will be asked to encode.
Learned memories and recurrent latent states can encode properties no designer specified, and whether they track the world is difficult to verify \cite{gonzalezlaiz2026limits}.
Generated frames are the most flexible option and lose whatever leaves the context window.
Hybrid designs that store static geometry and dynamic objects separately and retrieve them on revisit \cite{yu2026grome}, or that keep slowly changing layout and identities in an explicit layer and leave appearance and short-term dynamics to a learned one, combine these advantages at a cost of their own: geometry, object identity, uncertainty, and temporal updates must stay synchronized across layers, and coordinate frames and scales can silently diverge at the interface (Section~\ref{sec:p3_s7_4}).
The separation can be taken further.
State World Models keep a typed state ledger that evolves under executable rules written by a coding agent and use a frozen video model only to render each resulting event; on counterfactual twin pairs that share an initial frame and action but differ in hidden state, this harness produces the rule-required outcome in 92\% of pairs, where the best of eleven video models prompted with the same rules reaches 32\% \cite{chan2026stateworld}.
The result suggests that what fails in generator-only designs is the maintenance of state the current frame does not show, rather than rendering.
The appropriate architecture follows from the requirements of state maintenance, how long each part must persist and whether it must be inspected, rather than from a general preference for explicit or learned form.

World states should remain persistent without becoming rigid and revisable without losing established structure, and evaluation currently measures neither property.
Most benchmarks refer to a single moment, and the protocols that would expose persistence failures are not yet standard (Section~\ref{sec:p3_s7_5}): streaming protocols that withhold future frames \cite{li2026ovosbench}, revisit protocols that move the camera away and back \cite{wang2026worldprobe}, and change protocols that alter the scene behind the agent.
A more informative evaluation reads out the state itself, as Theory of Space does at each exploration step \cite{zhang2026theory}, and scores separately what the system has observed, what it remembers, and what it predicts, since conflating the three is how a stale belief passes as a current one.
Deployed systems could likewise report how far their world state lags behind the world it describes, a quantity that digital twins call the liveness gap (Section~\ref{sec:p4_s11_1}) and that is rarely measured elsewhere.
A state that remains consistent over hours can nonetheless describe appearance accurately and physical behavior poorly.


\section{Predictive Validity and Physical Grounding}
\label{sec:p5_validity}
\label{sec:p5_plausible}

\FramedInitial{A}{} video generator asked to continue a clip of a ball rolling off a table usually produces a smooth, well-lit continuation in which the ball falls, often along the wrong trajectory or with an implausible bounce.
Benchmarks of physical plausibility find that the physical accuracy of video generators is poor and does not improve with visual quality (Section~\ref{sec:p3_s7_5}), and Chapter~\ref{sec:generation} reported the same pattern across its subfields: shapes that change between views, feet that slide over the ground, and motions that violate conservation (Section~\ref{sec:p3_s8_8}).
A plausible continuation reproduces the appearance of a physical event without predicting its trajectory, contacts, or consequences.
Yet the same class of models learns substantial scene structure, since image generators adapted for perception segment, estimate depth, and ground referring expressions at the level of specialist models \cite{gabeur2026image}.
Structure and dynamics are apparently learned to different degrees from the same data.

World-model correctness has several distinct meanings, and the criteria are often run together.
Visual plausibility concerns realistic appearance; geometric consistency concerns coherent spatial structure across views and time; dynamical fidelity concerns how the world evolves under actions; decision reliability concerns whether actions planned in the model succeed in reality.
The criteria are related but not interchangeable, and they are ordered by how much they depend on action.
Common video metrics measure the first and do not register physically impossible motion (Section~\ref{subsec:p3_s8_3_6}); view-synthesis metrics measure the second and cannot tell a solid object from a cut-out that matches the training views (Section~\ref{subsec:p3_s8_2_1}); the third has no standard protocol (Section~\ref{subsec:p3_s8_6_discussion}); the fourth is tested, when it is tested at all, by deploying the resulting policy.
A model can score well on the first two and fail the last two, and vision-language models used as judges share the spatial blind spots of the models they evaluate (Section~\ref{sec:p3_s8_8}).

Learning how scenes typically evolve from observational data does not guarantee correct predictions under novel interventions.
In the action-conditioned definition of Chapter~\ref{chapter:definition}, $f^{\mathcal{S}}_{\text{pred}}$ maps a state and an action to the next state, but a generator trained on video sees only the actions that happened to be recorded, mostly everyday activity and camera motion, and sees them correlated with the states and circumstances in which they occurred.
A model that has watched thousands of doors being pushed open has learned what usually follows a push, not what determines the outcome, and so has little basis for predicting the locked door, the obstructed door, or the push applied at an unusual angle.
The regularities it learns are real and useful for geometry, appearance, and common motion; they constrain little about the consequences of an action outside the data.
An interventional test, which fixes the initial state, varies the action, and checks whether predicted outcomes vary as real outcomes do, is the most direct way to distinguish a model of what tends to happen from a model of what an action causes.

Accurate geometry does not imply an actionable physical model.
Mass, friction, articulation, and compliance are often visually unobservable yet critical for manipulation (Sections~\ref{subsec:p3_s7_1_2} and~\ref{sec:p4_s9_1}), and a Gaussian-splat replica of a room can be photometrically faithful while encoding nothing about which drawers open or how much force a door requires (Section~\ref{sec:p4_s11_1}).
Real-to-sim pipelines estimate such quantities by manipulating objects and observing the response \cite{pfaff2025scalable}, thereby estimating properties needed for prediction, and world-action models that predict compact encodings of future observations rather than pixels avoid spending capacity on appearance that control does not use \cite{yu2026jam}.
Confidence sets calibrated under one data-collection policy may become unreliable under policy shifts \cite{mittal2026calibrated}, and model-based action selection may induce policy shifts that invalidate this calibration.

A tractable form of the question is where, rather than whether, a world model can be trusted.
Treated like a measuring instrument, a world model would be calibrated against the environment it claims to simulate, carry a stated range of states and actions within which it is valid, and be used with caution outside it.
Establishing that range requires tests under controlled interventions and unfamiliar actions, under environmental change, and against downstream decisions, in place of average visual or predictive accuracy.
World models expressed as code, whose rules are explicit and correctable when they conflict with observation, admit such tests more readily than learned ones (Section~\ref{subsubsec:p3_s8_6_code}), and the natural sciences already hold generative models to this standard, counting a proposed material only once it has been synthesized (Section~\ref{sec:p4_s12_2}).
When an action falls outside a model's established range, the agent must decide how to obtain further evidence, which actions are safe meanwhile, and how to use the outcome to extend the model.


\section{Situated and Interactive Spatial Intelligence}
\label{sec:p5_agency}
\label{sec:p5_competence}
\label{sec:p5_people}

\FramedInitial{S}{patial} competence is not exhausted by interpreting observations or predicting future states.
In an unfamiliar environment an agent must decide where to look, what to investigate, and when to act under uncertainty, and models that perform well on provided evidence can perform poorly when responsible for gathering it.
Theory of Space makes the comparison directly: foundation models asked to explore until they knew where objects were did worse when they collected the observations themselves than when the same observations were supplied, and explored less systematically than simple scripted strategies \cite{zhang2026theory}.
On VABench, models that judge spatial relations well still fail many manipulation tasks and almost all long-horizon ones, and the gap narrows when they may move the camera and look again \cite{zhang2026vabench}.
These results locate part of the weakness in deciding what to do, rather than only in recognizing what is seen.

Section~\ref{sec:p5_observability} asked which properties can be made identifiable by acting; the question here is which action to choose.
Information gathering should be guided not only by uncertainty but by the consequences of that uncertainty for the decision at hand.
An observation that sharpens a reconstruction may be worth less than one that settles whether a passage is traversable or a grasp is safe, and the value of a view depends on the goals, costs, and risks of the actions that follow it.
Exploration policies learned in simulation formalize this for navigation objectives (Section~\ref{subsec:p3_s7_2_5}); the counterpart for the richer states of Sections~\ref{sec:p5_sufficiency} and~\ref{sec:p5_persistence}, where the quantity to resolve may be an object's identity, a person's belief, or a door's state, is largely untrained.
Reporting accuracy as a function of an interaction budget, the number of views, moves, or queries a model may spend, would place passive and active evaluation on one scale and reward agents that reach correct decisions with few observations.

Passive data, simulation, and real-world interaction provide complementary forms of experience.
Internet video, the largest source, supports perceptual and physical priors and already enables robot planning after modest fine-tuning (Section~\ref{sec:p3_s8_6}); simulation and generated worlds permit controlled exploration without risk (Section~\ref{subsec:p3_s7_2_5}); real interaction alone exposes the consequences of the agent's own decisions, which passive video does not contain and on which both the update policy of Section~\ref{sec:p5_persistence} and the predictor of Section~\ref{sec:p5_validity} depend.
Each source carries its own distortions.
The simulation gap is widest where contact, materials, lighting, and sensor noise matter (Sections~\ref{sec:p4_s9_1} and~\ref{sec:p4_s11_3}), and experience generated inside a world model is useful only within the model's range of validity.
Learning inside a world model, with rollbacks to earlier decision points and correction of the model from real transitions \cite{luo2026rollback}, lowers the cost of real interaction without removing the need for it.
How to combine the three sources without transferring incorrect assumptions from simulated or observational distributions into real-world action remains unresolved, and existing robot datasets, concentrated on tabletop success, offer little of the failures and recoveries such a combination requires (Section~\ref{subsec:p3_s7_3_2}).
The consequences of an action, moreover, are not confined to the agent's own observations: in a shared space, acting also changes what other agents observe, believe, and do.

Different agents observe and understand the same environment differently, and a shared physical world does not imply shared knowledge of it.
A system must relate its own reference frame to those of others, since a location on the user's left may lie on the robot's right and people switch between egocentric and allocentric descriptions without noticing (Section~\ref{subsec:p3_s7_3_1}).
It must distinguish the physical state from what a person believes about it, since a user who left the kitchen before the kettle was moved still believes it is on the stove.
It must anticipate how its own actions change what others perceive and do.
Mental World Modeling maintains joint physical and mental states for this reason and finds both necessary to predict human decisions (Section~\ref{subsubsec:p3_s8_6_mwm}) \cite{fei2026mental}.
People, on this view, belong in the world state on the same footing as walls and doors, with positions, viewpoints, and beliefs treated as quantities to be estimated.

In human environments, spatial actions are inherently interactive.
Robots that model people only as moving obstacles freeze in dense crowds once every path appears blocked, unable to anticipate that people will make room (Section~\ref{sec:p4_s13_1}), and an assistant's motion, speech, and help change what the people around it do next, so a plan that ignores this is planning against a stale state.
Persistent spatial memory carries a second consequence: a world state that includes people records where they live, what they own, and how they spend their days, and the design choices of what to store, at what detail, and for how long matter more than the access controls applied afterward (Sections~\ref{sec:p4_s10_1} and~\ref{sec:p4_s13_1}).
Such systems also act on their users' own spatial abilities, which weaken under habitual reliance on navigation aids and strengthen with practice \cite{dahmani2020habitual}.
Reliable spatial agency must therefore account for other agents' responses and for the consequences of maintaining and using spatial knowledge about them.

Agents that acquire the information they need, act despite incomplete knowledge, and revise their world models from the consequences of interaction would address each of the preceding questions in operational form: action changes what is identifiable, exposes which information the representation lacked, supplies the observations that keep the state current, and tests the predictions the model made.
Spatial intelligence, on this view, is a closed loop in which perception, representation, prediction, and action continually constrain and improve one another.
A usable model of physical space is one that remains open to new evidence, explicit about its limits, and corrected by what its own actions reveal.

\chapter{Research Gaps and Future Directions}
\label{chapter:p5_challenges}


% \lettrine[lines=3]{\initfamily\textcolor{darkgreen}{F}}{uture}
% Vision-based ranger involves leveraging geometric principles to calculate distance.
% A vision-based ranger could use one or multiple cameras.
% The monocular camera ranger requires prior knowledge (\eg the size of a reference object), and can only estimating in a rough distance.

\FramedInitial{T}{he} problems in this chapter are widely acknowledged and concern missing methods, inadequate measurements, and systems that work in the laboratory but not yet in open environments.
Section~\ref{sec:p5_gaps} describes research gaps along the loop of Figure~\ref{fig:p1_s2_1_intro}, Section~\ref{sec:p5_benchmarks} gaps in how progress is measured, Section~\ref{sec:p5_deployment} bottlenecks in deployment, and Section~\ref{sec:p5_directions} directions for future work.

Chapter~\ref{chapter:p5_discussions} asked why these problems are hard in principle; this chapter asks where current systems fail, how the failures can be measured, and which methods could be tested against them.
Four themes are its own: competence at tabletop or room scale does not transfer to buildings and cities; correct spatial judgments do not reliably become completed long-horizon actions; much evaluation does not measure the abilities it is named for; and deployed systems face latency, distribution shift, security, data scarcity, and multi-agent synchronization that research systems rarely confront.
Where a failure has a principled explanation in the preceding chapter, we point to it rather than repeat it.

\section{Research Gaps along the Loop}
\label{sec:p5_gaps}

\FramedInitial{M}{ost} of the following gaps are invisible in single-step evaluations and appear only when a system must integrate many views, revise its memory, act over many steps, or operate beyond room scale.

\paragraph{Models recognize scenes without relating views to one another.}
Multimodal models can name nearly every object in a room yet cannot reliably say where the objects are; Chapter~\ref{chapter:3:sec:6} found them strong on contextual and functional relations and weak on distances, sizes, and angles (Section~\ref{sec:spatial_relation_extraction}).
Given several views of a room, they rarely notice that one was captured after an object moved \cite{cao2026samespace}, and over routes spanning several streets they answer by recognizing landmarks and matching text, with little evidence of path integration \cite{mahendran2026kilometervision}.
These failures admit more than one explanation: the model may hold no account of the scene against which a later view could be checked, or it may hold spatial information that it cannot associate across viewpoints, update, or use in a metric computation.
Current benchmarks do not separate these cases (Section~\ref{sec:p5_benchmarks}).

SLAM registers each new view against a map and treats disagreement as a signal to correct one or the other (Section~\ref{sec:mapping_slam}).
Whether registering views in a common frame before language is involved (Section~\ref{cross_modality_vision_vision}), for instance with feed-forward models that regress point maps from image collections (Section~\ref{sec:visual_odometry}), closes the gap for multimodal models is unsettled, and metric estimates must in any case survive unknown or wrong camera calibration, the main open problem even for specialist pipelines (Section~\ref{sec:height_distance}).

\paragraph{Memory breaks when it has to be revised.}
Section~\ref{sec:p5_persistence} argued that persistence requires revising beliefs, not only retaining observations.
The gap is that few learned systems implement a write policy deciding what enters the world state, with what confidence, and when an entry is revised or removed.
Work on spatial memory has concentrated on recall, through longer contexts, better retrieval, and richer memory structures (Section~\ref{subsec:p3_s7_3_6}), whereas the reported failures concern revision: outdated beliefs kept after the world changes \cite{zhang2026theory}, append-only memories that accumulate contradictions, and agents that drift as their context fills (Section~\ref{subsec:p3_s7_3_7}).
Applications make the requirement concrete: a manipulator must keep an object's pose current while its own gripper occludes it (Section~\ref{sec:p4_s9_1}), and a home assistant must track object locations over weeks as the household moves things (Section~\ref{sec:p4_s13_1}).
SLAM makes a version of these decisions when it selects keyframes, merges landmarks, and discards map regions that no longer match (Section~\ref{sec:mapping_slam}); we are not aware of a learned counterpart for the richer world states that applications now require.

\paragraph{Pixel prediction is mismatched with the needs of control.}
Video world models predict pixels, while a controller needs the quantities that decide whether an action succeeds, such as contact, mass, and friction, which are often invisible in images (Section~\ref{subsec:p3_s7_1_2}).
A pixel loss spends most of a model's capacity on lighting and texture, pixel predictions drift over long horizons (Section~\ref{subsec:p3_s8_3_6}), and a model can score well on tasks with incorrect dynamics \cite{gonzalezlaiz2026limits}.
The objective rewards appearance; the use requires outcomes.
World-action models that predict compact encodings of future observations match generative ones at much lower cost \cite{yu2026jam}, following the predictive-embedding principle of the JEPA family (Section~\ref{subsec:p3_s8_6_rep}), and a further step is to predict physical quantities directly, estimated through pushing, lifting, or probing.
Whether such predictions hold under intervention is the question of Section~\ref{sec:p5_validity}; the gap here is that these quantities are rarely prediction targets at all.

\paragraph{Correct perception does not carry through to completed actions.}
Chapter~\ref{sec:reasoning} described the gap between spatial judgment and action as instructions interpreted correctly and executed incorrectly (Section~\ref{subsec:p3_s7_3_2}).
On VABench, models that judge spatial relations well still fail many manipulation tasks and almost all long-horizon ones, and degrade markedly on unfamiliar shapes \cite{zhang2026vabench}; generalist robot policies improve rapidly, yet their spatial representations drift over long action sequences (Section~\ref{subsec:p3_s7_3_2}).
Small errors compound over many steps, and in contact-rich tasks a single misjudged grasp can end an episode, so the gap is one of execution rather than judgment and grows with the horizon.
An agent must keep observing and correcting until the task is complete, as the gains from letting a model look again suggest (Section~\ref{sec:p5_agency}); this is the third open tension of Chapter~\ref{chapter:structure}, the integration of spatial reasoning into long-horizon embodied agents.
Closing it requires training and evaluating the act--observe--correct loop as a whole, with benchmarks that reward recovery as well as first-attempt success.

\paragraph{Room-scale results say little about larger spaces.}
Most results in this handbook concern single rooms or tabletops.
At street and city scale, models revert to landmark recognition even over a single kilometer \cite{mahendran2026kilometervision}, geospatial models trained on one region or sensor degrade on another (Section~\ref{sec:p4_s11_2}), aerial and ground robots must align views from very different heights (Section~\ref{sec:p4_s9_4}), and skills learned on one robot body transfer poorly to another (Section~\ref{sec:p4_s9_1}).
Scale changes more than the size of the numbers: the dominant representation shifts from local metric geometry to routes and topology and then to georeferenced coordinates, the sensing from egocentric views to aerial imagery, most of the state from seen to remembered or inferred, and the rate of change from seconds to years.

Geographic information systems handle scale by nesting reference frames, placing an object in a room, the room in a building, and the building in a georeferenced map, each level with its own coordinates and precision.
Learned spatial models rarely represent such nesting.
We expect nested frames to be part of the solution, together with benchmarks that pose the same question at several scales (Section~\ref{sec:p5_benchmarks}), and we regard cross-scale transfer as a challenge in its own right.

Each of these gaps appears where one part of the loop depends on another's output, which is why evaluating each part on its own misses them.




\section{Benchmarking Gaps}
\label{sec:p5_benchmarks}

\FramedInitial{T}{he} preface asked how the field can measure genuine progress without mistaking benchmark saturation for it.
Chapter~\ref{chapter:p5_discussions} described the properties a usable spatial model should have; this section asks how to measure them.
Most benchmarks compared in Section~\ref{sec:p3_s7_5} pose questions about a single image or video, and few require the model to act.
This format admits shortcuts, which models have repeatedly exploited: models that nearly solved synthetic visual reasoning datasets failed on novel combinations of familiar elements and on real images (Section~\ref{subsec:p3_s7_2_1}), and much of a referring-expression score survives when the words of the expression are shuffled (Section~\ref{sec:vl_grounding}).

Measured against the loop of Chapter~\ref{chapter:definition}, current evaluation covers some components far better than others.
Perception is evaluated frame by frame, through detection, segmentation, depth, and tracking accuracy, and rarely by whether its uncertainty is calibrated or its output stable over time.
Reasoning is evaluated by answer accuracy on static questions, rarely by whether the answer rests on a correct belief about the scene.
Prediction is evaluated by visual quality, rarely by correctness under actions absent from training.
Action is evaluated by task success in simulation, rarely by recovery after failure.
The update $\Psi_{\mathcal{S}}$ is rarely evaluated on its own, and neither is the effect a system has on the people who use it.

\paragraph{Score the belief behind the answer.}
A correct answer may come from language priors, regularities in the answer format, or a correct model of the scene, and accuracy alone cannot tell them apart, nor separate the explanations left open for the first gap of Section~\ref{sec:p5_gaps}.
Benchmarks can instead read out a model's belief about the scene and score it directly, as Theory of Space does at each step of an exploration \cite{zhang2026theory}, or probe hidden states, as was done for the Othello model of Section~\ref{sec:rep_language}.
A cheaper variant asks several questions about the same scene and checks that the answers are consistent with a single layout.
Controls that remove the image or shuffle the input are inexpensive and should be standard.

\paragraph{Gathering evidence is a separate ability.}
Static benchmarks supply the evidence; deployed systems must collect it.
Deciding where to look, when to move, and when to ask is a skill in its own right, and the active--passive gap of Section~\ref{sec:p5_agency} shows that current models lack it.
We propose reporting accuracy as a function of an interaction budget, the number of views, moves, or queries a model may use, which places passive and active evaluation on one scale and rewards models that reach correct answers with little effort.
In deployment the budget has a direct cost in time and energy.

\paragraph{Spatial failures unfold over time.}
Most benchmark questions refer to a single moment, while many failures in Section~\ref{sec:p5_gaps} emerge only over time.
Streaming protocols show the model only the video seen so far, so it cannot use future frames to answer about the past \cite{li2026ovosbench}; revisit protocols move the camera away and back to check whether a place is rendered consistently or changed correctly \cite{wang2026worldprobe}; change protocols alter the scene behind the agent and test whether its belief follows.
None is yet standard, and long-horizon tasks remain rare even among embodied benchmarks (Section~\ref{sec:p3_s7_5}).

\paragraph{The same question should be asked at several scales.}
Benchmarks cluster at room scale and are often built on a few indoor scan collections, while surgery, orbital operations, and underwater robotics have benchmarks of their own that rarely share protocols with them (Sections~\ref{sec:p4_s12_1} and~\ref{sec:p4_s11_3} and Chapter~\ref{sec:p4_s14}).
Posing the same spatial question at the scale of a table, a building, and a city would show where competence stops transferring and whether a model represents nested frames; a model that answers well about a desk and poorly about a building has learned the regularities of tabletop scenes, not a general spatial ability.

\paragraph{Generated worlds should be judged by what an agent can do in them.}
Generated and reconstructed worlds are still evaluated mainly by appearance, which Chapter~\ref{sec:generation} showed to be misleading.
Video metrics do not register physically impossible motion (Section~\ref{subsec:p3_s8_3_6}), rendered-view similarity cannot distinguish a solid object from a cut-out that matches the training views (Section~\ref{subsec:p3_s8_2_1}), no standard protocol exists for the physical correctness or long-term stability of world models (Section~\ref{subsec:p3_s8_6_discussion}), and vision-language judges share the blind spots of the models they evaluate (Section~\ref{sec:p3_s8_8}).
A more informative test asks whether an agent that plans or trains in the generated world succeeds in the real one, complemented by the interventional tests of Section~\ref{sec:p5_validity}.

\paragraph{Benchmarks should borrow construct validity from psychometrics.}
Spatial benchmarks adopted tasks from psychometrics, from mental rotation to perspective taking (Chapter~\ref{chapter:definition}), but not its standards.
A psychometric test reports the distribution of human performance, establishes reliability across repeated testing, and checks construct validity through agreement with other tests of the same ability and divergence from tests of different ones; most AI benchmarks report at most a single human score.
A benchmark whose questions can be answered from text alone does not measure spatial ability, and automatically generated benchmarks, now common, need such checks all the more.




\section{Deployment Bottlenecks}
\label{sec:p5_deployment}

\FramedInitial{M}{ost} spatial systems that operate reliably in deployment do so by restricting the world they must handle, a pattern that recurs across Part~\ref{part:applications}.
Driving systems are approved for a restricted operational design domain and add sensors and fallbacks to stay within it (Section~\ref{sec:p4_s9_2}); warehouse robots work in sites with fixed item sets and layouts (Section~\ref{sec:p4_s9_5}); surgical robots are mostly teleoperated and planetary rovers act within ground-supervised workflows (Sections~\ref{sec:p4_s12_1} and~\ref{sec:p4_s11_3}); headsets track most reliably in previously mapped rooms, and the first home robots rely on remote operators for tasks they cannot yet perform (Sections~\ref{sec:p4_s10_1} and~\ref{sec:p4_s13_1}).
When the domain is widened, the world state is typically the first component to fail.

\paragraph{Guarantees end at the edge of the model.}
A constraint-based planner certifies collision avoidance only for the obstacles it knows about, and a driving system is validated only on anticipated scenarios (Sections~\ref{sec:p4_s9_1} and~\ref{sec:p4_s9_2}).
Learning moves these boundaries without notice: a model accurate on average can be confidently wrong in rare situations, its uncertainty estimates can lose validity when its use changes (Section~\ref{sec:p5_validity}), and an update that improves one behavior can silently degrade another (Section~\ref{sec:p4_s9_2}).
No practical method yet states and verifies a property of a learned spatial system over a realistic domain.
Legal responsibility is bounded the same way, since under conditional automation the manufacturer is responsible for driving only within a narrow domain (Section~\ref{sec:p4_s9_2}), and each extension reopens the question.

\paragraph{The large model cannot sit in the inner loop.}
A headset must respond to head motion within tens of milliseconds, a robot arm must update its state estimate at control rate, a drone must detect small objects in real time on onboard hardware, and a spacecraft must do similar work on radiation-hardened processors far weaker than ground computers (Sections~\ref{sec:p4_s10_1}, \ref{sec:p4_s9_1}, \ref{sec:p4_s9_4}, and~\ref{sec:p4_s11_3}).
Large foundation models are much slower, even after work that accelerates them by predicting geometry internally and compressing spatial tokens \cite{li2026geothinkerf}.
A practical architecture lets a small, fast component keep the world state current at sensor rate and lets the large model consult it when time permits, mirroring the fast--slow division that Section~\ref{subsec:p3_s7_3_5} described for tool-augmented reasoning.
Whether that state is explicit, learned, or mixed is the question of Section~\ref{sec:p5_persistence}; the deployment constraint is only that it be maintained at sensor rate and remain consultable.

\paragraph{A map starts going stale as soon as it is made.}
Furniture is moved, shop displays change, buildings are renovated, crops grow, and tissue shifts during surgery, so a deployed map must be updated continually, which Part~\ref{part:applications} found difficult in every domain.
How a state should be revised is the algorithmic question of Section~\ref{sec:p5_persistence}; the deployment problem is the cost and latency of keeping it revised and the synchronization of its copies across devices.
Digital twins call the discrepancy between model and asset the liveness gap and treat closing it at scale as open (Section~\ref{sec:p4_s11_1}); lifelong SLAM must withstand months of appearance change (Section~\ref{sec:mapping_slam}); shared augmented-reality maps are hard to keep current (Section~\ref{sec:entertainment}); surgical navigation loses its preoperative registration once the brain shifts (Section~\ref{sec:p4_s12_1}).
We suggest that deployed systems report how far their world state lags behind the world it describes, a quantity that is rarely measured.

\paragraph{Real data are scarcest where simulation is least trustworthy.}
Spacecraft perception is trained mostly on simulated images, and simulators do not reproduce the lighting of orbit, reflections from spacecraft materials, or radiation noise (Section~\ref{sec:p4_s11_3}); touch, force, and sound have the least data of all modalities and are also hard to simulate (Sections~\ref{subsec:p3_s7_3_8} and~\ref{sec:p4_s9_1}); clinical data are scarce because of privacy regulation and the cost of expert annotation (Section~\ref{sec:p4_s12_1}).
In the first two cases simulation cannot compensate, since the simulator is least accurate for exactly the effects the missing data would capture.
Simulations built from real measurements, data engines that target rare cases, and cross-institution training without sharing raw data all help, but none yet operates at the breadth that general spatial competence requires.

\paragraph{A persistent map is also an attack surface.}
Geospatial data raise issues of sovereignty and adversarial manipulation (Section~\ref{sec:p4_s11_2}), a compromised digital-twin server exposes the layout and vulnerabilities of the asset it mirrors (Section~\ref{sec:p4_s11_1}), and world-action models can be attacked through their visual inputs, with small perturbations producing large changes in behavior \cite{wu2026pasakv}.
These risks grow with the persistence of a map and the number of parties sharing it, turning the privacy concerns of Section~\ref{sec:p5_agency} into security concerns.




\section{Future Directions}
\label{sec:p5_directions}

\FramedInitial{T}{he} following directions build on the gaps above and on the questions left open in Chapter~\ref{chapter:p5_discussions}.
Each is a hypothesis rather than a recommendation: where it takes an architectural position, we name the alternative it must beat and an experiment that would decide.

\paragraph{Persistence and rendering can be separated.}
Interactive video models remember a place by regenerating it, which ties memory to the context window and to the renderer.
A factorized design would keep what must persist, such as layout, object identities, and their uncertainty, in an explicit layer, and use a learned layer to predict change and render what the agent needs to see, with $\Psi_{\mathcal{S}}$ promoting confirmed predictions into the explicit layer and revising it when observations contradict it.
The alternatives are a latent state with learned memory, as in models that factorize memory inside the generator \cite{yu2026grome}, and a growing context with retrieval.
The hypothesis is that at equal compute the factorized design stays consistent over longer horizons and more changes; it is refuted if a latent-only model matches it on revisit and change protocols while passing the faithfulness tests of Section~\ref{sec:p5_validity}.
The decisive experiment is an agent that leaves a building, returns hours later, detects what changed, and acts accordingly; for the factorized variant, the open question is the interface that decides when a prediction enters the explicit layer and when an entry is overridden.

\paragraph{The medium of spatial reasoning can be learned per query.}
Section~\ref{sec:p5_sufficiency} found that spatial answers improve when a model can work with maps, imagined views, or geometric tools.
The choice of medium can be learned per query from whether it led to a correct answer and at what cost, so that simple questions are answered in language and a map or simulation is invoked only when needed, with the intermediate steps later distilled into the model (Section~\ref{subsec:p3_s7_3_8}).
The comparison is with the same model reasoning in text alone and with a policy that always invokes the tool, on map building and perspective taking, with image-removed controls; the hypothesis fails if the learned selector does not beat the always-tool policy on cost at equal accuracy.

\paragraph{The update policy can be trained from experience.}
Passive data leave the update policy of Section~\ref{sec:p5_persistence} largely untrained (Section~\ref{sec:p5_agency}).
It can be trained directly, by rewarding observations that add information to the map and penalizing beliefs kept after the evidence has changed, with learning inside a world model lowering the cost by trying alternatives in imagination first (Section~\ref{subsec:p3_s7_3_7}); benchmarks that read out a model's belief (Section~\ref{sec:p5_benchmarks}) supply the training signal as well as the evaluation.
The target is an agent that builds as accurate a world state from its own exploration as from a curated set of views, and the scripted strategies that currently outperform foundation models \cite{zhang2026theory} are the baseline to beat.

\paragraph{Verification is feasible where part of the state is explicit.}
Neural networks are difficult to verify, whereas objects, poses, and programs can be checked against constraints.
The Algebraic World Model turns exploration into hypotheses checked by machine-verified proofs \cite{pan2026algebraic}, and uncertainty can be verified once its calibration is matched to the policy actually executed \cite{mittal2026calibrated}.
The alternative is to verify latent states indirectly, through probes and faithfulness tests \cite{gonzalezlaiz2026limits}, which is cheaper and yields statistical rather than logical guarantees.
In either case the obstacle is perception, since any guarantee about the state inherits the errors of its estimate from sensor data.
A machine-checked property, such as a bound on collision probability, holding for a learned agent over a deployment-scale domain would change the situation of Section~\ref{sec:p5_deployment}; no current system provides one.

\paragraph{A shared world state needs the machinery of version control.}
When several agents and people maintain one world state, they need a common reference frame, stable object identities, a record of changes, a procedure for resolving conflicts, and permissions on who may read and write what.
Version control solved the analogous problems for shared code, and multi-agent systems that build cognitive maps already record changes as atomic commits and cross-check one another \cite{zhang2026cocosi}.
Embodied teams add difficulty, since ground and aerial robots must first align very different viewpoints (Section~\ref{sec:p4_s9_4}) and city-scale and planetary digital twins supply the georeferenced frames into which local maps must fit (Sections~\ref{sec:p4_s11_1} and~\ref{sec:p4_s11_2}); permissions are also where the privacy requirements of Section~\ref{sec:p5_agency} can be enforced.
A suitable test is a team of heterogeneous robots and a person who maintain one consistent world state of a building over several days, each able to query how the others currently perceive it.

\paragraph{A theory of spatial representation has a concrete agenda.}
Section~\ref{sec:p5_sufficiency} proposed sufficiency relative to tasks and actions as the criterion for a representation.
Turning it into a theory requires identifiability results stating which world-state variables are recoverable from a given set of modalities and actions, extending the sensor-specific analyses of Chapter~\ref{chapter:3:sec:6}; sufficiency results stating which variables must be preserved for a task family under every action the agent may take, of which equivariance to rigid motion (Section~\ref{subsec:p3_s7_2_3}) is the one case already worked out; faithfulness tests deciding, for any architecture, whether a latent state's dynamics follow the environment's \cite{gonzalezlaiz2026limits}; and a specification of which parts of a representation must be explicit for inspection and verification.
Such results would let the world-centered property required by Chapter~\ref{chapter:definition} be tested for any system, regardless of architecture.

\chapter{Conclusion}
\label{chapter:p5_conclusion}


% \lettrine[lines=3]{\initfamily\textcolor{darkgreen}{C}}{onclusion}
% Vision-based ranger involves leveraging geometric principles to calculate distance.
% A vision-based ranger could use one or multiple cameras.
% The monocular camera ranger requires prior knowledge (\eg the size of a reference object), and can only estimating in a rough distance.

\FramedInitial{P}{oincar\'e} wrote in 1902 that to localize an object is to represent to ourselves the movements needed to reach it, and this handbook opened with that sentence.
The work surveyed here supports his view, as the systems that act most reliably are those whose representations are most closely tied to prediction and action, and their main weakness is keeping those representations consistent over time.

This handbook set out to provide a common map of a fragmented field, reading work from SLAM and computer vision to geographic information science and the cognitive science of space through the representations each of them constructs.
Part~I traced five eras of spatial representation and proposed a definition in which a system and its environment are coupled through a world-centered representation that perception writes, action reads, reasoning and prediction operate on, and the update $\Psi_{\mathcal{S}}$ keeps current.
Parts~II to~IV followed these functions to the state of the art.
In the terms of Chapter~\ref{chapter:definition}, different communities built different parts of this loop: vision and graphics the representation, robotics the operators that act on it, SLAM an update with geometric guarantees; few systems close the entire loop.

Part~V located the open problems mainly at the junctions between these parts.
The grocery-handling robot introduced at the beginning of this part fails when its perception, memory, prediction, and action cease to agree about the kitchen, even though each capability works in isolation.
Chapter~\ref{chapter:p5_discussions} asked what such a shared state must provide to be usable: what can be identified from limited evidence, what must be represented, how beliefs are kept current, when predictions can be trusted, and how knowledge is acquired through interaction in space shared with others.
Chapter~\ref{chapter:p5_challenges} showed the same state failing to persist in research systems, rarely measured by benchmarks, and kept manageable in deployment only by restricting the world, and stated directions for closing these gaps as hypotheses that can be tested.

The preface posed four questions without claiming final answers, and Part~V suggests provisional ones.
A system is spatially intelligent, beyond its performance on particular tasks, when it maintains a world state that stays consistent as it moves and acts and that can serve tasks it was not designed for.
Geometric, semantic, and commonsense structure integrate best when each part of the state is stored according to how long it must persist and whether it must be inspected, and modalities as different as pixels, point clouds, language, and audio contribute as complementary constraints on one underdetermined state, each resolving ambiguities the others leave open; registering views in a common frame before language is involved appears to be the most promising route to cross-view alignment, though whether it closes the gap remains to be shown.
Spatial knowledge is most likely to transfer between environments when it is anchored in nested frames, from objects to rooms to maps.
Planning and action emerge from spatial understanding through prediction that is tested under intervention and through a loop in which the agent keeps observing and correcting.
Genuine progress shows in how well a system gathers evidence, keeps its beliefs current, and uses them, more than in its answers to static questions.

Chapter~\ref{chapter:definition} called its definition a definition in motion, and Part~V suggests three revisions to it.
The definition refers to a single system $\mathcal{S}$, whereas deployed systems increasingly share spaces and world states with other agents and with people, so it requires a plural form.
It requires the representation to support perception, reasoning, prediction, and action, and deployment adds the requirement that planners, regulators, and the people it describes can inspect it.
It admits social settings only as a property of the environment, whereas Chapter~\ref{chapter:p5_discussions} argued that people belong in the state itself, with their positions, viewpoints, and beliefs predicted and respected alongside walls and doors.

As the preface anticipated, some of the systems discussed in this handbook will soon be superseded, and some of the positions taken here will need revision as new evidence arrives.
We hope that a clear account of the open problems makes that revision easier, and that the next account of spatial intelligence finds more of the loop closed.




\clearpage
\section*{Acknowledge}
\label{sec:acknowledge-acknowledge}

Thanks for the support from xxx.
% bib
\bibliographystyle{unsrtnat}
\bibliography{main}  %%% Uncomment this line and comment out the 

%\printbibliography

\clearpage
\section*{Author Contributions}
\label{sec:author-contribution-list-author-contributions}
\markboth{Author Contributions}{Author Contributions}

Contributions are listed in the order of the author list on the title page. For each author who drafted text, the parts of the handbook they wrote are given in book order, and a co-authored part names its co-author. Authors who reviewed the manuscript and advised on its direction are marked as such.

% Order: the author list on the title page (main.tex). Attribution of written parts follows the former
% \RespondAuthor{} markers (see SECTION_OWNERS.md); parts without a marker are attributed to Shengqiong Wu.
\begin{itemize}[leftmargin=*]
    \item \textbf{Shengqiong Wu}, University of Oxford:
        \begin{itemize}
            \item Part~I, Introduction and Background, Chapters~\ref{chapter:intro}--\ref{chapter:structure} (with Hao Fei)
            \item Audio, Section~\ref{sec:rep_audio}
            \item Object Tracking, Section~\ref{sec:object_tracking}
            \item Physical Reasoning, Section~\ref{subsec:p3_s7_1_2}
            \item Differentiable Optimization Layers for Spatial Reasoning, Section~\ref{subsec:p3_s7_2_4}
            \item Reinforcement Learning with Spatial Inductive Biases, Section~\ref{subsec:p3_s7_2_5}
            \item Hybrid Spatial Reasoning, Section~\ref{sec:p3_s7_4}
            \item Spatial Reasoning Benchmarks, Section~\ref{sec:p3_s7_5}
            \item 3D Object Synthesis, Section~\ref{subsec:p3_s8_2_1}
            \item 3D Scene Reconstruction, Section~\ref{subsec:p3_s8_2_3}
            \item 3D Scene Editing, Section~\ref{subsec:p3_s8_2_5}
            \item Summary, Challenges, and Future Directions of Video Spatial Generation, Section~\ref{subsec:p3_s8_3_6}
            \item Introduction to Motion Generation, Section~\ref{sec:p3_s8_4}
            \item Programmatic and Symbolic Spatial Generation, Section~\ref{sec:p3_s8_7}
            \item Summary and Outlook of Spatial Generation, Section~\ref{sec:p3_s8_8}
            \item Aerospace Engineering, Section~\ref{sec:p4_s11_3}
            \item Spatial Intelligence for Others, Chapter~\ref{sec:p4_s14}
            \item Part~V, Discussions and Futures, Chapters~\ref{chapter:p5_discussions}--\ref{chapter:p5_conclusion}
            \item Introductions to Parts~II--IV, to Chapters~\ref{chapter:single_modality} and~\ref{chapter:3:sec:6}--\ref{sec:p4_s14}, and to the sections of Chapters~\ref{chapter:3:sec:6}--\ref{sec:generation}
        \end{itemize}
    \item \textbf{Yu Zhao}, National University of Singapore:
        \begin{itemize}
            \item Height and Distance Estimation, Section~\ref{sec:height_distance}
            \item Layout Estimation and Room Parsing, Section~\ref{sec:layout_estimation}
            \item Mapping and SLAM, Section~\ref{sec:mapping_slam}
            \item Surface Normal Estimation, Section~\ref{sec:surface_normal}
            \item Visual Odometry and Camera Pose Estimation, Section~\ref{sec:visual_odometry}
            \item Scene Flow Estimation, Section~\ref{sec:scene_flow}
            \item Motion Generation, Sections~\ref{subsec:p3_s8_4_1}--\ref{subsec:p3_s8_4_4}
            \item Autonomous Driving, Section~\ref{sec:p4_s9_2}
            \item GeoAI, Section~\ref{sec:p4_s11_2}
        \end{itemize}
    \item \textbf{Tianjie Ju}, National University of Singapore:
        \begin{itemize}
            \item Foundations: What Spatial Intelligence Requires, Chapter~\ref{chapter:p5_discussions}
            \item Research Gaps and Future Directions, Chapter~\ref{chapter:p5_challenges}
            \item Content Screening
        \end{itemize}
    \item \textbf{Jundong Xu}, National University of Singapore:
        \begin{itemize}
            \item Geometric and Algebraic Reasoning, Section~\ref{subsec:p3_s7_1_1}
            \item Topological and Graph-based Reasoning, Section~\ref{subsec:p3_s7_1_3}
            \item Constraint-based Spatial Solving, Section~\ref{subsec:p3_s7_1_4}
            \item Rule-based and Programmatic Reasoning, Section~\ref{subsec:p3_s7_1_5}
            \item Causal Reasoning, Section~\ref{subsec:p3_s7_1_6}
            \item Warehouse and Industrial Logistics, Section~\ref{sec:p4_s9_5}
        \end{itemize}
    \item \textbf{Bobo Li}, National University of Singapore:
        \begin{itemize}
            \item Temporal Localization, Section~\ref{sec:temporal_localization}
            \item Multi-agent Spatial Reasoning, Section~\ref{subsec:p3_s7_2_6}
            \item Text-to-image Generation, Section~\ref{subsec:t2i}
            \item Entertainment, Section~\ref{sec:entertainment}
            \item AR/VR/XR, Section~\ref{sec:p4_s10_1} (with Xin Kong)
        \end{itemize}
    \item \textbf{Jingkang Yang and Shulin Tian}, Nanyang Technological University, Singapore:
        \begin{itemize}
            \item Chain-of-Thought for Spatial Reasoning, Section~\ref{subsec:p3_s7_3_3}
            \item Tool-augmented Reasoning, Section~\ref{subsec:p3_s7_3_5}
            \item Retrieval- \& Memory-augmented Spatial Reasoning, Section~\ref{subsec:p3_s7_3_6}
            \item Agentic Spatial Reasoning, Section~\ref{subsec:p3_s7_3_7}
            \item Assistant, Section~\ref{sec:p4_s13_1}
        \end{itemize}
    \item \textbf{Daoan Zhang}, University of Rochester:
        \begin{itemize}
            \item Audio-Visual Spatial Localization, Section~\ref{sec:av_localization}
            \item 2D Spatial-aware Image Generation, Section~\ref{subsec:p3_s8_1_2}
            \item 2D Spatial Inpainting \& Editing, Section~\ref{subsec:p3_s8_1_3}
            \item Medicine and Clinical Healthcare, Section~\ref{sec:p4_s12_1}
        \end{itemize}
    \item \textbf{Yanguang Zhao}, National University of Singapore:
        \begin{itemize}
            \item Keypoint and Landmark Detection, Section~\ref{sec:keypoint_landmark}
            \item Depth Estimation, Section~\ref{sec:depth_estimation}
            \item Natural Science, Section~\ref{sec:Natural Science}
        \end{itemize}
    \item \textbf{Xuan Ju}, The Chinese University of Hong Kong:
        \begin{itemize}
            \item Text-to-Video Generation, Section~\ref{subsec:p3_s8_3_1}
            \item Spatial Controlled Video Generation, Section~\ref{subsec:p3_s8_3_2}
            \item Temporal Controlled Video Generation, Section~\ref{subsec:p3_s8_3_3}
            \item Other Controlled Video Generation, Section~\ref{subsec:p3_s8_3_4}
            \item Video Editing, Section~\ref{subsec:p3_s8_3_5}
        \end{itemize}
    \item \textbf{Haoyu Zhao}, The Chinese University of Hong Kong:
        \begin{itemize}
            \item 3D Representations, Section~\ref{sec:rep_3d}
            \item 4D Spatial Generation, Section~\ref{sec:p3_s8_5}
        \end{itemize}
    \item \textbf{Yuchen Zhou}, National University of Singapore:
        \begin{itemize}
            \item Spatial Relation Extraction, Section~\ref{sec:spatial_relation_extraction}
            \item Scene Graph Generation, Section~\ref{sec:scene_graph_generation}
            \item Action Detection, Section~\ref{sec:action_detection}
            \item LLM/VLM-centric Spatial Reasoning, Section~\ref{subsec:p3_s7_3_1}
            \item VLA-based Spatial Reasoning, Section~\ref{subsec:p3_s7_3_2}
        \end{itemize}
    \item \textbf{Xiangtai Li}, Nanyang Technological University, Singapore:
        \begin{itemize}
            \item 2D Representations, Section~\ref{sec:rep_2d}
            \item Object Detection, Section~\ref{sec:6-1-2}
            \item Segmentation, Section~\ref{sec:segmentation}
        \end{itemize}
    \item \textbf{Zhedong Zheng}, University of Macau:
        \begin{itemize}
            \item Aerial Agents, Section~\ref{sec:p4_s9_4}
        \end{itemize}
    \item \textbf{Joy Hsu}, Stanford University:
        \begin{itemize}
            \item Domain-Specific Representations, Section~\ref{sec:rep_other}
            \item Cross-Modal Spatial Representation, Chapter~\ref{chapter:cross_modality}
        \end{itemize}
    \item \textbf{Xin Kong}, Imperial College London:
        \begin{itemize}
            \item 4D Representations, Section~\ref{sec:rep_4d} (with Chuanxia Zheng)
            \item 3D Scene Synthesis, Section~\ref{subsec:p3_s8_2_2}
            \item AR/VR/XR, Section~\ref{sec:p4_s10_1} (with Bettina Schlager)
        \end{itemize}
    \item \textbf{Ruining Li}, University of Oxford:
        \begin{itemize}
            \item World-Model-based Spatial Reasoning, Section~\ref{subsec:p3_s7_3_8} (with Hao Fei)
            \item Digital Twins, Section~\ref{sec:p4_s11_1}
        \end{itemize}
    \item \textbf{Cheng Zhang}, Monash University:
        \begin{itemize}
            \item 3D Panorama Synthesis, Section~\ref{subsec:p3_s8_2_4}
        \end{itemize}
    \item \textbf{Chengzu Li}, University of Cambridge:
        \begin{itemize}
            \item Scene Classification, Section~\ref{sec:scene_classification} (with Hao Fei)
            \item Implicit Neural Spatial Reasoning, Section~\ref{subsec:p3_s7_2_1} (with Hao Fei)
            \item Explicit Structured Neural Reasoning, Section~\ref{subsec:p3_s7_2_2} (with Hao Fei)
            \item Geometry-aware Neural Spatial Reasoning, Section~\ref{subsec:p3_s7_2_3} (with Hao Fei)
            \item Code/DSL Generation for Spatial Reasoning, Section~\ref{subsec:p3_s7_3_4} (with Hao Fei)
        \end{itemize}
    \item \textbf{Jiafei Duan}, National University of Singapore:
        \begin{itemize}
            \item Robotic Manipulation, Section~\ref{sec:p4_s9_1}
        \end{itemize}
    \item \textbf{Yu (Bryan) Zhou}, University of California, Los Angeles:
        \begin{itemize}
            \item Vision-Language Grounding, Section~\ref{sec:vl_grounding}
        \end{itemize}
    \item \textbf{Kenneth Li}, Harvard University:
        \begin{itemize}
            \item Language Representations, Section~\ref{sec:rep_language}
        \end{itemize}
    \item \textbf{Chuanxia Zheng}, Nanyang Technological University, Singapore:
        \begin{itemize}
            \item 4D Representations, Section~\ref{sec:rep_4d} (with Xin Kong)
        \end{itemize}
    \item \textbf{Lei Cui}, Microsoft Research: Reviewing \& Advisory, providing comments and feedback that shaped the direction of the handbook
    \item \textbf{Furu Wei}, Microsoft Research: Reviewing \& Advisory, providing comments and feedback that shaped the direction of the handbook
    \item \textbf{Ziwei Liu}, Nanyang Technological University, Singapore: Reviewing \& Advisory, providing comments and feedback that shaped the direction of the handbook
    \item \textbf{Steven K. Feiner}, Columbia University: Reviewing \& Advisory, providing comments and feedback that shaped the direction of the handbook
    \item \textbf{Jiebo Luo}, University of Rochester: Reviewing \& Advisory, providing comments and feedback that shaped the direction of the handbook
    \item \textbf{Jianfei Cai}, Monash University: Reviewing \& Advisory, providing comments and feedback that shaped the direction of the handbook
    \item \textbf{Shuicheng Yan}, National University of Singapore: Reviewing \& Advisory, providing comments and feedback that shaped the direction of the handbook
    \item \textbf{Ming-Hsuan Yang}, University of California, Merced: Reviewing \& Advisory, providing comments and feedback that shaped the direction of the handbook
    \item \textbf{Hao Fei}, University of Oxford:
        \begin{itemize}
            \item Part~I, Introduction and Background, Chapters~\ref{chapter:intro}--\ref{chapter:structure} (with Shengqiong Wu)
            \item Scene Classification, Section~\ref{sec:scene_classification} (with Chengzu Li)
            \item Pose Estimation, Section~\ref{sec:object_pose}
            \item Implicit Neural Spatial Reasoning, Section~\ref{subsec:p3_s7_2_1} (with Chengzu Li)
            \item Explicit Structured Neural Reasoning, Section~\ref{subsec:p3_s7_2_2} (with Chengzu Li)
            \item Geometry-aware Neural Spatial Reasoning, Section~\ref{subsec:p3_s7_2_3} (with Chengzu Li)
            \item Code/DSL Generation for Spatial Reasoning, Section~\ref{subsec:p3_s7_3_4} (with Chengzu Li)
            \item World-Model-based Spatial Reasoning, Section~\ref{subsec:p3_s7_3_8} (with Ruining Li)
            \item World Modeling, Section~\ref{sec:worldmodeling}
            \item Wearable AI System, Section~\ref{sec:p4_s9_3}
            \item Figures: drew part of the handbook's illustrations, most notably the teaser figure
        \end{itemize}
\end{itemize}


\end{document}
